\chapter{Compiler error codes}%
\label{chap:codes}

Every diagnostic the compiler raises carries a code, printed next to its
severity:

\begin{lstlisting}[style=bashVerb]
Error[E4020] : class A has no method named foo
 --> main.yr:(8,15)
\end{lstlisting}

A code is allocated once, and never reused nor renumbered: rewording a message
keeps its code, and a message deleted from the compiler retires its code
(Section~\ref{sec:codes:retired}). Its first digit names the stage of the
compiler that raises it: 1 for lexing, 2 for parsing, 3 for declaration, 4 for
validation, and 5 for lowering to the intermediate language.

This appendix gives each code a page, grouped by theme. A page quotes the
message, where \errhole{} stands for a value the diagnostic fills in, says what
the message means, shows a program raising it together with the compiler's
output, and says how to fix it. Messages that only accompany another
diagnostic, such as ``when validating \ldots'' or ``candidate \ldots'', carry no
code and have no page.

\minitoc%

%% Generated by tools/gen_error_codes.py from docs/errors/ of the compiler;
%% edit the compiler's pages or tools/error_themes.txt, then rerun it.
%% Generated by tools/gen_error_codes.py from docs/errors/ of the compiler;
%% edit the compiler's pages or tools/error_themes.txt, then rerun it.

\clearpage
\section{Lexical structure}
\label{sec:codes:lexical}

\begin{longtable}{@{}l p{0.78\linewidth} r@{}}
\toprule Code & Message & Page\\ \midrule \endhead
\hyperref[sec:codes:e1001]{E1001} & unterminated string literal \errhole{} & \pageref{sec:codes:e1001}\\
\hyperref[sec:codes:e1002]{E1002} & unterminated comment block \errhole{} & \pageref{sec:codes:e1002}\\
\hyperref[sec:codes:e1004]{E1004} & line break encountered inside a single-line string literal & \pageref{sec:codes:e1004}\\
\hyperref[sec:codes:e1005]{E1005} & escaping a line break is prohibited; use a multiline string instead """ & \pageref{sec:codes:e1005}\\
\hyperref[sec:codes:e1006]{E1006} & escaping a line break is prohibited; use a multiline comment instead /**/ & \pageref{sec:codes:e1006}\\
\hyperref[sec:codes:e1007]{E1007} & escape character found at end of string literal & \pageref{sec:codes:e1007}\\
\hyperref[sec:codes:e1008]{E1008} & escape character found at end of comment block & \pageref{sec:codes:e1008}\\
\hyperref[sec:codes:e1009]{E1009} & invalid numeric value formatted in hexadecimal & \pageref{sec:codes:e1009}\\
\hyperref[sec:codes:e1010]{E1010} & invalid numeric value formatted in decimal & \pageref{sec:codes:e1010}\\
\hyperref[sec:codes:e1011]{E1011} & invalid numeric value formatted in binary & \pageref{sec:codes:e1011}\\
\hyperref[sec:codes:e1012]{E1012} & invalid numeric value formatted in octal & \pageref{sec:codes:e1012}\\
\hyperref[sec:codes:e1013]{E1013} & expected sign information in hexadecimal floating-point literal & \pageref{sec:codes:e1013}\\
\hyperref[sec:codes:e1014]{E1014} & expected sign information in scientific floating-point literal & \pageref{sec:codes:e1014}\\
\hyperref[sec:codes:e1015]{E1015} & \errhole{} is not a valid identifier; use backquotes for special identifiers(e.g. `\errhole{}`) & \pageref{sec:codes:e1015}\\
\hyperref[sec:codes:e1016]{E1016} & binary number is empty & \pageref{sec:codes:e1016}\\
\hyperref[sec:codes:e1017]{E1017} & octal number is empty & \pageref{sec:codes:e1017}\\
\hyperref[sec:codes:e1018]{E1018} & hexadecimal number is empty & \pageref{sec:codes:e1018}\\
\hyperref[sec:codes:e1019]{E1019} & the exponent part of the floating-point number is empty & \pageref{sec:codes:e1019}\\
\hyperref[sec:codes:e1020]{E1020} & the fractional part of the floating-point number is empty & \pageref{sec:codes:e1020}\\
\hyperref[sec:codes:e1021]{E1021} & missing the exponent part in a floating-point number formatted in hexadecimal & \pageref{sec:codes:e1021}\\
\hyperref[sec:codes:e1022]{E1022} & unterminated version block: missing \#end & \pageref{sec:codes:e1022}\\
\hyperref[sec:codes:e1023]{E1023} & \#\errhole{} without matching \#if & \pageref{sec:codes:e1023}\\
\hyperref[sec:codes:e1024]{E1024} & missing version condition: expected identifier or backquoted identifier & \pageref{sec:codes:e1024}\\
\hyperref[sec:codes:e1025]{E1025} & undefined version directive; expecting if, else or end & \pageref{sec:codes:e1025}\\
\bottomrule
\end{longtable}

\errorcode{E1001}{unterminated string literal \errhole{}}
\label{sec:codes:e1001}

Raised during \textbf{lexing}, from \texttt{LexingErrorMessage::\allowbreak{}UNTERMINATED\_\allowbreak{}STRING} (\textit{src/\allowbreak{}ymirc/\allowbreak{}lexing/\allowbreak{}errors.yr}).

\subsubsection*{What it means}

The lexer reached the end of the file while still inside a string literal: the opening quote it reports has no matching closing quote. The quote it was looking for is shown in the message, so a \tokennolst{"} and a \tokennolst{"""} are told apart.

\subsubsection*{Example}

\textit{test\_\allowbreak{}resources/\allowbreak{}lexing/\allowbreak{}test12.yr}:

%% check: skip
\begin{lstlisting}[style=coloredverbatimError]
let x = """abc
\end{lstlisting}

\begin{lstlisting}[style=bashVerb, breaklines=true, escapechar=~]
Error[E1001] : unterminated string literal """
 --> main.yr:(1,9)
 1  ~\lstbox{\textSFxi{}}~ let x = """abc
    ~\lstbox{\textSFv{}}~         ^^^
    ~\lstbox{\textSFii{}}~~\lstbox{\textSFx{}}~~\lstbox{\textSFx{}}~~\lstbox{\textSFx{}}~~\lstbox{\textSFx{}}~~\lstbox{\textSFx{}}~~\lstbox{\textSFx{}}~~\lstbox{\textSFx{}}~
\end{lstlisting}

\subsubsection*{How to fix it}

Close the literal with the same delimiter that opened it. A literal that is meant to span several lines needs the \tokennolst{"""} form, since a plain \tokennolst{"} literal cannot cross a line break (see \hyperref[sec:codes:e1004]{E1004}).

\errorcode{E1002}{unterminated comment block \errhole{}}
\label{sec:codes:e1002}

Raised during \textbf{lexing}, from \texttt{LexingErrorMessage::\allowbreak{}UNTERMINATED\_\allowbreak{}COMMENT} (\textit{src/\allowbreak{}ymirc/\allowbreak{}lexing/\allowbreak{}errors.yr}).

\subsubsection*{What it means}

The lexer reached the end of the file while still inside a comment block. The opening marker it reports was never matched by its closing one.

\subsubsection*{Example}

\textit{test\_\allowbreak{}resources/\allowbreak{}lexing/\allowbreak{}test13.yr}:

%% check: skip
\begin{lstlisting}[style=coloredverbatimError]
fn main () {
    /* unterminated
\end{lstlisting}

\begin{lstlisting}[style=bashVerb, breaklines=true, escapechar=~]
Error[E1002] : unterminated comment block */
 --> main.yr:(2,5)
 2  ~\lstbox{\textSFxi{}}~     /* unterminated
    ~\lstbox{\textSFv{}}~     ^^
    ~\lstbox{\textSFii{}}~~\lstbox{\textSFx{}}~~\lstbox{\textSFx{}}~~\lstbox{\textSFx{}}~~\lstbox{\textSFx{}}~~\lstbox{\textSFx{}}~~\lstbox{\textSFx{}}~~\lstbox{\textSFx{}}~
\end{lstlisting}

\subsubsection*{How to fix it}

Close the comment. \tokennolst{/*} is closed by \tokennolst{*/} and \tokennolst{/**} by \tokennolst{*/} as well; a comment opened inside another one has to be closed too, since block comments nest.

\errorcode{E1004}{line break encountered inside a single-line string literal}
\label{sec:codes:e1004}

Raised during \textbf{lexing}, from \texttt{LexingErrorMessage::\allowbreak{}RETURN\_\allowbreak{}IN\_\allowbreak{}SINGLE\_\allowbreak{}LINE\_\allowbreak{}STRING} (\textit{src/\allowbreak{}ymirc/\allowbreak{}lexing/\allowbreak{}errors.yr}).

\subsubsection*{What it means}

A \tokennolst{"} string literal contains a raw line break. A single-line literal ends at the end of its line, so the closing quote can never be reached.

\subsubsection*{Example}

\textit{test\_\allowbreak{}resources/\allowbreak{}lexing/\allowbreak{}test14.yr}:

%% check: skip
\begin{lstlisting}[style=coloredverbatimError]
let x = "abc
let y = 1;
\end{lstlisting}

\begin{lstlisting}[style=bashVerb, breaklines=true, escapechar=~]
Error[E1004] : line break encountered inside a single-line string literal
 --> main.yr:(1,9)
 1  ~\lstbox{\textSFxi{}}~ let x = "abc

    ~\lstbox{\textSFv{}}~         ^   ^
    ~\lstbox{\textSFii{}}~~\lstbox{\textSFx{}}~~\lstbox{\textSFx{}}~~\lstbox{\textSFx{}}~~\lstbox{\textSFx{}}~~\lstbox{\textSFx{}}~~\lstbox{\textSFx{}}~~\lstbox{\textSFx{}}~
\end{lstlisting}

\subsubsection*{How to fix it}

Either close the literal before the end of the line and concatenate the pieces, or use the multiline form \tokennolst{""" ... """}, which accepts line breaks. \tokennolst{\textbackslash{}n} inside a single-line literal also works when a line break is all that is wanted.

\errorcode{E1005}{escaping a line break is prohibited; use a multiline string instead """}
\label{sec:codes:e1005}

Raised during \textbf{lexing}, from \texttt{LexingErrorMessage::\allowbreak{}ESCAPE\_\allowbreak{}RETURN\_\allowbreak{}PROHIBITED\_\allowbreak{}STRING} (\textit{src/\allowbreak{}ymirc/\allowbreak{}lexing/\allowbreak{}errors.yr}).

\subsubsection*{What it means}

A backslash was used to continue a \tokennolst{"} string literal onto the next line. Ymir does not accept an escaped line break inside a string: a literal either fits on one line or uses the multiline form.

\subsubsection*{Example}

\textit{test\_\allowbreak{}resources/\allowbreak{}lexing/\allowbreak{}test15.yr}:

%% check: skip
\begin{lstlisting}[style=coloredverbatimError]
let x = "abc\
def";
\end{lstlisting}

\begin{lstlisting}[style=bashVerb, breaklines=true, escapechar=~]
Error[E1005] : escaping a line break is prohibited; use a multiline string instead """
 --> main.yr:(1,9)
 1  ~\lstbox{\textSFxi{}}~ let x = "abc\
    ~\lstbox{\textSFv{}}~         ^   ^
    ~\lstbox{\textSFii{}}~~\lstbox{\textSFx{}}~~\lstbox{\textSFx{}}~~\lstbox{\textSFx{}}~~\lstbox{\textSFx{}}~~\lstbox{\textSFx{}}~~\lstbox{\textSFx{}}~~\lstbox{\textSFx{}}~
\end{lstlisting}

\subsubsection*{How to fix it}

Write the literal as \tokennolst{""" ... """}, which spans lines without any escaping, or put \tokennolst{\textbackslash{}n} in a single-line literal where the break is wanted.

\errorcode{E1006}{escaping a line break is prohibited; use a multiline comment instead /**/}
\label{sec:codes:e1006}

Raised during \textbf{lexing}, from \texttt{LexingErrorMessage::\allowbreak{}ESCAPE\_\allowbreak{}RETURN\_\allowbreak{}PROHIBITED\_\allowbreak{}COMMENT} (\textit{src/\allowbreak{}ymirc/\allowbreak{}lexing/\allowbreak{}errors.yr}).

\subsubsection*{What it means}

A backslash was used to continue a single-line comment onto the next line. Escaping a line break has no meaning in a comment.

\subsubsection*{Example}

\textit{test\_\allowbreak{}resources/\allowbreak{}lexing/\allowbreak{}test16.yr}:

%% check: skip
\begin{lstlisting}[style=coloredverbatimError]
// abc\
def
\end{lstlisting}

\begin{lstlisting}[style=bashVerb, breaklines=true, escapechar=~]
Error[E1006] : escaping a line break is prohibited; use a multiline comment instead /**/
 --> main.yr:(1,1)
 1  ~\lstbox{\textSFxi{}}~ // abc\
    ~\lstbox{\textSFv{}}~ ^^    ^
    ~\lstbox{\textSFii{}}~~\lstbox{\textSFx{}}~~\lstbox{\textSFx{}}~~\lstbox{\textSFx{}}~~\lstbox{\textSFx{}}~~\lstbox{\textSFx{}}~~\lstbox{\textSFx{}}~~\lstbox{\textSFx{}}~
\end{lstlisting}

\subsubsection*{How to fix it}

Use a block comment (\tokennolst{/* ... */}) for a comment that spans several lines, or repeat \tokennolst{//} on each line.

\errorcode{E1007}{escape character found at end of string literal}
\label{sec:codes:e1007}

Raised during \textbf{lexing}, from \texttt{LexingErrorMessage::\allowbreak{}ESCAPE\_\allowbreak{}AT\_\allowbreak{}END\_\allowbreak{}OF\_\allowbreak{}STRING} (\textit{src/\allowbreak{}ymirc/\allowbreak{}lexing/\allowbreak{}errors.yr}).

\subsubsection*{What it means}

The string literal ends on a backslash: the escape sequence it starts has nothing to escape, because the closing quote is the next character.

\subsubsection*{Example}

\textit{test\_\allowbreak{}resources/\allowbreak{}lexing/\allowbreak{}test17.yr}:

%% check: skip
\begin{lstlisting}[style=coloredverbatimError]
let x = "abc\
\end{lstlisting}

\begin{lstlisting}[style=bashVerb, breaklines=true, escapechar=~]
Error[E1007] : escape character found at end of string literal
 --> main.yr:(1,9)
 1  ~\lstbox{\textSFxi{}}~ let x = "abc\
    ~\lstbox{\textSFv{}}~         ^   ^
    ~\lstbox{\textSFii{}}~~\lstbox{\textSFx{}}~~\lstbox{\textSFx{}}~~\lstbox{\textSFx{}}~~\lstbox{\textSFx{}}~~\lstbox{\textSFx{}}~~\lstbox{\textSFx{}}~~\lstbox{\textSFx{}}~
\end{lstlisting}

\subsubsection*{How to fix it}

Write \tokennolst{\textbackslash{}\textbackslash{}} if a literal backslash is wanted, or complete the escape sequence with the character it applies to.

\errorcode{E1008}{escape character found at end of comment block}
\label{sec:codes:e1008}

Raised during \textbf{lexing}, from \texttt{LexingErrorMessage::\allowbreak{}ESCAPE\_\allowbreak{}AT\_\allowbreak{}END\_\allowbreak{}OF\_\allowbreak{}COMMENT} (\textit{src/\allowbreak{}ymirc/\allowbreak{}lexing/\allowbreak{}errors.yr}).

\subsubsection*{What it means}

The comment block ends on a backslash, so the escape sequence it starts has nothing to escape.

\subsubsection*{Example}

\textit{test\_\allowbreak{}resources/\allowbreak{}lexing/\allowbreak{}test18.yr}:

%% check: skip
\begin{lstlisting}[style=coloredverbatimError]
/* abc\
\end{lstlisting}

\begin{lstlisting}[style=bashVerb, breaklines=true, escapechar=~]
Error[E1008] : escape character found at end of comment block
 --> main.yr:(1,1)
 1  ~\lstbox{\textSFxi{}}~ /* abc\
    ~\lstbox{\textSFv{}}~ ^^    ^
    ~\lstbox{\textSFii{}}~~\lstbox{\textSFx{}}~~\lstbox{\textSFx{}}~~\lstbox{\textSFx{}}~~\lstbox{\textSFx{}}~~\lstbox{\textSFx{}}~~\lstbox{\textSFx{}}~~\lstbox{\textSFx{}}~
\end{lstlisting}

\subsubsection*{How to fix it}

Remove the trailing backslash, or double it if a literal backslash is wanted.

\errorcode{E1009}{invalid numeric value formatted in hexadecimal}
\label{sec:codes:e1009}

Raised during \textbf{lexing}, from \texttt{LexingErrorMessage::\allowbreak{}INVALID\_\allowbreak{}HEX\_\allowbreak{}NUMBER} (\textit{src/\allowbreak{}ymirc/\allowbreak{}lexing/\allowbreak{}errors.yr}).

\subsubsection*{What it means}

A literal introduced by \tokennolst{0x} contains a character that is not a hexadecimal digit. Hexadecimal digits are \tokennolst{0} to \tokennolst{9}, \tokennolst{a} to \tokennolst{f} and \tokennolst{A} to \tokennolst{F}, plus \tokennolst{\_} as a separator.

\subsubsection*{Example}

\textit{test\_\allowbreak{}resources/\allowbreak{}lexing/\allowbreak{}test23.yr}:

%% check: skip
\begin{lstlisting}[style=coloredverbatimError]
let x = 0xZ1;
\end{lstlisting}

\begin{lstlisting}[style=bashVerb, breaklines=true, escapechar=~]
Error[E1009] : invalid numeric value formatted in hexadecimal
 --> main.yr:(1,9)
 1  ~\lstbox{\textSFxi{}}~ let x = 0xZ1;
    ~\lstbox{\textSFv{}}~         ^^^^
    ~\lstbox{\textSFxi{}}~ Note :
    ~\lstbox{\textSFxi{}}~  --> main.yr:(1,11)
    ~\lstbox{\textSFxi{}}~  1  ~\lstbox{\textSFxi{}}~ let x = 0xZ1;
    ~\lstbox{\textSFxi{}}~     ~\lstbox{\textSFv{}}~           ^
    ~\lstbox{\textSFxi{}}~     ~\lstbox{\textSFii{}}~~\lstbox{\textSFx{}}~~\lstbox{\textSFx{}}~~\lstbox{\textSFx{}}~~\lstbox{\textSFx{}}~~\lstbox{\textSFx{}}~~\lstbox{\textSFx{}}~~\lstbox{\textSFx{}}~
    ~\lstbox{\textSFii{}}~~\lstbox{\textSFx{}}~~\lstbox{\textSFx{}}~~\lstbox{\textSFx{}}~~\lstbox{\textSFx{}}~~\lstbox{\textSFx{}}~~\lstbox{\textSFvii{}}~~\lstbox{\textSFx{}}~
\end{lstlisting}

\subsubsection*{How to fix it}

Remove the offending character. If it was meant to be a type suffix, check its spelling: a suffix the lexer does not know is read as part of the number.

\errorcode{E1010}{invalid numeric value formatted in decimal}
\label{sec:codes:e1010}

Raised during \textbf{lexing}, from \texttt{LexingErrorMessage::\allowbreak{}INVALID\_\allowbreak{}DEC\_\allowbreak{}NUMBER} (\textit{src/\allowbreak{}ymirc/\allowbreak{}lexing/\allowbreak{}errors.yr}).

\subsubsection*{What it means}

A decimal literal contains a character that is neither a digit nor a valid suffix.

\subsubsection*{Example}

\textit{test\_\allowbreak{}resources/\allowbreak{}lexing/\allowbreak{}test25.yr}:

%% check: skip
\begin{lstlisting}[style=coloredverbatimError]
let x = 12a4;
\end{lstlisting}

\begin{lstlisting}[style=bashVerb, breaklines=true, escapechar=~]
Error[E1010] : invalid numeric value formatted in decimal
 --> main.yr:(1,9)
 1  ~\lstbox{\textSFxi{}}~ let x = 12a4;
    ~\lstbox{\textSFv{}}~         ^^^^
    ~\lstbox{\textSFxi{}}~ Note :
    ~\lstbox{\textSFxi{}}~  --> main.yr:(1,11)
    ~\lstbox{\textSFxi{}}~  1  ~\lstbox{\textSFxi{}}~ let x = 12a4;
    ~\lstbox{\textSFxi{}}~     ~\lstbox{\textSFv{}}~           ^
    ~\lstbox{\textSFxi{}}~     ~\lstbox{\textSFii{}}~~\lstbox{\textSFx{}}~~\lstbox{\textSFx{}}~~\lstbox{\textSFx{}}~~\lstbox{\textSFx{}}~~\lstbox{\textSFx{}}~~\lstbox{\textSFx{}}~~\lstbox{\textSFx{}}~
    ~\lstbox{\textSFii{}}~~\lstbox{\textSFx{}}~~\lstbox{\textSFx{}}~~\lstbox{\textSFx{}}~~\lstbox{\textSFx{}}~~\lstbox{\textSFx{}}~~\lstbox{\textSFvii{}}~~\lstbox{\textSFx{}}~
\end{lstlisting}

\subsubsection*{How to fix it}

Remove the offending character, or separate the number from what follows it. \tokennolst{\_} is accepted inside a number as a digit separator (\tokennolst{1\_000\_000}).

\errorcode{E1011}{invalid numeric value formatted in binary}
\label{sec:codes:e1011}

Raised during \textbf{lexing}, from \texttt{LexingErrorMessage::\allowbreak{}INVALID\_\allowbreak{}BIN\_\allowbreak{}NUMBER} (\textit{src/\allowbreak{}ymirc/\allowbreak{}lexing/\allowbreak{}errors.yr}).

\subsubsection*{What it means}

A literal introduced by \tokennolst{0b} contains a character other than \tokennolst{0} or \tokennolst{1}.

\subsubsection*{Example}

\textit{test\_\allowbreak{}resources/\allowbreak{}lexing/\allowbreak{}test19.yr}:

%% check: skip
\begin{lstlisting}[style=coloredverbatimError]
let x = 0b1021;
\end{lstlisting}

\begin{lstlisting}[style=bashVerb, breaklines=true, escapechar=~]
Error[E1011] : invalid numeric value formatted in binary
 --> main.yr:(1,9)
 1  ~\lstbox{\textSFxi{}}~ let x = 0b1021;
    ~\lstbox{\textSFv{}}~         ^^^^^^
    ~\lstbox{\textSFxi{}}~ Note :
    ~\lstbox{\textSFxi{}}~  --> main.yr:(1,13)
    ~\lstbox{\textSFxi{}}~  1  ~\lstbox{\textSFxi{}}~ let x = 0b1021;
    ~\lstbox{\textSFxi{}}~     ~\lstbox{\textSFv{}}~             ^
    ~\lstbox{\textSFxi{}}~     ~\lstbox{\textSFii{}}~~\lstbox{\textSFx{}}~~\lstbox{\textSFx{}}~~\lstbox{\textSFx{}}~~\lstbox{\textSFx{}}~~\lstbox{\textSFx{}}~~\lstbox{\textSFx{}}~~\lstbox{\textSFx{}}~
    ~\lstbox{\textSFii{}}~~\lstbox{\textSFx{}}~~\lstbox{\textSFx{}}~~\lstbox{\textSFx{}}~~\lstbox{\textSFx{}}~~\lstbox{\textSFx{}}~~\lstbox{\textSFvii{}}~~\lstbox{\textSFx{}}~
\end{lstlisting}

\subsubsection*{How to fix it}

Use only \tokennolst{0} and \tokennolst{1} in a binary literal, with \tokennolst{\_} as an optional separator. Write the value in another base if binary is not what was wanted.

\errorcode{E1012}{invalid numeric value formatted in octal}
\label{sec:codes:e1012}

Raised during \textbf{lexing}, from \texttt{LexingErrorMessage::\allowbreak{}INVALID\_\allowbreak{}OCT\_\allowbreak{}NUMBER} (\textit{src/\allowbreak{}ymirc/\allowbreak{}lexing/\allowbreak{}errors.yr}).

\subsubsection*{What it means}

A literal introduced by \tokennolst{0o} contains a digit outside \tokennolst{0} to \tokennolst{7}.

\subsubsection*{Example}

\textit{test\_\allowbreak{}resources/\allowbreak{}lexing/\allowbreak{}test21.yr}:

%% check: skip
\begin{lstlisting}[style=coloredverbatimError]
let x = 0o1892;
\end{lstlisting}

\begin{lstlisting}[style=bashVerb, breaklines=true, escapechar=~]
Error[E1012] : invalid numeric value formatted in octal
 --> main.yr:(1,9)
 1  ~\lstbox{\textSFxi{}}~ let x = 0o1892;
    ~\lstbox{\textSFv{}}~         ^^^^^^
    ~\lstbox{\textSFxi{}}~ Note :
    ~\lstbox{\textSFxi{}}~  --> main.yr:(1,12)
    ~\lstbox{\textSFxi{}}~  1  ~\lstbox{\textSFxi{}}~ let x = 0o1892;
    ~\lstbox{\textSFxi{}}~     ~\lstbox{\textSFv{}}~            ^
    ~\lstbox{\textSFxi{}}~     ~\lstbox{\textSFii{}}~~\lstbox{\textSFx{}}~~\lstbox{\textSFx{}}~~\lstbox{\textSFx{}}~~\lstbox{\textSFx{}}~~\lstbox{\textSFx{}}~~\lstbox{\textSFx{}}~~\lstbox{\textSFx{}}~
    ~\lstbox{\textSFii{}}~~\lstbox{\textSFx{}}~~\lstbox{\textSFx{}}~~\lstbox{\textSFx{}}~~\lstbox{\textSFx{}}~~\lstbox{\textSFx{}}~~\lstbox{\textSFvii{}}~~\lstbox{\textSFx{}}~
\end{lstlisting}

\subsubsection*{How to fix it}

Use only octal digits. A leading zero does not make a number octal in Ymir -- the \tokennolst{0o} prefix does -- so \tokennolst{09} is a decimal literal, not an octal one.

\errorcode{E1013}{expected sign information in hexadecimal floating-point literal}
\label{sec:codes:e1013}

Raised during \textbf{lexing}, from \texttt{LexingErrorMessage::\allowbreak{}MISSING\_\allowbreak{}HEX\_\allowbreak{}FLOAT\_\allowbreak{}SIGN} (\textit{src/\allowbreak{}ymirc/\allowbreak{}lexing/\allowbreak{}errors.yr}).

\subsubsection*{What it means}

A hexadecimal floating-point literal has a \tokennolst{p} exponent marker that is not followed by a sign. The exponent of a hexadecimal float is written with an explicit \tokennolst{+} or \tokennolst{-}.

\subsubsection*{Example}

\textit{test\_\allowbreak{}resources/\allowbreak{}lexing/\allowbreak{}test28.yr}:

%% check: skip
\begin{lstlisting}[style=coloredverbatimError]
let x = 0x1.8p;
\end{lstlisting}

\begin{lstlisting}[style=bashVerb, breaklines=true, escapechar=~]
Error[E1013] : expected sign information in hexadecimal floating-point literal
 --> main.yr:(1,9)
 1  ~\lstbox{\textSFxi{}}~ let x = 0x1.8p;
    ~\lstbox{\textSFv{}}~         ^^^
    ~\lstbox{\textSFii{}}~~\lstbox{\textSFx{}}~~\lstbox{\textSFx{}}~~\lstbox{\textSFx{}}~~\lstbox{\textSFx{}}~~\lstbox{\textSFx{}}~~\lstbox{\textSFx{}}~~\lstbox{\textSFx{}}~
\end{lstlisting}

\subsubsection*{How to fix it}

Write the sign: \tokennolst{0x1.8p+3} rather than \tokennolst{0x1.8p3}.

\errorcode{E1014}{expected sign information in scientific floating-point literal}
\label{sec:codes:e1014}

Raised during \textbf{lexing}, from \texttt{LexingErrorMessage::\allowbreak{}MISSING\_\allowbreak{}SCI\_\allowbreak{}FLOAT\_\allowbreak{}SIGN} (\textit{src/\allowbreak{}ymirc/\allowbreak{}lexing/\allowbreak{}errors.yr}).

\subsubsection*{What it means}

A floating-point literal in scientific notation has an \tokennolst{e} exponent marker that is not followed by a sign.

\subsubsection*{Example}

\textit{test\_\allowbreak{}resources/\allowbreak{}lexing/\allowbreak{}test27.yr}:

%% check: skip
\begin{lstlisting}[style=coloredverbatimError]
let x = 1.0e;
\end{lstlisting}

\begin{lstlisting}[style=bashVerb, breaklines=true, escapechar=~]
Error[E1014] : expected sign information in scientific floating-point literal
 --> main.yr:(1,9)
 1  ~\lstbox{\textSFxi{}}~ let x = 1.0e;
    ~\lstbox{\textSFv{}}~         ^
    ~\lstbox{\textSFii{}}~~\lstbox{\textSFx{}}~~\lstbox{\textSFx{}}~~\lstbox{\textSFx{}}~~\lstbox{\textSFx{}}~~\lstbox{\textSFx{}}~~\lstbox{\textSFx{}}~~\lstbox{\textSFx{}}~
\end{lstlisting}

\subsubsection*{How to fix it}

Write the sign: \tokennolst{1.5e+10} rather than \tokennolst{1.5e10}.

\errorcode{E1015}{\errhole{} is not a valid identifier; use backquotes for special identifiers(e.g. `\errhole{}`)}
\label{sec:codes:e1015}

Raised during \textbf{lexing}, from \texttt{LexingErrorMessage::\allowbreak{}INVALID\_\allowbreak{}IDENTIFIER} (\textit{src/\allowbreak{}ymirc/\allowbreak{}lexing/\allowbreak{}errors.yr}).

\subsubsection*{What it means}

The name reported is not a valid identifier: it either starts with a character identifiers cannot start with, or it is a reserved keyword. Ymir can still name such a symbol, but the name has to be quoted.

\subsubsection*{Example}

\textit{test\_\allowbreak{}resources/\allowbreak{}lexing/\allowbreak{}test30.yr}:

%% check: skip
\begin{lstlisting}[style=coloredverbatimError]
let __ = 1;
\end{lstlisting}

\begin{lstlisting}[style=bashVerb, breaklines=true, escapechar=~]
Error[E1015] : __ is not a valid identifier; use backquotes for special identifiers(e.g. `__`)
 --> main.yr:(1,5)
 1  ~\lstbox{\textSFxi{}}~ let __ = 1;
    ~\lstbox{\textSFv{}}~     ^^
    ~\lstbox{\textSFii{}}~~\lstbox{\textSFx{}}~~\lstbox{\textSFx{}}~~\lstbox{\textSFx{}}~~\lstbox{\textSFx{}}~~\lstbox{\textSFx{}}~~\lstbox{\textSFx{}}~~\lstbox{\textSFx{}}~
\end{lstlisting}

\subsubsection*{How to fix it}

Rename the symbol, or write it between backquotes -- \tokennolst{`if`} -- which makes any text usable as an identifier. This is how a symbol coming from another language whose name collides with an Ymir keyword is referred to.

\errorcode{E1016}{binary number is empty}
\label{sec:codes:e1016}

Raised during \textbf{lexing}, from \texttt{LexingErrorMessage::\allowbreak{}EMPTY\_\allowbreak{}BIN\_\allowbreak{}NUMBER} (\textit{src/\allowbreak{}ymirc/\allowbreak{}lexing/\allowbreak{}errors.yr}).

\subsubsection*{What it means}

A \tokennolst{0b} prefix is not followed by any binary digit, so the literal has no value.

\subsubsection*{Example}

\textit{test\_\allowbreak{}resources/\allowbreak{}lexing/\allowbreak{}test20.yr}:

%% check: skip
\begin{lstlisting}[style=coloredverbatimError]
let x = 0b_;
\end{lstlisting}

\begin{lstlisting}[style=bashVerb, breaklines=true, escapechar=~]
Error[E1016] : binary number is empty
 --> main.yr:(1,9)
 1  ~\lstbox{\textSFxi{}}~ let x = 0b_;
    ~\lstbox{\textSFv{}}~         ^^^
    ~\lstbox{\textSFii{}}~~\lstbox{\textSFx{}}~~\lstbox{\textSFx{}}~~\lstbox{\textSFx{}}~~\lstbox{\textSFx{}}~~\lstbox{\textSFx{}}~~\lstbox{\textSFx{}}~~\lstbox{\textSFx{}}~
\end{lstlisting}

\subsubsection*{How to fix it}

Write the digits after the prefix, or remove the prefix if a plain \tokennolst{0} was meant.

\errorcode{E1017}{octal number is empty}
\label{sec:codes:e1017}

Raised during \textbf{lexing}, from \texttt{LexingErrorMessage::\allowbreak{}EMPTY\_\allowbreak{}OCT\_\allowbreak{}NUMBER} (\textit{src/\allowbreak{}ymirc/\allowbreak{}lexing/\allowbreak{}errors.yr}).

\subsubsection*{What it means}

A \tokennolst{0o} prefix is not followed by any octal digit, so the literal has no value.

\subsubsection*{Example}

\textit{test\_\allowbreak{}resources/\allowbreak{}lexing/\allowbreak{}test22.yr}:

%% check: skip
\begin{lstlisting}[style=coloredverbatimError]
let x = 0o_;
\end{lstlisting}

\begin{lstlisting}[style=bashVerb, breaklines=true, escapechar=~]
Error[E1017] : octal number is empty
 --> main.yr:(1,9)
 1  ~\lstbox{\textSFxi{}}~ let x = 0o_;
    ~\lstbox{\textSFv{}}~         ^^^
    ~\lstbox{\textSFii{}}~~\lstbox{\textSFx{}}~~\lstbox{\textSFx{}}~~\lstbox{\textSFx{}}~~\lstbox{\textSFx{}}~~\lstbox{\textSFx{}}~~\lstbox{\textSFx{}}~~\lstbox{\textSFx{}}~
\end{lstlisting}

\subsubsection*{How to fix it}

Write the digits after the prefix, or remove the prefix if a plain \tokennolst{0} was meant.

\errorcode{E1018}{hexadecimal number is empty}
\label{sec:codes:e1018}

Raised during \textbf{lexing}, from \texttt{LexingErrorMessage::\allowbreak{}EMPTY\_\allowbreak{}HEX\_\allowbreak{}NUMBER} (\textit{src/\allowbreak{}ymirc/\allowbreak{}lexing/\allowbreak{}errors.yr}).

\subsubsection*{What it means}

A \tokennolst{0x} prefix is not followed by any hexadecimal digit, so the literal has no value.

\subsubsection*{Example}

\textit{test\_\allowbreak{}resources/\allowbreak{}lexing/\allowbreak{}test24.yr}:

%% check: skip
\begin{lstlisting}[style=coloredverbatimError]
let x = 0x_;
\end{lstlisting}

\begin{lstlisting}[style=bashVerb, breaklines=true, escapechar=~]
Error[E1018] : hexadecimal number is empty
 --> main.yr:(1,9)
 1  ~\lstbox{\textSFxi{}}~ let x = 0x_;
    ~\lstbox{\textSFv{}}~         ^^^
    ~\lstbox{\textSFii{}}~~\lstbox{\textSFx{}}~~\lstbox{\textSFx{}}~~\lstbox{\textSFx{}}~~\lstbox{\textSFx{}}~~\lstbox{\textSFx{}}~~\lstbox{\textSFx{}}~~\lstbox{\textSFx{}}~
\end{lstlisting}

\subsubsection*{How to fix it}

Write the digits after the prefix, or remove the prefix if a plain \tokennolst{0} was meant.

\errorcode{E1019}{the exponent part of the floating-point number is empty}
\label{sec:codes:e1019}

Raised during \textbf{lexing}, from \texttt{LexingErrorMessage::\allowbreak{}EMPTY\_\allowbreak{}EXP\_\allowbreak{}FLOAT\_\allowbreak{}NUMBER} (\textit{src/\allowbreak{}ymirc/\allowbreak{}lexing/\allowbreak{}errors.yr}).

\subsubsection*{What it means}

The exponent marker of a floating-point literal is followed by a sign but by no digit, so the exponent has no value.

\subsubsection*{Example}

\textit{test\_\allowbreak{}resources/\allowbreak{}lexing/\allowbreak{}test5.yr}:

%% check: skip
\begin{lstlisting}[style=coloredverbatimError]
1.e+
\end{lstlisting}

\begin{lstlisting}[style=bashVerb, breaklines=true, escapechar=~]
Error[E1019] : the exponent part of the floating-point number is empty
 --> main.yr:(1,1)
 1  ~\lstbox{\textSFxi{}}~ 1.e+
    ~\lstbox{\textSFv{}}~ ^
    ~\lstbox{\textSFii{}}~~\lstbox{\textSFx{}}~~\lstbox{\textSFx{}}~~\lstbox{\textSFx{}}~~\lstbox{\textSFx{}}~~\lstbox{\textSFx{}}~~\lstbox{\textSFx{}}~~\lstbox{\textSFx{}}~
\end{lstlisting}

\subsubsection*{How to fix it}

Write the exponent digits, as in \tokennolst{1.0e+10}, or drop the exponent part entirely.

\errorcode{E1020}{the fractional part of the floating-point number is empty}
\label{sec:codes:e1020}

Raised during \textbf{lexing}, from \texttt{LexingErrorMessage::\allowbreak{}EMPTY\_\allowbreak{}DEC\_\allowbreak{}FLOAT\_\allowbreak{}NUMBER} (\textit{src/\allowbreak{}ymirc/\allowbreak{}lexing/\allowbreak{}errors.yr}).

\subsubsection*{What it means}

A floating-point literal has a decimal point with no digit after it.

\subsubsection*{Example}

\textit{test\_\allowbreak{}resources/\allowbreak{}lexing/\allowbreak{}test26.yr}:

%% check: skip
\begin{lstlisting}[style=coloredverbatimError]
let x = 1._;
\end{lstlisting}

\begin{lstlisting}[style=bashVerb, breaklines=true, escapechar=~]
Error[E1020] : the fractional part of the floating-point number is empty
 --> main.yr:(1,9)
 1  ~\lstbox{\textSFxi{}}~ let x = 1._;
    ~\lstbox{\textSFv{}}~         ^
    ~\lstbox{\textSFxi{}}~ Note :
    ~\lstbox{\textSFxi{}}~  --> main.yr:(1,11)
    ~\lstbox{\textSFxi{}}~  1  ~\lstbox{\textSFxi{}}~ let x = 1._;
    ~\lstbox{\textSFxi{}}~     ~\lstbox{\textSFv{}}~           ^
    ~\lstbox{\textSFxi{}}~     ~\lstbox{\textSFii{}}~~\lstbox{\textSFx{}}~~\lstbox{\textSFx{}}~~\lstbox{\textSFx{}}~~\lstbox{\textSFx{}}~~\lstbox{\textSFx{}}~~\lstbox{\textSFx{}}~~\lstbox{\textSFx{}}~
    ~\lstbox{\textSFii{}}~~\lstbox{\textSFx{}}~~\lstbox{\textSFx{}}~~\lstbox{\textSFx{}}~~\lstbox{\textSFx{}}~~\lstbox{\textSFx{}}~~\lstbox{\textSFvii{}}~~\lstbox{\textSFx{}}~
\end{lstlisting}

\subsubsection*{How to fix it}

Write the fractional digits (\tokennolst{1.0}), or drop the point if an integer was meant. A trailing point is not a valid way to write a float in Ymir.

\errorcode{E1021}{missing the exponent part in a floating-point number formatted in hexadecimal}
\label{sec:codes:e1021}

Raised during \textbf{lexing}, from \texttt{LexingErrorMessage::\allowbreak{}MISSING\_\allowbreak{}P\_\allowbreak{}IN\_\allowbreak{}HEX\_\allowbreak{}FLOAT} (\textit{src/\allowbreak{}ymirc/\allowbreak{}lexing/\allowbreak{}errors.yr}).

\subsubsection*{What it means}

A hexadecimal literal contains a decimal point, which makes it a floating-point literal, but it has no \tokennolst{p} exponent part. A hexadecimal float always carries its binary exponent explicitly.

\subsubsection*{Example}

\textit{test\_\allowbreak{}resources/\allowbreak{}lexing/\allowbreak{}test29.yr}:

%% check: skip
\begin{lstlisting}[style=coloredverbatimError]
let x = 0x1.8;
\end{lstlisting}

\begin{lstlisting}[style=bashVerb, breaklines=true, escapechar=~]
Error[E1021] : missing the exponent part in a floating-point number formatted in hexadecimal
 --> main.yr:(1,9)
 1  ~\lstbox{\textSFxi{}}~ let x = 0x1.8;
    ~\lstbox{\textSFv{}}~         ^^^
    ~\lstbox{\textSFxi{}}~ Note :
    ~\lstbox{\textSFxi{}}~  --> main.yr:(1,12)
    ~\lstbox{\textSFxi{}}~  1  ~\lstbox{\textSFxi{}}~ let x = 0x1.8;
    ~\lstbox{\textSFxi{}}~     ~\lstbox{\textSFv{}}~            ^
    ~\lstbox{\textSFxi{}}~     ~\lstbox{\textSFii{}}~~\lstbox{\textSFx{}}~~\lstbox{\textSFx{}}~~\lstbox{\textSFx{}}~~\lstbox{\textSFx{}}~~\lstbox{\textSFx{}}~~\lstbox{\textSFx{}}~~\lstbox{\textSFx{}}~
    ~\lstbox{\textSFii{}}~~\lstbox{\textSFx{}}~~\lstbox{\textSFx{}}~~\lstbox{\textSFx{}}~~\lstbox{\textSFx{}}~~\lstbox{\textSFx{}}~~\lstbox{\textSFvii{}}~~\lstbox{\textSFx{}}~
\end{lstlisting}

\subsubsection*{How to fix it}

Add the exponent: \tokennolst{0x1.8p+0}. Use a decimal literal instead if the exponent notation is not what was wanted.

\errorcode{E1022}{unterminated version block: missing \#end}
\label{sec:codes:e1022}

Raised during \textbf{lexing}, from \texttt{LexingErrorMessage::\allowbreak{}UNTERMINATED\_\allowbreak{}VERSION\_\allowbreak{}BLOCK} (\textit{src/\allowbreak{}ymirc/\allowbreak{}lexing/\allowbreak{}errors.yr}).

\subsubsection*{What it means}

A \tokennolst{\#if} conditional compilation block was opened and the file ended before the matching \tokennolst{\#end} was found.

\subsubsection*{Example}

\textit{test\_\allowbreak{}resources/\allowbreak{}lexing/\allowbreak{}test9.yr}:

%% check: skip
\begin{lstlisting}[style=coloredverbatimError]
#if FOO
\end{lstlisting}

\begin{lstlisting}[style=bashVerb, breaklines=true, escapechar=~]
Error[E1022] : unterminated version block: missing #end
 --> main.yr:(1,2)
 1  ~\lstbox{\textSFxi{}}~ #if FOO
    ~\lstbox{\textSFv{}}~  ^^
    ~\lstbox{\textSFii{}}~~\lstbox{\textSFx{}}~~\lstbox{\textSFx{}}~~\lstbox{\textSFx{}}~~\lstbox{\textSFx{}}~~\lstbox{\textSFx{}}~~\lstbox{\textSFx{}}~~\lstbox{\textSFx{}}~
\end{lstlisting}

\subsubsection*{How to fix it}

Close the block with \tokennolst{\#end}. Every \tokennolst{\#if} needs one, including the ones whose \tokennolst{\#else} branch is empty.

\errorcode{E1023}{\#\errhole{} without matching \#if}
\label{sec:codes:e1023}

Raised during \textbf{lexing}, from \texttt{LexingErrorMessage::\allowbreak{}UNMATCHED\_\allowbreak{}VERSION\_\allowbreak{}END} (\textit{src/\allowbreak{}ymirc/\allowbreak{}lexing/\allowbreak{}errors.yr}).

\subsubsection*{What it means}

A \tokennolst{\#else} or \tokennolst{\#end} directive was found with no \tokennolst{\#if} open at that point -- either the \tokennolst{\#if} is missing, or an earlier \tokennolst{\#end} already closed it.

\subsubsection*{Example}

\textit{test\_\allowbreak{}resources/\allowbreak{}lexing/\allowbreak{}test8.yr}:

%% check: skip
\begin{lstlisting}[style=coloredverbatimError]
#else
\end{lstlisting}

\begin{lstlisting}[style=bashVerb, breaklines=true, escapechar=~]
Error[E1023] : #else without matching #if
 --> main.yr:(1,2)
 1  ~\lstbox{\textSFxi{}}~ #else
    ~\lstbox{\textSFv{}}~  ^^^^
    ~\lstbox{\textSFii{}}~~\lstbox{\textSFx{}}~~\lstbox{\textSFx{}}~~\lstbox{\textSFx{}}~~\lstbox{\textSFx{}}~~\lstbox{\textSFx{}}~~\lstbox{\textSFx{}}~~\lstbox{\textSFx{}}~
\end{lstlisting}

\subsubsection*{How to fix it}

Add the missing \tokennolst{\#if}, or remove the stray directive. Check the nesting of the enclosing blocks: an \tokennolst{\#end} written one level too early closes the outer block and leaves the next directive unmatched.

\errorcode{E1024}{missing version condition: expected identifier or backquoted identifier}
\label{sec:codes:e1024}

Raised during \textbf{lexing}, from \texttt{LexingErrorMessage::\allowbreak{}MISSING\_\allowbreak{}VERSION\_\allowbreak{}CONDITION} (\textit{src/\allowbreak{}ymirc/\allowbreak{}lexing/\allowbreak{}errors.yr}).

\subsubsection*{What it means}

A \tokennolst{\#if} directive is not followed by the version identifier it should test.

\subsubsection*{Example}

\textit{test\_\allowbreak{}resources/\allowbreak{}lexing/\allowbreak{}test7.yr}:

%% check: skip
\begin{lstlisting}[style=coloredverbatimError]
#if
foo
#end
\end{lstlisting}

\begin{lstlisting}[style=bashVerb, breaklines=true, escapechar=~]
Error[E1024] : missing version condition: expected identifier or backquoted identifier
 --> main.yr:(1,2)
 1  ~\lstbox{\textSFxi{}}~ #if
    ~\lstbox{\textSFv{}}~  ^^
    ~\lstbox{\textSFii{}}~~\lstbox{\textSFx{}}~~\lstbox{\textSFx{}}~~\lstbox{\textSFx{}}~~\lstbox{\textSFx{}}~~\lstbox{\textSFx{}}~~\lstbox{\textSFx{}}~~\lstbox{\textSFx{}}~
\end{lstlisting}

\subsubsection*{How to fix it}

Name the version to test after the directive, as in \tokennolst{\#if linux}. An identifier that is not a plain word has to be backquoted.

\errorcode{E1025}{undefined version directive; expecting if, else or end}
\label{sec:codes:e1025}

Raised during \textbf{lexing}, from \texttt{LexingErrorMessage::\allowbreak{}UNDEFINED\_\allowbreak{}VERSION\_\allowbreak{}DIRECTIVE} (\textit{src/\allowbreak{}ymirc/\allowbreak{}lexing/\allowbreak{}errors.yr}).

\subsubsection*{What it means}

The word following \tokennolst{\#} is not one of the directives the lexer knows. The only conditional compilation directives are \tokennolst{\#if}, \tokennolst{\#else} and \tokennolst{\#end}.

\subsubsection*{Example}

\textit{test\_\allowbreak{}resources/\allowbreak{}lexing/\allowbreak{}test10.yr}:

%% check: skip
\begin{lstlisting}[style=coloredverbatimError]
#foo
\end{lstlisting}

\begin{lstlisting}[style=bashVerb, breaklines=true, escapechar=~]
Error[E1025] : undefined version directive; expecting if, else or end
 --> main.yr:(1,1)
 1  ~\lstbox{\textSFxi{}}~ #foo
    ~\lstbox{\textSFv{}}~ ^
    ~\lstbox{\textSFxi{}}~ Note :
    ~\lstbox{\textSFxi{}}~  --> main.yr:(1,2)
    ~\lstbox{\textSFxi{}}~  1  ~\lstbox{\textSFxi{}}~ #foo
    ~\lstbox{\textSFxi{}}~     ~\lstbox{\textSFv{}}~  ^^^
    ~\lstbox{\textSFxi{}}~     ~\lstbox{\textSFii{}}~~\lstbox{\textSFx{}}~~\lstbox{\textSFx{}}~~\lstbox{\textSFx{}}~~\lstbox{\textSFx{}}~~\lstbox{\textSFx{}}~~\lstbox{\textSFx{}}~~\lstbox{\textSFx{}}~
    ~\lstbox{\textSFii{}}~~\lstbox{\textSFx{}}~~\lstbox{\textSFx{}}~~\lstbox{\textSFx{}}~~\lstbox{\textSFx{}}~~\lstbox{\textSFx{}}~~\lstbox{\textSFvii{}}~~\lstbox{\textSFx{}}~
\end{lstlisting}

\subsubsection*{How to fix it}

Correct the spelling of the directive, or remove the \tokennolst{\#} if the line was not meant to be a directive at all.

%% Generated by tools/gen_error_codes.py from docs/errors/ of the compiler;
%% edit the compiler's pages or tools/error_themes.txt, then rerun it.

\clearpage
\section{Syntax}
\label{sec:codes:syntax}

\begin{longtable}{@{}l p{0.78\linewidth} r@{}}
\toprule Code & Message & Page\\ \midrule \endhead
\hyperref[sec:codes:e2001]{E2001} & block is never closed & \pageref{sec:codes:e2001}\\
\hyperref[sec:codes:e2002]{E2002} & decorator \errhole{} is defined multiple times & \pageref{sec:codes:e2002}\\
\hyperref[sec:codes:e2004]{E2004} & scope guard \errhole{} is declared multiple times & \pageref{sec:codes:e2004}\\
\hyperref[sec:codes:e2006]{E2006} & decorator \errhole{} cannot be applied to a subpattern variable & \pageref{sec:codes:e2006}\\
\hyperref[sec:codes:e2007]{E2007} & unexpected \errhole{} & \pageref{sec:codes:e2007}\\
\hyperref[sec:codes:e2008]{E2008} & unexpected attribute list \errhole{} & \pageref{sec:codes:e2008}\\
\hyperref[sec:codes:e2009]{E2009} & read \errhole{}, but expected \errhole{} & \pageref{sec:codes:e2009}\\
\hyperref[sec:codes:e2012]{E2012} & a named move ctor requires an argument list & \pageref{sec:codes:e2012}\\
\hyperref[sec:codes:e2013]{E2013} & missing the value moved by the operator & \pageref{sec:codes:e2013}\\
\hyperref[sec:codes:e2014]{E2014} & attribute \errhole{} takes no argument & \pageref{sec:codes:e2014}\\
\hyperref[sec:codes:e2015]{E2015} & attribute \errhole{} takes \errhole{} argument(s), but \errhole{} were given & \pageref{sec:codes:e2015}\\
\bottomrule
\end{longtable}

\errorcode{E2001}{block is never closed}
\label{sec:codes:e2001}

Raised during \textbf{parsing}, from \texttt{SyntaxErrorMessage::\allowbreak{}BLOCK\_\allowbreak{}NEVER\_\allowbreak{}CLOSED} (\textit{src/\allowbreak{}ymirc/\allowbreak{}syntax/\allowbreak{}errors.yr}).

\subsubsection*{What it means}

The parser reached the end of the file, or a token that cannot appear there, while a block opened earlier was still waiting for its closing brace. The diagnostic underlines the whole of what the block swallowed, from the opening brace to the end of the file: the imbalance starts at the brace, and everything after it was read as part of the block.

\subsubsection*{Example}

\textit{test\_\allowbreak{}resources/\allowbreak{}diagnostics/\allowbreak{}test87.yr}:

%% check: skip
\begin{lstlisting}[style=coloredverbatimError]
// a function body never closed: the region runs from the '{' to the end of the file
fn main () {
    let a = 12;
    foo (a);
\end{lstlisting}

\begin{lstlisting}[style=bashVerb, breaklines=true, escapechar=~]
Error[E2001] : block is never closed
 --> test_resources/diagnostics/test87.yr:(2,12)
 2  ~\lstbox{\textSFxi{}}~   fn main () {
    ~\lstbox{\textSFv{}}~ ~\lstbox{\textSFi{}}~~\lstbox{\textSFx{}}~~\lstbox{\textSFx{}}~~\lstbox{\textSFx{}}~~\lstbox{\textSFx{}}~~\lstbox{\textSFx{}}~~\lstbox{\textSFx{}}~~\lstbox{\textSFx{}}~~\lstbox{\textSFx{}}~~\lstbox{\textSFx{}}~~\lstbox{\textSFx{}}~~\lstbox{\textSFx{}}~~\lstbox{\textSFx{}}~^
 3  ~\lstbox{\textSFxi{}}~ ~\lstbox{\textSFxi{}}~     let a = 12;
 4  ~\lstbox{\textSFxi{}}~ ~\lstbox{\textSFxi{}}~     foo (a);
    ~\lstbox{\textSFv{}}~ ~\lstbox{\textSFii{}}~~\lstbox{\textSFx{}}~~\lstbox{\textSFx{}}~~\lstbox{\textSFx{}}~~\lstbox{\textSFx{}}~~\lstbox{\textSFx{}}~~\lstbox{\textSFx{}}~~\lstbox{\textSFx{}}~~\lstbox{\textSFx{}}~~\lstbox{\textSFx{}}~~\lstbox{\textSFx{}}~~\lstbox{\textSFx{}}~~\lstbox{\textSFx{}}~^
    ~\lstbox{\textSFii{}}~~\lstbox{\textSFx{}}~~\lstbox{\textSFx{}}~~\lstbox{\textSFx{}}~~\lstbox{\textSFx{}}~~\lstbox{\textSFx{}}~~\lstbox{\textSFx{}}~~\lstbox{\textSFx{}}~
Error[E2009] : read (, but expected {
 --> test_resources/diagnostics/test87.yr:(2,9)
 2  ~\lstbox{\textSFxi{}}~ fn main () {
    ~\lstbox{\textSFv{}}~         ^  note: when reading template symbol
    ~\lstbox{\textSFii{}}~~\lstbox{\textSFx{}}~~\lstbox{\textSFx{}}~~\lstbox{\textSFx{}}~~\lstbox{\textSFx{}}~~\lstbox{\textSFx{}}~~\lstbox{\textSFx{}}~~\lstbox{\textSFx{}}~
\end{lstlisting}

\subsubsection*{How to fix it}

Close the block. When the braces look balanced, the culprit is usually a brace opened inside a string literal or a comment further up, or a missing one on an earlier declaration that swallowed the rest of the file.

\errorcode{E2002}{decorator \errhole{} is defined multiple times}
\label{sec:codes:e2002}

Raised during \textbf{parsing}, from \texttt{SyntaxErrorMessage::\allowbreak{}MULTIPLE\_\allowbreak{}ATTRS} (\textit{src/\allowbreak{}ymirc/\allowbreak{}syntax/\allowbreak{}errors.yr}).

\subsubsection*{What it means}

The same decorator was written twice on one declaration. A decorator states a property; stating it twice says nothing more and is more likely a mistake than an intention.

\subsubsection*{Example}

\textit{test\_\allowbreak{}resources/\allowbreak{}diagnostics/\allowbreak{}test1.yr}:

%% check: skip
\begin{lstlisting}[style=coloredverbatimError]
@final @final
fn foo ()-> i32 { 0 }

fn main () { foo (); }
\end{lstlisting}

\begin{lstlisting}[style=bashVerb, breaklines=true, escapechar=~]
Error[E2002] : decorator final is defined multiple times
 --> test_resources/diagnostics/test1.yr:(1,2)
 1  ~\lstbox{\textSFxi{}}~ @final @final
    ~\lstbox{\textSFv{}}~  ^^^^^  ^^^^^
    ~\lstbox{\textSFii{}}~~\lstbox{\textSFx{}}~~\lstbox{\textSFx{}}~~\lstbox{\textSFx{}}~~\lstbox{\textSFx{}}~~\lstbox{\textSFx{}}~~\lstbox{\textSFx{}}~~\lstbox{\textSFx{}}~
\end{lstlisting}

\subsubsection*{How to fix it}

Remove the duplicate. If two different decorators were meant, check the spelling of the second one.

\errorcode{E2004}{scope guard \errhole{} is declared multiple times}
\label{sec:codes:e2004}

Raised during \textbf{parsing}, from \texttt{SyntaxErrorMessage::\allowbreak{}MULTIPLE\_\allowbreak{}SAME\_\allowbreak{}GUARD} (\textit{src/\allowbreak{}ymirc/\allowbreak{}syntax/\allowbreak{}errors.yr}).

\subsubsection*{What it means}

The same kind of scope guard was declared twice in one block. A block has at most one \tokennolst{exit}, one \tokennolst{success} and one \tokennolst{failure} guard.

\subsubsection*{Example}

\textit{test\_\allowbreak{}resources/\allowbreak{}diagnostics/\allowbreak{}test2.yr}:

%% check: skip
\begin{lstlisting}[style=coloredverbatimError]
fn main () {
    {
    } exit {
    } exit {
    }
}
\end{lstlisting}

\begin{lstlisting}[style=bashVerb, breaklines=true, escapechar=~]
Error[E2004] : scope guard exit is declared multiple times
 --> test_resources/diagnostics/test2.yr:(3,7)
 3  ~\lstbox{\textSFxi{}}~     } exit {
    ~\lstbox{\textSFv{}}~       ^^^^
   ...
 --> test_resources/diagnostics/test2.yr:(4,7)
 4  ~\lstbox{\textSFxi{}}~     } exit {
    ~\lstbox{\textSFv{}}~       ^^^^
    ~\lstbox{\textSFii{}}~~\lstbox{\textSFx{}}~~\lstbox{\textSFx{}}~~\lstbox{\textSFx{}}~~\lstbox{\textSFx{}}~~\lstbox{\textSFx{}}~~\lstbox{\textSFx{}}~~\lstbox{\textSFx{}}~
\end{lstlisting}

\subsubsection*{How to fix it}

Merge the two bodies into a single guard. Two guards of the same kind would have no defined order between them, which is why the second is rejected rather than chained.

\errorcode{E2006}{decorator \errhole{} cannot be applied to a subpattern variable}
\label{sec:codes:e2006}

Raised during \textbf{parsing}, from \texttt{SyntaxErrorMessage::\allowbreak{}UNDEF\_\allowbreak{}DECORATOR\_\allowbreak{}IN\_\allowbreak{}MATCH} (\textit{src/\allowbreak{}ymirc/\allowbreak{}syntax/\allowbreak{}errors.yr}).

\subsubsection*{What it means}

A decorator was written on a variable of a subpattern inside a pattern match. The subpattern variables of a match are bound by the pattern itself; they take their mutability and their storage from the value being matched, so a decorator on them has nothing to apply to.

\subsubsection*{Example}

\textit{test\_\allowbreak{}resources/\allowbreak{}lazyness/\allowbreak{}test14.yr}:

%% check: skip
\begin{lstlisting}[style=coloredverbatimError]
class A {
    pub let x : i32 = 0;
    pub self () {}
}

fn main () {
    let lazy  A(x-> x) = lazy A();
}
\end{lstlisting}

\begin{lstlisting}[style=bashVerb, breaklines=true, escapechar=~]
Error[E2006] : decorator lazy cannot be applied to a subpattern variable
 --> test_resources/lazyness/test14.yr:(7,9)
 7  ~\lstbox{\textSFxi{}}~     let lazy  A(x-> x) = lazy A();
    ~\lstbox{\textSFv{}}~         ^^^^   ^  note: failed to read pattern variable declaration
    ~\lstbox{\textSFii{}}~~\lstbox{\textSFx{}}~~\lstbox{\textSFx{}}~~\lstbox{\textSFx{}}~~\lstbox{\textSFx{}}~~\lstbox{\textSFx{}}~~\lstbox{\textSFx{}}~~\lstbox{\textSFx{}}~
\end{lstlisting}

\subsubsection*{How to fix it}

Remove the decorator. When the binding needs to be mutable, the decorator belongs on the outermost variable of the pattern, not on the nested one.

\errorcode{E2007}{unexpected \errhole{}}
\label{sec:codes:e2007}

Raised during \textbf{parsing}, from \texttt{SyntaxErrorMessage::\allowbreak{}UNEXPECTED} (\textit{src/\allowbreak{}ymirc/\allowbreak{}syntax/\allowbreak{}errors.yr}).

\subsubsection*{What it means}

The parser read a token that cannot appear at that point of the grammar and had no more specific expectation to report. The token itself is in the message, a file stopping there naming itself ``end of file''.

\subsubsection*{Example}

\textit{test\_\allowbreak{}resources/\allowbreak{}diagnostics/\allowbreak{}test3.yr}:

%% check: skip
\begin{lstlisting}[style=coloredverbatimError]
class Foo {
    @final let x: i32;
    pub self (x: i32) with x = x {}
}

fn main () {}
\end{lstlisting}

\begin{lstlisting}[style=bashVerb, breaklines=true, escapechar=~]
Error[E2007] : unexpected @
 --> test_resources/diagnostics/test3.yr:(2,5)
 2  ~\lstbox{\textSFxi{}}~     @final let x: i32;
    ~\lstbox{\textSFv{}}~     ^  note: when reading template symbol
    ~\lstbox{\textSFii{}}~~\lstbox{\textSFx{}}~~\lstbox{\textSFx{}}~~\lstbox{\textSFx{}}~~\lstbox{\textSFx{}}~~\lstbox{\textSFx{}}~~\lstbox{\textSFx{}}~~\lstbox{\textSFx{}}~
\end{lstlisting}

\subsubsection*{How to fix it}

Read the reported token in its context: this is usually a missing separator, a missing closing delimiter earlier on, or a keyword used where an expression was expected.

\errorcode{E2008}{unexpected attribute list \errhole{}}
\label{sec:codes:e2008}

Raised during \textbf{parsing}, from \texttt{SyntaxErrorMessage::\allowbreak{}UNEXPECTED\_\allowbreak{}ATTRIBUTES} (\textit{src/\allowbreak{}ymirc/\allowbreak{}syntax/\allowbreak{}errors.yr}).

\subsubsection*{What it means}

An attribute list was written on a declaration that does not take one.

\subsubsection*{Example}

\textit{test\_\allowbreak{}resources/\allowbreak{}diagnostics/\allowbreak{}test3.yr}:

%% check: skip
\begin{lstlisting}[style=coloredverbatimError]
class Foo {
    @final let x: i32;
    pub self (x: i32) with x = x {}
}

fn main () {}
\end{lstlisting}

\begin{lstlisting}[style=bashVerb, breaklines=true, escapechar=~]
Error[E2008] : unexpected attribute list {final}
 --> test_resources/diagnostics/test3.yr:(2,5)
 2  ~\lstbox{\textSFxi{}}~     @final let x: i32;
    ~\lstbox{\textSFv{}}~     ^
    ~\lstbox{\textSFxi{}}~ Note :
    ~\lstbox{\textSFxi{}}~  --> test_resources/diagnostics/test3.yr:(2,12)
    ~\lstbox{\textSFxi{}}~  2  ~\lstbox{\textSFxi{}}~     @final let x: i32;
    ~\lstbox{\textSFxi{}}~     ~\lstbox{\textSFv{}}~            ^^^
    ~\lstbox{\textSFxi{}}~     ~\lstbox{\textSFii{}}~~\lstbox{\textSFx{}}~~\lstbox{\textSFx{}}~~\lstbox{\textSFx{}}~~\lstbox{\textSFx{}}~~\lstbox{\textSFx{}}~~\lstbox{\textSFx{}}~~\lstbox{\textSFx{}}~
    ~\lstbox{\textSFii{}}~~\lstbox{\textSFx{}}~~\lstbox{\textSFx{}}~~\lstbox{\textSFx{}}~~\lstbox{\textSFx{}}~~\lstbox{\textSFx{}}~~\lstbox{\textSFvii{}}~~\lstbox{\textSFx{}}~
\end{lstlisting}

\subsubsection*{How to fix it}

Remove the list, or move it onto the declaration it was meant for. Attribute lists apply to declarations, not to the statements inside them.

\errorcode{E2009}{read \errhole{}, but expected \errhole{}}
\label{sec:codes:e2009}

Raised during \textbf{parsing}, from \texttt{SyntaxErrorMessage::\allowbreak{}UNEXPECTED\_\allowbreak{}BUT\_\allowbreak{}LST} (\textit{src/\allowbreak{}ymirc/\allowbreak{}syntax/\allowbreak{}errors.yr}).

\subsubsection*{What it means}

The parser expected one of a known set of tokens and read something else. Both what was read and what was expected are in the message, so the set of possibilities at that point is explicit. A file stopping before any of them is read names itself ``end of file''.

\subsubsection*{Example}

\textit{test\_\allowbreak{}resources/\allowbreak{}lit\_\allowbreak{}class/\allowbreak{}test7.yr}:

%% check: skip
\begin{lstlisting}[style=coloredverbatimError]
class Z {
    pub self ()-> &Z {}
}
\end{lstlisting}

\begin{lstlisting}[style=bashVerb, breaklines=true, escapechar=~]
Error[E2009] : read ->, but expected {
 --> test_resources/lit_class/test7.yr:(2,16)
 2  ~\lstbox{\textSFxi{}}~     pub self ()-> &Z {}
    ~\lstbox{\textSFv{}}~                ^^
    ~\lstbox{\textSFii{}}~~\lstbox{\textSFx{}}~~\lstbox{\textSFx{}}~~\lstbox{\textSFx{}}~~\lstbox{\textSFx{}}~~\lstbox{\textSFx{}}~~\lstbox{\textSFx{}}~~\lstbox{\textSFx{}}~
\end{lstlisting}

\subsubsection*{How to fix it}

Replace the token with one of the listed alternatives. A long list of expectations usually means the construct itself was mis-started a few tokens earlier.

\errorcode{E2012}{a named move ctor requires an argument list}
\label{sec:codes:e2012}

Raised during \textbf{parsing}, from \texttt{SyntaxErrorMessage::\allowbreak{}MOVE\_\allowbreak{}NAMED\_\allowbreak{}CTOR\_\allowbreak{}PARAMS} (\textit{src/\allowbreak{}ymirc/\allowbreak{}syntax/\allowbreak{}errors.yr}).

\subsubsection*{What it means}

A \tokennolst{move} naming the ctor that resets the moved value was not followed by an argument list. The parentheses are what separate the arguments of the ctor from the value being moved, so they are mandatory as soon as a name is given -- even when the ctor takes nothing.

\subsubsection*{Example}

\textit{test\_\allowbreak{}resources/\allowbreak{}diagnostics/\allowbreak{}test90.yr}:

%% check: skip
\begin{lstlisting}[style=coloredverbatimError]
entity Foo {
    pub self open () {}
}

fn main () {
    let mut a = Foo::open ();
    let mut _b_ = move::open a;
}
\end{lstlisting}

\begin{lstlisting}[style=bashVerb, breaklines=true, escapechar=~]
Error[E2012] : a named move ctor requires an argument list
 --> test_resources/diagnostics/test90.yr:(8,23)
 8  ~\lstbox{\textSFxi{}}~     let mut _b_ = move::open a;
    ~\lstbox{\textSFv{}}~                       ^^----  note: failed to read standard variable declaration
    ~\lstbox{\textSFii{}}~~\lstbox{\textSFx{}}~~\lstbox{\textSFx{}}~~\lstbox{\textSFx{}}~~\lstbox{\textSFx{}}~~\lstbox{\textSFx{}}~~\lstbox{\textSFx{}}~~\lstbox{\textSFx{}}~
\end{lstlisting}

\subsubsection*{How to fix it}

Write the empty argument list: \tokennolst{move::open () a}.

\errorcode{E2013}{missing the value moved by the operator}
\label{sec:codes:e2013}

Raised during \textbf{parsing}, from \texttt{SyntaxErrorMessage::\allowbreak{}MOVE\_\allowbreak{}MISSING\_\allowbreak{}OPERAND} (\textit{src/\allowbreak{}ymirc/\allowbreak{}syntax/\allowbreak{}errors.yr}).

\subsubsection*{What it means}

A \tokennolst{(} directly following \tokennolst{move} always opens the argument list of the reset ctor, so what looks like the moved value is read as an argument and the operator is left without an operand.

\subsubsection*{Example}

\textit{test\_\allowbreak{}resources/\allowbreak{}diagnostics/\allowbreak{}test89.yr}:

%% check: skip
\begin{lstlisting}[style=coloredverbatimError]
entity Foo {
    pub self () {}
}

fn foo ()-> Foo;

fn main () {
    let mut _a_ = move (foo ());
}
\end{lstlisting}

\begin{lstlisting}[style=bashVerb, breaklines=true, escapechar=~]
Error[E2013] : missing the value moved by the operator
 --> test_resources/diagnostics/test89.yr:(9,19)
 9  ~\lstbox{\textSFxi{}}~     let mut _a_ = move (foo ());
    ~\lstbox{\textSFv{}}~                   ^^^^ -       -
    ~\lstbox{\textSFxi{}}~                      ~\lstbox{\textSFii{}}~~\lstbox{\textSFx{}}~ note: failed to read standard variable declaration
    ~\lstbox{\textSFxi{}}~                        ~\lstbox{\textSFii{}}~~\lstbox{\textSFx{}}~ note: the argument list of move was read instead, write move () (...) to move a parenthesized value
    ~\lstbox{\textSFxi{}}~                                ~\lstbox{\textSFii{}}~~\lstbox{\textSFx{}}~ error: unexpected ;
    ~\lstbox{\textSFii{}}~~\lstbox{\textSFx{}}~~\lstbox{\textSFx{}}~~\lstbox{\textSFx{}}~~\lstbox{\textSFx{}}~~\lstbox{\textSFx{}}~~\lstbox{\textSFx{}}~~\lstbox{\textSFx{}}~
\end{lstlisting}

\subsubsection*{How to fix it}

Drop the parentheses around the moved value, or -- when they are needed to enclose a complex expression -- put an explicit argument list in front of them: \tokennolst{move () (foo ())}.

\errorcode{E2014}{attribute \errhole{} takes no argument}
\label{sec:codes:e2014}

Raised during \textbf{parsing}, from \texttt{SyntaxErrorMessage::\allowbreak{}ATTRIBUTE\_\allowbreak{}WITHOUT\_\allowbreak{}ARGS} (\textit{src/\allowbreak{}ymirc/\allowbreak{}syntax/\allowbreak{}errors.yr}).

\subsubsection*{What it means}

An attribute declares how many arguments it takes, and most take none. The set of attributes is closed, so this is checked as the name is read, next to the error raised for an attribute that does not exist at all.

Only \tokennolst{@parameters} takes an argument today.

\subsubsection*{Example}

\textit{test\_\allowbreak{}resources/\allowbreak{}diagnostics/\allowbreak{}test91.yr}:

%% check: skip
\begin{lstlisting}[style=coloredverbatimError]
@final(3)
fn main () {}
\end{lstlisting}

\begin{lstlisting}[style=bashVerb, breaklines=true, escapechar=~]
Error[E2014] : attribute final takes no argument
 --> test_resources/diagnostics/test91.yr:(2,2)
 2  ~\lstbox{\textSFxi{}}~ @final(3)
    ~\lstbox{\textSFv{}}~  ^^^^^
    ~\lstbox{\textSFii{}}~~\lstbox{\textSFx{}}~~\lstbox{\textSFx{}}~~\lstbox{\textSFx{}}~~\lstbox{\textSFx{}}~~\lstbox{\textSFx{}}~~\lstbox{\textSFx{}}~~\lstbox{\textSFx{}}~
\end{lstlisting}

\subsubsection*{How to fix it}

Drop the argument list.

\errorcode{E2015}{attribute \errhole{} takes \errhole{} argument(s), but \errhole{} were given}
\label{sec:codes:e2015}

Raised during \textbf{parsing}, from \texttt{SyntaxErrorMessage::\allowbreak{}ATTRIBUTE\_\allowbreak{}ARITY} (\textit{src/\allowbreak{}ymirc/\allowbreak{}syntax/\allowbreak{}errors.yr}).

\subsubsection*{What it means}

An attribute taking arguments takes a fixed number of them, and was written with a different number -- or, as below, with none at all.

\subsubsection*{Example}

\textit{test\_\allowbreak{}resources/\allowbreak{}diagnostics/\allowbreak{}test92.yr}:

%% check: skip
\begin{lstlisting}[style=coloredverbatimError]
@parameters
__test (value: i32) {
    assert (value > 0);
}
\end{lstlisting}

\begin{lstlisting}[style=bashVerb, breaklines=true, escapechar=~]
Error[E2015] : attribute parameters takes 1 argument(s), but 0 were given
 --> test_resources/diagnostics/test92.yr:(2,2)
 2  ~\lstbox{\textSFxi{}}~ @parameters
    ~\lstbox{\textSFv{}}~  ^^^^^^^^^^
    ~\lstbox{\textSFii{}}~~\lstbox{\textSFx{}}~~\lstbox{\textSFx{}}~~\lstbox{\textSFx{}}~~\lstbox{\textSFx{}}~~\lstbox{\textSFx{}}~~\lstbox{\textSFx{}}~~\lstbox{\textSFx{}}~
\end{lstlisting}

\subsubsection*{How to fix it}

Give the attribute the arguments it takes. \tokennolst{@parameters} takes the one value producing the parameter sets of the test.

%% Generated by tools/gen_error_codes.py from docs/errors/ of the compiler;
%% edit the compiler's pages or tools/error_themes.txt, then rerun it.

\clearpage
\section{Attributes and decorators}
\label{sec:codes:attributes}

\begin{longtable}{@{}l p{0.78\linewidth} r@{}}
\toprule Code & Message & Page\\ \midrule \endhead
\hyperref[sec:codes:e3019]{E3019} & custom attribute \errhole{} is not valid on classes & \pageref{sec:codes:e3019}\\
\hyperref[sec:codes:e3020]{E3020} & custom attribute \errhole{} is not valid on constructors & \pageref{sec:codes:e3020}\\
\hyperref[sec:codes:e3021]{E3021} & custom attribute \errhole{} is not valid on functions & \pageref{sec:codes:e3021}\\
\hyperref[sec:codes:e3022]{E3022} & custom attribute \errhole{} is not valid on global variables & \pageref{sec:codes:e3022}\\
\hyperref[sec:codes:e3023]{E3023} & custom attribute \errhole{} is not valid on records & \pageref{sec:codes:e3023}\\
\hyperref[sec:codes:e3024]{E3024} & custom attribute \errhole{} is not valid on template methods & \pageref{sec:codes:e3024}\\
\hyperref[sec:codes:e3025]{E3025} & custom attribute \errhole{} may only be applied to classes & \pageref{sec:codes:e3025}\\
\hyperref[sec:codes:e3026]{E3026} & custom attribute \errhole{} may only be applied to records & \pageref{sec:codes:e3026}\\
\hyperref[sec:codes:e3035]{E3035} & a ctor cannot be both no\_move and only\_move & \pageref{sec:codes:e3035}\\
\hyperref[sec:codes:e3036]{E3036} & custom attribute \errhole{} may only be applied to the ctor of an entity & \pageref{sec:codes:e3036}\\
\hyperref[sec:codes:e4041]{E4041} & decorator is not applicable on inner types, already deeply mutable & \pageref{sec:codes:e4041}\\
\hyperref[sec:codes:e4240]{E4240} & decorator \errhole{} is not applicable in that context & \pageref{sec:codes:e4240}\\
\bottomrule
\end{longtable}

\errorcode{E3019}{custom attribute \errhole{} is not valid on classes}
\label{sec:codes:e3019}

Raised during \textbf{declaration}, from \texttt{DeclareErrorMessage::\allowbreak{}UNDEFINED\_\allowbreak{}ATTRIBUTE\_\allowbreak{}FOR\_\allowbreak{}CLASS} (\textit{src/\allowbreak{}ymirc/\allowbreak{}semantic/\allowbreak{}declarator/\allowbreak{}errors.yr}).

\subsubsection*{What it means}

A custom attribute was applied to a class, and the attribute does not accept classes.

\subsubsection*{Example}

\textit{test\_\allowbreak{}resources/\allowbreak{}data\_\allowbreak{}record/\allowbreak{}attrs/\allowbreak{}test2.yr}:

%% check: skip
\begin{lstlisting}[style=coloredverbatimError]
@data
class Foo {
    pub self () {}
}

fn main() {}
\end{lstlisting}

\begin{lstlisting}[style=bashVerb, breaklines=true, escapechar=~]
Error[E3019] : custom attribute data is not valid on classes
 --> test_resources/data_record/attrs/test2.yr:(1,2)
 1  ~\lstbox{\textSFxi{}}~ @data
    ~\lstbox{\textSFv{}}~  ^^^^
    ~\lstbox{\textSFxi{}}~ Note :
    ~\lstbox{\textSFxi{}}~  --> test_resources/data_record/attrs/test2.yr:(2,1)
    ~\lstbox{\textSFxi{}}~  2  ~\lstbox{\textSFxi{}}~ class Foo {
    ~\lstbox{\textSFxi{}}~     ~\lstbox{\textSFv{}}~ ^^^^^
    ~\lstbox{\textSFxi{}}~     ~\lstbox{\textSFii{}}~~\lstbox{\textSFx{}}~~\lstbox{\textSFx{}}~~\lstbox{\textSFx{}}~~\lstbox{\textSFx{}}~~\lstbox{\textSFx{}}~~\lstbox{\textSFx{}}~~\lstbox{\textSFx{}}~
    ~\lstbox{\textSFii{}}~~\lstbox{\textSFx{}}~~\lstbox{\textSFx{}}~~\lstbox{\textSFx{}}~~\lstbox{\textSFx{}}~~\lstbox{\textSFx{}}~~\lstbox{\textSFvii{}}~~\lstbox{\textSFx{}}~
\end{lstlisting}

\subsubsection*{How to fix it}

Remove the attribute, or apply it to a declaration it accepts. The declarations an attribute accepts are stated where the attribute itself is declared.

\errorcode{E3020}{custom attribute \errhole{} is not valid on constructors}
\label{sec:codes:e3020}

Raised during \textbf{declaration}, from \texttt{DeclareErrorMessage::\allowbreak{}UNDEFINED\_\allowbreak{}ATTRIBUTE\_\allowbreak{}FOR\_\allowbreak{}CTOR} (\textit{src/\allowbreak{}ymirc/\allowbreak{}semantic/\allowbreak{}declarator/\allowbreak{}errors.yr}).

\subsubsection*{What it means}

A custom attribute was applied to a constructor, and the attribute does not accept constructors.

\subsubsection*{Example}

\textit{test\_\allowbreak{}resources/\allowbreak{}diagnostics/\allowbreak{}test9.yr}:

%% check: skip
\begin{lstlisting}[style=coloredverbatimError]
class Foo {
    @final pub self () {}
}

fn main () {}
\end{lstlisting}

\begin{lstlisting}[style=bashVerb, breaklines=true, escapechar=~]
Error[E3020] : custom attribute final is not valid on constructors
 --> test_resources/diagnostics/test9.yr:(2,6)
 2  ~\lstbox{\textSFxi{}}~     @final pub self () {}
    ~\lstbox{\textSFv{}}~      ^^^^^
    ~\lstbox{\textSFii{}}~~\lstbox{\textSFx{}}~~\lstbox{\textSFx{}}~~\lstbox{\textSFx{}}~~\lstbox{\textSFx{}}~~\lstbox{\textSFx{}}~~\lstbox{\textSFx{}}~~\lstbox{\textSFx{}}~
\end{lstlisting}

\subsubsection*{How to fix it}

Same as \hyperref[sec:codes:e3019]{E3019}.

\errorcode{E3021}{custom attribute \errhole{} is not valid on functions}
\label{sec:codes:e3021}

Raised during \textbf{declaration}, from \texttt{DeclareErrorMessage::\allowbreak{}UNDEFINED\_\allowbreak{}ATTRIBUTE\_\allowbreak{}FOR\_\allowbreak{}FUNCTION} (\textit{src/\allowbreak{}ymirc/\allowbreak{}semantic/\allowbreak{}declarator/\allowbreak{}errors.yr}).

\subsubsection*{What it means}

A custom attribute was applied to a function, and the attribute does not accept functions.

\subsubsection*{Example}

\textit{test\_\allowbreak{}resources/\allowbreak{}diagnostics/\allowbreak{}test10.yr}:

%% check: skip
\begin{lstlisting}[style=coloredverbatimError]
@abstract
fn foo () {}

fn main () { foo (); }
\end{lstlisting}

\begin{lstlisting}[style=bashVerb, breaklines=true, escapechar=~]
Error[E3021] : custom attribute abstract is not valid on functions
 --> test_resources/diagnostics/test10.yr:(1,2)
 1  ~\lstbox{\textSFxi{}}~ @abstract
    ~\lstbox{\textSFv{}}~  ^^^^^^^^
    ~\lstbox{\textSFii{}}~~\lstbox{\textSFx{}}~~\lstbox{\textSFx{}}~~\lstbox{\textSFx{}}~~\lstbox{\textSFx{}}~~\lstbox{\textSFx{}}~~\lstbox{\textSFx{}}~~\lstbox{\textSFx{}}~
\end{lstlisting}

\subsubsection*{How to fix it}

Same as \hyperref[sec:codes:e3019]{E3019}.

\errorcode{E3022}{custom attribute \errhole{} is not valid on global variables}
\label{sec:codes:e3022}

Raised during \textbf{declaration}, from \texttt{DeclareErrorMessage::\allowbreak{}UNDEFINED\_\allowbreak{}ATTRIBUTE\_\allowbreak{}FOR\_\allowbreak{}GLOBAL} (\textit{src/\allowbreak{}ymirc/\allowbreak{}semantic/\allowbreak{}declarator/\allowbreak{}errors.yr}).

\subsubsection*{What it means}

A custom attribute was applied to a global variable, and the attribute does not accept global variables.

\subsubsection*{Example}

\textit{test\_\allowbreak{}resources/\allowbreak{}diagnostics/\allowbreak{}test11.yr}:

%% check: skip
\begin{lstlisting}[style=coloredverbatimError]
@final lazy X : i32 = 0;

fn main () { X; }
\end{lstlisting}

\begin{lstlisting}[style=bashVerb, breaklines=true, escapechar=~]
Error[E3022] : custom attribute final is not valid on global variables
 --> test_resources/diagnostics/test11.yr:(1,2)
 1  ~\lstbox{\textSFxi{}}~ @final lazy X : i32 = 0;
    ~\lstbox{\textSFv{}}~  ^^^^^
    ~\lstbox{\textSFxi{}}~ Note :
    ~\lstbox{\textSFxi{}}~  --> test_resources/diagnostics/test11.yr:(1,8)
    ~\lstbox{\textSFxi{}}~  1  ~\lstbox{\textSFxi{}}~ @final lazy X : i32 = 0;
    ~\lstbox{\textSFxi{}}~     ~\lstbox{\textSFv{}}~        ^^^^
    ~\lstbox{\textSFxi{}}~     ~\lstbox{\textSFii{}}~~\lstbox{\textSFx{}}~~\lstbox{\textSFx{}}~~\lstbox{\textSFx{}}~~\lstbox{\textSFx{}}~~\lstbox{\textSFx{}}~~\lstbox{\textSFx{}}~~\lstbox{\textSFx{}}~
    ~\lstbox{\textSFii{}}~~\lstbox{\textSFx{}}~~\lstbox{\textSFx{}}~~\lstbox{\textSFx{}}~~\lstbox{\textSFx{}}~~\lstbox{\textSFx{}}~~\lstbox{\textSFvii{}}~~\lstbox{\textSFx{}}~
\end{lstlisting}

\subsubsection*{How to fix it}

Same as \hyperref[sec:codes:e3019]{E3019}.

\errorcode{E3023}{custom attribute \errhole{} is not valid on records}
\label{sec:codes:e3023}

Raised during \textbf{declaration}, from \texttt{DeclareErrorMessage::\allowbreak{}UNDEFINED\_\allowbreak{}ATTRIBUTE\_\allowbreak{}FOR\_\allowbreak{}RECORD} (\textit{src/\allowbreak{}ymirc/\allowbreak{}semantic/\allowbreak{}declarator/\allowbreak{}errors.yr}).

\subsubsection*{What it means}

A custom attribute was applied to a record, and the attribute does not accept records.

\subsubsection*{Example}

\textit{test\_\allowbreak{}resources/\allowbreak{}diagnostics/\allowbreak{}test12.yr}:

%% check: skip
\begin{lstlisting}[style=coloredverbatimError]
@inline
record Foo {
    let x: i32;
}

fn main () { let _ = Foo (1); }
\end{lstlisting}

\begin{lstlisting}[style=bashVerb, breaklines=true, escapechar=~]
Error[E3023] : custom attribute inline is not valid on records
 --> test_resources/diagnostics/test12.yr:(1,2)
 1  ~\lstbox{\textSFxi{}}~ @inline
    ~\lstbox{\textSFv{}}~  ^^^^^^
    ~\lstbox{\textSFxi{}}~ Note :
    ~\lstbox{\textSFxi{}}~  --> test_resources/diagnostics/test12.yr:(2,1)
    ~\lstbox{\textSFxi{}}~  2  ~\lstbox{\textSFxi{}}~ record Foo {
    ~\lstbox{\textSFxi{}}~     ~\lstbox{\textSFv{}}~ ^^^^^^
    ~\lstbox{\textSFxi{}}~     ~\lstbox{\textSFii{}}~~\lstbox{\textSFx{}}~~\lstbox{\textSFx{}}~~\lstbox{\textSFx{}}~~\lstbox{\textSFx{}}~~\lstbox{\textSFx{}}~~\lstbox{\textSFx{}}~~\lstbox{\textSFx{}}~
    ~\lstbox{\textSFii{}}~~\lstbox{\textSFx{}}~~\lstbox{\textSFx{}}~~\lstbox{\textSFx{}}~~\lstbox{\textSFx{}}~~\lstbox{\textSFx{}}~~\lstbox{\textSFvii{}}~~\lstbox{\textSFx{}}~
\end{lstlisting}

\subsubsection*{How to fix it}

Same as \hyperref[sec:codes:e3019]{E3019}.

\errorcode{E3024}{custom attribute \errhole{} is not valid on template methods}
\label{sec:codes:e3024}

Raised during \textbf{declaration}, from \texttt{DeclareErrorMessage::\allowbreak{}UNDEFINED\_\allowbreak{}ATTRIBUTE\_\allowbreak{}FOR\_\allowbreak{}TEMPLATE\_\allowbreak{}METHOD} (\textit{src/\allowbreak{}ymirc/\allowbreak{}semantic/\allowbreak{}declarator/\allowbreak{}errors.yr}).

\subsubsection*{What it means}

A custom attribute was applied to a template method, and the attribute does not accept template methods.

\subsubsection*{Example}

\textit{test\_\allowbreak{}resources/\allowbreak{}diagnostics/\allowbreak{}test13.yr}:

%% check: skip
\begin{lstlisting}[style=coloredverbatimError]
class Foo {
    pub self () {}
    @packed pub fn bar {T} (self, x: T)-> T { x }
}

fn main () { let f = copy Foo (); f.bar (1); }
\end{lstlisting}

\begin{lstlisting}[style=bashVerb, breaklines=true, escapechar=~]
Error[E3024] : custom attribute packed is not valid on template methods
 --> test_resources/diagnostics/test13.yr:(3,6)
 3  ~\lstbox{\textSFxi{}}~     @packed pub fn bar {T} (self, x: T)-> T { x }
    ~\lstbox{\textSFv{}}~      ^^^^^^     --  note: declared here
    ~\lstbox{\textSFxi{}}~
    = when validating -- (test_resources/diagnostics/test13.yr:(1,7))
    ~\lstbox{\textSFii{}}~~\lstbox{\textSFx{}}~~\lstbox{\textSFx{}}~~\lstbox{\textSFx{}}~~\lstbox{\textSFx{}}~~\lstbox{\textSFx{}}~~\lstbox{\textSFx{}}~~\lstbox{\textSFx{}}~
Error[E4104] : symbol test13::Foo has errors
 --> test_resources/diagnostics/test13.yr:(6,27)
 6  ~\lstbox{\textSFxi{}}~ fn main () { let f = copy Foo (); f.bar (1); }
    ~\lstbox{\textSFv{}}~                           ^^^
 --> test_resources/diagnostics/test13.yr:(3,6)
 3  ~\lstbox{\textSFxi{}}~     @packed pub fn bar {T} (self, x: T)-> T { x }
    ~\lstbox{\textSFv{}}~      ------     --
    ~\lstbox{\textSFxi{}}~           ~\lstbox{\textSFii{}}~~\lstbox{\textSFx{}}~ error: custom attribute packed is not valid on template methods
    ~\lstbox{\textSFxi{}}~                  ~\lstbox{\textSFii{}}~~\lstbox{\textSFx{}}~ note: declared here
    ~\lstbox{\textSFxi{}}~
    = when validating -- (test_resources/diagnostics/test13.yr:(1,7))
    = when validating test13::main -- (test_resources/diagnostics/test13.yr:(6,4))
    ~\lstbox{\textSFii{}}~~\lstbox{\textSFx{}}~~\lstbox{\textSFx{}}~~\lstbox{\textSFx{}}~~\lstbox{\textSFx{}}~~\lstbox{\textSFx{}}~~\lstbox{\textSFx{}}~~\lstbox{\textSFx{}}~
Error[E4224] : undefined operator . for types error and bar
 --> test_resources/diagnostics/test13.yr:(6,36)
 6  ~\lstbox{\textSFxi{}}~ fn main () { let f = copy Foo (); f.bar (1); }
    ~\lstbox{\textSFv{}}~                                    ^
    ~\lstbox{\textSFxi{}}~
    = when validating test13::main -- (test_resources/diagnostics/test13.yr:(6,4))
    ~\lstbox{\textSFii{}}~~\lstbox{\textSFx{}}~~\lstbox{\textSFx{}}~~\lstbox{\textSFx{}}~~\lstbox{\textSFx{}}~~\lstbox{\textSFx{}}~~\lstbox{\textSFx{}}~~\lstbox{\textSFx{}}~
\end{lstlisting}

\subsubsection*{How to fix it}

Same as \hyperref[sec:codes:e3019]{E3019}. A template method is only a declaration until it is instantiated, which is why attributes that act on a concrete symbol cannot apply to it.

\errorcode{E3025}{custom attribute \errhole{} may only be applied to classes}
\label{sec:codes:e3025}

Raised during \textbf{declaration}, from \texttt{DeclareErrorMessage::\allowbreak{}UNDEFINED\_\allowbreak{}ATTRIBUTE\_\allowbreak{}NO\_\allowbreak{}CLASS} (\textit{src/\allowbreak{}ymirc/\allowbreak{}semantic/\allowbreak{}declarator/\allowbreak{}errors.yr}).

\subsubsection*{What it means}

The attribute may only be applied to classes, and it was written on something else.

\subsubsection*{Example}

\textit{test\_\allowbreak{}resources/\allowbreak{}diagnostics/\allowbreak{}test14.yr}:

%% check: skip
\begin{lstlisting}[style=coloredverbatimError]
@abstract
record Foo {
    let x: i32;
}

fn main () { let _ = Foo (1); }
\end{lstlisting}

\begin{lstlisting}[style=bashVerb, breaklines=true, escapechar=~]
Error[E3025] : custom attribute abstract may only be applied to classes
 --> test_resources/diagnostics/test14.yr:(1,2)
 1  ~\lstbox{\textSFxi{}}~ @abstract
    ~\lstbox{\textSFv{}}~  ^^^^^^^^
    ~\lstbox{\textSFxi{}}~ Note :
    ~\lstbox{\textSFxi{}}~  --> test_resources/diagnostics/test14.yr:(2,1)
    ~\lstbox{\textSFxi{}}~  2  ~\lstbox{\textSFxi{}}~ record Foo {
    ~\lstbox{\textSFxi{}}~     ~\lstbox{\textSFv{}}~ ^^^^^^
    ~\lstbox{\textSFxi{}}~     ~\lstbox{\textSFii{}}~~\lstbox{\textSFx{}}~~\lstbox{\textSFx{}}~~\lstbox{\textSFx{}}~~\lstbox{\textSFx{}}~~\lstbox{\textSFx{}}~~\lstbox{\textSFx{}}~~\lstbox{\textSFx{}}~
    ~\lstbox{\textSFii{}}~~\lstbox{\textSFx{}}~~\lstbox{\textSFx{}}~~\lstbox{\textSFx{}}~~\lstbox{\textSFx{}}~~\lstbox{\textSFx{}}~~\lstbox{\textSFvii{}}~~\lstbox{\textSFx{}}~
\end{lstlisting}

\subsubsection*{How to fix it}

Move the attribute onto the class it was meant for, or remove it.

\errorcode{E3026}{custom attribute \errhole{} may only be applied to records}
\label{sec:codes:e3026}

Raised during \textbf{declaration}, from \texttt{DeclareErrorMessage::\allowbreak{}UNDEFINED\_\allowbreak{}ATTRIBUTE\_\allowbreak{}NO\_\allowbreak{}RECORD} (\textit{src/\allowbreak{}ymirc/\allowbreak{}semantic/\allowbreak{}declarator/\allowbreak{}errors.yr}).

\subsubsection*{What it means}

The attribute may only be applied to records, and it was written on something else.

\subsubsection*{Example}

\textit{test\_\allowbreak{}resources/\allowbreak{}diagnostics/\allowbreak{}test15.yr}:

%% check: skip
\begin{lstlisting}[style=coloredverbatimError]
@packed
entity Foo {
    let x: i32 = 0;
}

fn main () { let _ = Foo (); }
\end{lstlisting}

\begin{lstlisting}[style=bashVerb, breaklines=true, escapechar=~]
Error[E3026] : custom attribute packed may only be applied to records
 --> test_resources/diagnostics/test15.yr:(1,2)
 1  ~\lstbox{\textSFxi{}}~ @packed
    ~\lstbox{\textSFv{}}~  ^^^^^^
    ~\lstbox{\textSFxi{}}~ Note :
    ~\lstbox{\textSFxi{}}~  --> test_resources/diagnostics/test15.yr:(2,1)
    ~\lstbox{\textSFxi{}}~  2  ~\lstbox{\textSFxi{}}~ entity Foo {
    ~\lstbox{\textSFxi{}}~     ~\lstbox{\textSFv{}}~ ^^^^^^
    ~\lstbox{\textSFxi{}}~     ~\lstbox{\textSFii{}}~~\lstbox{\textSFx{}}~~\lstbox{\textSFx{}}~~\lstbox{\textSFx{}}~~\lstbox{\textSFx{}}~~\lstbox{\textSFx{}}~~\lstbox{\textSFx{}}~~\lstbox{\textSFx{}}~
    ~\lstbox{\textSFii{}}~~\lstbox{\textSFx{}}~~\lstbox{\textSFx{}}~~\lstbox{\textSFx{}}~~\lstbox{\textSFx{}}~~\lstbox{\textSFx{}}~~\lstbox{\textSFvii{}}~~\lstbox{\textSFx{}}~
\end{lstlisting}

\subsubsection*{How to fix it}

Move the attribute onto the record it was meant for, or remove it.

\errorcode{E3035}{a ctor cannot be both no\_move and only\_move}
\label{sec:codes:e3035}

Raised during \textbf{declaration}, from \texttt{DeclareErrorMessage::\allowbreak{}EXCLUSIVE\_\allowbreak{}MOVE\_\allowbreak{}ATTRIBUTES} (\textit{src/\allowbreak{}ymirc/\allowbreak{}semantic/\allowbreak{}declarator/\allowbreak{}errors.yr}).

\subsubsection*{What it means}

\tokennolst{@no\_move} withdraws a ctor from the resolution of the \tokennolst{move} operator, and \tokennolst{@only\_move} restricts it to that resolution. A ctor carrying both would be selectable from nowhere.

\subsubsection*{Example}

\textit{test\_\allowbreak{}resources/\allowbreak{}entities/\allowbreak{}test15.yr}:

%% check: skip
\begin{lstlisting}[style=coloredverbatimError]
entity Foo {
    @{no_move, only_move}
    pub self () {}
}
\end{lstlisting}

\begin{lstlisting}[style=bashVerb, breaklines=true, escapechar=~]
Error[E3035] : a ctor cannot be both no_move and only_move
 --> test_resources/entities/test15.yr:(4,16)
 4  ~\lstbox{\textSFxi{}}~     @{no_move, only_move}
    ~\lstbox{\textSFv{}}~                ^^^^^^^^^
    ~\lstbox{\textSFii{}}~~\lstbox{\textSFx{}}~~\lstbox{\textSFx{}}~~\lstbox{\textSFx{}}~~\lstbox{\textSFx{}}~~\lstbox{\textSFx{}}~~\lstbox{\textSFx{}}~~\lstbox{\textSFx{}}~
\end{lstlisting}

\subsubsection*{How to fix it}

Keep the one that says which side of the \tokennolst{move} operator the ctor belongs to, or neither -- a ctor with no move attribute is a candidate on both sides.

\errorcode{E3036}{custom attribute \errhole{} may only be applied to the ctor of an entity}
\label{sec:codes:e3036}

Raised during \textbf{declaration}, from \texttt{DeclareErrorMessage::\allowbreak{}MOVE\_\allowbreak{}ATTRIBUTE\_\allowbreak{}NO\_\allowbreak{}ENTITY} (\textit{src/\allowbreak{}ymirc/\allowbreak{}semantic/\allowbreak{}declarator/\allowbreak{}errors.yr}).

\subsubsection*{What it means}

\tokennolst{@no\_move} and \tokennolst{@only\_move} say which side of the \tokennolst{move} operator a ctor may be selected from. Only an entity is movable, so the attributes mean nothing on the ctor of a class or of a record.

\subsubsection*{Example}

\textit{test\_\allowbreak{}resources/\allowbreak{}entities/\allowbreak{}test16.yr}:

%% check: skip
\begin{lstlisting}[style=coloredverbatimError]
class Foo {
    @no_move
    pub self () {}
}
\end{lstlisting}

\begin{lstlisting}[style=bashVerb, breaklines=true, escapechar=~]
Error[E3036] : custom attribute no_move may only be applied to the ctor of an entity
 --> test_resources/entities/test16.yr:(4,6)
 4  ~\lstbox{\textSFxi{}}~     @no_move
    ~\lstbox{\textSFv{}}~      ^^^^^^^
    ~\lstbox{\textSFii{}}~~\lstbox{\textSFx{}}~~\lstbox{\textSFx{}}~~\lstbox{\textSFx{}}~~\lstbox{\textSFx{}}~~\lstbox{\textSFx{}}~~\lstbox{\textSFx{}}~~\lstbox{\textSFx{}}~
\end{lstlisting}

\subsubsection*{How to fix it}

Drop the attribute, or make the type an \tokennolst{entity} if it is meant to own a resource that is moved rather than copied.

\errorcode{E4041}{decorator is not applicable on inner types, already deeply mutable}
\label{sec:codes:e4041}

Raised during \textbf{validation}, from \texttt{ValidateErrorMessage::\allowbreak{}DEEPLY\_\allowbreak{}INNER\_\allowbreak{}TYPE} (\textit{src/\allowbreak{}ymirc/\allowbreak{}semantic/\allowbreak{}validator/\allowbreak{}errors.yr}).

\subsubsection*{What it means}

A mutability decorator was applied to the inner part of a type that is already deeply mutable, so it states something already true.

\subsubsection*{Example}

\textit{test\_\allowbreak{}resources/\allowbreak{}diagnostics/\allowbreak{}test99.yr}:

%% check: skip
\begin{lstlisting}[style=coloredverbatimError]
// a mut decorator inside a type that is already deeply mutable
fn main () {
    let mut _x_ : dmut [mut i32] = copy [1];
}
\end{lstlisting}

\begin{lstlisting}[style=bashVerb, breaklines=true, escapechar=~]
Error[E4041] : decorator is not applicable on inner types, already deeply mutable
 --> test_resources/diagnostics/test99.yr:(3,25)
 3  ~\lstbox{\textSFxi{}}~     let mut _x_ : dmut [mut i32] = copy [1];
    ~\lstbox{\textSFv{}}~                         ^^^
    ~\lstbox{\textSFxi{}}~
    = when validating test99::main -- (test_resources/diagnostics/test99.yr:(2,4))
    ~\lstbox{\textSFii{}}~~\lstbox{\textSFx{}}~~\lstbox{\textSFx{}}~~\lstbox{\textSFx{}}~~\lstbox{\textSFx{}}~~\lstbox{\textSFx{}}~~\lstbox{\textSFx{}}~~\lstbox{\textSFx{}}~
\end{lstlisting}

\subsubsection*{How to fix it}

Remove the decorator.

\errorcode{E4240}{decorator \errhole{} is not applicable in that context}
\label{sec:codes:e4240}

Raised during \textbf{validation}, from \texttt{ValidateErrorMessage::\allowbreak{}UNDEF\_\allowbreak{}DECORATOR\_\allowbreak{}HERE} (\textit{src/\allowbreak{}ymirc/\allowbreak{}semantic/\allowbreak{}validator/\allowbreak{}errors.yr}).

\subsubsection*{What it means}

A decorator is valid in general but not where it was written.

\subsubsection*{Example}

\textit{test\_\allowbreak{}resources/\allowbreak{}pattern\_\allowbreak{}match/\allowbreak{}test19.yr}:

%% check: skip
\begin{lstlisting}[style=coloredverbatimError]
fn main () {
    let (lazy x,) = (1,);
    x;
}
\end{lstlisting}

\begin{lstlisting}[style=bashVerb, breaklines=true, escapechar=~]
Error[E4240] : decorator lazy is not applicable in that context
 --> test_resources/pattern_match/test19.yr:(2,10)
 2  ~\lstbox{\textSFxi{}}~     let (lazy x,) = (1,);
    ~\lstbox{\textSFv{}}~          ^^^^
    ~\lstbox{\textSFxi{}}~ Note : for pattern field access 0 on type (mut i32,)
    ~\lstbox{\textSFxi{}}~  --> test_resources/pattern_match/test19.yr:(2,9)
    ~\lstbox{\textSFxi{}}~  2  ~\lstbox{\textSFxi{}}~     let (lazy x,) = (1,);
    ~\lstbox{\textSFxi{}}~     ~\lstbox{\textSFv{}}~         ^
    ~\lstbox{\textSFxi{}}~     ~\lstbox{\textSFii{}}~~\lstbox{\textSFx{}}~~\lstbox{\textSFx{}}~~\lstbox{\textSFx{}}~~\lstbox{\textSFx{}}~~\lstbox{\textSFx{}}~~\lstbox{\textSFx{}}~~\lstbox{\textSFx{}}~
    ~\lstbox{\textSFii{}}~~\lstbox{\textSFx{}}~~\lstbox{\textSFx{}}~~\lstbox{\textSFx{}}~~\lstbox{\textSFx{}}~~\lstbox{\textSFx{}}~~\lstbox{\textSFvii{}}~~\lstbox{\textSFx{}}~
\end{lstlisting}

\subsubsection*{How to fix it}

Move it onto something that accepts it.

%% Generated by tools/gen_error_codes.py from docs/errors/ of the compiler;
%% edit the compiler's pages or tools/error_themes.txt, then rerun it.

\clearpage
\section{Modules and visibility}
\label{sec:codes:modules}

The book's chapter \textit{Global constructions} specifies the rules behind these codes.

\begin{longtable}{@{}l p{0.78\linewidth} r@{}}
\toprule Code & Message & Page\\ \midrule \endhead
\hyperref[sec:codes:e3002]{E3002} & module is ambiguously defined in both \errhole{} and \errhole{} & \pageref{sec:codes:e3002}\\
\hyperref[sec:codes:e3009]{E3009} & core module is malformed & \pageref{sec:codes:e3009}\\
\hyperref[sec:codes:e3012]{E3012} & no module named \errhole{} found inside module \errhole{} & \pageref{sec:codes:e3012}\\
\hyperref[sec:codes:e3013]{E3013} & module should be located in \errhole{} or \errhole{}, but no such file exists or access is denied & \pageref{sec:codes:e3013}\\
\hyperref[sec:codes:e3014]{E3014} & module should be located in \errhole{}, but no such file exists or access is denied & \pageref{sec:codes:e3014}\\
\hyperref[sec:codes:e3016]{E3016} & \errhole{} is private in this context & \pageref{sec:codes:e3016}\\
\hyperref[sec:codes:e3017]{E3017} & the \errhole{} module is imported implicitly; an explicit use is unnecessary & \pageref{sec:codes:e3017}\\
\hyperref[sec:codes:e3029]{E3029} & unexpected declaration inside an extern block & \pageref{sec:codes:e3029}\\
\hyperref[sec:codes:e3034]{E3034} & module \errhole{} must be defined in a file named \errhole{} & \pageref{sec:codes:e3034}\\
\hyperref[sec:codes:e3038]{E3038} & no symbol is resolved through the use of \errhole{} & \pageref{sec:codes:e3038}\\
\hyperref[sec:codes:e3039]{E3039} & only \errhole{} is resolved through the use of \errhole{} & \pageref{sec:codes:e3039}\\
\hyperref[sec:codes:e4116]{E4116} & core module is malformed & \pageref{sec:codes:e4116}\\
\hyperref[sec:codes:e4282]{E4282} & not supported in version ymir \errhole{} & \pageref{sec:codes:e4282}\\
\bottomrule
\end{longtable}

\errorcode{E3002}{module is ambiguously defined in both \errhole{} and \errhole{}}
\label{sec:codes:e3002}

Raised during \textbf{declaration}, from \texttt{DeclareErrorMessage::\allowbreak{}CONFLICT\_\allowbreak{}MODULES} (\textit{src/\allowbreak{}ymirc/\allowbreak{}semantic/\allowbreak{}declarator/\allowbreak{}errors.yr}).

\subsubsection*{What it means}

Two files claim to define the same module, so an import of that module has no single answer. Both paths are reported.

\subsubsection*{Example}

\textit{test\_\allowbreak{}resources/\allowbreak{}diagnostics/\allowbreak{}test93.yr}:

%% check: skip
\begin{lstlisting}[style=coloredverbatimError]
// a module found both next to the importing file and in its directory
mod dup;

fn main () {}
\end{lstlisting}

\begin{lstlisting}[style=bashVerb, breaklines=true, escapechar=~]
Error[E3002] : module is ambiguously defined in both /bootstrap/test_resources/diagnostics/test93/dup.yr and /bootstrap/test_resources/diagnostics/dup.yr
 --> test_resources/diagnostics/test93.yr:(2,5)
 2  ~\lstbox{\textSFxi{}}~ mod dup;
    ~\lstbox{\textSFv{}}~     ^^^
    ~\lstbox{\textSFii{}}~~\lstbox{\textSFx{}}~~\lstbox{\textSFx{}}~~\lstbox{\textSFx{}}~~\lstbox{\textSFx{}}~~\lstbox{\textSFx{}}~~\lstbox{\textSFx{}}~~\lstbox{\textSFx{}}~
\end{lstlisting}

\subsubsection*{How to fix it}

Remove or rename one of them. This most often comes from the same package being reachable twice through the include path, once directly and once through a dependency.

\errorcode{E3009}{core module is malformed}
\label{sec:codes:e3009}

Raised during \textbf{declaration}, from \texttt{DeclareErrorMessage::\allowbreak{}MALFORMED\_\allowbreak{}CORE} (\textit{src/\allowbreak{}ymirc/\allowbreak{}semantic/\allowbreak{}declarator/\allowbreak{}errors.yr}).

\subsubsection*{What it means}

The \tokennolst{core} module the compiler needs does not have the shape it expects: a symbol the compiler relies on is missing or has the wrong kind. This is an installation problem rather than a problem with the code being compiled.

\subsubsection*{How to fix it}

Check that the standard library found on the include path matches the compiler version. A \tokennolst{core} from another release is the usual cause.

\errorcode{E3012}{no module named \errhole{} found inside module \errhole{}}
\label{sec:codes:e3012}

Raised during \textbf{declaration}, from \texttt{DeclareErrorMessage::\allowbreak{}NO\_\allowbreak{}MODULE\_\allowbreak{}NAMED} (\textit{src/\allowbreak{}ymirc/\allowbreak{}semantic/\allowbreak{}declarator/\allowbreak{}errors.yr}).

\subsubsection*{What it means}

A \tokennolst{use} names a module that does not exist inside the module it was looked up in. Both names are reported.

\subsubsection*{Example}

\textit{test\_\allowbreak{}resources/\allowbreak{}diagnostics/\allowbreak{}test7.yr}:

%% check: skip
\begin{lstlisting}[style=coloredverbatimError]
use std::nosuchmodulehere;

fn main () {}
\end{lstlisting}

\begin{lstlisting}[style=bashVerb, breaklines=true, escapechar=~]
Error[E3012] : no module named nosuchmodulehere found inside module std
 --> test_resources/diagnostics/test7.yr:(1,10)
 1  ~\lstbox{\textSFxi{}}~ use std::nosuchmodulehere;
    ~\lstbox{\textSFv{}}~          ^^^^^^^^^^^^^^^^
    ~\lstbox{\textSFii{}}~~\lstbox{\textSFx{}}~~\lstbox{\textSFx{}}~~\lstbox{\textSFx{}}~~\lstbox{\textSFx{}}~~\lstbox{\textSFx{}}~~\lstbox{\textSFx{}}~~\lstbox{\textSFx{}}~
\end{lstlisting}

\subsubsection*{How to fix it}

Check the spelling and the path of the import. A module that exists on disk but is not reachable this way is usually a package-root problem: the include path has to contain the directory the package root sits in.

\errorcode{E3013}{module should be located in \errhole{} or \errhole{}, but no such file exists or access is denied}
\label{sec:codes:e3013}

Raised during \textbf{declaration}, from \texttt{DeclareErrorMessage::\allowbreak{}NO\_\allowbreak{}SUCH\_\allowbreak{}FILE} (\textit{src/\allowbreak{}ymirc/\allowbreak{}semantic/\allowbreak{}declarator/\allowbreak{}errors.yr}).

\subsubsection*{What it means}

A module was imported but no file defines it. The compiler reports both file paths it would have accepted -- the directory form and the plain file form.

\subsubsection*{Example}

\textit{test\_\allowbreak{}resources/\allowbreak{}diagnostics/\allowbreak{}test94.yr}:

%% check: skip
\begin{lstlisting}[style=coloredverbatimError]
// a module found neither next to the importing file nor in its directory
mod nothere;

fn main () {}
\end{lstlisting}

\begin{lstlisting}[style=bashVerb, breaklines=true, escapechar=~]
Error[E3013] : module should be located in /bootstrap/test_resources/diagnostics/test94/nothere.yr or /bootstrap/test_resources/diagnostics/nothere.yr, but no such file exists or access is denied
 --> test_resources/diagnostics/test94.yr:(2,5)
 2  ~\lstbox{\textSFxi{}}~ mod nothere;
    ~\lstbox{\textSFv{}}~     ^^^^^^^
    ~\lstbox{\textSFii{}}~~\lstbox{\textSFx{}}~~\lstbox{\textSFx{}}~~\lstbox{\textSFx{}}~~\lstbox{\textSFx{}}~~\lstbox{\textSFx{}}~~\lstbox{\textSFx{}}~~\lstbox{\textSFx{}}~
\end{lstlisting}

\subsubsection*{How to fix it}

Create the file at one of the two reported paths, or fix the import. An existing file at one of them points at a permission problem instead.

\errorcode{E3014}{module should be located in \errhole{}, but no such file exists or access is denied}
\label{sec:codes:e3014}

Raised during \textbf{declaration}, from \texttt{DeclareErrorMessage::\allowbreak{}NO\_\allowbreak{}SUCH\_\allowbreak{}FILE\_\allowbreak{}ONE} (\textit{src/\allowbreak{}ymirc/\allowbreak{}semantic/\allowbreak{}declarator/\allowbreak{}errors.yr}).

\subsubsection*{What it means}

A module was imported but the single file that could define it does not exist or cannot be read.

\subsubsection*{Example}

\textit{test\_\allowbreak{}resources/\allowbreak{}diagnostics/\allowbreak{}test95.yr}:

%% check: skip
\begin{lstlisting}[style=coloredverbatimError]
// a child module with no file in the directory of the importing module
mod ::nothere;

fn main () {}
\end{lstlisting}

\begin{lstlisting}[style=bashVerb, breaklines=true, escapechar=~]
Error[E3014] : module should be located in /bootstrap/test_resources/diagnostics/test95/nothere.yr, but no such file exists or access is denied
 --> test_resources/diagnostics/test95.yr:(2,5)
 2  ~\lstbox{\textSFxi{}}~ mod ::nothere;
    ~\lstbox{\textSFv{}}~     ^^-------
    ~\lstbox{\textSFii{}}~~\lstbox{\textSFx{}}~~\lstbox{\textSFx{}}~~\lstbox{\textSFx{}}~~\lstbox{\textSFx{}}~~\lstbox{\textSFx{}}~~\lstbox{\textSFx{}}~~\lstbox{\textSFx{}}~
\end{lstlisting}

\subsubsection*{How to fix it}

Same as \hyperref[sec:codes:e3013]{E3013}, with only one candidate path.

\errorcode{E3016}{\errhole{} is private in this context}
\label{sec:codes:e3016}

Raised during \textbf{declaration}, from \texttt{DeclareErrorMessage::\allowbreak{}PRIVATE\_\allowbreak{}IN\_\allowbreak{}THIS\_\allowbreak{}CONTEXT} (\textit{src/\allowbreak{}ymirc/\allowbreak{}semantic/\allowbreak{}declarator/\allowbreak{}errors.yr}).

\subsubsection*{What it means}

The symbol exists but its protection does not let this module see it.

\subsubsection*{Example}

\textit{test\_\allowbreak{}resources/\allowbreak{}data\_\allowbreak{}record/\allowbreak{}fields/\allowbreak{}test2.yr}:

%% check: skip
\begin{lstlisting}[style=coloredverbatimError]
@data
record Foo {
    let x: i32;
    prv let _y: i32 = 8;
}

fn main() {
    let a = Foo(1);
    a._y;
}
\end{lstlisting}

\begin{lstlisting}[style=bashVerb, breaklines=true, escapechar=~]
Error[E3016] : _y is private in this context
 --> test_resources/data_record/fields/test2.yr:(9,7)
 9  ~\lstbox{\textSFxi{}}~     a._y;
    ~\lstbox{\textSFv{}}~       ^^  note: class test2::Foo has no field access method named _y
    ~\lstbox{\textSFxi{}}~        ~\lstbox{\textSFii{}}~~\lstbox{\textSFx{}}~ note: class test2::Foo has no method named _y
    ~\lstbox{\textSFxi{}}~ Note : declared here
    ~\lstbox{\textSFxi{}}~  --> test_resources/data_record/fields/test2.yr:(4,13)
    ~\lstbox{\textSFxi{}}~  4  ~\lstbox{\textSFxi{}}~     prv let _y: i32 = 8;
    ~\lstbox{\textSFxi{}}~     ~\lstbox{\textSFv{}}~             ^^
    ~\lstbox{\textSFxi{}}~     ~\lstbox{\textSFii{}}~~\lstbox{\textSFx{}}~~\lstbox{\textSFx{}}~~\lstbox{\textSFx{}}~~\lstbox{\textSFx{}}~~\lstbox{\textSFx{}}~~\lstbox{\textSFx{}}~~\lstbox{\textSFx{}}~
    ~\lstbox{\textSFii{}}~~\lstbox{\textSFx{}}~~\lstbox{\textSFx{}}~~\lstbox{\textSFx{}}~~\lstbox{\textSFx{}}~~\lstbox{\textSFx{}}~~\lstbox{\textSFvii{}}~~\lstbox{\textSFx{}}~
\end{lstlisting}

\subsubsection*{How to fix it}

Import it from a module that may see it, or widen its protection at the definition. Inside a package, \tokennolst{prot} (protected) makes a symbol visible to the package without exporting it.

\errorcode{E3017}{the \errhole{} module is imported implicitly; an explicit use is unnecessary}
\label{sec:codes:e3017}

Raised during \textbf{declaration}, from \texttt{DeclareErrorMessage::\allowbreak{}SELF\_\allowbreak{}USE} (\textit{src/\allowbreak{}ymirc/\allowbreak{}semantic/\allowbreak{}declarator/\allowbreak{}errors.yr}).

\subsubsection*{What it means}

A module imports itself, directly or through its own package root. That import is always implicit, so writing it does nothing.

\subsubsection*{Example}

\textit{test\_\allowbreak{}resources/\allowbreak{}warnings/\allowbreak{}test2.yr}:

%% check: skip
\begin{lstlisting}[style=coloredverbatimError]
// a module importing itself only warns in debug mode
use test2;

fn main () {
}
\end{lstlisting}

\begin{lstlisting}[style=bashVerb, breaklines=true, escapechar=~]
Warning[E3017] : the test2 module is imported implicitly; an explicit use is unnecessary
 --> test_resources/warnings/test2.yr:(2,5)
 2  ~\lstbox{\textSFxi{}}~ use test2;
    ~\lstbox{\textSFv{}}~     ^^^^^
    ~\lstbox{\textSFii{}}~~\lstbox{\textSFx{}}~~\lstbox{\textSFx{}}~~\lstbox{\textSFx{}}~~\lstbox{\textSFx{}}~~\lstbox{\textSFx{}}~~\lstbox{\textSFx{}}~~\lstbox{\textSFx{}}~
\end{lstlisting}

\subsubsection*{How to fix it}

Remove the \tokennolst{use}.

\errorcode{E3029}{unexpected declaration inside an extern block}
\label{sec:codes:e3029}

Raised during \textbf{declaration}, from \texttt{DeclareErrorMessage::\allowbreak{}UNEXPECTED\_\allowbreak{}IN\_\allowbreak{}EXTERN} (\textit{src/\allowbreak{}ymirc/\allowbreak{}semantic/\allowbreak{}declarator/\allowbreak{}errors.yr}).

\subsubsection*{What it means}

A declaration appeared inside an \tokennolst{extern} block that such a block does not accept. An \tokennolst{extern} block declares symbols defined elsewhere, so it takes prototypes and variable declarations, not definitions.

\subsubsection*{Example}

\textit{test\_\allowbreak{}resources/\allowbreak{}diagnostics/\allowbreak{}test16.yr}:

%% check: skip
\begin{lstlisting}[style=coloredverbatimError]
extern (C) {
    class Foo {
        pub self () {}
    }
}

fn main () {}
\end{lstlisting}

\begin{lstlisting}[style=bashVerb, breaklines=true, escapechar=~]
Error[E3029] : unexpected declaration inside an extern block
 --> test_resources/diagnostics/test16.yr:(2,5)
 2  ~\lstbox{\textSFxi{}}~     class Foo {
    ~\lstbox{\textSFv{}}~     ^^^^^
    ~\lstbox{\textSFii{}}~~\lstbox{\textSFx{}}~~\lstbox{\textSFx{}}~~\lstbox{\textSFx{}}~~\lstbox{\textSFx{}}~~\lstbox{\textSFx{}}~~\lstbox{\textSFx{}}~~\lstbox{\textSFx{}}~
\end{lstlisting}

\subsubsection*{How to fix it}

Move the declaration out of the block.

\errorcode{E3034}{module \errhole{} must be defined in a file named \errhole{}}
\label{sec:codes:e3034}

Raised during \textbf{declaration}, from \texttt{DeclareErrorMessage::\allowbreak{}WRONG\_\allowbreak{}MODULE\_\allowbreak{}NAME} (\textit{src/\allowbreak{}ymirc/\allowbreak{}semantic/\allowbreak{}declarator/\allowbreak{}errors.yr}).

\subsubsection*{What it means}

The module name declared by the \tokennolst{in}/\tokennolst{mod} header does not match the file it is written in. Module names are resolved to paths, so the two have to agree for an import to find the file.

\subsubsection*{Example}

\textit{test\_\allowbreak{}resources/\allowbreak{}diagnostics/\allowbreak{}test17.yr}:

%% check: skip
\begin{lstlisting}[style=coloredverbatimError]
in someothername;

fn main () {}
\end{lstlisting}

\begin{lstlisting}[style=bashVerb, breaklines=true, escapechar=~]
Error[E3034] : module someothername --> /home/emile/ymir/gcc/gcc-src/gcc/ymir/bootstrap/test_resources/diagnostics/test17.yr:(1,4) must be defined in a file named someothername.yr
 --> test_resources/diagnostics/test17.yr:(1,4)
 1  ~\lstbox{\textSFxi{}}~ in someothername;
    ~\lstbox{\textSFv{}}~    ^^^^^^^^^^^^^
    ~\lstbox{\textSFii{}}~~\lstbox{\textSFx{}}~~\lstbox{\textSFx{}}~~\lstbox{\textSFx{}}~~\lstbox{\textSFx{}}~~\lstbox{\textSFx{}}~~\lstbox{\textSFx{}}~~\lstbox{\textSFx{}}~
\end{lstlisting}

\subsubsection*{How to fix it}

Rename the file to the reported name, or change the declared module name to match the file.

\errorcode{E3038}{no symbol is resolved through the use of \errhole{}}
\label{sec:codes:e3038}

Raised from \tokennolst{DeclareErrorMessage::UNUSED\_USE} (\textit{src/\allowbreak{}ymirc/\allowbreak{}semantic/\allowbreak{}declarator/\allowbreak{}errors.yr}), by the check that runs once the whole package is validated.

\subsubsection*{What it means}

\tokennolst{use} does not decide what a package can reach -- package resolution and \tokennolst{mod} do that, and every symbol is in the tree from the start, subject to protection. A \tokennolst{use} only lets a symbol be named by its short name instead of its full path.

So this is redundant-shorthand detection, and ``used'' has an exact meaning: a resolution \textit{selected} a symbol that this \tokennolst{use}, and nothing else, made nameable. Two consequences follow, and both are the rule rather than an approximation of it:

\begin{itemize}
\item a candidate that lost overload resolution contributed nothing to the compiled program, so the \tokennolst{use} that offered it is still unused;
\item a symbol the scope could name anyway -- written out as a full path, supplied by another \tokennolst{use}, declared in an enclosing scope, or provided by a core module -- was never owed to this \tokennolst{use} either.
\end{itemize}

\subsubsection*{Example}

\textit{test\_\allowbreak{}resources/\allowbreak{}warnings/\allowbreak{}test9.yr}, where \tokennolst{pick} is declared in both imported modules and the call selects the one taking an \tokennolst{i32}:

%% check: skip
\begin{lstlisting}[style=coloredverbatimError]
in test9;

mod ::alpha;
mod ::beta;

use test9::alpha;
use test9::beta;

fn main () {
    pick (1);
}
\end{lstlisting}

\begin{lstlisting}[style=bashVerb, breaklines=true, escapechar=~]
Warning[E3038] : no symbol is resolved through the use of test9::beta
 --> test_resources/warnings/test9.yr:(10,12)
10  ~\lstbox{\textSFxi{}}~ use test9::beta;
    ~\lstbox{\textSFv{}}~            ^^^^
    ~\lstbox{\textSFii{}}~~\lstbox{\textSFx{}}~~\lstbox{\textSFx{}}~~\lstbox{\textSFx{}}~~\lstbox{\textSFx{}}~~\lstbox{\textSFx{}}~~\lstbox{\textSFx{}}~~\lstbox{\textSFx{}}~
\end{lstlisting}

\subsubsection*{How to fix it}

Remove the \tokennolst{use}. Nothing it names is reached, so nothing else changes.

\errorcode{E3039}{only \errhole{} is resolved through the use of \errhole{}}
\label{sec:codes:e3039}

Raised from \tokennolst{DeclareErrorMessage::SINGLE\_WILDCARD\_USE} (\textit{src/\allowbreak{}ymirc/\allowbreak{}semantic/\allowbreak{}declarator/\allowbreak{}errors.yr}), by the check that runs once the whole package is validated.

\subsubsection*{What it means}

A path ending with \tokennolst{\_} makes every submodule of the module it names visible at once. When a single one of them turns out to be resolved through, the \tokennolst{use} reaches much further than it needs to: every name the other modules declare takes part in the resolutions of the scope, and can win one.

A \tokennolst{use} nothing at all is resolved through is reported as \hyperref[sec:codes:e3038]{E3038} instead.

\subsubsection*{Example}

\textit{test\_\allowbreak{}resources/\allowbreak{}warnings/\allowbreak{}test6.yr}, where \tokennolst{group} holds \tokennolst{first} and \tokennolst{second}:

%% check: skip
\begin{lstlisting}[style=coloredverbatimError]
in test6;

pub mod ::group;

use test6::group::_;

fn main () {
    firstValue ();
}
\end{lstlisting}

\begin{lstlisting}[style=bashVerb, breaklines=true, escapechar=~]
Warning[E3039] : only test6::group::first is resolved through the use of test6::group::_
 --> test_resources/warnings/test6.yr:(7,19)
 7  ~\lstbox{\textSFxi{}}~ use test6::group::_;
    ~\lstbox{\textSFv{}}~                   ^  note: name it on its own, with use test6::group::first;
    ~\lstbox{\textSFii{}}~~\lstbox{\textSFx{}}~~\lstbox{\textSFx{}}~~\lstbox{\textSFx{}}~~\lstbox{\textSFx{}}~~\lstbox{\textSFx{}}~~\lstbox{\textSFx{}}~~\lstbox{\textSFx{}}~
\end{lstlisting}

\subsubsection*{How to fix it}

Name the module the scope actually resolves through, as the note spells it out.

\errorcode{E4116}{core module is malformed}
\label{sec:codes:e4116}

Raised during \textbf{validation}, from \texttt{ValidateErrorMessage::\allowbreak{}MALFORMED\_\allowbreak{}CORE} (\textit{src/\allowbreak{}ymirc/\allowbreak{}semantic/\allowbreak{}validator/\allowbreak{}errors.yr}).

\subsubsection*{What it means}

The \tokennolst{core} module does not have the shape the validator expects, cf. \hyperref[sec:codes:e3009]{E3009}.

\subsubsection*{How to fix it}

Check that the standard library matches the compiler version.

\errorcode{E4282}{not supported in version ymir \errhole{}}
\label{sec:codes:e4282}

Raised during \textbf{validation}, from \texttt{ValidateErrorMessage::\allowbreak{}NOT\_\allowbreak{}SUPPORTED\_\allowbreak{}IN} (\textit{src/\allowbreak{}ymirc/\allowbreak{}semantic/\allowbreak{}validator/\allowbreak{}errors.yr}).

\subsubsection*{What it means}

The construct used is not supported by this version of the compiler. The version is in the message.

\tokennolst{await} is what raises it today: the expression parses, but nothing validates it.

\subsubsection*{Example}

\textit{test\_\allowbreak{}resources/\allowbreak{}syntax/\allowbreak{}test15.yr}:

%% check: skip
\begin{lstlisting}[style=coloredverbatimError]
fn main () {
    let th = spawn compute (10);
    let _aw_ = await th;
}
\end{lstlisting}

\begin{lstlisting}[style=bashVerb, breaklines=true, escapechar=~]
Error[E4282] : not supported in version ymir 1.6
 --> test_resources/syntax/test15.yr:(27,16)
27  ~\lstbox{\textSFxi{}}~     let _aw_ = await th;
    ~\lstbox{\textSFv{}}~                ^^^^^
    ~\lstbox{\textSFxi{}}~
    = when validating test15::main -- (test_resources/syntax/test15.yr:(25,4))
    ~\lstbox{\textSFii{}}~~\lstbox{\textSFx{}}~~\lstbox{\textSFx{}}~~\lstbox{\textSFx{}}~~\lstbox{\textSFx{}}~~\lstbox{\textSFx{}}~~\lstbox{\textSFx{}}~~\lstbox{\textSFx{}}~
\end{lstlisting}

\subsubsection*{How to fix it}

Use a construct this version supports, or compile with a version that implements it.

%% Generated by tools/gen_error_codes.py from docs/errors/ of the compiler;
%% edit the compiler's pages or tools/error_themes.txt, then rerun it.

\clearpage
\section{Symbols and variables}
\label{sec:codes:symbols}

The book's chapter \textit{Global constructions} specifies the rules behind these codes.

\begin{longtable}{@{}l p{0.78\linewidth} r@{}}
\toprule Code & Message & Page\\ \midrule \endhead
\hyperref[sec:codes:e3003]{E3003} & only external global variables may be declared static & \pageref{sec:codes:e3003}\\
\hyperref[sec:codes:e3004]{E3004} & only local global variables may be declared thread-local & \pageref{sec:codes:e3004}\\
\hyperref[sec:codes:e3018]{E3018} & declaration of \errhole{} shadows a previous declaration & \pageref{sec:codes:e3018}\\
\hyperref[sec:codes:e3027]{E3027} & undefined symbol \errhole{} & \pageref{sec:codes:e3027}\\
\hyperref[sec:codes:e4001]{E4001} & a name alias cannot be created to alias class ctors & \pageref{sec:codes:e4001}\\
\hyperref[sec:codes:e4038]{E4038} & the symbol \errhole{} was declared but never used & \pageref{sec:codes:e4038}\\
\hyperref[sec:codes:e4039]{E4039} & cannot declare a variable of type \errhole{} & \pageref{sec:codes:e4039}\\
\hyperref[sec:codes:e4058]{E4058} & external global variable cannot be declared with a value & \pageref{sec:codes:e4058}\\
\hyperref[sec:codes:e4068]{E4068} & the name alias \errhole{} cannot be validated due to forward reference & \pageref{sec:codes:e4068}\\
\hyperref[sec:codes:e4069]{E4069} & the type cannot be validated due to forward reference & \pageref{sec:codes:e4069}\\
\hyperref[sec:codes:e4080]{E4080} & local global variable must be declared with a value & \pageref{sec:codes:e4080}\\
\hyperref[sec:codes:e4104]{E4104} & symbol \errhole{} has errors & \pageref{sec:codes:e4104}\\
\hyperref[sec:codes:e4106]{E4106} & the identifier \errhole{} describes a native type but has no meaning within this context & \pageref{sec:codes:e4106}\\
\hyperref[sec:codes:e4135]{E4135} & expression refers to multiple symbols that generate types & \pageref{sec:codes:e4135}\\
\hyperref[sec:codes:e4199]{E4199} & declaration of \errhole{} shadows another declaration & \pageref{sec:codes:e4199}\\
\hyperref[sec:codes:e4281]{E4281} & var declaration \errhole{} outside a function context & \pageref{sec:codes:e4281}\\
\hyperref[sec:codes:e4290]{E4290} & none of the symbols \errhole{} can be used as a value & \pageref{sec:codes:e4290}\\
\bottomrule
\end{longtable}

\errorcode{E3003}{only external global variables may be declared static}
\label{sec:codes:e3003}

Raised during \textbf{declaration}, from \texttt{DeclareErrorMessage::\allowbreak{}DECL\_\allowbreak{}GLOB\_\allowbreak{}STATIC\_\allowbreak{}NON\_\allowbreak{}EXTERN} (\textit{src/\allowbreak{}ymirc/\allowbreak{}semantic/\allowbreak{}declarator/\allowbreak{}errors.yr}).

\subsubsection*{What it means}

A global variable is declared \tokennolst{static} without being \tokennolst{extern}. \tokennolst{static} says the variable is defined elsewhere and merely referenced here, which only makes sense for an external symbol.

\subsubsection*{Example}

\textit{test\_\allowbreak{}resources/\allowbreak{}diagnostics/\allowbreak{}test5.yr}:

%% check: skip
\begin{lstlisting}[style=coloredverbatimError]
static X : i32;

fn main () { X; }
\end{lstlisting}

\begin{lstlisting}[style=bashVerb, breaklines=true, escapechar=~]
Error[E3003] : only external global variables may be declared static
 --> test_resources/diagnostics/test5.yr:(1,1)
 1  ~\lstbox{\textSFxi{}}~ static X : i32;
    ~\lstbox{\textSFv{}}~ ^^^^^^
    ~\lstbox{\textSFii{}}~~\lstbox{\textSFx{}}~~\lstbox{\textSFx{}}~~\lstbox{\textSFx{}}~~\lstbox{\textSFx{}}~~\lstbox{\textSFx{}}~~\lstbox{\textSFx{}}~~\lstbox{\textSFx{}}~
\end{lstlisting}

\subsubsection*{How to fix it}

Drop \tokennolst{static} for a variable this module defines, or declare it \tokennolst{extern} if it is defined by another object file.

\errorcode{E3004}{only local global variables may be declared thread-local}
\label{sec:codes:e3004}

Raised during \textbf{declaration}, from \texttt{DeclareErrorMessage::\allowbreak{}DECL\_\allowbreak{}GLOB\_\allowbreak{}THREAD\_\allowbreak{}EXTERN} (\textit{src/\allowbreak{}ymirc/\allowbreak{}semantic/\allowbreak{}declarator/\allowbreak{}errors.yr}).

\subsubsection*{What it means}

A global variable is declared both thread-local and \tokennolst{extern}. Thread-local storage is allocated by the module that defines the variable, so an external declaration cannot ask for it.

\subsubsection*{Example}

\textit{test\_\allowbreak{}resources/\allowbreak{}global/\allowbreak{}test10.yr}:

%% check: skip
\begin{lstlisting}[style=coloredverbatimError]
@thread
extern lazy i : i32;
\end{lstlisting}

\begin{lstlisting}[style=bashVerb, breaklines=true, escapechar=~]
Error[E3004] : only local global variables may be declared thread-local
 --> test_resources/global/test10.yr:(1,2)
 1  ~\lstbox{\textSFxi{}}~ @thread
    ~\lstbox{\textSFv{}}~  ^^^^^^
    ~\lstbox{\textSFxi{}}~ Note :
    ~\lstbox{\textSFxi{}}~  --> test_resources/global/test10.yr:(2,8)
    ~\lstbox{\textSFxi{}}~  2  ~\lstbox{\textSFxi{}}~ extern lazy i : i32;
    ~\lstbox{\textSFxi{}}~     ~\lstbox{\textSFv{}}~        ^^^^
    ~\lstbox{\textSFxi{}}~     ~\lstbox{\textSFii{}}~~\lstbox{\textSFx{}}~~\lstbox{\textSFx{}}~~\lstbox{\textSFx{}}~~\lstbox{\textSFx{}}~~\lstbox{\textSFx{}}~~\lstbox{\textSFx{}}~~\lstbox{\textSFx{}}~
    ~\lstbox{\textSFii{}}~~\lstbox{\textSFx{}}~~\lstbox{\textSFx{}}~~\lstbox{\textSFx{}}~~\lstbox{\textSFx{}}~~\lstbox{\textSFx{}}~~\lstbox{\textSFvii{}}~~\lstbox{\textSFx{}}~
\end{lstlisting}

\subsubsection*{How to fix it}

Drop the thread-local decorator on the \tokennolst{extern} declaration; it belongs on the definition.

\errorcode{E3018}{declaration of \errhole{} shadows a previous declaration}
\label{sec:codes:e3018}

Raised during \textbf{declaration}, from \texttt{DeclareErrorMessage::\allowbreak{}SHADOWING\_\allowbreak{}DECL} (\textit{src/\allowbreak{}ymirc/\allowbreak{}semantic/\allowbreak{}declarator/\allowbreak{}errors.yr}).

\subsubsection*{What it means}

A declaration hides an earlier one with the same name that is still in scope. Ymir rejects the shadowing rather than picking one, so a reader never has to work out which of two same-named symbols a use refers to.

\subsubsection*{Example}

\textit{test\_\allowbreak{}resources/\allowbreak{}unittest/\allowbreak{}test2.yr}:

%% check: skip
\begin{lstlisting}[style=coloredverbatimError]
__test foo {
    assert(false, "Test foo");
}
__test foo {
    assert(false, "Test foo");
}
\end{lstlisting}

\begin{lstlisting}[style=bashVerb, breaklines=true, escapechar=~]
Error[E3018] : declaration of foo shadows a previous declaration
 --> test_resources/unittest/test2.yr:(4,8)
 4  ~\lstbox{\textSFxi{}}~ __test foo {
    ~\lstbox{\textSFv{}}~        ^^^
    ~\lstbox{\textSFxi{}}~ Note :
    ~\lstbox{\textSFxi{}}~  --> test_resources/unittest/test2.yr:(1,8)
    ~\lstbox{\textSFxi{}}~  1  ~\lstbox{\textSFxi{}}~ __test foo {
    ~\lstbox{\textSFxi{}}~     ~\lstbox{\textSFv{}}~        ^^^
    ~\lstbox{\textSFxi{}}~     ~\lstbox{\textSFii{}}~~\lstbox{\textSFx{}}~~\lstbox{\textSFx{}}~~\lstbox{\textSFx{}}~~\lstbox{\textSFx{}}~~\lstbox{\textSFx{}}~~\lstbox{\textSFx{}}~~\lstbox{\textSFx{}}~
    ~\lstbox{\textSFii{}}~~\lstbox{\textSFx{}}~~\lstbox{\textSFx{}}~~\lstbox{\textSFx{}}~~\lstbox{\textSFx{}}~~\lstbox{\textSFx{}}~~\lstbox{\textSFvii{}}~~\lstbox{\textSFx{}}~
\end{lstlisting}

\subsubsection*{How to fix it}

Rename one of the two. When the outer one is no longer needed at that point, moving its declaration into a narrower block also resolves it.

\errorcode{E3027}{undefined symbol \errhole{}}
\label{sec:codes:e3027}

Raised during \textbf{declaration}, from \texttt{DeclareErrorMessage::\allowbreak{}UNDEF\_\allowbreak{}VAR} (\textit{src/\allowbreak{}ymirc/\allowbreak{}semantic/\allowbreak{}declarator/\allowbreak{}errors.yr}).

\subsubsection*{What it means}

A name used in a declaration resolves to nothing: no symbol of that name is visible from this module.

\subsubsection*{Example}

\textit{test\_\allowbreak{}resources/\allowbreak{}regressions/\allowbreak{}use\_\allowbreak{}is\_\allowbreak{}file\_\allowbreak{}local/\allowbreak{}test1/\allowbreak{}foo.yr}:

%% check: skip
\begin{lstlisting}[style=coloredverbatimError]
in foo;
use std::io;

mod bar;
\end{lstlisting}

\begin{lstlisting}[style=bashVerb, breaklines=true, escapechar=~]
Error[E3027] : undefined symbol println
 --> test_resources/regressions/use_is_file_local/test1/foo/bar.yr:(5,5)
 5  ~\lstbox{\textSFxi{}}~     println ("in bar");
    ~\lstbox{\textSFv{}}~     ^^^^^^^
    ~\lstbox{\textSFii{}}~~\lstbox{\textSFx{}}~~\lstbox{\textSFx{}}~~\lstbox{\textSFx{}}~~\lstbox{\textSFx{}}~~\lstbox{\textSFx{}}~~\lstbox{\textSFx{}}~~\lstbox{\textSFx{}}~
\end{lstlisting}

\subsubsection*{How to fix it}

Check the spelling, and check that the module defining the symbol is imported. A symbol that exists but is not visible gives \hyperref[sec:codes:e3016]{E3016} instead, which distinguishes a missing import from a protection problem.

\errorcode{E4001}{a name alias cannot be created to alias class ctors}
\label{sec:codes:e4001}

Raised during \textbf{validation}, from \texttt{ValidateErrorMessage::\allowbreak{}AKA\_\allowbreak{}CLASS\_\allowbreak{}CTOR} (\textit{src/\allowbreak{}ymirc/\allowbreak{}semantic/\allowbreak{}validator/\allowbreak{}errors.yr}).

\subsubsection*{What it means}

A name alias (\tokennolst{def}) was made to point at the constructors of a class. A constructor is not a symbol that can be referred to on its own: it is reached through the class it belongs to.

\subsubsection*{Example}

\textit{test\_\allowbreak{}resources/\allowbreak{}aka/\allowbreak{}test14.yr}:

%% check: skip
\begin{lstlisting}[style=coloredverbatimError]
class A {
    pub self new () {}
}

def B = A::new;
def C = A;
\end{lstlisting}

\begin{lstlisting}[style=bashVerb, breaklines=true, escapechar=~]
Error[E4001] : a name alias cannot be created to alias class ctors
 --> test_resources/aka/test14.yr:(5,5)
 5  ~\lstbox{\textSFxi{}}~ def B = A::new;
    ~\lstbox{\textSFv{}}~     ^    ^^  note: maybe a type name alias was meant (i.e. using the token : instead of =)
    ~\lstbox{\textSFii{}}~~\lstbox{\textSFx{}}~~\lstbox{\textSFx{}}~~\lstbox{\textSFx{}}~~\lstbox{\textSFx{}}~~\lstbox{\textSFx{}}~~\lstbox{\textSFx{}}~~\lstbox{\textSFx{}}~
\end{lstlisting}

\subsubsection*{How to fix it}

Alias the class instead and construct through the alias.

\errorcode{E4038}{the symbol \errhole{} was declared but never used}
\label{sec:codes:e4038}

Raised during \textbf{validation}, from \texttt{ValidateErrorMessage::\allowbreak{}DECLARED\_\allowbreak{}NOT\_\allowbreak{}USED} (\textit{src/\allowbreak{}ymirc/\allowbreak{}semantic/\allowbreak{}validator/\allowbreak{}errors.yr}).

\subsubsection*{What it means}

A symbol was declared and never read. Since an unused declaration is almost always a leftover or a typo, the compiler reports it rather than letting it pass silently.

\subsubsection*{Example}

\textit{test\_\allowbreak{}resources/\allowbreak{}control\_\allowbreak{}flow/\allowbreak{}test1.yr}:

%% check: skip
\begin{lstlisting}[style=coloredverbatimError]
fn main () {
    let _a_ = 12;
    let a = 12;
    let _a = 12;
    let a_ = 12;
    let _ = 12;
}
\end{lstlisting}

\begin{lstlisting}[style=bashVerb, breaklines=true, escapechar=~]
Warning[E4038] : the symbol a was declared but never used
 --> test_resources/control_flow/test1.yr:(3,9)
 3  ~\lstbox{\textSFxi{}}~     let a = 12;
    ~\lstbox{\textSFv{}}~         ^
    ~\lstbox{\textSFii{}}~~\lstbox{\textSFx{}}~~\lstbox{\textSFx{}}~~\lstbox{\textSFx{}}~~\lstbox{\textSFx{}}~~\lstbox{\textSFx{}}~~\lstbox{\textSFx{}}~~\lstbox{\textSFx{}}~
\end{lstlisting}

\subsubsection*{How to fix it}

Remove the declaration, or use it. A value deliberately ignored can be bound to \tokennolst{\_}, which is never reported.

\errorcode{E4039}{cannot declare a variable of type \errhole{}}
\label{sec:codes:e4039}

Raised during \textbf{validation}, from \texttt{ValidateErrorMessage::\allowbreak{}DECLARE\_\allowbreak{}VOID\_\allowbreak{}VAR} (\textit{src/\allowbreak{}ymirc/\allowbreak{}semantic/\allowbreak{}validator/\allowbreak{}errors.yr}).

\subsubsection*{What it means}

A variable was declared with a type that has no values to hold -- \tokennolst{void}, or a type reduced to nothing.

\subsubsection*{Example}

\textit{test\_\allowbreak{}resources/\allowbreak{}lit\_\allowbreak{}slices/\allowbreak{}test21.yr}:

%% check: skip
\begin{lstlisting}[style=coloredverbatimError]
fn foo () {
    for i in copy [] {}
}
\end{lstlisting}

\begin{lstlisting}[style=bashVerb, breaklines=true, escapechar=~]
Error[E4039] : cannot declare a variable of type void
 --> test_resources/lit_slices/test21.yr:(2,9)
 2  ~\lstbox{\textSFxi{}}~     for i in copy [] {}
    ~\lstbox{\textSFv{}}~         ^
    ~\lstbox{\textSFii{}}~~\lstbox{\textSFx{}}~~\lstbox{\textSFx{}}~~\lstbox{\textSFx{}}~~\lstbox{\textSFx{}}~~\lstbox{\textSFx{}}~~\lstbox{\textSFx{}}~~\lstbox{\textSFx{}}~
\end{lstlisting}

\subsubsection*{How to fix it}

Remove the declaration. When the value came from a call, the function returns nothing and the call should be a statement rather than an initializer.

\errorcode{E4058}{external global variable cannot be declared with a value}
\label{sec:codes:e4058}

Raised during \textbf{validation}, from \texttt{ValidateErrorMessage::\allowbreak{}EXTERN\_\allowbreak{}GLOBAL\_\allowbreak{}WITH\_\allowbreak{}VALUE} (\textit{src/\allowbreak{}ymirc/\allowbreak{}semantic/\allowbreak{}validator/\allowbreak{}errors.yr}).

\subsubsection*{What it means}

An \tokennolst{extern} global variable was given an initial value. \tokennolst{extern} says the variable is defined elsewhere; the initial value belongs to that definition.

\subsubsection*{Example}

\textit{test\_\allowbreak{}resources/\allowbreak{}diagnostics/\allowbreak{}test102.yr}:

%% check: skip
\begin{lstlisting}[style=coloredverbatimError]
// an external global variable given a value
extern (C) static X : i32 = 1;

fn main () {
    let _ = X;
}
\end{lstlisting}

\begin{lstlisting}[style=bashVerb, breaklines=true, escapechar=~]
Error[E4058] : external global variable cannot be declared with a value
 --> test_resources/diagnostics/test102.yr:(2,19)
 2  ~\lstbox{\textSFxi{}}~ extern (C) static X : i32 = 1;
    ~\lstbox{\textSFv{}}~                   ^         -
    ~\lstbox{\textSFxi{}}~
    = when validating -- (test_resources/diagnostics/test102.yr:(2,19))
    ~\lstbox{\textSFii{}}~~\lstbox{\textSFx{}}~~\lstbox{\textSFx{}}~~\lstbox{\textSFx{}}~~\lstbox{\textSFx{}}~~\lstbox{\textSFx{}}~~\lstbox{\textSFx{}}~~\lstbox{\textSFx{}}~
Error[E4104] : symbol test102::X has errors
 --> test_resources/diagnostics/test102.yr:(5,13)
 5  ~\lstbox{\textSFxi{}}~     let _ = X;
    ~\lstbox{\textSFv{}}~             ^
 --> test_resources/diagnostics/test102.yr:(2,19)
 2  ~\lstbox{\textSFxi{}}~ extern (C) static X : i32 = 1;
    ~\lstbox{\textSFv{}}~                   -  error: external global variable cannot be declared with a value
    ~\lstbox{\textSFxi{}}~
    = when validating -- (test_resources/diagnostics/test102.yr:(2,19))
    = when validating test102::main -- (test_resources/diagnostics/test102.yr:(4,4))
    ~\lstbox{\textSFii{}}~~\lstbox{\textSFx{}}~~\lstbox{\textSFx{}}~~\lstbox{\textSFx{}}~~\lstbox{\textSFx{}}~~\lstbox{\textSFx{}}~~\lstbox{\textSFx{}}~~\lstbox{\textSFx{}}~
\end{lstlisting}

\subsubsection*{How to fix it}

Remove the value, or remove \tokennolst{extern} if this module is meant to define it.

\errorcode{E4068}{the name alias \errhole{} cannot be validated due to forward reference}
\label{sec:codes:e4068}

Raised during \textbf{validation}, from \texttt{ValidateErrorMessage::\allowbreak{}FORWARD\_\allowbreak{}REFERENCE\_\allowbreak{}AKA} (\textit{src/\allowbreak{}ymirc/\allowbreak{}semantic/\allowbreak{}validator/\allowbreak{}errors.yr}).

\subsubsection*{What it means}

A name alias (\tokennolst{def}) refers, directly or through other aliases, to something that is not yet defined at the point the compiler needs it.

\subsubsection*{Example}

\textit{test\_\allowbreak{}resources/\allowbreak{}aka/\allowbreak{}test9.yr}:

%% check: skip
\begin{lstlisting}[style=coloredverbatimError]
def A = V;
def V = A;
\end{lstlisting}

\begin{lstlisting}[style=bashVerb, breaklines=true, escapechar=~]
Error[E4104] : symbol test9::V has errors
 --> test_resources/aka/test9.yr:(1,9)
 1  ~\lstbox{\textSFxi{}}~ def A = V;
    ~\lstbox{\textSFv{}}~     -   ^  error: the name alias test9::A cannot be validated due to forward reference
 2  ~\lstbox{\textSFxi{}}~ def V = A;
    ~\lstbox{\textSFv{}}~         -  error: symbol test9::A has errors
    ~\lstbox{\textSFxi{}}~
    = when validating test9::V -- (test_resources/aka/test9.yr:(2,5))
    = when validating test9::A -- (test_resources/aka/test9.yr:(1,5))
    ~\lstbox{\textSFii{}}~~\lstbox{\textSFx{}}~~\lstbox{\textSFx{}}~~\lstbox{\textSFx{}}~~\lstbox{\textSFx{}}~~\lstbox{\textSFx{}}~~\lstbox{\textSFx{}}~~\lstbox{\textSFx{}}~
Error[E4104] : symbol test9::A has errors
 --> test_resources/aka/test9.yr:(2,9)
 1  ~\lstbox{\textSFxi{}}~ def A = V;
    ~\lstbox{\textSFv{}}~     -  error: the name alias test9::A cannot be validated due to forward reference
 2  ~\lstbox{\textSFxi{}}~ def V = A;
    ~\lstbox{\textSFv{}}~         ^
    ~\lstbox{\textSFxi{}}~
    = when validating test9::V -- (test_resources/aka/test9.yr:(2,5))
    ~\lstbox{\textSFii{}}~~\lstbox{\textSFx{}}~~\lstbox{\textSFx{}}~~\lstbox{\textSFx{}}~~\lstbox{\textSFx{}}~~\lstbox{\textSFx{}}~~\lstbox{\textSFx{}}~~\lstbox{\textSFx{}}~
\end{lstlisting}

\subsubsection*{How to fix it}

Reorder the declarations, or break the cycle the aliases form.

\errorcode{E4069}{the type cannot be validated due to forward reference}
\label{sec:codes:e4069}

Raised during \textbf{validation}, from \texttt{ValidateErrorMessage::\allowbreak{}FORWARD\_\allowbreak{}REFERENCE\_\allowbreak{}TYPE} (\textit{src/\allowbreak{}ymirc/\allowbreak{}semantic/\allowbreak{}validator/\allowbreak{}errors.yr}).

\subsubsection*{What it means}

A type refers to itself, or to a type still being validated, through a path the compiler must resolve first.

\subsubsection*{Example}

\textit{test\_\allowbreak{}resources/\allowbreak{}diagnostics/\allowbreak{}test135.yr}:

%% check: skip
\begin{lstlisting}[style=coloredverbatimError]
// two globals whose types depend on each other
lazy X : i32 = Y;
lazy Y : i32 = X;

fn main () {
    let _ = X;
}
\end{lstlisting}

\begin{lstlisting}[style=bashVerb, breaklines=true, escapechar=~]
Error[E4104] : symbol test135::Y has errors
 --> test_resources/diagnostics/test135.yr:(2,16)
 2  ~\lstbox{\textSFxi{}}~ lazy X : i32 = Y;
    ~\lstbox{\textSFv{}}~      -         ^  error: the type cannot be validated due to forward reference
 3  ~\lstbox{\textSFxi{}}~ lazy Y : i32 = X;
    ~\lstbox{\textSFv{}}~                -  error: symbol test135::X has errors
    ~\lstbox{\textSFxi{}}~
    = when validating -- (test_resources/diagnostics/test135.yr:(3,6))
    = when validating -- (test_resources/diagnostics/test135.yr:(2,6))
    ~\lstbox{\textSFii{}}~~\lstbox{\textSFx{}}~~\lstbox{\textSFx{}}~~\lstbox{\textSFx{}}~~\lstbox{\textSFx{}}~~\lstbox{\textSFx{}}~~\lstbox{\textSFx{}}~~\lstbox{\textSFx{}}~
Error[E4104] : symbol test135::X has errors
 --> test_resources/diagnostics/test135.yr:(3,16)
 2  ~\lstbox{\textSFxi{}}~ lazy X : i32 = Y;
    ~\lstbox{\textSFv{}}~      -  error: the type cannot be validated due to forward reference
 3  ~\lstbox{\textSFxi{}}~ lazy Y : i32 = X;
    ~\lstbox{\textSFv{}}~                ^
    ~\lstbox{\textSFxi{}}~
    = when validating -- (test_resources/diagnostics/test135.yr:(3,6))
    ~\lstbox{\textSFii{}}~~\lstbox{\textSFx{}}~~\lstbox{\textSFx{}}~~\lstbox{\textSFx{}}~~\lstbox{\textSFx{}}~~\lstbox{\textSFx{}}~~\lstbox{\textSFx{}}~~\lstbox{\textSFx{}}~
Error[E4104] : symbol test135::X has errors
 --> test_resources/diagnostics/test135.yr:(6,13)
 6  ~\lstbox{\textSFxi{}}~     let _ = X;
    ~\lstbox{\textSFv{}}~             ^
 --> test_resources/diagnostics/test135.yr:(2,6)
 2  ~\lstbox{\textSFxi{}}~ lazy X : i32 = Y;
    ~\lstbox{\textSFv{}}~      -         -
    ~\lstbox{\textSFxi{}}~      ~\lstbox{\textSFii{}}~~\lstbox{\textSFx{}}~ error: the type cannot be validated due to forward reference
    ~\lstbox{\textSFxi{}}~                ~\lstbox{\textSFii{}}~~\lstbox{\textSFx{}}~ error: symbol test135::Y has errors
 3  ~\lstbox{\textSFxi{}}~ lazy Y : i32 = X;
    ~\lstbox{\textSFv{}}~                -  error: symbol test135::X has errors
    ~\lstbox{\textSFxi{}}~
    = when validating -- (test_resources/diagnostics/test135.yr:(3,6))
    = when validating -- (test_resources/diagnostics/test135.yr:(2,6))
    = when validating test135::main -- (test_resources/diagnostics/test135.yr:(5,4))
    ~\lstbox{\textSFii{}}~~\lstbox{\textSFx{}}~~\lstbox{\textSFx{}}~~\lstbox{\textSFx{}}~~\lstbox{\textSFx{}}~~\lstbox{\textSFx{}}~~\lstbox{\textSFx{}}~~\lstbox{\textSFx{}}~
\end{lstlisting}

\subsubsection*{How to fix it}

Break the cycle by holding the self-reference behind a class or a pointer, cf. \hyperref[sec:codes:e4017]{E4017}.

\errorcode{E4080}{local global variable must be declared with a value}
\label{sec:codes:e4080}

Raised during \textbf{validation}, from \texttt{ValidateErrorMessage::\allowbreak{}GLOBAL\_\allowbreak{}NO\_\allowbreak{}VALUE} (\textit{src/\allowbreak{}ymirc/\allowbreak{}semantic/\allowbreak{}validator/\allowbreak{}errors.yr}).

\subsubsection*{What it means}

A global variable local to this module was declared without a value. Unlike an \tokennolst{extern} one, it is defined here and must be initialized.

\subsubsection*{Example}

\textit{test\_\allowbreak{}resources/\allowbreak{}diagnostics/\allowbreak{}test38.yr}:

%% check: skip
\begin{lstlisting}[style=coloredverbatimError]
lazy X : i32;

fn main () { X; }
\end{lstlisting}

\begin{lstlisting}[style=bashVerb, breaklines=true, escapechar=~]
Error[E4080] : local global variable must be declared with a value
 --> test_resources/diagnostics/test38.yr:(1,1)
 1  ~\lstbox{\textSFxi{}}~ lazy X : i32;
    ~\lstbox{\textSFv{}}~ ^^^^ -
    ~\lstbox{\textSFxi{}}~
    = when validating -- (test_resources/diagnostics/test38.yr:(1,6))
    ~\lstbox{\textSFii{}}~~\lstbox{\textSFx{}}~~\lstbox{\textSFx{}}~~\lstbox{\textSFx{}}~~\lstbox{\textSFx{}}~~\lstbox{\textSFx{}}~~\lstbox{\textSFx{}}~~\lstbox{\textSFx{}}~
Error[E4104] : symbol test38::X has errors
 --> test_resources/diagnostics/test38.yr:(3,14)
 1  ~\lstbox{\textSFxi{}}~ lazy X : i32;
    ~\lstbox{\textSFv{}}~ ----  error: local global variable must be declared with a value
 2  ~\lstbox{\textSFxi{}}~
 3  ~\lstbox{\textSFxi{}}~ fn main () { X; }
    ~\lstbox{\textSFv{}}~              ^
    ~\lstbox{\textSFxi{}}~
    = when validating -- (test_resources/diagnostics/test38.yr:(1,6))
    = when validating test38::main -- (test_resources/diagnostics/test38.yr:(3,4))
    ~\lstbox{\textSFii{}}~~\lstbox{\textSFx{}}~~\lstbox{\textSFx{}}~~\lstbox{\textSFx{}}~~\lstbox{\textSFx{}}~~\lstbox{\textSFx{}}~~\lstbox{\textSFx{}}~~\lstbox{\textSFx{}}~
\end{lstlisting}

\subsubsection*{How to fix it}

Give it an initial value, or declare it \tokennolst{extern} if it is defined elsewhere.

\errorcode{E4104}{symbol \errhole{} has errors}
\label{sec:codes:e4104}

Raised during \textbf{validation}, from \texttt{ValidateErrorMessage::\allowbreak{}INVALID\_\allowbreak{}SYMBOLS} (\textit{src/\allowbreak{}ymirc/\allowbreak{}semantic/\allowbreak{}validator/\allowbreak{}errors.yr}).

\subsubsection*{What it means}

A symbol this code uses could not be validated, so the compiler cannot say anything about the use. The error on the symbol itself is reported under this one.

\subsubsection*{Example}

\textit{test\_\allowbreak{}resources/\allowbreak{}aka/\allowbreak{}test9.yr}:

%% check: skip
\begin{lstlisting}[style=coloredverbatimError]
def A = V;
def V = A;
\end{lstlisting}

\begin{lstlisting}[style=bashVerb, breaklines=true, escapechar=~]
Error[E4104] : symbol test9::V has errors
 --> test_resources/aka/test9.yr:(1,9)
 1  ~\lstbox{\textSFxi{}}~ def A = V;
    ~\lstbox{\textSFv{}}~     -   ^  error: the name alias test9::A cannot be validated due to forward reference
 2  ~\lstbox{\textSFxi{}}~ def V = A;
    ~\lstbox{\textSFv{}}~         -  error: symbol test9::A has errors
    ~\lstbox{\textSFxi{}}~
    = when validating test9::V -- (test_resources/aka/test9.yr:(2,5))
    = when validating test9::A -- (test_resources/aka/test9.yr:(1,5))
    ~\lstbox{\textSFii{}}~~\lstbox{\textSFx{}}~~\lstbox{\textSFx{}}~~\lstbox{\textSFx{}}~~\lstbox{\textSFx{}}~~\lstbox{\textSFx{}}~~\lstbox{\textSFx{}}~~\lstbox{\textSFx{}}~
Error[E4104] : symbol test9::A has errors
 --> test_resources/aka/test9.yr:(2,9)
 1  ~\lstbox{\textSFxi{}}~ def A = V;
    ~\lstbox{\textSFv{}}~     -  error: the name alias test9::A cannot be validated due to forward reference
 2  ~\lstbox{\textSFxi{}}~ def V = A;
    ~\lstbox{\textSFv{}}~         ^
    ~\lstbox{\textSFxi{}}~
    = when validating test9::V -- (test_resources/aka/test9.yr:(2,5))
    ~\lstbox{\textSFii{}}~~\lstbox{\textSFx{}}~~\lstbox{\textSFx{}}~~\lstbox{\textSFx{}}~~\lstbox{\textSFx{}}~~\lstbox{\textSFx{}}~~\lstbox{\textSFx{}}~~\lstbox{\textSFx{}}~
\end{lstlisting}

\subsubsection*{How to fix it}

Fix the reported errors on the symbol. This diagnostic disappears with them.

\errorcode{E4106}{the identifier \errhole{} describes a native type but has no meaning within this context}
\label{sec:codes:e4106}

Raised during \textbf{validation}, from \texttt{ValidateErrorMessage::\allowbreak{}IS\_\allowbreak{}NATIVE\_\allowbreak{}TYPE} (\textit{src/\allowbreak{}ymirc/\allowbreak{}semantic/\allowbreak{}validator/\allowbreak{}errors.yr}).

\subsubsection*{What it means}

The name of a native type was used where it has no meaning -- as a value, for instance, rather than as a type.

\subsubsection*{Example}

\textit{test\_\allowbreak{}resources/\allowbreak{}diagnostics/\allowbreak{}test108.yr}:

%% check: skip
\begin{lstlisting}[style=coloredverbatimError]
// a native type used as a value
fn main () {
    let _ = i32;
}
\end{lstlisting}

\begin{lstlisting}[style=bashVerb, breaklines=true, escapechar=~]
Error[E3027] : undefined symbol i32
 --> test_resources/diagnostics/test108.yr:(3,13)
 3  ~\lstbox{\textSFxi{}}~     let _ = i32;
    ~\lstbox{\textSFv{}}~             ^^^
    ~\lstbox{\textSFxi{}}~
    = when validating test108::main -- (test_resources/diagnostics/test108.yr:(2,4))
    ~\lstbox{\textSFii{}}~~\lstbox{\textSFx{}}~~\lstbox{\textSFx{}}~~\lstbox{\textSFx{}}~~\lstbox{\textSFx{}}~~\lstbox{\textSFx{}}~~\lstbox{\textSFx{}}~~\lstbox{\textSFx{}}~
Error[E4106] : the identifier i32 describes a native type but has no meaning within this context
 --> test_resources/diagnostics/test108.yr:(3,13)
 3  ~\lstbox{\textSFxi{}}~     let _ = i32;
    ~\lstbox{\textSFv{}}~             ^^^
    ~\lstbox{\textSFxi{}}~
    = when validating test108::main -- (test_resources/diagnostics/test108.yr:(2,4))
    ~\lstbox{\textSFii{}}~~\lstbox{\textSFx{}}~~\lstbox{\textSFx{}}~~\lstbox{\textSFx{}}~~\lstbox{\textSFx{}}~~\lstbox{\textSFx{}}~~\lstbox{\textSFx{}}~~\lstbox{\textSFx{}}~
\end{lstlisting}

\subsubsection*{How to fix it}

Use the identifier where a type is expected, or use a value of that type.

\errorcode{E4135}{expression refers to multiple symbols that generate types}
\label{sec:codes:e4135}

Raised during \textbf{validation}, from \texttt{ValidateErrorMessage::\allowbreak{}MULTIPLE\_\allowbreak{}SYMBOL\_\allowbreak{}TYPES} (\textit{src/\allowbreak{}ymirc/\allowbreak{}semantic/\allowbreak{}validator/\allowbreak{}errors.yr}).

\subsubsection*{What it means}

A name resolves to several symbols that all denote types, so the expression has no single type to mean.

\subsubsection*{Example}

\textit{test\_\allowbreak{}resources/\allowbreak{}diagnostics/\allowbreak{}test114.yr}:

%% check: skip
\begin{lstlisting}[style=coloredverbatimError]
// a type name brought by two imported modules
mod first {
    pub class Point {
        pub self () {}
    }
}

mod second {
    pub class Point {
        pub self () {}
    }
}

use first;
use second;

fn main () {
    let _p_ : &Point = copy first::Point ();
}
\end{lstlisting}

\begin{lstlisting}[style=bashVerb, breaklines=true, escapechar=~]
Error[E4135] : expression refers to multiple symbols that generate types
 --> test_resources/diagnostics/test114.yr:(18,16)
18  ~\lstbox{\textSFxi{}}~     let _p_ : &Point = copy first::Point ();
    ~\lstbox{\textSFv{}}~                ^^^^^
 --> test_resources/diagnostics/test114.yr:(3,15)
 3  ~\lstbox{\textSFxi{}}~     pub class Point {
    ~\lstbox{\textSFv{}}~               -----  note: candidate test114::first::Point
   ...
 9  ~\lstbox{\textSFxi{}}~     pub class Point {
    ~\lstbox{\textSFv{}}~               -----  note: candidate test114::second::Point
    ~\lstbox{\textSFxi{}}~
    = when validating test114::main -- (test_resources/diagnostics/test114.yr:(17,4))
    ~\lstbox{\textSFii{}}~~\lstbox{\textSFx{}}~~\lstbox{\textSFx{}}~~\lstbox{\textSFx{}}~~\lstbox{\textSFx{}}~~\lstbox{\textSFx{}}~~\lstbox{\textSFx{}}~~\lstbox{\textSFx{}}~
\end{lstlisting}

\subsubsection*{How to fix it}

Qualify the name with the module it should come from, or remove the import that brings in the other one.

\errorcode{E4199}{declaration of \errhole{} shadows another declaration}
\label{sec:codes:e4199}

Raised during \textbf{validation}, from \texttt{ValidateErrorMessage::\allowbreak{}SHADOWING\_\allowbreak{}DECL} (\textit{src/\allowbreak{}ymirc/\allowbreak{}semantic/\allowbreak{}validator/\allowbreak{}errors.yr}).

\subsubsection*{What it means}

A declaration hides an earlier one still in scope, cf. \hyperref[sec:codes:e3018]{E3018}.

\subsubsection*{Example}

\textit{test\_\allowbreak{}resources/\allowbreak{}enumeration/\allowbreak{}test17.yr}:

%% check: skip
\begin{lstlisting}[style=coloredverbatimError]
enum X
| A = 1
| A = 2;
\end{lstlisting}

\begin{lstlisting}[style=bashVerb, breaklines=true, escapechar=~]
Error[E4199] : declaration of A shadows another declaration
 --> test_resources/enumeration/test17.yr:(3,1)
 3  ~\lstbox{\textSFxi{}}~ | A = 2;
    ~\lstbox{\textSFv{}}~ ^
    ~\lstbox{\textSFxi{}}~ Note : defined here
    ~\lstbox{\textSFxi{}}~  --> test_resources/enumeration/test17.yr:(2,7)
    ~\lstbox{\textSFxi{}}~  2  ~\lstbox{\textSFxi{}}~ | A = 1
    ~\lstbox{\textSFxi{}}~     ~\lstbox{\textSFv{}}~       ^
    ~\lstbox{\textSFxi{}}~     ~\lstbox{\textSFii{}}~~\lstbox{\textSFx{}}~~\lstbox{\textSFx{}}~~\lstbox{\textSFx{}}~~\lstbox{\textSFx{}}~~\lstbox{\textSFx{}}~~\lstbox{\textSFx{}}~~\lstbox{\textSFx{}}~
    ~\lstbox{\textSFii{}}~~\lstbox{\textSFx{}}~~\lstbox{\textSFx{}}~~\lstbox{\textSFx{}}~~\lstbox{\textSFx{}}~~\lstbox{\textSFx{}}~~\lstbox{\textSFvii{}}~~\lstbox{\textSFx{}}~
\end{lstlisting}

\subsubsection*{How to fix it}

Rename one of the two.

\errorcode{E4281}{var declaration \errhole{} outside a function context}
\label{sec:codes:e4281}

Raised during \textbf{validation}, from \texttt{ValidateErrorMessage::\allowbreak{}DECLARE\_\allowbreak{}VARIABLE\_\allowbreak{}NO\_\allowbreak{}FUNCTION} (\textit{src/\allowbreak{}ymirc/\allowbreak{}semantic/\allowbreak{}validator/\allowbreak{}errors.yr}).

\subsubsection*{What it means}

A \tokennolst{var} declaration was written outside any function. Variables live in a function frame; a declaration at module level is a global variable, declared differently.

\subsubsection*{How to fix it}

Move the declaration into a function, or declare it as a global.

\errorcode{E4290}{none of the symbols \errhole{} can be used as a value}
\label{sec:codes:e4290}

Raised during \textbf{validation}, from \texttt{ValidateErrorMessage::\allowbreak{}UNUSABLE\_\allowbreak{}SYMBOLS} (\textit{src/\allowbreak{}ymirc/\allowbreak{}semantic/\allowbreak{}validator/\allowbreak{}errors.yr}).

\subsubsection*{What it means}

A name designates several symbols, and each of them was rejected as a value -- a type, a trait, a method referenced without an instance. The reason each was rejected is listed under the diagnostic.

\subsubsection*{Example}

\textit{test\_\allowbreak{}resources/\allowbreak{}diagnostics/\allowbreak{}test142.yr}:

%% check: skip
\begin{lstlisting}[style=coloredverbatimError]
// two imported traits of the same name, used as a value
mod a;
mod b;

use test142::a;
use test142::b;

fn main () {
    let _ = T;
}
\end{lstlisting}

\begin{lstlisting}[style=bashVerb, breaklines=true, escapechar=~]
Error[E4290] : none of the symbols test142::a::T, test142::b::T can be used as a value
 --> test_resources/diagnostics/test142.yr:(9,13)
 9  ~\lstbox{\textSFxi{}}~     let _ = T;
    ~\lstbox{\textSFv{}}~             ^  error: expression produces the type test142::a::T, but is used as a value
    ~\lstbox{\textSFxi{}}~             ~\lstbox{\textSFii{}}~~\lstbox{\textSFx{}}~ error: expression produces the type test142::b::T, but is used as a value
    ~\lstbox{\textSFxi{}}~
    = when validating test142::main -- (test_resources/diagnostics/test142.yr:(8,4))
    ~\lstbox{\textSFii{}}~~\lstbox{\textSFx{}}~~\lstbox{\textSFx{}}~~\lstbox{\textSFx{}}~~\lstbox{\textSFx{}}~~\lstbox{\textSFx{}}~~\lstbox{\textSFx{}}~~\lstbox{\textSFx{}}~
\end{lstlisting}

\subsubsection*{How to fix it}

Fix the use according to the reason given for the symbol that was meant, and qualify it by its module (\tokennolst{a::T}) when the others are not.

%% Generated by tools/gen_error_codes.py from docs/errors/ of the compiler;
%% edit the compiler's pages or tools/error_themes.txt, then rerun it.

\clearpage
\section{Literals and arithmetic}
\label{sec:codes:literals}

The book's chapter \textit{Native scalar types} specifies the rules behind these codes.

\begin{longtable}{@{}l p{0.78\linewidth} r@{}}
\toprule Code & Message & Page\\ \midrule \endhead
\hyperref[sec:codes:e4046]{E4046} & division by 0 & \pageref{sec:codes:e4046}\\
\hyperref[sec:codes:e4050]{E4050} & empty interpolation, expected an expression between \$\{ and \} & \pageref{sec:codes:e4050}\\
\hyperref[sec:codes:e4115]{E4115} & malformed literal, number of \errhole{} is \errhole{} & \pageref{sec:codes:e4115}\\
\hyperref[sec:codes:e4117]{E4117} & malformed float literal \errhole{} & \pageref{sec:codes:e4117}\\
\hyperref[sec:codes:e4118]{E4118} & malformed int literal \errhole{} & \pageref{sec:codes:e4118}\\
\hyperref[sec:codes:e4178]{E4178} & overflow capacity for type \errhole{} = \errhole{}, maximum value is \errhole{} & \pageref{sec:codes:e4178}\\
\hyperref[sec:codes:e4180]{E4180} & overflow capacity for type \errhole{} = \errhole{}, minimum value is \errhole{} & \pageref{sec:codes:e4180}\\
\hyperref[sec:codes:e4230]{E4230} & undefined escape sequence & \pageref{sec:codes:e4230}\\
\hyperref[sec:codes:e4264]{E4264} & unterminated escape sequence & \pageref{sec:codes:e4264}\\
\hyperref[sec:codes:e4265]{E4265} & unterminated interpolation, missing closing \} & \pageref{sec:codes:e4265}\\
\bottomrule
\end{longtable}

\errorcode{E4046}{division by 0}
\label{sec:codes:e4046}

Raised during \textbf{validation}, from \texttt{ValidateErrorMessage::\allowbreak{}DIVISION\_\allowbreak{}BY\_\allowbreak{}ZERO} (\textit{src/\allowbreak{}ymirc/\allowbreak{}semantic/\allowbreak{}validator/\allowbreak{}errors.yr}).

\subsubsection*{What it means}

A division or a modulo by a zero the compiler could compute.

\subsubsection*{Example}

\textit{test\_\allowbreak{}resources/\allowbreak{}binary\_\allowbreak{}int/\allowbreak{}test5.yr}:

%% check: skip
\begin{lstlisting}[style=coloredverbatimError]
in test5;

fn main () {
    cte {1 % 0};

    1u32 + 3;

    let a = 12;
    a = 9;

    3 + 1i32;

    let mut b = 1;
    b = 1us;

}
\end{lstlisting}

\begin{lstlisting}[style=bashVerb, breaklines=true, escapechar=~]
Error[E4046] : division by 0
 --> test_resources/binary_int/test5.yr:(4,12)
 4  ~\lstbox{\textSFxi{}}~     cte {1 % 0};
    ~\lstbox{\textSFv{}}~            ^
    ~\lstbox{\textSFii{}}~~\lstbox{\textSFx{}}~~\lstbox{\textSFx{}}~~\lstbox{\textSFx{}}~~\lstbox{\textSFx{}}~~\lstbox{\textSFx{}}~~\lstbox{\textSFx{}}~~\lstbox{\textSFx{}}~
\end{lstlisting}

\subsubsection*{How to fix it}

Fix the divisor, or test it before dividing when it is not statically known.

\errorcode{E4050}{empty interpolation, expected an expression between \$\{ and \}}
\label{sec:codes:e4050}

Raised during \textbf{validation}, from \texttt{ValidateErrorMessage::\allowbreak{}EMPTY\_\allowbreak{}FORMAT\_\allowbreak{}EMBED} (\textit{src/\allowbreak{}ymirc/\allowbreak{}semantic/\allowbreak{}validator/\allowbreak{}errors.yr}).

\subsubsection*{What it means}

A string interpolation \tokennolst{\$\{\}} contains nothing between its braces.

\subsubsection*{Example}

\textit{test\_\allowbreak{}resources/\allowbreak{}format\_\allowbreak{}string/\allowbreak{}test8.yr}:

%% check: skip
\begin{lstlisting}[style=coloredverbatimError]
in test8;

fn main () {
    let str = f"empty: ${}";
    str;
}
\end{lstlisting}

\begin{lstlisting}[style=bashVerb, breaklines=true, escapechar=~]
Error[E4050] : empty interpolation, expected an expression between ${ and }
 --> test_resources/format_string/test8.yr:(4,26)
 4  ~\lstbox{\textSFxi{}}~     let str = f"empty: ${}";
    ~\lstbox{\textSFv{}}~                          ^
    ~\lstbox{\textSFii{}}~~\lstbox{\textSFx{}}~~\lstbox{\textSFx{}}~~\lstbox{\textSFx{}}~~\lstbox{\textSFx{}}~~\lstbox{\textSFx{}}~~\lstbox{\textSFx{}}~~\lstbox{\textSFx{}}~
\end{lstlisting}

\subsubsection*{How to fix it}

Write the expression to interpolate, or remove the interpolation.

\errorcode{E4115}{malformed literal, number of \errhole{} is \errhole{}}
\label{sec:codes:e4115}

Raised during \textbf{validation}, from \texttt{ValidateErrorMessage::\allowbreak{}MALFORMED\_\allowbreak{}CHAR} (\textit{src/\allowbreak{}ymirc/\allowbreak{}semantic/\allowbreak{}validator/\allowbreak{}errors.yr}).

\subsubsection*{What it means}

A character literal does not hold exactly one character -- the count the compiler found is in the message.

\subsubsection*{Example}

\textit{test\_\allowbreak{}resources/\allowbreak{}diagnostics/\allowbreak{}test50.yr}:

%% check: skip
\begin{lstlisting}[style=coloredverbatimError]
fn main () {
    let _ = 'ab'c8;
}
\end{lstlisting}

\begin{lstlisting}[style=bashVerb, breaklines=true, escapechar=~]
Error[E4115] : malformed literal, number of c8 is 2
 --> test_resources/diagnostics/test50.yr:(2,13)
 2  ~\lstbox{\textSFxi{}}~     let _ = 'ab'c8;
    ~\lstbox{\textSFv{}}~             ^
    ~\lstbox{\textSFii{}}~~\lstbox{\textSFx{}}~~\lstbox{\textSFx{}}~~\lstbox{\textSFx{}}~~\lstbox{\textSFx{}}~~\lstbox{\textSFx{}}~~\lstbox{\textSFx{}}~~\lstbox{\textSFx{}}~
\end{lstlisting}

\subsubsection*{How to fix it}

Write a single character between quotes, or use a string literal if several characters were meant. A multi-byte character needs the character type wide enough to hold it.

\errorcode{E4117}{malformed float literal \errhole{}}
\label{sec:codes:e4117}

Raised during \textbf{validation}, from \texttt{ValidateErrorMessage::\allowbreak{}MALFORMED\_\allowbreak{}FLOAT\_\allowbreak{}LITERAL} (\textit{src/\allowbreak{}ymirc/\allowbreak{}semantic/\allowbreak{}validator/\allowbreak{}errors.yr}).

\subsubsection*{What it means}

A floating-point literal cannot be represented: it is malformed, or outside what its type can hold.

\subsubsection*{Example}

\textit{test\_\allowbreak{}resources/\allowbreak{}diagnostics/\allowbreak{}test110.yr}:

%% check: skip
\begin{lstlisting}[style=coloredverbatimError]
// a float literal whose exponent does not fit
fn main () {
    let _ = 1.0e99999999999999999999;
}
\end{lstlisting}

\begin{lstlisting}[style=bashVerb, breaklines=true, escapechar=~]
Error[E4117] : malformed float literal 1.0e+99999999999999999999
 --> test_resources/diagnostics/test110.yr:(3,13)
 3  ~\lstbox{\textSFxi{}}~     let _ = 1.0e99999999999999999999;
    ~\lstbox{\textSFv{}}~             ^
    ~\lstbox{\textSFxi{}}~
    = when validating test110::main -- (test_resources/diagnostics/test110.yr:(2,4))
    ~\lstbox{\textSFii{}}~~\lstbox{\textSFx{}}~~\lstbox{\textSFx{}}~~\lstbox{\textSFx{}}~~\lstbox{\textSFx{}}~~\lstbox{\textSFx{}}~~\lstbox{\textSFx{}}~~\lstbox{\textSFx{}}~
\end{lstlisting}

\subsubsection*{How to fix it}

Correct the literal, or widen the type.

\errorcode{E4118}{malformed int literal \errhole{}}
\label{sec:codes:e4118}

Raised during \textbf{validation}, from \texttt{ValidateErrorMessage::\allowbreak{}MALFORMED\_\allowbreak{}INT\_\allowbreak{}LITERAL} (\textit{src/\allowbreak{}ymirc/\allowbreak{}semantic/\allowbreak{}validator/\allowbreak{}errors.yr}).

\subsubsection*{What it means}

An integer literal cannot be represented: it is malformed, or outside what its type can hold.

\subsubsection*{Example}

\textit{test\_\allowbreak{}resources/\allowbreak{}diagnostics/\allowbreak{}test51.yr}:

%% check: skip
\begin{lstlisting}[style=coloredverbatimError]
fn main () {
    let _ = 99999999999999999999999999999999999999;
}
\end{lstlisting}

\begin{lstlisting}[style=bashVerb, breaklines=true, escapechar=~]
Error[E4118] : malformed int literal 99999999999999999999999999999999999999
 --> test_resources/diagnostics/test51.yr:(2,13)
 2  ~\lstbox{\textSFxi{}}~     let _ = 99999999999999999999999999999999999999;
    ~\lstbox{\textSFv{}}~             ^^^^^^^^^^^^^^^^^^^^^^^^^^^^^^^^^^^^^^
    ~\lstbox{\textSFii{}}~~\lstbox{\textSFx{}}~~\lstbox{\textSFx{}}~~\lstbox{\textSFx{}}~~\lstbox{\textSFx{}}~~\lstbox{\textSFx{}}~~\lstbox{\textSFx{}}~~\lstbox{\textSFx{}}~
\end{lstlisting}

\subsubsection*{How to fix it}

Correct the literal, or use a wider integer type. A suffix on the literal selects its type explicitly.

\errorcode{E4178}{overflow capacity for type \errhole{} = \errhole{}, maximum value is \errhole{}}
\label{sec:codes:e4178}

Raised during \textbf{validation}, from \texttt{ValidateErrorMessage::\allowbreak{}OVERFLOW\_\allowbreak{}CAPACITY} (\textit{src/\allowbreak{}ymirc/\allowbreak{}semantic/\allowbreak{}validator/\allowbreak{}errors.yr}).

\subsubsection*{What it means}

A value the compiler computed does not fit in its type. The type, the value and the maximum it can hold are all in the message.

\subsubsection*{Example}

\textit{test\_\allowbreak{}resources/\allowbreak{}enumeration/\allowbreak{}test15.yr}:

%% check: skip
\begin{lstlisting}[style=coloredverbatimError]
enum X : i32
| A = 1
| B = 900;

fn main () {
    let mut a : i8 = X::A;
    a = X::B;

}
\end{lstlisting}

\begin{lstlisting}[style=bashVerb, breaklines=true, escapechar=~]
Error[E4178] : overflow capacity for type mut i8 = 900, maximum value is 127
 --> test_resources/enumeration/test15.yr:(3,7)
 3  ~\lstbox{\textSFxi{}}~ | B = 900;
    ~\lstbox{\textSFv{}}~       ^^^
    ~\lstbox{\textSFii{}}~~\lstbox{\textSFx{}}~~\lstbox{\textSFx{}}~~\lstbox{\textSFx{}}~~\lstbox{\textSFx{}}~~\lstbox{\textSFx{}}~~\lstbox{\textSFx{}}~~\lstbox{\textSFx{}}~
\end{lstlisting}

\subsubsection*{How to fix it}

Use a wider type, or reduce the value.

\errorcode{E4180}{overflow capacity for type \errhole{} = \errhole{}, minimum value is \errhole{}}
\label{sec:codes:e4180}

Raised during \textbf{validation}, from \texttt{ValidateErrorMessage::\allowbreak{}OVERFLOW\_\allowbreak{}CAPACITY\_\allowbreak{}MIN} (\textit{src/\allowbreak{}ymirc/\allowbreak{}semantic/\allowbreak{}validator/\allowbreak{}errors.yr}).

\subsubsection*{What it means}

A value the compiler computed is below the minimum of its type.

\subsubsection*{Example}

\textit{test\_\allowbreak{}resources/\allowbreak{}regressions/\allowbreak{}negative\_\allowbreak{}int\_\allowbreak{}literals/\allowbreak{}test2.yr}:

%% check: skip
\begin{lstlisting}[style=coloredverbatimError]
in test2;

/**
 * One past the minimum stays refused, and the error names the negative value
 * that was written rather than the magnitude behind it.
 */
fn main () {
    -2147483649i32;
    -129i8;
    -9223372036854775809i64;
}
\end{lstlisting}

\begin{lstlisting}[style=bashVerb, breaklines=true, escapechar=~]
Error[E4180] : overflow capacity for type i32 = -2147483649, minimum value is -2147483648
 --> test_resources/regressions/negative_int_literals/test2.yr:(8,6)
 8  ~\lstbox{\textSFxi{}}~     -2147483649i32;
    ~\lstbox{\textSFv{}}~      ^^^^^^^^^^
    ~\lstbox{\textSFii{}}~~\lstbox{\textSFx{}}~~\lstbox{\textSFx{}}~~\lstbox{\textSFx{}}~~\lstbox{\textSFx{}}~~\lstbox{\textSFx{}}~~\lstbox{\textSFx{}}~~\lstbox{\textSFx{}}~
\end{lstlisting}

\subsubsection*{How to fix it}

Use a signed or wider type, or correct the value.

\errorcode{E4230}{undefined escape sequence}
\label{sec:codes:e4230}

Raised during \textbf{validation}, from \texttt{ValidateErrorMessage::\allowbreak{}UNDEFINED\_\allowbreak{}ESCAPE} (\textit{src/\allowbreak{}ymirc/\allowbreak{}semantic/\allowbreak{}validator/\allowbreak{}errors.yr}).

\subsubsection*{What it means}

A backslash in a literal is followed by a character that is not an escape the language defines.

\subsubsection*{Example}

\textit{test\_\allowbreak{}resources/\allowbreak{}diagnostics/\allowbreak{}test75.yr}:

%% check: skip
\begin{lstlisting}[style=coloredverbatimError]
fn main () {
    let _ = "\q";
}
\end{lstlisting}

\begin{lstlisting}[style=bashVerb, breaklines=true, escapechar=~]
Error[E4230] : undefined escape sequence
 --> test_resources/diagnostics/test75.yr:(2,15)
 2  ~\lstbox{\textSFxi{}}~     let _ = "\q";
    ~\lstbox{\textSFv{}}~               ^
    ~\lstbox{\textSFii{}}~~\lstbox{\textSFx{}}~~\lstbox{\textSFx{}}~~\lstbox{\textSFx{}}~~\lstbox{\textSFx{}}~~\lstbox{\textSFx{}}~~\lstbox{\textSFx{}}~~\lstbox{\textSFx{}}~
\end{lstlisting}

\subsubsection*{How to fix it}

Use a defined escape, or double the backslash for a literal one.

\errorcode{E4264}{unterminated escape sequence}
\label{sec:codes:e4264}

Raised during \textbf{validation}, from \texttt{ValidateErrorMessage::\allowbreak{}UNTERMINATED\_\allowbreak{}ESCAPE} (\textit{src/\allowbreak{}ymirc/\allowbreak{}semantic/\allowbreak{}validator/\allowbreak{}errors.yr}).

\subsubsection*{What it means}

A literal ends inside an escape sequence, cf. \hyperref[sec:codes:e1007]{E1007}.

\subsubsection*{Example}

\textit{test\_\allowbreak{}resources/\allowbreak{}diagnostics/\allowbreak{}test125.yr}:

%% check: skip
\begin{lstlisting}[style=coloredverbatimError]
// a unicode escape with no code point
fn main () {
    let _ = "\u";
}
\end{lstlisting}

\begin{lstlisting}[style=bashVerb, breaklines=true, escapechar=~]
Error[E4264] : unterminated escape sequence
 --> test_resources/diagnostics/test125.yr:(3,16)
 3  ~\lstbox{\textSFxi{}}~     let _ = "\u";
    ~\lstbox{\textSFv{}}~                ^
    ~\lstbox{\textSFxi{}}~
    = when validating test125::main -- (test_resources/diagnostics/test125.yr:(2,4))
    ~\lstbox{\textSFii{}}~~\lstbox{\textSFx{}}~~\lstbox{\textSFx{}}~~\lstbox{\textSFx{}}~~\lstbox{\textSFx{}}~~\lstbox{\textSFx{}}~~\lstbox{\textSFx{}}~~\lstbox{\textSFx{}}~
\end{lstlisting}

\subsubsection*{How to fix it}

Complete the escape, or double the backslash.

\errorcode{E4265}{unterminated interpolation, missing closing \}}
\label{sec:codes:e4265}

Raised during \textbf{validation}, from \texttt{ValidateErrorMessage::\allowbreak{}UNTERMINATED\_\allowbreak{}FORMAT\_\allowbreak{}EMBED} (\textit{src/\allowbreak{}ymirc/\allowbreak{}semantic/\allowbreak{}validator/\allowbreak{}errors.yr}).

\subsubsection*{What it means}

A string interpolation was opened with \tokennolst{\$\{} and never closed.

\subsubsection*{Example}

\textit{test\_\allowbreak{}resources/\allowbreak{}format\_\allowbreak{}string/\allowbreak{}test5.yr}:

%% check: skip
\begin{lstlisting}[style=coloredverbatimError]
in test5;

fn main () {
    let str = f"unterminated ${1 + 2";
    str;
}
\end{lstlisting}

\begin{lstlisting}[style=bashVerb, breaklines=true, escapechar=~]
Error[E4265] : unterminated interpolation, missing closing }
 --> test_resources/format_string/test5.yr:(4,30)
 4  ~\lstbox{\textSFxi{}}~     let str = f"unterminated ${1 + 2";
    ~\lstbox{\textSFv{}}~                              ^^
    ~\lstbox{\textSFii{}}~~\lstbox{\textSFx{}}~~\lstbox{\textSFx{}}~~\lstbox{\textSFx{}}~~\lstbox{\textSFx{}}~~\lstbox{\textSFx{}}~~\lstbox{\textSFx{}}~~\lstbox{\textSFx{}}~
\end{lstlisting}

\subsubsection*{How to fix it}

Add the closing brace.

%% Generated by tools/gen_error_codes.py from docs/errors/ of the compiler;
%% edit the compiler's pages or tools/error_themes.txt, then rerun it.

\clearpage
\section{Types and conversions}
\label{sec:codes:types}

\begin{longtable}{@{}l p{0.78\linewidth} r@{}}
\toprule Code & Message & Page\\ \midrule \endhead
\hyperref[sec:codes:e4049]{E4049} & array inner type cannot be \errhole{} & \pageref{sec:codes:e4049}\\
\hyperref[sec:codes:e4095]{E4095} & incompatible types \errhole{} and \errhole{} & \pageref{sec:codes:e4095}\\
\hyperref[sec:codes:e4096]{E4096} & incompatible values \errhole{} and \errhole{} & \pageref{sec:codes:e4096}\\
\hyperref[sec:codes:e4097]{E4097} & the type \errhole{} is not complete & \pageref{sec:codes:e4097}\\
\hyperref[sec:codes:e4217]{E4217} & type \errhole{} has no fields & \pageref{sec:codes:e4217}\\
\hyperref[sec:codes:e4218]{E4218} & temporary type \errhole{} has no size & \pageref{sec:codes:e4218}\\
\hyperref[sec:codes:e4220]{E4220} & type \errhole{} has no field named \errhole{} & \pageref{sec:codes:e4220}\\
\hyperref[sec:codes:e4227]{E4227} & cannot cast value of type \errhole{} into a value of type \errhole{} & \pageref{sec:codes:e4227}\\
\hyperref[sec:codes:e4270]{E4270} & expression produces a value, but is used as a type & \pageref{sec:codes:e4270}\\
\hyperref[sec:codes:e4272]{E4272} & expression produces a value of type \errhole{}, but is used as a type & \pageref{sec:codes:e4272}\\
\hyperref[sec:codes:e4273]{E4273} & expression produces a type, but is used as a value & \pageref{sec:codes:e4273}\\
\hyperref[sec:codes:e4275]{E4275} & expression produces the type \errhole{}, but is used as a value & \pageref{sec:codes:e4275}\\
\hyperref[sec:codes:e4278]{E4278} & void expression cannot be used as a value & \pageref{sec:codes:e4278}\\
\hyperref[sec:codes:e4304]{E4304} & \errhole{} has no type of its own, it takes the type of what it is assigned to & \pageref{sec:codes:e4304}\\
\bottomrule
\end{longtable}

\errorcode{E4049}{array inner type cannot be \errhole{}}
\label{sec:codes:e4049}

Raised during \textbf{validation}, from \texttt{ValidateErrorMessage::\allowbreak{}EMPTY\_\allowbreak{}ARRAY\_\allowbreak{}INNER} (\textit{src/\allowbreak{}ymirc/\allowbreak{}semantic/\allowbreak{}validator/\allowbreak{}errors.yr}).

\subsubsection*{What it means}

An array was declared with an inner type that cannot be stored in an array -- a type with no size, such as \tokennolst{void}.

\subsubsection*{Example}

\textit{test\_\allowbreak{}resources/\allowbreak{}diagnostics/\allowbreak{}test101.yr}:

%% check: skip
\begin{lstlisting}[style=coloredverbatimError]
// a static array of void values
fn nothing () {}

fn main () {
    let _ = [nothing (); 3];
}
\end{lstlisting}

\begin{lstlisting}[style=bashVerb, breaklines=true, escapechar=~]
Error[E4049] : array inner type cannot be void
 --> test_resources/diagnostics/test101.yr:(5,22)
 5  ~\lstbox{\textSFxi{}}~     let _ = [nothing (); 3];
    ~\lstbox{\textSFv{}}~                      ^
    ~\lstbox{\textSFxi{}}~
    = when validating test101::main -- (test_resources/diagnostics/test101.yr:(4,4))
    ~\lstbox{\textSFii{}}~~\lstbox{\textSFx{}}~~\lstbox{\textSFx{}}~~\lstbox{\textSFx{}}~~\lstbox{\textSFx{}}~~\lstbox{\textSFx{}}~~\lstbox{\textSFx{}}~~\lstbox{\textSFx{}}~
\end{lstlisting}

\subsubsection*{How to fix it}

Use a type with values. If the element type came from a call, check what the call returns.

\errorcode{E4095}{incompatible types \errhole{} and \errhole{}}
\label{sec:codes:e4095}

Raised during \textbf{validation}, from \texttt{ValidateErrorMessage::\allowbreak{}INCOMPATIBLE\_\allowbreak{}TYPE} (\textit{src/\allowbreak{}ymirc/\allowbreak{}semantic/\allowbreak{}validator/\allowbreak{}errors.yr}).

\subsubsection*{What it means}

Two types meet where they have to agree -- the two sides of an operator, a value and the variable it is assigned to, the branches of a conditional -- and neither converts to the other.

\subsubsection*{Example}

\textit{test\_\allowbreak{}resources/\allowbreak{}enumeration/\allowbreak{}test3.yr}:

%% check: skip
\begin{lstlisting}[style=coloredverbatimError]
enum X : f32
| A = 1;
\end{lstlisting}

\begin{lstlisting}[style=bashVerb, breaklines=true, escapechar=~]
Error[E4095] : incompatible types mut f32 and i32
 --> test_resources/enumeration/test3.yr:(2,1)
 2  ~\lstbox{\textSFxi{}}~ | A = 1;
    ~\lstbox{\textSFv{}}~ ^     ^
    ~\lstbox{\textSFii{}}~~\lstbox{\textSFx{}}~~\lstbox{\textSFx{}}~~\lstbox{\textSFx{}}~~\lstbox{\textSFx{}}~~\lstbox{\textSFx{}}~~\lstbox{\textSFx{}}~~\lstbox{\textSFx{}}~
\end{lstlisting}

\subsubsection*{How to fix it}

Convert one of them explicitly with \tokennolst{cast}, or change one of the two declarations. Mutability is part of the type, so two types differing only in it are also incompatible here.

\errorcode{E4096}{incompatible values \errhole{} and \errhole{}}
\label{sec:codes:e4096}

Raised during \textbf{validation}, from \texttt{ValidateErrorMessage::\allowbreak{}INCOMPATIBLE\_\allowbreak{}VALUES} (\textit{src/\allowbreak{}ymirc/\allowbreak{}semantic/\allowbreak{}validator/\allowbreak{}errors.yr}).

\subsubsection*{What it means}

Two values meet where they have to agree and cannot be reconciled, cf. \hyperref[sec:codes:e4095]{E4095}, the difference being in the values rather than in their types.

\subsubsection*{Example}

\textit{test\_\allowbreak{}resources/\allowbreak{}templates/\allowbreak{}implicit/\allowbreak{}test5.yr}:

%% check: skip
\begin{lstlisting}[style=coloredverbatimError]
fn foo {U, N : usize, J = N * 3} (a : [U ; N], z : [U ; J]) {
    a;
    z;
}

fn main () {
    foo ([1, 2], [1, 2, 3, 4]);
}
\end{lstlisting}

\begin{lstlisting}[style=bashVerb, breaklines=true, escapechar=~]
Error[E4096] : incompatible values 6us and 4us
 --> test_resources/templates/implicit/test5.yr:(2,29)
 2  ~\lstbox{\textSFxi{}}~ fn foo {U, N : usize, J = N * 3} (a : [U ; N], z : [U ; J]) {
    ~\lstbox{\textSFv{}}~                             ^
    ~\lstbox{\textSFii{}}~~\lstbox{\textSFx{}}~~\lstbox{\textSFx{}}~~\lstbox{\textSFx{}}~~\lstbox{\textSFx{}}~~\lstbox{\textSFx{}}~~\lstbox{\textSFx{}}~~\lstbox{\textSFx{}}~
\end{lstlisting}

\subsubsection*{How to fix it}

Make the two values agree.

\errorcode{E4097}{the type \errhole{} is not complete}
\label{sec:codes:e4097}

Raised during \textbf{validation}, from \texttt{ValidateErrorMessage::\allowbreak{}INCOMPLETE\_\allowbreak{}TYPE} (\textit{src/\allowbreak{}ymirc/\allowbreak{}semantic/\allowbreak{}validator/\allowbreak{}errors.yr}).

\subsubsection*{What it means}

A type is used where its size or its layout is needed, and it is not finished being validated.

\subsubsection*{Example}

\textit{test\_\allowbreak{}resources/\allowbreak{}lit\_\allowbreak{}class/\allowbreak{}ctors/\allowbreak{}test26.yr}:

%% check: skip
\begin{lstlisting}[style=coloredverbatimError]
fn foo (_ : &A)-> i32;
fn bar (_ : &A)-> i32;

class A {
    let x : i32;

    pub self ()
        with x = foo (self)
    {}
}

class B over A {
    let z : i32;

    pub self ()
        with z = bar (super)
    {}
}
\end{lstlisting}

\begin{lstlisting}[style=bashVerb, breaklines=true, escapechar=~]
Error[E4097] : the type &(test26::A)::fields is not complete
 --> test_resources/lit_class/ctors/test26.yr:(7,9)
 7  ~\lstbox{\textSFxi{}}~     pub self ()
    ~\lstbox{\textSFv{}}~         ^^^^  note: when creating the value self of &(test26::A)::fields
   ...
 --> test_resources/lit_class/ctors/test26.yr:(8,23)
 8  ~\lstbox{\textSFxi{}}~         with x = foo (self)
    ~\lstbox{\textSFv{}}~                       ^^^^
    ~\lstbox{\textSFii{}}~~\lstbox{\textSFx{}}~~\lstbox{\textSFx{}}~~\lstbox{\textSFx{}}~~\lstbox{\textSFx{}}~~\lstbox{\textSFx{}}~~\lstbox{\textSFx{}}~~\lstbox{\textSFx{}}~
\end{lstlisting}

\subsubsection*{How to fix it}

Break the cycle, cf. \hyperref[sec:codes:e4017]{E4017}.

\errorcode{E4217}{type \errhole{} has no fields}
\label{sec:codes:e4217}

Raised during \textbf{validation}, from \texttt{ValidateErrorMessage::\allowbreak{}TYPE\_\allowbreak{}HAS\_\allowbreak{}NO\_\allowbreak{}FIELDINFO} (\textit{src/\allowbreak{}ymirc/\allowbreak{}semantic/\allowbreak{}validator/\allowbreak{}errors.yr}).

\subsubsection*{What it means}

Field information was asked of a type that has no fields.

\subsubsection*{Example}

\textit{test\_\allowbreak{}resources/\allowbreak{}diagnostics/\allowbreak{}test72.yr}:

%% check: skip
\begin{lstlisting}[style=coloredverbatimError]
fn main () {
    let _ = __pragma!field_infos (i32);
}
\end{lstlisting}

\begin{lstlisting}[style=bashVerb, breaklines=true, escapechar=~]
Error[E4217] : type i32 has no fields
 --> test_resources/diagnostics/test72.yr:(2,13)
 2  ~\lstbox{\textSFxi{}}~     let _ = __pragma!field_infos (i32);
    ~\lstbox{\textSFv{}}~             ^^^^^^^^
    ~\lstbox{\textSFii{}}~~\lstbox{\textSFx{}}~~\lstbox{\textSFx{}}~~\lstbox{\textSFx{}}~~\lstbox{\textSFx{}}~~\lstbox{\textSFx{}}~~\lstbox{\textSFx{}}~~\lstbox{\textSFx{}}~
\end{lstlisting}

\subsubsection*{How to fix it}

Ask it of a record, an entity or a class.

\errorcode{E4218}{temporary type \errhole{} has no size}
\label{sec:codes:e4218}

Raised during \textbf{validation}, from \texttt{ValidateErrorMessage::\allowbreak{}TYPE\_\allowbreak{}HAS\_\allowbreak{}NO\_\allowbreak{}SIZE} (\textit{src/\allowbreak{}ymirc/\allowbreak{}semantic/\allowbreak{}validator/\allowbreak{}errors.yr}).

\subsubsection*{What it means}

The size of a temporary type was asked for, and it has none.

\subsubsection*{Example}

\textit{test\_\allowbreak{}resources/\allowbreak{}struct/\allowbreak{}test1.yr}:

%% check: skip
\begin{lstlisting}[style=coloredverbatimError]
in test1;

record A {
    let a : A;
}

record B {
    let c : C;
}

record C {
    let b : B;
}
\end{lstlisting}

\begin{lstlisting}[style=bashVerb, breaklines=true, escapechar=~]
Error[E4218] : temporary type error has no size
 --> test_resources/struct/test1.yr:(8,9)
 8  ~\lstbox{\textSFxi{}}~     let c : C;
    ~\lstbox{\textSFv{}}~         ^
    ~\lstbox{\textSFxi{}}~ Note : for the field c
    ~\lstbox{\textSFxi{}}~  --> test_resources/struct/test1.yr:(8,9)
    ~\lstbox{\textSFxi{}}~  8  ~\lstbox{\textSFxi{}}~     let c : C;
    ~\lstbox{\textSFxi{}}~     ~\lstbox{\textSFv{}}~         ^
    ~\lstbox{\textSFxi{}}~     ~\lstbox{\textSFxi{}}~ Note : within type test1::B
    ~\lstbox{\textSFxi{}}~     ~\lstbox{\textSFxi{}}~  --> test_resources/struct/test1.yr:(7,8)
    ~\lstbox{\textSFxi{}}~     ~\lstbox{\textSFxi{}}~  7  ~\lstbox{\textSFxi{}}~ record B {
    ~\lstbox{\textSFxi{}}~     ~\lstbox{\textSFxi{}}~     ~\lstbox{\textSFv{}}~        ^
    ~\lstbox{\textSFxi{}}~     ~\lstbox{\textSFxi{}}~     ~\lstbox{\textSFii{}}~~\lstbox{\textSFx{}}~~\lstbox{\textSFx{}}~~\lstbox{\textSFx{}}~~\lstbox{\textSFx{}}~~\lstbox{\textSFx{}}~~\lstbox{\textSFx{}}~~\lstbox{\textSFx{}}~
    ~\lstbox{\textSFxi{}}~     ~\lstbox{\textSFii{}}~~\lstbox{\textSFx{}}~~\lstbox{\textSFx{}}~~\lstbox{\textSFx{}}~~\lstbox{\textSFx{}}~~\lstbox{\textSFx{}}~~\lstbox{\textSFvii{}}~~\lstbox{\textSFx{}}~
    ~\lstbox{\textSFxi{}}~ Note : for the field b
    ~\lstbox{\textSFxi{}}~  --> test_resources/struct/test1.yr:(12,9)
    ~\lstbox{\textSFxi{}}~ 12  ~\lstbox{\textSFxi{}}~     let b : B;
    ~\lstbox{\textSFxi{}}~     ~\lstbox{\textSFv{}}~         ^
    ~\lstbox{\textSFxi{}}~     ~\lstbox{\textSFxi{}}~ Note : within type test1::C
    ~\lstbox{\textSFxi{}}~     ~\lstbox{\textSFxi{}}~  --> test_resources/struct/test1.yr:(11,8)
    ~\lstbox{\textSFxi{}}~     ~\lstbox{\textSFxi{}}~ 11  ~\lstbox{\textSFxi{}}~ record C {
    ~\lstbox{\textSFxi{}}~     ~\lstbox{\textSFxi{}}~     ~\lstbox{\textSFv{}}~        ^
    ~\lstbox{\textSFxi{}}~     ~\lstbox{\textSFxi{}}~     ~\lstbox{\textSFii{}}~~\lstbox{\textSFx{}}~~\lstbox{\textSFx{}}~~\lstbox{\textSFx{}}~~\lstbox{\textSFx{}}~~\lstbox{\textSFx{}}~~\lstbox{\textSFx{}}~~\lstbox{\textSFx{}}~
    ~\lstbox{\textSFxi{}}~     ~\lstbox{\textSFii{}}~~\lstbox{\textSFx{}}~~\lstbox{\textSFx{}}~~\lstbox{\textSFx{}}~~\lstbox{\textSFx{}}~~\lstbox{\textSFx{}}~~\lstbox{\textSFvii{}}~~\lstbox{\textSFx{}}~
    ~\lstbox{\textSFii{}}~~\lstbox{\textSFx{}}~~\lstbox{\textSFx{}}~~\lstbox{\textSFx{}}~~\lstbox{\textSFx{}}~~\lstbox{\textSFx{}}~~\lstbox{\textSFvii{}}~~\lstbox{\textSFx{}}~
\end{lstlisting}

\subsubsection*{How to fix it}

Ask the size of a concrete type instead.

\errorcode{E4220}{type \errhole{} has no field named \errhole{}}
\label{sec:codes:e4220}

Raised during \textbf{validation}, from \texttt{ValidateErrorMessage::\allowbreak{}TYPE\_\allowbreak{}NO\_\allowbreak{}FIELD} (\textit{src/\allowbreak{}ymirc/\allowbreak{}semantic/\allowbreak{}validator/\allowbreak{}errors.yr}).

\subsubsection*{What it means}

A field name was used on a type that has no field of that name, cf. \hyperref[sec:codes:e4019]{E4019}.

\subsubsection*{Example}

\textit{test\_\allowbreak{}resources/\allowbreak{}diagnostics/\allowbreak{}test73.yr}:

%% check: skip
\begin{lstlisting}[style=coloredverbatimError]
fn main () {
    let a = [1, 2, 3];
    let _ = a.nosuchfield;
}
\end{lstlisting}

\begin{lstlisting}[style=bashVerb, breaklines=true, escapechar=~]
Error[E4220] : type [i32 ; 3us] has no field named nosuchfield
 --> test_resources/diagnostics/test73.yr:(3,15)
 3  ~\lstbox{\textSFxi{}}~     let _ = a.nosuchfield;
    ~\lstbox{\textSFv{}}~               ^^^^^^^^^^^
    ~\lstbox{\textSFii{}}~~\lstbox{\textSFx{}}~~\lstbox{\textSFx{}}~~\lstbox{\textSFx{}}~~\lstbox{\textSFx{}}~~\lstbox{\textSFx{}}~~\lstbox{\textSFx{}}~~\lstbox{\textSFx{}}~
\end{lstlisting}

\subsubsection*{How to fix it}

Check the spelling and the type.

\errorcode{E4227}{cannot cast value of type \errhole{} into a value of type \errhole{}}
\label{sec:codes:e4227}

Raised during \textbf{validation}, from \texttt{ValidateErrorMessage::\allowbreak{}UNDEFINED\_\allowbreak{}CAST\_\allowbreak{}OP} (\textit{src/\allowbreak{}ymirc/\allowbreak{}semantic/\allowbreak{}validator/\allowbreak{}errors.yr}).

\subsubsection*{What it means}

A \tokennolst{cast} was asked for between two types with no conversion between them.

\subsubsection*{Example}

\textit{test\_\allowbreak{}resources/\allowbreak{}char\_\allowbreak{}type/\allowbreak{}test3.yr}:

%% check: skip
\begin{lstlisting}[style=coloredverbatimError]
in test3;

fn main () {
    let a = 'a'c8;
    let b = 'b'c16;
    let c = 'c'c32;

    cast!u16 (a);
    cast!u8 (b);
    cast!u8 (c);
    cast!u32 (a);

    cast!c16 (12u8);
    cast!c32 (12u16);
    cast!c8 (12u32);
}
\end{lstlisting}

\begin{lstlisting}[style=bashVerb, breaklines=true, escapechar=~]
Error[E4227] : cannot cast value of type c8 into a value of type u16
 --> test_resources/char_type/test3.yr:(8,5)
 8  ~\lstbox{\textSFxi{}}~     cast!u16 (a);
    ~\lstbox{\textSFv{}}~     ^^^^
    ~\lstbox{\textSFii{}}~~\lstbox{\textSFx{}}~~\lstbox{\textSFx{}}~~\lstbox{\textSFx{}}~~\lstbox{\textSFx{}}~~\lstbox{\textSFx{}}~~\lstbox{\textSFx{}}~~\lstbox{\textSFx{}}~
\end{lstlisting}

\subsubsection*{How to fix it}

Go through a type both convert to, or define the cast operator on the source type.

\errorcode{E4270}{expression produces a value, but is used as a type}
\label{sec:codes:e4270}

Raised during \textbf{validation}, from \texttt{ValidateErrorMessage::\allowbreak{}USE\_\allowbreak{}AS\_\allowbreak{}TYPE} (\textit{src/\allowbreak{}ymirc/\allowbreak{}semantic/\allowbreak{}validator/\allowbreak{}errors.yr}).

\subsubsection*{What it means}

An expression producing a value was written where a type is expected.

\subsubsection*{Example}

\textit{test\_\allowbreak{}resources/\allowbreak{}diagnostics/\allowbreak{}test82.yr}:

%% check: skip
\begin{lstlisting}[style=coloredverbatimError]
fn main () {
    let _ : 12 = 1;
}
\end{lstlisting}

\begin{lstlisting}[style=bashVerb, breaklines=true, escapechar=~]
Error[E4270] : expression produces a value, but is used as a type
 --> test_resources/diagnostics/test82.yr:(2,13)
 2  ~\lstbox{\textSFxi{}}~     let _ : 12 = 1;
    ~\lstbox{\textSFv{}}~             ^^
    ~\lstbox{\textSFii{}}~~\lstbox{\textSFx{}}~~\lstbox{\textSFx{}}~~\lstbox{\textSFx{}}~~\lstbox{\textSFx{}}~~\lstbox{\textSFx{}}~~\lstbox{\textSFx{}}~~\lstbox{\textSFx{}}~
\end{lstlisting}

\subsubsection*{How to fix it}

Write the type. When the two are spelled the same, the name resolves to the value; qualify it so it resolves to the type.

\errorcode{E4272}{expression produces a value of type \errhole{}, but is used as a type}
\label{sec:codes:e4272}

Raised during \textbf{validation}, from \texttt{ValidateErrorMessage::\allowbreak{}USE\_\allowbreak{}AS\_\allowbreak{}TYPE\_\allowbreak{}VAL} (\textit{src/\allowbreak{}ymirc/\allowbreak{}semantic/\allowbreak{}validator/\allowbreak{}errors.yr}).

\subsubsection*{What it means}

An expression produces a value of the reported type where a type is expected, cf. \hyperref[sec:codes:e4270]{E4270}.

\subsubsection*{Example}

\textit{test\_\allowbreak{}resources/\allowbreak{}diagnostics/\allowbreak{}test126.yr}:

%% check: skip
\begin{lstlisting}[style=coloredverbatimError]
// a function name used as a type
fn nothing () {}

fn main () {
    let _x_ : nothing = 1;
}
\end{lstlisting}

\begin{lstlisting}[style=bashVerb, breaklines=true, escapechar=~]
Error[E4272] : expression produces a value of type test126::nothing, but is used as a type
 --> test_resources/diagnostics/test126.yr:(5,15)
 5  ~\lstbox{\textSFxi{}}~     let _x_ : nothing = 1;
    ~\lstbox{\textSFv{}}~               ^^^^^^^
    ~\lstbox{\textSFxi{}}~
    = when validating test126::main -- (test_resources/diagnostics/test126.yr:(4,4))
    ~\lstbox{\textSFii{}}~~\lstbox{\textSFx{}}~~\lstbox{\textSFx{}}~~\lstbox{\textSFx{}}~~\lstbox{\textSFx{}}~~\lstbox{\textSFx{}}~~\lstbox{\textSFx{}}~~\lstbox{\textSFx{}}~
\end{lstlisting}

\subsubsection*{How to fix it}

Write the type itself.

\errorcode{E4273}{expression produces a type, but is used as a value}
\label{sec:codes:e4273}

Raised during \textbf{validation}, from \texttt{ValidateErrorMessage::\allowbreak{}USE\_\allowbreak{}AS\_\allowbreak{}VALUE} (\textit{src/\allowbreak{}ymirc/\allowbreak{}semantic/\allowbreak{}validator/\allowbreak{}errors.yr}).

\subsubsection*{What it means}

A type was written where a value is expected.

\subsubsection*{Example}

\textit{test\_\allowbreak{}resources/\allowbreak{}diagnostics/\allowbreak{}test127.yr}:

%% check: skip
\begin{lstlisting}[style=coloredverbatimError]
// a function pointer type used as a value
fn main () {
    let _ = fn (i32)-> i32;
}
\end{lstlisting}

\begin{lstlisting}[style=bashVerb, breaklines=true, escapechar=~]
Error[E4273] : expression produces a type, but is used as a value
 --> test_resources/diagnostics/test127.yr:(3,13)
 3  ~\lstbox{\textSFxi{}}~     let _ = fn (i32)-> i32;
    ~\lstbox{\textSFv{}}~             ^^
    ~\lstbox{\textSFxi{}}~
    = when validating test127::main -- (test_resources/diagnostics/test127.yr:(2,4))
    ~\lstbox{\textSFii{}}~~\lstbox{\textSFx{}}~~\lstbox{\textSFx{}}~~\lstbox{\textSFx{}}~~\lstbox{\textSFx{}}~~\lstbox{\textSFx{}}~~\lstbox{\textSFx{}}~~\lstbox{\textSFx{}}~
\end{lstlisting}

\subsubsection*{How to fix it}

Construct a value of that type, or use a value that already exists.

\errorcode{E4275}{expression produces the type \errhole{}, but is used as a value}
\label{sec:codes:e4275}

Raised during \textbf{validation}, from \texttt{ValidateErrorMessage::\allowbreak{}USE\_\allowbreak{}AS\_\allowbreak{}VALUE\_\allowbreak{}TYPE} (\textit{src/\allowbreak{}ymirc/\allowbreak{}semantic/\allowbreak{}validator/\allowbreak{}errors.yr}).

\subsubsection*{What it means}

An expression produces the reported type where a value is expected, cf. \hyperref[sec:codes:e4273]{E4273}.

\subsubsection*{Example}

\textit{test\_\allowbreak{}resources/\allowbreak{}diagnostics/\allowbreak{}test137.yr}:

%% check: skip
\begin{lstlisting}[style=coloredverbatimError]
// a trait used as a value
trait Shape {}

fn main () {
    let _ = Shape;
}
\end{lstlisting}

\begin{lstlisting}[style=bashVerb, breaklines=true, escapechar=~]
Error[E4275] : expression produces the type test137::Shape, but is used as a value
 --> test_resources/diagnostics/test137.yr:(5,13)
 5  ~\lstbox{\textSFxi{}}~     let _ = Shape;
    ~\lstbox{\textSFv{}}~             ^^^^^
    ~\lstbox{\textSFxi{}}~
    = when validating test137::main -- (test_resources/diagnostics/test137.yr:(4,4))
    ~\lstbox{\textSFii{}}~~\lstbox{\textSFx{}}~~\lstbox{\textSFx{}}~~\lstbox{\textSFx{}}~~\lstbox{\textSFx{}}~~\lstbox{\textSFx{}}~~\lstbox{\textSFx{}}~~\lstbox{\textSFx{}}~
\end{lstlisting}

\subsubsection*{How to fix it}

Construct a value of the type.

\errorcode{E4278}{void expression cannot be used as a value}
\label{sec:codes:e4278}

Raised during \textbf{validation}, from \texttt{ValidateErrorMessage::\allowbreak{}VOID\_\allowbreak{}VALUE} (\textit{src/\allowbreak{}ymirc/\allowbreak{}semantic/\allowbreak{}validator/\allowbreak{}errors.yr}).

\subsubsection*{What it means}

An expression producing nothing was used where a value is required.

\subsubsection*{Example}

\textit{test\_\allowbreak{}resources/\allowbreak{}diagnostics/\allowbreak{}test128.yr}:

%% check: skip
\begin{lstlisting}[style=coloredverbatimError]
// an array literal of void values
fn nothing () {}

fn main () {
    let _ = [nothing (), nothing ()];
}
\end{lstlisting}

\begin{lstlisting}[style=bashVerb, breaklines=true, escapechar=~]
Error[E4278] : void expression cannot be used as a value
 --> test_resources/diagnostics/test128.yr:(5,22)
 5  ~\lstbox{\textSFxi{}}~     let _ = [nothing (), nothing ()];
    ~\lstbox{\textSFv{}}~                      ^
    ~\lstbox{\textSFxi{}}~
    = when validating test128::main -- (test_resources/diagnostics/test128.yr:(4,4))
    ~\lstbox{\textSFii{}}~~\lstbox{\textSFx{}}~~\lstbox{\textSFx{}}~~\lstbox{\textSFx{}}~~\lstbox{\textSFx{}}~~\lstbox{\textSFx{}}~~\lstbox{\textSFx{}}~~\lstbox{\textSFx{}}~
\end{lstlisting}

\subsubsection*{How to fix it}

Use a call that returns something, or make the surrounding construct a statement.

\errorcode{E4304}{\errhole{} has no type of its own, it takes the type of what it is assigned to}
\label{sec:codes:e4304}

Raised during \textbf{validation}, from \texttt{ValidateErrorMessage::\allowbreak{}TYPEOF\_\allowbreak{}UNTYPED} (\textit{src/\allowbreak{}ymirc/\allowbreak{}semantic/\allowbreak{}validator/\allowbreak{}errors.yr}).

\subsubsection*{What it means}

\tokennolst{typeof} is applied to an untyped \tokennolst{null} or \tokennolst{none}. Neither literal has a complete type: \tokennolst{null} becomes the pointer type it is assigned to, and \tokennolst{none} the option type it is assigned to, so there is no type for \tokennolst{typeof} to name until then.

\subsubsection*{Example}

\textit{test\_\allowbreak{}resources/\allowbreak{}diagnostics/\allowbreak{}test148.yr}:

%% check: skip
\begin{lstlisting}[style=coloredverbatimError]
fn main () {
    let _ : typeof (null) = null;
}
\end{lstlisting}

\begin{lstlisting}[style=bashVerb, breaklines=true, escapechar=~]
Error[E4304] : null has no type of its own, it takes the type of what it is assigned to
 --> test_resources/diagnostics/test148.yr:(2,21)
 2  ~\lstbox{\textSFxi{}}~     let _ : typeof (null) = null;
    ~\lstbox{\textSFv{}}~                     ^^^^
    ~\lstbox{\textSFxi{}}~
    = when validating test148::main -- (test_resources/diagnostics/test148.yr:(1,4))
    ~\lstbox{\textSFii{}}~~\lstbox{\textSFx{}}~~\lstbox{\textSFx{}}~~\lstbox{\textSFx{}}~~\lstbox{\textSFx{}}~~\lstbox{\textSFx{}}~~\lstbox{\textSFx{}}~~\lstbox{\textSFx{}}~
\end{lstlisting}

\subsubsection*{How to fix it}

Name the type the literal is meant to have, or take the type of a value that already has it:

%% check: skip
\begin{lstlisting}[style=coloredverbatim]
let x : *i32 = null;
let _ : typeof (x) = null;
\end{lstlisting}

To recognize the \tokennolst{null} literal in a macro, match it with the \tokennolst{\_\_key} rule rather than comparing its type.

%% Generated by tools/gen_error_codes.py from docs/errors/ of the compiler;
%% edit the compiler's pages or tools/error_themes.txt, then rerun it.

\clearpage
\section{Operators and indexing}
\label{sec:codes:operators}

\begin{longtable}{@{}l p{0.78\linewidth} r@{}}
\toprule Code & Message & Page\\ \midrule \endhead
\hyperref[sec:codes:e4005]{E4005} & slice access, index overflow \errhole{} (len) \textless{}= \errhole{} (index) & \pageref{sec:codes:e4005}\\
\hyperref[sec:codes:e4040]{E4040} & range index cannot be decreasing for type \errhole{}, but goes from \errhole{} to \errhole{} & \pageref{sec:codes:e4040}\\
\hyperref[sec:codes:e4047]{E4047} & '\$' is usable only inside an index operation & \pageref{sec:codes:e4047}\\
\hyperref[sec:codes:e4142]{E4142} & index cannot be negative for type \errhole{}, but is \errhole{} & \pageref{sec:codes:e4142}\\
\hyperref[sec:codes:e4195]{E4195} & the index operator on \errhole{} with a dynamic operand of type \errhole{} is allowed only in the context of a copy statement & \pageref{sec:codes:e4195}\\
\hyperref[sec:codes:e4222]{E4222} & undefined operator \errhole{} for type \errhole{} and field \errhole{} & \pageref{sec:codes:e4222}\\
\hyperref[sec:codes:e4223]{E4223} & undefined module operator \errhole{} for \errhole{} and field \errhole{} & \pageref{sec:codes:e4223}\\
\hyperref[sec:codes:e4224]{E4224} & undefined operator \errhole{} for types \errhole{} and \errhole{} & \pageref{sec:codes:e4224}\\
\hyperref[sec:codes:e4229]{E4229} & operator \$ is not defined for type \errhole{} & \pageref{sec:codes:e4229}\\
\hyperref[sec:codes:e4231]{E4231} & the index operator is not defined for type \errhole{} and \errhole{} & \pageref{sec:codes:e4231}\\
\hyperref[sec:codes:e4237]{E4237} & undefined operator \errhole{} for type \errhole{} & \pageref{sec:codes:e4237}\\
\hyperref[sec:codes:e5011]{E5011} & slice access always overflows, \errhole{} (len) \errhole{} \errhole{} (index) & \pageref{sec:codes:e5011}\\
\bottomrule
\end{longtable}

\errorcode{E4005}{slice access, index overflow \errhole{} (len) \textless{}= \errhole{} (index)}
\label{sec:codes:e4005}

Raised during \textbf{validation}, from \texttt{ValidateErrorMessage::\allowbreak{}ARRAY\_\allowbreak{}OVERFLOW} (\textit{src/\allowbreak{}ymirc/\allowbreak{}semantic/\allowbreak{}validator/\allowbreak{}errors.yr}).

\subsubsection*{What it means}

A slice is indexed at a position the compiler could compute, and that position is outside the slice. Both the length and the index are in the message.

\subsubsection*{Example}

\textit{test\_\allowbreak{}resources/\allowbreak{}lit\_\allowbreak{}macro/\allowbreak{}test29.yr}:

%% check: skip
\begin{lstlisting}[style=coloredverbatimError]
/*
 * An indexed access of a repeated capture, that is out of the bounds of
 * what the call captured
 */

macro pick {
    self (v=(a=__operand)*) skips (" ") #{
        #(v[5]::a)
    }
}

fn foo ()-> i32 {
    pick#{ 1 2 }
}
\end{lstlisting}

\begin{lstlisting}[style=bashVerb, breaklines=true, escapechar=~]
Error[E4005] : slice access, index overflow 2 (len) <= 5 (index)
 --> :(1,8)
 1  ~\lstbox{\textSFxi{}}~     #(v[5]::a)
    ~\lstbox{\textSFv{}}~        ^
    ~\lstbox{\textSFii{}}~~\lstbox{\textSFx{}}~~\lstbox{\textSFx{}}~~\lstbox{\textSFx{}}~~\lstbox{\textSFx{}}~~\lstbox{\textSFx{}}~~\lstbox{\textSFx{}}~~\lstbox{\textSFx{}}~
\end{lstlisting}

\subsubsection*{How to fix it}

Correct the index, or check the length at run time before indexing when it is not known statically.

\errorcode{E4040}{range index cannot be decreasing for type \errhole{}, but goes from \errhole{} to \errhole{}}
\label{sec:codes:e4040}

Raised during \textbf{validation}, from \texttt{ValidateErrorMessage::\allowbreak{}DECREASING\_\allowbreak{}RANGE\_\allowbreak{}ACCESS} (\textit{src/\allowbreak{}ymirc/\allowbreak{}semantic/\allowbreak{}validator/\allowbreak{}errors.yr}).

\subsubsection*{What it means}

A range index goes from a higher bound to a lower one on a type whose indices only increase.

\subsubsection*{Example}

\textit{test\_\allowbreak{}resources/\allowbreak{}diagnostics/\allowbreak{}test26.yr}:

%% check: skip
\begin{lstlisting}[style=coloredverbatimError]
fn main () {
    let a = [1, 2, 3];
    let _ = a [2 .. 0];
}
\end{lstlisting}

\begin{lstlisting}[style=bashVerb, breaklines=true, escapechar=~]
Error[E4040] : range index cannot be decreasing for type [i32 ; 3us], but goes from 2 to 0
 --> test_resources/diagnostics/test26.yr:(3,15)
 3  ~\lstbox{\textSFxi{}}~     let _ = a [2 .. 0];
    ~\lstbox{\textSFv{}}~               ^  ^^
    ~\lstbox{\textSFii{}}~~\lstbox{\textSFx{}}~~\lstbox{\textSFx{}}~~\lstbox{\textSFx{}}~~\lstbox{\textSFx{}}~~\lstbox{\textSFx{}}~~\lstbox{\textSFx{}}~~\lstbox{\textSFx{}}~
\end{lstlisting}

\subsubsection*{How to fix it}

Swap the bounds. To walk the values backwards, reverse the result rather than the range.

\errorcode{E4047}{'\$' is usable only inside an index operation}
\label{sec:codes:e4047}

Raised during \textbf{validation}, from \texttt{ValidateErrorMessage::\allowbreak{}DOLLAR\_\allowbreak{}OUTSIDE\_\allowbreak{}CONTEXT} (\textit{src/\allowbreak{}ymirc/\allowbreak{}semantic/\allowbreak{}validator/\allowbreak{}errors.yr}).

\subsubsection*{What it means}

\tokennolst{\$}, which stands for the length of the value being indexed, was used outside an index operation, where it has nothing to refer to.

\subsubsection*{Example}

\textit{test\_\allowbreak{}resources/\allowbreak{}diagnostics/\allowbreak{}test29.yr}:

%% check: skip
\begin{lstlisting}[style=coloredverbatimError]
fn main () {
    let _ = $;
}
\end{lstlisting}

\begin{lstlisting}[style=bashVerb, breaklines=true, escapechar=~]
Error[E4047] : '$' is usable only inside an index operation
 --> test_resources/diagnostics/test29.yr:(2,13)
 2  ~\lstbox{\textSFxi{}}~     let _ = $;
    ~\lstbox{\textSFv{}}~             ^
    ~\lstbox{\textSFii{}}~~\lstbox{\textSFx{}}~~\lstbox{\textSFx{}}~~\lstbox{\textSFx{}}~~\lstbox{\textSFx{}}~~\lstbox{\textSFx{}}~~\lstbox{\textSFx{}}~~\lstbox{\textSFx{}}~
\end{lstlisting}

\subsubsection*{How to fix it}

Use the length explicitly (\tokennolst{x.len}) outside an index.

\errorcode{E4142}{index cannot be negative for type \errhole{}, but is \errhole{}}
\label{sec:codes:e4142}

Raised during \textbf{validation}, from \texttt{ValidateErrorMessage::\allowbreak{}NEGATIVE\_\allowbreak{}INT\_\allowbreak{}INDEX} (\textit{src/\allowbreak{}ymirc/\allowbreak{}semantic/\allowbreak{}validator/\allowbreak{}errors.yr}).

\subsubsection*{What it means}

An index the compiler computed is negative, on a type whose indices start at zero.

\subsubsection*{Example}

\textit{test\_\allowbreak{}resources/\allowbreak{}diagnostics/\allowbreak{}test56.yr}:

%% check: skip
\begin{lstlisting}[style=coloredverbatimError]
fn main () {
    let a = [1, 2, 3];
    let _ = a [-1];
}
\end{lstlisting}

\begin{lstlisting}[style=bashVerb, breaklines=true, escapechar=~]
Error[E4142] : index cannot be negative for type [i32 ; 3us], but is -1
 --> test_resources/diagnostics/test56.yr:(3,15)
 3  ~\lstbox{\textSFxi{}}~     let _ = a [-1];
    ~\lstbox{\textSFv{}}~               ^^
    ~\lstbox{\textSFii{}}~~\lstbox{\textSFx{}}~~\lstbox{\textSFx{}}~~\lstbox{\textSFx{}}~~\lstbox{\textSFx{}}~~\lstbox{\textSFx{}}~~\lstbox{\textSFx{}}~~\lstbox{\textSFx{}}~
\end{lstlisting}

\subsubsection*{How to fix it}

Correct the index. To count from the end, use \tokennolst{\$} inside the index operation.

\errorcode{E4195}{the index operator on \errhole{} with a dynamic operand of type \errhole{} is allowed only in the context of a copy statement}
\label{sec:codes:e4195}

Raised during \textbf{validation}, from \texttt{ValidateErrorMessage::\allowbreak{}RANGE\_\allowbreak{}ON\_\allowbreak{}ARRAY\_\allowbreak{}NO\_\allowbreak{}COPY} (\textit{src/\allowbreak{}ymirc/\allowbreak{}semantic/\allowbreak{}validator/\allowbreak{}errors.yr}).

\subsubsection*{What it means}

A range index whose bounds are not known at compile time was applied outside a copy. Such an index produces a new value rather than a view, which only the copy context accepts.

\subsubsection*{Example}

\textit{test\_\allowbreak{}resources/\allowbreak{}diagnostics/\allowbreak{}test122.yr}:

%% check: skip
\begin{lstlisting}[style=coloredverbatimError]
// a range index with a runtime bound on a static array, outside a copy
fn main () {
    let a = [1, 2, 3];
    let mut i = 0us;
    i = 1us;
    let _ = a [i .. 2us];
}
\end{lstlisting}

\begin{lstlisting}[style=bashVerb, breaklines=true, escapechar=~]
Error[E4195] : the index operator on [i32 ; 3us] with a dynamic operand of type (..usize) is allowed only in the context of a copy statement
 --> test_resources/diagnostics/test122.yr:(6,15)
 6  ~\lstbox{\textSFxi{}}~     let _ = a [i .. 2us];
    ~\lstbox{\textSFv{}}~               ^  error: value of type (..usize) is needed but unknown at compilation time
    ~\lstbox{\textSFxi{}}~
    = when validating test122::main -- (test_resources/diagnostics/test122.yr:(2,4))
    ~\lstbox{\textSFii{}}~~\lstbox{\textSFx{}}~~\lstbox{\textSFx{}}~~\lstbox{\textSFx{}}~~\lstbox{\textSFx{}}~~\lstbox{\textSFx{}}~~\lstbox{\textSFx{}}~~\lstbox{\textSFx{}}~
\end{lstlisting}

\subsubsection*{How to fix it}

Enclose the index in a \tokennolst{copy}, or use bounds the compiler can compute.

\errorcode{E4222}{undefined operator \errhole{} for type \errhole{} and field \errhole{}}
\label{sec:codes:e4222}

Raised during \textbf{validation}, from \texttt{ValidateErrorMessage::\allowbreak{}UNDEFINED\_\allowbreak{}BIN\_\allowbreak{}ACC\_\allowbreak{}OP} (\textit{src/\allowbreak{}ymirc/\allowbreak{}semantic/\allowbreak{}validator/\allowbreak{}errors.yr}).

\subsubsection*{What it means}

A binary operator is not defined between the type and the field named.

\subsubsection*{Example}

\textit{test\_\allowbreak{}resources/\allowbreak{}regressions/\allowbreak{}float\_\allowbreak{}properties/\allowbreak{}test2.yr}:

%% check: skip
\begin{lstlisting}[style=coloredverbatimError]
in test2;

fn main () {

    f32::foo;
    f64::bar;
    f80::baz;
    fsize::bad;


}
\end{lstlisting}

\begin{lstlisting}[style=bashVerb, breaklines=true, escapechar=~]
Error[E4222] : undefined operator :: for type f32 and field foo
 --> test_resources/regressions/float_properties/test2.yr:(5,8)
 5  ~\lstbox{\textSFxi{}}~     f32::foo;
    ~\lstbox{\textSFv{}}~        ^^
    ~\lstbox{\textSFii{}}~~\lstbox{\textSFx{}}~~\lstbox{\textSFx{}}~~\lstbox{\textSFx{}}~~\lstbox{\textSFx{}}~~\lstbox{\textSFx{}}~~\lstbox{\textSFx{}}~~\lstbox{\textSFx{}}~
\end{lstlisting}

\subsubsection*{Notes}

When the field is \tokennolst{typeinfo}, the reason is reported under this diagnostic, as a note carrying no code of its own:

\begin{itemize}
\item \texttt{temporary type \textless{}\textgreater{} has no type information} (formerly E4219): \tokennolst{typeinfo} was requested on a type that has none, such as a trait.
\end{itemize}

\subsubsection*{How to fix it}

Check the field: a field method may be needed, or the operator may have to be defined on the field type.

\errorcode{E4223}{undefined module operator \errhole{} for \errhole{} and field \errhole{}}
\label{sec:codes:e4223}

Raised during \textbf{validation}, from \texttt{ValidateErrorMessage::\allowbreak{}UNDEFINED\_\allowbreak{}BIN\_\allowbreak{}MOD\_\allowbreak{}OP} (\textit{src/\allowbreak{}ymirc/\allowbreak{}semantic/\allowbreak{}validator/\allowbreak{}errors.yr}).

\subsubsection*{What it means}

A module-level binary operator is not defined for the operands given.

\subsubsection*{Example}

\textit{test\_\allowbreak{}resources/\allowbreak{}diagnostics/\allowbreak{}test62.yr}:

%% check: skip
\begin{lstlisting}[style=coloredverbatimError]
class Foo {
    pub self () {}
}

fn main () { let _ = copy Foo::named (); }
\end{lstlisting}

\begin{lstlisting}[style=bashVerb, breaklines=true, escapechar=~]
Error[E4223] : undefined module operator :: for test62::Foo::self and field named
 --> test_resources/diagnostics/test62.yr:(5,30)
 5  ~\lstbox{\textSFxi{}}~ fn main () { let _ = copy Foo::named (); }
    ~\lstbox{\textSFv{}}~                              ^^
    ~\lstbox{\textSFii{}}~~\lstbox{\textSFx{}}~~\lstbox{\textSFx{}}~~\lstbox{\textSFx{}}~~\lstbox{\textSFx{}}~~\lstbox{\textSFx{}}~~\lstbox{\textSFx{}}~~\lstbox{\textSFx{}}~
\end{lstlisting}

\subsubsection*{How to fix it}

Define the operator, or use one that exists for these types.

\errorcode{E4224}{undefined operator \errhole{} for types \errhole{} and \errhole{}}
\label{sec:codes:e4224}

Raised during \textbf{validation}, from \texttt{ValidateErrorMessage::\allowbreak{}UNDEFINED\_\allowbreak{}BIN\_\allowbreak{}OP} (\textit{src/\allowbreak{}ymirc/\allowbreak{}semantic/\allowbreak{}validator/\allowbreak{}errors.yr}).

\subsubsection*{What it means}

A binary operator is not defined for the two operand types. Both are in the message.

\subsubsection*{Example}

\textit{test\_\allowbreak{}resources/\allowbreak{}lit\_\allowbreak{}class/\allowbreak{}operators/\allowbreak{}test6.yr}:

%% check: skip
\begin{lstlisting}[style=coloredverbatimError]
class A {
    pub self () {}
}

fn main () {
    let a = copy A ();
    a * a;

}
\end{lstlisting}

\begin{lstlisting}[style=bashVerb, breaklines=true, escapechar=~]
Error[E4224] : undefined operator * for types &(test6::A) and &(test6::A)
 --> test_resources/lit_class/operators/test6.yr:(7,7)
 7  ~\lstbox{\textSFxi{}}~     a * a;
    ~\lstbox{\textSFv{}}~       ^  error: class &(test6::A) has no method named opBinary
    ~\lstbox{\textSFxi{}}~       ~\lstbox{\textSFii{}}~~\lstbox{\textSFx{}}~ error: class &(test6::A) has no method named opBinaryRight
    ~\lstbox{\textSFxi{}}~
    = when validating test6::main -- (test_resources/lit_class/operators/test6.yr:(5,4))
    ~\lstbox{\textSFii{}}~~\lstbox{\textSFx{}}~~\lstbox{\textSFx{}}~~\lstbox{\textSFx{}}~~\lstbox{\textSFx{}}~~\lstbox{\textSFx{}}~~\lstbox{\textSFx{}}~~\lstbox{\textSFx{}}~
\end{lstlisting}

\subsubsection*{How to fix it}

Convert one operand, or define the operator on the type. When both types look the same, the difference is in their mutability, which is part of the type.

\errorcode{E4229}{operator \$ is not defined for type \errhole{}}
\label{sec:codes:e4229}

Raised during \textbf{validation}, from \texttt{ValidateErrorMessage::\allowbreak{}UNDEFINED\_\allowbreak{}DOLLAR\_\allowbreak{}OP} (\textit{src/\allowbreak{}ymirc/\allowbreak{}semantic/\allowbreak{}validator/\allowbreak{}errors.yr}).

\subsubsection*{What it means}

\tokennolst{\$} was used on a type that does not define a length.

\subsubsection*{Example}

\textit{test\_\allowbreak{}resources/\allowbreak{}diagnostics/\allowbreak{}test74.yr}:

%% check: skip
\begin{lstlisting}[style=coloredverbatimError]
class Foo { pub self () {} }

fn main () {
    let f = copy Foo ();
    let _ = f [$];
}
\end{lstlisting}

\begin{lstlisting}[style=bashVerb, breaklines=true, escapechar=~]
Error[E4229] : operator $ is not defined for type &(test74::Foo)
 --> test_resources/diagnostics/test74.yr:(5,16)
 5  ~\lstbox{\textSFxi{}}~     let _ = f [$];
    ~\lstbox{\textSFv{}}~                ^  error: class &(test74::Foo) has no method named opDollar
    ~\lstbox{\textSFxi{}}~
    = when validating test74::main -- (test_resources/diagnostics/test74.yr:(3,4))
    ~\lstbox{\textSFii{}}~~\lstbox{\textSFx{}}~~\lstbox{\textSFx{}}~~\lstbox{\textSFx{}}~~\lstbox{\textSFx{}}~~\lstbox{\textSFx{}}~~\lstbox{\textSFx{}}~~\lstbox{\textSFx{}}~
\end{lstlisting}

\subsubsection*{How to fix it}

Use a type with a length, or write the bound explicitly.

\errorcode{E4231}{the index operator is not defined for type \errhole{} and \errhole{}}
\label{sec:codes:e4231}

Raised during \textbf{validation}, from \texttt{ValidateErrorMessage::\allowbreak{}UNDEFINED\_\allowbreak{}INDEX\_\allowbreak{}OP} (\textit{src/\allowbreak{}ymirc/\allowbreak{}semantic/\allowbreak{}validator/\allowbreak{}errors.yr}).

\subsubsection*{What it means}

The index operator is not defined between the indexed type and the index type. Both are in the message.

\subsubsection*{Example}

\textit{test\_\allowbreak{}resources/\allowbreak{}aka/\allowbreak{}test2.yr}:

%% check: skip
\begin{lstlisting}[style=coloredverbatimError]
def A : dmut [i32];

fn main ()
{
    let dmut x : A = copy [1, 2, 3];
    x [0] = 1;
}
\end{lstlisting}

\begin{lstlisting}[style=bashVerb, breaklines=true, escapechar=~]
Error[E4231] : the index operator is not defined for type error and {i32}
 --> test_resources/aka/test2.yr:(6,7)
 6  ~\lstbox{\textSFxi{}}~     x [0] = 1;
    ~\lstbox{\textSFv{}}~       ^
    ~\lstbox{\textSFii{}}~~\lstbox{\textSFx{}}~~\lstbox{\textSFx{}}~~\lstbox{\textSFx{}}~~\lstbox{\textSFx{}}~~\lstbox{\textSFx{}}~~\lstbox{\textSFx{}}~~\lstbox{\textSFx{}}~
\end{lstlisting}

\subsubsection*{How to fix it}

Index with the type the operation expects, or define the index operator on the type.

\errorcode{E4237}{undefined operator \errhole{} for type \errhole{}}
\label{sec:codes:e4237}

Raised during \textbf{validation}, from \texttt{ValidateErrorMessage::\allowbreak{}UNDEFINED\_\allowbreak{}UN\_\allowbreak{}OP} (\textit{src/\allowbreak{}ymirc/\allowbreak{}semantic/\allowbreak{}validator/\allowbreak{}errors.yr}).

\subsubsection*{What it means}

A unary operator is not defined for the operand type.

\subsubsection*{Example}

\textit{test\_\allowbreak{}resources/\allowbreak{}char\_\allowbreak{}type/\allowbreak{}test4.yr}:

%% check: skip
\begin{lstlisting}[style=coloredverbatimError]
in test4;


fn main () {
    let a = 'r'c8;
    let b = 'r'c16;
    let c = 'r'c32;


    *a;
    !a;
    -a;
    &a;

    *b;
    !b;
    -b;
    &b;

    *c;
    !c;
    -c;
    &c;
}
\end{lstlisting}

\begin{lstlisting}[style=bashVerb, breaklines=true, escapechar=~]
Error[E4237] : undefined operator * for type c8
 --> test_resources/char_type/test4.yr:(10,5)
10  ~\lstbox{\textSFxi{}}~     *a;
    ~\lstbox{\textSFv{}}~     ^
    ~\lstbox{\textSFii{}}~~\lstbox{\textSFx{}}~~\lstbox{\textSFx{}}~~\lstbox{\textSFx{}}~~\lstbox{\textSFx{}}~~\lstbox{\textSFx{}}~~\lstbox{\textSFx{}}~~\lstbox{\textSFx{}}~
\end{lstlisting}

\subsubsection*{How to fix it}

Define the operator on the type, or convert the operand.

\errorcode{E5011}{slice access always overflows, \errhole{} (len) \errhole{} \errhole{} (index)}
\label{sec:codes:e5011}

Raised during \textbf{lowering (YIL)}, from \tokennolst{OptimizerErrorMessage::ALWAYS\_OVERFLOW} (\textit{src/\allowbreak{}ymirc/\allowbreak{}lint/\allowbreak{}optimizer/\allowbreak{}errors.yr}), by the bounds check pass of the optimizer.

\subsubsection*{What it means}

An access to a slice or an array is out of its bounds on every execution that reaches it: the check guarding it is always taken, and the program panics there.

The validator already rejects an index it can compute on a length it can compute (\hyperref[sec:codes:e4005]{E4005}). This one is found later, on what the control flow of the function says: the length of a mutable slice after the last assignment to it, or the length a slice was given on every path leading to the access. The optimizer only runs from \tokennolst{-O1}, so the diagnostic is never raised at \tokennolst{-O0}.

It is a warning: the access may sit on a path that never runs. Like any warning, it only lets the compilation carry on in debug mode. The check and the panic it guards are kept, so a program that reaches the access still panics there.

\subsubsection*{Example}

\textit{test\_\allowbreak{}resources/\allowbreak{}warnings/\allowbreak{}test31.yr}:

%% check: skip
\begin{lstlisting}[style=coloredverbatimError]
// an access the optimizer proves out of bounds warns, and panics as it did
fn foo ()-> i32 {
    let dmut a = copy [1, 2, 3];

    a [10]
}
\end{lstlisting}

\begin{lstlisting}[style=bashVerb, breaklines=true, escapechar=~]
Warning[E5011] : slice access always overflows, 3 (len) <= 10 (index)
 --> test_resources/warnings/test31.yr:(5,7)
 5  ~\lstbox{\textSFxi{}}~     a [10]
    ~\lstbox{\textSFv{}}~       ^--
    ~\lstbox{\textSFii{}}~~\lstbox{\textSFx{}}~~\lstbox{\textSFx{}}~~\lstbox{\textSFx{}}~~\lstbox{\textSFx{}}~~\lstbox{\textSFx{}}~~\lstbox{\textSFx{}}~~\lstbox{\textSFx{}}~
\end{lstlisting}

\subsubsection*{How to fix it}

Correct the index, or the length the slice is given before it is read.

%% Generated by tools/gen_error_codes.py from docs/errors/ of the compiler;
%% edit the compiler's pages or tools/error_themes.txt, then rerun it.

\clearpage
\section{Tuples, arrays, maps and options}
\label{sec:codes:compound}

The book's chapter \textit{Native compound types} specifies the rules behind these codes.

\begin{longtable}{@{}l p{0.78\linewidth} r@{}}
\toprule Code & Message & Page\\ \midrule \endhead
\hyperref[sec:codes:e4036]{E4036} & using a deep copy instead of a one level copy is prohibited when constructing a map \errhole{} & \pageref{sec:codes:e4036}\\
\hyperref[sec:codes:e4037]{E4037} & creating an instance of map \errhole{} on the stack is prohibited & \pageref{sec:codes:e4037}\\
\hyperref[sec:codes:e4121]{E4121} & the type \errhole{} cannot be used as a key, it is not comparable & \pageref{sec:codes:e4121}\\
\hyperref[sec:codes:e4122]{E4122} & the type \errhole{} cannot be used as a key, it is not hashable & \pageref{sec:codes:e4122}\\
\hyperref[sec:codes:e4131]{E4131} & mismatch tuple arity \errhole{} and \errhole{} & \pageref{sec:codes:e4131}\\
\hyperref[sec:codes:e4163]{E4163} & \errhole{} is not a tuple type & \pageref{sec:codes:e4163}\\
\hyperref[sec:codes:e4175]{E4175} & option of type \errhole{} has no error & \pageref{sec:codes:e4175}\\
\hyperref[sec:codes:e4176]{E4176} & option of type \errhole{} has no value & \pageref{sec:codes:e4176}\\
\hyperref[sec:codes:e4179]{E4179} & static array size \errhole{} exceeds limits \errhole{} & \pageref{sec:codes:e4179}\\
\hyperref[sec:codes:e4213]{E4213} & tuple left affectation is only defined for the operator \errhole{} & \pageref{sec:codes:e4213}\\
\hyperref[sec:codes:e4214]{E4214} & tuple access out of bound(\errhole{}), tuple arity is \errhole{} & \pageref{sec:codes:e4214}\\
\hyperref[sec:codes:e4267]{E4267} & operator \errhole{} is useless, the left operand of type \errhole{} always has a value & \pageref{sec:codes:e4267}\\
\bottomrule
\end{longtable}

\errorcode{E4036}{using a deep copy instead of a one level copy is prohibited when constructing a map \errhole{}}
\label{sec:codes:e4036}

Raised during \textbf{validation}, from \texttt{ValidateErrorMessage::\allowbreak{}CTOR\_\allowbreak{}MAP\_\allowbreak{}DCOPY} (\textit{src/\allowbreak{}ymirc/\allowbreak{}semantic/\allowbreak{}validator/\allowbreak{}errors.yr}).

\subsubsection*{What it means}

A map was constructed with \tokennolst{dcopy} where \tokennolst{copy} is required, cf. \hyperref[sec:codes:retired]{E4033}.

\subsubsection*{Example}

\textit{test\_\allowbreak{}resources/\allowbreak{}lit\_\allowbreak{}map/\allowbreak{}test12.yr}:

%% check: skip
\begin{lstlisting}[style=coloredverbatimError]
// a map is built one level deep, so a deep copy of a map literal is refused
fn main () {
    let _m_ = dcopy [1 => 2];
}
\end{lstlisting}

\begin{lstlisting}[style=bashVerb, breaklines=true, escapechar=~]
Error[E4036] : using a deep copy instead of a one level copy is prohibited when constructing a map mut [i32 => mut i32]
 --> test_resources/lit_map/test12.yr:(3,21)
 3  ~\lstbox{\textSFxi{}}~     let _m_ = dcopy [1 => 2];
    ~\lstbox{\textSFv{}}~                     ^
    ~\lstbox{\textSFxi{}}~
    = help: replace dcopy by the simple copy operator
    ~\lstbox{\textSFxi{}}~
 3  -     let _m_ = dcopy [1 => 2];
 3  +     let _m_ = copy [1 => 2];
    ~\lstbox{\textSFxi{}}~
    = when validating test12::main -- (test_resources/lit_map/test12.yr:(2,4))
    ~\lstbox{\textSFii{}}~~\lstbox{\textSFx{}}~~\lstbox{\textSFx{}}~~\lstbox{\textSFx{}}~~\lstbox{\textSFx{}}~~\lstbox{\textSFx{}}~~\lstbox{\textSFx{}}~~\lstbox{\textSFx{}}~
\end{lstlisting}

\subsubsection*{How to fix it}

Use \tokennolst{copy} to construct the map.

\errorcode{E4037}{creating an instance of map \errhole{} on the stack is prohibited}
\label{sec:codes:e4037}

Raised during \textbf{validation}, from \texttt{ValidateErrorMessage::\allowbreak{}CTOR\_\allowbreak{}MAP\_\allowbreak{}STACK} (\textit{src/\allowbreak{}ymirc/\allowbreak{}semantic/\allowbreak{}validator/\allowbreak{}errors.yr}).

\subsubsection*{What it means}

A map was created on the stack. Like classes, maps are reference types.

\subsubsection*{Example}

\textit{test\_\allowbreak{}resources/\allowbreak{}lit\_\allowbreak{}map/\allowbreak{}test3.yr}:

%% check: skip
\begin{lstlisting}[style=coloredverbatimError]
fn main () {
    let dmut _ = [1 => 2];
    let dmut _ = ["foo" => copy [1, 2, 3]];

    let dmut _ = ["foo" => [1, 2], "bar" => [1]];
    let dmut _ = ["foo" => 1, 2 => 3];
}
\end{lstlisting}

\begin{lstlisting}[style=bashVerb, breaklines=true, escapechar=~]
Error[E4037] : creating an instance of map mut [i32 => mut i32] on the stack is prohibited
 --> test_resources/lit_map/test3.yr:(2,18)
 2  ~\lstbox{\textSFxi{}}~     let dmut _ = [1 => 2];
    ~\lstbox{\textSFv{}}~                  ^
    ~\lstbox{\textSFxi{}}~
    = help: did you mean to enclose it within a copy?
    ~\lstbox{\textSFxi{}}~
 2  -     let dmut _ = [1 => 2];
 2  +     let dmut _ = copy [1 => 2];
    ~\lstbox{\textSFii{}}~~\lstbox{\textSFx{}}~~\lstbox{\textSFx{}}~~\lstbox{\textSFx{}}~~\lstbox{\textSFx{}}~~\lstbox{\textSFx{}}~~\lstbox{\textSFx{}}~~\lstbox{\textSFx{}}~
\end{lstlisting}

\subsubsection*{How to fix it}

Construct it with \tokennolst{copy}.

\errorcode{E4121}{the type \errhole{} cannot be used as a key, it is not comparable}
\label{sec:codes:e4121}

Raised during \textbf{validation}, from \texttt{ValidateErrorMessage::\allowbreak{}MAP\_\allowbreak{}KEY\_\allowbreak{}NOT\_\allowbreak{}COMPARABLE} (\textit{src/\allowbreak{}ymirc/\allowbreak{}semantic/\allowbreak{}validator/\allowbreak{}errors.yr}).

\subsubsection*{What it means}

A type is used as a map key and defines no comparison, so the map could not tell two keys apart.

\subsubsection*{Example}

\textit{test\_\allowbreak{}resources/\allowbreak{}diagnostics/\allowbreak{}test111.yr}:

%% check: skip
\begin{lstlisting}[style=coloredverbatimError]
// a class used as a map key, hashable but not comparable
class Key {
    pub self () {}
    impl Hashable;
}

fn main () {
    let _m_ : [&Key => i32] = copy [copy Key () => 1];
}
\end{lstlisting}

\begin{lstlisting}[style=bashVerb, breaklines=true, escapechar=~]
Error[E4121] : the type &(test111::Key) cannot be used as a key, it is not comparable
 --> test_resources/diagnostics/test111.yr:(2,7)
 2  ~\lstbox{\textSFxi{}}~ class Key {
    ~\lstbox{\textSFv{}}~       ^^^
 --> test_resources/diagnostics/test111.yr:(8,36)
 8  ~\lstbox{\textSFxi{}}~     let _m_ : [&Key => i32] = copy [copy Key () => 1];
    ~\lstbox{\textSFv{}}~                                    -  error: undefined operator == for types &(test111::Key) and &(test111::Key)
    ~\lstbox{\textSFxi{}}~                                    ~\lstbox{\textSFviii{}}~~\lstbox{\textSFx{}}~ error: class &(test111::Key) has no method named opEquals
    ~\lstbox{\textSFxi{}}~                                    ~\lstbox{\textSFii{}}~~\lstbox{\textSFx{}}~ error: class &(test111::Key) has no method named opCmp
    ~\lstbox{\textSFxi{}}~
    = when validating -- (test_resources/diagnostics/test111.yr:(8,36))
    = when validating test111::main -- (test_resources/diagnostics/test111.yr:(7,4))
    ~\lstbox{\textSFii{}}~~\lstbox{\textSFx{}}~~\lstbox{\textSFx{}}~~\lstbox{\textSFx{}}~~\lstbox{\textSFx{}}~~\lstbox{\textSFx{}}~~\lstbox{\textSFx{}}~~\lstbox{\textSFx{}}~
\end{lstlisting}

\subsubsection*{How to fix it}

Define the comparison operator on the type, or key the map on something that already has one.

\errorcode{E4122}{the type \errhole{} cannot be used as a key, it is not hashable}
\label{sec:codes:e4122}

Raised during \textbf{validation}, from \texttt{ValidateErrorMessage::\allowbreak{}MAP\_\allowbreak{}KEY\_\allowbreak{}NOT\_\allowbreak{}HASHABLE} (\textit{src/\allowbreak{}ymirc/\allowbreak{}semantic/\allowbreak{}validator/\allowbreak{}errors.yr}).

\subsubsection*{What it means}

A type is used as a map key and cannot be hashed.

\subsubsection*{Example}

\textit{test\_\allowbreak{}resources/\allowbreak{}diagnostics/\allowbreak{}test112.yr}:

%% check: skip
\begin{lstlisting}[style=coloredverbatimError]
// a class used as a map key without being hashable
class Key {
    pub self () {}
}

fn main () {
    let _m_ : [&Key => i32] = copy [copy Key () => 1];
}
\end{lstlisting}

\begin{lstlisting}[style=bashVerb, breaklines=true, escapechar=~]
Error[E4122] : the type &(test112::Key) cannot be used as a key, it is not hashable
 --> test_resources/diagnostics/test112.yr:(2,7)
 2  ~\lstbox{\textSFxi{}}~ class Key {
    ~\lstbox{\textSFv{}}~       ^^^  error: class type test112::Key does not implement trait core::types::Hashable
 --> test_resources/diagnostics/test112.yr:(7,36)
 7  ~\lstbox{\textSFxi{}}~     let _m_ : [&Key => i32] = copy [copy Key () => 1];
    ~\lstbox{\textSFv{}}~                                    -
    ~\lstbox{\textSFxi{}}~
    = when validating -- (test_resources/diagnostics/test112.yr:(7,36))
    = when validating test112::main -- (test_resources/diagnostics/test112.yr:(6,4))
    ~\lstbox{\textSFii{}}~~\lstbox{\textSFx{}}~~\lstbox{\textSFx{}}~~\lstbox{\textSFx{}}~~\lstbox{\textSFx{}}~~\lstbox{\textSFx{}}~~\lstbox{\textSFx{}}~~\lstbox{\textSFx{}}~
\end{lstlisting}

\subsubsection*{How to fix it}

Make the type hashable, or key the map on a value derived from it that already is.

\errorcode{E4131}{mismatch tuple arity \errhole{} and \errhole{}}
\label{sec:codes:e4131}

Raised during \textbf{validation}, from \texttt{ValidateErrorMessage::\allowbreak{}MISMATCH\_\allowbreak{}TUPLE\_\allowbreak{}ARITY} (\textit{src/\allowbreak{}ymirc/\allowbreak{}semantic/\allowbreak{}validator/\allowbreak{}errors.yr}).

\subsubsection*{What it means}

Two tuples meet where they have to agree and do not hold the same number of elements. Both arities are in the message.

\subsubsection*{Example}

\textit{test\_\allowbreak{}resources/\allowbreak{}lit\_\allowbreak{}tuples/\allowbreak{}test21.yr}:

%% check: skip
\begin{lstlisting}[style=coloredverbatimError]
fn foo ()-> (i32, i32);

fn main () {
    let mut a = 0, mut b = 0, mut c = 0;

    (a, b, c) = foo ();
}
\end{lstlisting}

\begin{lstlisting}[style=bashVerb, breaklines=true, escapechar=~]
Error[E4131] : mismatch tuple arity 3 and 2
 --> test_resources/lit_tuples/test21.yr:(6,5)
 6  ~\lstbox{\textSFxi{}}~     (a, b, c) = foo ();
    ~\lstbox{\textSFv{}}~     ^               ^
    ~\lstbox{\textSFii{}}~~\lstbox{\textSFx{}}~~\lstbox{\textSFx{}}~~\lstbox{\textSFx{}}~~\lstbox{\textSFx{}}~~\lstbox{\textSFx{}}~~\lstbox{\textSFx{}}~~\lstbox{\textSFx{}}~
\end{lstlisting}

\subsubsection*{How to fix it}

Make the two tuples the same length. A destructuring that binds fewer names than the tuple has also produces this.

\errorcode{E4163}{\errhole{} is not a tuple type}
\label{sec:codes:e4163}

Raised during \textbf{validation}, from \texttt{ValidateErrorMessage::\allowbreak{}NOT\_\allowbreak{}A\_\allowbreak{}TUPLE} (\textit{src/\allowbreak{}ymirc/\allowbreak{}semantic/\allowbreak{}validator/\allowbreak{}errors.yr}).

\subsubsection*{What it means}

An operation defined only on tuples was applied to something else.

\subsubsection*{Example}

\textit{test\_\allowbreak{}resources/\allowbreak{}lit\_\allowbreak{}tuples/\allowbreak{}test22.yr}:

%% check: skip
\begin{lstlisting}[style=coloredverbatimError]
fn main () {
    let mut a = 0, mut b = 0;

    (a, b) = 12;
}
\end{lstlisting}

\begin{lstlisting}[style=bashVerb, breaklines=true, escapechar=~]
Error[E4163] : i32 is not a tuple type
 --> test_resources/lit_tuples/test22.yr:(4,12)
 4  ~\lstbox{\textSFxi{}}~     (a, b) = 12;
    ~\lstbox{\textSFv{}}~            ^ ^^
    ~\lstbox{\textSFii{}}~~\lstbox{\textSFx{}}~~\lstbox{\textSFx{}}~~\lstbox{\textSFx{}}~~\lstbox{\textSFx{}}~~\lstbox{\textSFx{}}~~\lstbox{\textSFx{}}~~\lstbox{\textSFx{}}~
\end{lstlisting}

\subsubsection*{How to fix it}

Use a tuple value.

\errorcode{E4175}{option of type \errhole{} has no error}
\label{sec:codes:e4175}

Raised during \textbf{validation}, from \texttt{ValidateErrorMessage::\allowbreak{}OPTION\_\allowbreak{}HAS\_\allowbreak{}NO\_\allowbreak{}ERROR} (\textit{src/\allowbreak{}ymirc/\allowbreak{}semantic/\allowbreak{}validator/\allowbreak{}errors.yr}).

\subsubsection*{What it means}

The error side of an option was read on an option that never holds one.

\subsubsection*{Example}

\textit{test\_\allowbreak{}resources/\allowbreak{}diagnostics/\allowbreak{}test116.yr}:

%% check: skip
\begin{lstlisting}[style=coloredverbatimError]
// the error of an option known to hold a value
fn main () {
    let x : (i32)? = (1)?;
    let _ = x.error;
}
\end{lstlisting}

\begin{lstlisting}[style=bashVerb, breaklines=true, escapechar=~]
Error[E4175] : option of type (i32)? has no error
 --> test_resources/diagnostics/test116.yr:(4,14)
 4  ~\lstbox{\textSFxi{}}~     let _ = x.error;
    ~\lstbox{\textSFv{}}~              ^
    ~\lstbox{\textSFxi{}}~
    = when validating test116::main -- (test_resources/diagnostics/test116.yr:(2,4))
    ~\lstbox{\textSFii{}}~~\lstbox{\textSFx{}}~~\lstbox{\textSFx{}}~~\lstbox{\textSFx{}}~~\lstbox{\textSFx{}}~~\lstbox{\textSFx{}}~~\lstbox{\textSFx{}}~~\lstbox{\textSFx{}}~
\end{lstlisting}

\subsubsection*{How to fix it}

Read the value instead. An option built from a non-throwing computation has no error, cf. \hyperref[sec:codes:e4078]{E4078}.

\errorcode{E4176}{option of type \errhole{} has no value}
\label{sec:codes:e4176}

Raised during \textbf{validation}, from \texttt{ValidateErrorMessage::\allowbreak{}OPTION\_\allowbreak{}HAS\_\allowbreak{}NO\_\allowbreak{}VALUE} (\textit{src/\allowbreak{}ymirc/\allowbreak{}semantic/\allowbreak{}validator/\allowbreak{}errors.yr}).

\subsubsection*{What it means}

The value side of an option was read on an option that never holds one.

\subsubsection*{Example}

\textit{test\_\allowbreak{}resources/\allowbreak{}diagnostics/\allowbreak{}test117.yr}:

%% check: skip
\begin{lstlisting}[style=coloredverbatimError]
// the value of an option known to be empty
fn main () {
    let x : (i32)? = none;
    let _ = x.value;
}
\end{lstlisting}

\begin{lstlisting}[style=bashVerb, breaklines=true, escapechar=~]
Error[E4176] : option of type (i32)? has no value
 --> test_resources/diagnostics/test117.yr:(4,14)
 4  ~\lstbox{\textSFxi{}}~     let _ = x.value;
    ~\lstbox{\textSFv{}}~              ^
    ~\lstbox{\textSFxi{}}~
    = when validating test117::main -- (test_resources/diagnostics/test117.yr:(2,4))
    ~\lstbox{\textSFii{}}~~\lstbox{\textSFx{}}~~\lstbox{\textSFx{}}~~\lstbox{\textSFx{}}~~\lstbox{\textSFx{}}~~\lstbox{\textSFx{}}~~\lstbox{\textSFx{}}~~\lstbox{\textSFx{}}~
\end{lstlisting}

\subsubsection*{How to fix it}

Read the error instead, or build the option from a computation that can succeed.

\errorcode{E4179}{static array size \errhole{} exceeds limits \errhole{}}
\label{sec:codes:e4179}

Raised during \textbf{validation}, from \texttt{ValidateErrorMessage::\allowbreak{}OVERFLOW\_\allowbreak{}CAPACITY\_\allowbreak{}ARRAY} (\textit{src/\allowbreak{}ymirc/\allowbreak{}semantic/\allowbreak{}validator/\allowbreak{}errors.yr}).

\subsubsection*{What it means}

A static array is declared with a size beyond what the compiler accepts.

\subsubsection*{Example}

\textit{test\_\allowbreak{}resources/\allowbreak{}diagnostics/\allowbreak{}test118.yr}:

%% check: skip
\begin{lstlisting}[style=coloredverbatimError]
// a static array too large for the stack
fn main () {
    let _x_ : [i32 ; 100000000] = [0 ; 100000000];
}
\end{lstlisting}

\begin{lstlisting}[style=bashVerb, breaklines=true, escapechar=~]
Error[E4179] : static array size 100000000 exceeds limits 1048576
 --> test_resources/diagnostics/test118.yr:(3,40)
 3  ~\lstbox{\textSFxi{}}~     let _x_ : [i32 ; 100000000] = [0 ; 100000000];
    ~\lstbox{\textSFv{}}~                                        ^^^^^^^^^
    ~\lstbox{\textSFxi{}}~
    = when validating test118::main -- (test_resources/diagnostics/test118.yr:(2,4))
    ~\lstbox{\textSFii{}}~~\lstbox{\textSFx{}}~~\lstbox{\textSFx{}}~~\lstbox{\textSFx{}}~~\lstbox{\textSFx{}}~~\lstbox{\textSFx{}}~~\lstbox{\textSFx{}}~~\lstbox{\textSFx{}}~
\end{lstlisting}

\subsubsection*{How to fix it}

Reduce the size, or allocate a slice at run time instead.

\errorcode{E4213}{tuple left affectation is only defined for the operator \errhole{}}
\label{sec:codes:e4213}

Raised during \textbf{validation}, from \texttt{ValidateErrorMessage::\allowbreak{}TUPLE\_\allowbreak{}AFFECT\_\allowbreak{}OPERATOR} (\textit{src/\allowbreak{}ymirc/\allowbreak{}semantic/\allowbreak{}validator/\allowbreak{}errors.yr}).

\subsubsection*{What it means}

A tuple on the left of an assignment was used with an operator other than the plain one. Destructuring assignment is defined for \tokennolst{=} alone.

\subsubsection*{Example}

\textit{test\_\allowbreak{}resources/\allowbreak{}lit\_\allowbreak{}tuples/\allowbreak{}test23.yr}:

%% check: skip
\begin{lstlisting}[style=coloredverbatimError]
fn foo ()-> (i32, i32);

fn main () {
    let mut a = 0, mut b = 0;

    (a, b) += foo ();
}
\end{lstlisting}

\begin{lstlisting}[style=bashVerb, breaklines=true, escapechar=~]
Error[E4213] : tuple left affectation is only defined for the operator =
 --> test_resources/lit_tuples/test23.yr:(6,5)
 6  ~\lstbox{\textSFxi{}}~     (a, b) += foo ();
    ~\lstbox{\textSFv{}}~     ^      ^^
    ~\lstbox{\textSFii{}}~~\lstbox{\textSFx{}}~~\lstbox{\textSFx{}}~~\lstbox{\textSFx{}}~~\lstbox{\textSFx{}}~~\lstbox{\textSFx{}}~~\lstbox{\textSFx{}}~~\lstbox{\textSFx{}}~
\end{lstlisting}

\subsubsection*{How to fix it}

Split the compound assignment into a read and a write.

\errorcode{E4214}{tuple access out of bound(\errhole{}), tuple arity is \errhole{}}
\label{sec:codes:e4214}

Raised during \textbf{validation}, from \texttt{ValidateErrorMessage::\allowbreak{}TUPLE\_\allowbreak{}ARITY\_\allowbreak{}OVERFLOW} (\textit{src/\allowbreak{}ymirc/\allowbreak{}semantic/\allowbreak{}validator/\allowbreak{}errors.yr}).

\subsubsection*{What it means}

A tuple element was accessed beyond the arity of the tuple. Both the index and the arity are in the message.

\subsubsection*{Example}

\textit{test\_\allowbreak{}resources/\allowbreak{}diagnostics/\allowbreak{}test71.yr}:

%% check: skip
\begin{lstlisting}[style=coloredverbatimError]
fn main () {
    let t = (1, 2);
    let _ = t._5;
}
\end{lstlisting}

\begin{lstlisting}[style=bashVerb, breaklines=true, escapechar=~]
Error[E4214] : tuple access out of bound(5), tuple arity is 2
 --> test_resources/diagnostics/test71.yr:(3,14)
 3  ~\lstbox{\textSFxi{}}~     let _ = t._5;
    ~\lstbox{\textSFv{}}~              ^
    ~\lstbox{\textSFii{}}~~\lstbox{\textSFx{}}~~\lstbox{\textSFx{}}~~\lstbox{\textSFx{}}~~\lstbox{\textSFx{}}~~\lstbox{\textSFx{}}~~\lstbox{\textSFx{}}~~\lstbox{\textSFx{}}~
\end{lstlisting}

\subsubsection*{How to fix it}

Correct the index; tuple elements are numbered from zero.

\errorcode{E4267}{operator \errhole{} is useless, the left operand of type \errhole{} always has a value}
\label{sec:codes:e4267}

Raised during \textbf{validation}, from \texttt{ValidateErrorMessage::\allowbreak{}USELESS\_\allowbreak{}OPTION\_\allowbreak{}OR\_\allowbreak{}ELSE} (\textit{src/\allowbreak{}ymirc/\allowbreak{}semantic/\allowbreak{}validator/\allowbreak{}errors.yr}).

\subsubsection*{What it means}

An option fallback operator was applied to a value that always holds one, so the fallback can never run.

\subsubsection*{Example}

\textit{test\_\allowbreak{}resources/\allowbreak{}lit\_\allowbreak{}options/\allowbreak{}test21.yr}:

%% check: skip
\begin{lstlisting}[style=coloredverbatimError]
// the left operand always has a value, the right one would never be evaluated
fn main () {
    let _a_ : i32 = 12 ?? 42;
}
\end{lstlisting}

\begin{lstlisting}[style=bashVerb, breaklines=true, escapechar=~]
Error[E4267] : operator ?? is useless, the left operand of type i32 always has a value
 --> test_resources/lit_options/test21.yr:(3,24)
 3  ~\lstbox{\textSFxi{}}~     let _a_ : i32 = 12 ?? 42;
    ~\lstbox{\textSFv{}}~                        ^^
    ~\lstbox{\textSFii{}}~~\lstbox{\textSFx{}}~~\lstbox{\textSFx{}}~~\lstbox{\textSFx{}}~~\lstbox{\textSFx{}}~~\lstbox{\textSFx{}}~~\lstbox{\textSFx{}}~~\lstbox{\textSFx{}}~
\end{lstlisting}

\subsubsection*{How to fix it}

Remove the operator and use the value directly.

%% Generated by tools/gen_error_codes.py from docs/errors/ of the compiler;
%% edit the compiler's pages or tools/error_themes.txt, then rerun it.

\clearpage
\section{Mutability, aliasing and moves}
\label{sec:codes:memory}

The book's chapter \textit{Variables and memory management} specifies the rules behind these codes.

\begin{longtable}{@{}l p{0.78\linewidth} r@{}}
\toprule Code & Message & Page\\ \midrule \endhead
\hyperref[sec:codes:e4002]{E4002} & class alias operator :[ is prohibited on the left operand of an affectation & \pageref{sec:codes:e4002}\\
\hyperref[sec:codes:e4014]{E4014} & class alias operator :. is not usable to access fields & \pageref{sec:codes:e4014}\\
\hyperref[sec:codes:e4015]{E4015} & alias operator :. is only usable on class, record, entity, map or generator types, not \errhole{} & \pageref{sec:codes:e4015}\\
\hyperref[sec:codes:e4016]{E4016} & class alias operator :[ is only usable on class, record or entity types, not \errhole{} & \pageref{sec:codes:e4016}\\
\hyperref[sec:codes:e4029]{E4029} & movable type \errhole{} cannot be embedded within another type & \pageref{sec:codes:e4029}\\
\hyperref[sec:codes:e4045]{E4045} & discarding the constant property is prohibited & \pageref{sec:codes:e4045}\\
\hyperref[sec:codes:e4082]{E4082} & left operand of type \errhole{} is immutable & \pageref{sec:codes:e4082}\\
\hyperref[sec:codes:e4083]{E4083} & type must be immutable & \pageref{sec:codes:e4083}\\
\hyperref[sec:codes:e4088]{E4088} & implicit move of type \errhole{} is prohibited & \pageref{sec:codes:e4088}\\
\hyperref[sec:codes:e4089]{E4089} & ref of type \errhole{} must be mutable, but is const & \pageref{sec:codes:e4089}\\
\hyperref[sec:codes:e4092]{E4092} & cannot pass value of type \errhole{} as ref & \pageref{sec:codes:e4092}\\
\hyperref[sec:codes:e4099]{E4099} & index assignment of a value of type \errhole{} requires a mutable value & \pageref{sec:codes:e4099}\\
\hyperref[sec:codes:e4125]{E4125} & cannot alias an expand value, maybe alias and expand keywords are inverted? & \pageref{sec:codes:e4125}\\
\hyperref[sec:codes:e4126]{E4126} & cannot alias a lazy value, maybe alias and lazy keywords are inverted? & \pageref{sec:codes:e4126}\\
\hyperref[sec:codes:e4127]{E4127} & cannot copy an expand value, maybe copy and expand keywords are inverted? & \pageref{sec:codes:e4127}\\
\hyperref[sec:codes:e4128]{E4128} & cannot copy a lazy value, maybe copy and lazy keywords are inverted? & \pageref{sec:codes:e4128}\\
\hyperref[sec:codes:e4129]{E4129} & cannot copy an expand value, maybe dcopy and expand keywords are inverted? & \pageref{sec:codes:e4129}\\
\hyperref[sec:codes:e4130]{E4130} & cannot copy a lazy value, maybe dcopy and lazy keywords are inverted? & \pageref{sec:codes:e4130}\\
\hyperref[sec:codes:e4137]{E4137} & an iterator cannot be mutable, if it is not a reference or does not borrow mutable data & \pageref{sec:codes:e4137}\\
\hyperref[sec:codes:e4138]{E4138} & a parameter cannot be mutable, if it is not a reference or does not borrow mutable data & \pageref{sec:codes:e4138}\\
\hyperref[sec:codes:e4140]{E4140} & a lazy variable cannot be mutable, if it does not borrow mutable data & \pageref{sec:codes:e4140}\\
\hyperref[sec:codes:e4146]{E4146} & \errhole{} is not an aliasable type & \pageref{sec:codes:e4146}\\
\hyperref[sec:codes:e4151]{E4151} & not a lvalue & \pageref{sec:codes:e4151}\\
\hyperref[sec:codes:e4153]{E4153} & \errhole{} is a value iterator, and cannot be used as a lvalue & \pageref{sec:codes:e4153}\\
\hyperref[sec:codes:e4154]{E4154} & \errhole{} is a lazy variable, and cannot be used as a lvalue & \pageref{sec:codes:e4154}\\
\hyperref[sec:codes:e4155]{E4155} & \errhole{} is a map access, and cannot be used as a lvalue & \pageref{sec:codes:e4155}\\
\hyperref[sec:codes:e4156]{E4156} & \errhole{} is a value parameter, and cannot be used as a lvalue & \pageref{sec:codes:e4156}\\
\hyperref[sec:codes:e4157]{E4157} & value of type \errhole{} is not a lvalue & \pageref{sec:codes:e4157}\\
\hyperref[sec:codes:e4165]{E4165} & type \errhole{} is not a movable type & \pageref{sec:codes:e4165}\\
\hyperref[sec:codes:e4167]{E4167} & no copy exists for type \errhole{} & \pageref{sec:codes:e4167}\\
\hyperref[sec:codes:e4247]{E4247} & aliasing the value of type \errhole{} to create a constant borrowing is prohibited & \pageref{sec:codes:e4247}\\
\hyperref[sec:codes:e4248]{E4248} & aliasing the value of type \errhole{} to call begin and end iterator constant methods is useless & \pageref{sec:codes:e4248}\\
\hyperref[sec:codes:e4251]{E4251} & the move has no effect & \pageref{sec:codes:e4251}\\
\hyperref[sec:codes:e4252]{E4252} & referencing the value has no effect & \pageref{sec:codes:e4252}\\
\bottomrule
\end{longtable}

\errorcode{E4002}{class alias operator :[ is prohibited on the left operand of an affectation}
\label{sec:codes:e4002}

Raised during \textbf{validation}, from \texttt{ValidateErrorMessage::\allowbreak{}ALIAS\_\allowbreak{}INDEX\_\allowbreak{}AFFECT} (\textit{src/\allowbreak{}ymirc/\allowbreak{}semantic/\allowbreak{}validator/\allowbreak{}errors.yr}).

\subsubsection*{What it means}

The class alias index operator \tokennolst{:[} was used on the left of an assignment. That operator produces a borrowed view of the indexed value; the left of an assignment is storage, not a value, so there is nothing for the alias to borrow.

\subsubsection*{Example}

\textit{test\_\allowbreak{}resources/\allowbreak{}lit\_\allowbreak{}class/\allowbreak{}operators/\allowbreak{}test16.yr}:

%% check: skip
\begin{lstlisting}[style=coloredverbatimError]
class A {

    pub self () {}

    pub fn opIndexAssign (self, _ : i32, _ : i32) {}

}

fn main () {
    let dmut a = copy A ();
    a:[0] = 0;    
    a:.opIndexAssign (0, 0);
}
\end{lstlisting}

\begin{lstlisting}[style=bashVerb, breaklines=true, escapechar=~]
Error[E4002] : class alias operator :[ is prohibited on the left operand of an affectation
 --> test_resources/lit_class/operators/test16.yr:(11,6)
11  ~\lstbox{\textSFxi{}}~     a:[0] = 0;
    ~\lstbox{\textSFv{}}~      ^^
    ~\lstbox{\textSFxi{}}~
    = help: the index assignment already aliases the value implicitly, use [ instead
    ~\lstbox{\textSFxi{}}~
11  -     a:[0] = 0;
11  +     a[0] = 0;
    ~\lstbox{\textSFii{}}~~\lstbox{\textSFx{}}~~\lstbox{\textSFx{}}~~\lstbox{\textSFx{}}~~\lstbox{\textSFx{}}~~\lstbox{\textSFx{}}~~\lstbox{\textSFx{}}~~\lstbox{\textSFx{}}~
\end{lstlisting}

\subsubsection*{How to fix it}

Use the plain \tokennolst{[} index operator on the left of the assignment. The assignment already aliases the value it writes, which is the edit the diagnostic proposes.

\errorcode{E4014}{class alias operator :. is not usable to access fields}
\label{sec:codes:e4014}

Raised during \textbf{validation}, from \texttt{ValidateErrorMessage::\allowbreak{}CLASS\_\allowbreak{}ALIAS\_\allowbreak{}FIELD} (\textit{src/\allowbreak{}ymirc/\allowbreak{}semantic/\allowbreak{}validator/\allowbreak{}errors.yr}).

\subsubsection*{What it means}

The class alias operator \tokennolst{:.} was used to access a field. That operator produces a borrowed view for calling a method; a field access takes its mutability from the value itself, so the operator has nothing to do.

\subsubsection*{Example}

\textit{test\_\allowbreak{}resources/\allowbreak{}diagnostics/\allowbreak{}test98.yr}:

%% check: skip
\begin{lstlisting}[style=coloredverbatimError]
// the alias operator used to read a field of a map
fn main () {
    let mut m = copy [1 => 2];
    let _ = m:.len;
}
\end{lstlisting}

\begin{lstlisting}[style=bashVerb, breaklines=true, escapechar=~]
Error[E4014] : class alias operator :. is not usable to access fields
 --> test_resources/diagnostics/test98.yr:(4,14)
 4  ~\lstbox{\textSFxi{}}~     let _ = m:.len;
    ~\lstbox{\textSFv{}}~              ^^---
    ~\lstbox{\textSFxi{}}~
    = when validating test98::main -- (test_resources/diagnostics/test98.yr:(2,4))
    ~\lstbox{\textSFii{}}~~\lstbox{\textSFx{}}~~\lstbox{\textSFx{}}~~\lstbox{\textSFx{}}~~\lstbox{\textSFx{}}~~\lstbox{\textSFx{}}~~\lstbox{\textSFx{}}~~\lstbox{\textSFx{}}~
\end{lstlisting}

\subsubsection*{How to fix it}

Use a plain \tokennolst{.} to access the field.

\errorcode{E4015}{alias operator :. is only usable on class, record, entity, map or generator types, not \errhole{}}
\label{sec:codes:e4015}

Raised during \textbf{validation}, from \texttt{ValidateErrorMessage::\allowbreak{}CLASS\_\allowbreak{}ALIAS\_\allowbreak{}FIELD\_\allowbreak{}NO\_\allowbreak{}CLASS} (\textit{src/\allowbreak{}ymirc/\allowbreak{}semantic/\allowbreak{}validator/\allowbreak{}errors.yr}).

\subsubsection*{What it means}

The alias operator \tokennolst{:.} was applied to a type that has no method to alias -- it is only defined for classes, records, entities, maps and generators.

\subsubsection*{Example}

\textit{test\_\allowbreak{}resources/\allowbreak{}diagnostics/\allowbreak{}test23.yr}:

%% check: skip
\begin{lstlisting}[style=coloredverbatimError]
fn main () {
    let x = 1;
    x:.foo ();
}
\end{lstlisting}

\begin{lstlisting}[style=bashVerb, breaklines=true, escapechar=~]
Error[E4015] : alias operator :. is only usable on class, record, entity, map or generator types, not x
 --> test_resources/diagnostics/test23.yr:(3,5)
 3  ~\lstbox{\textSFxi{}}~     x:.foo ();
    ~\lstbox{\textSFv{}}~     ^
    ~\lstbox{\textSFii{}}~~\lstbox{\textSFx{}}~~\lstbox{\textSFx{}}~~\lstbox{\textSFx{}}~~\lstbox{\textSFx{}}~~\lstbox{\textSFx{}}~~\lstbox{\textSFx{}}~~\lstbox{\textSFx{}}~
\end{lstlisting}

\subsubsection*{How to fix it}

Drop the operator, or apply it to a value of one of those types.

\errorcode{E4016}{class alias operator :[ is only usable on class, record or entity types, not \errhole{}}
\label{sec:codes:e4016}

Raised during \textbf{validation}, from \texttt{ValidateErrorMessage::\allowbreak{}CLASS\_\allowbreak{}ALIAS\_\allowbreak{}INDEX\_\allowbreak{}NO\_\allowbreak{}CLASS} (\textit{src/\allowbreak{}ymirc/\allowbreak{}semantic/\allowbreak{}validator/\allowbreak{}errors.yr}).

\subsubsection*{What it means}

The class alias index operator \tokennolst{:[} was applied to a type it is not defined for; it only applies to classes, records and entities.

\subsubsection*{Example}

\textit{test\_\allowbreak{}resources/\allowbreak{}diagnostics/\allowbreak{}test24.yr}:

%% check: skip
\begin{lstlisting}[style=coloredverbatimError]
fn main () {
    let x = 1;
    x:[0];
}
\end{lstlisting}

\begin{lstlisting}[style=bashVerb, breaklines=true, escapechar=~]
Error[E4016] : class alias operator :[ is only usable on class, record or entity types, not x
 --> test_resources/diagnostics/test24.yr:(3,5)
 3  ~\lstbox{\textSFxi{}}~     x:[0];
    ~\lstbox{\textSFv{}}~     ^
    ~\lstbox{\textSFii{}}~~\lstbox{\textSFx{}}~~\lstbox{\textSFx{}}~~\lstbox{\textSFx{}}~~\lstbox{\textSFx{}}~~\lstbox{\textSFx{}}~~\lstbox{\textSFx{}}~~\lstbox{\textSFx{}}~
\end{lstlisting}

\subsubsection*{How to fix it}

Use the plain \tokennolst{[} index operator.

\errorcode{E4029}{movable type \errhole{} cannot be embedded within another type}
\label{sec:codes:e4029}

Raised during \textbf{validation}, from \texttt{ValidateErrorMessage::\allowbreak{}CONTAIN\_\allowbreak{}MOVABLE\_\allowbreak{}TYPE} (\textit{src/\allowbreak{}ymirc/\allowbreak{}semantic/\allowbreak{}validator/\allowbreak{}errors.yr}).

\subsubsection*{What it means}

A movable type was embedded in another type. A movable type has an identity that survives being moved, which embedding it by value would break.

\subsubsection*{Example}

\textit{test\_\allowbreak{}resources/\allowbreak{}union/\allowbreak{}test4.yr}:

%% check: skip
\begin{lstlisting}[style=coloredverbatimError]
entity Z {
    pub self () {}
}

@overlaid
record A {
    let z : Z;
}
\end{lstlisting}

\begin{lstlisting}[style=bashVerb, breaklines=true, escapechar=~]
Error[E4029] : movable type test4::Z cannot be embedded within another type
 --> test_resources/union/test4.yr:(7,5)
 7  ~\lstbox{\textSFxi{}}~     let z : Z;
    ~\lstbox{\textSFv{}}~     ^^^     ^
    ~\lstbox{\textSFxi{}}~ Note : the structure is defined as a record but must be an entity to store movable values
    ~\lstbox{\textSFxi{}}~  --> test_resources/union/test4.yr:(6,8)
    ~\lstbox{\textSFxi{}}~  6  ~\lstbox{\textSFxi{}}~ record A {
    ~\lstbox{\textSFxi{}}~     ~\lstbox{\textSFv{}}~        ^
    ~\lstbox{\textSFxi{}}~     ~\lstbox{\textSFii{}}~~\lstbox{\textSFx{}}~~\lstbox{\textSFx{}}~~\lstbox{\textSFx{}}~~\lstbox{\textSFx{}}~~\lstbox{\textSFx{}}~~\lstbox{\textSFx{}}~~\lstbox{\textSFx{}}~
    ~\lstbox{\textSFii{}}~~\lstbox{\textSFx{}}~~\lstbox{\textSFx{}}~~\lstbox{\textSFx{}}~~\lstbox{\textSFx{}}~~\lstbox{\textSFx{}}~~\lstbox{\textSFvii{}}~~\lstbox{\textSFx{}}~
\end{lstlisting}

\subsubsection*{How to fix it}

Hold it behind a class or a pointer, rather than by value.

\errorcode{E4045}{discarding the constant property is prohibited}
\label{sec:codes:e4045}

Raised during \textbf{validation}, from \texttt{ValidateErrorMessage::\allowbreak{}DISCARD\_\allowbreak{}CONST} (\textit{src/\allowbreak{}ymirc/\allowbreak{}semantic/\allowbreak{}validator/\allowbreak{}errors.yr}).

\subsubsection*{What it means}

An operation would drop the constant property of a value -- turning something borrowed as constant into something writable.

\subsubsection*{Example}

\textit{test\_\allowbreak{}resources/\allowbreak{}lit\_\allowbreak{}tuples/\allowbreak{}test13.yr}:

%% check: skip
\begin{lstlisting}[style=coloredverbatimError]
fn main () {
    let dmut a = (copy [1, 2],);
    let dmut _ = a.0;
}
\end{lstlisting}

\begin{lstlisting}[style=bashVerb, breaklines=true, escapechar=~]
Error[E4045] : discarding the constant property is prohibited
 --> test_resources/lit_tuples/test13.yr:(3,14)
 3  ~\lstbox{\textSFxi{}}~     let dmut _ = a.0;
    ~\lstbox{\textSFv{}}~              ^
    ~\lstbox{\textSFxi{}}~ Note : implicit alias of type mut [mut i32] is prohibited, it implicitly gives up on borrowings of mutable values
    ~\lstbox{\textSFxi{}}~  --> test_resources/lit_tuples/test13.yr:(3,19)
    ~\lstbox{\textSFxi{}}~  3  ~\lstbox{\textSFxi{}}~     let dmut _ = a.0;
    ~\lstbox{\textSFxi{}}~     ~\lstbox{\textSFv{}}~                   ^
    ~\lstbox{\textSFxi{}}~     ~\lstbox{\textSFii{}}~~\lstbox{\textSFx{}}~~\lstbox{\textSFx{}}~~\lstbox{\textSFx{}}~~\lstbox{\textSFx{}}~~\lstbox{\textSFx{}}~~\lstbox{\textSFx{}}~~\lstbox{\textSFx{}}~
    ~\lstbox{\textSFii{}}~~\lstbox{\textSFx{}}~~\lstbox{\textSFx{}}~~\lstbox{\textSFx{}}~~\lstbox{\textSFx{}}~~\lstbox{\textSFx{}}~~\lstbox{\textSFvii{}}~~\lstbox{\textSFx{}}~
\end{lstlisting}

\subsubsection*{How to fix it}

Take a copy if the value has to be modified, or keep the operation constant. The mutability of a borrowed value is part of its type and cannot be cast away.

\errorcode{E4082}{left operand of type \errhole{} is immutable}
\label{sec:codes:e4082}

Raised during \textbf{validation}, from \texttt{ValidateErrorMessage::\allowbreak{}IMMUTABLE\_\allowbreak{}LVALUE} (\textit{src/\allowbreak{}ymirc/\allowbreak{}semantic/\allowbreak{}validator/\allowbreak{}errors.yr}).

\subsubsection*{What it means}

An assignment writes through a value whose type is immutable.

\subsubsection*{Example}

\textit{test\_\allowbreak{}resources/\allowbreak{}for\_\allowbreak{}loops/\allowbreak{}tuples/\allowbreak{}test2.yr}:

%% check: skip
\begin{lstlisting}[style=coloredverbatimError]
fn main () {
    let a = (1, 2, 3);
    for x in ref a {
        x;
    }
}
\end{lstlisting}

\begin{lstlisting}[style=bashVerb, breaklines=true, escapechar=~]
Error[E4082] : left operand of type (i32, i32, i32) is immutable
 --> test_resources/for_loops/tuples/test2.yr:(3,14)
 3  ~\lstbox{\textSFxi{}}~     for x in ref a {
    ~\lstbox{\textSFv{}}~              ^^^ ^
    ~\lstbox{\textSFii{}}~~\lstbox{\textSFx{}}~~\lstbox{\textSFx{}}~~\lstbox{\textSFx{}}~~\lstbox{\textSFx{}}~~\lstbox{\textSFx{}}~~\lstbox{\textSFx{}}~~\lstbox{\textSFx{}}~
\end{lstlisting}

\subsubsection*{How to fix it}

Declare the variable, or the borrow it came from, as mutable. Mutability travels with the type, so a value reached through a constant borrow stays constant.

\errorcode{E4083}{type must be immutable}
\label{sec:codes:e4083}

Raised during \textbf{validation}, from \texttt{ValidateErrorMessage::\allowbreak{}IMMUTABLE\_\allowbreak{}PARENT} (\textit{src/\allowbreak{}ymirc/\allowbreak{}semantic/\allowbreak{}validator/\allowbreak{}errors.yr}).

\subsubsection*{What it means}

A type used in this position has to be immutable and is not -- typically the aliased type of a \tokennolst{def}, or a type used where sharing requires immutability.

\subsubsection*{Example}

\textit{test\_\allowbreak{}resources/\allowbreak{}aka/\allowbreak{}test2.yr}:

%% check: skip
\begin{lstlisting}[style=coloredverbatimError]
def A : dmut [i32];

fn main ()
{
    let dmut x : A = copy [1, 2, 3];
    x [0] = 1;
}
\end{lstlisting}

\begin{lstlisting}[style=bashVerb, breaklines=true, escapechar=~]
Error[E4083] : type must be immutable
 --> test_resources/aka/test2.yr:(1,9)
 1  ~\lstbox{\textSFxi{}}~ def A : dmut [i32];
    ~\lstbox{\textSFv{}}~         ^^^^
    ~\lstbox{\textSFii{}}~~\lstbox{\textSFx{}}~~\lstbox{\textSFx{}}~~\lstbox{\textSFx{}}~~\lstbox{\textSFx{}}~~\lstbox{\textSFx{}}~~\lstbox{\textSFx{}}~~\lstbox{\textSFx{}}~
\end{lstlisting}

\subsubsection*{How to fix it}

Remove the mutability decorator from the type.

\errorcode{E4088}{implicit move of type \errhole{} is prohibited}
\label{sec:codes:e4088}

Raised during \textbf{validation}, from \texttt{ValidateErrorMessage::\allowbreak{}IMPLICIT\_\allowbreak{}MOVE} (\textit{src/\allowbreak{}ymirc/\allowbreak{}semantic/\allowbreak{}validator/\allowbreak{}errors.yr}).

\subsubsection*{What it means}

A value would be moved implicitly. A move leaves its source unusable, so it is always written out.

\subsubsection*{Example}

\textit{test\_\allowbreak{}resources/\allowbreak{}diagnostics/\allowbreak{}test106.yr}:

%% check: skip
\begin{lstlisting}[style=coloredverbatimError]
// an entity bound to a second variable without a move
entity Counter {
    pub self () {}
}

fn main () {
    let c = Counter ();
    let _d_ = c;
}
\end{lstlisting}

\begin{lstlisting}[style=bashVerb, breaklines=true, escapechar=~]
Error[E4088] : implicit move of type test106::Counter is prohibited
 --> test_resources/diagnostics/test106.yr:(8,9)
 8  ~\lstbox{\textSFxi{}}~     let _d_ = c;
    ~\lstbox{\textSFv{}}~         ^^^   -
    ~\lstbox{\textSFxi{}}~
    = when validating test106::main -- (test_resources/diagnostics/test106.yr:(6,4))
    ~\lstbox{\textSFii{}}~~\lstbox{\textSFx{}}~~\lstbox{\textSFx{}}~~\lstbox{\textSFx{}}~~\lstbox{\textSFx{}}~~\lstbox{\textSFx{}}~~\lstbox{\textSFx{}}~~\lstbox{\textSFx{}}~
\end{lstlisting}

\subsubsection*{How to fix it}

Write \tokennolst{move} before the value, or pass a copy instead.

\errorcode{E4089}{ref of type \errhole{} must be mutable, but is const}
\label{sec:codes:e4089}

Raised during \textbf{validation}, from \texttt{ValidateErrorMessage::\allowbreak{}IMPLICIT\_\allowbreak{}MUT\_\allowbreak{}REFERENCE} (\textit{src/\allowbreak{}ymirc/\allowbreak{}semantic/\allowbreak{}validator/\allowbreak{}errors.yr}).

\subsubsection*{What it means}

A \tokennolst{ref} parameter is declared mutable and a constant value was passed.

\subsubsection*{Example}

\textit{test\_\allowbreak{}resources/\allowbreak{}for\_\allowbreak{}loops/\allowbreak{}tuples/\allowbreak{}test5.yr}:

%% check: skip
\begin{lstlisting}[style=coloredverbatimError]
fn main () {
    let dmut a = (copy [1, 2],);
    for ref mut x in a {
        x;
    }
}
\end{lstlisting}

\begin{lstlisting}[style=bashVerb, breaklines=true, escapechar=~]
Error[E4089] : ref of type mut [mut i32] must be mutable, but is const
 --> test_resources/for_loops/tuples/test5.yr:(3,17)
 3  ~\lstbox{\textSFxi{}}~     for ref mut x in a {
    ~\lstbox{\textSFv{}}~                 ^
    ~\lstbox{\textSFii{}}~~\lstbox{\textSFx{}}~~\lstbox{\textSFx{}}~~\lstbox{\textSFx{}}~~\lstbox{\textSFx{}}~~\lstbox{\textSFx{}}~~\lstbox{\textSFx{}}~~\lstbox{\textSFx{}}~
\end{lstlisting}

\subsubsection*{How to fix it}

Pass a mutable value, or declare the parameter constant.

\errorcode{E4092}{cannot pass value of type \errhole{} as ref}
\label{sec:codes:e4092}

Raised during \textbf{validation}, from \texttt{ValidateErrorMessage::\allowbreak{}IMPLICIT\_\allowbreak{}REFERENCE} (\textit{src/\allowbreak{}ymirc/\allowbreak{}semantic/\allowbreak{}validator/\allowbreak{}errors.yr}).

\subsubsection*{What it means}

A value was passed as \tokennolst{ref} that has no storage to borrow -- a temporary, or the result of an expression.

\subsubsection*{Example}

\textit{test\_\allowbreak{}resources/\allowbreak{}control\_\allowbreak{}flow/\allowbreak{}test21.yr}:

%% check: skip
\begin{lstlisting}[style=coloredverbatimError]
fn main () {
    let (ref x,) = (1,);
}
\end{lstlisting}

\begin{lstlisting}[style=bashVerb, breaklines=true, escapechar=~]
Error[E4092] : cannot pass value of type i32 as ref
 --> test_resources/control_flow/test21.yr:(2,14)
 2  ~\lstbox{\textSFxi{}}~     let (ref x,) = (1,);
    ~\lstbox{\textSFv{}}~              ^      ^
    ~\lstbox{\textSFxi{}}~ Note : for pattern field access 0 on type (mut i32,)
    ~\lstbox{\textSFxi{}}~  --> test_resources/control_flow/test21.yr:(2,9)
    ~\lstbox{\textSFxi{}}~  2  ~\lstbox{\textSFxi{}}~     let (ref x,) = (1,);
    ~\lstbox{\textSFxi{}}~     ~\lstbox{\textSFv{}}~         ^
    ~\lstbox{\textSFxi{}}~     ~\lstbox{\textSFii{}}~~\lstbox{\textSFx{}}~~\lstbox{\textSFx{}}~~\lstbox{\textSFx{}}~~\lstbox{\textSFx{}}~~\lstbox{\textSFx{}}~~\lstbox{\textSFx{}}~~\lstbox{\textSFx{}}~
    ~\lstbox{\textSFii{}}~~\lstbox{\textSFx{}}~~\lstbox{\textSFx{}}~~\lstbox{\textSFx{}}~~\lstbox{\textSFx{}}~~\lstbox{\textSFx{}}~~\lstbox{\textSFvii{}}~~\lstbox{\textSFx{}}~
\end{lstlisting}

\subsubsection*{How to fix it}

Bind the value to a variable first and pass that.

\errorcode{E4099}{index assignment of a value of type \errhole{} requires a mutable value}
\label{sec:codes:e4099}

Raised during \textbf{validation}, from \texttt{ValidateErrorMessage::\allowbreak{}INDEX\_\allowbreak{}ASSIGN\_\allowbreak{}NOT\_\allowbreak{}MUTABLE} (\textit{src/\allowbreak{}ymirc/\allowbreak{}semantic/\allowbreak{}validator/\allowbreak{}errors.yr}).

\subsubsection*{What it means}

An index assignment writes into a value that is not mutable.

\subsubsection*{Example}

\textit{test\_\allowbreak{}resources/\allowbreak{}lit\_\allowbreak{}class/\allowbreak{}operators/\allowbreak{}test17.yr}:

%% check: skip
\begin{lstlisting}[style=coloredverbatimError]
class A {
    pub self () {}

    pub fn opIndexAssign (mut self, _ : i32, _ : i32) {}

}


fn main () {
    let dmut a = copy A ();
    a [0] = 0;
    a:[0] = 0;

    let b = copy A ();
    b [0] = 0;
    b:[0] = 0;
}
\end{lstlisting}

\begin{lstlisting}[style=bashVerb, breaklines=true, escapechar=~]
Error[E4099] : index assignment of a value of type &(test17::A) requires a mutable value
 --> test_resources/lit_class/operators/test17.yr:(15,5)
15  ~\lstbox{\textSFxi{}}~     b [0] = 0;
    ~\lstbox{\textSFv{}}~     ^ ^  note: the index assignment calls the mutable method opIndexAssign, and thus implicitly aliases the value
    ~\lstbox{\textSFii{}}~~\lstbox{\textSFx{}}~~\lstbox{\textSFx{}}~~\lstbox{\textSFx{}}~~\lstbox{\textSFx{}}~~\lstbox{\textSFx{}}~~\lstbox{\textSFx{}}~~\lstbox{\textSFx{}}~
\end{lstlisting}

\subsubsection*{How to fix it}

Declare the indexed value mutable, cf. \hyperref[sec:codes:e4082]{E4082}.

\errorcode{E4125}{cannot alias an expand value, maybe alias and expand keywords are inverted?}
\label{sec:codes:e4125}

Raised during \textbf{validation}, from \texttt{ValidateErrorMessage::\allowbreak{}MISMATCH\_\allowbreak{}ALIAS\_\allowbreak{}EXPAND} (\textit{src/\allowbreak{}ymirc/\allowbreak{}semantic/\allowbreak{}validator/\allowbreak{}errors.yr}).

\subsubsection*{What it means}

\tokennolst{alias} was applied to an \tokennolst{expand} value. The two keywords are in the order that aliases the expansion rather than expanding the alias, which is almost never what is meant.

\subsubsection*{Example}

\textit{test\_\allowbreak{}resources/\allowbreak{}diagnostics/\allowbreak{}test138.yr}:

%% check: skip
\begin{lstlisting}[style=coloredverbatimError]
// alias and expand written the wrong way round
fn take (dmut _ : [i32], _ : [i32]) {}

fn main () {
    let mut z : (mut [mut i32], [i32]) = alias (copy [1, 2], copy [3, 4]);
    take (alias expand z);
}
\end{lstlisting}

\begin{lstlisting}[style=bashVerb, breaklines=true, escapechar=~]
Error[E4125] : cannot alias an expand value, maybe alias and expand keywords are inverted?
 --> test_resources/diagnostics/test138.yr:(6,11)
 6  ~\lstbox{\textSFxi{}}~     take (alias expand z);
    ~\lstbox{\textSFv{}}~           ^^^^^ ------
    ~\lstbox{\textSFxi{}}~
    = help: invert the two keywords
    ~\lstbox{\textSFxi{}}~
 6  -     take (alias expand z);
 6  +     take (expand alias z);
    ~\lstbox{\textSFxi{}}~
    = when validating test138::main -- (test_resources/diagnostics/test138.yr:(4,4))
    ~\lstbox{\textSFii{}}~~\lstbox{\textSFx{}}~~\lstbox{\textSFx{}}~~\lstbox{\textSFx{}}~~\lstbox{\textSFx{}}~~\lstbox{\textSFx{}}~~\lstbox{\textSFx{}}~~\lstbox{\textSFx{}}~
\end{lstlisting}

\subsubsection*{How to fix it}

Swap them: \tokennolst{expand alias v}.

\errorcode{E4126}{cannot alias a lazy value, maybe alias and lazy keywords are inverted?}
\label{sec:codes:e4126}

Raised during \textbf{validation}, from \texttt{ValidateErrorMessage::\allowbreak{}MISMATCH\_\allowbreak{}ALIAS\_\allowbreak{}LAZY} (\textit{src/\allowbreak{}ymirc/\allowbreak{}semantic/\allowbreak{}validator/\allowbreak{}errors.yr}).

\subsubsection*{What it means}

\tokennolst{alias} was applied to a \tokennolst{lazy} value; the keywords are inverted.

\subsubsection*{Example}

\textit{test\_\allowbreak{}resources/\allowbreak{}lazyness/\allowbreak{}test6.yr}:

%% check: skip
\begin{lstlisting}[style=coloredverbatimError]
fn main () {
    let dmut i = copy [1, 2, 3];
    let lazy dmut a = copy (lazy i);
    let lazy dmut b = alias (lazy i);
}
\end{lstlisting}

\begin{lstlisting}[style=bashVerb, breaklines=true, escapechar=~]
Error[E4126] : cannot alias a lazy value, maybe alias and lazy keywords are inverted?
 --> test_resources/lazyness/test6.yr:(4,23)
 4  ~\lstbox{\textSFxi{}}~     let lazy dmut b = alias (lazy i);
    ~\lstbox{\textSFv{}}~                       ^^^^^  ----
    ~\lstbox{\textSFxi{}}~
    = help: invert the two keywords
    ~\lstbox{\textSFxi{}}~
 4  -     let lazy dmut b = alias (lazy i);
 4  +     let lazy dmut b = lazy (alias i);
    ~\lstbox{\textSFii{}}~~\lstbox{\textSFx{}}~~\lstbox{\textSFx{}}~~\lstbox{\textSFx{}}~~\lstbox{\textSFx{}}~~\lstbox{\textSFx{}}~~\lstbox{\textSFx{}}~~\lstbox{\textSFx{}}~
\end{lstlisting}

\subsubsection*{How to fix it}

Swap them: \tokennolst{lazy alias v}.

\errorcode{E4127}{cannot copy an expand value, maybe copy and expand keywords are inverted?}
\label{sec:codes:e4127}

Raised during \textbf{validation}, from \texttt{ValidateErrorMessage::\allowbreak{}MISMATCH\_\allowbreak{}COPY\_\allowbreak{}EXPAND} (\textit{src/\allowbreak{}ymirc/\allowbreak{}semantic/\allowbreak{}validator/\allowbreak{}errors.yr}).

\subsubsection*{What it means}

\tokennolst{copy} was applied to an \tokennolst{expand} value; the keywords are inverted.

\subsubsection*{Example}

\textit{test\_\allowbreak{}resources/\allowbreak{}diagnostics/\allowbreak{}test139.yr}:

%% check: skip
\begin{lstlisting}[style=coloredverbatimError]
// copy and expand written the wrong way round
fn take (_ : i32, _ : i32) {}

fn main () {
    let x = [1, 2];
    take (copy expand x);
}
\end{lstlisting}

\begin{lstlisting}[style=bashVerb, breaklines=true, escapechar=~]
Error[E4127] : cannot copy an expand value, maybe copy and expand keywords are inverted?
 --> test_resources/diagnostics/test139.yr:(6,11)
 6  ~\lstbox{\textSFxi{}}~     take (copy expand x);
    ~\lstbox{\textSFv{}}~           ^^^^ ------
    ~\lstbox{\textSFxi{}}~
    = help: invert the two keywords
    ~\lstbox{\textSFxi{}}~
 6  -     take (copy expand x);
 6  +     take (expand copy x);
    ~\lstbox{\textSFxi{}}~
    = when validating test139::main -- (test_resources/diagnostics/test139.yr:(4,4))
    ~\lstbox{\textSFii{}}~~\lstbox{\textSFx{}}~~\lstbox{\textSFx{}}~~\lstbox{\textSFx{}}~~\lstbox{\textSFx{}}~~\lstbox{\textSFx{}}~~\lstbox{\textSFx{}}~~\lstbox{\textSFx{}}~
\end{lstlisting}

\subsubsection*{How to fix it}

Swap them: \tokennolst{expand copy v}.

\errorcode{E4128}{cannot copy a lazy value, maybe copy and lazy keywords are inverted?}
\label{sec:codes:e4128}

Raised during \textbf{validation}, from \texttt{ValidateErrorMessage::\allowbreak{}MISMATCH\_\allowbreak{}COPY\_\allowbreak{}LAZY} (\textit{src/\allowbreak{}ymirc/\allowbreak{}semantic/\allowbreak{}validator/\allowbreak{}errors.yr}).

\subsubsection*{What it means}

\tokennolst{copy} was applied to a \tokennolst{lazy} value; the keywords are inverted.

\subsubsection*{Example}

\textit{test\_\allowbreak{}resources/\allowbreak{}lazyness/\allowbreak{}test6.yr}:

%% check: skip
\begin{lstlisting}[style=coloredverbatimError]
fn main () {
    let dmut i = copy [1, 2, 3];
    let lazy dmut a = copy (lazy i);
    let lazy dmut b = alias (lazy i);
}
\end{lstlisting}

\begin{lstlisting}[style=bashVerb, breaklines=true, escapechar=~]
Error[E4128] : cannot copy a lazy value, maybe copy and lazy keywords are inverted?
 --> test_resources/lazyness/test6.yr:(3,23)
 3  ~\lstbox{\textSFxi{}}~     let lazy dmut a = copy (lazy i);
    ~\lstbox{\textSFv{}}~                       ^^^^  ----
    ~\lstbox{\textSFxi{}}~
    = help: invert the two keywords
    ~\lstbox{\textSFxi{}}~
 3  -     let lazy dmut a = copy (lazy i);
 3  +     let lazy dmut a = lazy (copy i);
    ~\lstbox{\textSFii{}}~~\lstbox{\textSFx{}}~~\lstbox{\textSFx{}}~~\lstbox{\textSFx{}}~~\lstbox{\textSFx{}}~~\lstbox{\textSFx{}}~~\lstbox{\textSFx{}}~~\lstbox{\textSFx{}}~
\end{lstlisting}

\subsubsection*{How to fix it}

Swap them: \tokennolst{lazy copy v}.

\errorcode{E4129}{cannot copy an expand value, maybe dcopy and expand keywords are inverted?}
\label{sec:codes:e4129}

Raised during \textbf{validation}, from \texttt{ValidateErrorMessage::\allowbreak{}MISMATCH\_\allowbreak{}DCOPY\_\allowbreak{}EXPAND} (\textit{src/\allowbreak{}ymirc/\allowbreak{}semantic/\allowbreak{}validator/\allowbreak{}errors.yr}).

\subsubsection*{What it means}

\tokennolst{dcopy} was applied to an \tokennolst{expand} value; the keywords are inverted.

\subsubsection*{Example}

\textit{test\_\allowbreak{}resources/\allowbreak{}diagnostics/\allowbreak{}test140.yr}:

%% check: skip
\begin{lstlisting}[style=coloredverbatimError]
// dcopy and expand written the wrong way round
fn take (_ : i32, _ : i32) {}

fn main () {
    let x = [1, 2];
    take (dcopy expand x);
}
\end{lstlisting}

\begin{lstlisting}[style=bashVerb, breaklines=true, escapechar=~]
Error[E4129] : cannot copy an expand value, maybe dcopy and expand keywords are inverted?
 --> test_resources/diagnostics/test140.yr:(6,11)
 6  ~\lstbox{\textSFxi{}}~     take (dcopy expand x);
    ~\lstbox{\textSFv{}}~           ^^^^^ ------
    ~\lstbox{\textSFxi{}}~
    = help: invert the two keywords
    ~\lstbox{\textSFxi{}}~
 6  -     take (dcopy expand x);
 6  +     take (expand dcopy x);
    ~\lstbox{\textSFxi{}}~
    = when validating test140::main -- (test_resources/diagnostics/test140.yr:(4,4))
    ~\lstbox{\textSFii{}}~~\lstbox{\textSFx{}}~~\lstbox{\textSFx{}}~~\lstbox{\textSFx{}}~~\lstbox{\textSFx{}}~~\lstbox{\textSFx{}}~~\lstbox{\textSFx{}}~~\lstbox{\textSFx{}}~
\end{lstlisting}

\subsubsection*{How to fix it}

Swap them: \tokennolst{expand dcopy v}.

\errorcode{E4130}{cannot copy a lazy value, maybe dcopy and lazy keywords are inverted?}
\label{sec:codes:e4130}

Raised during \textbf{validation}, from \texttt{ValidateErrorMessage::\allowbreak{}MISMATCH\_\allowbreak{}DCOPY\_\allowbreak{}LAZY} (\textit{src/\allowbreak{}ymirc/\allowbreak{}semantic/\allowbreak{}validator/\allowbreak{}errors.yr}).

\subsubsection*{What it means}

\tokennolst{dcopy} was applied to a \tokennolst{lazy} value; the keywords are inverted.

\subsubsection*{Example}

\textit{test\_\allowbreak{}resources/\allowbreak{}lazyness/\allowbreak{}test19.yr}:

%% check: skip
\begin{lstlisting}[style=coloredverbatimError]
// a deep copy of a lazy value reads the wrong way round, the same way a
// shallow one does
fn main () {
    let dmut i = copy [1, 2, 3];
    let lazy dmut a = dcopy (lazy i);
    a;
}
\end{lstlisting}

\begin{lstlisting}[style=bashVerb, breaklines=true, escapechar=~]
Error[E4130] : cannot copy a lazy value, maybe dcopy and lazy keywords are inverted?
 --> test_resources/lazyness/test19.yr:(5,23)
 5  ~\lstbox{\textSFxi{}}~     let lazy dmut a = dcopy (lazy i);
    ~\lstbox{\textSFv{}}~                       ^^^^^  ----
    ~\lstbox{\textSFxi{}}~
    = help: invert the two keywords
    ~\lstbox{\textSFxi{}}~
 5  -     let lazy dmut a = dcopy (lazy i);
 5  +     let lazy dmut a = lazy (dcopy i);
    ~\lstbox{\textSFxi{}}~
    = when validating test19::main -- (test_resources/lazyness/test19.yr:(3,4))
    ~\lstbox{\textSFii{}}~~\lstbox{\textSFx{}}~~\lstbox{\textSFx{}}~~\lstbox{\textSFx{}}~~\lstbox{\textSFx{}}~~\lstbox{\textSFx{}}~~\lstbox{\textSFx{}}~~\lstbox{\textSFx{}}~
\end{lstlisting}

\subsubsection*{How to fix it}

Swap them: \tokennolst{lazy dcopy v}.

\errorcode{E4137}{an iterator cannot be mutable, if it is not a reference or does not borrow mutable data}
\label{sec:codes:e4137}

Raised during \textbf{validation}, from \texttt{ValidateErrorMessage::\allowbreak{}MUTABLE\_\allowbreak{}CONST\_\allowbreak{}ITERATOR} (\textit{src/\allowbreak{}ymirc/\allowbreak{}semantic/\allowbreak{}validator/\allowbreak{}errors.yr}).

\subsubsection*{What it means}

A loop iterator is declared mutable while it neither is a reference nor borrows mutable data. Mutating it would change a copy, which the loop would then discard.

\subsubsection*{Example}

\textit{test\_\allowbreak{}resources/\allowbreak{}for\_\allowbreak{}loops/\allowbreak{}arrays/\allowbreak{}test2.yr}:

%% check: skip
\begin{lstlisting}[style=coloredverbatimError]
in test2;

fn main () {
    let dmut a = [1, 2, 3, 4];
    for mut elem : i32 in a {
        assert (a[i] == elem);
    }
}
\end{lstlisting}

\begin{lstlisting}[style=bashVerb, breaklines=true, escapechar=~]
Error[E4137] : an iterator cannot be mutable, if it is not a reference or does not borrow mutable data
 --> test_resources/for_loops/arrays/test2.yr:(5,9)
 5  ~\lstbox{\textSFxi{}}~     for mut elem : i32 in a {
    ~\lstbox{\textSFv{}}~         ^^^
    ~\lstbox{\textSFii{}}~~\lstbox{\textSFx{}}~~\lstbox{\textSFx{}}~~\lstbox{\textSFx{}}~~\lstbox{\textSFx{}}~~\lstbox{\textSFx{}}~~\lstbox{\textSFx{}}~~\lstbox{\textSFx{}}~
\end{lstlisting}

\subsubsection*{How to fix it}

Drop the \tokennolst{mut}, or iterate by reference if the elements are to be modified.

\errorcode{E4138}{a parameter cannot be mutable, if it is not a reference or does not borrow mutable data}
\label{sec:codes:e4138}

Raised during \textbf{validation}, from \texttt{ValidateErrorMessage::\allowbreak{}MUTABLE\_\allowbreak{}CONST\_\allowbreak{}PARAM} (\textit{src/\allowbreak{}ymirc/\allowbreak{}semantic/\allowbreak{}validator/\allowbreak{}errors.yr}).

\subsubsection*{What it means}

A parameter is declared mutable while it neither is a reference nor borrows mutable data, so the mutation would be invisible to the caller.

\subsubsection*{Example}

\textit{test\_\allowbreak{}resources/\allowbreak{}lazyness/\allowbreak{}test11.yr}:

%% check: skip
\begin{lstlisting}[style=coloredverbatimError]
fn foo (lazy mut a : i32) {}
\end{lstlisting}

\begin{lstlisting}[style=bashVerb, breaklines=true, escapechar=~]
Error[E4138] : a parameter cannot be mutable, if it is not a reference or does not borrow mutable data
 --> test_resources/lazyness/test11.yr:(1,14)
 1  ~\lstbox{\textSFxi{}}~ fn foo (lazy mut a : i32) {}
    ~\lstbox{\textSFv{}}~              ^^^
    ~\lstbox{\textSFii{}}~~\lstbox{\textSFx{}}~~\lstbox{\textSFx{}}~~\lstbox{\textSFx{}}~~\lstbox{\textSFx{}}~~\lstbox{\textSFx{}}~~\lstbox{\textSFx{}}~~\lstbox{\textSFx{}}~
\end{lstlisting}

\subsubsection*{How to fix it}

Drop the \tokennolst{mut}, or declare the parameter \tokennolst{ref} if the caller's value is to be modified.

\errorcode{E4140}{a lazy variable cannot be mutable, if it does not borrow mutable data}
\label{sec:codes:e4140}

Raised during \textbf{validation}, from \texttt{ValidateErrorMessage::\allowbreak{}MUTABLE\_\allowbreak{}LAZY\_\allowbreak{}VAR} (\textit{src/\allowbreak{}ymirc/\allowbreak{}semantic/\allowbreak{}validator/\allowbreak{}errors.yr}).

\subsubsection*{What it means}

A lazy binding is declared mutable while it does not borrow mutable data. It names a computation rather than a location: each read produces the value again in a fresh slot, so there is nothing a write could persist to.

A lazy \textit{global} is not in that case -- its value is stored inside the global itself and only computed on first read -- so \tokennolst{lazy mut A : i32 = 12;} is accepted and \tokennolst{A = 42;} writes to that storage.

\subsubsection*{Example}

\textit{test\_\allowbreak{}resources/\allowbreak{}lazyness/\allowbreak{}test20.yr}:

%% check: skip
\begin{lstlisting}[style=coloredverbatimError]
fn main () {
    let lazy mut a = lazy 1;
    a;
}
\end{lstlisting}

\begin{lstlisting}[style=bashVerb, breaklines=true, escapechar=~]
Error[E4140] : a lazy variable cannot be mutable, if it does not borrow mutable data
 --> test_resources/lazyness/test20.yr:(3,14)
 3  ~\lstbox{\textSFxi{}}~     let lazy mut a = lazy 1;
    ~\lstbox{\textSFv{}}~              ^^^
    ~\lstbox{\textSFii{}}~~\lstbox{\textSFx{}}~~\lstbox{\textSFx{}}~~\lstbox{\textSFx{}}~~\lstbox{\textSFx{}}~~\lstbox{\textSFx{}}~~\lstbox{\textSFx{}}~~\lstbox{\textSFx{}}~
\end{lstlisting}

\subsubsection*{How to fix it}

Drop the \tokennolst{mut}, or bind the value to an ordinary variable.

\errorcode{E4146}{\errhole{} is not an aliasable type}
\label{sec:codes:e4146}

Raised during \textbf{validation}, from \texttt{ValidateErrorMessage::\allowbreak{}NOT\_\allowbreak{}ALIASABLE} (\textit{src/\allowbreak{}ymirc/\allowbreak{}semantic/\allowbreak{}validator/\allowbreak{}errors.yr}).

\subsubsection*{What it means}

\tokennolst{alias} was applied to a type that does not borrow anything, so there is nothing to alias.

\subsubsection*{Example}

\textit{test\_\allowbreak{}resources/\allowbreak{}templates/\allowbreak{}implicit/\allowbreak{}test16.yr}:

%% check: skip
\begin{lstlisting}[style=coloredverbatimError]
fn foo {alias T} (a : T) {
    a;
}

fn main () {
    foo ([1, 2, 3]);
    foo ((1, 2, 3));
}
\end{lstlisting}

\begin{lstlisting}[style=bashVerb, breaklines=true, escapechar=~]
Error[E4146] : [i32 ; 3us] is not an aliasable type
 --> test_resources/templates/implicit/test16.yr:(1,15)
 1  ~\lstbox{\textSFxi{}}~ fn foo {alias T} (a : T) {
    ~\lstbox{\textSFv{}}~               ^
   ...
 --> test_resources/templates/implicit/test16.yr:(6,10)
 6  ~\lstbox{\textSFxi{}}~     foo ([1, 2, 3]);
    ~\lstbox{\textSFv{}}~          ^
    ~\lstbox{\textSFii{}}~~\lstbox{\textSFx{}}~~\lstbox{\textSFx{}}~~\lstbox{\textSFx{}}~~\lstbox{\textSFx{}}~~\lstbox{\textSFx{}}~~\lstbox{\textSFx{}}~~\lstbox{\textSFx{}}~
\end{lstlisting}

\subsubsection*{How to fix it}

Remove the \tokennolst{alias}.

\errorcode{E4151}{not a lvalue}
\label{sec:codes:e4151}

Raised during \textbf{validation}, from \texttt{ValidateErrorMessage::\allowbreak{}NOT\_\allowbreak{}A\_\allowbreak{}LVALUE} (\textit{src/\allowbreak{}ymirc/\allowbreak{}semantic/\allowbreak{}validator/\allowbreak{}errors.yr}).

\subsubsection*{What it means}

An expression was used where storage is required -- the left of an assignment, the target of a \tokennolst{ref}, the operand of an address-of -- and it has no storage of its own.

\subsubsection*{Example}

\textit{test\_\allowbreak{}resources/\allowbreak{}lit\_\allowbreak{}tuples/\allowbreak{}test20.yr}:

%% check: skip
\begin{lstlisting}[style=coloredverbatimError]
fn foo ()-> (i32, i32);

fn main () {
    let mut a = 0;

    (a, 1) = foo ();
}
\end{lstlisting}

\begin{lstlisting}[style=bashVerb, breaklines=true, escapechar=~]
Error[E4151] : not a lvalue
 --> test_resources/lit_tuples/test20.yr:(6,9)
 6  ~\lstbox{\textSFxi{}}~     (a, 1) = foo ();
    ~\lstbox{\textSFv{}}~         ^  ^
    ~\lstbox{\textSFii{}}~~\lstbox{\textSFx{}}~~\lstbox{\textSFx{}}~~\lstbox{\textSFx{}}~~\lstbox{\textSFx{}}~~\lstbox{\textSFx{}}~~\lstbox{\textSFx{}}~~\lstbox{\textSFx{}}~
\end{lstlisting}

\subsubsection*{How to fix it}

Bind the value to a variable and use that. The more specific diagnostics E4152 to E4157 say why a particular value has no storage.

\errorcode{E4153}{\errhole{} is a value iterator, and cannot be used as a lvalue}
\label{sec:codes:e4153}

Raised during \textbf{validation}, from \texttt{ValidateErrorMessage::\allowbreak{}NOT\_\allowbreak{}A\_\allowbreak{}LVALUE\_\allowbreak{}ITERATOR} (\textit{src/\allowbreak{}ymirc/\allowbreak{}semantic/\allowbreak{}validator/\allowbreak{}errors.yr}).

\subsubsection*{What it means}

A \tokennolst{for} loop iterator that yields values, not references, was used as an lvalue. Each iteration produces a fresh value, which is discarded at the end of the round.

\subsubsection*{Example}

\textit{test\_\allowbreak{}resources/\allowbreak{}for\_\allowbreak{}loops/\allowbreak{}tuples/\allowbreak{}test9.yr}:

%% check: skip
\begin{lstlisting}[style=coloredverbatimError]
fn main () {
    let dmut a = (copy [1, 2],);
    for dmut x in alias a {
        x = copy [8, 3];
    }
}
\end{lstlisting}

\begin{lstlisting}[style=bashVerb, breaklines=true, escapechar=~]
Error[E4153] : x is a value iterator, and cannot be used as a lvalue
 --> test_resources/for_loops/tuples/test9.yr:(4,9)
 4  ~\lstbox{\textSFxi{}}~         x = copy [8, 3];
    ~\lstbox{\textSFv{}}~         ^ ^
    ~\lstbox{\textSFii{}}~~\lstbox{\textSFx{}}~~\lstbox{\textSFx{}}~~\lstbox{\textSFx{}}~~\lstbox{\textSFx{}}~~\lstbox{\textSFx{}}~~\lstbox{\textSFx{}}~~\lstbox{\textSFx{}}~
\end{lstlisting}

\subsubsection*{How to fix it}

Iterate by reference so the loop yields the elements themselves, or collect the new values into a result.

\errorcode{E4154}{\errhole{} is a lazy variable, and cannot be used as a lvalue}
\label{sec:codes:e4154}

Raised during \textbf{validation}, from \texttt{ValidateErrorMessage::\allowbreak{}NOT\_\allowbreak{}A\_\allowbreak{}LVALUE\_\allowbreak{}LAZY} (\textit{src/\allowbreak{}ymirc/\allowbreak{}semantic/\allowbreak{}validator/\allowbreak{}errors.yr}).

\subsubsection*{What it means}

A lazy parameter or binding was used as an lvalue. It names a computation, not a location: each read produces the value again in a fresh slot, so there is no storage a write could persist to. A lazy global is not concerned -- its value is stored inside the global itself -- and a mutable one can be assigned.

\subsubsection*{Example}

\textit{test\_\allowbreak{}resources/\allowbreak{}lazyness/\allowbreak{}test7.yr}:

%% check: skip
\begin{lstlisting}[style=coloredverbatimError]
fn main () {
    let lazy dmut a = lazy copy [1, 2, 3];

    a = copy [1, 2, 3];
    a = lazy copy [2, 3, 4];
}
\end{lstlisting}

\begin{lstlisting}[style=bashVerb, breaklines=true, escapechar=~]
Error[E4154] : a is a lazy variable, and cannot be used as a lvalue
 --> test_resources/lazyness/test7.yr:(4,5)
 4  ~\lstbox{\textSFxi{}}~     a = copy [1, 2, 3];
    ~\lstbox{\textSFv{}}~     ^ ^
    ~\lstbox{\textSFii{}}~~\lstbox{\textSFx{}}~~\lstbox{\textSFx{}}~~\lstbox{\textSFx{}}~~\lstbox{\textSFx{}}~~\lstbox{\textSFx{}}~~\lstbox{\textSFx{}}~~\lstbox{\textSFx{}}~
\end{lstlisting}

\subsubsection*{How to fix it}

Bind the lazy value to an ordinary variable first.

\errorcode{E4155}{\errhole{} is a map access, and cannot be used as a lvalue}
\label{sec:codes:e4155}

Raised during \textbf{validation}, from \texttt{ValidateErrorMessage::\allowbreak{}NOT\_\allowbreak{}A\_\allowbreak{}LVALUE\_\allowbreak{}MAP\_\allowbreak{}ACCESS} (\textit{src/\allowbreak{}ymirc/\allowbreak{}semantic/\allowbreak{}validator/\allowbreak{}errors.yr}).

\subsubsection*{What it means}

A map access was used as an lvalue. A map lookup produces the value stored, not the slot it lives in, so writing to it would not reach the map.

\subsubsection*{Example}

\textit{test\_\allowbreak{}resources/\allowbreak{}lit\_\allowbreak{}map/\allowbreak{}test10.yr}:

%% check: skip
\begin{lstlisting}[style=coloredverbatimError]
fn main () {
    let dmut a : [[c8] => [i32]] = copy [];

    a ["foo"] = copy [0];
    if let Ok (ref dmut z) = ref a ["foo"] {
        z ~= [8];
    }
}
\end{lstlisting}

\begin{lstlisting}[style=bashVerb, breaklines=true, escapechar=~]
Error[E4155] : a ["foo"s8] is a map access, and cannot be used as a lvalue
 --> test_resources/lit_map/test10.yr:(5,30)
 5  ~\lstbox{\textSFxi{}}~     if let Ok (ref dmut z) = ref a ["foo"] {
    ~\lstbox{\textSFv{}}~                              ^^^   ^
    ~\lstbox{\textSFii{}}~~\lstbox{\textSFx{}}~~\lstbox{\textSFx{}}~~\lstbox{\textSFx{}}~~\lstbox{\textSFx{}}~~\lstbox{\textSFx{}}~~\lstbox{\textSFx{}}~~\lstbox{\textSFx{}}~
\end{lstlisting}

\subsubsection*{How to fix it}

Use the map's insertion operation to write, rather than assigning to the lookup.

\errorcode{E4156}{\errhole{} is a value parameter, and cannot be used as a lvalue}
\label{sec:codes:e4156}

Raised during \textbf{validation}, from \texttt{ValidateErrorMessage::\allowbreak{}NOT\_\allowbreak{}A\_\allowbreak{}LVALUE\_\allowbreak{}PARAM} (\textit{src/\allowbreak{}ymirc/\allowbreak{}semantic/\allowbreak{}validator/\allowbreak{}errors.yr}).

\subsubsection*{What it means}

A parameter passed by value was used as an lvalue in a way that expects the caller's storage.

\subsubsection*{Example}

\textit{test\_\allowbreak{}resources/\allowbreak{}diagnostics/\allowbreak{}test59.yr}:

%% check: skip
\begin{lstlisting}[style=coloredverbatimError]
fn foo (x: i32) {
    x = 1;
}

fn main () { foo (1); }
\end{lstlisting}

\begin{lstlisting}[style=bashVerb, breaklines=true, escapechar=~]
Error[E4156] : x is a value parameter, and cannot be used as a lvalue
 --> test_resources/diagnostics/test59.yr:(2,5)
 2  ~\lstbox{\textSFxi{}}~     x = 1;
    ~\lstbox{\textSFv{}}~     ^ ^
    ~\lstbox{\textSFii{}}~~\lstbox{\textSFx{}}~~\lstbox{\textSFx{}}~~\lstbox{\textSFx{}}~~\lstbox{\textSFx{}}~~\lstbox{\textSFx{}}~~\lstbox{\textSFx{}}~~\lstbox{\textSFx{}}~
\end{lstlisting}

\subsubsection*{How to fix it}

Declare the parameter \tokennolst{ref} if the caller's value has to be written, or bind the value to a local variable.

\errorcode{E4157}{value of type \errhole{} is not a lvalue}
\label{sec:codes:e4157}

Raised during \textbf{validation}, from \texttt{ValidateErrorMessage::\allowbreak{}NOT\_\allowbreak{}A\_\allowbreak{}LVALUE\_\allowbreak{}TYPE} (\textit{src/\allowbreak{}ymirc/\allowbreak{}semantic/\allowbreak{}validator/\allowbreak{}errors.yr}).

\subsubsection*{What it means}

A value of this type never has storage of its own -- a temporary, or the result of an expression.

\subsubsection*{Example}

\textit{test\_\allowbreak{}resources/\allowbreak{}function/\allowbreak{}test2.yr}:

%% check: skip
\begin{lstlisting}[style=coloredverbatimError]
fn foo (mut x : [mut i32]);

fn main () {
    let a = copy [1];
    foo (alias a);
}
\end{lstlisting}

\begin{lstlisting}[style=bashVerb, breaklines=true, escapechar=~]
Error[E4157] : value of type [i32] is not a lvalue
 --> test_resources/function/test2.yr:(5,10)
 5  ~\lstbox{\textSFxi{}}~     foo (alias a);
    ~\lstbox{\textSFv{}}~          ^^^^^ ^
    ~\lstbox{\textSFii{}}~~\lstbox{\textSFx{}}~~\lstbox{\textSFx{}}~~\lstbox{\textSFx{}}~~\lstbox{\textSFx{}}~~\lstbox{\textSFx{}}~~\lstbox{\textSFx{}}~~\lstbox{\textSFx{}}~
\end{lstlisting}

\subsubsection*{How to fix it}

Bind it to a variable and use that, cf. \hyperref[sec:codes:e4151]{E4151}.

\errorcode{E4165}{type \errhole{} is not a movable type}
\label{sec:codes:e4165}

Raised during \textbf{validation}, from \texttt{ValidateErrorMessage::\allowbreak{}NOT\_\allowbreak{}MOVABLE} (\textit{src/\allowbreak{}ymirc/\allowbreak{}semantic/\allowbreak{}validator/\allowbreak{}errors.yr}).

\subsubsection*{What it means}

\tokennolst{move} was applied to a type that is not movable. Moving is defined only for types that carry an identity worth transferring.

\subsubsection*{Example}

\textit{test\_\allowbreak{}resources/\allowbreak{}diagnostics/\allowbreak{}test60.yr}:

%% check: skip
\begin{lstlisting}[style=coloredverbatimError]
fn main () {
    let x = 1;
    let _ = move x;
}
\end{lstlisting}

\begin{lstlisting}[style=bashVerb, breaklines=true, escapechar=~]
Error[E4165] : type i32 is not a movable type
 --> test_resources/diagnostics/test60.yr:(3,13)
 3  ~\lstbox{\textSFxi{}}~     let _ = move x;
    ~\lstbox{\textSFv{}}~             ^^^^
    ~\lstbox{\textSFii{}}~~\lstbox{\textSFx{}}~~\lstbox{\textSFx{}}~~\lstbox{\textSFx{}}~~\lstbox{\textSFx{}}~~\lstbox{\textSFx{}}~~\lstbox{\textSFx{}}~~\lstbox{\textSFx{}}~
\end{lstlisting}

\subsubsection*{How to fix it}

Copy the value instead, or make the type movable.

\errorcode{E4167}{no copy exists for type \errhole{}}
\label{sec:codes:e4167}

Raised during \textbf{validation}, from \texttt{ValidateErrorMessage::\allowbreak{}NO\_\allowbreak{}COPY\_\allowbreak{}EXIST} (\textit{src/\allowbreak{}ymirc/\allowbreak{}semantic/\allowbreak{}validator/\allowbreak{}errors.yr}).

\subsubsection*{What it means}

\tokennolst{copy} was applied to a type that defines no copy.

\subsubsection*{Example}

\textit{test\_\allowbreak{}resources/\allowbreak{}diagnostics/\allowbreak{}test115.yr}:

%% check: skip
\begin{lstlisting}[style=coloredverbatimError]
// a copy of a value that holds nothing to copy
fn main () {
    let _ = copy 1;
}
\end{lstlisting}

\begin{lstlisting}[style=bashVerb, breaklines=true, escapechar=~]
Error[E4167] : no copy exists for type i32
 --> test_resources/diagnostics/test115.yr:(3,13)
 3  ~\lstbox{\textSFxi{}}~     let _ = copy 1;
    ~\lstbox{\textSFv{}}~             ^^^^
    ~\lstbox{\textSFxi{}}~
    = when validating test115::main -- (test_resources/diagnostics/test115.yr:(2,4))
    ~\lstbox{\textSFii{}}~~\lstbox{\textSFx{}}~~\lstbox{\textSFx{}}~~\lstbox{\textSFx{}}~~\lstbox{\textSFx{}}~~\lstbox{\textSFx{}}~~\lstbox{\textSFx{}}~~\lstbox{\textSFx{}}~
\end{lstlisting}

\subsubsection*{How to fix it}

Define the copy operation on the type, or pass the value without copying if a borrow is enough.

\errorcode{E4247}{aliasing the value of type \errhole{} to create a constant borrowing is prohibited}
\label{sec:codes:e4247}

Raised during \textbf{validation}, from \texttt{ValidateErrorMessage::\allowbreak{}UNECESSARY\_\allowbreak{}ALIAS} (\textit{src/\allowbreak{}ymirc/\allowbreak{}semantic/\allowbreak{}validator/\allowbreak{}errors.yr}).

\subsubsection*{What it means}

\tokennolst{alias} was used to produce a constant borrow. Aliasing exists to carry mutability across a borrow; producing a constant one from it says the borrow is not needed.

\subsubsection*{Example}

\textit{test\_\allowbreak{}resources/\allowbreak{}function/\allowbreak{}test3.yr}:

%% check: skip
\begin{lstlisting}[style=coloredverbatimError]
fn foo (x : [i32]);

fn main () {
    let dmut a = copy [1];
    foo (alias a);
}
\end{lstlisting}

\begin{lstlisting}[style=bashVerb, breaklines=true, escapechar=~]
Error[E4247] : aliasing the value of type [i32] to create a constant borrowing is prohibited
 --> test_resources/function/test3.yr:(5,10)
 5  ~\lstbox{\textSFxi{}}~     foo (alias a);
    ~\lstbox{\textSFv{}}~          ^^^^^  note: for parameter x: [i32]
    ~\lstbox{\textSFii{}}~~\lstbox{\textSFx{}}~~\lstbox{\textSFx{}}~~\lstbox{\textSFx{}}~~\lstbox{\textSFx{}}~~\lstbox{\textSFx{}}~~\lstbox{\textSFx{}}~~\lstbox{\textSFx{}}~
\end{lstlisting}

\subsubsection*{How to fix it}

Drop the \tokennolst{alias}: a constant borrow is what passing the value already gives.

\errorcode{E4248}{aliasing the value of type \errhole{} to call begin and end iterator constant methods is useless}
\label{sec:codes:e4248}

Raised during \textbf{validation}, from \texttt{ValidateErrorMessage::\allowbreak{}UNECESSARY\_\allowbreak{}ALIAS\_\allowbreak{}CPTR\_\allowbreak{}LOOP} (\textit{src/\allowbreak{}ymirc/\allowbreak{}semantic/\allowbreak{}validator/\allowbreak{}errors.yr}).

\subsubsection*{What it means}

\tokennolst{alias} was used to call the constant \tokennolst{begin} and \tokennolst{end} iterator methods of a loop, which do not need the mutability the alias carries.

\subsubsection*{How to fix it}

Drop the \tokennolst{alias} on the loop subject.

\errorcode{E4251}{the move has no effect}
\label{sec:codes:e4251}

Raised during \textbf{validation}, from \texttt{ValidateErrorMessage::\allowbreak{}UNECESSARY\_\allowbreak{}MOVE} (\textit{src/\allowbreak{}ymirc/\allowbreak{}semantic/\allowbreak{}validator/\allowbreak{}errors.yr}).

\subsubsection*{What it means}

\tokennolst{move} was applied where the source is not used afterwards anyway, so the move changes nothing.

\subsubsection*{Example}

\textit{test\_\allowbreak{}resources/\allowbreak{}regressions/\allowbreak{}move\_\allowbreak{}unnecessary\_\allowbreak{}ctor\_\allowbreak{}call/\allowbreak{}test1.yr}:

%% check: skip
\begin{lstlisting}[style=coloredverbatimError]
entity A {
    pub self () {}

    __dtor (mut self) {}
}

fn foo ()-> A {
    move A ()
}

fn main () {
    let _ = foo ();
}
\end{lstlisting}

\begin{lstlisting}[style=bashVerb, breaklines=true, escapechar=~]
Error[E4251] : the move has no effect
 --> test_resources/regressions/move_unnecessary_ctor_call/test1.yr:(8,5)
 8  ~\lstbox{\textSFxi{}}~     move A ()
    ~\lstbox{\textSFv{}}~     ^^^^   ^
    ~\lstbox{\textSFii{}}~~\lstbox{\textSFx{}}~~\lstbox{\textSFx{}}~~\lstbox{\textSFx{}}~~\lstbox{\textSFx{}}~~\lstbox{\textSFx{}}~~\lstbox{\textSFx{}}~~\lstbox{\textSFx{}}~
\end{lstlisting}

\subsubsection*{How to fix it}

Drop the \tokennolst{move}.

\errorcode{E4252}{referencing the value has no effect}
\label{sec:codes:e4252}

Raised during \textbf{validation}, from \texttt{ValidateErrorMessage::\allowbreak{}UNECESSARY\_\allowbreak{}REFERENCE} (\textit{src/\allowbreak{}ymirc/\allowbreak{}semantic/\allowbreak{}validator/\allowbreak{}errors.yr}).

\subsubsection*{What it means}

A reference was taken where the value is used as a value anyway.

\subsubsection*{Example}

\textit{test\_\allowbreak{}resources/\allowbreak{}for\_\allowbreak{}loops/\allowbreak{}arrays/\allowbreak{}test13.yr}:

%% check: skip
\begin{lstlisting}[style=coloredverbatimError]
fn main () {
    let dmut a = [1, 2, 3, 4];
    for x in ref a {
        x;
    }
}
\end{lstlisting}

\begin{lstlisting}[style=bashVerb, breaklines=true, escapechar=~]
Error[E4252] : referencing the value has no effect
 --> test_resources/for_loops/arrays/test13.yr:(3,14)
 3  ~\lstbox{\textSFxi{}}~     for x in ref a {
    ~\lstbox{\textSFv{}}~              ^^^
    ~\lstbox{\textSFii{}}~~\lstbox{\textSFx{}}~~\lstbox{\textSFx{}}~~\lstbox{\textSFx{}}~~\lstbox{\textSFx{}}~~\lstbox{\textSFx{}}~~\lstbox{\textSFx{}}~~\lstbox{\textSFx{}}~
\end{lstlisting}

\subsubsection*{How to fix it}

Drop the reference.

%% Generated by tools/gen_error_codes.py from docs/errors/ of the compiler;
%% edit the compiler's pages or tools/error_themes.txt, then rerun it.

\clearpage
\section{Unsafe code}
\label{sec:codes:unsafe}

\begin{longtable}{@{}l p{0.78\linewidth} r@{}}
\toprule Code & Message & Page\\ \midrule \endhead
\hyperref[sec:codes:e4136]{E4136} & context is already unsafe & \pageref{sec:codes:e4136}\\
\hyperref[sec:codes:e4166]{E4166} & an unsafe context is entered, but no unsafe operations are made & \pageref{sec:codes:e4166}\\
\hyperref[sec:codes:e4262]{E4262} & call of unsafe function outside unsafe context & \pageref{sec:codes:e4262}\\
\hyperref[sec:codes:e4263]{E4263} & unsafe operation outside unsafe context & \pageref{sec:codes:e4263}\\
\bottomrule
\end{longtable}

\errorcode{E4136}{context is already unsafe}
\label{sec:codes:e4136}

Raised during \textbf{validation}, from \texttt{ValidateErrorMessage::\allowbreak{}MULTIPLE\_\allowbreak{}UNSAFE} (\textit{src/\allowbreak{}ymirc/\allowbreak{}semantic/\allowbreak{}validator/\allowbreak{}errors.yr}).

\subsubsection*{What it means}

An \tokennolst{unsafe} block was opened inside another one, so it adds nothing.

\subsubsection*{Example}

\textit{test\_\allowbreak{}resources/\allowbreak{}diagnostics/\allowbreak{}test54.yr}:

%% check: skip
\begin{lstlisting}[style=coloredverbatimError]
fn main () {
    unsafe {
        unsafe {
        }
    }
}
\end{lstlisting}

\begin{lstlisting}[style=bashVerb, breaklines=true, escapechar=~]
Error[E4136] : context is already unsafe
 --> test_resources/diagnostics/test54.yr:(3,9)
 3  ~\lstbox{\textSFxi{}}~         unsafe {
    ~\lstbox{\textSFv{}}~         ^^^^^^
    ~\lstbox{\textSFxi{}}~ Note :
    ~\lstbox{\textSFxi{}}~  --> test_resources/diagnostics/test54.yr:(2,5)
    ~\lstbox{\textSFxi{}}~  2  ~\lstbox{\textSFxi{}}~     unsafe {
    ~\lstbox{\textSFxi{}}~     ~\lstbox{\textSFv{}}~     ^^^^^^
    ~\lstbox{\textSFxi{}}~     ~\lstbox{\textSFii{}}~~\lstbox{\textSFx{}}~~\lstbox{\textSFx{}}~~\lstbox{\textSFx{}}~~\lstbox{\textSFx{}}~~\lstbox{\textSFx{}}~~\lstbox{\textSFx{}}~~\lstbox{\textSFx{}}~
    ~\lstbox{\textSFii{}}~~\lstbox{\textSFx{}}~~\lstbox{\textSFx{}}~~\lstbox{\textSFx{}}~~\lstbox{\textSFx{}}~~\lstbox{\textSFx{}}~~\lstbox{\textSFvii{}}~~\lstbox{\textSFx{}}~
\end{lstlisting}

\subsubsection*{How to fix it}

Remove the inner block.

\errorcode{E4166}{an unsafe context is entered, but no unsafe operations are made}
\label{sec:codes:e4166}

Raised during \textbf{validation}, from \texttt{ValidateErrorMessage::\allowbreak{}NOT\_\allowbreak{}UNSAFE} (\textit{src/\allowbreak{}ymirc/\allowbreak{}semantic/\allowbreak{}validator/\allowbreak{}errors.yr}).

\subsubsection*{What it means}

An \tokennolst{unsafe} block contains no operation that needed it. A block that says it is unsafe but is not tells a reader the wrong thing.

\subsubsection*{Example}

\textit{test\_\allowbreak{}resources/\allowbreak{}diagnostics/\allowbreak{}test80.yr}:

%% check: skip
\begin{lstlisting}[style=coloredverbatimError]
@unsafe
fn foo () {}

fn main () { foo (); }
\end{lstlisting}

\begin{lstlisting}[style=bashVerb, breaklines=true, escapechar=~]
Error[E4166] : an unsafe context is entered, but no unsafe operations are made
 --> test_resources/diagnostics/test80.yr:(1,2)
 1  ~\lstbox{\textSFxi{}}~ @unsafe
    ~\lstbox{\textSFv{}}~  ^^^^^^
    ~\lstbox{\textSFii{}}~~\lstbox{\textSFx{}}~~\lstbox{\textSFx{}}~~\lstbox{\textSFx{}}~~\lstbox{\textSFx{}}~~\lstbox{\textSFx{}}~~\lstbox{\textSFx{}}~~\lstbox{\textSFx{}}~
\end{lstlisting}

\subsubsection*{How to fix it}

Remove the \tokennolst{unsafe}, or narrow it to the operation that actually needs it.

\errorcode{E4262}{call of unsafe function outside unsafe context}
\label{sec:codes:e4262}

Raised during \textbf{validation}, from \texttt{ValidateErrorMessage::\allowbreak{}UNSAFE\_\allowbreak{}CALL} (\textit{src/\allowbreak{}ymirc/\allowbreak{}semantic/\allowbreak{}validator/\allowbreak{}errors.yr}).

\subsubsection*{What it means}

An unsafe function was called outside an \tokennolst{unsafe} context. Calling it is allowed, but the call has to say so.

\subsubsection*{Example}

\textit{test\_\allowbreak{}resources/\allowbreak{}diagnostics/\allowbreak{}test80.yr}:

%% check: skip
\begin{lstlisting}[style=coloredverbatimError]
@unsafe
fn foo () {}

fn main () { foo (); }
\end{lstlisting}

\begin{lstlisting}[style=bashVerb, breaklines=true, escapechar=~]
Error[E4262] : call of unsafe function outside unsafe context
 --> test_resources/diagnostics/test80.yr:(2,4)
 2  ~\lstbox{\textSFxi{}}~ fn foo () {}
    ~\lstbox{\textSFv{}}~    ^^^
    ~\lstbox{\textSFii{}}~~\lstbox{\textSFx{}}~~\lstbox{\textSFx{}}~~\lstbox{\textSFx{}}~~\lstbox{\textSFx{}}~~\lstbox{\textSFx{}}~~\lstbox{\textSFx{}}~~\lstbox{\textSFx{}}~
\end{lstlisting}

\subsubsection*{How to fix it}

Wrap the call in an \tokennolst{unsafe} block, or mark the calling function unsafe too.

\errorcode{E4263}{unsafe operation outside unsafe context}
\label{sec:codes:e4263}

Raised during \textbf{validation}, from \texttt{ValidateErrorMessage::\allowbreak{}UNSAFE\_\allowbreak{}OPERATION} (\textit{src/\allowbreak{}ymirc/\allowbreak{}semantic/\allowbreak{}validator/\allowbreak{}errors.yr}).

\subsubsection*{What it means}

An operation the language considers unsafe -- a raw pointer dereference, for instance -- was written outside an \tokennolst{unsafe} context.

\subsubsection*{Example}

\textit{test\_\allowbreak{}resources/\allowbreak{}diagnostics/\allowbreak{}test124.yr}:

%% check: skip
\begin{lstlisting}[style=coloredverbatimError]
// a pointer dereferenced outside an unsafe block
fn main () {
    let x = 1;
    let p = &x;
    let _ = *p;
}
\end{lstlisting}

\begin{lstlisting}[style=bashVerb, breaklines=true, escapechar=~]
Error[E4263] : unsafe operation outside unsafe context
 --> test_resources/diagnostics/test124.yr:(5,13)
 5  ~\lstbox{\textSFxi{}}~     let _ = *p;
    ~\lstbox{\textSFv{}}~             ^
    ~\lstbox{\textSFxi{}}~
    = when validating test124::main -- (test_resources/diagnostics/test124.yr:(2,4))
    ~\lstbox{\textSFii{}}~~\lstbox{\textSFx{}}~~\lstbox{\textSFx{}}~~\lstbox{\textSFx{}}~~\lstbox{\textSFx{}}~~\lstbox{\textSFx{}}~~\lstbox{\textSFx{}}~~\lstbox{\textSFx{}}~
\end{lstlisting}

\subsubsection*{How to fix it}

Wrap it in an \tokennolst{unsafe} block, cf. \hyperref[sec:codes:e4262]{E4262}.

%% Generated by tools/gen_error_codes.py from docs/errors/ of the compiler;
%% edit the compiler's pages or tools/error_themes.txt, then rerun it.

\clearpage
\section{Control flow and pattern matching}
\label{sec:codes:flow}

The book's chapter \textit{Control flows} specifies the rules behind these codes.

\begin{longtable}{@{}l p{0.78\linewidth} r@{}}
\toprule Code & Message & Page\\ \midrule \endhead
\hyperref[sec:codes:e4006]{E4006} & array pattern match len \errhole{} never matches the value len \errhole{} & \pageref{sec:codes:e4006}\\
\hyperref[sec:codes:e4007]{E4007} & break statement is not within a loop & \pageref{sec:codes:e4007}\\
\hyperref[sec:codes:e4008]{E4008} & cannot break loop inside a scope guard & \pageref{sec:codes:e4008}\\
\hyperref[sec:codes:e4009]{E4009} & continue statement is not within a loop & \pageref{sec:codes:e4009}\\
\hyperref[sec:codes:e4010]{E4010} & cannot continue loop inside a scope guard & \pageref{sec:codes:e4010}\\
\hyperref[sec:codes:e4059]{E4059} & expected a name expression & \pageref{sec:codes:e4059}\\
\hyperref[sec:codes:e4067]{E4067} & decorators are forbidden for index iterator, on \errhole{} iteration & \pageref{sec:codes:e4067}\\
\hyperref[sec:codes:e4072]{E4072} & overriding the for loop operation on a class type \errhole{} & \pageref{sec:codes:e4072}\\
\hyperref[sec:codes:e4081]{E4081} & the if condition has no else value, so it must produce a void value not a \errhole{} & \pageref{sec:codes:e4081}\\
\hyperref[sec:codes:e4100]{E4100} & infinite loop with always true test is not allowed, rather use a loop construct & \pageref{sec:codes:e4100}\\
\hyperref[sec:codes:e4123]{E4123} & the pattern matching has no default case, so each branch must produce a void value not a \errhole{} & \pageref{sec:codes:e4123}\\
\hyperref[sec:codes:e4143]{E4143} & loop test is always false, loop is never entered & \pageref{sec:codes:e4143}\\
\hyperref[sec:codes:e4177]{E4177} & option matcher \errhole{} only applies to option values not \errhole{} & \pageref{sec:codes:e4177}\\
\hyperref[sec:codes:e4193]{E4193} & the pattern \errhole{} with value \errhole{} is refutable & \pageref{sec:codes:e4193}\\
\hyperref[sec:codes:e4197]{E4197} & return statement is not within a function & \pageref{sec:codes:e4197}\\
\hyperref[sec:codes:e4198]{E4198} & cannot return inside a scope guard & \pageref{sec:codes:e4198}\\
\hyperref[sec:codes:e4216]{E4216} & pattern \errhole{} only matches tuple values, not \errhole{} & \pageref{sec:codes:e4216}\\
\hyperref[sec:codes:e4239]{E4239} & undefined cte for loop operator with \errhole{} iterator for type \errhole{} & \pageref{sec:codes:e4239}\\
\hyperref[sec:codes:e4244]{E4244} & undefined for loop operator with \errhole{} iterator for type \errhole{} & \pageref{sec:codes:e4244}\\
\hyperref[sec:codes:e4259]{E4259} & matcher expression is never evaluated & \pageref{sec:codes:e4259}\\
\hyperref[sec:codes:e4260]{E4260} & unreachable statement & \pageref{sec:codes:e4260}\\
\hyperref[sec:codes:e4266]{E4266} & conditional test is always evaluated to false, branch is never entered & \pageref{sec:codes:e4266}\\
\hyperref[sec:codes:e4268]{E4268} & conditional test is always evaluated to true, branch is always entered & \pageref{sec:codes:e4268}\\
\hyperref[sec:codes:e4269]{E4269} & useless runtime assertion for a test that is always true & \pageref{sec:codes:e4269}\\
\hyperref[sec:codes:e4318]{E4318} & the pattern matching over \errhole{} has no default case, so it must cover each of its members & \pageref{sec:codes:e4318}\\
\bottomrule
\end{longtable}

\errorcode{E4006}{array pattern match len \errhole{} never matches the value len \errhole{}}
\label{sec:codes:e4006}

Raised during \textbf{validation}, from \texttt{ValidateErrorMessage::\allowbreak{}ARRAY\_\allowbreak{}PATTERN\_\allowbreak{}SIZE} (\textit{src/\allowbreak{}ymirc/\allowbreak{}semantic/\allowbreak{}validator/\allowbreak{}errors.yr}).

\subsubsection*{What it means}

An array pattern matches a fixed number of elements, and the value it is matched against has a different, statically known length, so the pattern can never match.

\subsubsection*{Example}

\textit{test\_\allowbreak{}resources/\allowbreak{}pattern\_\allowbreak{}match/\allowbreak{}test13.yr}:

%% check: skip
\begin{lstlisting}[style=coloredverbatimError]
fn main () {
    let x = [1, 2, 3];
    let [a = 1, b] = x;
}
\end{lstlisting}

\begin{lstlisting}[style=bashVerb, breaklines=true, escapechar=~]
Error[E4006] : array pattern match len 2us never matches the value len 3us
 --> test_resources/pattern_match/test13.yr:(4,9)
 4  ~\lstbox{\textSFxi{}}~     let [a = 1, b] = x;
    ~\lstbox{\textSFv{}}~         ^
    ~\lstbox{\textSFii{}}~~\lstbox{\textSFx{}}~~\lstbox{\textSFx{}}~~\lstbox{\textSFx{}}~~\lstbox{\textSFx{}}~~\lstbox{\textSFx{}}~~\lstbox{\textSFx{}}~~\lstbox{\textSFx{}}~
\end{lstlisting}

\subsubsection*{How to fix it}

Give the pattern the length of the value, or use a slice pattern if a variable number of elements was meant.

\errorcode{E4007}{break statement is not within a loop}
\label{sec:codes:e4007}

Raised during \textbf{validation}, from \texttt{ValidateErrorMessage::\allowbreak{}BREAK\_\allowbreak{}NO\_\allowbreak{}LOOP} (\textit{src/\allowbreak{}ymirc/\allowbreak{}semantic/\allowbreak{}validator/\allowbreak{}errors.yr}).

\subsubsection*{What it means}

A \tokennolst{break} was written outside any loop.

\subsubsection*{Example}

\textit{test\_\allowbreak{}resources/\allowbreak{}diagnostics/\allowbreak{}test18.yr}:

%% check: skip
\begin{lstlisting}[style=coloredverbatimError]
fn main () {
    break;
}
\end{lstlisting}

\begin{lstlisting}[style=bashVerb, breaklines=true, escapechar=~]
Error[E4007] : break statement is not within a loop
 --> test_resources/diagnostics/test18.yr:(2,5)
 2  ~\lstbox{\textSFxi{}}~     break;
    ~\lstbox{\textSFv{}}~     ^^^^^
    ~\lstbox{\textSFii{}}~~\lstbox{\textSFx{}}~~\lstbox{\textSFx{}}~~\lstbox{\textSFx{}}~~\lstbox{\textSFx{}}~~\lstbox{\textSFx{}}~~\lstbox{\textSFx{}}~~\lstbox{\textSFx{}}~
\end{lstlisting}

\subsubsection*{How to fix it}

Remove it, or move it inside the loop it was meant to leave. To leave a function early, use \tokennolst{return}.

\errorcode{E4008}{cannot break loop inside a scope guard}
\label{sec:codes:e4008}

Raised during \textbf{validation}, from \texttt{ValidateErrorMessage::\allowbreak{}BREAK\_\allowbreak{}SCOPE\_\allowbreak{}GUARD} (\textit{src/\allowbreak{}ymirc/\allowbreak{}semantic/\allowbreak{}validator/\allowbreak{}errors.yr}).

\subsubsection*{What it means}

A \tokennolst{break} inside a scope guard would leave the loop while the guard is running. A guard runs as part of leaving the scope it guards, so jumping out of it would abandon a scope exit half-way through.

\subsubsection*{Example}

\textit{test\_\allowbreak{}resources/\allowbreak{}diagnostics/\allowbreak{}test19.yr}:

%% check: skip
\begin{lstlisting}[style=coloredverbatimError]
fn main () {
    loop {
        {
        } exit {
            break;
        }
    }
}
\end{lstlisting}

\begin{lstlisting}[style=bashVerb, breaklines=true, escapechar=~]
Error[E4008] : cannot break loop inside a scope guard
 --> test_resources/diagnostics/test19.yr:(4,11)
 4  ~\lstbox{\textSFxi{}}~         } exit {
    ~\lstbox{\textSFv{}}~           ^^^^
   ...
 --> test_resources/diagnostics/test19.yr:(5,13)
 5  ~\lstbox{\textSFxi{}}~             break;
    ~\lstbox{\textSFv{}}~             ^^^^^
    ~\lstbox{\textSFii{}}~~\lstbox{\textSFx{}}~~\lstbox{\textSFx{}}~~\lstbox{\textSFx{}}~~\lstbox{\textSFx{}}~~\lstbox{\textSFx{}}~~\lstbox{\textSFx{}}~~\lstbox{\textSFx{}}~
\end{lstlisting}

\subsubsection*{How to fix it}

Move the \tokennolst{break} out of the guard body. A guard that has to influence the loop should set a flag the loop tests instead.

\errorcode{E4009}{continue statement is not within a loop}
\label{sec:codes:e4009}

Raised during \textbf{validation}, from \texttt{ValidateErrorMessage::\allowbreak{}CONTINUE\_\allowbreak{}NO\_\allowbreak{}LOOP} (\textit{src/\allowbreak{}ymirc/\allowbreak{}semantic/\allowbreak{}validator/\allowbreak{}errors.yr}).

\subsubsection*{What it means}

A \tokennolst{continue} was written outside any loop.

\subsubsection*{Example}

\textit{test\_\allowbreak{}resources/\allowbreak{}diagnostics/\allowbreak{}test20.yr}:

%% check: skip
\begin{lstlisting}[style=coloredverbatimError]
fn main () {
    continue;
}
\end{lstlisting}

\begin{lstlisting}[style=bashVerb, breaklines=true, escapechar=~]
Error[E4009] : continue statement is not within a loop
 --> test_resources/diagnostics/test20.yr:(2,5)
 2  ~\lstbox{\textSFxi{}}~     continue;
    ~\lstbox{\textSFv{}}~     ^^^^^^^^
    ~\lstbox{\textSFii{}}~~\lstbox{\textSFx{}}~~\lstbox{\textSFx{}}~~\lstbox{\textSFx{}}~~\lstbox{\textSFx{}}~~\lstbox{\textSFx{}}~~\lstbox{\textSFx{}}~~\lstbox{\textSFx{}}~
\end{lstlisting}

\subsubsection*{How to fix it}

Remove it, or move it inside the loop it was meant to restart.

\errorcode{E4010}{cannot continue loop inside a scope guard}
\label{sec:codes:e4010}

Raised during \textbf{validation}, from \texttt{ValidateErrorMessage::\allowbreak{}CONTINUE\_\allowbreak{}SCOPE\_\allowbreak{}GUARD} (\textit{src/\allowbreak{}ymirc/\allowbreak{}semantic/\allowbreak{}validator/\allowbreak{}errors.yr}).

\subsubsection*{What it means}

A \tokennolst{continue} inside a scope guard would jump out of the guard while it is running, cf. \hyperref[sec:codes:e4008]{E4008}.

\subsubsection*{Example}

\textit{test\_\allowbreak{}resources/\allowbreak{}diagnostics/\allowbreak{}test21.yr}:

%% check: skip
\begin{lstlisting}[style=coloredverbatimError]
fn main () {
    loop {
        {
        } exit {
            continue;
        }
    }
}
\end{lstlisting}

\begin{lstlisting}[style=bashVerb, breaklines=true, escapechar=~]
Error[E4010] : cannot continue loop inside a scope guard
 --> test_resources/diagnostics/test21.yr:(4,11)
 4  ~\lstbox{\textSFxi{}}~         } exit {
    ~\lstbox{\textSFv{}}~           ^^^^
   ...
 --> test_resources/diagnostics/test21.yr:(5,13)
 5  ~\lstbox{\textSFxi{}}~             continue;
    ~\lstbox{\textSFv{}}~             ^^^^^^^^
    ~\lstbox{\textSFii{}}~~\lstbox{\textSFx{}}~~\lstbox{\textSFx{}}~~\lstbox{\textSFx{}}~~\lstbox{\textSFx{}}~~\lstbox{\textSFx{}}~~\lstbox{\textSFx{}}~~\lstbox{\textSFx{}}~
\end{lstlisting}

\subsubsection*{How to fix it}

Move the \tokennolst{continue} out of the guard body.

\errorcode{E4059}{expected a name expression}
\label{sec:codes:e4059}

Raised during \textbf{validation}, from \texttt{ValidateErrorMessage::\allowbreak{}FIELD\_\allowbreak{}MATCHER\_\allowbreak{}NO\_\allowbreak{}NAME} (\textit{src/\allowbreak{}ymirc/\allowbreak{}semantic/\allowbreak{}validator/\allowbreak{}errors.yr}).

\subsubsection*{What it means}

A field pattern expected the name of a field and got something that is not a name.

\subsubsection*{Example}

\textit{test\_\allowbreak{}resources/\allowbreak{}diagnostics/\allowbreak{}test131.yr}:

%% check: skip
\begin{lstlisting}[style=coloredverbatimError]
// a record pattern with a field given by position
record Point {
    pub let x : i32;
    pub self (x : i32) with x = x {}
}

fn main () {
    let p = Point (1);
    match p {
        Point (1) => {}
        _ => {}
    }
}
\end{lstlisting}

\begin{lstlisting}[style=bashVerb, breaklines=true, escapechar=~]
Error[E4059] : expected a name expression
 --> test_resources/diagnostics/test131.yr:(10,16)
10  ~\lstbox{\textSFxi{}}~         Point (1) => {}
    ~\lstbox{\textSFv{}}~                ^
    ~\lstbox{\textSFxi{}}~
    = when validating pattern matching with value p of type test131::Point -- (test_resources/diagnostics/test131.yr:(9,5))
    = when validating test131::main -- (test_resources/diagnostics/test131.yr:(7,4))
    ~\lstbox{\textSFii{}}~~\lstbox{\textSFx{}}~~\lstbox{\textSFx{}}~~\lstbox{\textSFx{}}~~\lstbox{\textSFx{}}~~\lstbox{\textSFx{}}~~\lstbox{\textSFx{}}~~\lstbox{\textSFx{}}~
\end{lstlisting}

\subsubsection*{How to fix it}

Write the field name. A pattern matching by position rather than by name uses the tuple form instead.

\errorcode{E4067}{decorators are forbidden for index iterator, on \errhole{} iteration}
\label{sec:codes:e4067}

Raised during \textbf{validation}, from \texttt{ValidateErrorMessage::\allowbreak{}FORBID\_\allowbreak{}DECO\_\allowbreak{}FOR\_\allowbreak{}LOOP} (\textit{src/\allowbreak{}ymirc/\allowbreak{}semantic/\allowbreak{}validator/\allowbreak{}errors.yr}).

\subsubsection*{What it means}

A decorator was written on the index variable of a \tokennolst{for} loop. The index is produced by the loop itself, so its mutability is not the loop body's to choose.

\subsubsection*{Example}

\textit{test\_\allowbreak{}resources/\allowbreak{}for\_\allowbreak{}loops/\allowbreak{}arrays/\allowbreak{}test12.yr}:

%% check: skip
\begin{lstlisting}[style=coloredverbatimError]
fn main () {
    let a = [1, 2, 3];
    for mut i, _ in a {
        i;
    }
}
\end{lstlisting}

\begin{lstlisting}[style=bashVerb, breaklines=true, escapechar=~]
Error[E4067] : decorators are forbidden for index iterator, on [i32] iteration
 --> test_resources/for_loops/arrays/test12.yr:(3,9)
 3  ~\lstbox{\textSFxi{}}~     for mut i, _ in a {
    ~\lstbox{\textSFv{}}~         ^^^ ^
    ~\lstbox{\textSFii{}}~~\lstbox{\textSFx{}}~~\lstbox{\textSFx{}}~~\lstbox{\textSFx{}}~~\lstbox{\textSFx{}}~~\lstbox{\textSFx{}}~~\lstbox{\textSFx{}}~~\lstbox{\textSFx{}}~
\end{lstlisting}

\subsubsection*{How to fix it}

Remove the decorator. Copy the index into a local variable if a mutable one is needed.

\errorcode{E4072}{overriding the for loop operation on a class type \errhole{}}
\label{sec:codes:e4072}

Raised during \textbf{validation}, from \texttt{ValidateErrorMessage::\allowbreak{}FOR\_\allowbreak{}LOOP\_\allowbreak{}CPTR} (\textit{src/\allowbreak{}ymirc/\allowbreak{}semantic/\allowbreak{}validator/\allowbreak{}errors.yr}).

\subsubsection*{What it means}

The \tokennolst{for} loop operators are being overridden on a class type; this reports the context in which the override is validated.

\subsubsection*{How to fix it}

Check the \tokennolst{begin}, \tokennolst{next} and \tokennolst{get} methods the loop needs, which the notes under this diagnostic point at.

\errorcode{E4081}{the if condition has no else value, so it must produce a void value not a \errhole{}}
\label{sec:codes:e4081}

Raised during \textbf{validation}, from \texttt{ValidateErrorMessage::\allowbreak{}IF\_\allowbreak{}COND\_\allowbreak{}NOT\_\allowbreak{}COMPLETE} (\textit{src/\allowbreak{}ymirc/\allowbreak{}semantic/\allowbreak{}validator/\allowbreak{}errors.yr}).

\subsubsection*{What it means}

An \tokennolst{if} used as an expression produces a value on the branch that is taken, but has no \tokennolst{else}, so there would be nothing to produce when the condition is false.

\subsubsection*{Example}

\textit{test\_\allowbreak{}resources/\allowbreak{}control\_\allowbreak{}flow/\allowbreak{}test16.yr}:

%% check: skip
\begin{lstlisting}[style=coloredverbatimError]
fn main () {
    let mut cond = false;
    let a = if cond {
        12
    }; // error, no else value

    let b = if cond {
        12
    } else {
        "str"
    }; // error, incompatible i32 and [c8]

    let d = if cond {
        12
    } else {
        11u32
    }; // ok, implicit cast of 11u32 to i32
}
\end{lstlisting}

\begin{lstlisting}[style=bashVerb, breaklines=true, escapechar=~]
Error[E4081] : the if condition has no else value, so it must produce a void value not a i32
 --> test_resources/control_flow/test16.yr:(3,13)
 3  ~\lstbox{\textSFxi{}}~     let a = if cond {
    ~\lstbox{\textSFv{}}~             ^^
   ...
 --> test_resources/control_flow/test16.yr:(4,9)
 4  ~\lstbox{\textSFxi{}}~         12
    ~\lstbox{\textSFv{}}~         ^^
    ~\lstbox{\textSFii{}}~~\lstbox{\textSFx{}}~~\lstbox{\textSFx{}}~~\lstbox{\textSFx{}}~~\lstbox{\textSFx{}}~~\lstbox{\textSFx{}}~~\lstbox{\textSFx{}}~~\lstbox{\textSFx{}}~
\end{lstlisting}

\subsubsection*{How to fix it}

Add an \tokennolst{else} branch producing a value of the same type, or make the \tokennolst{if} a statement by ending the branch with something that produces nothing.

\errorcode{E4100}{infinite loop with always true test is not allowed, rather use a loop construct}
\label{sec:codes:e4100}

Raised during \textbf{validation}, from \texttt{ValidateErrorMessage::\allowbreak{}INFINITE\_\allowbreak{}LOOP} (\textit{src/\allowbreak{}ymirc/\allowbreak{}semantic/\allowbreak{}validator/\allowbreak{}errors.yr}).

\subsubsection*{What it means}

A \tokennolst{while} loop has a condition that is always true. Written that way, the condition suggests a test that does not exist.

\subsubsection*{Example}

\textit{test\_\allowbreak{}resources/\allowbreak{}control\_\allowbreak{}flow/\allowbreak{}test3.yr}:

%% check: skip
\begin{lstlisting}[style=coloredverbatimError]
fn main () {
    let a = [1, 2, 3];
    if let [x, y, z] = a { // error, len of /a/ is /cte/
        x;
        y;
        z;
    }

    if (a.len <= 3) {
    }

    if (a.len >= 4) {
    }
    
    while a.len <= 3 {
    }

    while a.len >= 4 {
    }

}
\end{lstlisting}

\begin{lstlisting}[style=bashVerb, breaklines=true, escapechar=~]
Error[E4100] : infinite loop with always true test is not allowed, rather use a loop construct
 --> test_resources/control_flow/test3.yr:(15,17)
15  ~\lstbox{\textSFxi{}}~     while a.len <= 3 {
    ~\lstbox{\textSFv{}}~                 ^^
    ~\lstbox{\textSFii{}}~~\lstbox{\textSFx{}}~~\lstbox{\textSFx{}}~~\lstbox{\textSFx{}}~~\lstbox{\textSFx{}}~~\lstbox{\textSFx{}}~~\lstbox{\textSFx{}}~~\lstbox{\textSFx{}}~
\end{lstlisting}

\subsubsection*{How to fix it}

Use \tokennolst{loop}, which says the loop is infinite on purpose.

\errorcode{E4123}{the pattern matching has no default case, so each branch must produce a void value not a \errhole{}}
\label{sec:codes:e4123}

Raised during \textbf{validation}, from \texttt{ValidateErrorMessage::\allowbreak{}MATCH\_\allowbreak{}NOT\_\allowbreak{}COMPLETE} (\textit{src/\allowbreak{}ymirc/\allowbreak{}semantic/\allowbreak{}validator/\allowbreak{}errors.yr}).

\subsubsection*{What it means}

A pattern matching used as an expression has no default case, so a value not covered by any pattern would have nothing to produce.

\subsubsection*{Example}

\textit{test\_\allowbreak{}resources/\allowbreak{}diagnostics/\allowbreak{}test52.yr}:

%% check: skip
\begin{lstlisting}[style=coloredverbatimError]
fn main () {
    let x = 1;
    let _ = match x {
        1 => { 2 }
    };
}
\end{lstlisting}

\begin{lstlisting}[style=bashVerb, breaklines=true, escapechar=~]
Error[E4123] : the pattern matching has no default case, so each branch must produce a void value not a i32
 --> test_resources/diagnostics/test52.yr:(4,14)
 4  ~\lstbox{\textSFxi{}}~         1 => { 2 }
    ~\lstbox{\textSFv{}}~              ^
    ~\lstbox{\textSFxi{}}~
    = when validating pattern matching with value x of type i32 -- (test_resources/diagnostics/test52.yr:(3,13))
    = when validating test52::main -- (test_resources/diagnostics/test52.yr:(1,4))
    ~\lstbox{\textSFii{}}~~\lstbox{\textSFx{}}~~\lstbox{\textSFx{}}~~\lstbox{\textSFx{}}~~\lstbox{\textSFx{}}~~\lstbox{\textSFx{}}~~\lstbox{\textSFx{}}~~\lstbox{\textSFx{}}~
\end{lstlisting}

\subsubsection*{How to fix it}

Add a default case producing a value of the same type, or make the match a statement by having its branches produce nothing.

A match over a sealed type -- a union or an enum -- is not reported here: it must cover the type whatever its arms produce, and is reported by \hyperref[sec:codes:e4318]{E4318} instead.

\errorcode{E4143}{loop test is always false, loop is never entered}
\label{sec:codes:e4143}

Raised during \textbf{validation}, from \texttt{ValidateErrorMessage::\allowbreak{}NEVER\_\allowbreak{}ENTERED\_\allowbreak{}LOOP} (\textit{src/\allowbreak{}ymirc/\allowbreak{}semantic/\allowbreak{}validator/\allowbreak{}errors.yr}).

\subsubsection*{What it means}

A loop condition is always false, so the body never runs.

\subsubsection*{Example}

\textit{test\_\allowbreak{}resources/\allowbreak{}control\_\allowbreak{}flow/\allowbreak{}test3.yr}:

%% check: skip
\begin{lstlisting}[style=coloredverbatimError]
fn main () {
    let a = [1, 2, 3];
    if let [x, y, z] = a { // error, len of /a/ is /cte/
        x;
        y;
        z;
    }

    if (a.len <= 3) {
    }

    if (a.len >= 4) {
    }
    
    while a.len <= 3 {
    }

    while a.len >= 4 {
    }

}
\end{lstlisting}

\begin{lstlisting}[style=bashVerb, breaklines=true, escapechar=~]
Error[E4143] : loop test is always false, loop is never entered
 --> test_resources/control_flow/test3.yr:(18,17)
18  ~\lstbox{\textSFxi{}}~     while a.len >= 4 {
    ~\lstbox{\textSFv{}}~                 ^^
    ~\lstbox{\textSFii{}}~~\lstbox{\textSFx{}}~~\lstbox{\textSFx{}}~~\lstbox{\textSFx{}}~~\lstbox{\textSFx{}}~~\lstbox{\textSFx{}}~~\lstbox{\textSFx{}}~~\lstbox{\textSFx{}}~
\end{lstlisting}

\subsubsection*{How to fix it}

Remove the loop, or fix the condition. This most often means the variable the condition tests is not the one the loop updates.

\errorcode{E4177}{option matcher \errhole{} only applies to option values not \errhole{}}
\label{sec:codes:e4177}

Raised during \textbf{validation}, from \texttt{ValidateErrorMessage::\allowbreak{}OPTION\_\allowbreak{}MATCHER} (\textit{src/\allowbreak{}ymirc/\allowbreak{}semantic/\allowbreak{}validator/\allowbreak{}errors.yr}).

\subsubsection*{What it means}

An option pattern was matched against something that is not an option.

\subsubsection*{Example}

\textit{test\_\allowbreak{}resources/\allowbreak{}diagnostics/\allowbreak{}test132.yr}:

%% check: skip
\begin{lstlisting}[style=coloredverbatimError]
// an option pattern matched against a value that is not an option
fn main () {
    match 1 {
        Ok (x) => { let _ = x; }
        _ => {}
    }
}
\end{lstlisting}

\begin{lstlisting}[style=bashVerb, breaklines=true, escapechar=~]
Error[E4177] : option matcher Ok --> test_resources/diagnostics/test132.yr:(4,9) only applies to option values not i32
 --> test_resources/diagnostics/test132.yr:(4,9)
 4  ~\lstbox{\textSFxi{}}~         Ok (x) => { let _ = x; }
    ~\lstbox{\textSFv{}}~         ^^
    ~\lstbox{\textSFxi{}}~
    = when validating pattern matching with value 1 of type i32 -- (test_resources/diagnostics/test132.yr:(3,5))
    = when validating test132::main -- (test_resources/diagnostics/test132.yr:(2,4))
    ~\lstbox{\textSFii{}}~~\lstbox{\textSFx{}}~~\lstbox{\textSFx{}}~~\lstbox{\textSFx{}}~~\lstbox{\textSFx{}}~~\lstbox{\textSFx{}}~~\lstbox{\textSFx{}}~~\lstbox{\textSFx{}}~
\end{lstlisting}

\subsubsection*{How to fix it}

Match the option itself, or use a pattern for the type at hand.

\errorcode{E4193}{the pattern \errhole{} with value \errhole{} is refutable}
\label{sec:codes:e4193}

Raised during \textbf{validation}, from \texttt{ValidateErrorMessage::\allowbreak{}PATTERN\_\allowbreak{}IS\_\allowbreak{}REFUTABLE} (\textit{src/\allowbreak{}ymirc/\allowbreak{}semantic/\allowbreak{}validator/\allowbreak{}errors.yr}).

\subsubsection*{What it means}

A pattern is used where every value has to match -- a destructuring binding, for instance -- and this one can fail.

\subsubsection*{Example}

\textit{test\_\allowbreak{}resources/\allowbreak{}control\_\allowbreak{}flow/\allowbreak{}test8.yr}:

%% check: skip
\begin{lstlisting}[style=coloredverbatimError]
fn foo ()-> i32?;

fn main () {
    let Ok (x : i32) = foo (); // error, refutable since it may fail if /foo/ returns /none/
}
\end{lstlisting}

\begin{lstlisting}[style=bashVerb, breaklines=true, escapechar=~]
Error[E4193] : the pattern Ok(x: i32) with value test8::foo () is refutable
 --> test_resources/control_flow/test8.yr:(4,5)
 4  ~\lstbox{\textSFxi{}}~     let Ok (x : i32) = foo (); // error, refutable since it may fail if /foo/ returns /none/
    ~\lstbox{\textSFv{}}~     ^^^
    ~\lstbox{\textSFii{}}~~\lstbox{\textSFx{}}~~\lstbox{\textSFx{}}~~\lstbox{\textSFx{}}~~\lstbox{\textSFx{}}~~\lstbox{\textSFx{}}~~\lstbox{\textSFx{}}~~\lstbox{\textSFx{}}~
\end{lstlisting}

\subsubsection*{How to fix it}

Use a pattern that always matches, or move the destructuring into a \tokennolst{match} with a default case.

\errorcode{E4197}{return statement is not within a function}
\label{sec:codes:e4197}

Raised during \textbf{validation}, from \texttt{ValidateErrorMessage::\allowbreak{}RETURN\_\allowbreak{}NO\_\allowbreak{}FUNCTION} (\textit{src/\allowbreak{}ymirc/\allowbreak{}semantic/\allowbreak{}validator/\allowbreak{}errors.yr}).

\subsubsection*{What it means}

A \tokennolst{return} was written outside any function.

\subsubsection*{Example}

\textit{test\_\allowbreak{}resources/\allowbreak{}diagnostics/\allowbreak{}test69.yr}:

%% check: skip
\begin{lstlisting}[style=coloredverbatimError]
lazy X : i32 = { return 1; };

fn main () { X; }
\end{lstlisting}

\begin{lstlisting}[style=bashVerb, breaklines=true, escapechar=~]
Error[E4197] : return statement is not within a function
 --> test_resources/diagnostics/test69.yr:(1,18)
 1  ~\lstbox{\textSFxi{}}~ lazy X : i32 = { return 1; };
    ~\lstbox{\textSFv{}}~                  ^^^^^^
    ~\lstbox{\textSFxi{}}~
    = when validating -- (test_resources/diagnostics/test69.yr:(1,6))
    ~\lstbox{\textSFii{}}~~\lstbox{\textSFx{}}~~\lstbox{\textSFx{}}~~\lstbox{\textSFx{}}~~\lstbox{\textSFx{}}~~\lstbox{\textSFx{}}~~\lstbox{\textSFx{}}~~\lstbox{\textSFx{}}~
Error[E4104] : symbol test69::X has errors
 --> test_resources/diagnostics/test69.yr:(3,14)
 1  ~\lstbox{\textSFxi{}}~ lazy X : i32 = { return 1; };
    ~\lstbox{\textSFv{}}~                  ------  error: return statement is not within a function
 2  ~\lstbox{\textSFxi{}}~
 3  ~\lstbox{\textSFxi{}}~ fn main () { X; }
    ~\lstbox{\textSFv{}}~              ^
    ~\lstbox{\textSFxi{}}~
    = when validating -- (test_resources/diagnostics/test69.yr:(1,6))
    = when validating test69::main -- (test_resources/diagnostics/test69.yr:(3,4))
    ~\lstbox{\textSFii{}}~~\lstbox{\textSFx{}}~~\lstbox{\textSFx{}}~~\lstbox{\textSFx{}}~~\lstbox{\textSFx{}}~~\lstbox{\textSFx{}}~~\lstbox{\textSFx{}}~~\lstbox{\textSFx{}}~
\end{lstlisting}

\subsubsection*{How to fix it}

Remove it, or move it into the function it was meant to leave.

\errorcode{E4198}{cannot return inside a scope guard}
\label{sec:codes:e4198}

Raised during \textbf{validation}, from \texttt{ValidateErrorMessage::\allowbreak{}RETURN\_\allowbreak{}SCOPE\_\allowbreak{}GUARD} (\textit{src/\allowbreak{}ymirc/\allowbreak{}semantic/\allowbreak{}validator/\allowbreak{}errors.yr}).

\subsubsection*{What it means}

A \tokennolst{return} inside a scope guard would leave the function while the guard is running, abandoning a scope exit half-way through, cf. \hyperref[sec:codes:e4008]{E4008}.

\subsubsection*{Example}

\textit{test\_\allowbreak{}resources/\allowbreak{}scope\_\allowbreak{}guards/\allowbreak{}test7.yr}:

%% check: skip
\begin{lstlisting}[style=coloredverbatimError]
fn main () {
    {

    } exit {
        {
            throw copy AssertError ("");
        } catch {
            _ => return ;
        }
    }
}
\end{lstlisting}

\begin{lstlisting}[style=bashVerb, breaklines=true, escapechar=~]
Error[E4198] : cannot return inside a scope guard
 --> test_resources/scope_guards/test7.yr:(4,7)
 4  ~\lstbox{\textSFxi{}}~     } exit {
    ~\lstbox{\textSFv{}}~       ^^^^
 --> test_resources/scope_guards/test7.yr:(8,18)
 8  ~\lstbox{\textSFxi{}}~             _ => return ;
    ~\lstbox{\textSFv{}}~                  ------
    ~\lstbox{\textSFxi{}}~
    = when validating catching pattern with type &(core::exception::assertion::AssertError) -- (test_resources/scope_guards/test7.yr:(7,11))
    = when validating test7::main -- (test_resources/scope_guards/test7.yr:(1,4))
    ~\lstbox{\textSFii{}}~~\lstbox{\textSFx{}}~~\lstbox{\textSFx{}}~~\lstbox{\textSFx{}}~~\lstbox{\textSFx{}}~~\lstbox{\textSFx{}}~~\lstbox{\textSFx{}}~~\lstbox{\textSFx{}}~
\end{lstlisting}

\subsubsection*{How to fix it}

Move the \tokennolst{return} out of the guard body.

\errorcode{E4216}{pattern \errhole{} only matches tuple values, not \errhole{}}
\label{sec:codes:e4216}

Raised during \textbf{validation}, from \texttt{ValidateErrorMessage::\allowbreak{}TUPLE\_\allowbreak{}PATTERN} (\textit{src/\allowbreak{}ymirc/\allowbreak{}semantic/\allowbreak{}validator/\allowbreak{}errors.yr}).

\subsubsection*{What it means}

A pattern that destructures a tuple was applied to a value that is not one. Both the pattern and the type it was matched against are in the message.

\subsubsection*{Example}

\textit{test\_\allowbreak{}resources/\allowbreak{}diagnostics/\allowbreak{}test133.yr}:

%% check: skip
\begin{lstlisting}[style=coloredverbatimError]
// a tuple pattern matched against a value that is not a tuple
fn main () {
    match 1 {
        (x, y) => { let _ = (x, y); }
        _ => {}
    }
}
\end{lstlisting}

\begin{lstlisting}[style=bashVerb, breaklines=true, escapechar=~]
Error[E4216] : pattern (x, y) only matches tuple values, not i32
 --> test_resources/diagnostics/test133.yr:(4,9)
 4  ~\lstbox{\textSFxi{}}~         (x, y) => { let _ = (x, y); }
    ~\lstbox{\textSFv{}}~         ^
    ~\lstbox{\textSFxi{}}~
    = when validating pattern matching with value 1 of type i32 -- (test_resources/diagnostics/test133.yr:(3,5))
    = when validating test133::main -- (test_resources/diagnostics/test133.yr:(2,4))
    ~\lstbox{\textSFii{}}~~\lstbox{\textSFx{}}~~\lstbox{\textSFx{}}~~\lstbox{\textSFx{}}~~\lstbox{\textSFx{}}~~\lstbox{\textSFx{}}~~\lstbox{\textSFx{}}~~\lstbox{\textSFx{}}~
\end{lstlisting}

\subsubsection*{How to fix it}

Match a tuple, or use a pattern defined for the type at hand. A record is not a tuple: its fields are matched by name, not by position.

\errorcode{E4239}{undefined cte for loop operator with \errhole{} iterator for type \errhole{}}
\label{sec:codes:e4239}

Raised during \textbf{validation}, from \texttt{ValidateErrorMessage::\allowbreak{}UNDEF\_\allowbreak{}CTE\_\allowbreak{}FOR\_\allowbreak{}LOOP\_\allowbreak{}OPERATOR} (\textit{src/\allowbreak{}ymirc/\allowbreak{}semantic/\allowbreak{}validator/\allowbreak{}errors.yr}).

\subsubsection*{What it means}

A \tokennolst{cte for} loop was written over a type that defines no compile time iteration.

\subsubsection*{Example}

\textit{test\_\allowbreak{}resources/\allowbreak{}diagnostics/\allowbreak{}test123.yr}:

%% check: skip
\begin{lstlisting}[style=coloredverbatimError]
// a compile time loop over a slice, whose length is not known at compile time
fn main () {
    cte for i in copy [1, 2] {
        let _ = i;
    }
}
\end{lstlisting}

\begin{lstlisting}[style=bashVerb, breaklines=true, escapechar=~]
Error[E4239] : undefined cte for loop operator with 1 iterator for type mut [mut i32]
 --> test_resources/diagnostics/test123.yr:(3,9)
 3  ~\lstbox{\textSFxi{}}~     cte for i in copy [1, 2] {
    ~\lstbox{\textSFv{}}~         ^^^
    ~\lstbox{\textSFxi{}}~
    = when validating test123::main -- (test_resources/diagnostics/test123.yr:(2,4))
    ~\lstbox{\textSFii{}}~~\lstbox{\textSFx{}}~~\lstbox{\textSFx{}}~~\lstbox{\textSFx{}}~~\lstbox{\textSFx{}}~~\lstbox{\textSFx{}}~~\lstbox{\textSFx{}}~~\lstbox{\textSFx{}}~
\end{lstlisting}

\subsubsection*{How to fix it}

Iterate over a type that does -- a tuple, an array, or a compile time range.

\errorcode{E4244}{undefined for loop operator with \errhole{} iterator for type \errhole{}}
\label{sec:codes:e4244}

Raised during \textbf{validation}, from \texttt{ValidateErrorMessage::\allowbreak{}UNDEF\_\allowbreak{}FOR\_\allowbreak{}LOOP\_\allowbreak{}OPERATOR} (\textit{src/\allowbreak{}ymirc/\allowbreak{}semantic/\allowbreak{}validator/\allowbreak{}errors.yr}).

\subsubsection*{What it means}

A \tokennolst{for} loop was written over a type that defines no iteration. A type is iterable when it provides the \tokennolst{begin}, \tokennolst{next} and \tokennolst{get} operators.

\subsubsection*{Example}

\textit{test\_\allowbreak{}resources/\allowbreak{}diagnostics/\allowbreak{}test77.yr}:

%% check: skip
\begin{lstlisting}[style=coloredverbatimError]
fn main () {
    for i in 1 { i; }
}
\end{lstlisting}

\begin{lstlisting}[style=bashVerb, breaklines=true, escapechar=~]
Error[E4244] : undefined for loop operator with 1 iterator for type i32
 --> test_resources/diagnostics/test77.yr:(2,5)
 2  ~\lstbox{\textSFxi{}}~     for i in 1 { i; }
    ~\lstbox{\textSFv{}}~     ^^^
    ~\lstbox{\textSFii{}}~~\lstbox{\textSFx{}}~~\lstbox{\textSFx{}}~~\lstbox{\textSFx{}}~~\lstbox{\textSFx{}}~~\lstbox{\textSFx{}}~~\lstbox{\textSFx{}}~~\lstbox{\textSFx{}}~
\end{lstlisting}

\subsubsection*{How to fix it}

Define those operators on the type, or iterate over something that already has them.

\errorcode{E4259}{matcher expression is never evaluated}
\label{sec:codes:e4259}

Raised during \textbf{validation}, from \texttt{ValidateErrorMessage::\allowbreak{}UNREACHABLE\_\allowbreak{}MATCHER} (\textit{src/\allowbreak{}ymirc/\allowbreak{}semantic/\allowbreak{}validator/\allowbreak{}errors.yr}).

\subsubsection*{What it means}

A pattern in a \tokennolst{match} is covered by an earlier one, so it can never run. A default case raises it too when the arms before it already cover every member of the union, or every field of the enum, being matched.

\subsubsection*{How to fix it}

Remove it, or move it before the pattern that shadows it.

\errorcode{E4260}{unreachable statement}
\label{sec:codes:e4260}

Raised during \textbf{validation}, from \texttt{ValidateErrorMessage::\allowbreak{}UNREACHBLE\_\allowbreak{}STATEMENT} (\textit{src/\allowbreak{}ymirc/\allowbreak{}semantic/\allowbreak{}validator/\allowbreak{}errors.yr}).

\subsubsection*{What it means}

A statement follows one that never completes -- a \tokennolst{return}, a \tokennolst{break}, a \tokennolst{throw} -- so it can never run.

\subsubsection*{Example}

\textit{test\_\allowbreak{}resources/\allowbreak{}lit\_\allowbreak{}options/\allowbreak{}test16.yr}:

%% check: skip
\begin{lstlisting}[style=coloredverbatimError]
fn main () {
    loop {
        let _a_ = ({
            break 12;
        })?;
    };
}
\end{lstlisting}

\begin{lstlisting}[style=bashVerb, breaklines=true, escapechar=~]
Error[E4260] : unreachable statement
 --> test_resources/lit_options/test16.yr:(5,11)
 5  ~\lstbox{\textSFxi{}}~         })?;
    ~\lstbox{\textSFv{}}~           ^
    ~\lstbox{\textSFxi{}}~ Note : exiting scope here
    ~\lstbox{\textSFxi{}}~  --> test_resources/lit_options/test16.yr:(4,13)
    ~\lstbox{\textSFxi{}}~  4  ~\lstbox{\textSFxi{}}~             break 12;
    ~\lstbox{\textSFxi{}}~     ~\lstbox{\textSFv{}}~             ^^^^^
    ~\lstbox{\textSFxi{}}~     ~\lstbox{\textSFii{}}~~\lstbox{\textSFx{}}~~\lstbox{\textSFx{}}~~\lstbox{\textSFx{}}~~\lstbox{\textSFx{}}~~\lstbox{\textSFx{}}~~\lstbox{\textSFx{}}~~\lstbox{\textSFx{}}~
    ~\lstbox{\textSFii{}}~~\lstbox{\textSFx{}}~~\lstbox{\textSFx{}}~~\lstbox{\textSFx{}}~~\lstbox{\textSFx{}}~~\lstbox{\textSFx{}}~~\lstbox{\textSFvii{}}~~\lstbox{\textSFx{}}~
\end{lstlisting}

\subsubsection*{How to fix it}

Remove the statement, or move it before the one that leaves the block.

\errorcode{E4266}{conditional test is always evaluated to false, branch is never entered}
\label{sec:codes:e4266}

Raised during \textbf{validation}, from \texttt{ValidateErrorMessage::\allowbreak{}USELESS\_\allowbreak{}COND\_\allowbreak{}FALSE} (\textit{src/\allowbreak{}ymirc/\allowbreak{}semantic/\allowbreak{}validator/\allowbreak{}errors.yr}).

\subsubsection*{What it means}

A condition the compiler could evaluate is always false, so the branch it guards is dead.

\subsubsection*{Example}

\textit{test\_\allowbreak{}resources/\allowbreak{}control\_\allowbreak{}flow/\allowbreak{}test15.yr}:

%% check: skip
\begin{lstlisting}[style=coloredverbatimError]
fn main () {
    let a = [1, 2, 3];
    if (a.len <= 3) { // error, len of /a/ is /cte/
        println (a [2]);
    }

    if (a.len >= 4) {
        println (a [3]); // error, dead code
    }
}
\end{lstlisting}

\begin{lstlisting}[style=bashVerb, breaklines=true, escapechar=~]
Error[E4266] : conditional test is always evaluated to false, branch is never entered
 --> test_resources/control_flow/test15.yr:(7,15)
 7  ~\lstbox{\textSFxi{}}~     if (a.len >= 4) {
    ~\lstbox{\textSFv{}}~               ^^
    ~\lstbox{\textSFii{}}~~\lstbox{\textSFx{}}~~\lstbox{\textSFx{}}~~\lstbox{\textSFx{}}~~\lstbox{\textSFx{}}~~\lstbox{\textSFx{}}~~\lstbox{\textSFx{}}~~\lstbox{\textSFx{}}~
\end{lstlisting}

\subsubsection*{How to fix it}

Remove the branch, or fix the condition.

\errorcode{E4268}{conditional test is always evaluated to true, branch is always entered}
\label{sec:codes:e4268}

Raised during \textbf{validation}, from \texttt{ValidateErrorMessage::\allowbreak{}USELESS\_\allowbreak{}COND\_\allowbreak{}TRUE} (\textit{src/\allowbreak{}ymirc/\allowbreak{}semantic/\allowbreak{}validator/\allowbreak{}errors.yr}).

\subsubsection*{What it means}

A condition the compiler could evaluate is always true, so the branch it guards always runs and the test says nothing.

\subsubsection*{Example}

\textit{test\_\allowbreak{}resources/\allowbreak{}control\_\allowbreak{}flow/\allowbreak{}test15.yr}:

%% check: skip
\begin{lstlisting}[style=coloredverbatimError]
fn main () {
    let a = [1, 2, 3];
    if (a.len <= 3) { // error, len of /a/ is /cte/
        println (a [2]);
    }

    if (a.len >= 4) {
        println (a [3]); // error, dead code
    }
}
\end{lstlisting}

\begin{lstlisting}[style=bashVerb, breaklines=true, escapechar=~]
Error[E4268] : conditional test is always evaluated to true, branch is always entered
 --> test_resources/control_flow/test15.yr:(3,15)
 3  ~\lstbox{\textSFxi{}}~     if (a.len <= 3) { // error, len of /a/ is /cte/
    ~\lstbox{\textSFv{}}~               ^^
    ~\lstbox{\textSFii{}}~~\lstbox{\textSFx{}}~~\lstbox{\textSFx{}}~~\lstbox{\textSFx{}}~~\lstbox{\textSFx{}}~~\lstbox{\textSFx{}}~~\lstbox{\textSFx{}}~~\lstbox{\textSFx{}}~
\end{lstlisting}

\subsubsection*{How to fix it}

Remove the condition, or fix it.

\errorcode{E4269}{useless runtime assertion for a test that is always true}
\label{sec:codes:e4269}

Raised during \textbf{validation}, from \texttt{ValidateErrorMessage::\allowbreak{}USELESS\_\allowbreak{}RUNTIME\_\allowbreak{}ASSERT} (\textit{src/\allowbreak{}ymirc/\allowbreak{}semantic/\allowbreak{}validator/\allowbreak{}errors.yr}).

\subsubsection*{What it means}

A run time assertion tests something the compiler already proved true, so it can never fire.

\subsubsection*{Example}

\textit{test\_\allowbreak{}resources/\allowbreak{}diagnostics/\allowbreak{}test81.yr}:

%% check: skip
\begin{lstlisting}[style=coloredverbatimError]
fn main () {
    assert (true);
}
\end{lstlisting}

\begin{lstlisting}[style=bashVerb, breaklines=true, escapechar=~]
Error[E4269] : useless runtime assertion for a test that is always true
 --> test_resources/diagnostics/test81.yr:(2,5)
 2  ~\lstbox{\textSFxi{}}~     assert (true);
    ~\lstbox{\textSFv{}}~     ^^^^^^  ^^^^
    ~\lstbox{\textSFii{}}~~\lstbox{\textSFx{}}~~\lstbox{\textSFx{}}~~\lstbox{\textSFx{}}~~\lstbox{\textSFx{}}~~\lstbox{\textSFx{}}~~\lstbox{\textSFx{}}~~\lstbox{\textSFx{}}~
\end{lstlisting}

\subsubsection*{How to fix it}

Remove the assertion, or make it a \tokennolst{cte assert} if the compile time check is what was wanted.

\errorcode{E4318}{the pattern matching over \errhole{} has no default case, so it must cover each of its members}
\label{sec:codes:e4318}

Raised during \textbf{validation}, from \texttt{ValidateErrorMessage::\allowbreak{}MATCH\_\allowbreak{}SEALED\_\allowbreak{}NOT\_\allowbreak{}COMPLETE} (\textit{src/\allowbreak{}ymirc/\allowbreak{}semantic/\allowbreak{}validator/\allowbreak{}errors.yr}).

\subsubsection*{What it means}

A \tokennolst{match} over a union or an enum has no default case and leaves out a member of the union, or a field of the enum. A union and an enum list everything they can hold, so such a match must cover that list, whether its arms produce a value or not: adding a member to the type is then an error at every match that forgets it, instead of a value that silently falls through. The notes name each member or field no arm covers.

An arm carrying a guard, or a pattern filtering the fields of the member, decides on more than the member it names, so it does not count as covering it.

\subsubsection*{Example}

\textit{test\_\allowbreak{}resources/\allowbreak{}type\_\allowbreak{}union/\allowbreak{}test83.yr}:

%% check: skip
\begin{lstlisting}[style=coloredverbatimError]
// A match over a union must be complete even when its arms produce nothing, so
// the member no arm names is reported.

record X {
    pub let a : i32;

    pub self (a : i32)
        with a = a
    {}
}

record Y {
    pub let b : i32;

    pub self (b : i32)
        with b = b
    {}
}

union U
| X
| Y;

fn run (u : U) {
    match u {
        _ : X => {}
    }
}

fn main () {
    run (X (1));
}
\end{lstlisting}

\begin{lstlisting}[style=bashVerb, breaklines=true, escapechar=~]
Error[E4318] : the pattern matching over test83::U has no default case, so it must cover each of its members
 --> test_resources/type_union/test83.yr:(25,5)
25  ~\lstbox{\textSFxi{}}~     match u {
    ~\lstbox{\textSFv{}}~     ^^^^^
    ~\lstbox{\textSFxi{}}~
    = when validating pattern matching with value u of type test83::U -- (test_resources/type_union/test83.yr:(25,5))
    = when validating test83::run -- (test_resources/type_union/test83.yr:(24,4))
    ~\lstbox{\textSFii{}}~~\lstbox{\textSFx{}}~~\lstbox{\textSFx{}}~~\lstbox{\textSFx{}}~~\lstbox{\textSFx{}}~~\lstbox{\textSFx{}}~~\lstbox{\textSFx{}}~~\lstbox{\textSFx{}}~
Note : the member test83::Y is never matched
 --> test_resources/type_union/test83.yr:(25,5)
25  ~\lstbox{\textSFxi{}}~     match u {
    ~\lstbox{\textSFv{}}~     ^^^^^
    ~\lstbox{\textSFxi{}}~
    = when validating pattern matching with value u of type test83::U -- (test_resources/type_union/test83.yr:(25,5))
    = when validating test83::run -- (test_resources/type_union/test83.yr:(24,4))
    ~\lstbox{\textSFii{}}~~\lstbox{\textSFx{}}~~\lstbox{\textSFx{}}~~\lstbox{\textSFx{}}~~\lstbox{\textSFx{}}~~\lstbox{\textSFx{}}~~\lstbox{\textSFx{}}~~\lstbox{\textSFx{}}~
\end{lstlisting}

\subsubsection*{How to fix it}

Add an arm for each member or field the notes name, or a default case (\tokennolst{\_ =\textgreater{}}) when leaving them out is deliberate.

%% Generated by tools/gen_error_codes.py from docs/errors/ of the compiler;
%% edit the compiler's pages or tools/error_themes.txt, then rerun it.

\clearpage
\section{Functions, calls and closures}
\label{sec:codes:functions}

\begin{longtable}{@{}l p{0.78\linewidth} r@{}}
\toprule Code & Message & Page\\ \midrule \endhead
\hyperref[sec:codes:e4023]{E4023} & colliding function definitions \errhole{} and \errhole{} & \pageref{sec:codes:e4023}\\
\hyperref[sec:codes:e4042]{E4042} & parameter with a default value must have a referenceable name & \pageref{sec:codes:e4042}\\
\hyperref[sec:codes:e4043]{E4043} & parameter with a default value cannot be a reference & \pageref{sec:codes:e4043}\\
\hyperref[sec:codes:e4073]{E4073} & cannot create a function pointer, entities are not copiable & \pageref{sec:codes:e4073}\\
\hyperref[sec:codes:e4074]{E4074} & cannot create a function pointer, it would enclose the record by reference & \pageref{sec:codes:e4074}\\
\hyperref[sec:codes:e4075]{E4075} & using a deep copy instead of a one level copy is prohibited when enclosing a record & \pageref{sec:codes:e4075}\\
\hyperref[sec:codes:e4076]{E4076} & cannot create a function pointer from the throwing function \errhole{} & \pageref{sec:codes:e4076}\\
\hyperref[sec:codes:e4077]{E4077} & cannot create a function pointer from the unsafe function \errhole{} & \pageref{sec:codes:e4077}\\
\hyperref[sec:codes:e4078]{E4078} & cannot create an option delegate from the non throwing function \errhole{} & \pageref{sec:codes:e4078}\\
\hyperref[sec:codes:e4084]{E4084} & cannot create a function pointer from \errhole{} implicitly & \pageref{sec:codes:e4084}\\
\hyperref[sec:codes:e4085]{E4085} & cannot create a delegate from the method \errhole{} implicitly & \pageref{sec:codes:e4085}\\
\hyperref[sec:codes:e4105]{E4105} & \errhole{} is a method and must be called from an instance & \pageref{sec:codes:e4105}\\
\hyperref[sec:codes:e4107]{E4107} & call operator is not defined for value \errhole{} & \pageref{sec:codes:e4107}\\
\hyperref[sec:codes:e4113]{E4113} & main function takes at most one argument & \pageref{sec:codes:e4113}\\
\hyperref[sec:codes:e4114]{E4114} & main function cannot be inline & \pageref{sec:codes:e4114}\\
\hyperref[sec:codes:e4134]{E4134} & named parameter \errhole{} is set multiple times & \pageref{sec:codes:e4134}\\
\hyperref[sec:codes:e4171]{E4171} & no parameter is named \errhole{} & \pageref{sec:codes:e4171}\\
\hyperref[sec:codes:e4212]{E4212} & \errhole{} parameters \errhole{} expected, but \errhole{} \errhole{} provided & \pageref{sec:codes:e4212}\\
\hyperref[sec:codes:e4226]{E4226} & the call operator is not defined for \errhole{} and \errhole{} & \pageref{sec:codes:e4226}\\
\hyperref[sec:codes:e4246]{E4246} & the creation of a delegate from a method has no effect & \pageref{sec:codes:e4246}\\
\hyperref[sec:codes:e4250]{E4250} & copy the lambda closure has no effect, no values are enclosed & \pageref{sec:codes:e4250}\\
\hyperref[sec:codes:e4280]{E4280} & \errhole{} called with \errhole{} works with multiple candidates & \pageref{sec:codes:e4280}\\
\hyperref[sec:codes:e4283]{E4283} & the parameter \errhole{} has a default value, it can only be passed by name & \pageref{sec:codes:e4283}\\
\hyperref[sec:codes:e4284]{E4284} & the parameter \errhole{} can only be passed by name & \pageref{sec:codes:e4284}\\
\hyperref[sec:codes:e4285]{E4285} & parameter with a default value cannot be marked \errhole{}, it can already only be passed by name & \pageref{sec:codes:e4285}\\
\hyperref[sec:codes:e4298]{E4298} & the method \errhole{} is not a value, it can only be called & \pageref{sec:codes:e4298}\\
\bottomrule
\end{longtable}

\errorcode{E4023}{colliding function definitions \errhole{} and \errhole{}}
\label{sec:codes:e4023}

Raised during \textbf{validation}, from \texttt{ValidateErrorMessage::\allowbreak{}COLLIDING\_\allowbreak{}FUNCTION\_\allowbreak{}DEFINITION} (\textit{src/\allowbreak{}ymirc/\allowbreak{}semantic/\allowbreak{}validator/\allowbreak{}errors.yr}).

\subsubsection*{What it means}

Two functions have the same name and signatures that overload resolution cannot tell apart, so a call naming them would be ambiguous.

\subsubsection*{Example}

\textit{test\_\allowbreak{}resources/\allowbreak{}function/\allowbreak{}test15.yr}:

%% check: skip
\begin{lstlisting}[style=coloredverbatimError]
fn foo () {}
fn foo () {}

fn bar (i : i32) {i;}
fn bar () {}


fn baz (i : i32, j : i32) { i; j; }
fn baz (i : i32, ref j : i32) {
    i; j;
}
\end{lstlisting}

\begin{lstlisting}[style=bashVerb, breaklines=true, escapechar=~]
Error[E4023] : colliding function definitions test15::foo ()-> void and test15::foo ()-> void
 --> test_resources/function/test15.yr:(1,4)
 1  ~\lstbox{\textSFxi{}}~ fn foo () {}
    ~\lstbox{\textSFv{}}~    ^^^
   ...
 --> test_resources/function/test15.yr:(2,4)
 2  ~\lstbox{\textSFxi{}}~ fn foo () {}
    ~\lstbox{\textSFv{}}~    ^^^
    ~\lstbox{\textSFii{}}~~\lstbox{\textSFx{}}~~\lstbox{\textSFx{}}~~\lstbox{\textSFx{}}~~\lstbox{\textSFx{}}~~\lstbox{\textSFx{}}~~\lstbox{\textSFx{}}~~\lstbox{\textSFx{}}~
\end{lstlisting}

\subsubsection*{How to fix it}

Rename one of them, or change the parameters so the two differ in a way a call can select on.

\errorcode{E4042}{parameter with a default value must have a referenceable name}
\label{sec:codes:e4042}

Raised during \textbf{validation}, from \texttt{ValidateErrorMessage::\allowbreak{}DEFAULT\_\allowbreak{}VAR\_\allowbreak{}NO\_\allowbreak{}NAME} (\textit{src/\allowbreak{}ymirc/\allowbreak{}semantic/\allowbreak{}validator/\allowbreak{}errors.yr}).

\subsubsection*{What it means}

A parameter has a default value but no name a caller could use. Default values are only reachable by naming the parameter at the call site, so an unnamed parameter could never take its default.

\subsubsection*{Example}

\textit{test\_\allowbreak{}resources/\allowbreak{}diagnostics/\allowbreak{}test27.yr}:

%% check: skip
\begin{lstlisting}[style=coloredverbatimError]
fn foo (_: i32 = 0)-> i32 { 1 }

fn main () { foo (); }
\end{lstlisting}

\begin{lstlisting}[style=bashVerb, breaklines=true, escapechar=~]
Error[E4042] : parameter with a default value must have a referenceable name
 --> test_resources/diagnostics/test27.yr:(1,9)
 1  ~\lstbox{\textSFxi{}}~ fn foo (_: i32 = 0)-> i32 { 1 }
    ~\lstbox{\textSFv{}}~         ^        -
    ~\lstbox{\textSFxi{}}~
    = when validating test27::foo -- (test_resources/diagnostics/test27.yr:(1,4))
    ~\lstbox{\textSFii{}}~~\lstbox{\textSFx{}}~~\lstbox{\textSFx{}}~~\lstbox{\textSFx{}}~~\lstbox{\textSFx{}}~~\lstbox{\textSFx{}}~~\lstbox{\textSFx{}}~~\lstbox{\textSFx{}}~
Error[E4104] : symbol test27::foo has errors
 --> test_resources/diagnostics/test27.yr:(3,14)
 1  ~\lstbox{\textSFxi{}}~ fn foo (_: i32 = 0)-> i32 { 1 }
    ~\lstbox{\textSFv{}}~         -        -  error: parameter with a default value must have a referenceable name
 2  ~\lstbox{\textSFxi{}}~
 3  ~\lstbox{\textSFxi{}}~ fn main () { foo (); }
    ~\lstbox{\textSFv{}}~              ^^^
    ~\lstbox{\textSFxi{}}~
    = when validating test27::foo -- (test_resources/diagnostics/test27.yr:(1,4))
    = when validating test27::main -- (test_resources/diagnostics/test27.yr:(3,4))
    ~\lstbox{\textSFii{}}~~\lstbox{\textSFx{}}~~\lstbox{\textSFx{}}~~\lstbox{\textSFx{}}~~\lstbox{\textSFx{}}~~\lstbox{\textSFx{}}~~\lstbox{\textSFx{}}~~\lstbox{\textSFx{}}~
\end{lstlisting}

\subsubsection*{How to fix it}

Give the parameter a name, or remove the default value.

\errorcode{E4043}{parameter with a default value cannot be a reference}
\label{sec:codes:e4043}

Raised during \textbf{validation}, from \texttt{ValidateErrorMessage::\allowbreak{}DEFAULT\_\allowbreak{}VAR\_\allowbreak{}REF} (\textit{src/\allowbreak{}ymirc/\allowbreak{}semantic/\allowbreak{}validator/\allowbreak{}errors.yr}).

\subsubsection*{What it means}

A parameter is declared \tokennolst{ref} and given a default value. A reference parameter borrows the caller's storage; a default value has no storage in the caller to borrow.

\subsubsection*{Example}

\textit{test\_\allowbreak{}resources/\allowbreak{}diagnostics/\allowbreak{}test28.yr}:

%% check: skip
\begin{lstlisting}[style=coloredverbatimError]
fn foo (ref x: i32 = 0)-> i32 { x }

fn main () { foo (); }
\end{lstlisting}

\begin{lstlisting}[style=bashVerb, breaklines=true, escapechar=~]
Error[E4043] : parameter with a default value cannot be a reference
 --> test_resources/diagnostics/test28.yr:(1,9)
 1  ~\lstbox{\textSFxi{}}~ fn foo (ref x: i32 = 0)-> i32 { x }
    ~\lstbox{\textSFv{}}~         ^^^ -        -
    ~\lstbox{\textSFxi{}}~
    = when validating test28::foo -- (test_resources/diagnostics/test28.yr:(1,4))
    ~\lstbox{\textSFii{}}~~\lstbox{\textSFx{}}~~\lstbox{\textSFx{}}~~\lstbox{\textSFx{}}~~\lstbox{\textSFx{}}~~\lstbox{\textSFx{}}~~\lstbox{\textSFx{}}~~\lstbox{\textSFx{}}~
Error[E4104] : symbol test28::foo has errors
 --> test_resources/diagnostics/test28.yr:(3,14)
 1  ~\lstbox{\textSFxi{}}~ fn foo (ref x: i32 = 0)-> i32 { x }
    ~\lstbox{\textSFv{}}~         ---          -  error: parameter with a default value cannot be a reference
 2  ~\lstbox{\textSFxi{}}~
 3  ~\lstbox{\textSFxi{}}~ fn main () { foo (); }
    ~\lstbox{\textSFv{}}~              ^^^
    ~\lstbox{\textSFxi{}}~
    = when validating test28::foo -- (test_resources/diagnostics/test28.yr:(1,4))
    = when validating test28::main -- (test_resources/diagnostics/test28.yr:(3,4))
    ~\lstbox{\textSFii{}}~~\lstbox{\textSFx{}}~~\lstbox{\textSFx{}}~~\lstbox{\textSFx{}}~~\lstbox{\textSFx{}}~~\lstbox{\textSFx{}}~~\lstbox{\textSFx{}}~~\lstbox{\textSFx{}}~
\end{lstlisting}

\subsubsection*{How to fix it}

Drop either the \tokennolst{ref} or the default value.

\errorcode{E4073}{cannot create a function pointer, entities are not copiable}
\label{sec:codes:e4073}

Raised during \textbf{validation}, from \texttt{ValidateErrorMessage::\allowbreak{}FUNCTION\_\allowbreak{}TO\_\allowbreak{}FPTR\_\allowbreak{}ENTITY} (\textit{src/\allowbreak{}ymirc/\allowbreak{}semantic/\allowbreak{}validator/\allowbreak{}errors.yr}).

\subsubsection*{What it means}

A function pointer would have to capture an entity, and entities are not copiable, so there is no way to store one in the closure environment.

\subsubsection*{Example}

\textit{test\_\allowbreak{}resources/\allowbreak{}diagnostics/\allowbreak{}test103.yr}:

%% check: skip
\begin{lstlisting}[style=coloredverbatimError]
// a function pointer made from the method of an entity
entity Counter {
    pub self () {}
    pub fn get (self)-> i32 { 1 }
}

fn main () {
    let c = Counter ();
    let _ = &c.get;
}
\end{lstlisting}

\begin{lstlisting}[style=bashVerb, breaklines=true, escapechar=~]
Error[E4073] : cannot create a function pointer, entities are not copiable
 --> test_resources/diagnostics/test103.yr:(9,13)
 9  ~\lstbox{\textSFxi{}}~     let _ = &c.get;
    ~\lstbox{\textSFv{}}~             ^
    ~\lstbox{\textSFxi{}}~
    = when validating test103::main -- (test_resources/diagnostics/test103.yr:(7,4))
    ~\lstbox{\textSFii{}}~~\lstbox{\textSFx{}}~~\lstbox{\textSFx{}}~~\lstbox{\textSFx{}}~~\lstbox{\textSFx{}}~~\lstbox{\textSFx{}}~~\lstbox{\textSFx{}}~~\lstbox{\textSFx{}}~
\end{lstlisting}

\subsubsection*{How to fix it}

Pass the entity as a parameter of the function rather than capturing it.

\errorcode{E4074}{cannot create a function pointer, it would enclose the record by reference}
\label{sec:codes:e4074}

Raised during \textbf{validation}, from \texttt{ValidateErrorMessage::\allowbreak{}FUNCTION\_\allowbreak{}TO\_\allowbreak{}FPTR\_\allowbreak{}RECORD} (\textit{src/\allowbreak{}ymirc/\allowbreak{}semantic/\allowbreak{}validator/\allowbreak{}errors.yr}).

\subsubsection*{What it means}

A function pointer would capture a record by reference, which would outlive the record.

\subsubsection*{Example}

\textit{test\_\allowbreak{}resources/\allowbreak{}diagnostics/\allowbreak{}test104.yr}:

%% check: skip
\begin{lstlisting}[style=coloredverbatimError]
// a function pointer made from the method of a record, outside a copy
record Counter {
    pub self () {}
    pub fn get (self)-> i32 { 1 }
}

fn main () {
    let c = Counter ();
    let _ = &c.get;
}
\end{lstlisting}

\begin{lstlisting}[style=bashVerb, breaklines=true, escapechar=~]
Error[E4074] : cannot create a function pointer, it would enclose the record by reference
 --> test_resources/diagnostics/test104.yr:(9,13)
 9  ~\lstbox{\textSFxi{}}~     let _ = &c.get;
    ~\lstbox{\textSFv{}}~             ^  note: did you mean to enclose it within a copy?
    ~\lstbox{\textSFxi{}}~
    = when validating test104::main -- (test_resources/diagnostics/test104.yr:(7,4))
    ~\lstbox{\textSFii{}}~~\lstbox{\textSFx{}}~~\lstbox{\textSFx{}}~~\lstbox{\textSFx{}}~~\lstbox{\textSFx{}}~~\lstbox{\textSFx{}}~~\lstbox{\textSFx{}}~~\lstbox{\textSFx{}}~
\end{lstlisting}

\subsubsection*{How to fix it}

Capture a copy of the record, or pass it as a parameter.

\errorcode{E4075}{using a deep copy instead of a one level copy is prohibited when enclosing a record}
\label{sec:codes:e4075}

Raised during \textbf{validation}, from \texttt{ValidateErrorMessage::\allowbreak{}FUNCTION\_\allowbreak{}TO\_\allowbreak{}FPTR\_\allowbreak{}RECORD\_\allowbreak{}DCOPY} (\textit{src/\allowbreak{}ymirc/\allowbreak{}semantic/\allowbreak{}validator/\allowbreak{}errors.yr}).

\subsubsection*{What it means}

A record was captured with \tokennolst{dcopy} where the one level \tokennolst{copy} is required, cf. \hyperref[sec:codes:retired]{E4033}.

\subsubsection*{Example}

\textit{test\_\allowbreak{}resources/\allowbreak{}diagnostics/\allowbreak{}test105.yr}:

%% check: skip
\begin{lstlisting}[style=coloredverbatimError]
// a function pointer made from the method of a record, inside a deep copy
record Counter {
    pub self () {}
    pub fn get (self)-> i32 { 1 }
}

fn main () {
    let c = Counter ();
    let _ = dcopy &c.get;
}
\end{lstlisting}

\begin{lstlisting}[style=bashVerb, breaklines=true, escapechar=~]
Error[E4075] : using a deep copy instead of a one level copy is prohibited when enclosing a record
 --> test_resources/diagnostics/test105.yr:(9,19)
 9  ~\lstbox{\textSFxi{}}~     let _ = dcopy &c.get;
    ~\lstbox{\textSFv{}}~                   ^
    ~\lstbox{\textSFxi{}}~
    = help: replace dcopy by the simple copy operator
    ~\lstbox{\textSFxi{}}~
 9  -     let _ = dcopy &c.get;
 9  +     let _ = copy &c.get;
    ~\lstbox{\textSFxi{}}~
    = when validating test105::main -- (test_resources/diagnostics/test105.yr:(7,4))
    ~\lstbox{\textSFii{}}~~\lstbox{\textSFx{}}~~\lstbox{\textSFx{}}~~\lstbox{\textSFx{}}~~\lstbox{\textSFx{}}~~\lstbox{\textSFx{}}~~\lstbox{\textSFx{}}~~\lstbox{\textSFx{}}~
\end{lstlisting}

\subsubsection*{How to fix it}

Use \tokennolst{copy}.

\errorcode{E4076}{cannot create a function pointer from the throwing function \errhole{}}
\label{sec:codes:e4076}

Raised during \textbf{validation}, from \texttt{ValidateErrorMessage::\allowbreak{}FUNCTION\_\allowbreak{}TO\_\allowbreak{}FPTR\_\allowbreak{}THROWING} (\textit{src/\allowbreak{}ymirc/\allowbreak{}semantic/\allowbreak{}validator/\allowbreak{}errors.yr}).

\subsubsection*{What it means}

A function pointer was created from a function that can throw. A function pointer has no place to declare what it throws, so a caller could not know it has to handle an exception.

\subsubsection*{Example}

\textit{test\_\allowbreak{}resources/\allowbreak{}function/\allowbreak{}test11.yr}:

%% check: skip
\begin{lstlisting}[style=coloredverbatimError]
fn foo ()
    throws AssertError
{
    assert (false);
}

fn main () {
    let a = &foo;
    a ();
}
\end{lstlisting}

\begin{lstlisting}[style=bashVerb, breaklines=true, escapechar=~]
Error[E4076] : cannot create a function pointer from the throwing function test11::foo ()-> void
 --> test_resources/function/test11.yr:(8,13)
 8  ~\lstbox{\textSFxi{}}~     let a = &foo;
    ~\lstbox{\textSFv{}}~             ^
    ~\lstbox{\textSFxi{}}~ Note : throwing core::exception::assertion::AssertError
    ~\lstbox{\textSFxi{}}~  --> test_resources/function/test11.yr:(2,12)
    ~\lstbox{\textSFxi{}}~  2  ~\lstbox{\textSFxi{}}~     throws AssertError
    ~\lstbox{\textSFxi{}}~     ~\lstbox{\textSFv{}}~            ^^^^^^^^^^^
    ~\lstbox{\textSFxi{}}~     ~\lstbox{\textSFii{}}~~\lstbox{\textSFx{}}~~\lstbox{\textSFx{}}~~\lstbox{\textSFx{}}~~\lstbox{\textSFx{}}~~\lstbox{\textSFx{}}~~\lstbox{\textSFx{}}~~\lstbox{\textSFx{}}~
    ~\lstbox{\textSFii{}}~~\lstbox{\textSFx{}}~~\lstbox{\textSFx{}}~~\lstbox{\textSFx{}}~~\lstbox{\textSFx{}}~~\lstbox{\textSFx{}}~~\lstbox{\textSFvii{}}~~\lstbox{\textSFx{}}~
\end{lstlisting}

\subsubsection*{How to fix it}

Wrap the call in a lambda that catches the exception, or use a delegate that carries the throw information.

\errorcode{E4077}{cannot create a function pointer from the unsafe function \errhole{}}
\label{sec:codes:e4077}

Raised during \textbf{validation}, from \texttt{ValidateErrorMessage::\allowbreak{}FUNCTION\_\allowbreak{}TO\_\allowbreak{}FPTR\_\allowbreak{}UNSAFE} (\textit{src/\allowbreak{}ymirc/\allowbreak{}semantic/\allowbreak{}validator/\allowbreak{}errors.yr}).

\subsubsection*{What it means}

A function pointer was created from an unsafe function. Calling through the pointer would leave no trace that the call is unsafe.

\subsubsection*{Example}

\textit{test\_\allowbreak{}resources/\allowbreak{}function/\allowbreak{}test13.yr}:

%% check: skip
\begin{lstlisting}[style=coloredverbatimError]
@unsafe
fn foo ();

fn main () {
    let a = &foo;
    a ();
}
\end{lstlisting}

\begin{lstlisting}[style=bashVerb, breaklines=true, escapechar=~]
Error[E4077] : cannot create a function pointer from the unsafe function test13::foo ()-> void
 --> test_resources/function/test13.yr:(5,13)
 5  ~\lstbox{\textSFxi{}}~     let a = &foo;
    ~\lstbox{\textSFv{}}~             ^
    ~\lstbox{\textSFii{}}~~\lstbox{\textSFx{}}~~\lstbox{\textSFx{}}~~\lstbox{\textSFx{}}~~\lstbox{\textSFx{}}~~\lstbox{\textSFx{}}~~\lstbox{\textSFx{}}~~\lstbox{\textSFx{}}~
\end{lstlisting}

\subsubsection*{How to fix it}

Wrap the call in a safe function that enters an \tokennolst{unsafe} block itself, and take a pointer to that.

\errorcode{E4078}{cannot create an option delegate from the non throwing function \errhole{}}
\label{sec:codes:e4078}

Raised during \textbf{validation}, from \texttt{ValidateErrorMessage::\allowbreak{}FUNCTION\_\allowbreak{}TO\_\allowbreak{}OPTION\_\allowbreak{}NO\_\allowbreak{}THROW} (\textit{src/\allowbreak{}ymirc/\allowbreak{}semantic/\allowbreak{}validator/\allowbreak{}errors.yr}).

\subsubsection*{What it means}

The \tokennolst{\&?} operator, which wraps a throwing function into one returning an option, was applied to a function that does not throw. There is no error for the option to carry.

\subsubsection*{Example}

\textit{test\_\allowbreak{}resources/\allowbreak{}function/\allowbreak{}test19.yr}:

%% check: skip
\begin{lstlisting}[style=coloredverbatimError]
fn foo (i : i32)-> i32 {
    i + 1
}

fn main () {
    let _f_ = &?foo; // nothing to catch, '&' is what creates a pointer on 'foo'
}
\end{lstlisting}

\begin{lstlisting}[style=bashVerb, breaklines=true, escapechar=~]
Error[E4078] : cannot create an option delegate from the non throwing function test19::foo (i: i32)-> i32
 --> test_resources/function/test19.yr:(6,15)
 6  ~\lstbox{\textSFxi{}}~     let _f_ = &?foo; // nothing to catch, '&' is what creates a pointer on 'foo'
    ~\lstbox{\textSFv{}}~               ^^
    ~\lstbox{\textSFxi{}}~
    = help: replace the operator &? by the simple & operator
    ~\lstbox{\textSFxi{}}~
 6  -     let _f_ = &?foo; // nothing to catch, '&' is what creates a pointer on 'foo'
 6  +     let _f_ = &foo; // nothing to catch, '&' is what creates a pointer on 'foo'
    ~\lstbox{\textSFii{}}~~\lstbox{\textSFx{}}~~\lstbox{\textSFx{}}~~\lstbox{\textSFx{}}~~\lstbox{\textSFx{}}~~\lstbox{\textSFx{}}~~\lstbox{\textSFx{}}~~\lstbox{\textSFx{}}~
\end{lstlisting}

\subsubsection*{How to fix it}

Use a plain \tokennolst{\&} to take the address of a non-throwing function.

\errorcode{E4084}{cannot create a function pointer from \errhole{} implicitly}
\label{sec:codes:e4084}

Raised during \textbf{validation}, from \texttt{ValidateErrorMessage::\allowbreak{}IMPLICIT\_\allowbreak{}ADDRESS\_\allowbreak{}FUNCTION} (\textit{src/\allowbreak{}ymirc/\allowbreak{}semantic/\allowbreak{}validator/\allowbreak{}errors.yr}).

\subsubsection*{What it means}

A function was used where a function pointer is expected, without asking for its address. Ymir does not take the address of a function implicitly, so a function name used as a value is a mistake rather than a conversion.

\subsubsection*{Example}

\textit{test\_\allowbreak{}resources/\allowbreak{}diagnostics/\allowbreak{}test39.yr}:

%% check: skip
\begin{lstlisting}[style=coloredverbatimError]
fn foo ()-> i32 { 0 }

fn main () {
    let _: fn ()-> i32 = foo;
}
\end{lstlisting}

\begin{lstlisting}[style=bashVerb, breaklines=true, escapechar=~]
Error[E4084] : cannot create a function pointer from test39::foo ()-> i32 implicitly
 --> test_resources/diagnostics/test39.yr:(1,4)
 1  ~\lstbox{\textSFxi{}}~ fn foo ()-> i32 { 0 }
    ~\lstbox{\textSFv{}}~    ^^^
   ...
 --> test_resources/diagnostics/test39.yr:(4,26)
 4  ~\lstbox{\textSFxi{}}~     let _: fn ()-> i32 = foo;
    ~\lstbox{\textSFv{}}~                          ---  note: did you forget ()?
    ~\lstbox{\textSFxi{}}~
    = help: or add the missing &
    ~\lstbox{\textSFxi{}}~
 4  -     let _: fn ()-> i32 = foo;
 4  +     let _: fn ()-> i32 = &foo;
    ~\lstbox{\textSFii{}}~~\lstbox{\textSFx{}}~~\lstbox{\textSFx{}}~~\lstbox{\textSFx{}}~~\lstbox{\textSFx{}}~~\lstbox{\textSFx{}}~~\lstbox{\textSFx{}}~~\lstbox{\textSFx{}}~
\end{lstlisting}

\subsubsection*{How to fix it}

Write \tokennolst{\&} before the function name.

\errorcode{E4085}{cannot create a delegate from the method \errhole{} implicitly}
\label{sec:codes:e4085}

Raised during \textbf{validation}, from \texttt{ValidateErrorMessage::\allowbreak{}IMPLICIT\_\allowbreak{}ADDRESS\_\allowbreak{}METHOD} (\textit{src/\allowbreak{}ymirc/\allowbreak{}semantic/\allowbreak{}validator/\allowbreak{}errors.yr}).

\subsubsection*{What it means}

A method was used where a delegate is expected, without asking for its address.

\subsubsection*{Example}

\textit{test\_\allowbreak{}resources/\allowbreak{}diagnostics/\allowbreak{}test40.yr}:

%% check: skip
\begin{lstlisting}[style=coloredverbatimError]
class Foo {
    pub self () {}
    pub fn bar (self)-> i32 { 0 }
}

fn main () {
    let f = copy Foo ();
    let _: dg ()-> i32 = f.bar;
}
\end{lstlisting}

\begin{lstlisting}[style=bashVerb, breaklines=true, escapechar=~]
Error[E4085] : cannot create a delegate from the method f.test40::Foo::bar ()-> i32 implicitly
 --> test_resources/diagnostics/test40.yr:(8,27)
 8  ~\lstbox{\textSFxi{}}~     let _: dg ()-> i32 = f.bar;
    ~\lstbox{\textSFv{}}~                           ^  note: did you forget ()?
    ~\lstbox{\textSFxi{}}~
    = help: or add the missing &
    ~\lstbox{\textSFxi{}}~
 8  -     let _: dg ()-> i32 = f.bar;
 8  +     let _: dg ()-> i32 = &f.bar;
    ~\lstbox{\textSFii{}}~~\lstbox{\textSFx{}}~~\lstbox{\textSFx{}}~~\lstbox{\textSFx{}}~~\lstbox{\textSFx{}}~~\lstbox{\textSFx{}}~~\lstbox{\textSFx{}}~~\lstbox{\textSFx{}}~
\end{lstlisting}

\subsubsection*{How to fix it}

Write \tokennolst{\&} before the method access.

\errorcode{E4105}{\errhole{} is a method and must be called from an instance}
\label{sec:codes:e4105}

Raised during \textbf{validation}, from \texttt{ValidateErrorMessage::\allowbreak{}IS\_\allowbreak{}A\_\allowbreak{}METHOD} (\textit{src/\allowbreak{}ymirc/\allowbreak{}semantic/\allowbreak{}validator/\allowbreak{}errors.yr}).

\subsubsection*{What it means}

A method was used as if it were a free function. A method needs the instance it runs on.

\subsubsection*{Example}

\textit{test\_\allowbreak{}resources/\allowbreak{}lit\_\allowbreak{}class/\allowbreak{}errors/\allowbreak{}test1.yr}:

%% check: skip
\begin{lstlisting}[style=coloredverbatimError]
class A {
    pub self () {}

    pub fn bar (self) {}
    pub fn foo (self) {
        bar ();
    }
}
\end{lstlisting}

\begin{lstlisting}[style=bashVerb, breaklines=true, escapechar=~]
Error[E4105] : test1::A::bar is a method and must be called from an instance
 --> test_resources/lit_class/errors/test1.yr:(8,9)
 6  ~\lstbox{\textSFxi{}}~     pub fn bar (self) {}
    ~\lstbox{\textSFv{}}~            ---  note: defined here
 7  ~\lstbox{\textSFxi{}}~     pub fn foo (self) {
 8  ~\lstbox{\textSFxi{}}~         bar ();
    ~\lstbox{\textSFv{}}~         ^^^
    ~\lstbox{\textSFxi{}}~
    = when validating test1::A::foo -- (test_resources/lit_class/errors/test1.yr:(7,12))
    = when validating -- (test_resources/lit_class/errors/test1.yr:(3,7))
    ~\lstbox{\textSFii{}}~~\lstbox{\textSFx{}}~~\lstbox{\textSFx{}}~~\lstbox{\textSFx{}}~~\lstbox{\textSFx{}}~~\lstbox{\textSFx{}}~~\lstbox{\textSFx{}}~~\lstbox{\textSFx{}}~
\end{lstlisting}

\subsubsection*{How to fix it}

Call it on an instance (\tokennolst{x.m ()}), or take a delegate with \tokennolst{\&x.m} if the call is to happen later.

\errorcode{E4107}{call operator is not defined for value \errhole{}}
\label{sec:codes:e4107}

Raised during \textbf{validation}, from \texttt{ValidateErrorMessage::\allowbreak{}IS\_\allowbreak{}NOT\_\allowbreak{}CALLABLE} (\textit{src/\allowbreak{}ymirc/\allowbreak{}semantic/\allowbreak{}validator/\allowbreak{}errors.yr}).

\subsubsection*{What it means}

A value was called and its type defines no call operator.

\subsubsection*{Example}

\textit{test\_\allowbreak{}resources/\allowbreak{}diagnostics/\allowbreak{}test48.yr}:

%% check: skip
\begin{lstlisting}[style=coloredverbatimError]
fn main () {
    let x = 1;
    x ();
}
\end{lstlisting}

\begin{lstlisting}[style=bashVerb, breaklines=true, escapechar=~]
Error[E4107] : call operator is not defined for value 1
 --> test_resources/diagnostics/test48.yr:(2,13)
 2  ~\lstbox{\textSFxi{}}~     let x = 1;
    ~\lstbox{\textSFv{}}~             ^
    ~\lstbox{\textSFii{}}~~\lstbox{\textSFx{}}~~\lstbox{\textSFx{}}~~\lstbox{\textSFx{}}~~\lstbox{\textSFx{}}~~\lstbox{\textSFx{}}~~\lstbox{\textSFx{}}~~\lstbox{\textSFx{}}~
\end{lstlisting}

\subsubsection*{How to fix it}

Check that the value is the function, the function pointer or the delegate intended. A class can be made callable by defining the call operator on it.

\errorcode{E4113}{main function takes at most one argument}
\label{sec:codes:e4113}

Raised during \textbf{validation}, from \texttt{ValidateErrorMessage::\allowbreak{}MAIN\_\allowbreak{}FUNCTION\_\allowbreak{}ONE\_\allowbreak{}ARG} (\textit{src/\allowbreak{}ymirc/\allowbreak{}semantic/\allowbreak{}validator/\allowbreak{}errors.yr}).

\subsubsection*{What it means}

The \tokennolst{main} function declares more than one parameter. It takes either nothing or the command line arguments.

\subsubsection*{Example}

\textit{test\_\allowbreak{}resources/\allowbreak{}diagnostics/\allowbreak{}test141.yr}:

%% check: skip
\begin{lstlisting}[style=coloredverbatimError]
// a main function taking a second parameter
fn main (args : [[c8]], count : i32) {
    let _ = (args, count);
}
\end{lstlisting}

\begin{lstlisting}[style=bashVerb, breaklines=true, escapechar=~]
Error[E4113] : main function takes at most one argument
 --> test_resources/diagnostics/test141.yr:(2,25)
 2  ~\lstbox{\textSFxi{}}~ fn main (args : [[c8]], count : i32) {
    ~\lstbox{\textSFv{}}~                         ^^^^^
    ~\lstbox{\textSFxi{}}~
    = when validating test141::main -- (test_resources/diagnostics/test141.yr:(2,4))
    ~\lstbox{\textSFii{}}~~\lstbox{\textSFx{}}~~\lstbox{\textSFx{}}~~\lstbox{\textSFx{}}~~\lstbox{\textSFx{}}~~\lstbox{\textSFx{}}~~\lstbox{\textSFx{}}~~\lstbox{\textSFx{}}~
\end{lstlisting}

\subsubsection*{How to fix it}

Reduce it to at most one parameter.

\errorcode{E4114}{main function cannot be inline}
\label{sec:codes:e4114}

Raised during \textbf{validation}, from \texttt{ValidateErrorMessage::\allowbreak{}MAIN\_\allowbreak{}INLINE} (\textit{src/\allowbreak{}ymirc/\allowbreak{}semantic/\allowbreak{}validator/\allowbreak{}errors.yr}).

\subsubsection*{What it means}

The \tokennolst{main} function is marked \tokennolst{inline}. It is the entry point the runtime calls, so it has to exist as a symbol.

\subsubsection*{Example}

\textit{test\_\allowbreak{}resources/\allowbreak{}diagnostics/\allowbreak{}test49.yr}:

%% check: skip
\begin{lstlisting}[style=coloredverbatimError]
@inline
fn main () {}
\end{lstlisting}

\begin{lstlisting}[style=bashVerb, breaklines=true, escapechar=~]
Error[E4114] : main function cannot be inline
 --> test_resources/diagnostics/test49.yr:(1,2)
 1  ~\lstbox{\textSFxi{}}~ @inline
    ~\lstbox{\textSFv{}}~  ^^^^^^
    ~\lstbox{\textSFii{}}~~\lstbox{\textSFx{}}~~\lstbox{\textSFx{}}~~\lstbox{\textSFx{}}~~\lstbox{\textSFx{}}~~\lstbox{\textSFx{}}~~\lstbox{\textSFx{}}~~\lstbox{\textSFx{}}~
\end{lstlisting}

\subsubsection*{How to fix it}

Remove \tokennolst{inline}.

\errorcode{E4134}{named parameter \errhole{} is set multiple times}
\label{sec:codes:e4134}

Raised during \textbf{validation}, from \texttt{ValidateErrorMessage::\allowbreak{}MULTIPLE\_\allowbreak{}NAMED\_\allowbreak{}PARAM} (\textit{src/\allowbreak{}ymirc/\allowbreak{}semantic/\allowbreak{}validator/\allowbreak{}errors.yr}).

\subsubsection*{What it means}

The same parameter was given twice by name at a call site.

\subsubsection*{Example}

\textit{test\_\allowbreak{}resources/\allowbreak{}diagnostics/\allowbreak{}test53.yr}:

%% check: skip
\begin{lstlisting}[style=coloredverbatimError]
fn foo (a: i32 = 0, b: i32 = 0)-> i32 { a + b }

fn main () { foo (a-> 1, a-> 2); }
\end{lstlisting}

\begin{lstlisting}[style=bashVerb, breaklines=true, escapechar=~]
Error[E4134] : named parameter a is set multiple times
 --> test_resources/diagnostics/test53.yr:(3,23)
 3  ~\lstbox{\textSFxi{}}~ fn main () { foo (a-> 1, a-> 2); }
    ~\lstbox{\textSFv{}}~                       ^   ^^
    ~\lstbox{\textSFii{}}~~\lstbox{\textSFx{}}~~\lstbox{\textSFx{}}~~\lstbox{\textSFx{}}~~\lstbox{\textSFx{}}~~\lstbox{\textSFx{}}~~\lstbox{\textSFx{}}~~\lstbox{\textSFx{}}~
\end{lstlisting}

\subsubsection*{How to fix it}

Remove the duplicate. A parameter given both positionally and by name also produces this.

\errorcode{E4171}{no parameter is named \errhole{}}
\label{sec:codes:e4171}

Raised during \textbf{validation}, from \texttt{ValidateErrorMessage::\allowbreak{}NO\_\allowbreak{}PARAMETER\_\allowbreak{}NAMED} (\textit{src/\allowbreak{}ymirc/\allowbreak{}semantic/\allowbreak{}validator/\allowbreak{}errors.yr}).

\subsubsection*{What it means}

A named argument at a call site names a parameter the callee does not have.

\subsubsection*{Example}

\textit{test\_\allowbreak{}resources/\allowbreak{}default\_\allowbreak{}ctor/\allowbreak{}test3.yr}:

%% check: skip
\begin{lstlisting}[style=coloredverbatimError]
// A field that is not public is not nameable at the call site, it is not part of
// the signature of the generated ctor.
class Foo {
   pub {
     let x: i32;
     let mut y: i32;
   }

   let mut z: i32 = 98;

   pub self;
}

fn main () {
    let a = copy Foo (y-> 2, 1, z-> 3);
    a.x;
}
\end{lstlisting}

\begin{lstlisting}[style=bashVerb, breaklines=true, escapechar=~]
Error[E4171] : no parameter is named z
 --> test_resources/default_ctor/test3.yr:(15,34)
15  ~\lstbox{\textSFxi{}}~     let a = copy Foo (y-> 2, 1, z-> 3);
    ~\lstbox{\textSFv{}}~                                  ^^
    ~\lstbox{\textSFii{}}~~\lstbox{\textSFx{}}~~\lstbox{\textSFx{}}~~\lstbox{\textSFx{}}~~\lstbox{\textSFx{}}~~\lstbox{\textSFx{}}~~\lstbox{\textSFx{}}~~\lstbox{\textSFx{}}~
\end{lstlisting}

\subsubsection*{How to fix it}

Check the spelling against the prototype. Note that a parameter with a default value can only be passed by name, so an unnamed parameter cannot be reached this way.

\errorcode{E4212}{\errhole{} parameters \errhole{} expected, but \errhole{} \errhole{} provided}
\label{sec:codes:e4212}

Raised during \textbf{validation}, from \texttt{ValidateErrorMessage::\allowbreak{}TOO\_\allowbreak{}FEW\_\allowbreak{}PARAMETERS} (\textit{src/\allowbreak{}ymirc/\allowbreak{}semantic/\allowbreak{}validator/\allowbreak{}errors.yr}).

\subsubsection*{What it means}

A call provides a different number of arguments than the callee takes. Both counts are in the message.

A parameter that declares a default value is filled by name or left to its default, so it is not part of the count a positional call is compared against -- \tokennolst{foo (x: i32, y: i32 = 8)} expects one argument, not two. When the extra arguments are explained by such a parameter, \hyperref[sec:codes:e4283]{E4283} hangs under this error and names it.

\subsubsection*{Example}

\textit{test\_\allowbreak{}resources/\allowbreak{}data\_\allowbreak{}record/\allowbreak{}ctors/\allowbreak{}test2.yr}:

%% check: skip
\begin{lstlisting}[style=coloredverbatimError]
@data
record Foo {
    let x: i32;
}

fn main() {
    let a = Foo();
    let b = Foo(x-> 34);
    a.x;
    b.x;
}
\end{lstlisting}

\begin{lstlisting}[style=bashVerb, breaklines=true, escapechar=~]
Error[E4212] : 1 parameters is expected, but 0 are provided
 --> test_resources/data_record/ctors/test2.yr:(7,16)
 7  ~\lstbox{\textSFxi{}}~     let a = Foo();
    ~\lstbox{\textSFv{}}~                ^
    ~\lstbox{\textSFii{}}~~\lstbox{\textSFx{}}~~\lstbox{\textSFx{}}~~\lstbox{\textSFx{}}~~\lstbox{\textSFx{}}~~\lstbox{\textSFx{}}~~\lstbox{\textSFx{}}~~\lstbox{\textSFx{}}~
\end{lstlisting}

\subsubsection*{How to fix it}

Pass the missing arguments. If the count looks right, check for a parameter with a default value: it is reachable by name only, so it is not one of the arguments a positional call is expected to provide.

\errorcode{E4226}{the call operator is not defined for \errhole{} and \errhole{}}
\label{sec:codes:e4226}

Raised during \textbf{validation}, from \texttt{ValidateErrorMessage::\allowbreak{}UNDEFINED\_\allowbreak{}CALL\_\allowbreak{}OP} (\textit{src/\allowbreak{}ymirc/\allowbreak{}semantic/\allowbreak{}validator/\allowbreak{}errors.yr}).

\subsubsection*{What it means}

A value was called with arguments no candidate accepts. The value and the arguments are in the message, and the candidates are listed under it.

\subsubsection*{Example}

\textit{test\_\allowbreak{}resources/\allowbreak{}diagnostics/\allowbreak{}test48.yr}:

%% check: skip
\begin{lstlisting}[style=coloredverbatimError]
fn main () {
    let x = 1;
    x ();
}
\end{lstlisting}

\begin{lstlisting}[style=bashVerb, breaklines=true, escapechar=~]
Error[E4226] : the call operator is not defined for x: i32 and {}
 --> test_resources/diagnostics/test48.yr:(3,7)
 3  ~\lstbox{\textSFxi{}}~     x ();
    ~\lstbox{\textSFv{}}~       ^
    ~\lstbox{\textSFxi{}}~ Error[E4107] : call operator is not defined for value 1
    ~\lstbox{\textSFxi{}}~  --> test_resources/diagnostics/test48.yr:(2,13)
    ~\lstbox{\textSFxi{}}~  2  ~\lstbox{\textSFxi{}}~     let x = 1;
    ~\lstbox{\textSFxi{}}~     ~\lstbox{\textSFv{}}~             ^
    ~\lstbox{\textSFxi{}}~     ~\lstbox{\textSFii{}}~~\lstbox{\textSFx{}}~~\lstbox{\textSFx{}}~~\lstbox{\textSFx{}}~~\lstbox{\textSFx{}}~~\lstbox{\textSFx{}}~~\lstbox{\textSFx{}}~~\lstbox{\textSFx{}}~
    ~\lstbox{\textSFii{}}~~\lstbox{\textSFx{}}~~\lstbox{\textSFx{}}~~\lstbox{\textSFx{}}~~\lstbox{\textSFx{}}~~\lstbox{\textSFx{}}~~\lstbox{\textSFvii{}}~~\lstbox{\textSFx{}}~
\end{lstlisting}

\subsubsection*{Notes}

A candidate is refused with a note under this diagnostic. Besides the notes of \hyperref[sec:codes:e4236]{E4236} for a template candidate, one carries no code of its own:

\begin{itemize}
\item \texttt{ref of type \textless{}\textgreater{} must be const, but is mutable} (formerly E4086): a mutable reference was passed to a \tokennolst{ref} parameter that is not \tokennolst{mut}.
\end{itemize}

\subsubsection*{How to fix it}

Match one of the listed candidates, cf. \hyperref[sec:codes:e4107]{E4107}.

\errorcode{E4246}{the creation of a delegate from a method has no effect}
\label{sec:codes:e4246}

Raised during \textbf{validation}, from \texttt{ValidateErrorMessage::\allowbreak{}UNECESSARY\_\allowbreak{}ADDRESS\_\allowbreak{}METHOD} (\textit{src/\allowbreak{}ymirc/\allowbreak{}semantic/\allowbreak{}validator/\allowbreak{}errors.yr}).

\subsubsection*{What it means}

A delegate was taken of a method in a place where the delegate is used immediately, so building it changes nothing.

\subsubsection*{Example}

\textit{test\_\allowbreak{}resources/\allowbreak{}diagnostics/\allowbreak{}test78.yr}:

%% check: skip
\begin{lstlisting}[style=coloredverbatimError]
class Foo {
    pub self () {}
    pub fn bar (self) {}
}

fn main () {
    let f = copy Foo ();
    (&f.bar) ();
}
\end{lstlisting}

\begin{lstlisting}[style=bashVerb, breaklines=true, escapechar=~]
Error[E4246] : the creation of a delegate from a method has no effect
 --> test_resources/diagnostics/test78.yr:(8,6)
 8  ~\lstbox{\textSFxi{}}~     (&f.bar) ();
    ~\lstbox{\textSFv{}}~      ^
    ~\lstbox{\textSFii{}}~~\lstbox{\textSFx{}}~~\lstbox{\textSFx{}}~~\lstbox{\textSFx{}}~~\lstbox{\textSFx{}}~~\lstbox{\textSFx{}}~~\lstbox{\textSFx{}}~~\lstbox{\textSFx{}}~
\end{lstlisting}

\subsubsection*{How to fix it}

Call the method directly.

\errorcode{E4250}{copy the lambda closure has no effect, no values are enclosed}
\label{sec:codes:e4250}

Raised during \textbf{validation}, from \texttt{ValidateErrorMessage::\allowbreak{}UNECESSARY\_\allowbreak{}LMBD\_\allowbreak{}COPY} (\textit{src/\allowbreak{}ymirc/\allowbreak{}semantic/\allowbreak{}validator/\allowbreak{}errors.yr}).

\subsubsection*{What it means}

\tokennolst{copy} was applied to a lambda that captures nothing, so there is no environment to duplicate.

\subsubsection*{Example}

\textit{test\_\allowbreak{}resources/\allowbreak{}diagnostics/\allowbreak{}test79.yr}:

%% check: skip
\begin{lstlisting}[style=coloredverbatimError]
fn main () {
    let f = copy (|| => { 1 });
    f ();
}
\end{lstlisting}

\begin{lstlisting}[style=bashVerb, breaklines=true, escapechar=~]
Error[E4250] : copy the lambda closure has no effect, no values are enclosed
 --> test_resources/diagnostics/test79.yr:(2,13)
 2  ~\lstbox{\textSFxi{}}~     let f = copy (|| => { 1 });
    ~\lstbox{\textSFv{}}~             ^^^^
    ~\lstbox{\textSFii{}}~~\lstbox{\textSFx{}}~~\lstbox{\textSFx{}}~~\lstbox{\textSFx{}}~~\lstbox{\textSFx{}}~~\lstbox{\textSFx{}}~~\lstbox{\textSFx{}}~~\lstbox{\textSFx{}}~
\end{lstlisting}

\subsubsection*{How to fix it}

Drop the \tokennolst{copy}.

\errorcode{E4280}{\errhole{} called with \errhole{} works with multiple candidates}
\label{sec:codes:e4280}

Raised during \textbf{validation}, from \texttt{ValidateErrorMessage::\allowbreak{}WORKS\_\allowbreak{}WITH\_\allowbreak{}BOTH} (\textit{src/\allowbreak{}ymirc/\allowbreak{}semantic/\allowbreak{}validator/\allowbreak{}errors.yr}).

\subsubsection*{What it means}

A call matches several candidates equally well, so overload resolution has no single answer. The candidates are listed under the diagnostic.

\subsubsection*{Example}

\textit{test\_\allowbreak{}resources/\allowbreak{}root\_\allowbreak{}path/\allowbreak{}test7.yr}:

%% check: skip
\begin{lstlisting}[style=coloredverbatimError]
mod A {
    pub fn foo ()-> i32 { 12 }
}

mod B {
    mod A {
        pub fn foo ()-> f64 { 12.0 }
    }

    pub fn bar ()-> f64 {
        A::foo ()
    }
}

fn main () {
    let a = B::bar ();
    a;
}
\end{lstlisting}

\begin{lstlisting}[style=bashVerb, breaklines=true, escapechar=~]
Error[E4280] : foo called with {} works with multiple candidates
 --> test_resources/root_path/test7.yr:(11,16)
11  ~\lstbox{\textSFxi{}}~         A::foo ()
    ~\lstbox{\textSFv{}}~                ^  note: candidate test7::A::foo ()-> i32 -- (test_resources/root_path/test7.yr:(2,12))
    ~\lstbox{\textSFxi{}}~                ~\lstbox{\textSFii{}}~~\lstbox{\textSFx{}}~ note: candidate test7::B::A::foo ()-> f64 -- (test_resources/root_path/test7.yr:(7,16))
    ~\lstbox{\textSFii{}}~~\lstbox{\textSFx{}}~~\lstbox{\textSFx{}}~~\lstbox{\textSFx{}}~~\lstbox{\textSFx{}}~~\lstbox{\textSFx{}}~~\lstbox{\textSFx{}}~~\lstbox{\textSFx{}}~
\end{lstlisting}

\subsubsection*{How to fix it}

Make the call select one: convert an argument explicitly, pass an argument by name, or qualify the symbol with its module.

\errorcode{E4283}{the parameter \errhole{} has a default value, it can only be passed by name}
\label{sec:codes:e4283}

Raised during \textbf{validation}, from \texttt{ValidateErrorMessage::\allowbreak{}DEFAULT\_\allowbreak{}PARAM\_\allowbreak{}BY\_\allowbreak{}POSITION} (\textit{src/\allowbreak{}ymirc/\allowbreak{}semantic/\allowbreak{}validator/\allowbreak{}errors.yr}).

\subsubsection*{What it means}

A parameter that declares a default value is filled by name or left to its default -- it is never a slot a positional argument lands in. A call that provides one argument too many for such a prototype was written expecting the opposite, so this hangs under the \hyperref[sec:codes:e4212]{E4212} arity error and names the parameter responsible for the difference between the two counts.

It is not raised when a call is \textit{short} of an argument: the missing parameter is then one that has no default, and naming a defaulted one would point at the wrong place.

A tuple cast to a record is expanded into a positional call, so it reaches this error too. A tuple has no name syntax, so the fix there is to drop the element and let the field take its default.

\subsubsection*{Example}

\textit{test\_\allowbreak{}resources/\allowbreak{}diagnostics/\allowbreak{}test83.yr}:

%% check: skip
\begin{lstlisting}[style=coloredverbatimError]
fn foo (x : i32, y : i32 = 8)-> i32 {
    x + y
}

fn main () {
    foo (1, 2);
}
\end{lstlisting}

\begin{lstlisting}[style=bashVerb, breaklines=true, escapechar=~]
Error[E4212] : 1 parameters is expected, but 2 are provided
 --> test_resources/diagnostics/test83.yr:(6,9)
 6  ~\lstbox{\textSFxi{}}~     foo (1, 2);
    ~\lstbox{\textSFv{}}~         ^
    ~\lstbox{\textSFxi{}}~ Error[E4283] : the parameter y has a default value, it can only be passed by name
    ~\lstbox{\textSFxi{}}~  --> test_resources/diagnostics/test83.yr:(1,18)
    ~\lstbox{\textSFxi{}}~  1  ~\lstbox{\textSFxi{}}~ fn foo (x : i32, y : i32 = 8)-> i32 {
    ~\lstbox{\textSFxi{}}~     ~\lstbox{\textSFv{}}~                  ^
    ~\lstbox{\textSFxi{}}~     ~\lstbox{\textSFii{}}~~\lstbox{\textSFx{}}~~\lstbox{\textSFx{}}~~\lstbox{\textSFx{}}~~\lstbox{\textSFx{}}~~\lstbox{\textSFx{}}~~\lstbox{\textSFx{}}~~\lstbox{\textSFx{}}~
    ~\lstbox{\textSFii{}}~~\lstbox{\textSFx{}}~~\lstbox{\textSFx{}}~~\lstbox{\textSFx{}}~~\lstbox{\textSFx{}}~~\lstbox{\textSFx{}}~~\lstbox{\textSFvii{}}~~\lstbox{\textSFx{}}~
\end{lstlisting}

\subsubsection*{How to fix it}

Pass the parameter by name:

%% check: skip
\begin{lstlisting}[style=coloredverbatim]
foo (1, y-> 2);
\end{lstlisting}

or drop the argument and let the default apply:

%% check: skip
\begin{lstlisting}[style=coloredverbatim]
foo (1);
\end{lstlisting}

\errorcode{E4284}{the parameter \errhole{} can only be passed by name}
\label{sec:codes:e4284}

Raised during \textbf{validation}, from \texttt{ValidateErrorMessage::\allowbreak{}NAMED\_\allowbreak{}PARAM\_\allowbreak{}BY\_\allowbreak{}POSITION} (\textit{src/\allowbreak{}ymirc/\allowbreak{}semantic/\allowbreak{}validator/\allowbreak{}errors.yr}).

\subsubsection*{What it means}

A parameter marked \tokennolst{@} is forced named: a call fills it by name, in any position relative to the positional arguments, and never by position. A call that leaves it unnamed is rejected, one error per such parameter, each pointing at the parameter's declaration.

When the positional arguments are otherwise the right number, this is the only reason the call fails. When they are not, it comes with the \hyperref[sec:codes:e4212]{E4212} arity error it explains: the positional values that were written for the marked parameters are the ones counted too many or too few.

The call is rejected as a candidate, so these errors are listed under the candidate note of \hyperref[sec:codes:e4226]{E4226}.

\subsubsection*{Example}

\textit{test\_\allowbreak{}resources/\allowbreak{}diagnostics/\allowbreak{}test145.yr}:

%% check: skip
\begin{lstlisting}[style=coloredverbatimError]
fn foo (@a: i32, b: i32)-> i32 { a + b }

fn main () { foo (1); }
\end{lstlisting}

\begin{lstlisting}[style=bashVerb, breaklines=true, escapechar=~]
Error[E4226] : the call operator is not defined for test145::foo (@a: i32, b: i32)-> i32 and {1: i32}
 --> test_resources/diagnostics/test145.yr:(3,18)
 1  ~\lstbox{\textSFxi{}}~ fn foo (@a: i32, b: i32)-> i32 { a + b }
    ~\lstbox{\textSFv{}}~    ---  note: candidate test145::foo (@a: i32, b: i32)-> i32
    ~\lstbox{\textSFxi{}}~      ~\lstbox{\textSFii{}}~~\lstbox{\textSFx{}}~ error: the parameter a can only be passed by name
 2  ~\lstbox{\textSFxi{}}~
 3  ~\lstbox{\textSFxi{}}~ fn main () { foo (1); }
    ~\lstbox{\textSFv{}}~                  ^
    ~\lstbox{\textSFxi{}}~
    = when validating test145::main -- (test_resources/diagnostics/test145.yr:(3,4))
    ~\lstbox{\textSFii{}}~~\lstbox{\textSFx{}}~~\lstbox{\textSFx{}}~~\lstbox{\textSFx{}}~~\lstbox{\textSFx{}}~~\lstbox{\textSFx{}}~~\lstbox{\textSFx{}}~~\lstbox{\textSFx{}}~
\end{lstlisting}

\subsubsection*{How to fix it}

Name the parameter at the call site:

%% check: skip
\begin{lstlisting}[style=coloredverbatim]
foo (a-> 1, 2);
\end{lstlisting}

\errorcode{E4285}{parameter with a default value cannot be marked \errhole{}, it can already only be passed by name}
\label{sec:codes:e4285}

Raised during \textbf{validation}, from \texttt{ValidateErrorMessage::\allowbreak{}NAMED\_\allowbreak{}PARAM\_\allowbreak{}WITH\_\allowbreak{}DEFAULT} (\textit{src/\allowbreak{}ymirc/\allowbreak{}semantic/\allowbreak{}validator/\allowbreak{}errors.yr}).

\subsubsection*{What it means}

The \tokennolst{@} marker forces a parameter to be passed by name. A parameter that declares a default value is already filled by name or left to its default, never by position (cf. \hyperref[sec:codes:e4283]{E4283}), so the marker has nothing left to ask and the two cannot be combined. The error points at the parameter and at its default value.

\subsubsection*{Example}

\textit{test\_\allowbreak{}resources/\allowbreak{}diagnostics/\allowbreak{}test84.yr}:

%% check: skip
\begin{lstlisting}[style=coloredverbatimError]
fn foo (@a: i32 = 2)-> i32 { a }

fn main () { foo (); }
\end{lstlisting}

\begin{lstlisting}[style=bashVerb, breaklines=true, escapechar=~]
Error[E4285] : parameter with a default value cannot be marked @, it can already only be passed by name
 --> test_resources/diagnostics/test84.yr:(1,10)
 1  ~\lstbox{\textSFxi{}}~ fn foo (@a: i32 = 2)-> i32 { a }
    ~\lstbox{\textSFv{}}~          ^        -
    ~\lstbox{\textSFxi{}}~
    = when validating test84::foo -- (test_resources/diagnostics/test84.yr:(1,4))
    ~\lstbox{\textSFii{}}~~\lstbox{\textSFx{}}~~\lstbox{\textSFx{}}~~\lstbox{\textSFx{}}~~\lstbox{\textSFx{}}~~\lstbox{\textSFx{}}~~\lstbox{\textSFx{}}~~\lstbox{\textSFx{}}~
Error[E4104] : symbol test84::foo has errors
 --> test_resources/diagnostics/test84.yr:(3,14)
 1  ~\lstbox{\textSFxi{}}~ fn foo (@a: i32 = 2)-> i32 { a }
    ~\lstbox{\textSFv{}}~          -        -  error: parameter with a default value cannot be marked @, it can already only be passed by name
 2  ~\lstbox{\textSFxi{}}~
 3  ~\lstbox{\textSFxi{}}~ fn main () { foo (); }
    ~\lstbox{\textSFv{}}~              ^^^
    ~\lstbox{\textSFxi{}}~
    = when validating test84::foo -- (test_resources/diagnostics/test84.yr:(1,4))
    = when validating test84::main -- (test_resources/diagnostics/test84.yr:(3,4))
    ~\lstbox{\textSFii{}}~~\lstbox{\textSFx{}}~~\lstbox{\textSFx{}}~~\lstbox{\textSFx{}}~~\lstbox{\textSFx{}}~~\lstbox{\textSFx{}}~~\lstbox{\textSFx{}}~~\lstbox{\textSFx{}}~
\end{lstlisting}

\subsubsection*{How to fix it}

Drop the marker, the parameter stays passed by name:

%% check: skip
\begin{lstlisting}[style=coloredverbatim]
fn foo (a: i32 = 2)-> i32 { a }
\end{lstlisting}

or drop the default value, if the parameter has to be provided at every call.

\errorcode{E4298}{the method \errhole{} is not a value, it can only be called}
\label{sec:codes:e4298}

Raised during \textbf{validation}, from \texttt{ValidateErrorMessage::\allowbreak{}FAKE\_\allowbreak{}METHOD\_\allowbreak{}CALL} (\textit{src/\allowbreak{}ymirc/\allowbreak{}semantic/\allowbreak{}validator/\allowbreak{}errors.yr}).

\subsubsection*{What it means}

A fake method is not a field, it is a method syntax the compiler rewrites when it is called. There is nothing to refer to before that rewriting, so the access cannot be used as a value of its own. They are \tokennolst{map:.remove}, \tokennolst{range.stepBy}, \tokennolst{range.reverse} and \tokennolst{generator.next}.

\subsubsection*{Example}

\textit{test\_\allowbreak{}resources/\allowbreak{}generators/\allowbreak{}test32.yr}:

%% check: skip
\begin{lstlisting}[style=coloredverbatimError]
// the fake method next is not a value
yield range (i: i32, j: i32)-> i32 {
    for n in i .. j {
        yield n;
    }
}

fn main () {
    let r = range (0, 2);
    let n = r.next;
    n;
}
\end{lstlisting}

\begin{lstlisting}[style=bashVerb, breaklines=true, escapechar=~]
Error[E4298] : the method next is not a value, it can only be called
 --> test_resources/generators/test32.yr:(10,14)
10  ~\lstbox{\textSFxi{}}~     let n = r.next;
    ~\lstbox{\textSFv{}}~              ^----
    ~\lstbox{\textSFxi{}}~
    = when validating test32::main -- (test_resources/generators/test32.yr:(8,4))
    ~\lstbox{\textSFii{}}~~\lstbox{\textSFx{}}~~\lstbox{\textSFx{}}~~\lstbox{\textSFx{}}~~\lstbox{\textSFx{}}~~\lstbox{\textSFx{}}~~\lstbox{\textSFx{}}~~\lstbox{\textSFx{}}~
\end{lstlisting}

\subsubsection*{How to fix it}

Call the method. \tokennolst{r.next ()} returns the next value of the generator wrapped in an option, \tokennolst{none} once the generator is over. \tokennolst{r:.next ()} borrows that value mutably instead of copying it constantly, on a mutable generator yielding borrowed data.

%% Generated by tools/gen_error_codes.py from docs/errors/ of the compiler;
%% edit the compiler's pages or tools/error_themes.txt, then rerun it.

\clearpage
\section{Classes, records and traits}
\label{sec:codes:classes}

\begin{longtable}{@{}l p{0.78\linewidth} r@{}}
\toprule Code & Message & Page\\ \midrule \endhead
\hyperref[sec:codes:e3001]{E3001} & a class cannot be both abstract and final & \pageref{sec:codes:e3001}\\
\hyperref[sec:codes:e3006]{E3006} & the ctor built from the fields is already declared & \pageref{sec:codes:e3006}\\
\hyperref[sec:codes:e3007]{E3007} & the ctor built from the fields is implicit in a data record & \pageref{sec:codes:e3007}\\
\hyperref[sec:codes:e3010]{E3010} & a class may only define one destructor & \pageref{sec:codes:e3010}\\
\hyperref[sec:codes:e3015]{E3015} & a record cannot be both overlaid and packed & \pageref{sec:codes:e3015}\\
\hyperref[sec:codes:e3030]{E3030} & unexpected declaration inside an impl block & \pageref{sec:codes:e3030}\\
\hyperref[sec:codes:e3033]{E3033} & unexpected declaration inside a trait & \pageref{sec:codes:e3033}\\
\hyperref[sec:codes:e4011]{E4011} & cannot cast derived type \errhole{} to base class type \errhole{} in that context & \pageref{sec:codes:e4011}\\
\hyperref[sec:codes:e4017]{E4017} & field \errhole{} has not yet been validated & \pageref{sec:codes:e4017}\\
\hyperref[sec:codes:e4018]{E4018} & class type \errhole{} does not implement trait \errhole{} & \pageref{sec:codes:e4018}\\
\hyperref[sec:codes:e4019]{E4019} & class \errhole{} has no field named \errhole{} & \pageref{sec:codes:e4019}\\
\hyperref[sec:codes:e4020]{E4020} & class \errhole{} has no method named \errhole{} & \pageref{sec:codes:e4020}\\
\hyperref[sec:codes:e4024]{E4024} & colliding method definitions \errhole{} and \errhole{} & \pageref{sec:codes:e4024}\\
\hyperref[sec:codes:e4026]{E4026} & cannot construct an instance of class \errhole{} that is abstract & \pageref{sec:codes:e4026}\\
\hyperref[sec:codes:e4027]{E4027} & method \errhole{} is defined as constant & \pageref{sec:codes:e4027}\\
\hyperref[sec:codes:e4035]{E4035} & infinite constructor redirection calls when calling \errhole{} & \pageref{sec:codes:e4035}\\
\hyperref[sec:codes:e4044]{E4044} & field must have a referenceable name & \pageref{sec:codes:e4044}\\
\hyperref[sec:codes:e4051]{E4051} & cannot perform a direct call of the method \errhole{} that is empty & \pageref{sec:codes:e4051}\\
\hyperref[sec:codes:e4060]{E4060} & field access method must return a value & \pageref{sec:codes:e4060}\\
\hyperref[sec:codes:e4062]{E4062} & field assign method must be mutable & \pageref{sec:codes:e4062}\\
\hyperref[sec:codes:e4063]{E4063} & field assign method must not return a value & \pageref{sec:codes:e4063}\\
\hyperref[sec:codes:e4065]{E4065} & field method \errhole{} is mutable & \pageref{sec:codes:e4065}\\
\hyperref[sec:codes:e4066]{E4066} & field method cannot have more than one parameter & \pageref{sec:codes:e4066}\\
\hyperref[sec:codes:e4079]{E4079} & the field \errhole{} is not public, and has no initial value to be constructed from & \pageref{sec:codes:e4079}\\
\hyperref[sec:codes:e4090]{E4090} & method \errhole{} implicitly overrides method \errhole{} & \pageref{sec:codes:e4090}\\
\hyperref[sec:codes:e4091]{E4091} & implicit override of method \errhole{} by trait impl & \pageref{sec:codes:e4091}\\
\hyperref[sec:codes:e4093]{E4093} & trait \errhole{} is implemented multiple times by \errhole{} & \pageref{sec:codes:e4093}\\
\hyperref[sec:codes:e4094]{E4094} & impl statement must be followed by a trait, not \errhole{} & \pageref{sec:codes:e4094}\\
\hyperref[sec:codes:e4101]{E4101} & the base class \errhole{} is marked as final & \pageref{sec:codes:e4101}\\
\hyperref[sec:codes:e4102]{E4102} & the base of a class must be a class, not a \errhole{} & \pageref{sec:codes:e4102}\\
\hyperref[sec:codes:e4103]{E4103} & method \errhole{} cannot be inline as it is virtual(overridable or inherited) & \pageref{sec:codes:e4103}\\
\hyperref[sec:codes:e4133]{E4133} & field \errhole{} is initialized multiple times & \pageref{sec:codes:e4133}\\
\hyperref[sec:codes:e4141]{E4141} & method \errhole{} is defined as mutable & \pageref{sec:codes:e4141}\\
\hyperref[sec:codes:e4144]{E4144} & class \errhole{} is not abstract but has empty methods & \pageref{sec:codes:e4144}\\
\hyperref[sec:codes:e4147]{E4147} & \errhole{} is not an ancestor type of \errhole{} & \pageref{sec:codes:e4147}\\
\hyperref[sec:codes:e4150]{E4150} & \errhole{} is not a class type & \pageref{sec:codes:e4150}\\
\hyperref[sec:codes:e4162]{E4162} & \errhole{} is not a trait & \pageref{sec:codes:e4162}\\
\hyperref[sec:codes:e4168]{E4168} & no constructor found for class \errhole{} & \pageref{sec:codes:e4168}\\
\hyperref[sec:codes:e4169]{E4169} & no constructor named \errhole{} found for class \errhole{} & \pageref{sec:codes:e4169}\\
\hyperref[sec:codes:e4172]{E4172} & record or entity type \errhole{} has no size due to forward references & \pageref{sec:codes:e4172}\\
\hyperref[sec:codes:e4173]{E4173} & class \errhole{} has no ancestor & \pageref{sec:codes:e4173}\\
\hyperref[sec:codes:e4181]{E4181} & method \errhole{} must have a body to override \errhole{} & \pageref{sec:codes:e4181}\\
\hyperref[sec:codes:e4182]{E4182} & method \errhole{} must be a field method to override \errhole{} & \pageref{sec:codes:e4182}\\
\hyperref[sec:codes:e4183]{E4183} & cannot override final method \errhole{} & \pageref{sec:codes:e4183}\\
\hyperref[sec:codes:e4184]{E4184} & the return type of the overriding method \errhole{} is not compatible with the return type of the ancestor method \errhole{} & \pageref{sec:codes:e4184}\\
\hyperref[sec:codes:e4185]{E4185} & the protection \errhole{} of the overriding method \errhole{} does not match the definition in the ancestor class \errhole{} & \pageref{sec:codes:e4185}\\
\hyperref[sec:codes:e4186]{E4186} & the throwers of the overriding method \errhole{} are not compatible with the throwers of the ancestor method \errhole{} & \pageref{sec:codes:e4186}\\
\hyperref[sec:codes:e4187]{E4187} & trait method \errhole{} was already overridden & \pageref{sec:codes:e4187}\\
\hyperref[sec:codes:e4188]{E4188} & cannot override a non trait method \errhole{} with \errhole{} inside impl & \pageref{sec:codes:e4188}\\
\hyperref[sec:codes:e4189]{E4189} & method \errhole{} overrides nothing & \pageref{sec:codes:e4189}\\
\hyperref[sec:codes:e4190]{E4190} & field method \errhole{} cannot override the non field method \errhole{} & \pageref{sec:codes:e4190}\\
\hyperref[sec:codes:e4192]{E4192} & cannot override trait method \errhole{} with \errhole{} outside impl & \pageref{sec:codes:e4192}\\
\hyperref[sec:codes:e4196]{E4196} & ancestor cycle found when validating \errhole{} & \pageref{sec:codes:e4196}\\
\hyperref[sec:codes:e4233]{E4233} & no constructor is callable with the parameters \errhole{} & \pageref{sec:codes:e4233}\\
\hyperref[sec:codes:e4253]{E4253} & the field \errhole{} has no initial value & \pageref{sec:codes:e4253}\\
\hyperref[sec:codes:e4319]{E4319} & method \errhole{} must be a generator to override the generator \errhole{} & \pageref{sec:codes:e4319}\\
\hyperref[sec:codes:e4320]{E4320} & generator \errhole{} cannot override the non generator method \errhole{} & \pageref{sec:codes:e4320}\\
\bottomrule
\end{longtable}

\errorcode{E3001}{a class cannot be both abstract and final}
\label{sec:codes:e3001}

Raised during \textbf{declaration}, from \texttt{DeclareErrorMessage::\allowbreak{}ABSTRACT\_\allowbreak{}AND\_\allowbreak{}FINAL} (\textit{src/\allowbreak{}ymirc/\allowbreak{}semantic/\allowbreak{}declarator/\allowbreak{}errors.yr}).

\subsubsection*{What it means}

A class is declared both \tokennolst{abstract} and \tokennolst{final}. \tokennolst{abstract} means the class cannot be instantiated and must be derived from; \tokennolst{final} means it cannot be derived from. Together they describe a class nothing can ever do anything with.

\subsubsection*{Example}

\textit{test\_\allowbreak{}resources/\allowbreak{}diagnostics/\allowbreak{}test4.yr}:

%% check: skip
\begin{lstlisting}[style=coloredverbatimError]
@abstract @final
class Foo {
    pub self () {}
}

fn main () {}
\end{lstlisting}

\begin{lstlisting}[style=bashVerb, breaklines=true, escapechar=~]
Error[E3001] : a class cannot be both abstract and final
 --> test_resources/diagnostics/test4.yr:(1,12)
 1  ~\lstbox{\textSFxi{}}~ @abstract @final
    ~\lstbox{\textSFv{}}~            ^^^^^
    ~\lstbox{\textSFxi{}}~ Note :
    ~\lstbox{\textSFxi{}}~  --> test_resources/diagnostics/test4.yr:(1,2)
    ~\lstbox{\textSFxi{}}~  1  ~\lstbox{\textSFxi{}}~ @abstract @final
    ~\lstbox{\textSFxi{}}~     ~\lstbox{\textSFv{}}~  ^^^^^^^^
    ~\lstbox{\textSFxi{}}~     ~\lstbox{\textSFii{}}~~\lstbox{\textSFx{}}~~\lstbox{\textSFx{}}~~\lstbox{\textSFx{}}~~\lstbox{\textSFx{}}~~\lstbox{\textSFx{}}~~\lstbox{\textSFx{}}~~\lstbox{\textSFx{}}~
    ~\lstbox{\textSFxi{}}~ Note :
    ~\lstbox{\textSFxi{}}~  --> test_resources/diagnostics/test4.yr:(2,1)
    ~\lstbox{\textSFxi{}}~  2  ~\lstbox{\textSFxi{}}~ class Foo {
    ~\lstbox{\textSFxi{}}~     ~\lstbox{\textSFv{}}~ ^^^^^
    ~\lstbox{\textSFxi{}}~     ~\lstbox{\textSFii{}}~~\lstbox{\textSFx{}}~~\lstbox{\textSFx{}}~~\lstbox{\textSFx{}}~~\lstbox{\textSFx{}}~~\lstbox{\textSFx{}}~~\lstbox{\textSFx{}}~~\lstbox{\textSFx{}}~
    ~\lstbox{\textSFii{}}~~\lstbox{\textSFx{}}~~\lstbox{\textSFx{}}~~\lstbox{\textSFx{}}~~\lstbox{\textSFx{}}~~\lstbox{\textSFx{}}~~\lstbox{\textSFvii{}}~~\lstbox{\textSFx{}}~
\end{lstlisting}

\subsubsection*{How to fix it}

Drop whichever of the two does not apply.

\errorcode{E3006}{the ctor built from the fields is already declared}
\label{sec:codes:e3006}

Raised during \textbf{declaration}, from \texttt{DeclareErrorMessage::\allowbreak{}DEFAULT\_\allowbreak{}CTOR\_\allowbreak{}ALREADY\_\allowbreak{}DECLARED} (\textit{src/\allowbreak{}ymirc/\allowbreak{}semantic/\allowbreak{}declarator/\allowbreak{}errors.yr}).

\subsubsection*{What it means}

The constructor built from the fields of the type is declared explicitly while the compiler also generates it. Both would have the same signature.

\subsubsection*{Example}

\textit{test\_\allowbreak{}resources/\allowbreak{}default\_\allowbreak{}ctor/\allowbreak{}test11.yr}:

%% check: skip
\begin{lstlisting}[style=coloredverbatimError]
// The request is a declaration of its own, so a class makes it once.
class Foo {
   pub let x: i32;

   pub self;
   prv self;
}

fn main () {
    let a = copy Foo (1);
    a.x;
}
\end{lstlisting}

\begin{lstlisting}[style=bashVerb, breaklines=true, escapechar=~]
Error[E3006] : the ctor built from the fields is already declared
 --> test_resources/default_ctor/test11.yr:(6,8)
 6  ~\lstbox{\textSFxi{}}~    prv self;
    ~\lstbox{\textSFv{}}~        ^^^^
    ~\lstbox{\textSFxi{}}~ Note : declared here
    ~\lstbox{\textSFxi{}}~  --> test_resources/default_ctor/test11.yr:(5,8)
    ~\lstbox{\textSFxi{}}~  5  ~\lstbox{\textSFxi{}}~    pub self;
    ~\lstbox{\textSFxi{}}~     ~\lstbox{\textSFv{}}~        ^^^^
    ~\lstbox{\textSFxi{}}~     ~\lstbox{\textSFii{}}~~\lstbox{\textSFx{}}~~\lstbox{\textSFx{}}~~\lstbox{\textSFx{}}~~\lstbox{\textSFx{}}~~\lstbox{\textSFx{}}~~\lstbox{\textSFx{}}~~\lstbox{\textSFx{}}~
    ~\lstbox{\textSFii{}}~~\lstbox{\textSFx{}}~~\lstbox{\textSFx{}}~~\lstbox{\textSFx{}}~~\lstbox{\textSFx{}}~~\lstbox{\textSFx{}}~~\lstbox{\textSFvii{}}~~\lstbox{\textSFx{}}~
\end{lstlisting}

\subsubsection*{How to fix it}

Remove the explicit constructor and rely on the generated one, or change its parameters so the two no longer collide.

\errorcode{E3007}{the ctor built from the fields is implicit in a data record}
\label{sec:codes:e3007}

Raised during \textbf{declaration}, from \texttt{DeclareErrorMessage::\allowbreak{}DEFAULT\_\allowbreak{}CTOR\_\allowbreak{}IN\_\allowbreak{}DATA} (\textit{src/\allowbreak{}ymirc/\allowbreak{}semantic/\allowbreak{}declarator/\allowbreak{}errors.yr}).

\subsubsection*{What it means}

A \tokennolst{data} record generates its field constructor implicitly, and this declaration writes it out again.

\subsubsection*{Example}

\textit{test\_\allowbreak{}resources/\allowbreak{}default\_\allowbreak{}ctor/\allowbreak{}test10.yr}:

%% check: skip
\begin{lstlisting}[style=coloredverbatimError]
// A `@data` record asks for the ctor by its attribute, the request is implicit
// and is not redeclarable.
@data
record Foo {
    let x: i32;

    pub self;
}

fn main () {
    let a = Foo (1);
    a.x;
}
\end{lstlisting}

\begin{lstlisting}[style=bashVerb, breaklines=true, escapechar=~]
Error[E3007] : the ctor built from the fields is implicit in a data record
 --> test_resources/default_ctor/test10.yr:(7,9)
 7  ~\lstbox{\textSFxi{}}~     pub self;
    ~\lstbox{\textSFv{}}~         ^^^^
    ~\lstbox{\textSFxi{}}~ Note : declared here
    ~\lstbox{\textSFxi{}}~  --> test_resources/default_ctor/test10.yr:(4,8)
    ~\lstbox{\textSFxi{}}~  4  ~\lstbox{\textSFxi{}}~ record Foo {
    ~\lstbox{\textSFxi{}}~     ~\lstbox{\textSFv{}}~        ^^^
    ~\lstbox{\textSFxi{}}~     ~\lstbox{\textSFii{}}~~\lstbox{\textSFx{}}~~\lstbox{\textSFx{}}~~\lstbox{\textSFx{}}~~\lstbox{\textSFx{}}~~\lstbox{\textSFx{}}~~\lstbox{\textSFx{}}~~\lstbox{\textSFx{}}~
    ~\lstbox{\textSFii{}}~~\lstbox{\textSFx{}}~~\lstbox{\textSFx{}}~~\lstbox{\textSFx{}}~~\lstbox{\textSFx{}}~~\lstbox{\textSFx{}}~~\lstbox{\textSFvii{}}~~\lstbox{\textSFx{}}~
\end{lstlisting}

\subsubsection*{How to fix it}

Remove the explicit constructor: a \tokennolst{data} record always has it.

\errorcode{E3010}{a class may only define one destructor}
\label{sec:codes:e3010}

Raised during \textbf{declaration}, from \texttt{DeclareErrorMessage::\allowbreak{}MULTIPLE\_\allowbreak{}DESTRUCTOR} (\textit{src/\allowbreak{}ymirc/\allowbreak{}semantic/\allowbreak{}declarator/\allowbreak{}errors.yr}).

\subsubsection*{What it means}

A class declares more than one destructor. A destructor takes no parameter, so there is nothing to overload on.

\subsubsection*{Example}

\textit{test\_\allowbreak{}resources/\allowbreak{}diagnostics/\allowbreak{}test6.yr}:

%% check: skip
\begin{lstlisting}[style=coloredverbatimError]
class Foo {
    pub self () {}
    __dtor (mut self) {}
    __dtor (mut self) {}
}

fn main () {}
\end{lstlisting}

\begin{lstlisting}[style=bashVerb, breaklines=true, escapechar=~]
Error[E3010] : a class may only define one destructor
 --> test_resources/diagnostics/test6.yr:(4,5)
 4  ~\lstbox{\textSFxi{}}~     __dtor (mut self) {}
    ~\lstbox{\textSFv{}}~     ^^^^^^
    ~\lstbox{\textSFxi{}}~ Note :
    ~\lstbox{\textSFxi{}}~  --> test_resources/diagnostics/test6.yr:(3,5)
    ~\lstbox{\textSFxi{}}~  3  ~\lstbox{\textSFxi{}}~     __dtor (mut self) {}
    ~\lstbox{\textSFxi{}}~     ~\lstbox{\textSFv{}}~     ^^^^^^
    ~\lstbox{\textSFxi{}}~     ~\lstbox{\textSFii{}}~~\lstbox{\textSFx{}}~~\lstbox{\textSFx{}}~~\lstbox{\textSFx{}}~~\lstbox{\textSFx{}}~~\lstbox{\textSFx{}}~~\lstbox{\textSFx{}}~~\lstbox{\textSFx{}}~
    ~\lstbox{\textSFii{}}~~\lstbox{\textSFx{}}~~\lstbox{\textSFx{}}~~\lstbox{\textSFx{}}~~\lstbox{\textSFx{}}~~\lstbox{\textSFx{}}~~\lstbox{\textSFvii{}}~~\lstbox{\textSFx{}}~
\end{lstlisting}

\subsubsection*{How to fix it}

Keep one destructor and merge the bodies.

\errorcode{E3015}{a record cannot be both overlaid and packed}
\label{sec:codes:e3015}

Raised during \textbf{declaration}, from \texttt{DeclareErrorMessage::\allowbreak{}OVERLAID\_\allowbreak{}AND\_\allowbreak{}PACKED} (\textit{src/\allowbreak{}ymirc/\allowbreak{}semantic/\allowbreak{}declarator/\allowbreak{}errors.yr}).

\subsubsection*{What it means}

A record is declared both overlaid and packed. Overlaying places every field at the same offset; packing removes the padding between successive fields. There is no layout satisfying both.

\subsubsection*{Example}

\textit{test\_\allowbreak{}resources/\allowbreak{}diagnostics/\allowbreak{}test8.yr}:

%% check: skip
\begin{lstlisting}[style=coloredverbatimError]
@overlaid @packed
record Foo {
    let x: i32;
    let y: i32;
}

fn main () { let _ = Foo (1, 2); }
\end{lstlisting}

\begin{lstlisting}[style=bashVerb, breaklines=true, escapechar=~]
Error[E3015] : a record cannot be both overlaid and packed
 --> test_resources/diagnostics/test8.yr:(1,12)
 1  ~\lstbox{\textSFxi{}}~ @overlaid @packed
    ~\lstbox{\textSFv{}}~            ^^^^^^
    ~\lstbox{\textSFxi{}}~ Note :
    ~\lstbox{\textSFxi{}}~  --> test_resources/diagnostics/test8.yr:(1,2)
    ~\lstbox{\textSFxi{}}~  1  ~\lstbox{\textSFxi{}}~ @overlaid @packed
    ~\lstbox{\textSFxi{}}~     ~\lstbox{\textSFv{}}~  ^^^^^^^^
    ~\lstbox{\textSFxi{}}~     ~\lstbox{\textSFii{}}~~\lstbox{\textSFx{}}~~\lstbox{\textSFx{}}~~\lstbox{\textSFx{}}~~\lstbox{\textSFx{}}~~\lstbox{\textSFx{}}~~\lstbox{\textSFx{}}~~\lstbox{\textSFx{}}~
    ~\lstbox{\textSFxi{}}~ Note :
    ~\lstbox{\textSFxi{}}~  --> test_resources/diagnostics/test8.yr:(2,1)
    ~\lstbox{\textSFxi{}}~  2  ~\lstbox{\textSFxi{}}~ record Foo {
    ~\lstbox{\textSFxi{}}~     ~\lstbox{\textSFv{}}~ ^^^^^^
    ~\lstbox{\textSFxi{}}~     ~\lstbox{\textSFii{}}~~\lstbox{\textSFx{}}~~\lstbox{\textSFx{}}~~\lstbox{\textSFx{}}~~\lstbox{\textSFx{}}~~\lstbox{\textSFx{}}~~\lstbox{\textSFx{}}~~\lstbox{\textSFx{}}~
    ~\lstbox{\textSFii{}}~~\lstbox{\textSFx{}}~~\lstbox{\textSFx{}}~~\lstbox{\textSFx{}}~~\lstbox{\textSFx{}}~~\lstbox{\textSFx{}}~~\lstbox{\textSFvii{}}~~\lstbox{\textSFx{}}~
\end{lstlisting}

\subsubsection*{How to fix it}

Keep whichever describes the layout wanted. A union-like type is overlaid; a type mirroring a wire format is packed.

\errorcode{E3030}{unexpected declaration inside an impl block}
\label{sec:codes:e3030}

Raised during \textbf{declaration}, from \texttt{DeclareErrorMessage::\allowbreak{}UNEXPECTED\_\allowbreak{}IN\_\allowbreak{}IMPL} (\textit{src/\allowbreak{}ymirc/\allowbreak{}semantic/\allowbreak{}declarator/\allowbreak{}errors.yr}).

\subsubsection*{What it means}

A declaration appeared inside an \tokennolst{impl} block that such a block does not accept. An \tokennolst{impl} block provides the methods of a trait for a type.

\subsubsection*{Example}

\textit{test\_\allowbreak{}resources/\allowbreak{}diagnostics/\allowbreak{}test96.yr}:

%% check: skip
\begin{lstlisting}[style=coloredverbatimError]
// a template method inside an impl block
trait Shape {
    pub fn area (self)-> i32;
}

class Square {
    pub self () {}

    impl Shape {
        pub over area {T} (self)-> i32 { 1 }
    }
}

fn main () {
    let _ = copy Square ();
}
\end{lstlisting}

\begin{lstlisting}[style=bashVerb, breaklines=true, escapechar=~]
Error[E3030] : unexpected declaration inside an impl block
 --> test_resources/diagnostics/test96.yr:(10,18)
10  ~\lstbox{\textSFxi{}}~         pub over area {T} (self)-> i32 { 1 }
    ~\lstbox{\textSFv{}}~                  ^^^^
    ~\lstbox{\textSFii{}}~~\lstbox{\textSFx{}}~~\lstbox{\textSFx{}}~~\lstbox{\textSFx{}}~~\lstbox{\textSFx{}}~~\lstbox{\textSFx{}}~~\lstbox{\textSFx{}}~~\lstbox{\textSFx{}}~
\end{lstlisting}

\subsubsection*{How to fix it}

Move the declaration out of the block. A method that is not part of the trait belongs in the type body rather than in the \tokennolst{impl}.

\errorcode{E3033}{unexpected declaration inside a trait}
\label{sec:codes:e3033}

Raised during \textbf{declaration}, from \texttt{DeclareErrorMessage::\allowbreak{}UNEXPECTED\_\allowbreak{}IN\_\allowbreak{}TRAIT} (\textit{src/\allowbreak{}ymirc/\allowbreak{}semantic/\allowbreak{}declarator/\allowbreak{}errors.yr}).

\subsubsection*{What it means}

A declaration appeared inside a trait that a trait body does not accept. A trait declares methods, with or without a default body.

\subsubsection*{Example}

\textit{test\_\allowbreak{}resources/\allowbreak{}diagnostics/\allowbreak{}test97.yr}:

%% check: skip
\begin{lstlisting}[style=coloredverbatimError]
// a use declaration inside a trait
trait Shape {
    use std::io;
}

fn main () {}
\end{lstlisting}

\begin{lstlisting}[style=bashVerb, breaklines=true, escapechar=~]
Error[E3033] : unexpected declaration inside a trait
 --> test_resources/diagnostics/test97.yr:(3,5)
 3  ~\lstbox{\textSFxi{}}~     use std::io;
    ~\lstbox{\textSFv{}}~     ^^^
    ~\lstbox{\textSFii{}}~~\lstbox{\textSFx{}}~~\lstbox{\textSFx{}}~~\lstbox{\textSFx{}}~~\lstbox{\textSFx{}}~~\lstbox{\textSFx{}}~~\lstbox{\textSFx{}}~~\lstbox{\textSFx{}}~
\end{lstlisting}

\subsubsection*{How to fix it}

Move the declaration out of the trait. A trait cannot hold fields: the state a method needs belongs to the type implementing it.

\errorcode{E4011}{cannot cast derived type \errhole{} to base class type \errhole{} in that context}
\label{sec:codes:e4011}

Raised during \textbf{validation}, from \texttt{ValidateErrorMessage::\allowbreak{}CANNOT\_\allowbreak{}CASTTO\_\allowbreak{}BASE\_\allowbreak{}CLASS} (\textit{src/\allowbreak{}ymirc/\allowbreak{}semantic/\allowbreak{}validator/\allowbreak{}errors.yr}).

\subsubsection*{What it means}

A value of a derived class type is cast to one of its base classes in a context that does not allow it. The cast itself is meaningful, but not where it was written.

\subsubsection*{Example}

\textit{test\_\allowbreak{}resources/\allowbreak{}diagnostics/\allowbreak{}test129.yr}:

%% check: skip
\begin{lstlisting}[style=coloredverbatimError]
// a derived instance bound by reference to its base class in a pattern
class Base {
    pub self () {}
}

class Derived over Base {
    pub self () with super () {}
}

fn main () {
    let d = copy Derived ();
    match d {
        ref b : &Base => { let _ = b; }
        _ => {}
    }
}
\end{lstlisting}

\begin{lstlisting}[style=bashVerb, breaklines=true, escapechar=~]
Error[E4011] : cannot cast derived type &(test129::Derived) to base class type &(test129::Base) in that context
 --> test_resources/diagnostics/test129.yr:(12,11)
12  ~\lstbox{\textSFxi{}}~     match d {
    ~\lstbox{\textSFv{}}~           ^
13  ~\lstbox{\textSFxi{}}~         ref b : &Base => { let _ = b; }
    ~\lstbox{\textSFv{}}~         --- -
    ~\lstbox{\textSFxi{}}~
    = when validating pattern matching with value d of type &(test129::Derived) -- (test_resources/diagnostics/test129.yr:(12,5))
    = when validating test129::main -- (test_resources/diagnostics/test129.yr:(10,4))
    ~\lstbox{\textSFii{}}~~\lstbox{\textSFx{}}~~\lstbox{\textSFx{}}~~\lstbox{\textSFx{}}~~\lstbox{\textSFx{}}~~\lstbox{\textSFx{}}~~\lstbox{\textSFx{}}~~\lstbox{\textSFx{}}~
\end{lstlisting}

\subsubsection*{How to fix it}

Bind the value to a variable of the base class type, which performs the conversion where it is allowed, and use that variable.

\errorcode{E4017}{field \errhole{} has not yet been validated}
\label{sec:codes:e4017}

Raised during \textbf{validation}, from \texttt{ValidateErrorMessage::\allowbreak{}CLASS\_\allowbreak{}FIELD\_\allowbreak{}NOT\_\allowbreak{}VALIDATED\_\allowbreak{}YET} (\textit{src/\allowbreak{}ymirc/\allowbreak{}semantic/\allowbreak{}validator/\allowbreak{}errors.yr}).

\subsubsection*{What it means}

A field is used while the type it belongs to is still being validated, so its type is not known yet. This happens when a type refers to itself through a path the compiler has to resolve before it can finish the type.

\subsubsection*{Example}

\textit{test\_\allowbreak{}resources/\allowbreak{}lit\_\allowbreak{}class/\allowbreak{}ctors/\allowbreak{}test23.yr}:

%% check: skip
\begin{lstlisting}[style=coloredverbatimError]
class A {
    let x : i32;
    let y : i32;

    pub self ()
        with y = self.x
        , x = 9
    {}
}
\end{lstlisting}

\begin{lstlisting}[style=bashVerb, breaklines=true, escapechar=~]
Error[E4017] : field x has not yet been validated
 --> test_resources/lit_class/ctors/test23.yr:(6,23)
 6  ~\lstbox{\textSFxi{}}~         with y = self.x
    ~\lstbox{\textSFv{}}~                       ^
    ~\lstbox{\textSFxi{}}~ Note : declared here
    ~\lstbox{\textSFxi{}}~  --> test_resources/lit_class/ctors/test23.yr:(2,9)
    ~\lstbox{\textSFxi{}}~  2  ~\lstbox{\textSFxi{}}~     let x : i32;
    ~\lstbox{\textSFxi{}}~     ~\lstbox{\textSFv{}}~         ^
    ~\lstbox{\textSFxi{}}~     ~\lstbox{\textSFii{}}~~\lstbox{\textSFx{}}~~\lstbox{\textSFx{}}~~\lstbox{\textSFx{}}~~\lstbox{\textSFx{}}~~\lstbox{\textSFx{}}~~\lstbox{\textSFx{}}~~\lstbox{\textSFx{}}~
    ~\lstbox{\textSFii{}}~~\lstbox{\textSFx{}}~~\lstbox{\textSFx{}}~~\lstbox{\textSFx{}}~~\lstbox{\textSFx{}}~~\lstbox{\textSFx{}}~~\lstbox{\textSFvii{}}~~\lstbox{\textSFx{}}~
\end{lstlisting}

\subsubsection*{How to fix it}

Break the cycle: hold the self-reference behind a class (which is a reference) or a pointer, rather than by value.

\errorcode{E4018}{class type \errhole{} does not implement trait \errhole{}}
\label{sec:codes:e4018}

Raised during \textbf{validation}, from \texttt{ValidateErrorMessage::\allowbreak{}CLASS\_\allowbreak{}NOT\_\allowbreak{}IMPL} (\textit{src/\allowbreak{}ymirc/\allowbreak{}semantic/\allowbreak{}validator/\allowbreak{}errors.yr}).

\subsubsection*{What it means}

A class is used where the trait named is required, and the class does not implement it.

\subsubsection*{Example}

\textit{test\_\allowbreak{}resources/\allowbreak{}templates/\allowbreak{}explicit/\allowbreak{}test13.yr}:

%% check: skip
\begin{lstlisting}[style=coloredverbatimError]
trait R {I} {}

class B  {
    pub self () {}
    impl R!i32;
}

fn foo {T impl X!{I}, X, I of f32} () {
}

fn main () {
    foo!(&B) ();
}
\end{lstlisting}

\begin{lstlisting}[style=bashVerb, breaklines=true, escapechar=~]
Error[E4018] : class type &(test13::B) does not implement trait X!{I}
 --> test_resources/templates/explicit/test13.yr:(8,9)
 8  ~\lstbox{\textSFxi{}}~ fn foo {T impl X!{I}, X, I of f32} () {
    ~\lstbox{\textSFv{}}~         ^
    ~\lstbox{\textSFxi{}}~ Error[E4095] : incompatible types f32 and i32
    ~\lstbox{\textSFxi{}}~  --> test_resources/templates/explicit/test13.yr:(5,12)
    ~\lstbox{\textSFxi{}}~  5  ~\lstbox{\textSFxi{}}~     impl R!i32;
    ~\lstbox{\textSFxi{}}~     ~\lstbox{\textSFv{}}~            ^^^
    ~\lstbox{\textSFxi{}}~     ~\lstbox{\textSFxi{}}~ Note :
    ~\lstbox{\textSFxi{}}~     ~\lstbox{\textSFxi{}}~  --> test_resources/templates/explicit/test13.yr:(8,31)
    ~\lstbox{\textSFxi{}}~     ~\lstbox{\textSFxi{}}~  8  ~\lstbox{\textSFxi{}}~ fn foo {T impl X!{I}, X, I of f32} () {
    ~\lstbox{\textSFxi{}}~     ~\lstbox{\textSFxi{}}~     ~\lstbox{\textSFv{}}~                               ^^^
    ~\lstbox{\textSFxi{}}~     ~\lstbox{\textSFxi{}}~     ~\lstbox{\textSFii{}}~~\lstbox{\textSFx{}}~~\lstbox{\textSFx{}}~~\lstbox{\textSFx{}}~~\lstbox{\textSFx{}}~~\lstbox{\textSFx{}}~~\lstbox{\textSFx{}}~~\lstbox{\textSFx{}}~
    ~\lstbox{\textSFxi{}}~     ~\lstbox{\textSFii{}}~~\lstbox{\textSFx{}}~~\lstbox{\textSFx{}}~~\lstbox{\textSFx{}}~~\lstbox{\textSFx{}}~~\lstbox{\textSFx{}}~~\lstbox{\textSFvii{}}~~\lstbox{\textSFx{}}~
    ~\lstbox{\textSFii{}}~~\lstbox{\textSFx{}}~~\lstbox{\textSFx{}}~~\lstbox{\textSFx{}}~~\lstbox{\textSFx{}}~~\lstbox{\textSFx{}}~~\lstbox{\textSFvii{}}~~\lstbox{\textSFx{}}~
\end{lstlisting}

\subsubsection*{How to fix it}

Add an \tokennolst{impl} block for the trait to the class, or use a type that already implements it.

\errorcode{E4019}{class \errhole{} has no field named \errhole{}}
\label{sec:codes:e4019}

Raised during \textbf{validation}, from \texttt{ValidateErrorMessage::\allowbreak{}CLASS\_\allowbreak{}NO\_\allowbreak{}FIELD} (\textit{src/\allowbreak{}ymirc/\allowbreak{}semantic/\allowbreak{}validator/\allowbreak{}errors.yr}).

\subsubsection*{What it means}

A field name was used on a class that has no field of that name.

\subsubsection*{Example}

\textit{test\_\allowbreak{}resources/\allowbreak{}diagnostics/\allowbreak{}test63.yr}:

%% check: skip
\begin{lstlisting}[style=coloredverbatimError]
class Foo {
    pub self () {
        self.super ();
    }
}

fn main () { let _ = copy Foo (); }
\end{lstlisting}

\begin{lstlisting}[style=bashVerb, breaklines=true, escapechar=~]
Error[E4019] : class mut &(mut test63::Foo) has no field named super
 --> test_resources/diagnostics/test63.yr:(3,14)
 3  ~\lstbox{\textSFxi{}}~         self.super ();
    ~\lstbox{\textSFv{}}~              ^^^^^  note: class mut &(mut test63::Foo) has no method named super
    ~\lstbox{\textSFxi{}}~                  ~\lstbox{\textSFii{}}~~\lstbox{\textSFx{}}~ note: class mut &(mut test63::Foo) has no field access method named super
    ~\lstbox{\textSFxi{}}~ Error[E4173] : class &(test63::Foo) has no ancestor
    ~\lstbox{\textSFxi{}}~  --> test_resources/diagnostics/test63.yr:(3,9)
    ~\lstbox{\textSFxi{}}~  3  ~\lstbox{\textSFxi{}}~         self.super ();
    ~\lstbox{\textSFxi{}}~     ~\lstbox{\textSFv{}}~         ^^^^ ^^^^^
    ~\lstbox{\textSFxi{}}~     ~\lstbox{\textSFii{}}~~\lstbox{\textSFx{}}~~\lstbox{\textSFx{}}~~\lstbox{\textSFx{}}~~\lstbox{\textSFx{}}~~\lstbox{\textSFx{}}~~\lstbox{\textSFx{}}~~\lstbox{\textSFx{}}~
    ~\lstbox{\textSFii{}}~~\lstbox{\textSFx{}}~~\lstbox{\textSFx{}}~~\lstbox{\textSFx{}}~~\lstbox{\textSFx{}}~~\lstbox{\textSFx{}}~~\lstbox{\textSFvii{}}~~\lstbox{\textSFx{}}~
\end{lstlisting}

\subsubsection*{Notes}

When the name is a method rather than a field, the reason it cannot be reached is reported under this diagnostic, as a note carrying no code of its own:

\begin{itemize}
\item \tokennolst{method \textless{}\textgreater{} is a field method} (formerly E4061): a field method was reached through the alias operator \tokennolst{:.}.
\item \texttt{the superclass proxy of \textless{}\textgreater{} is only accessible through self} (formerly E4200): \tokennolst{super} was accessed on a value other than \tokennolst{self}.
\end{itemize}

\subsubsection*{How to fix it}

Check the spelling. A field that exists but is not visible from here gives a protection error instead, so this really means no such field. A field method of that name is reached the same way, so check that it is declared too.

\errorcode{E4020}{class \errhole{} has no method named \errhole{}}
\label{sec:codes:e4020}

Raised during \textbf{validation}, from \texttt{ValidateErrorMessage::\allowbreak{}CLASS\_\allowbreak{}NO\_\allowbreak{}METHOD} (\textit{src/\allowbreak{}ymirc/\allowbreak{}semantic/\allowbreak{}validator/\allowbreak{}errors.yr}).

\subsubsection*{What it means}

A method name was used on a class that has no method of that name.

\subsubsection*{Example}

\textit{test\_\allowbreak{}resources/\allowbreak{}lit\_\allowbreak{}class/\allowbreak{}operators/\allowbreak{}test13.yr}:

%% check: skip
\begin{lstlisting}[style=coloredverbatimError]
class A {
    pub self () {}

    pub fn opIndex (self, _ : i32)-> i32 {
        1
    }

}


class B  {
    pub self () {}
}



fn main () {
    let a = copy A ();
    let _ = a ['i'];

    let b = copy B ();
    let _ = b [1];
}
\end{lstlisting}

\begin{lstlisting}[style=bashVerb, breaklines=true, escapechar=~]
Error[E4020] : class &(test13::B) has no method named opIndex
 --> test_resources/lit_class/operators/test13.yr:(22,15)
22  ~\lstbox{\textSFxi{}}~     let _ = b [1];
    ~\lstbox{\textSFv{}}~               ^
    ~\lstbox{\textSFxi{}}~
    = when validating test13::main -- (test_resources/lit_class/operators/test13.yr:(17,4))
    ~\lstbox{\textSFii{}}~~\lstbox{\textSFx{}}~~\lstbox{\textSFx{}}~~\lstbox{\textSFx{}}~~\lstbox{\textSFx{}}~~\lstbox{\textSFx{}}~~\lstbox{\textSFx{}}~~\lstbox{\textSFx{}}~
\end{lstlisting}

\subsubsection*{How to fix it}

Check the spelling and the class. When the method comes from a trait, the class needs an \tokennolst{impl} block for it, cf. \hyperref[sec:codes:e4018]{E4018}.

\errorcode{E4024}{colliding method definitions \errhole{} and \errhole{}}
\label{sec:codes:e4024}

Raised during \textbf{validation}, from \texttt{ValidateErrorMessage::\allowbreak{}COLLIDING\_\allowbreak{}METHOD\_\allowbreak{}DEFINITION} (\textit{src/\allowbreak{}ymirc/\allowbreak{}semantic/\allowbreak{}validator/\allowbreak{}errors.yr}).

\subsubsection*{What it means}

Two methods of a class have the same name and indistinguishable signatures, cf. \hyperref[sec:codes:e4023]{E4023}.

\subsubsection*{Example}

\textit{test\_\allowbreak{}resources/\allowbreak{}default\_\allowbreak{}ctor/\allowbreak{}test6.yr}:

%% check: skip
\begin{lstlisting}[style=coloredverbatimError]
// A `self;` declares the ctor built from the fields, so a ctor written by hand
// with the same signature declares it twice.
class Foo {
   pub let x: i32;

   pub self (x: i32) with x = x + 1 {}
   pub self;
}

fn main () {
    let a = copy Foo (1);
    a.x;
}
\end{lstlisting}

\begin{lstlisting}[style=bashVerb, breaklines=true, escapechar=~]
Error[E4024] : colliding method definitions test6::Foo::self and test6::Foo::self
 --> test_resources/default_ctor/test6.yr:(6,8)
 6  ~\lstbox{\textSFxi{}}~    pub self (x: i32) with x = x + 1 {}
    ~\lstbox{\textSFv{}}~        ^^^^
 7  ~\lstbox{\textSFxi{}}~    pub self;
    ~\lstbox{\textSFv{}}~        ----
    ~\lstbox{\textSFxi{}}~
    = when validating -- (test_resources/default_ctor/test6.yr:(3,7))
    ~\lstbox{\textSFii{}}~~\lstbox{\textSFx{}}~~\lstbox{\textSFx{}}~~\lstbox{\textSFx{}}~~\lstbox{\textSFx{}}~~\lstbox{\textSFx{}}~~\lstbox{\textSFx{}}~~\lstbox{\textSFx{}}~
Error[E4104] : symbol test6::Foo has errors
 --> test_resources/default_ctor/test6.yr:(11,18)
11  ~\lstbox{\textSFxi{}}~     let a = copy Foo (1);
    ~\lstbox{\textSFv{}}~                  ^^^
 --> test_resources/default_ctor/test6.yr:(6,8)
 6  ~\lstbox{\textSFxi{}}~    pub self (x: i32) with x = x + 1 {}
    ~\lstbox{\textSFv{}}~        ----  error: colliding method definitions test6::Foo::self and test6::Foo::self
    ~\lstbox{\textSFxi{}}~
    = when validating -- (test_resources/default_ctor/test6.yr:(3,7))
    = when validating test6::main -- (test_resources/default_ctor/test6.yr:(10,4))
    ~\lstbox{\textSFii{}}~~\lstbox{\textSFx{}}~~\lstbox{\textSFx{}}~~\lstbox{\textSFx{}}~~\lstbox{\textSFx{}}~~\lstbox{\textSFx{}}~~\lstbox{\textSFx{}}~~\lstbox{\textSFx{}}~
Error[E4224] : undefined operator . for types error and x
 --> test_resources/default_ctor/test6.yr:(12,6)
12  ~\lstbox{\textSFxi{}}~     a.x;
    ~\lstbox{\textSFv{}}~      ^
    ~\lstbox{\textSFxi{}}~
    = when validating test6::main -- (test_resources/default_ctor/test6.yr:(10,4))
    ~\lstbox{\textSFii{}}~~\lstbox{\textSFx{}}~~\lstbox{\textSFx{}}~~\lstbox{\textSFx{}}~~\lstbox{\textSFx{}}~~\lstbox{\textSFx{}}~~\lstbox{\textSFx{}}~~\lstbox{\textSFx{}}~
\end{lstlisting}

\subsubsection*{How to fix it}

Rename one of them, or change the parameters.

\errorcode{E4026}{cannot construct an instance of class \errhole{} that is abstract}
\label{sec:codes:e4026}

Raised during \textbf{validation}, from \texttt{ValidateErrorMessage::\allowbreak{}CONSTRUCT\_\allowbreak{}ABSTRACT\_\allowbreak{}CLASS} (\textit{src/\allowbreak{}ymirc/\allowbreak{}semantic/\allowbreak{}validator/\allowbreak{}errors.yr}).

\subsubsection*{What it means}

An instance of an abstract class was constructed. An abstract class has methods with no body, so an instance of it would have nothing to call.

\subsubsection*{Example}

\textit{test\_\allowbreak{}resources/\allowbreak{}diagnostics/\allowbreak{}test134.yr}:

%% check: skip
\begin{lstlisting}[style=coloredverbatimError]
// an abstract class constructed
@abstract
class Shape {
    pub self () {}
}

fn main () {
    let _ = copy Shape ();
}
\end{lstlisting}

\begin{lstlisting}[style=bashVerb, breaklines=true, escapechar=~]
Error[E4026] : cannot construct an instance of class test134::Shape that is abstract
 --> test_resources/diagnostics/test134.yr:(8,18)
 8  ~\lstbox{\textSFxi{}}~     let _ = copy Shape ();
    ~\lstbox{\textSFv{}}~                  ^^^^^
    ~\lstbox{\textSFxi{}}~
    = when validating test134::main -- (test_resources/diagnostics/test134.yr:(7,4))
    ~\lstbox{\textSFii{}}~~\lstbox{\textSFx{}}~~\lstbox{\textSFx{}}~~\lstbox{\textSFx{}}~~\lstbox{\textSFx{}}~~\lstbox{\textSFx{}}~~\lstbox{\textSFx{}}~~\lstbox{\textSFx{}}~
\end{lstlisting}

\subsubsection*{How to fix it}

Construct a derived class that implements the empty methods, or give the class bodies for them and drop \tokennolst{abstract}.

\errorcode{E4027}{method \errhole{} is defined as constant}
\label{sec:codes:e4027}

Raised during \textbf{validation}, from \texttt{ValidateErrorMessage::\allowbreak{}CONST\_\allowbreak{}METHOD} (\textit{src/\allowbreak{}ymirc/\allowbreak{}semantic/\allowbreak{}validator/\allowbreak{}errors.yr}).

\subsubsection*{What it means}

A method declared constant is used where a mutable one is needed -- typically to write through the instance it is called on.

\subsubsection*{Example}

\textit{test\_\allowbreak{}resources/\allowbreak{}lit\_\allowbreak{}class/\allowbreak{}operators/\allowbreak{}test15.yr}:

%% check: skip
\begin{lstlisting}[style=coloredverbatimError]
class A {

    pub self () {}

    pub fn opIndex (self, _ : i32) {}

}

fn main () {
    let dmut a = copy A ();
    let _ = a:[0];
    let _ = a:.opIndex (0);
}
\end{lstlisting}

\begin{lstlisting}[style=bashVerb, breaklines=true, escapechar=~]
Error[E4027] : method test15::A::opIndex (_: i32)-> void is defined as constant
 --> test_resources/lit_class/operators/test15.yr:(5,12)
 5  ~\lstbox{\textSFxi{}}~     pub fn opIndex (self, _ : i32) {}
    ~\lstbox{\textSFv{}}~            ^^^^^^^
 --> test_resources/lit_class/operators/test15.yr:(11,14)
11  ~\lstbox{\textSFxi{}}~     let _ = a:[0];
    ~\lstbox{\textSFv{}}~              --  note: aliasing the value of type mut &(mut test15::A) to create a constant borrowing is prohibited
    ~\lstbox{\textSFxi{}}~
    = when validating test15::main -- (test_resources/lit_class/operators/test15.yr:(9,4))
    ~\lstbox{\textSFii{}}~~\lstbox{\textSFx{}}~~\lstbox{\textSFx{}}~~\lstbox{\textSFx{}}~~\lstbox{\textSFx{}}~~\lstbox{\textSFx{}}~~\lstbox{\textSFx{}}~~\lstbox{\textSFx{}}~
\end{lstlisting}

\subsubsection*{How to fix it}

Make the method mutable (\tokennolst{self} declared \tokennolst{mut}), or call a mutable method instead.

\errorcode{E4035}{infinite constructor redirection calls when calling \errhole{}}
\label{sec:codes:e4035}

Raised during \textbf{validation}, from \texttt{ValidateErrorMessage::\allowbreak{}CTOR\_\allowbreak{}INFINITE\_\allowbreak{}REDIRECTION} (\textit{src/\allowbreak{}ymirc/\allowbreak{}semantic/\allowbreak{}validator/\allowbreak{}errors.yr}).

\subsubsection*{What it means}

A chain of constructor redirections comes back to a constructor already in the chain, so construction would never terminate.

\subsubsection*{Example}

\textit{test\_\allowbreak{}resources/\allowbreak{}templates/\allowbreak{}class/\allowbreak{}test3.yr}:

%% check: skip
\begin{lstlisting}[style=coloredverbatimError]
class A {
    pub self (a : i32)
        with self (a)
    {}
}
\end{lstlisting}

\begin{lstlisting}[style=bashVerb, breaklines=true, escapechar=~]
Error[E4035] : infinite constructor redirection calls when calling test3::A::self (a: i32)-> mut &(mut test3::A)
 --> test_resources/templates/class/test3.yr:(3,14)
 3  ~\lstbox{\textSFxi{}}~         with self (a)
    ~\lstbox{\textSFv{}}~              ^^^^
    ~\lstbox{\textSFii{}}~~\lstbox{\textSFx{}}~~\lstbox{\textSFx{}}~~\lstbox{\textSFx{}}~~\lstbox{\textSFx{}}~~\lstbox{\textSFx{}}~~\lstbox{\textSFx{}}~~\lstbox{\textSFx{}}~
\end{lstlisting}

\subsubsection*{How to fix it}

Break the cycle: one constructor of the chain has to initialize the fields itself instead of redirecting.

\errorcode{E4044}{field must have a referenceable name}
\label{sec:codes:e4044}

Raised during \textbf{validation}, from \texttt{ValidateErrorMessage::\allowbreak{}CLASS\_\allowbreak{}FIELD\_\allowbreak{}NO\_\allowbreak{}NAME} (\textit{src/\allowbreak{}ymirc/\allowbreak{}semantic/\allowbreak{}validator/\allowbreak{}errors.yr}).

\subsubsection*{What it means}

A field was declared without a name that can be referred to.

\subsubsection*{Example}

\textit{test\_\allowbreak{}resources/\allowbreak{}diagnostics/\allowbreak{}test100.yr}:

%% check: skip
\begin{lstlisting}[style=coloredverbatimError]
// a class field named _
class Point {
    let _ : i32 = 1;
    pub self () {}
}

fn main () {
    let _ = copy Point ();
}
\end{lstlisting}

\begin{lstlisting}[style=bashVerb, breaklines=true, escapechar=~]
Error[E4044] : field must have a referenceable name
 --> test_resources/diagnostics/test100.yr:(3,9)
 3  ~\lstbox{\textSFxi{}}~     let _ : i32 = 1;
    ~\lstbox{\textSFv{}}~         ^         -
    ~\lstbox{\textSFxi{}}~
    = when validating -- (test_resources/diagnostics/test100.yr:(2,7))
    ~\lstbox{\textSFii{}}~~\lstbox{\textSFx{}}~~\lstbox{\textSFx{}}~~\lstbox{\textSFx{}}~~\lstbox{\textSFx{}}~~\lstbox{\textSFx{}}~~\lstbox{\textSFx{}}~~\lstbox{\textSFx{}}~
Error[E4104] : symbol test100::Point has errors
 --> test_resources/diagnostics/test100.yr:(8,18)
 8  ~\lstbox{\textSFxi{}}~     let _ = copy Point ();
    ~\lstbox{\textSFv{}}~                  ^^^^^
 --> test_resources/diagnostics/test100.yr:(3,9)
 3  ~\lstbox{\textSFxi{}}~     let _ : i32 = 1;
    ~\lstbox{\textSFv{}}~         -         -  error: field must have a referenceable name
    ~\lstbox{\textSFxi{}}~
    = when validating -- (test_resources/diagnostics/test100.yr:(2,7))
    = when validating test100::main -- (test_resources/diagnostics/test100.yr:(7,4))
    ~\lstbox{\textSFii{}}~~\lstbox{\textSFx{}}~~\lstbox{\textSFx{}}~~\lstbox{\textSFx{}}~~\lstbox{\textSFx{}}~~\lstbox{\textSFx{}}~~\lstbox{\textSFx{}}~~\lstbox{\textSFx{}}~
\end{lstlisting}

\subsubsection*{How to fix it}

Give the field a name.

\errorcode{E4051}{cannot perform a direct call of the method \errhole{} that is empty}
\label{sec:codes:e4051}

Raised during \textbf{validation}, from \texttt{ValidateErrorMessage::\allowbreak{}EMPTY\_\allowbreak{}METHOD\_\allowbreak{}CALL} (\textit{src/\allowbreak{}ymirc/\allowbreak{}semantic/\allowbreak{}validator/\allowbreak{}errors.yr}).

\subsubsection*{What it means}

A method with no body was called directly rather than through the virtual dispatch of an instance. An empty method has no code to jump to.

\subsubsection*{Example}

\textit{test\_\allowbreak{}resources/\allowbreak{}diagnostics/\allowbreak{}test30.yr}:

%% check: skip
\begin{lstlisting}[style=coloredverbatimError]
@abstract
class Foo {
    pub self () {}
    pub fn bar (self);
}

class Bar over Foo {
    pub self () {}
    pub over bar (self) { self.super.bar (); }
}

fn main () { let _ = copy Bar (); }
\end{lstlisting}

\begin{lstlisting}[style=bashVerb, breaklines=true, escapechar=~]
Note : when using uniform function call syntax
 --> test_resources/diagnostics/test30.yr:(9,38)
 9  ~\lstbox{\textSFxi{}}~     pub over bar (self) { self.super.bar (); }
    ~\lstbox{\textSFv{}}~              ---                     ^^^
    ~\lstbox{\textSFxi{}}~                ~\lstbox{\textSFii{}}~~\lstbox{\textSFx{}}~ note: defined here
    ~\lstbox{\textSFxi{}}~                                        ~\lstbox{\textSFii{}}~~\lstbox{\textSFx{}}~ error: test30::Bar::bar is a method and must be called from an instance
    ~\lstbox{\textSFxi{}}~
    = when validating test30::Bar::bar -- (test_resources/diagnostics/test30.yr:(9,14))
    = when validating -- (test_resources/diagnostics/test30.yr:(7,7))
    ~\lstbox{\textSFii{}}~~\lstbox{\textSFx{}}~~\lstbox{\textSFx{}}~~\lstbox{\textSFx{}}~~\lstbox{\textSFx{}}~~\lstbox{\textSFx{}}~~\lstbox{\textSFx{}}~~\lstbox{\textSFx{}}~
Error[E4019] : class &(test30::Foo) has no field named bar
 --> test_resources/diagnostics/test30.yr:(9,38)
 9  ~\lstbox{\textSFxi{}}~     pub over bar (self) { self.super.bar (); }
    ~\lstbox{\textSFv{}}~                                ----- ^^^
    ~\lstbox{\textSFxi{}}~                                    ~\lstbox{\textSFii{}}~~\lstbox{\textSFx{}}~ note: from an proxied instance of class &(test30::Foo) that is abstract
    ~\lstbox{\textSFxi{}}~                                        ~\lstbox{\textSFii{}}~~\lstbox{\textSFx{}}~ note: class &(test30::Foo) has no field access method named bar
 --> test_resources/diagnostics/test30.yr:(4,12)
 4  ~\lstbox{\textSFxi{}}~     pub fn bar (self);
    ~\lstbox{\textSFv{}}~            ---  error: cannot perform a direct call of the method test30::Foo::bar ()-> void that is empty
    ~\lstbox{\textSFxi{}}~
    = when validating test30::Bar::bar -- (test_resources/diagnostics/test30.yr:(9,14))
    = when validating -- (test_resources/diagnostics/test30.yr:(7,7))
    ~\lstbox{\textSFii{}}~~\lstbox{\textSFx{}}~~\lstbox{\textSFx{}}~~\lstbox{\textSFx{}}~~\lstbox{\textSFx{}}~~\lstbox{\textSFx{}}~~\lstbox{\textSFx{}}~~\lstbox{\textSFx{}}~
\end{lstlisting}

\subsubsection*{How to fix it}

Call it on an instance whose dynamic type provides a body, or give the method a body.

\errorcode{E4060}{field access method must return a value}
\label{sec:codes:e4060}

Raised during \textbf{validation}, from \texttt{ValidateErrorMessage::\allowbreak{}FIELD\_\allowbreak{}METHOD\_\allowbreak{}ACCESS\_\allowbreak{}NO\_\allowbreak{}RETURN} (\textit{src/\allowbreak{}ymirc/\allowbreak{}semantic/\allowbreak{}validator/\allowbreak{}errors.yr}).

\subsubsection*{What it means}

A field access method has no return value. A field access method stands in for reading a field, so it has to produce one.

\subsubsection*{Example}

\textit{test\_\allowbreak{}resources/\allowbreak{}diagnostics/\allowbreak{}test34.yr}:

%% check: skip
\begin{lstlisting}[style=coloredverbatimError]
class Foo {
    pub self () {}
    @field pub fn bar (self) {}
}

fn main () { let f = copy Foo (); f.bar; }
\end{lstlisting}

\begin{lstlisting}[style=bashVerb, breaklines=true, escapechar=~]
Error[E4060] : field access method must return a value
 --> test_resources/diagnostics/test34.yr:(3,16)
 3  ~\lstbox{\textSFxi{}}~     @field pub fn bar (self) {}
    ~\lstbox{\textSFv{}}~      -----     ^^
    ~\lstbox{\textSFxi{}}~
    = when validating test34::Foo::bar -- (test_resources/diagnostics/test34.yr:(3,19))
    = when validating -- (test_resources/diagnostics/test34.yr:(1,7))
    ~\lstbox{\textSFii{}}~~\lstbox{\textSFx{}}~~\lstbox{\textSFx{}}~~\lstbox{\textSFx{}}~~\lstbox{\textSFx{}}~~\lstbox{\textSFx{}}~~\lstbox{\textSFx{}}~~\lstbox{\textSFx{}}~
Error[E4104] : symbol test34::Foo has errors
 --> test_resources/diagnostics/test34.yr:(6,27)
 6  ~\lstbox{\textSFxi{}}~ fn main () { let f = copy Foo (); f.bar; }
    ~\lstbox{\textSFv{}}~                           ^^^
 --> test_resources/diagnostics/test34.yr:(3,6)
 3  ~\lstbox{\textSFxi{}}~     @field pub fn bar (self) {}
    ~\lstbox{\textSFv{}}~      -----     --  error: field access method must return a value
    ~\lstbox{\textSFxi{}}~
    = when validating test34::Foo::bar -- (test_resources/diagnostics/test34.yr:(3,19))
    = when validating -- (test_resources/diagnostics/test34.yr:(1,7))
    = when validating test34::main -- (test_resources/diagnostics/test34.yr:(6,4))
    ~\lstbox{\textSFii{}}~~\lstbox{\textSFx{}}~~\lstbox{\textSFx{}}~~\lstbox{\textSFx{}}~~\lstbox{\textSFx{}}~~\lstbox{\textSFx{}}~~\lstbox{\textSFx{}}~~\lstbox{\textSFx{}}~
Error[E4224] : undefined operator . for types error and bar
 --> test_resources/diagnostics/test34.yr:(6,36)
 6  ~\lstbox{\textSFxi{}}~ fn main () { let f = copy Foo (); f.bar; }
    ~\lstbox{\textSFv{}}~                                    ^
    ~\lstbox{\textSFxi{}}~
    = when validating test34::main -- (test_resources/diagnostics/test34.yr:(6,4))
    ~\lstbox{\textSFii{}}~~\lstbox{\textSFx{}}~~\lstbox{\textSFx{}}~~\lstbox{\textSFx{}}~~\lstbox{\textSFx{}}~~\lstbox{\textSFx{}}~~\lstbox{\textSFx{}}~~\lstbox{\textSFx{}}~
\end{lstlisting}

\subsubsection*{How to fix it}

Return the value the pseudo-field stands for.

\errorcode{E4062}{field assign method must be mutable}
\label{sec:codes:e4062}

Raised during \textbf{validation}, from \texttt{ValidateErrorMessage::\allowbreak{}FIELD\_\allowbreak{}METHOD\_\allowbreak{}ASSIGN\_\allowbreak{}CONST} (\textit{src/\allowbreak{}ymirc/\allowbreak{}semantic/\allowbreak{}validator/\allowbreak{}errors.yr}).

\subsubsection*{What it means}

A field assign method is declared constant. It stands in for writing a field, which modifies the instance.

\subsubsection*{Example}

\textit{test\_\allowbreak{}resources/\allowbreak{}diagnostics/\allowbreak{}test35.yr}:

%% check: skip
\begin{lstlisting}[style=coloredverbatimError]
class Foo {
    let mut x: i32 = 0;
    pub self () {}
    @field pub fn bar (self, v: i32) { self.x = v; }
}

fn main () { let dmut f = copy Foo (); f:.bar = 1; }
\end{lstlisting}

\begin{lstlisting}[style=bashVerb, breaklines=true, escapechar=~]
Error[E4062] : field assign method must be mutable
 --> test_resources/diagnostics/test35.yr:(4,16)
 4  ~\lstbox{\textSFxi{}}~     @field pub fn bar (self, v: i32) { self.x = v; }
    ~\lstbox{\textSFv{}}~      -----     ^^
    ~\lstbox{\textSFxi{}}~
    = when validating test35::Foo::bar -- (test_resources/diagnostics/test35.yr:(4,19))
    = when validating -- (test_resources/diagnostics/test35.yr:(1,7))
    ~\lstbox{\textSFii{}}~~\lstbox{\textSFx{}}~~\lstbox{\textSFx{}}~~\lstbox{\textSFx{}}~~\lstbox{\textSFx{}}~~\lstbox{\textSFx{}}~~\lstbox{\textSFx{}}~~\lstbox{\textSFx{}}~
Error[E4104] : symbol test35::Foo has errors
 --> test_resources/diagnostics/test35.yr:(7,32)
 7  ~\lstbox{\textSFxi{}}~ fn main () { let dmut f = copy Foo (); f:.bar = 1; }
    ~\lstbox{\textSFv{}}~                                ^^^
 --> test_resources/diagnostics/test35.yr:(4,6)
 4  ~\lstbox{\textSFxi{}}~     @field pub fn bar (self, v: i32) { self.x = v; }
    ~\lstbox{\textSFv{}}~      -----     --  error: field assign method must be mutable
    ~\lstbox{\textSFxi{}}~
    = when validating test35::Foo::bar -- (test_resources/diagnostics/test35.yr:(4,19))
    = when validating -- (test_resources/diagnostics/test35.yr:(1,7))
    = when validating test35::main -- (test_resources/diagnostics/test35.yr:(7,4))
    ~\lstbox{\textSFii{}}~~\lstbox{\textSFx{}}~~\lstbox{\textSFx{}}~~\lstbox{\textSFx{}}~~\lstbox{\textSFx{}}~~\lstbox{\textSFx{}}~~\lstbox{\textSFx{}}~~\lstbox{\textSFx{}}~
Error[E4015] : alias operator :. is only usable on class, record, entity, map or generator types, not f
 --> test_resources/diagnostics/test35.yr:(7,40)
 7  ~\lstbox{\textSFxi{}}~ fn main () { let dmut f = copy Foo (); f:.bar = 1; }
    ~\lstbox{\textSFv{}}~                                        ^
    ~\lstbox{\textSFxi{}}~
    = when validating test35::main -- (test_resources/diagnostics/test35.yr:(7,4))
    ~\lstbox{\textSFii{}}~~\lstbox{\textSFx{}}~~\lstbox{\textSFx{}}~~\lstbox{\textSFx{}}~~\lstbox{\textSFx{}}~~\lstbox{\textSFx{}}~~\lstbox{\textSFx{}}~~\lstbox{\textSFx{}}~
\end{lstlisting}

\subsubsection*{How to fix it}

Declare \tokennolst{self} mutable on the method.

\errorcode{E4063}{field assign method must not return a value}
\label{sec:codes:e4063}

Raised during \textbf{validation}, from \texttt{ValidateErrorMessage::\allowbreak{}FIELD\_\allowbreak{}METHOD\_\allowbreak{}ASSIGN\_\allowbreak{}WITH\_\allowbreak{}RETURN} (\textit{src/\allowbreak{}ymirc/\allowbreak{}semantic/\allowbreak{}validator/\allowbreak{}errors.yr}).

\subsubsection*{What it means}

A field assign method returns a value. It stands in for an assignment, which produces nothing.

\subsubsection*{Example}

\textit{test\_\allowbreak{}resources/\allowbreak{}diagnostics/\allowbreak{}test36.yr}:

%% check: skip
\begin{lstlisting}[style=coloredverbatimError]
class Foo {
    let mut x: i32 = 0;
    pub self () {}
    @field pub fn bar (mut self, v: i32)-> i32 { self.x = v; v }
}

fn main () { let dmut f = copy Foo (); f:.bar = 1; }
\end{lstlisting}

\begin{lstlisting}[style=bashVerb, breaklines=true, escapechar=~]
Error[E4063] : field assign method must not return a value
 --> test_resources/diagnostics/test36.yr:(4,16)
 4  ~\lstbox{\textSFxi{}}~     @field pub fn bar (mut self, v: i32)-> i32 { self.x = v; v }
    ~\lstbox{\textSFv{}}~      -----     ^^
    ~\lstbox{\textSFxi{}}~
    = when validating test36::Foo::bar -- (test_resources/diagnostics/test36.yr:(4,19))
    = when validating -- (test_resources/diagnostics/test36.yr:(1,7))
    ~\lstbox{\textSFii{}}~~\lstbox{\textSFx{}}~~\lstbox{\textSFx{}}~~\lstbox{\textSFx{}}~~\lstbox{\textSFx{}}~~\lstbox{\textSFx{}}~~\lstbox{\textSFx{}}~~\lstbox{\textSFx{}}~
Error[E4104] : symbol test36::Foo has errors
 --> test_resources/diagnostics/test36.yr:(7,32)
 7  ~\lstbox{\textSFxi{}}~ fn main () { let dmut f = copy Foo (); f:.bar = 1; }
    ~\lstbox{\textSFv{}}~                                ^^^
 --> test_resources/diagnostics/test36.yr:(4,6)
 4  ~\lstbox{\textSFxi{}}~     @field pub fn bar (mut self, v: i32)-> i32 { self.x = v; v }
    ~\lstbox{\textSFv{}}~      -----     --  error: field assign method must not return a value
    ~\lstbox{\textSFxi{}}~
    = when validating test36::Foo::bar -- (test_resources/diagnostics/test36.yr:(4,19))
    = when validating -- (test_resources/diagnostics/test36.yr:(1,7))
    = when validating test36::main -- (test_resources/diagnostics/test36.yr:(7,4))
    ~\lstbox{\textSFii{}}~~\lstbox{\textSFx{}}~~\lstbox{\textSFx{}}~~\lstbox{\textSFx{}}~~\lstbox{\textSFx{}}~~\lstbox{\textSFx{}}~~\lstbox{\textSFx{}}~~\lstbox{\textSFx{}}~
Error[E4015] : alias operator :. is only usable on class, record, entity, map or generator types, not f
 --> test_resources/diagnostics/test36.yr:(7,40)
 7  ~\lstbox{\textSFxi{}}~ fn main () { let dmut f = copy Foo (); f:.bar = 1; }
    ~\lstbox{\textSFv{}}~                                        ^
    ~\lstbox{\textSFxi{}}~
    = when validating test36::main -- (test_resources/diagnostics/test36.yr:(7,4))
    ~\lstbox{\textSFii{}}~~\lstbox{\textSFx{}}~~\lstbox{\textSFx{}}~~\lstbox{\textSFx{}}~~\lstbox{\textSFx{}}~~\lstbox{\textSFx{}}~~\lstbox{\textSFx{}}~~\lstbox{\textSFx{}}~
\end{lstlisting}

\subsubsection*{How to fix it}

Remove the return value.

\errorcode{E4065}{field method \errhole{} is mutable}
\label{sec:codes:e4065}

Raised during \textbf{validation}, from \texttt{ValidateErrorMessage::\allowbreak{}FIELD\_\allowbreak{}METHOD\_\allowbreak{}MUTABLE\_\allowbreak{}ALIAS} (\textit{src/\allowbreak{}ymirc/\allowbreak{}semantic/\allowbreak{}validator/\allowbreak{}errors.yr}).

\subsubsection*{What it means}

A field method is mutable and is being used where a constant one is required -- typically through a value borrowed as constant.

\subsubsection*{Example}

\textit{test\_\allowbreak{}resources/\allowbreak{}lit\_\allowbreak{}class/\allowbreak{}methods/\allowbreak{}test8.yr}:

%% check: skip
\begin{lstlisting}[style=coloredverbatimError]
record A {
    let mut _x : i32 = 0;

    pub self () {}

    @field
    pub fn foo (self)-> i32 {
        self._x
    }

    @field
    pub fn foo (mut self, i : i32) {
        self._x = i;
    }

}


fn main () {
    let a = A ();
    a.foo += 1;
}
\end{lstlisting}

\begin{lstlisting}[style=bashVerb, breaklines=true, escapechar=~]
Error[E4065] : field method test8::A::foo (i: i32)-> void is mutable
 --> test_resources/lit_class/methods/test8.yr:(12,12)
12  ~\lstbox{\textSFxi{}}~     pub fn foo (mut self, i : i32) {
    ~\lstbox{\textSFv{}}~            ^^^
    ~\lstbox{\textSFxi{}}~ Error[E4082] : left operand of type test8::A is immutable
    ~\lstbox{\textSFxi{}}~  --> test_resources/lit_class/methods/test8.yr:(21,5)
    ~\lstbox{\textSFxi{}}~ 21  ~\lstbox{\textSFxi{}}~     a.foo += 1;
    ~\lstbox{\textSFxi{}}~     ~\lstbox{\textSFv{}}~     ^^
    ~\lstbox{\textSFxi{}}~     ~\lstbox{\textSFii{}}~~\lstbox{\textSFx{}}~~\lstbox{\textSFx{}}~~\lstbox{\textSFx{}}~~\lstbox{\textSFx{}}~~\lstbox{\textSFx{}}~~\lstbox{\textSFx{}}~~\lstbox{\textSFx{}}~
    ~\lstbox{\textSFii{}}~~\lstbox{\textSFx{}}~~\lstbox{\textSFx{}}~~\lstbox{\textSFx{}}~~\lstbox{\textSFx{}}~~\lstbox{\textSFx{}}~~\lstbox{\textSFvii{}}~~\lstbox{\textSFx{}}~
\end{lstlisting}

\subsubsection*{How to fix it}

Use a constant field method for reading, and reserve the mutable one for writing.

\errorcode{E4066}{field method cannot have more than one parameter}
\label{sec:codes:e4066}

Raised during \textbf{validation}, from \texttt{ValidateErrorMessage::\allowbreak{}FIELD\_\allowbreak{}METHOD\_\allowbreak{}PARAMS} (\textit{src/\allowbreak{}ymirc/\allowbreak{}semantic/\allowbreak{}validator/\allowbreak{}errors.yr}).

\subsubsection*{What it means}

A field method declares more than one parameter besides \tokennolst{self}. A field access method takes none and a field assign method takes exactly the value being assigned.

\subsubsection*{Example}

\textit{test\_\allowbreak{}resources/\allowbreak{}diagnostics/\allowbreak{}test37.yr}:

%% check: skip
\begin{lstlisting}[style=coloredverbatimError]
class Foo {
    pub self () {}
    @field pub fn bar (self, a: i32, b: i32)-> i32 { a + b }
}

fn main () { let f = copy Foo (); f.bar; }
\end{lstlisting}

\begin{lstlisting}[style=bashVerb, breaklines=true, escapechar=~]
Error[E4066] : field method cannot have more than one parameter
 --> test_resources/diagnostics/test37.yr:(3,16)
 3  ~\lstbox{\textSFxi{}}~     @field pub fn bar (self, a: i32, b: i32)-> i32 { a + b }
    ~\lstbox{\textSFv{}}~      -----     ^^
    ~\lstbox{\textSFxi{}}~
    = when validating test37::Foo::bar -- (test_resources/diagnostics/test37.yr:(3,19))
    = when validating -- (test_resources/diagnostics/test37.yr:(1,7))
    ~\lstbox{\textSFii{}}~~\lstbox{\textSFx{}}~~\lstbox{\textSFx{}}~~\lstbox{\textSFx{}}~~\lstbox{\textSFx{}}~~\lstbox{\textSFx{}}~~\lstbox{\textSFx{}}~~\lstbox{\textSFx{}}~
Error[E4104] : symbol test37::Foo has errors
 --> test_resources/diagnostics/test37.yr:(6,27)
 6  ~\lstbox{\textSFxi{}}~ fn main () { let f = copy Foo (); f.bar; }
    ~\lstbox{\textSFv{}}~                           ^^^
 --> test_resources/diagnostics/test37.yr:(3,6)
 3  ~\lstbox{\textSFxi{}}~     @field pub fn bar (self, a: i32, b: i32)-> i32 { a + b }
    ~\lstbox{\textSFv{}}~      -----     --  error: field method cannot have more than one parameter
    ~\lstbox{\textSFxi{}}~
    = when validating test37::Foo::bar -- (test_resources/diagnostics/test37.yr:(3,19))
    = when validating -- (test_resources/diagnostics/test37.yr:(1,7))
    = when validating test37::main -- (test_resources/diagnostics/test37.yr:(6,4))
    ~\lstbox{\textSFii{}}~~\lstbox{\textSFx{}}~~\lstbox{\textSFx{}}~~\lstbox{\textSFx{}}~~\lstbox{\textSFx{}}~~\lstbox{\textSFx{}}~~\lstbox{\textSFx{}}~~\lstbox{\textSFx{}}~
Error[E4224] : undefined operator . for types error and bar
 --> test_resources/diagnostics/test37.yr:(6,36)
 6  ~\lstbox{\textSFxi{}}~ fn main () { let f = copy Foo (); f.bar; }
    ~\lstbox{\textSFv{}}~                                    ^
    ~\lstbox{\textSFxi{}}~
    = when validating test37::main -- (test_resources/diagnostics/test37.yr:(6,4))
    ~\lstbox{\textSFii{}}~~\lstbox{\textSFx{}}~~\lstbox{\textSFx{}}~~\lstbox{\textSFx{}}~~\lstbox{\textSFx{}}~~\lstbox{\textSFx{}}~~\lstbox{\textSFx{}}~~\lstbox{\textSFx{}}~
\end{lstlisting}

\subsubsection*{How to fix it}

Reduce the parameter list. A method genuinely needing several arguments is an ordinary method, not a field one.

\errorcode{E4079}{the field \errhole{} is not public, and has no initial value to be constructed from}
\label{sec:codes:e4079}

Raised during \textbf{validation}, from \texttt{ValidateErrorMessage::\allowbreak{}GEN\_\allowbreak{}CTOR\_\allowbreak{}HIDDEN\_\allowbreak{}FIELD\_\allowbreak{}NO\_\allowbreak{}VALUE} (\textit{src/\allowbreak{}ymirc/\allowbreak{}semantic/\allowbreak{}validator/\allowbreak{}errors.yr}).

\subsubsection*{What it means}

The constructor built from the fields cannot initialize the reported field: it is not public, so a caller could not pass it, and it has no initial value to fall back on.

\subsubsection*{Example}

\textit{test\_\allowbreak{}resources/\allowbreak{}data\_\allowbreak{}record/\allowbreak{}ctors/\allowbreak{}test3.yr}:

%% check: skip
\begin{lstlisting}[style=coloredverbatimError]
@data
record Foo {
    prv let _x: i32;
    let y: i32;
}

fn main() {
    let a = Foo(1);
    a.y;
}
\end{lstlisting}

\begin{lstlisting}[style=bashVerb, breaklines=true, escapechar=~]
Error[E4079] : the field _x is not public, and has no initial value to be constructed from
 --> test_resources/data_record/ctors/test3.yr:(3,13)
 3  ~\lstbox{\textSFxi{}}~     prv let _x: i32;
    ~\lstbox{\textSFv{}}~             ^^
    ~\lstbox{\textSFxi{}}~
    = when validating -- (test_resources/data_record/ctors/test3.yr:(2,8))
    ~\lstbox{\textSFii{}}~~\lstbox{\textSFx{}}~~\lstbox{\textSFx{}}~~\lstbox{\textSFx{}}~~\lstbox{\textSFx{}}~~\lstbox{\textSFx{}}~~\lstbox{\textSFx{}}~~\lstbox{\textSFx{}}~
Error[E4104] : symbol test3::Foo has errors
 --> test_resources/data_record/ctors/test3.yr:(8,13)
 8  ~\lstbox{\textSFxi{}}~     let a = Foo(1);
    ~\lstbox{\textSFv{}}~             ^^^
 --> test_resources/data_record/ctors/test3.yr:(3,13)
 3  ~\lstbox{\textSFxi{}}~     prv let _x: i32;
    ~\lstbox{\textSFv{}}~             --  error: the field _x is not public, and has no initial value to be constructed from
    ~\lstbox{\textSFxi{}}~
    = when validating -- (test_resources/data_record/ctors/test3.yr:(2,8))
    = when validating test3::main -- (test_resources/data_record/ctors/test3.yr:(7,4))
    ~\lstbox{\textSFii{}}~~\lstbox{\textSFx{}}~~\lstbox{\textSFx{}}~~\lstbox{\textSFx{}}~~\lstbox{\textSFx{}}~~\lstbox{\textSFx{}}~~\lstbox{\textSFx{}}~~\lstbox{\textSFx{}}~
Error[E4224] : undefined operator . for types error and y
 --> test_resources/data_record/ctors/test3.yr:(9,6)
 9  ~\lstbox{\textSFxi{}}~     a.y;
    ~\lstbox{\textSFv{}}~      ^
    ~\lstbox{\textSFxi{}}~
    = when validating test3::main -- (test_resources/data_record/ctors/test3.yr:(7,4))
    ~\lstbox{\textSFii{}}~~\lstbox{\textSFx{}}~~\lstbox{\textSFx{}}~~\lstbox{\textSFx{}}~~\lstbox{\textSFx{}}~~\lstbox{\textSFx{}}~~\lstbox{\textSFx{}}~~\lstbox{\textSFx{}}~
\end{lstlisting}

\subsubsection*{How to fix it}

Give the field an initial value, make it public, or write the constructor explicitly.

\errorcode{E4090}{method \errhole{} implicitly overrides method \errhole{}}
\label{sec:codes:e4090}

Raised during \textbf{validation}, from \texttt{ValidateErrorMessage::\allowbreak{}IMPLICIT\_\allowbreak{}OVERRIDE} (\textit{src/\allowbreak{}ymirc/\allowbreak{}semantic/\allowbreak{}validator/\allowbreak{}errors.yr}).

\subsubsection*{What it means}

A method has the same signature as one of an ancestor class but is not marked \tokennolst{override}. Ymir requires the marker so that adding a method to a base class cannot silently change what a derived class does.

\subsubsection*{Example}

\textit{test\_\allowbreak{}resources/\allowbreak{}diagnostics/\allowbreak{}test42.yr}:

%% check: skip
\begin{lstlisting}[style=coloredverbatimError]
class A {
    pub self () {}
    pub fn foo (self) {}
}

class B over A {
    pub self () {}
    pub fn foo (self) {}
}

fn main () { let _ = copy B (); }
\end{lstlisting}

\begin{lstlisting}[style=bashVerb, breaklines=true, escapechar=~]
Error[E4090] : method test42::B::foo ()-> void implicitly overrides method test42::A::foo ()-> void
 --> test_resources/diagnostics/test42.yr:(3,12)
 3  ~\lstbox{\textSFxi{}}~     pub fn foo (self) {}
    ~\lstbox{\textSFv{}}~            ^^^
 --> test_resources/diagnostics/test42.yr:(8,12)
 8  ~\lstbox{\textSFxi{}}~     pub fn foo (self) {}
    ~\lstbox{\textSFv{}}~            ---
    ~\lstbox{\textSFxi{}}~
    = when validating -- (test_resources/diagnostics/test42.yr:(6,7))
    ~\lstbox{\textSFii{}}~~\lstbox{\textSFx{}}~~\lstbox{\textSFx{}}~~\lstbox{\textSFx{}}~~\lstbox{\textSFx{}}~~\lstbox{\textSFx{}}~~\lstbox{\textSFx{}}~~\lstbox{\textSFx{}}~
Error[E4104] : symbol test42::B has errors
 --> test_resources/diagnostics/test42.yr:(11,27)
11  ~\lstbox{\textSFxi{}}~ fn main () { let _ = copy B (); }
    ~\lstbox{\textSFv{}}~                           ^
 --> test_resources/diagnostics/test42.yr:(3,12)
 3  ~\lstbox{\textSFxi{}}~     pub fn foo (self) {}
    ~\lstbox{\textSFv{}}~            ---  error: method test42::B::foo ()-> void implicitly overrides method test42::A::foo ()-> void
    ~\lstbox{\textSFxi{}}~
    = when validating -- (test_resources/diagnostics/test42.yr:(6,7))
    = when validating test42::main -- (test_resources/diagnostics/test42.yr:(11,4))
    ~\lstbox{\textSFii{}}~~\lstbox{\textSFx{}}~~\lstbox{\textSFx{}}~~\lstbox{\textSFx{}}~~\lstbox{\textSFx{}}~~\lstbox{\textSFx{}}~~\lstbox{\textSFx{}}~~\lstbox{\textSFx{}}~
\end{lstlisting}

\subsubsection*{How to fix it}

Add \tokennolst{override} if overriding was intended, or rename the method if it was meant to be a new one.

\errorcode{E4091}{implicit override of method \errhole{} by trait impl}
\label{sec:codes:e4091}

Raised during \textbf{validation}, from \texttt{ValidateErrorMessage::\allowbreak{}IMPLICIT\_\allowbreak{}OVERRIDE\_\allowbreak{}BY\_\allowbreak{}TRAIT} (\textit{src/\allowbreak{}ymirc/\allowbreak{}semantic/\allowbreak{}validator/\allowbreak{}errors.yr}).

\subsubsection*{What it means}

A trait \tokennolst{impl} provides a method that overrides an ancestor method, without the override being declared.

\subsubsection*{Example}

\textit{test\_\allowbreak{}resources/\allowbreak{}diagnostics/\allowbreak{}test107.yr}:

%% check: skip
\begin{lstlisting}[style=coloredverbatimError]
// a trait implemented by a class whose ancestor already has a method of that name
trait Shape {
    pub fn area (self)-> i32 { 1 }
}

class Base {
    pub self () {}
    pub fn area (self)-> i32 { 2 }
}

class Square over Base {
    pub self () with super () {}
    impl Shape;
}

fn main () {
    let _ = copy Square ();
}
\end{lstlisting}

\begin{lstlisting}[style=bashVerb, breaklines=true, escapechar=~]
Error[E4091] : implicit override of method test107::Base::area by trait impl
 --> test_resources/diagnostics/test107.yr:(13,5)
13  ~\lstbox{\textSFxi{}}~     impl Shape;
    ~\lstbox{\textSFv{}}~     ^^^^
 --> test_resources/diagnostics/test107.yr:(3,12)
 3  ~\lstbox{\textSFxi{}}~     pub fn area (self)-> i32 { 1 }
    ~\lstbox{\textSFv{}}~            ----
    ~\lstbox{\textSFxi{}}~
    = when validating -- (test_resources/diagnostics/test107.yr:(11,7))
    ~\lstbox{\textSFii{}}~~\lstbox{\textSFx{}}~~\lstbox{\textSFx{}}~~\lstbox{\textSFx{}}~~\lstbox{\textSFx{}}~~\lstbox{\textSFx{}}~~\lstbox{\textSFx{}}~~\lstbox{\textSFx{}}~
Error[E4104] : symbol test107::Square has errors
 --> test_resources/diagnostics/test107.yr:(17,18)
17  ~\lstbox{\textSFxi{}}~     let _ = copy Square ();
    ~\lstbox{\textSFv{}}~                  ^^^^^^
 --> test_resources/diagnostics/test107.yr:(3,12)
 3  ~\lstbox{\textSFxi{}}~     pub fn area (self)-> i32 { 1 }
    ~\lstbox{\textSFv{}}~            ----
   ...
13  ~\lstbox{\textSFxi{}}~     impl Shape;
    ~\lstbox{\textSFv{}}~     ----  error: implicit override of method test107::Base::area by trait impl
    ~\lstbox{\textSFxi{}}~
    = when validating -- (test_resources/diagnostics/test107.yr:(11,7))
    = when validating test107::main -- (test_resources/diagnostics/test107.yr:(16,4))
    ~\lstbox{\textSFii{}}~~\lstbox{\textSFx{}}~~\lstbox{\textSFx{}}~~\lstbox{\textSFx{}}~~\lstbox{\textSFx{}}~~\lstbox{\textSFx{}}~~\lstbox{\textSFx{}}~~\lstbox{\textSFx{}}~
\end{lstlisting}

\subsubsection*{How to fix it}

Mark the method \tokennolst{override} in the \tokennolst{impl}, or rename it.

\errorcode{E4093}{trait \errhole{} is implemented multiple times by \errhole{}}
\label{sec:codes:e4093}

Raised during \textbf{validation}, from \texttt{ValidateErrorMessage::\allowbreak{}IMPL\_\allowbreak{}MULTIPLE\_\allowbreak{}TIMES} (\textit{src/\allowbreak{}ymirc/\allowbreak{}semantic/\allowbreak{}validator/\allowbreak{}errors.yr}).

\subsubsection*{What it means}

The same trait is implemented twice by one type, so a call to one of its methods would have two candidates.

\subsubsection*{Example}

\textit{test\_\allowbreak{}resources/\allowbreak{}diagnostics/\allowbreak{}test43.yr}:

%% check: skip
\begin{lstlisting}[style=coloredverbatimError]
trait T { pub fn f (self); }

class Foo {
    pub self () {}
    impl T { pub over f (self) {} }
    impl T { pub over f (self) {} }
}

fn main () {}
\end{lstlisting}

\begin{lstlisting}[style=bashVerb, breaklines=true, escapechar=~]
Error[E4093] : trait test43::T is implemented multiple times by test43::Foo
 --> test_resources/diagnostics/test43.yr:(6,5)
 6  ~\lstbox{\textSFxi{}}~     impl T { pub over f (self) {} }
    ~\lstbox{\textSFv{}}~     ^^^^
    ~\lstbox{\textSFxi{}}~ Note : here
    ~\lstbox{\textSFxi{}}~  --> test_resources/diagnostics/test43.yr:(1,7)
    ~\lstbox{\textSFxi{}}~  1  ~\lstbox{\textSFxi{}}~ trait T { pub fn f (self); }
    ~\lstbox{\textSFxi{}}~     ~\lstbox{\textSFv{}}~       ^
    ~\lstbox{\textSFxi{}}~     ~\lstbox{\textSFii{}}~~\lstbox{\textSFx{}}~~\lstbox{\textSFx{}}~~\lstbox{\textSFx{}}~~\lstbox{\textSFx{}}~~\lstbox{\textSFx{}}~~\lstbox{\textSFx{}}~~\lstbox{\textSFx{}}~
    ~\lstbox{\textSFii{}}~~\lstbox{\textSFx{}}~~\lstbox{\textSFx{}}~~\lstbox{\textSFx{}}~~\lstbox{\textSFx{}}~~\lstbox{\textSFx{}}~~\lstbox{\textSFvii{}}~~\lstbox{\textSFx{}}~
\end{lstlisting}

\subsubsection*{How to fix it}

Merge the two \tokennolst{impl} blocks.

\errorcode{E4094}{impl statement must be followed by a trait, not \errhole{}}
\label{sec:codes:e4094}

Raised during \textbf{validation}, from \texttt{ValidateErrorMessage::\allowbreak{}IMPL\_\allowbreak{}NO\_\allowbreak{}TRAIT} (\textit{src/\allowbreak{}ymirc/\allowbreak{}semantic/\allowbreak{}validator/\allowbreak{}errors.yr}).

\subsubsection*{What it means}

An \tokennolst{impl} statement names something that is not a trait.

\subsubsection*{Example}

\textit{test\_\allowbreak{}resources/\allowbreak{}diagnostics/\allowbreak{}test44.yr}:

%% check: skip
\begin{lstlisting}[style=coloredverbatimError]
class Bar { pub self () {} }

class Foo {
    pub self () {}
    impl Bar {}
}

fn main () {}
\end{lstlisting}

\begin{lstlisting}[style=bashVerb, breaklines=true, escapechar=~]
Error[E4094] : impl statement must be followed by a trait, not test44::Bar
 --> test_resources/diagnostics/test44.yr:(1,7)
 1  ~\lstbox{\textSFxi{}}~ class Bar { pub self () {} }
    ~\lstbox{\textSFv{}}~       ^^^
    ~\lstbox{\textSFii{}}~~\lstbox{\textSFx{}}~~\lstbox{\textSFx{}}~~\lstbox{\textSFx{}}~~\lstbox{\textSFx{}}~~\lstbox{\textSFx{}}~~\lstbox{\textSFx{}}~~\lstbox{\textSFx{}}~
\end{lstlisting}

\subsubsection*{How to fix it}

Name a trait. To add methods to a type without a trait, declare them in the type body.

\errorcode{E4101}{the base class \errhole{} is marked as final}
\label{sec:codes:e4101}

Raised during \textbf{validation}, from \texttt{ValidateErrorMessage::\allowbreak{}INHERIT\_\allowbreak{}FINAL\_\allowbreak{}CLASS} (\textit{src/\allowbreak{}ymirc/\allowbreak{}semantic/\allowbreak{}validator/\allowbreak{}errors.yr}).

\subsubsection*{What it means}

A class inherits from a class marked \tokennolst{final}. \tokennolst{final} says the class is a leaf of its hierarchy.

\subsubsection*{Example}

\textit{test\_\allowbreak{}resources/\allowbreak{}diagnostics/\allowbreak{}test45.yr}:

%% check: skip
\begin{lstlisting}[style=coloredverbatimError]
@final
class A { pub self () {} }

class B over A { pub self () {} }

fn main () { let _ = copy B (); }
\end{lstlisting}

\begin{lstlisting}[style=bashVerb, breaklines=true, escapechar=~]
Error[E4101] : the base class test45::A is marked as final
 --> test_resources/diagnostics/test45.yr:(4,14)
 4  ~\lstbox{\textSFxi{}}~ class B over A { pub self () {} }
    ~\lstbox{\textSFv{}}~              ^
    ~\lstbox{\textSFxi{}}~
    = when validating -- (test_resources/diagnostics/test45.yr:(4,7))
    ~\lstbox{\textSFii{}}~~\lstbox{\textSFx{}}~~\lstbox{\textSFx{}}~~\lstbox{\textSFx{}}~~\lstbox{\textSFx{}}~~\lstbox{\textSFx{}}~~\lstbox{\textSFx{}}~~\lstbox{\textSFx{}}~
Error[E4104] : symbol test45::B has errors
 --> test_resources/diagnostics/test45.yr:(6,27)
 4  ~\lstbox{\textSFxi{}}~ class B over A { pub self () {} }
    ~\lstbox{\textSFv{}}~              -  error: the base class test45::A is marked as final
 5  ~\lstbox{\textSFxi{}}~
 6  ~\lstbox{\textSFxi{}}~ fn main () { let _ = copy B (); }
    ~\lstbox{\textSFv{}}~                           ^
    ~\lstbox{\textSFxi{}}~
    = when validating -- (test_resources/diagnostics/test45.yr:(4,7))
    = when validating test45::main -- (test_resources/diagnostics/test45.yr:(6,4))
    ~\lstbox{\textSFii{}}~~\lstbox{\textSFx{}}~~\lstbox{\textSFx{}}~~\lstbox{\textSFx{}}~~\lstbox{\textSFx{}}~~\lstbox{\textSFx{}}~~\lstbox{\textSFx{}}~~\lstbox{\textSFx{}}~
\end{lstlisting}

\subsubsection*{How to fix it}

Remove \tokennolst{final} from the base class, or inherit from one of its ancestors.

\errorcode{E4102}{the base of a class must be a class, not a \errhole{}}
\label{sec:codes:e4102}

Raised during \textbf{validation}, from \texttt{ValidateErrorMessage::\allowbreak{}INHERIT\_\allowbreak{}NO\_\allowbreak{}CLASS} (\textit{src/\allowbreak{}ymirc/\allowbreak{}semantic/\allowbreak{}validator/\allowbreak{}errors.yr}).

\subsubsection*{What it means}

The base of a class is not a class -- a record, a trait or a plain type was named instead.

\subsubsection*{Example}

\textit{test\_\allowbreak{}resources/\allowbreak{}diagnostics/\allowbreak{}test46.yr}:

%% check: skip
\begin{lstlisting}[style=coloredverbatimError]
record R { let x: i32; }

class B over R { pub self () {} }

fn main () { let _ = copy B (); }
\end{lstlisting}

\begin{lstlisting}[style=bashVerb, breaklines=true, escapechar=~]
Error[E4102] : the base of a class must be a class, not a test46::R
 --> test_resources/diagnostics/test46.yr:(3,14)
 3  ~\lstbox{\textSFxi{}}~ class B over R { pub self () {} }
    ~\lstbox{\textSFv{}}~              ^
    ~\lstbox{\textSFxi{}}~
    = when validating -- (test_resources/diagnostics/test46.yr:(3,7))
    ~\lstbox{\textSFii{}}~~\lstbox{\textSFx{}}~~\lstbox{\textSFx{}}~~\lstbox{\textSFx{}}~~\lstbox{\textSFx{}}~~\lstbox{\textSFx{}}~~\lstbox{\textSFx{}}~~\lstbox{\textSFx{}}~
Error[E4104] : symbol test46::B has errors
 --> test_resources/diagnostics/test46.yr:(5,27)
 3  ~\lstbox{\textSFxi{}}~ class B over R { pub self () {} }
    ~\lstbox{\textSFv{}}~              -  error: the base of a class must be a class, not a test46::R
 4  ~\lstbox{\textSFxi{}}~
 5  ~\lstbox{\textSFxi{}}~ fn main () { let _ = copy B (); }
    ~\lstbox{\textSFv{}}~                           ^
    ~\lstbox{\textSFxi{}}~
    = when validating -- (test_resources/diagnostics/test46.yr:(3,7))
    = when validating test46::main -- (test_resources/diagnostics/test46.yr:(5,4))
    ~\lstbox{\textSFii{}}~~\lstbox{\textSFx{}}~~\lstbox{\textSFx{}}~~\lstbox{\textSFx{}}~~\lstbox{\textSFx{}}~~\lstbox{\textSFx{}}~~\lstbox{\textSFx{}}~~\lstbox{\textSFx{}}~
\end{lstlisting}

\subsubsection*{How to fix it}

Inherit from a class. To reuse a trait, use an \tokennolst{impl} block; to reuse a record, hold it as a field.

\errorcode{E4103}{method \errhole{} cannot be inline as it is virtual(overridable or inherited)}
\label{sec:codes:e4103}

Raised during \textbf{validation}, from \texttt{ValidateErrorMessage::\allowbreak{}INLINE\_\allowbreak{}VIRTUAL\_\allowbreak{}METHOD} (\textit{src/\allowbreak{}ymirc/\allowbreak{}semantic/\allowbreak{}validator/\allowbreak{}errors.yr}).

\subsubsection*{What it means}

A method is marked \tokennolst{inline} and is virtual. A virtual call is resolved at run time through the vtable, so there is no single body to inline.

\subsubsection*{Example}

\textit{test\_\allowbreak{}resources/\allowbreak{}diagnostics/\allowbreak{}test47.yr}:

%% check: skip
\begin{lstlisting}[style=coloredverbatimError]
class A {
    pub self () {}
    @inline pub fn foo (self) {}
}

fn main () { let _ = copy A (); }
\end{lstlisting}

\begin{lstlisting}[style=bashVerb, breaklines=true, escapechar=~]
Error[E4103] : method test47::A::foo cannot be inline as it is virtual(overridable or inherited)
 --> test_resources/diagnostics/test47.yr:(3,6)
 3  ~\lstbox{\textSFxi{}}~     @inline pub fn foo (self) {}
    ~\lstbox{\textSFv{}}~      ^^^^^^
    ~\lstbox{\textSFxi{}}~
    = when validating test47::A::foo -- (test_resources/diagnostics/test47.yr:(3,20))
    = when validating -- (test_resources/diagnostics/test47.yr:(1,7))
    ~\lstbox{\textSFii{}}~~\lstbox{\textSFx{}}~~\lstbox{\textSFx{}}~~\lstbox{\textSFx{}}~~\lstbox{\textSFx{}}~~\lstbox{\textSFx{}}~~\lstbox{\textSFx{}}~~\lstbox{\textSFx{}}~
Error[E4104] : symbol test47::A has errors
 --> test_resources/diagnostics/test47.yr:(6,27)
 6  ~\lstbox{\textSFxi{}}~ fn main () { let _ = copy A (); }
    ~\lstbox{\textSFv{}}~                           ^
 --> test_resources/diagnostics/test47.yr:(3,6)
 3  ~\lstbox{\textSFxi{}}~     @inline pub fn foo (self) {}
    ~\lstbox{\textSFv{}}~      ------  error: method test47::A::foo cannot be inline as it is virtual(overridable or inherited)
    ~\lstbox{\textSFxi{}}~
    = when validating test47::A::foo -- (test_resources/diagnostics/test47.yr:(3,20))
    = when validating -- (test_resources/diagnostics/test47.yr:(1,7))
    = when validating test47::main -- (test_resources/diagnostics/test47.yr:(6,4))
    ~\lstbox{\textSFii{}}~~\lstbox{\textSFx{}}~~\lstbox{\textSFx{}}~~\lstbox{\textSFx{}}~~\lstbox{\textSFx{}}~~\lstbox{\textSFx{}}~~\lstbox{\textSFx{}}~~\lstbox{\textSFx{}}~
\end{lstlisting}

\subsubsection*{How to fix it}

Drop \tokennolst{inline}, or make the method \tokennolst{final} so it is no longer overridable.

\errorcode{E4133}{field \errhole{} is initialized multiple times}
\label{sec:codes:e4133}

Raised during \textbf{validation}, from \texttt{ValidateErrorMessage::\allowbreak{}MULTIPLE\_\allowbreak{}FIELD\_\allowbreak{}INIT} (\textit{src/\allowbreak{}ymirc/\allowbreak{}semantic/\allowbreak{}validator/\allowbreak{}errors.yr}).

\subsubsection*{What it means}

A constructor assigns the same field twice.

\subsubsection*{Example}

\textit{test\_\allowbreak{}resources/\allowbreak{}union/\allowbreak{}test3.yr}:

%% check: skip
\begin{lstlisting}[style=coloredverbatimError]
@overlaid
record A {
    let mut a : i32;
    let mut b : i64 = 0;

    pub self ()
        with a = 12
        , a = 9
    {}
}

fn main () {
    let _ = A ();
}
\end{lstlisting}

\begin{lstlisting}[style=bashVerb, breaklines=true, escapechar=~]
Error[E4133] : field a is initialized multiple times
 --> test_resources/union/test3.yr:(8,11)
 8  ~\lstbox{\textSFxi{}}~         , a = 9
    ~\lstbox{\textSFv{}}~           ^
    ~\lstbox{\textSFxi{}}~ Note : here
    ~\lstbox{\textSFxi{}}~  --> test_resources/union/test3.yr:(7,14)
    ~\lstbox{\textSFxi{}}~  7  ~\lstbox{\textSFxi{}}~         with a = 12
    ~\lstbox{\textSFxi{}}~     ~\lstbox{\textSFv{}}~              ^
    ~\lstbox{\textSFxi{}}~     ~\lstbox{\textSFii{}}~~\lstbox{\textSFx{}}~~\lstbox{\textSFx{}}~~\lstbox{\textSFx{}}~~\lstbox{\textSFx{}}~~\lstbox{\textSFx{}}~~\lstbox{\textSFx{}}~~\lstbox{\textSFx{}}~
    ~\lstbox{\textSFii{}}~~\lstbox{\textSFx{}}~~\lstbox{\textSFx{}}~~\lstbox{\textSFx{}}~~\lstbox{\textSFx{}}~~\lstbox{\textSFx{}}~~\lstbox{\textSFvii{}}~~\lstbox{\textSFx{}}~
\end{lstlisting}

\subsubsection*{How to fix it}

Keep one assignment. When the two came from a field constructor and an explicit assignment, drop the explicit one.

\errorcode{E4141}{method \errhole{} is defined as mutable}
\label{sec:codes:e4141}

Raised during \textbf{validation}, from \texttt{ValidateErrorMessage::\allowbreak{}MUTABLE\_\allowbreak{}METHOD} (\textit{src/\allowbreak{}ymirc/\allowbreak{}semantic/\allowbreak{}validator/\allowbreak{}errors.yr}).

\subsubsection*{What it means}

A mutable method is called on a value that is not mutable.

\subsubsection*{Example}

\textit{test\_\allowbreak{}resources/\allowbreak{}diagnostics/\allowbreak{}test55.yr}:

%% check: skip
\begin{lstlisting}[style=coloredverbatimError]
class Foo {
    let mut x: i32 = 0;
    pub self () {}
    pub fn bar (mut self) { self.x = 1; }
}

fn main () {
    let f = copy Foo ();
    f.bar ();
}
\end{lstlisting}

\begin{lstlisting}[style=bashVerb, breaklines=true, escapechar=~]
Error[E4141] : method test55::Foo::bar ()-> void is defined as mutable
 --> test_resources/diagnostics/test55.yr:(4,12)
 4  ~\lstbox{\textSFxi{}}~     pub fn bar (mut self) { self.x = 1; }
    ~\lstbox{\textSFv{}}~            ^^^
    ~\lstbox{\textSFxi{}}~ Note : implicit alias of type &(test55::Foo) is prohibited, it implicitly gives up on borrowings of mutable values
    ~\lstbox{\textSFxi{}}~  --> test_resources/diagnostics/test55.yr:(9,5)
    ~\lstbox{\textSFxi{}}~  9  ~\lstbox{\textSFxi{}}~     f.bar ();
    ~\lstbox{\textSFxi{}}~     ~\lstbox{\textSFv{}}~     ^
    ~\lstbox{\textSFxi{}}~     ~\lstbox{\textSFii{}}~~\lstbox{\textSFx{}}~~\lstbox{\textSFx{}}~~\lstbox{\textSFx{}}~~\lstbox{\textSFx{}}~~\lstbox{\textSFx{}}~~\lstbox{\textSFx{}}~~\lstbox{\textSFx{}}~
    ~\lstbox{\textSFii{}}~~\lstbox{\textSFx{}}~~\lstbox{\textSFx{}}~~\lstbox{\textSFx{}}~~\lstbox{\textSFx{}}~~\lstbox{\textSFx{}}~~\lstbox{\textSFvii{}}~~\lstbox{\textSFx{}}~
\end{lstlisting}

\subsubsection*{How to fix it}

Make the value mutable, or call a constant method instead, cf. \hyperref[sec:codes:e4082]{E4082}.

\errorcode{E4144}{class \errhole{} is not abstract but has empty methods}
\label{sec:codes:e4144}

Raised during \textbf{validation}, from \texttt{ValidateErrorMessage::\allowbreak{}NON\_\allowbreak{}ABSTRACT\_\allowbreak{}NOT\_\allowbreak{}COMPLETE} (\textit{src/\allowbreak{}ymirc/\allowbreak{}semantic/\allowbreak{}validator/\allowbreak{}errors.yr}).

\subsubsection*{What it means}

A class has methods with no body but is not declared \tokennolst{abstract}, so it could be instantiated with nothing to call.

\subsubsection*{Example}

\textit{test\_\allowbreak{}resources/\allowbreak{}implement/\allowbreak{}test3.yr}:

%% check: skip
\begin{lstlisting}[style=coloredverbatimError]
trait X {
    pub fn foo (self);
}

class A {
    pub self () {}
    impl X ;
}

record B {
    pub self () {}
    impl X ;
}
\end{lstlisting}

\begin{lstlisting}[style=bashVerb, breaklines=true, escapechar=~]
Error[E4144] : class test3::A is not abstract but has empty methods
 --> test_resources/implement/test3.yr:(5,7)
 5  ~\lstbox{\textSFxi{}}~ class A {
    ~\lstbox{\textSFv{}}~       ^  note: X::foo ()-> void is empty -- (test_resources/implement/test3.yr:(2,12))
    ~\lstbox{\textSFii{}}~~\lstbox{\textSFx{}}~~\lstbox{\textSFx{}}~~\lstbox{\textSFx{}}~~\lstbox{\textSFx{}}~~\lstbox{\textSFx{}}~~\lstbox{\textSFx{}}~~\lstbox{\textSFx{}}~
\end{lstlisting}

\subsubsection*{How to fix it}

Give the methods bodies, or declare the class \tokennolst{abstract}, cf. \hyperref[sec:codes:e4026]{E4026}.

\errorcode{E4147}{\errhole{} is not an ancestor type of \errhole{}}
\label{sec:codes:e4147}

Raised during \textbf{validation}, from \texttt{ValidateErrorMessage::\allowbreak{}NOT\_\allowbreak{}ANCESTOR} (\textit{src/\allowbreak{}ymirc/\allowbreak{}semantic/\allowbreak{}validator/\allowbreak{}errors.yr}).

\subsubsection*{What it means}

A type is used where an ancestor of another type is required, and it is not one of its ancestors.

\subsubsection*{Example}

\textit{test\_\allowbreak{}resources/\allowbreak{}regressions/\allowbreak{}throws\_\allowbreak{}without\_\allowbreak{}ampersand/\allowbreak{}test3.yr}:

%% check: skip
\begin{lstlisting}[style=coloredverbatimError]
fn main ()
    throws i32
{}

fn foo () {
    throw 12;
}
\end{lstlisting}

\begin{lstlisting}[style=bashVerb, breaklines=true, escapechar=~]
Error[E4147] : mut &(mut core::exception::Exception) is not an ancestor type of i32
 --> test_resources/regressions/throws_without_ampersand/test3.yr:(6,5)
 6  ~\lstbox{\textSFxi{}}~     throw 12;
    ~\lstbox{\textSFv{}}~     ^^^^^
    ~\lstbox{\textSFii{}}~~\lstbox{\textSFx{}}~~\lstbox{\textSFx{}}~~\lstbox{\textSFx{}}~~\lstbox{\textSFx{}}~~\lstbox{\textSFx{}}~~\lstbox{\textSFx{}}~~\lstbox{\textSFx{}}~
\end{lstlisting}

\subsubsection*{How to fix it}

Name a class the other one actually derives from.

\errorcode{E4150}{\errhole{} is not a class type}
\label{sec:codes:e4150}

Raised during \textbf{validation}, from \texttt{ValidateErrorMessage::\allowbreak{}NOT\_\allowbreak{}A\_\allowbreak{}CLASS} (\textit{src/\allowbreak{}ymirc/\allowbreak{}semantic/\allowbreak{}validator/\allowbreak{}errors.yr}).

\subsubsection*{What it means}

An operation defined only on classes was applied to something else.

\subsubsection*{Example}

\textit{test\_\allowbreak{}resources/\allowbreak{}regressions/\allowbreak{}throws\_\allowbreak{}without\_\allowbreak{}ampersand/\allowbreak{}test3.yr}:

%% check: skip
\begin{lstlisting}[style=coloredverbatimError]
fn main ()
    throws i32
{}

fn foo () {
    throw 12;
}
\end{lstlisting}

\begin{lstlisting}[style=bashVerb, breaklines=true, escapechar=~]
Error[E4150] : i32 is not a class type
 --> test_resources/regressions/throws_without_ampersand/test3.yr:(2,12)
 2  ~\lstbox{\textSFxi{}}~     throws i32
    ~\lstbox{\textSFv{}}~            ^^^
    ~\lstbox{\textSFii{}}~~\lstbox{\textSFx{}}~~\lstbox{\textSFx{}}~~\lstbox{\textSFx{}}~~\lstbox{\textSFx{}}~~\lstbox{\textSFx{}}~~\lstbox{\textSFx{}}~~\lstbox{\textSFx{}}~
\end{lstlisting}

\subsubsection*{How to fix it}

Use a class value.

\errorcode{E4162}{\errhole{} is not a trait}
\label{sec:codes:e4162}

Raised during \textbf{validation}, from \texttt{ValidateErrorMessage::\allowbreak{}NOT\_\allowbreak{}A\_\allowbreak{}TRAIT} (\textit{src/\allowbreak{}ymirc/\allowbreak{}semantic/\allowbreak{}validator/\allowbreak{}errors.yr}).

\subsubsection*{What it means}

A name is used where a trait is required and does not denote one, cf. \hyperref[sec:codes:e4094]{E4094}.

\subsubsection*{How to fix it}

Name a trait.

\errorcode{E4168}{no constructor found for class \errhole{}}
\label{sec:codes:e4168}

Raised during \textbf{validation}, from \texttt{ValidateErrorMessage::\allowbreak{}NO\_\allowbreak{}CTOR\_\allowbreak{}FOUND} (\textit{src/\allowbreak{}ymirc/\allowbreak{}semantic/\allowbreak{}validator/\allowbreak{}errors.yr}).

\subsubsection*{What it means}

A class was constructed and it declares no constructor callable that way.

\subsubsection*{Example}

\textit{test\_\allowbreak{}resources/\allowbreak{}default\_\allowbreak{}ctor/\allowbreak{}test8.yr}:

%% check: skip
\begin{lstlisting}[style=coloredverbatimError]
// The generated ctor carries the protection of the `self;` asking for it.
class Foo {
   pub let x: i32;
   prv self;

   pub fn get (self)-> i32 { self.x }
}

fn main () {
    let a = copy Foo (1);
    a.get ();
}
\end{lstlisting}

\begin{lstlisting}[style=bashVerb, breaklines=true, escapechar=~]
Error[E4168] : no constructor found for class test8::Foo
 --> test_resources/default_ctor/test8.yr:(10,18)
10  ~\lstbox{\textSFxi{}}~     let a = copy Foo (1);
    ~\lstbox{\textSFv{}}~                  ^^^
 --> test_resources/default_ctor/test8.yr:(4,8)
 4  ~\lstbox{\textSFxi{}}~    prv self;
    ~\lstbox{\textSFv{}}~        ----  error: test8::Foo::self is private in this context
    ~\lstbox{\textSFxi{}}~
    = when validating test8::main -- (test_resources/default_ctor/test8.yr:(9,4))
    ~\lstbox{\textSFii{}}~~\lstbox{\textSFx{}}~~\lstbox{\textSFx{}}~~\lstbox{\textSFx{}}~~\lstbox{\textSFx{}}~~\lstbox{\textSFx{}}~~\lstbox{\textSFx{}}~~\lstbox{\textSFx{}}~
Error[E4224] : undefined operator . for types error and get
 --> test_resources/default_ctor/test8.yr:(11,6)
11  ~\lstbox{\textSFxi{}}~     a.get ();
    ~\lstbox{\textSFv{}}~      ^
    ~\lstbox{\textSFxi{}}~
    = when validating test8::main -- (test_resources/default_ctor/test8.yr:(9,4))
    ~\lstbox{\textSFii{}}~~\lstbox{\textSFx{}}~~\lstbox{\textSFx{}}~~\lstbox{\textSFx{}}~~\lstbox{\textSFx{}}~~\lstbox{\textSFx{}}~~\lstbox{\textSFx{}}~~\lstbox{\textSFx{}}~
\end{lstlisting}

\subsubsection*{How to fix it}

Declare a constructor matching the arguments given, or call one that exists; the candidates are listed under the diagnostic.

\errorcode{E4169}{no constructor named \errhole{} found for class \errhole{}}
\label{sec:codes:e4169}

Raised during \textbf{validation}, from \texttt{ValidateErrorMessage::\allowbreak{}NO\_\allowbreak{}CTOR\_\allowbreak{}FOUND\_\allowbreak{}NAME} (\textit{src/\allowbreak{}ymirc/\allowbreak{}semantic/\allowbreak{}validator/\allowbreak{}errors.yr}).

\subsubsection*{What it means}

A named constructor was called and the class has no constructor of that name.

\subsubsection*{Example}

\textit{test\_\allowbreak{}resources/\allowbreak{}diagnostics/\allowbreak{}test62.yr}:

%% check: skip
\begin{lstlisting}[style=coloredverbatimError]
class Foo {
    pub self () {}
}

fn main () { let _ = copy Foo::named (); }
\end{lstlisting}

\begin{lstlisting}[style=bashVerb, breaklines=true, escapechar=~]
Error[E4169] : no constructor named named found for class test62::Foo
 --> test_resources/diagnostics/test62.yr:(5,30)
 5  ~\lstbox{\textSFxi{}}~ fn main () { let _ = copy Foo::named (); }
    ~\lstbox{\textSFv{}}~                              ^^
    ~\lstbox{\textSFii{}}~~\lstbox{\textSFx{}}~~\lstbox{\textSFx{}}~~\lstbox{\textSFx{}}~~\lstbox{\textSFx{}}~~\lstbox{\textSFx{}}~~\lstbox{\textSFx{}}~~\lstbox{\textSFx{}}~
\end{lstlisting}

\subsubsection*{How to fix it}

Check the spelling, or use the unnamed constructor.

\errorcode{E4172}{record or entity type \errhole{} has no size due to forward references}
\label{sec:codes:e4172}

Raised during \textbf{validation}, from \texttt{ValidateErrorMessage::\allowbreak{}NO\_\allowbreak{}SIZE\_\allowbreak{}FORWARD\_\allowbreak{}REF} (\textit{src/\allowbreak{}ymirc/\allowbreak{}semantic/\allowbreak{}validator/\allowbreak{}errors.yr}).

\subsubsection*{What it means}

The size of a record or an entity cannot be computed because one of its fields refers to a type not yet defined.

\subsubsection*{Example}

\textit{test\_\allowbreak{}resources/\allowbreak{}struct/\allowbreak{}test1.yr}:

%% check: skip
\begin{lstlisting}[style=coloredverbatimError]
in test1;

record A {
    let a : A;
}

record B {
    let c : C;
}

record C {
    let b : B;
}
\end{lstlisting}

\begin{lstlisting}[style=bashVerb, breaklines=true, escapechar=~]
Error[E4172] : record or entity type test1::A has no size due to forward references
 --> test_resources/struct/test1.yr:(3,8)
 3  ~\lstbox{\textSFxi{}}~ record A {
    ~\lstbox{\textSFv{}}~        ^
    ~\lstbox{\textSFxi{}}~ Note : for the field a
    ~\lstbox{\textSFxi{}}~  --> test_resources/struct/test1.yr:(4,9)
    ~\lstbox{\textSFxi{}}~  4  ~\lstbox{\textSFxi{}}~     let a : A;
    ~\lstbox{\textSFxi{}}~     ~\lstbox{\textSFv{}}~         ^
    ~\lstbox{\textSFxi{}}~     ~\lstbox{\textSFxi{}}~ Note : within type test1::A
    ~\lstbox{\textSFxi{}}~     ~\lstbox{\textSFxi{}}~  --> test_resources/struct/test1.yr:(3,8)
    ~\lstbox{\textSFxi{}}~     ~\lstbox{\textSFxi{}}~  3  ~\lstbox{\textSFxi{}}~ record A {
    ~\lstbox{\textSFxi{}}~     ~\lstbox{\textSFxi{}}~     ~\lstbox{\textSFv{}}~        ^
    ~\lstbox{\textSFxi{}}~     ~\lstbox{\textSFxi{}}~     ~\lstbox{\textSFii{}}~~\lstbox{\textSFx{}}~~\lstbox{\textSFx{}}~~\lstbox{\textSFx{}}~~\lstbox{\textSFx{}}~~\lstbox{\textSFx{}}~~\lstbox{\textSFx{}}~~\lstbox{\textSFx{}}~
    ~\lstbox{\textSFxi{}}~     ~\lstbox{\textSFii{}}~~\lstbox{\textSFx{}}~~\lstbox{\textSFx{}}~~\lstbox{\textSFx{}}~~\lstbox{\textSFx{}}~~\lstbox{\textSFx{}}~~\lstbox{\textSFvii{}}~~\lstbox{\textSFx{}}~
    ~\lstbox{\textSFii{}}~~\lstbox{\textSFx{}}~~\lstbox{\textSFx{}}~~\lstbox{\textSFx{}}~~\lstbox{\textSFx{}}~~\lstbox{\textSFx{}}~~\lstbox{\textSFvii{}}~~\lstbox{\textSFx{}}~
\end{lstlisting}

\subsubsection*{How to fix it}

Break the cycle: hold the offending field behind a class or a pointer, cf. \hyperref[sec:codes:e4017]{E4017}.

\errorcode{E4173}{class \errhole{} has no ancestor}
\label{sec:codes:e4173}

Raised during \textbf{validation}, from \texttt{ValidateErrorMessage::\allowbreak{}NO\_\allowbreak{}SUPER\_\allowbreak{}CLASS} (\textit{src/\allowbreak{}ymirc/\allowbreak{}semantic/\allowbreak{}validator/\allowbreak{}errors.yr}).

\subsubsection*{What it means}

\tokennolst{super} was used in a class with no base class.

\subsubsection*{Example}

\textit{test\_\allowbreak{}resources/\allowbreak{}diagnostics/\allowbreak{}test63.yr}:

%% check: skip
\begin{lstlisting}[style=coloredverbatimError]
class Foo {
    pub self () {
        self.super ();
    }
}

fn main () { let _ = copy Foo (); }
\end{lstlisting}

\begin{lstlisting}[style=bashVerb, breaklines=true, escapechar=~]
Error[E4173] : class &(test63::Foo) has no ancestor
 --> test_resources/diagnostics/test63.yr:(3,9)
 3  ~\lstbox{\textSFxi{}}~         self.super ();
    ~\lstbox{\textSFv{}}~         ^^^^ ^^^^^
    ~\lstbox{\textSFii{}}~~\lstbox{\textSFx{}}~~\lstbox{\textSFx{}}~~\lstbox{\textSFx{}}~~\lstbox{\textSFx{}}~~\lstbox{\textSFx{}}~~\lstbox{\textSFx{}}~~\lstbox{\textSFx{}}~
\end{lstlisting}

\subsubsection*{How to fix it}

Give the class a base, or remove the \tokennolst{super} access.

\errorcode{E4181}{method \errhole{} must have a body to override \errhole{}}
\label{sec:codes:e4181}

Raised during \textbf{validation}, from \texttt{ValidateErrorMessage::\allowbreak{}OVERRIDE\_\allowbreak{}EMPTY} (\textit{src/\allowbreak{}ymirc/\allowbreak{}semantic/\allowbreak{}validator/\allowbreak{}errors.yr}).

\subsubsection*{What it means}

A method marked \tokennolst{override} has no body. Overriding replaces an implementation, so it has to provide one.

\subsubsection*{Example}

\textit{test\_\allowbreak{}resources/\allowbreak{}diagnostics/\allowbreak{}test64.yr}:

%% check: skip
\begin{lstlisting}[style=coloredverbatimError]
class A {
    pub self () {}
    pub fn foo (self) {}
}

class B over A {
    pub self () {}
    pub over foo (self);
}

fn main () { let _ = copy B (); }
\end{lstlisting}

\begin{lstlisting}[style=bashVerb, breaklines=true, escapechar=~]
Error[E4181] : method test64::B::foo ()-> void must have a body to override test64::A::foo ()-> void
 --> test_resources/diagnostics/test64.yr:(8,14)
 8  ~\lstbox{\textSFxi{}}~     pub over foo (self);
    ~\lstbox{\textSFv{}}~              ^^^
    ~\lstbox{\textSFxi{}}~
    = when validating -- (test_resources/diagnostics/test64.yr:(6,7))
    ~\lstbox{\textSFii{}}~~\lstbox{\textSFx{}}~~\lstbox{\textSFx{}}~~\lstbox{\textSFx{}}~~\lstbox{\textSFx{}}~~\lstbox{\textSFx{}}~~\lstbox{\textSFx{}}~~\lstbox{\textSFx{}}~
Error[E4104] : symbol test64::B has errors
 --> test_resources/diagnostics/test64.yr:(11,27)
11  ~\lstbox{\textSFxi{}}~ fn main () { let _ = copy B (); }
    ~\lstbox{\textSFv{}}~                           ^
 --> test_resources/diagnostics/test64.yr:(8,14)
 8  ~\lstbox{\textSFxi{}}~     pub over foo (self);
    ~\lstbox{\textSFv{}}~              ---  error: method test64::B::foo ()-> void must have a body to override test64::A::foo ()-> void
    ~\lstbox{\textSFxi{}}~
    = when validating -- (test_resources/diagnostics/test64.yr:(6,7))
    = when validating test64::main -- (test_resources/diagnostics/test64.yr:(11,4))
    ~\lstbox{\textSFii{}}~~\lstbox{\textSFx{}}~~\lstbox{\textSFx{}}~~\lstbox{\textSFx{}}~~\lstbox{\textSFx{}}~~\lstbox{\textSFx{}}~~\lstbox{\textSFx{}}~~\lstbox{\textSFx{}}~
\end{lstlisting}

\subsubsection*{How to fix it}

Give the method a body, or drop \tokennolst{override} and declare it in an abstract class as an empty method.

\errorcode{E4182}{method \errhole{} must be a field method to override \errhole{}}
\label{sec:codes:e4182}

Raised during \textbf{validation}, from \texttt{ValidateErrorMessage::\allowbreak{}OVERRIDE\_\allowbreak{}FIELD\_\allowbreak{}NO\_\allowbreak{}FIELD} (\textit{src/\allowbreak{}ymirc/\allowbreak{}semantic/\allowbreak{}validator/\allowbreak{}errors.yr}).

\subsubsection*{What it means}

An ordinary method tries to override a field method. The two are reached through different syntax, so one cannot stand in for the other.

\subsubsection*{Example}

\textit{test\_\allowbreak{}resources/\allowbreak{}lit\_\allowbreak{}class/\allowbreak{}methods/\allowbreak{}test3.yr}:

%% check: skip
\begin{lstlisting}[style=coloredverbatimError]
@abstract
class A {
    let _x : i32 = 0;
    let _y : i32 = 9;

    pub self () {}

    @field
    pub fn foo (self)-> i32;
}

class B over A {
    pub self () {}

    pub over foo (self)-> i32 {
        self._x
    }
}

fn main () {
    let b : &A = copy B ();
    let _z_ = b.foo;
    let _y_ = b.foo (1);
}
\end{lstlisting}

\begin{lstlisting}[style=bashVerb, breaklines=true, escapechar=~]
Error[E4182] : method test3::B::foo ()-> i32 must be a field method to override test3::A::foo ()-> i32
 --> test_resources/lit_class/methods/test3.yr:(15,14)
15  ~\lstbox{\textSFxi{}}~     pub over foo (self)-> i32 {
    ~\lstbox{\textSFv{}}~              ^^^
 --> test_resources/lit_class/methods/test3.yr:(8,6)
 8  ~\lstbox{\textSFxi{}}~     @field
    ~\lstbox{\textSFv{}}~      -----
    ~\lstbox{\textSFxi{}}~
    = when validating -- (test_resources/lit_class/methods/test3.yr:(12,7))
    ~\lstbox{\textSFii{}}~~\lstbox{\textSFx{}}~~\lstbox{\textSFx{}}~~\lstbox{\textSFx{}}~~\lstbox{\textSFx{}}~~\lstbox{\textSFx{}}~~\lstbox{\textSFx{}}~~\lstbox{\textSFx{}}~
Error[E4104] : symbol test3::B has errors
 --> test_resources/lit_class/methods/test3.yr:(21,23)
21  ~\lstbox{\textSFxi{}}~     let b : &A = copy B ();
    ~\lstbox{\textSFv{}}~                       ^
 --> test_resources/lit_class/methods/test3.yr:(8,6)
 8  ~\lstbox{\textSFxi{}}~     @field
    ~\lstbox{\textSFv{}}~      -----
   ...
15  ~\lstbox{\textSFxi{}}~     pub over foo (self)-> i32 {
    ~\lstbox{\textSFv{}}~              ---  error: method test3::B::foo ()-> i32 must be a field method to override test3::A::foo ()-> i32
    ~\lstbox{\textSFxi{}}~
    = when validating -- (test_resources/lit_class/methods/test3.yr:(12,7))
    = when validating test3::main -- (test_resources/lit_class/methods/test3.yr:(20,4))
    ~\lstbox{\textSFii{}}~~\lstbox{\textSFx{}}~~\lstbox{\textSFx{}}~~\lstbox{\textSFx{}}~~\lstbox{\textSFx{}}~~\lstbox{\textSFx{}}~~\lstbox{\textSFx{}}~~\lstbox{\textSFx{}}~
Error[E4224] : undefined operator . for types error and foo
 --> test_resources/lit_class/methods/test3.yr:(22,16)
22  ~\lstbox{\textSFxi{}}~     let _z_ = b.foo;
    ~\lstbox{\textSFv{}}~                ^
    ~\lstbox{\textSFxi{}}~
    = when validating test3::main -- (test_resources/lit_class/methods/test3.yr:(20,4))
    ~\lstbox{\textSFii{}}~~\lstbox{\textSFx{}}~~\lstbox{\textSFx{}}~~\lstbox{\textSFx{}}~~\lstbox{\textSFx{}}~~\lstbox{\textSFx{}}~~\lstbox{\textSFx{}}~~\lstbox{\textSFx{}}~
Error[E4224] : undefined operator . for types error and foo
 --> test_resources/lit_class/methods/test3.yr:(23,16)
23  ~\lstbox{\textSFxi{}}~     let _y_ = b.foo (1);
    ~\lstbox{\textSFv{}}~                ^
    ~\lstbox{\textSFxi{}}~
    = when validating test3::main -- (test_resources/lit_class/methods/test3.yr:(20,4))
    ~\lstbox{\textSFii{}}~~\lstbox{\textSFx{}}~~\lstbox{\textSFx{}}~~\lstbox{\textSFx{}}~~\lstbox{\textSFx{}}~~\lstbox{\textSFx{}}~~\lstbox{\textSFx{}}~~\lstbox{\textSFx{}}~
\end{lstlisting}

\subsubsection*{How to fix it}

Declare the overriding method as a field method too.

\errorcode{E4183}{cannot override final method \errhole{}}
\label{sec:codes:e4183}

Raised during \textbf{validation}, from \texttt{ValidateErrorMessage::\allowbreak{}OVERRIDE\_\allowbreak{}FINAL} (\textit{src/\allowbreak{}ymirc/\allowbreak{}semantic/\allowbreak{}validator/\allowbreak{}errors.yr}).

\subsubsection*{What it means}

A method marked \tokennolst{final} in an ancestor is being overridden. \tokennolst{final} says the implementation is the last word.

\subsubsection*{Example}

\textit{test\_\allowbreak{}resources/\allowbreak{}diagnostics/\allowbreak{}test65.yr}:

%% check: skip
\begin{lstlisting}[style=coloredverbatimError]
class A {
    pub self () {}
    @final pub fn foo (self) {}
}

class B over A {
    pub self () {}
    pub over foo (self) {}
}

fn main () { let _ = copy B (); }
\end{lstlisting}

\begin{lstlisting}[style=bashVerb, breaklines=true, escapechar=~]
Error[E4183] : cannot override final method test65::A::foo ()-> void
 --> test_resources/diagnostics/test65.yr:(3,19)
 3  ~\lstbox{\textSFxi{}}~     @final pub fn foo (self) {}
    ~\lstbox{\textSFv{}}~                   ^^^
 --> test_resources/diagnostics/test65.yr:(8,14)
 8  ~\lstbox{\textSFxi{}}~     pub over foo (self) {}
    ~\lstbox{\textSFv{}}~              ---
    ~\lstbox{\textSFxi{}}~
    = when validating -- (test_resources/diagnostics/test65.yr:(6,7))
    ~\lstbox{\textSFii{}}~~\lstbox{\textSFx{}}~~\lstbox{\textSFx{}}~~\lstbox{\textSFx{}}~~\lstbox{\textSFx{}}~~\lstbox{\textSFx{}}~~\lstbox{\textSFx{}}~~\lstbox{\textSFx{}}~
Error[E4104] : symbol test65::B has errors
 --> test_resources/diagnostics/test65.yr:(11,27)
11  ~\lstbox{\textSFxi{}}~ fn main () { let _ = copy B (); }
    ~\lstbox{\textSFv{}}~                           ^
 --> test_resources/diagnostics/test65.yr:(3,19)
 3  ~\lstbox{\textSFxi{}}~     @final pub fn foo (self) {}
    ~\lstbox{\textSFv{}}~                   ---  error: cannot override final method test65::A::foo ()-> void
    ~\lstbox{\textSFxi{}}~
    = when validating -- (test_resources/diagnostics/test65.yr:(6,7))
    = when validating test65::main -- (test_resources/diagnostics/test65.yr:(11,4))
    ~\lstbox{\textSFii{}}~~\lstbox{\textSFx{}}~~\lstbox{\textSFx{}}~~\lstbox{\textSFx{}}~~\lstbox{\textSFx{}}~~\lstbox{\textSFx{}}~~\lstbox{\textSFx{}}~~\lstbox{\textSFx{}}~
\end{lstlisting}

\subsubsection*{How to fix it}

Remove \tokennolst{final} from the ancestor method, or give the new method a different name.

\errorcode{E4184}{the return type of the overriding method \errhole{} is not compatible with the return type of the ancestor method \errhole{}}
\label{sec:codes:e4184}

Raised during \textbf{validation}, from \texttt{ValidateErrorMessage::\allowbreak{}OVERRIDE\_\allowbreak{}INCOMPATIBLE\_\allowbreak{}RETURN\_\allowbreak{}TYPE} (\textit{src/\allowbreak{}ymirc/\allowbreak{}semantic/\allowbreak{}validator/\allowbreak{}errors.yr}).

\subsubsection*{What it means}

An overriding method returns a type the ancestor method's return type does not accept. A caller holding the base type has to be able to use the result.

\subsubsection*{Example}

\textit{test\_\allowbreak{}resources/\allowbreak{}implement/\allowbreak{}test8.yr}:

%% check: skip
\begin{lstlisting}[style=coloredverbatimError]
trait X {T} {
    cte if template!{T}{U of i32} {
        pub fn foo (self)-> i32;
    } else {
        pub fn foo (self)-> f32;
    }
}

class Z {
    impl X!{i32} {
        pub over foo (self)-> f32 {
            1.0f
        }
    }
}

class Y {
    impl X!{f32} {
        pub over foo (self)-> i32 {
            1
        }
    }
}
\end{lstlisting}

\begin{lstlisting}[style=bashVerb, breaklines=true, escapechar=~]
Error[E4184] : the return type of the overriding method impl::foo ()-> f32 is not compatible with the return type of the ancestor method test8::X!{i32}::X::foo ()-> i32
 --> test_resources/implement/test8.yr:(11,18)
11  ~\lstbox{\textSFxi{}}~         pub over foo (self)-> f32 {
    ~\lstbox{\textSFv{}}~                  ^^^
    ~\lstbox{\textSFxi{}}~ Error[E4095] : incompatible types i32 and f32
    ~\lstbox{\textSFxi{}}~  --> test_resources/implement/test8.yr:(3,29)
    ~\lstbox{\textSFxi{}}~  3  ~\lstbox{\textSFxi{}}~         pub fn foo (self)-> i32;
    ~\lstbox{\textSFxi{}}~     ~\lstbox{\textSFv{}}~                             ^^^
    ~\lstbox{\textSFxi{}}~    ...
    ~\lstbox{\textSFxi{}}~  --> test_resources/implement/test8.yr:(11,31)
    ~\lstbox{\textSFxi{}}~ 11  ~\lstbox{\textSFxi{}}~         pub over foo (self)-> f32 {
    ~\lstbox{\textSFxi{}}~     ~\lstbox{\textSFv{}}~                               ^^^
    ~\lstbox{\textSFxi{}}~     ~\lstbox{\textSFii{}}~~\lstbox{\textSFx{}}~~\lstbox{\textSFx{}}~~\lstbox{\textSFx{}}~~\lstbox{\textSFx{}}~~\lstbox{\textSFx{}}~~\lstbox{\textSFx{}}~~\lstbox{\textSFx{}}~
    ~\lstbox{\textSFii{}}~~\lstbox{\textSFx{}}~~\lstbox{\textSFx{}}~~\lstbox{\textSFx{}}~~\lstbox{\textSFx{}}~~\lstbox{\textSFx{}}~~\lstbox{\textSFvii{}}~~\lstbox{\textSFx{}}~
\end{lstlisting}

\subsubsection*{How to fix it}

Return the ancestor's type, or a type that converts to it.

\errorcode{E4185}{the protection \errhole{} of the overriding method \errhole{} does not match the definition in the ancestor class \errhole{}}
\label{sec:codes:e4185}

Raised during \textbf{validation}, from \texttt{ValidateErrorMessage::\allowbreak{}OVERRIDE\_\allowbreak{}MISMATCH\_\allowbreak{}PROTECTION} (\textit{src/\allowbreak{}ymirc/\allowbreak{}semantic/\allowbreak{}validator/\allowbreak{}errors.yr}).

\subsubsection*{What it means}

An overriding method has a different protection from the one it overrides. A caller reaching the method through the base class would otherwise see a protection that does not hold for the actual instance.

\subsubsection*{Example}

\textit{test\_\allowbreak{}resources/\allowbreak{}diagnostics/\allowbreak{}test66.yr}:

%% check: skip
\begin{lstlisting}[style=coloredverbatimError]
class A {
    pub self () {}
    pub fn foo (self) {}
}

class B over A {
    pub self () {}
    prv over foo (self) {}
}

fn main () { let _ = copy B (); }
\end{lstlisting}

\begin{lstlisting}[style=bashVerb, breaklines=true, escapechar=~]
Error[E4185] : the protection prv of the overriding method test66::B::foo ()-> void does not match the definition in the ancestor class pub
 --> test_resources/diagnostics/test66.yr:(3,12)
 3  ~\lstbox{\textSFxi{}}~     pub fn foo (self) {}
    ~\lstbox{\textSFv{}}~            ^^^
 --> test_resources/diagnostics/test66.yr:(8,14)
 8  ~\lstbox{\textSFxi{}}~     prv over foo (self) {}
    ~\lstbox{\textSFv{}}~              ---
    ~\lstbox{\textSFxi{}}~
    = when validating -- (test_resources/diagnostics/test66.yr:(6,7))
    ~\lstbox{\textSFii{}}~~\lstbox{\textSFx{}}~~\lstbox{\textSFx{}}~~\lstbox{\textSFx{}}~~\lstbox{\textSFx{}}~~\lstbox{\textSFx{}}~~\lstbox{\textSFx{}}~~\lstbox{\textSFx{}}~
Error[E4104] : symbol test66::B has errors
 --> test_resources/diagnostics/test66.yr:(11,27)
11  ~\lstbox{\textSFxi{}}~ fn main () { let _ = copy B (); }
    ~\lstbox{\textSFv{}}~                           ^
 --> test_resources/diagnostics/test66.yr:(3,12)
 3  ~\lstbox{\textSFxi{}}~     pub fn foo (self) {}
    ~\lstbox{\textSFv{}}~            ---  error: the protection prv of the overriding method test66::B::foo ()-> void does not match the definition in the ancestor class pub
    ~\lstbox{\textSFxi{}}~
    = when validating -- (test_resources/diagnostics/test66.yr:(6,7))
    = when validating test66::main -- (test_resources/diagnostics/test66.yr:(11,4))
    ~\lstbox{\textSFii{}}~~\lstbox{\textSFx{}}~~\lstbox{\textSFx{}}~~\lstbox{\textSFx{}}~~\lstbox{\textSFx{}}~~\lstbox{\textSFx{}}~~\lstbox{\textSFx{}}~~\lstbox{\textSFx{}}~
\end{lstlisting}

\subsubsection*{How to fix it}

Give the overriding method the same protection as the ancestor one.

\errorcode{E4186}{the throwers of the overriding method \errhole{} are not compatible with the throwers of the ancestor method \errhole{}}
\label{sec:codes:e4186}

Raised during \textbf{validation}, from \texttt{ValidateErrorMessage::\allowbreak{}OVERRIDE\_\allowbreak{}MISMATCH\_\allowbreak{}THROWERS} (\textit{src/\allowbreak{}ymirc/\allowbreak{}semantic/\allowbreak{}validator/\allowbreak{}errors.yr}).

\subsubsection*{What it means}

An overriding method can throw an exception the ancestor method does not declare. A caller holding the base type would not know to handle it.

\subsubsection*{Example}

\textit{test\_\allowbreak{}resources/\allowbreak{}diagnostics/\allowbreak{}test67.yr}:

%% check: skip
\begin{lstlisting}[style=coloredverbatimError]
class A {
    pub self () {}
    pub fn foo (self) {}
}

class B over A {
    pub self () {}
    pub over foo (self) throws AssertError {
        throw copy AssertError ("x");
    }
}

fn main () { let _ = copy B (); }
\end{lstlisting}

\begin{lstlisting}[style=bashVerb, breaklines=true, escapechar=~]
Error[E4186] : the throwers of the overriding method test67::B::foo ()-> void are not compatible with the throwers of the ancestor method test67::A::foo ()-> void
 --> test_resources/diagnostics/test67.yr:(3,12)
 3  ~\lstbox{\textSFxi{}}~     pub fn foo (self) {}
    ~\lstbox{\textSFv{}}~            ^^^
 --> test_resources/diagnostics/test67.yr:(8,32)
 8  ~\lstbox{\textSFxi{}}~     pub over foo (self) throws AssertError {
    ~\lstbox{\textSFv{}}~                                -----------  error: the method might throw an exception of type core::exception::assertion::AssertError, but that is not covered by the method of the ancestor class
    ~\lstbox{\textSFxi{}}~
    = when validating -- (test_resources/diagnostics/test67.yr:(6,7))
    ~\lstbox{\textSFii{}}~~\lstbox{\textSFx{}}~~\lstbox{\textSFx{}}~~\lstbox{\textSFx{}}~~\lstbox{\textSFx{}}~~\lstbox{\textSFx{}}~~\lstbox{\textSFx{}}~~\lstbox{\textSFx{}}~
Error[E4104] : symbol test67::B has errors
 --> test_resources/diagnostics/test67.yr:(13,27)
13  ~\lstbox{\textSFxi{}}~ fn main () { let _ = copy B (); }
    ~\lstbox{\textSFv{}}~                           ^
 --> test_resources/diagnostics/test67.yr:(3,12)
 3  ~\lstbox{\textSFxi{}}~     pub fn foo (self) {}
    ~\lstbox{\textSFv{}}~            ---  error: the throwers of the overriding method test67::B::foo ()-> void are not compatible with the throwers of the ancestor method test67::A::foo ()-> void
   ...
 8  ~\lstbox{\textSFxi{}}~     pub over foo (self) throws AssertError {
    ~\lstbox{\textSFv{}}~                                -----------  error: the method might throw an exception of type core::exception::assertion::AssertError, but that is not covered by the method of the ancestor class
    ~\lstbox{\textSFxi{}}~
    = when validating -- (test_resources/diagnostics/test67.yr:(6,7))
    = when validating test67::main -- (test_resources/diagnostics/test67.yr:(13,4))
    ~\lstbox{\textSFii{}}~~\lstbox{\textSFx{}}~~\lstbox{\textSFx{}}~~\lstbox{\textSFx{}}~~\lstbox{\textSFx{}}~~\lstbox{\textSFx{}}~~\lstbox{\textSFx{}}~~\lstbox{\textSFx{}}~
\end{lstlisting}

\subsubsection*{How to fix it}

Catch the exception inside the overriding method, or declare it on the ancestor method.

\errorcode{E4187}{trait method \errhole{} was already overridden}
\label{sec:codes:e4187}

Raised during \textbf{validation}, from \texttt{ValidateErrorMessage::\allowbreak{}OVERRIDE\_\allowbreak{}MULTIPLE\_\allowbreak{}TIMES\_\allowbreak{}TRAIT} (\textit{src/\allowbreak{}ymirc/\allowbreak{}semantic/\allowbreak{}validator/\allowbreak{}errors.yr}).

\subsubsection*{What it means}

A trait method is overridden twice, so two bodies claim the same slot.

\subsubsection*{Example}

\textit{test\_\allowbreak{}resources/\allowbreak{}diagnostics/\allowbreak{}test119.yr}:

%% check: skip
\begin{lstlisting}[style=coloredverbatimError]
// a trait method overridden twice in the same impl block
trait Shape {
    pub fn area (self)-> i32;
}

class Square {
    pub self () {}

    impl Shape {
        pub over area (self)-> i32 { 1 }
        pub over area (self)-> i32 { 2 }
    }
}

fn main () {
    let _ = copy Square ();
}
\end{lstlisting}

\begin{lstlisting}[style=bashVerb, breaklines=true, escapechar=~]
Error[E4187] : trait method impl::area ()-> i32 was already overridden
 --> test_resources/diagnostics/test119.yr:(10,18)
10  ~\lstbox{\textSFxi{}}~         pub over area (self)-> i32 { 1 }
    ~\lstbox{\textSFv{}}~                  ^^^^
11  ~\lstbox{\textSFxi{}}~         pub over area (self)-> i32 { 2 }
    ~\lstbox{\textSFv{}}~                  ----
    ~\lstbox{\textSFxi{}}~
    = when validating -- (test_resources/diagnostics/test119.yr:(6,7))
    ~\lstbox{\textSFii{}}~~\lstbox{\textSFx{}}~~\lstbox{\textSFx{}}~~\lstbox{\textSFx{}}~~\lstbox{\textSFx{}}~~\lstbox{\textSFx{}}~~\lstbox{\textSFx{}}~~\lstbox{\textSFx{}}~
Error[E4104] : symbol test119::Square has errors
 --> test_resources/diagnostics/test119.yr:(16,18)
16  ~\lstbox{\textSFxi{}}~     let _ = copy Square ();
    ~\lstbox{\textSFv{}}~                  ^^^^^^
 --> test_resources/diagnostics/test119.yr:(10,18)
10  ~\lstbox{\textSFxi{}}~         pub over area (self)-> i32 { 1 }
    ~\lstbox{\textSFv{}}~                  ----  error: trait method impl::area ()-> i32 was already overridden
    ~\lstbox{\textSFxi{}}~
    = when validating -- (test_resources/diagnostics/test119.yr:(6,7))
    = when validating test119::main -- (test_resources/diagnostics/test119.yr:(15,4))
    ~\lstbox{\textSFii{}}~~\lstbox{\textSFx{}}~~\lstbox{\textSFx{}}~~\lstbox{\textSFx{}}~~\lstbox{\textSFx{}}~~\lstbox{\textSFx{}}~~\lstbox{\textSFx{}}~~\lstbox{\textSFx{}}~
\end{lstlisting}

\subsubsection*{How to fix it}

Keep one of them.

\errorcode{E4188}{cannot override a non trait method \errhole{} with \errhole{} inside impl}
\label{sec:codes:e4188}

Raised during \textbf{validation}, from \texttt{ValidateErrorMessage::\allowbreak{}OVERRIDE\_\allowbreak{}NON\_\allowbreak{}TRAIT\_\allowbreak{}INSIDE} (\textit{src/\allowbreak{}ymirc/\allowbreak{}semantic/\allowbreak{}validator/\allowbreak{}errors.yr}).

\subsubsection*{What it means}

A method inside an \tokennolst{impl} block overrides a method that is not part of the trait being implemented. An \tokennolst{impl} block provides the trait, and nothing else.

\subsubsection*{Example}

\textit{test\_\allowbreak{}resources/\allowbreak{}diagnostics/\allowbreak{}test120.yr}:

%% check: skip
\begin{lstlisting}[style=coloredverbatimError]
// an impl block overriding a method that is not part of the trait
trait Shape {
    pub fn area (self)-> i32;
}

class Square {
    pub self () {}
    pub fn side (self)-> i32 { 1 }

    impl Shape {
        pub over area (self)-> i32 { 1 }
        pub over side (self)-> i32 { 2 }
    }
}

fn main () {
    let _ = copy Square ();
}
\end{lstlisting}

\begin{lstlisting}[style=bashVerb, breaklines=true, escapechar=~]
Error[E4188] : cannot override a non trait method test120::Square::side ()-> i32 with impl::side ()-> i32 inside impl
 --> test_resources/diagnostics/test120.yr:(8,12)
 8  ~\lstbox{\textSFxi{}}~     pub fn side (self)-> i32 { 1 }
    ~\lstbox{\textSFv{}}~            ^^^^
 --> test_resources/diagnostics/test120.yr:(12,18)
12  ~\lstbox{\textSFxi{}}~         pub over side (self)-> i32 { 2 }
    ~\lstbox{\textSFv{}}~                  ----
    ~\lstbox{\textSFxi{}}~
    = when validating -- (test_resources/diagnostics/test120.yr:(6,7))
    ~\lstbox{\textSFii{}}~~\lstbox{\textSFx{}}~~\lstbox{\textSFx{}}~~\lstbox{\textSFx{}}~~\lstbox{\textSFx{}}~~\lstbox{\textSFx{}}~~\lstbox{\textSFx{}}~~\lstbox{\textSFx{}}~
Error[E4104] : symbol test120::Square has errors
 --> test_resources/diagnostics/test120.yr:(17,18)
17  ~\lstbox{\textSFxi{}}~     let _ = copy Square ();
    ~\lstbox{\textSFv{}}~                  ^^^^^^
 --> test_resources/diagnostics/test120.yr:(8,12)
 8  ~\lstbox{\textSFxi{}}~     pub fn side (self)-> i32 { 1 }
    ~\lstbox{\textSFv{}}~            ----  error: cannot override a non trait method test120::Square::side ()-> i32 with impl::side ()-> i32 inside impl
    ~\lstbox{\textSFxi{}}~
    = when validating -- (test_resources/diagnostics/test120.yr:(6,7))
    = when validating test120::main -- (test_resources/diagnostics/test120.yr:(16,4))
    ~\lstbox{\textSFii{}}~~\lstbox{\textSFx{}}~~\lstbox{\textSFx{}}~~\lstbox{\textSFx{}}~~\lstbox{\textSFx{}}~~\lstbox{\textSFx{}}~~\lstbox{\textSFx{}}~~\lstbox{\textSFx{}}~
\end{lstlisting}

\subsubsection*{How to fix it}

Move the method out of the \tokennolst{impl} block into the type body.

\errorcode{E4189}{method \errhole{} overrides nothing}
\label{sec:codes:e4189}

Raised during \textbf{validation}, from \texttt{ValidateErrorMessage::\allowbreak{}OVERRIDE\_\allowbreak{}NOTHING} (\textit{src/\allowbreak{}ymirc/\allowbreak{}semantic/\allowbreak{}validator/\allowbreak{}errors.yr}).

\subsubsection*{What it means}

A method is marked \tokennolst{override} and no ancestor method matches it, cf. \hyperref[sec:codes:retired]{E3011}.

\subsubsection*{Example}

\textit{test\_\allowbreak{}resources/\allowbreak{}diagnostics/\allowbreak{}test68.yr}:

%% check: skip
\begin{lstlisting}[style=coloredverbatimError]
class A {
    pub self () {}
    pub over foo (self) {}
}

fn main () { let _ = copy A (); }
\end{lstlisting}

\begin{lstlisting}[style=bashVerb, breaklines=true, escapechar=~]
Error[E4189] : method test68::A::foo overrides nothing
 --> test_resources/diagnostics/test68.yr:(3,14)
 3  ~\lstbox{\textSFxi{}}~     pub over foo (self) {}
    ~\lstbox{\textSFv{}}~              ^^^
    ~\lstbox{\textSFxi{}}~
    = when validating -- (test_resources/diagnostics/test68.yr:(1,7))
    ~\lstbox{\textSFii{}}~~\lstbox{\textSFx{}}~~\lstbox{\textSFx{}}~~\lstbox{\textSFx{}}~~\lstbox{\textSFx{}}~~\lstbox{\textSFx{}}~~\lstbox{\textSFx{}}~~\lstbox{\textSFx{}}~
Error[E4104] : symbol test68::A has errors
 --> test_resources/diagnostics/test68.yr:(6,27)
 6  ~\lstbox{\textSFxi{}}~ fn main () { let _ = copy A (); }
    ~\lstbox{\textSFv{}}~                           ^
 --> test_resources/diagnostics/test68.yr:(3,14)
 3  ~\lstbox{\textSFxi{}}~     pub over foo (self) {}
    ~\lstbox{\textSFv{}}~              ---  error: method test68::A::foo overrides nothing
    ~\lstbox{\textSFxi{}}~
    = when validating -- (test_resources/diagnostics/test68.yr:(1,7))
    = when validating test68::main -- (test_resources/diagnostics/test68.yr:(6,4))
    ~\lstbox{\textSFii{}}~~\lstbox{\textSFx{}}~~\lstbox{\textSFx{}}~~\lstbox{\textSFx{}}~~\lstbox{\textSFx{}}~~\lstbox{\textSFx{}}~~\lstbox{\textSFx{}}~~\lstbox{\textSFx{}}~
\end{lstlisting}

\subsubsection*{How to fix it}

Compare the signature with the ancestor method, or drop \tokennolst{override}.

\errorcode{E4190}{field method \errhole{} cannot override the non field method \errhole{}}
\label{sec:codes:e4190}

Raised during \textbf{validation}, from \texttt{ValidateErrorMessage::\allowbreak{}OVERRIDE\_\allowbreak{}NO\_\allowbreak{}FIELD\_\allowbreak{}BY\_\allowbreak{}FIELD} (\textit{src/\allowbreak{}ymirc/\allowbreak{}semantic/\allowbreak{}validator/\allowbreak{}errors.yr}).

\subsubsection*{What it means}

A field method tries to override an ordinary method, cf. \hyperref[sec:codes:e4182]{E4182}.

\subsubsection*{Example}

\textit{test\_\allowbreak{}resources/\allowbreak{}diagnostics/\allowbreak{}test121.yr}:

%% check: skip
\begin{lstlisting}[style=coloredverbatimError]
// a field method overriding a method that is not one
class Base {
    pub self () {}
    pub fn size (self)-> i32 { 1 }
}

class Derived over Base {
    pub self () with super () {}

    @field
    pub over size (self)-> i32 { 2 }
}

fn main () {
    let _ = copy Derived ();
}
\end{lstlisting}

\begin{lstlisting}[style=bashVerb, breaklines=true, escapechar=~]
Error[E4190] : field method test121::Derived::size ()-> i32 cannot override the non field method test121::Base::size ()-> i32
 --> test_resources/diagnostics/test121.yr:(10,6)
10  ~\lstbox{\textSFxi{}}~     @field
    ~\lstbox{\textSFv{}}~      ^^^^^
11  ~\lstbox{\textSFxi{}}~     pub over size (self)-> i32 { 2 }
    ~\lstbox{\textSFv{}}~              ----
    ~\lstbox{\textSFxi{}}~
    = when validating -- (test_resources/diagnostics/test121.yr:(7,7))
    ~\lstbox{\textSFii{}}~~\lstbox{\textSFx{}}~~\lstbox{\textSFx{}}~~\lstbox{\textSFx{}}~~\lstbox{\textSFx{}}~~\lstbox{\textSFx{}}~~\lstbox{\textSFx{}}~~\lstbox{\textSFx{}}~
Error[E4104] : symbol test121::Derived has errors
 --> test_resources/diagnostics/test121.yr:(15,18)
15  ~\lstbox{\textSFxi{}}~     let _ = copy Derived ();
    ~\lstbox{\textSFv{}}~                  ^^^^^^^
 --> test_resources/diagnostics/test121.yr:(10,6)
10  ~\lstbox{\textSFxi{}}~     @field
    ~\lstbox{\textSFv{}}~      -----  error: field method test121::Derived::size ()-> i32 cannot override the non field method test121::Base::size ()-> i32
    ~\lstbox{\textSFxi{}}~
    = when validating -- (test_resources/diagnostics/test121.yr:(7,7))
    = when validating test121::main -- (test_resources/diagnostics/test121.yr:(14,4))
    ~\lstbox{\textSFii{}}~~\lstbox{\textSFx{}}~~\lstbox{\textSFx{}}~~\lstbox{\textSFx{}}~~\lstbox{\textSFx{}}~~\lstbox{\textSFx{}}~~\lstbox{\textSFx{}}~~\lstbox{\textSFx{}}~
\end{lstlisting}

\subsubsection*{How to fix it}

Declare the overriding method as an ordinary method.

\errorcode{E4192}{cannot override trait method \errhole{} with \errhole{} outside impl}
\label{sec:codes:e4192}

Raised during \textbf{validation}, from \texttt{ValidateErrorMessage::\allowbreak{}OVERRIDE\_\allowbreak{}TRAIT\_\allowbreak{}OUTSIDE} (\textit{src/\allowbreak{}ymirc/\allowbreak{}semantic/\allowbreak{}validator/\allowbreak{}errors.yr}).

\subsubsection*{What it means}

A trait method is overridden outside an \tokennolst{impl} block. The trait states which methods it provides, and their bodies belong in the block implementing it.

\subsubsection*{Example}

\textit{test\_\allowbreak{}resources/\allowbreak{}implement/\allowbreak{}test4.yr}:

%% check: skip
\begin{lstlisting}[style=coloredverbatimError]
trait X {
    pub fn foo (self);
}

@abstract
class A {
    pub self () {}
    impl X ;
}

class B over A {
    pub over foo (self) {}
}
\end{lstlisting}

\begin{lstlisting}[style=bashVerb, breaklines=true, escapechar=~]
Error[E4192] : cannot override trait method X::foo ()-> void with test4::B::foo ()-> void outside impl
 --> test_resources/implement/test4.yr:(2,12)
 2  ~\lstbox{\textSFxi{}}~     pub fn foo (self);
    ~\lstbox{\textSFv{}}~            ^^^
   ...
 --> test_resources/implement/test4.yr:(12,14)
12  ~\lstbox{\textSFxi{}}~     pub over foo (self) {}
    ~\lstbox{\textSFv{}}~              ^^^
    ~\lstbox{\textSFii{}}~~\lstbox{\textSFx{}}~~\lstbox{\textSFx{}}~~\lstbox{\textSFx{}}~~\lstbox{\textSFx{}}~~\lstbox{\textSFx{}}~~\lstbox{\textSFx{}}~~\lstbox{\textSFx{}}~
\end{lstlisting}

\subsubsection*{How to fix it}

Move the method into the \tokennolst{impl} block for that trait.

\errorcode{E4196}{ancestor cycle found when validating \errhole{}}
\label{sec:codes:e4196}

Raised during \textbf{validation}, from \texttt{ValidateErrorMessage::\allowbreak{}RECURSIVE\_\allowbreak{}ANCESTOR} (\textit{src/\allowbreak{}ymirc/\allowbreak{}semantic/\allowbreak{}validator/\allowbreak{}errors.yr}).

\subsubsection*{What it means}

A class hierarchy loops back on itself: a class is, directly or transitively, its own ancestor.

\subsubsection*{Example}

\textit{test\_\allowbreak{}resources/\allowbreak{}diagnostics/\allowbreak{}test136.yr}:

%% check: skip
\begin{lstlisting}[style=coloredverbatimError]
// a class that is its own ancestor
class Node over Node {
    pub self () {}
}

fn main () {
    let _ = copy Node ();
}
\end{lstlisting}

\begin{lstlisting}[style=bashVerb, breaklines=true, escapechar=~]
Error[E4104] : symbol test136::Node has errors
 --> test_resources/diagnostics/test136.yr:(2,17)
 2  ~\lstbox{\textSFxi{}}~ class Node over Node {
    ~\lstbox{\textSFv{}}~       ----      ^^^^  error: ancestor cycle found when validating test136::Node
    ~\lstbox{\textSFxi{}}~
    = when validating -- (test_resources/diagnostics/test136.yr:(2,7))
    = when validating -- (test_resources/diagnostics/test136.yr:(2,7))
    ~\lstbox{\textSFii{}}~~\lstbox{\textSFx{}}~~\lstbox{\textSFx{}}~~\lstbox{\textSFx{}}~~\lstbox{\textSFx{}}~~\lstbox{\textSFx{}}~~\lstbox{\textSFx{}}~~\lstbox{\textSFx{}}~
Error[E4104] : symbol test136::Node has errors
 --> test_resources/diagnostics/test136.yr:(7,18)
 7  ~\lstbox{\textSFxi{}}~     let _ = copy Node ();
    ~\lstbox{\textSFv{}}~                  ^^^^
 --> test_resources/diagnostics/test136.yr:(2,7)
 2  ~\lstbox{\textSFxi{}}~ class Node over Node {
    ~\lstbox{\textSFv{}}~       ----      ----
    ~\lstbox{\textSFxi{}}~          ~\lstbox{\textSFii{}}~~\lstbox{\textSFx{}}~ error: ancestor cycle found when validating test136::Node
    ~\lstbox{\textSFxi{}}~                    ~\lstbox{\textSFii{}}~~\lstbox{\textSFx{}}~ error: symbol test136::Node has errors
    ~\lstbox{\textSFxi{}}~
    = when validating -- (test_resources/diagnostics/test136.yr:(2,7))
    = when validating -- (test_resources/diagnostics/test136.yr:(2,7))
    = when validating test136::main -- (test_resources/diagnostics/test136.yr:(6,4))
    ~\lstbox{\textSFii{}}~~\lstbox{\textSFx{}}~~\lstbox{\textSFx{}}~~\lstbox{\textSFx{}}~~\lstbox{\textSFx{}}~~\lstbox{\textSFx{}}~~\lstbox{\textSFx{}}~~\lstbox{\textSFx{}}~
\end{lstlisting}

\subsubsection*{How to fix it}

Break the cycle in the inheritance chain.

\errorcode{E4233}{no constructor is callable with the parameters \errhole{}}
\label{sec:codes:e4233}

Raised during \textbf{validation}, from \texttt{ValidateErrorMessage::\allowbreak{}UNDEFINED\_\allowbreak{}REDIRECT\_\allowbreak{}CALL\_\allowbreak{}CTOR} (\textit{src/\allowbreak{}ymirc/\allowbreak{}semantic/\allowbreak{}validator/\allowbreak{}errors.yr}).

\subsubsection*{What it means}

A constructor redirection names no constructor callable with the arguments given.

\subsubsection*{Example}

\textit{test\_\allowbreak{}resources/\allowbreak{}diagnostics/\allowbreak{}test76.yr}:

%% check: skip
\begin{lstlisting}[style=coloredverbatimError]
class Foo {
    let x: i32;
    pub self () with self (1, 2) {}
    pub self (x: i32) with x = x {}
}

fn main () { let _ = copy Foo (); }
\end{lstlisting}

\begin{lstlisting}[style=bashVerb, breaklines=true, escapechar=~]
Error[E4233] : no constructor is callable with the parameters {1, 2}
 --> test_resources/diagnostics/test76.yr:(3,22)
 3  ~\lstbox{\textSFxi{}}~     pub self () with self (1, 2) {}
    ~\lstbox{\textSFv{}}~                      ^^^^
    ~\lstbox{\textSFxi{}}~ Error[E4226] : the call operator is not defined for new and {1: i32, 2: i32}
    ~\lstbox{\textSFxi{}}~  --> test_resources/diagnostics/test76.yr:(3,22)
    ~\lstbox{\textSFxi{}}~  3  ~\lstbox{\textSFxi{}}~     pub self () with self (1, 2) {}
    ~\lstbox{\textSFxi{}}~     ~\lstbox{\textSFv{}}~                      ^^^^
    ~\lstbox{\textSFxi{}}~     ~\lstbox{\textSFxi{}}~ Note : candidate test76::Foo::self ()-> mut &(mut test76::Foo) -- (test_resources/diagnostics/test76.yr:(3,9))
    ~\lstbox{\textSFxi{}}~     ~\lstbox{\textSFxi{}}~     ~\lstbox{\textSFxi{}}~ Error[E4212] : 0 parameters are expected, but 2 are provided
    ~\lstbox{\textSFxi{}}~     ~\lstbox{\textSFxi{}}~     ~\lstbox{\textSFxi{}}~  --> test_resources/diagnostics/test76.yr:(3,22)
    ~\lstbox{\textSFxi{}}~     ~\lstbox{\textSFxi{}}~     ~\lstbox{\textSFxi{}}~  3  ~\lstbox{\textSFxi{}}~     pub self () with self (1, 2) {}
    ~\lstbox{\textSFxi{}}~     ~\lstbox{\textSFxi{}}~     ~\lstbox{\textSFxi{}}~     ~\lstbox{\textSFv{}}~                      ^^^^
    ~\lstbox{\textSFxi{}}~     ~\lstbox{\textSFxi{}}~     ~\lstbox{\textSFxi{}}~     ~\lstbox{\textSFii{}}~~\lstbox{\textSFx{}}~~\lstbox{\textSFx{}}~~\lstbox{\textSFx{}}~~\lstbox{\textSFx{}}~~\lstbox{\textSFx{}}~~\lstbox{\textSFx{}}~~\lstbox{\textSFx{}}~
    ~\lstbox{\textSFxi{}}~     ~\lstbox{\textSFxi{}}~     ~\lstbox{\textSFii{}}~~\lstbox{\textSFx{}}~~\lstbox{\textSFx{}}~~\lstbox{\textSFx{}}~~\lstbox{\textSFx{}}~~\lstbox{\textSFx{}}~~\lstbox{\textSFvii{}}~~\lstbox{\textSFx{}}~
    ~\lstbox{\textSFxi{}}~     ~\lstbox{\textSFxi{}}~ Note : candidate test76::Foo::self (x: i32)-> mut &(mut test76::Foo) -- (test_resources/diagnostics/test76.yr:(4,9))
    ~\lstbox{\textSFxi{}}~     ~\lstbox{\textSFxi{}}~     ~\lstbox{\textSFxi{}}~ Error[E4212] : 1 parameters is expected, but 2 are provided
    ~\lstbox{\textSFxi{}}~     ~\lstbox{\textSFxi{}}~     ~\lstbox{\textSFxi{}}~  --> test_resources/diagnostics/test76.yr:(3,22)
    ~\lstbox{\textSFxi{}}~     ~\lstbox{\textSFxi{}}~     ~\lstbox{\textSFxi{}}~  3  ~\lstbox{\textSFxi{}}~     pub self () with self (1, 2) {}
    ~\lstbox{\textSFxi{}}~     ~\lstbox{\textSFxi{}}~     ~\lstbox{\textSFxi{}}~     ~\lstbox{\textSFv{}}~                      ^^^^
    ~\lstbox{\textSFxi{}}~     ~\lstbox{\textSFxi{}}~     ~\lstbox{\textSFxi{}}~     ~\lstbox{\textSFii{}}~~\lstbox{\textSFx{}}~~\lstbox{\textSFx{}}~~\lstbox{\textSFx{}}~~\lstbox{\textSFx{}}~~\lstbox{\textSFx{}}~~\lstbox{\textSFx{}}~~\lstbox{\textSFx{}}~
    ~\lstbox{\textSFxi{}}~     ~\lstbox{\textSFxi{}}~     ~\lstbox{\textSFii{}}~~\lstbox{\textSFx{}}~~\lstbox{\textSFx{}}~~\lstbox{\textSFx{}}~~\lstbox{\textSFx{}}~~\lstbox{\textSFx{}}~~\lstbox{\textSFvii{}}~~\lstbox{\textSFx{}}~
    ~\lstbox{\textSFxi{}}~     ~\lstbox{\textSFii{}}~~\lstbox{\textSFx{}}~~\lstbox{\textSFx{}}~~\lstbox{\textSFx{}}~~\lstbox{\textSFx{}}~~\lstbox{\textSFx{}}~~\lstbox{\textSFvii{}}~~\lstbox{\textSFx{}}~
    ~\lstbox{\textSFii{}}~~\lstbox{\textSFx{}}~~\lstbox{\textSFx{}}~~\lstbox{\textSFx{}}~~\lstbox{\textSFx{}}~~\lstbox{\textSFx{}}~~\lstbox{\textSFvii{}}~~\lstbox{\textSFx{}}~
\end{lstlisting}

\subsubsection*{How to fix it}

Match the parameters of an existing constructor; the candidates are listed under the diagnostic.

\errorcode{E4253}{the field \errhole{} has no initial value}
\label{sec:codes:e4253}

Raised during \textbf{validation}, from \texttt{ValidateErrorMessage::\allowbreak{}UNINIT\_\allowbreak{}FIELD} (\textit{src/\allowbreak{}ymirc/\allowbreak{}semantic/\allowbreak{}validator/\allowbreak{}errors.yr}).

\subsubsection*{What it means}

A field has no initial value where one is required -- a type whose values can be created without naming every field needs every field to have a default.

\subsubsection*{Example}

\textit{test\_\allowbreak{}resources/\allowbreak{}lit\_\allowbreak{}class/\allowbreak{}ctors/\allowbreak{}test2.yr}:

%% check: skip
\begin{lstlisting}[style=coloredverbatimError]
class A {
    let i : i32;
    pub self () {}
}
\end{lstlisting}

\begin{lstlisting}[style=bashVerb, breaklines=true, escapechar=~]
Error[E4253] : the field i has no initial value
 --> test_resources/lit_class/ctors/test2.yr:(2,9)
 2  ~\lstbox{\textSFxi{}}~     let i : i32;
    ~\lstbox{\textSFv{}}~         ^
    ~\lstbox{\textSFii{}}~~\lstbox{\textSFx{}}~~\lstbox{\textSFx{}}~~\lstbox{\textSFx{}}~~\lstbox{\textSFx{}}~~\lstbox{\textSFx{}}~~\lstbox{\textSFx{}}~~\lstbox{\textSFx{}}~
\end{lstlisting}

\subsubsection*{How to fix it}

Give the field an initial value, or construct the type through a constructor that sets it.

\errorcode{E4319}{method \errhole{} must be a generator to override the generator \errhole{}}
\label{sec:codes:e4319}

Raised during \textbf{validation}, from \texttt{ValidateErrorMessage::\allowbreak{}OVERRIDE\_\allowbreak{}GENERATOR\_\allowbreak{}NO\_\allowbreak{}GENERATOR} (\textit{src/\allowbreak{}ymirc/\allowbreak{}semantic/\allowbreak{}validator/\allowbreak{}errors.yr}).

\subsubsection*{What it means}

An ordinary method tries to override a generator, declared by \tokennolst{yield}. Being a generator is part of the contract of a method, even when the overriding method returns the same handle type, cf. \hyperref[sec:codes:e4320]{E4320}.

\subsubsection*{Example}

\textit{test\_\allowbreak{}resources/\allowbreak{}generators/\allowbreak{}test85.yr}:

%% check: skip
\begin{lstlisting}[style=coloredverbatimError]
// a method overriding a generator must be a generator
yield gen ()-> i32 {
    yield 1;
}

class A {
    pub self () {}

    pub yield f (self)-> i32 {
        yield 1;
    }
}

class B over A {
    pub self () {}

    pub over f (self)-> (yield-> i32) {
        gen ()
    }
}

fn main () {
    let _ = B ();
}
\end{lstlisting}

\begin{lstlisting}[style=bashVerb, breaklines=true, escapechar=~]
Error[E4319] : method test85::B::f ()-> (yield-> i32) must be a generator to override the generator test85::A::f ()-> (yield-> i32)
 --> test_resources/generators/test85.yr:(17,14)
17  ~\lstbox{\textSFxi{}}~     pub over f (self)-> (yield-> i32) {
    ~\lstbox{\textSFv{}}~              ^
 --> test_resources/generators/test85.yr:(9,15)
 9  ~\lstbox{\textSFxi{}}~     pub yield f (self)-> i32 {
    ~\lstbox{\textSFv{}}~               -
    ~\lstbox{\textSFxi{}}~
    = when validating test85::B -- (test_resources/generators/test85.yr:(14,7))
    ~\lstbox{\textSFii{}}~~\lstbox{\textSFx{}}~~\lstbox{\textSFx{}}~~\lstbox{\textSFx{}}~~\lstbox{\textSFx{}}~~\lstbox{\textSFx{}}~~\lstbox{\textSFx{}}~~\lstbox{\textSFx{}}~
Error[E4104] : symbol test85::B has errors
 --> test_resources/generators/test85.yr:(23,13)
23  ~\lstbox{\textSFxi{}}~     let _ = B ();
    ~\lstbox{\textSFv{}}~             ^
 --> test_resources/generators/test85.yr:(9,15)
 9  ~\lstbox{\textSFxi{}}~     pub yield f (self)-> i32 {
    ~\lstbox{\textSFv{}}~               -
   ...
17  ~\lstbox{\textSFxi{}}~     pub over f (self)-> (yield-> i32) {
    ~\lstbox{\textSFv{}}~              -  error: method test85::B::f ()-> (yield-> i32) must be a generator to override the generator test85::A::f ()-> (yield-> i32)
    ~\lstbox{\textSFxi{}}~
    = when validating test85::B -- (test_resources/generators/test85.yr:(14,7))
    = when validating test85::main -- (test_resources/generators/test85.yr:(22,4))
    ~\lstbox{\textSFii{}}~~\lstbox{\textSFx{}}~~\lstbox{\textSFx{}}~~\lstbox{\textSFx{}}~~\lstbox{\textSFx{}}~~\lstbox{\textSFx{}}~~\lstbox{\textSFx{}}~~\lstbox{\textSFx{}}~
\end{lstlisting}

\subsubsection*{How to fix it}

Declare the overriding method with \tokennolst{over yield}, and yield its values.

\errorcode{E4320}{generator \errhole{} cannot override the non generator method \errhole{}}
\label{sec:codes:e4320}

Raised during \textbf{validation}, from \texttt{ValidateErrorMessage::\allowbreak{}OVERRIDE\_\allowbreak{}NO\_\allowbreak{}GENERATOR\_\allowbreak{}BY\_\allowbreak{}GENERATOR} (\textit{src/\allowbreak{}ymirc/\allowbreak{}semantic/\allowbreak{}validator/\allowbreak{}errors.yr}).

\subsubsection*{What it means}

A generator, declared by \tokennolst{over yield}, tries to override an ordinary method, even one returning the handle of a generator, cf. \hyperref[sec:codes:e4319]{E4319}.

\subsubsection*{Example}

\textit{test\_\allowbreak{}resources/\allowbreak{}generators/\allowbreak{}test86.yr}:

%% check: skip
\begin{lstlisting}[style=coloredverbatimError]
// a generator cannot override a method returning the handle of a generator
yield gen ()-> i32 {
    yield 1;
}

class A {
    pub self () {}

    pub fn f (self)-> (yield-> i32) {
        gen ()
    }
}

class B over A {
    pub self () {}

    pub over yield f (self)-> i32 {
        yield 2;
    }
}

fn main () {
    let _ = B ();
}
\end{lstlisting}

\begin{lstlisting}[style=bashVerb, breaklines=true, escapechar=~]
Error[E4320] : generator test86::B::f ()-> (yield-> i32) cannot override the non generator method test86::A::f ()-> (yield-> i32)
 --> test_resources/generators/test86.yr:(17,20)
17  ~\lstbox{\textSFxi{}}~     pub over yield f (self)-> i32 {
    ~\lstbox{\textSFv{}}~                    ^
 --> test_resources/generators/test86.yr:(9,12)
 9  ~\lstbox{\textSFxi{}}~     pub fn f (self)-> (yield-> i32) {
    ~\lstbox{\textSFv{}}~            -
    ~\lstbox{\textSFxi{}}~
    = when validating test86::B -- (test_resources/generators/test86.yr:(14,7))
    ~\lstbox{\textSFii{}}~~\lstbox{\textSFx{}}~~\lstbox{\textSFx{}}~~\lstbox{\textSFx{}}~~\lstbox{\textSFx{}}~~\lstbox{\textSFx{}}~~\lstbox{\textSFx{}}~~\lstbox{\textSFx{}}~
Error[E4104] : symbol test86::B has errors
 --> test_resources/generators/test86.yr:(23,13)
23  ~\lstbox{\textSFxi{}}~     let _ = B ();
    ~\lstbox{\textSFv{}}~             ^
 --> test_resources/generators/test86.yr:(9,12)
 9  ~\lstbox{\textSFxi{}}~     pub fn f (self)-> (yield-> i32) {
    ~\lstbox{\textSFv{}}~            -
   ...
17  ~\lstbox{\textSFxi{}}~     pub over yield f (self)-> i32 {
    ~\lstbox{\textSFv{}}~                    -  error: generator test86::B::f ()-> (yield-> i32) cannot override the non generator method test86::A::f ()-> (yield-> i32)
    ~\lstbox{\textSFxi{}}~
    = when validating test86::B -- (test_resources/generators/test86.yr:(14,7))
    = when validating test86::main -- (test_resources/generators/test86.yr:(22,4))
    ~\lstbox{\textSFii{}}~~\lstbox{\textSFx{}}~~\lstbox{\textSFx{}}~~\lstbox{\textSFx{}}~~\lstbox{\textSFx{}}~~\lstbox{\textSFx{}}~~\lstbox{\textSFx{}}~~\lstbox{\textSFx{}}~
\end{lstlisting}

\subsubsection*{How to fix it}

Declare the overriding method as an ordinary method, with \tokennolst{over}.

%% Generated by tools/gen_error_codes.py from docs/errors/ of the compiler;
%% edit the compiler's pages or tools/error_themes.txt, then rerun it.

\clearpage
\section{Unions}
\label{sec:codes:unions}

\begin{longtable}{@{}l p{0.78\linewidth} r@{}}
\toprule Code & Message & Page\\ \midrule \endhead
\hyperref[sec:codes:e4028]{E4028} & type \errhole{} cannot be embedded within a union type & \pageref{sec:codes:e4028}\\
\hyperref[sec:codes:e4306]{E4306} & over is prohibited on a union of records & \pageref{sec:codes:e4306}\\
\hyperref[sec:codes:e4307]{E4307} & union \errhole{} declares the member \errhole{} twice & \pageref{sec:codes:e4307}\\
\hyperref[sec:codes:e4308]{E4308} & union \errhole{} declares no member & \pageref{sec:codes:e4308}\\
\hyperref[sec:codes:e4309]{E4309} & a union cannot be a member of another union & \pageref{sec:codes:e4309}\\
\hyperref[sec:codes:e4310]{E4310} & the member of a union must be a record or a class type, not \errhole{} & \pageref{sec:codes:e4310}\\
\hyperref[sec:codes:e4311]{E4311} & union \errhole{} mixes record and class members & \pageref{sec:codes:e4311}\\
\hyperref[sec:codes:e4312]{E4312} & union \errhole{} has no method named \errhole{} & \pageref{sec:codes:e4312}\\
\hyperref[sec:codes:e4313]{E4313} & union \errhole{} cannot be converted to \errhole{} & \pageref{sec:codes:e4313}\\
\hyperref[sec:codes:e4314]{E4314} & union \errhole{} does not list the member \errhole{} & \pageref{sec:codes:e4314}\\
\hyperref[sec:codes:e4315]{E4315} & \errhole{} is not a union type & \pageref{sec:codes:e4315}\\
\hyperref[sec:codes:e4317]{E4317} & union \errhole{} does not implement trait \errhole{} & \pageref{sec:codes:e4317}\\
\bottomrule
\end{longtable}

\errorcode{E4028}{type \errhole{} cannot be embedded within a union type}
\label{sec:codes:e4028}

Raised during \textbf{validation}, from \texttt{ValidateErrorMessage::\allowbreak{}CONTAIN\_\allowbreak{}ALIASABLE\_\allowbreak{}TYPE} (\textit{src/\allowbreak{}ymirc/\allowbreak{}semantic/\allowbreak{}validator/\allowbreak{}errors.yr}).

\subsubsection*{What it means}

A type that can be aliased -- one that borrows memory -- was placed in a union. The fields of a union share their storage, so the compiler cannot know which of them owns the memory at a given moment.

\subsubsection*{Example}

\textit{test\_\allowbreak{}resources/\allowbreak{}union/\allowbreak{}test5.yr}:

%% check: skip
\begin{lstlisting}[style=coloredverbatimError]
@overlaid
record A {
    let a : [i32];
}
\end{lstlisting}

\begin{lstlisting}[style=bashVerb, breaklines=true, escapechar=~]
Error[E4028] : type [i32] cannot be embedded within a union type
 --> test_resources/union/test5.yr:(3,5)
 3  ~\lstbox{\textSFxi{}}~     let a : [i32];
    ~\lstbox{\textSFv{}}~     ^^^     ^
    ~\lstbox{\textSFxi{}}~ Note : union cannot contain aliasable types, with the exceptions of pointer types
    ~\lstbox{\textSFxi{}}~  --> test_resources/union/test5.yr:(2,8)
    ~\lstbox{\textSFxi{}}~  2  ~\lstbox{\textSFxi{}}~ record A {
    ~\lstbox{\textSFxi{}}~     ~\lstbox{\textSFv{}}~        ^
    ~\lstbox{\textSFxi{}}~     ~\lstbox{\textSFii{}}~~\lstbox{\textSFx{}}~~\lstbox{\textSFx{}}~~\lstbox{\textSFx{}}~~\lstbox{\textSFx{}}~~\lstbox{\textSFx{}}~~\lstbox{\textSFx{}}~~\lstbox{\textSFx{}}~
    ~\lstbox{\textSFii{}}~~\lstbox{\textSFx{}}~~\lstbox{\textSFx{}}~~\lstbox{\textSFx{}}~~\lstbox{\textSFx{}}~~\lstbox{\textSFx{}}~~\lstbox{\textSFvii{}}~~\lstbox{\textSFx{}}~
\end{lstlisting}

\subsubsection*{How to fix it}

Store the borrowing type outside the union, and keep only plain values in it.

\errorcode{E4306}{over is prohibited on a union of records}
\label{sec:codes:e4306}

Raised during \textbf{validation}, from \texttt{ValidateErrorMessage::\allowbreak{}UNION\_\allowbreak{}BASE\_\allowbreak{}ON\_\allowbreak{}RECORD} (\textit{src/\allowbreak{}ymirc/\allowbreak{}semantic/\allowbreak{}validator/\allowbreak{}errors.yr}).

\subsubsection*{What it means}

A union of records is perfectly legal -- what is rejected here is the \tokennolst{over} on it. \tokennolst{over B} states that a vtable suitable for every member already exists, so that a method of \tokennolst{B} can be called on the union whichever member it holds. Records have no vtable and no inheritance, so there is nothing for a union of records to declare a base of. The note names the member that makes it one.

A union of records still gets the methods its members have in common: the compiler synthesizes a table of the methods every member declares with an identical signature, without a base being named.

\subsubsection*{Example}

\textit{test\_\allowbreak{}resources/\allowbreak{}type\_\allowbreak{}union/\allowbreak{}test10.yr}:

%% check: skip
\begin{lstlisting}[style=coloredverbatimError]
record A {
    let a : i32;
}

record B {
    let b : i32;
}

class Base {
    pub self () {}
}

union U over Base
| A
| B;
\end{lstlisting}

\begin{lstlisting}[style=bashVerb, breaklines=true, escapechar=~]
Error[E4306] : over is prohibited on a union of records
 --> test_resources/type_union/test10.yr:(15,14)
15  ~\lstbox{\textSFxi{}}~ union U over Base
    ~\lstbox{\textSFv{}}~              ^^^^
16  ~\lstbox{\textSFxi{}}~ | A
    ~\lstbox{\textSFv{}}~   -  note: test10::A is a record type
    ~\lstbox{\textSFxi{}}~
    = when validating test10::U -- (test_resources/type_union/test10.yr:(15,7))
    ~\lstbox{\textSFii{}}~~\lstbox{\textSFx{}}~~\lstbox{\textSFx{}}~~\lstbox{\textSFx{}}~~\lstbox{\textSFx{}}~~\lstbox{\textSFx{}}~~\lstbox{\textSFx{}}~~\lstbox{\textSFx{}}~
\end{lstlisting}

\subsubsection*{How to fix it}

Drop the \tokennolst{over}, or make the members classes descending from the base.

\errorcode{E4307}{union \errhole{} declares the member \errhole{} twice}
\label{sec:codes:e4307}

Raised during \textbf{validation}, from \texttt{ValidateErrorMessage::\allowbreak{}UNION\_\allowbreak{}DUPLICATE\_\allowbreak{}MEMBER} (\textit{src/\allowbreak{}ymirc/\allowbreak{}semantic/\allowbreak{}validator/\allowbreak{}errors.yr}).

\subsubsection*{What it means}

A member is identified at runtime by its position in the member list, so two members naming the same type are two names for one case: nothing can ever produce the second, and the second of two arms matching it is unreachable. The duplicate is reported rather than dropped, the note pointing at the first occurrence.

\subsubsection*{Example}

\textit{test\_\allowbreak{}resources/\allowbreak{}type\_\allowbreak{}union/\allowbreak{}test7.yr}:

%% check: skip
\begin{lstlisting}[style=coloredverbatimError]
union U
| A
| B
| A;
\end{lstlisting}

\begin{lstlisting}[style=bashVerb, breaklines=true, escapechar=~]
Error[E4307] : union test7::U declares the member test7::A twice
 --> test_resources/type_union/test7.yr:(14,3)
12  ~\lstbox{\textSFxi{}}~ | A
    ~\lstbox{\textSFv{}}~   -
13  ~\lstbox{\textSFxi{}}~ | B
14  ~\lstbox{\textSFxi{}}~ | A;
    ~\lstbox{\textSFv{}}~   ^
    ~\lstbox{\textSFxi{}}~
    = when validating test7::U -- (test_resources/type_union/test7.yr:(11,7))
    ~\lstbox{\textSFii{}}~~\lstbox{\textSFx{}}~~\lstbox{\textSFx{}}~~\lstbox{\textSFx{}}~~\lstbox{\textSFx{}}~~\lstbox{\textSFx{}}~~\lstbox{\textSFx{}}~~\lstbox{\textSFx{}}~
\end{lstlisting}

\subsubsection*{How to fix it}

Remove the repeated member.

\errorcode{E4308}{union \errhole{} declares no member}
\label{sec:codes:e4308}

Raised during \textbf{validation}, from \texttt{ValidateErrorMessage::\allowbreak{}UNION\_\allowbreak{}EMPTY} (\textit{src/\allowbreak{}ymirc/\allowbreak{}semantic/\allowbreak{}validator/\allowbreak{}errors.yr}).

\subsubsection*{What it means}

A union holds exactly one of its members. With no member there is nothing to build such a value from, no arm that could match it, and no default value to fall back on, so the type is uninhabited and the declaration is rejected.

A union of a \textit{single} member is accepted: written by hand it is a redundant alias, but a template union can specialize down to one member.

\subsubsection*{Example}

\textit{test\_\allowbreak{}resources/\allowbreak{}type\_\allowbreak{}union/\allowbreak{}test6.yr}:

%% check: skip
\begin{lstlisting}[style=coloredverbatimError]
union U;
\end{lstlisting}

\begin{lstlisting}[style=bashVerb, breaklines=true, escapechar=~]
Error[E4308] : union test6::U declares no member
 --> test_resources/type_union/test6.yr:(3,7)
 3  ~\lstbox{\textSFxi{}}~ union U;
    ~\lstbox{\textSFv{}}~       ^
    ~\lstbox{\textSFxi{}}~
    = when validating test6::U -- (test_resources/type_union/test6.yr:(3,7))
    ~\lstbox{\textSFii{}}~~\lstbox{\textSFx{}}~~\lstbox{\textSFx{}}~~\lstbox{\textSFx{}}~~\lstbox{\textSFx{}}~~\lstbox{\textSFx{}}~~\lstbox{\textSFx{}}~~\lstbox{\textSFx{}}~
\end{lstlisting}

\subsubsection*{How to fix it}

List the members the union can hold.

\errorcode{E4309}{a union cannot be a member of another union}
\label{sec:codes:e4309}

Raised during \textbf{validation}, from \texttt{ValidateErrorMessage::\allowbreak{}UNION\_\allowbreak{}MEMBER\_\allowbreak{}IN\_\allowbreak{}UNION} (\textit{src/\allowbreak{}ymirc/\allowbreak{}semantic/\allowbreak{}validator/\allowbreak{}errors.yr}).

\subsubsection*{What it means}

A union listed as the member of another union is neither flattened nor nested. Flattening would make the member list of the outer union depend on a declaration written elsewhere, and a match over it would have no member to name, the inner union not being one of its cases.

\subsubsection*{Example}

\textit{test\_\allowbreak{}resources/\allowbreak{}type\_\allowbreak{}union/\allowbreak{}test8.yr}:

%% check: skip
\begin{lstlisting}[style=coloredverbatimError]
union U
| A
| B;

union V
| A
| U;
\end{lstlisting}

\begin{lstlisting}[style=bashVerb, breaklines=true, escapechar=~]
Error[E4309] : a union cannot be a member of another union
 --> test_resources/type_union/test8.yr:(17,3)
17  ~\lstbox{\textSFxi{}}~ | U;
    ~\lstbox{\textSFv{}}~   ^
 --> test_resources/type_union/test8.yr:(11,7)
11  ~\lstbox{\textSFxi{}}~ union U
    ~\lstbox{\textSFv{}}~       -
    ~\lstbox{\textSFxi{}}~
    = when validating test8::V -- (test_resources/type_union/test8.yr:(15,7))
    ~\lstbox{\textSFii{}}~~\lstbox{\textSFx{}}~~\lstbox{\textSFx{}}~~\lstbox{\textSFx{}}~~\lstbox{\textSFx{}}~~\lstbox{\textSFx{}}~~\lstbox{\textSFx{}}~~\lstbox{\textSFx{}}~
\end{lstlisting}

\subsubsection*{How to fix it}

List the members of the inner union directly.

\errorcode{E4310}{the member of a union must be a record or a class type, not \errhole{}}
\label{sec:codes:e4310}

Raised during \textbf{validation}, from \texttt{ValidateErrorMessage::\allowbreak{}UNION\_\allowbreak{}MEMBER\_\allowbreak{}NOT\_\allowbreak{}RECORD\_\allowbreak{}NOR\_\allowbreak{}CLASS} (\textit{src/\allowbreak{}ymirc/\allowbreak{}semantic/\allowbreak{}validator/\allowbreak{}errors.yr}).

\subsubsection*{What it means}

The members of a union are the named types it can hold, and only a record or a class is one. A scalar, a slice or an entity is rejected -- an entity carrying a note saying so, since it otherwise reads like a record.

\subsubsection*{Example}

\textit{test\_\allowbreak{}resources/\allowbreak{}type\_\allowbreak{}union/\allowbreak{}test12.yr}:

%% check: skip
\begin{lstlisting}[style=coloredverbatimError]
union U
| i32
| [c8];
\end{lstlisting}

\begin{lstlisting}[style=bashVerb, breaklines=true, escapechar=~]
Error[E4310] : the member of a union must be a record or a class type, not i32
 --> test_resources/type_union/test12.yr:(4,3)
 4  ~\lstbox{\textSFxi{}}~ | i32
    ~\lstbox{\textSFv{}}~   ^^^
    ~\lstbox{\textSFxi{}}~
    = when validating test12::U -- (test_resources/type_union/test12.yr:(3,7))
    ~\lstbox{\textSFii{}}~~\lstbox{\textSFx{}}~~\lstbox{\textSFx{}}~~\lstbox{\textSFx{}}~~\lstbox{\textSFx{}}~~\lstbox{\textSFx{}}~~\lstbox{\textSFx{}}~~\lstbox{\textSFx{}}~
\end{lstlisting}

\subsubsection*{How to fix it}

Wrap the value in a record, and list that record as the member.

\errorcode{E4311}{union \errhole{} mixes record and class members}
\label{sec:codes:e4311}

Raised during \textbf{validation}, from \texttt{ValidateErrorMessage::\allowbreak{}UNION\_\allowbreak{}MIXED\_\allowbreak{}MEMBERS} (\textit{src/\allowbreak{}ymirc/\allowbreak{}semantic/\allowbreak{}validator/\allowbreak{}errors.yr}).

\subsubsection*{What it means}

The members of a union share one overlay, every member sitting at the same offset. A record is held there as an inline value and a class as a pointer, so mixing the two would put the value semantics of one and the reference semantics of the other in a single slot. A union is therefore over records, or over classes, never both. The note points at the member that set the kind.

\subsubsection*{Example}

\textit{test\_\allowbreak{}resources/\allowbreak{}type\_\allowbreak{}union/\allowbreak{}test9.yr}:

%% check: skip
\begin{lstlisting}[style=coloredverbatimError]
record A {
    let a : i32;
}

class B {
    pub self () {}
}

union U
| A
| B;
\end{lstlisting}

\begin{lstlisting}[style=bashVerb, breaklines=true, escapechar=~]
Error[E4311] : union test9::U mixes record and class members
 --> test_resources/type_union/test9.yr:(13,3)
12  ~\lstbox{\textSFxi{}}~ | A
    ~\lstbox{\textSFv{}}~   -
13  ~\lstbox{\textSFxi{}}~ | B;
    ~\lstbox{\textSFv{}}~   ^
    ~\lstbox{\textSFxi{}}~
    = when validating test9::U -- (test_resources/type_union/test9.yr:(11,7))
    ~\lstbox{\textSFii{}}~~\lstbox{\textSFx{}}~~\lstbox{\textSFx{}}~~\lstbox{\textSFx{}}~~\lstbox{\textSFx{}}~~\lstbox{\textSFx{}}~~\lstbox{\textSFx{}}~~\lstbox{\textSFx{}}~
\end{lstlisting}

\subsubsection*{How to fix it}

Make every member a record, or every member a class.

\errorcode{E4312}{union \errhole{} has no method named \errhole{}}
\label{sec:codes:e4312}

Raised during \textbf{validation}, from \texttt{ValidateErrorMessage::\allowbreak{}UNION\_\allowbreak{}NO\_\allowbreak{}METHOD} (\textit{src/\allowbreak{}ymirc/\allowbreak{}semantic/\allowbreak{}validator/\allowbreak{}errors.yr}).

\subsubsection*{What it means}

A union holds one of its members, which one being known only at runtime, so a method call on it has to be dispatched whichever member it holds. What it answers to is therefore not the union of the methods of its members but their intersection: a method every member declares publicly, under the same name and with the same signature, the return type included. A union declaring a base answers to the public methods of that base instead, every member descending from it.

An operator applied to a union is such a call too, and the method it names is the one reported: \tokennolst{u [i]} calls \tokennolst{opIndex}, \tokennolst{u [i] = v} calls \tokennolst{opIndexAssign}, \tokennolst{u == v} calls \tokennolst{opEquals}, and \tokennolst{u \textless{} v} calls \tokennolst{opCmp}.

The notes name, member by member, what each declares under that name: the prototype it declares, that it declares none, or that the one it declares is private. Comparing them shows which member keeps the method out of the union.

\subsubsection*{Example}

\textit{test\_\allowbreak{}resources/\allowbreak{}type\_\allowbreak{}union/\allowbreak{}test23.yr}:

%% check: skip
\begin{lstlisting}[style=coloredverbatimError]
class A {
    pub self () {}
    pub fn foo (self)-> i32 { 1 }
}

class B {
    pub self () {}
}

union U
| A
| B;

fn bar (u : &U)-> i32 {
    u.foo ()
}
\end{lstlisting}

\begin{lstlisting}[style=bashVerb, breaklines=true, escapechar=~]
Error[E4312] : union &(test23::U) has no method named foo
 --> test_resources/type_union/test23.yr:(17,7)
17  ~\lstbox{\textSFxi{}}~     u.foo ()
    ~\lstbox{\textSFv{}}~       ^^^
 --> test_resources/type_union/test23.yr:(5,12)
 5  ~\lstbox{\textSFxi{}}~     pub fn foo (self)-> i32 { 1 }
    ~\lstbox{\textSFv{}}~            ---  note: the member test23::A declares test23::A::foo ()-> i32
   ...
14  ~\lstbox{\textSFxi{}}~ | B;
    ~\lstbox{\textSFv{}}~   -  note: the member test23::B declares no method named foo
    ~\lstbox{\textSFxi{}}~
    = when validating test23::bar -- (test_resources/type_union/test23.yr:(16,4))
    ~\lstbox{\textSFii{}}~~\lstbox{\textSFx{}}~~\lstbox{\textSFx{}}~~\lstbox{\textSFx{}}~~\lstbox{\textSFx{}}~~\lstbox{\textSFx{}}~~\lstbox{\textSFx{}}~~\lstbox{\textSFx{}}~
\end{lstlisting}

\subsubsection*{How to fix it}

Declare the method in every member, publicly and with the same signature, or reach the member through a match and call it there.

\errorcode{E4313}{union \errhole{} cannot be converted to \errhole{}}
\label{sec:codes:e4313}

Raised during \textbf{validation}, from \texttt{ValidateErrorMessage::\allowbreak{}UNION\_\allowbreak{}NARROWING} (\textit{src/\allowbreak{}ymirc/\allowbreak{}semantic/\allowbreak{}validator/\allowbreak{}errors.yr}).

\subsubsection*{What it means}

Conversion into a union is one way. A member widens into the union listing it wherever a conversion is allowed -- a variable initialization, an assignment, an argument, a returned value -- and the widening records which member it is, as the tag the union object carries.

Going back is not a conversion. The member held is only known at runtime, so recovering it is a test, and a test is what a pattern is: naming a member binds it with its own type when the tag says the union holds it, whether the pattern is written as an arm of a \tokennolst{match} or as an \tokennolst{if let}. Nothing about a union value says statically which member it is, so no conversion can produce one.

The base a union declares is no exception. \tokennolst{union U over B} says every member descends from \tokennolst{B}, which is what lets the methods of \tokennolst{B} be called on the union directly; it does not make the union an instance of \tokennolst{B}. The union object is its own aggregate -- a tag, a vtable and the overlay of the members -- and the \tokennolst{\&B} is inside that overlay, reached the same way as any other member.

\subsubsection*{Example}

\textit{test\_\allowbreak{}resources/\allowbreak{}type\_\allowbreak{}union/\allowbreak{}test37.yr}:

%% check: skip
\begin{lstlisting}[style=coloredverbatimError]
class A {
    pub self () {}
}

class B {
    pub self () {}
}

union U
| A
| B;

fn take (u : &U)-> &A {
    u
}
\end{lstlisting}

\begin{lstlisting}[style=bashVerb, breaklines=true, escapechar=~]
Error[E4313] : union &(test37::U) cannot be converted to &(test37::A)
 --> test_resources/type_union/test37.yr:(15,20)
15  ~\lstbox{\textSFxi{}}~ fn take (u : &U)-> &A {
    ~\lstbox{\textSFv{}}~                    ^
16  ~\lstbox{\textSFxi{}}~     u
    ~\lstbox{\textSFv{}}~     -
 --> test_resources/type_union/test37.yr:(11,7)
11  ~\lstbox{\textSFxi{}}~ union U
    ~\lstbox{\textSFv{}}~       -  note: the member held is only known at runtime
    ~\lstbox{\textSFxi{}}~
    = when validating test37::take -- (test_resources/type_union/test37.yr:(15,4))
    ~\lstbox{\textSFii{}}~~\lstbox{\textSFx{}}~~\lstbox{\textSFx{}}~~\lstbox{\textSFx{}}~~\lstbox{\textSFx{}}~~\lstbox{\textSFx{}}~~\lstbox{\textSFx{}}~~\lstbox{\textSFx{}}~
\end{lstlisting}

\subsubsection*{How to fix it}

Name the member in a pattern -- an arm of a \tokennolst{match} over the union, or an \tokennolst{if let} -- and work there, which is where the member has its own type.

\errorcode{E4314}{union \errhole{} does not list the member \errhole{}}
\label{sec:codes:e4314}

Raised during \textbf{validation}, from \texttt{ValidateErrorMessage::\allowbreak{}UNION\_\allowbreak{}MEMBER\_\allowbreak{}UNKNOWN} (\textit{src/\allowbreak{}ymirc/\allowbreak{}semantic/\allowbreak{}validator/\allowbreak{}errors.yr}).

\subsubsection*{What it means}

A pattern over a union names the member it tests, and the union holds one of the members it declares and nothing else. A type the union does not list can never be the one held, so the arm naming it matches nothing and is rejected rather than left as dead code.

The member is matched by identity, not by convertibility: a type that merely converts to a declared member is not that member either.

\subsubsection*{Example}

\textit{test\_\allowbreak{}resources/\allowbreak{}type\_\allowbreak{}union/\allowbreak{}test48.yr}:

%% check: skip
\begin{lstlisting}[style=coloredverbatimError]
record X {
    pub self () {}
}

record Y {
    pub self () {}
}

record Z {
    pub self () {}
}

union U
| X
| Y;

fn main () {
    let u : U = X ();
    match u {
        z : Z => { z; }
        _ => {}
    }
}
\end{lstlisting}

\begin{lstlisting}[style=bashVerb, breaklines=true, escapechar=~]
Error[E4314] : union test48::U does not list the member test48::Z
 --> test_resources/type_union/test48.yr:(22,13)
22  ~\lstbox{\textSFxi{}}~         z : Z => { z; }
    ~\lstbox{\textSFv{}}~             ^
 --> test_resources/type_union/test48.yr:(15,7)
15  ~\lstbox{\textSFxi{}}~ union U
    ~\lstbox{\textSFv{}}~       -
    ~\lstbox{\textSFxi{}}~
    = when validating pattern matching with value u of type test48::U -- (test_resources/type_union/test48.yr:(21,5))
    = when validating test48::main -- (test_resources/type_union/test48.yr:(19,4))
    ~\lstbox{\textSFii{}}~~\lstbox{\textSFx{}}~~\lstbox{\textSFx{}}~~\lstbox{\textSFx{}}~~\lstbox{\textSFx{}}~~\lstbox{\textSFx{}}~~\lstbox{\textSFx{}}~~\lstbox{\textSFx{}}~
\end{lstlisting}

\subsubsection*{How to fix it}

Name a member the union declares, or add the type to the union when it is meant to be one.

\errorcode{E4315}{\errhole{} is not a union type}
\label{sec:codes:e4315}

Raised during \textbf{validation}, from \texttt{ValidateErrorMessage::\allowbreak{}NOT\_\allowbreak{}A\_\allowbreak{}UNION} (\textit{src/\allowbreak{}ymirc/\allowbreak{}semantic/\allowbreak{}validator/\allowbreak{}errors.yr}).

\subsubsection*{What it means}

A template parameter written \tokennolst{union T} constrains the kind of type that may be bound to it, the way \tokennolst{class T} or \tokennolst{tuple T} do: it accepts a union and nothing else. The constraint is checked wherever the argument comes from -- written explicitly at the call site, or deduced from a parameter typed \tokennolst{T}.

What fails is the candidate, not the program: a template whose parameter cannot be bound is discarded, so this diagnostic is reported as the reason under the error naming the call, and the call fails only when no candidate is left.

\subsubsection*{Example}

\textit{test\_\allowbreak{}resources/\allowbreak{}type\_\allowbreak{}union/\allowbreak{}test62.yr}:

%% check: skip
\begin{lstlisting}[style=coloredverbatimError]
fn keep {union T} (u : T)-> T {
    u
}

fn main () {
    keep (12);
}
\end{lstlisting}

\begin{lstlisting}[style=bashVerb, breaklines=true, escapechar=~]
Error[E4226] : the call operator is not defined for test62::keep {union T}(u: T)-> T and {12: i32}
 --> test_resources/type_union/test62.yr:(8,10)
 8  ~\lstbox{\textSFxi{}}~     keep (12);
    ~\lstbox{\textSFv{}}~          ^
 --> test_resources/type_union/test62.yr:(3,4)
 3  ~\lstbox{\textSFxi{}}~ fn keep {union T} (u : T)-> T {
    ~\lstbox{\textSFv{}}~    ----  note: candidate test62::keep {union T}(u: T)-> T
    ~\lstbox{\textSFxi{}}~       ~\lstbox{\textSFii{}}~~\lstbox{\textSFx{}}~ error: i32 is not a union type
    ~\lstbox{\textSFxi{}}~
    = when validating test62::main -- (test_resources/type_union/test62.yr:(7,4))
    ~\lstbox{\textSFii{}}~~\lstbox{\textSFx{}}~~\lstbox{\textSFx{}}~~\lstbox{\textSFx{}}~~\lstbox{\textSFx{}}~~\lstbox{\textSFx{}}~~\lstbox{\textSFx{}}~~\lstbox{\textSFx{}}~
\end{lstlisting}

\subsubsection*{How to fix it}

Pass a union, or drop the \tokennolst{union} constraint when the parameter is meant to accept any type.

\errorcode{E4317}{union \errhole{} does not implement trait \errhole{}}
\label{sec:codes:e4317}

Raised during \textbf{validation}, from \texttt{ValidateErrorMessage::\allowbreak{}UNION\_\allowbreak{}NOT\_\allowbreak{}IMPL} (\textit{src/\allowbreak{}ymirc/\allowbreak{}semantic/\allowbreak{}validator/\allowbreak{}errors.yr}).

\subsubsection*{What it means}

A union is used where the trait named is required, and the union does not implement it.

A union implements a trait when every one of its members implements it, so that the methods of the trait can be dispatched to whichever member the union holds. A union declaring a base with \tokennolst{over} is dispatched through the vtable of that base, so it implements only the traits the base implements, even when every member implements more. The notes name the members, or the base, that do not implement the trait.

\subsubsection*{Example}

\textit{test\_\allowbreak{}resources/\allowbreak{}type\_\allowbreak{}union/\allowbreak{}test69.yr}:

%% check: skip
\begin{lstlisting}[style=coloredverbatimError]
// A union having a member that does not implement a trait does not implement it.

use std::stream;

record A {
    impl Streamable;
}

record C {}

union Y
| A
| C;

fn foo {T impl Streamable} (_ : T) {}

fn bar (y : Y) {
    foo (y);
}
\end{lstlisting}

\begin{lstlisting}[style=bashVerb, breaklines=true, escapechar=~]
Error[E4226] : the call operator is not defined for test69::foo {T impl Streamable}(_: T) and {y: test69::Y}
 --> test_resources/type_union/test69.yr:(18,9)
18  ~\lstbox{\textSFxi{}}~     foo (y);
    ~\lstbox{\textSFv{}}~         ^
 --> test_resources/type_union/test69.yr:(15,4)
15  ~\lstbox{\textSFxi{}}~ fn foo {T impl Streamable} (_ : T) {}
    ~\lstbox{\textSFv{}}~    ---  note: candidate test69::foo {T impl Streamable}(_: T)
    ~\lstbox{\textSFxi{}}~      ~\lstbox{\textSFviii{}}~~\lstbox{\textSFx{}}~ error: union test69::Y does not implement trait Streamable
    ~\lstbox{\textSFxi{}}~      ~\lstbox{\textSFii{}}~~\lstbox{\textSFx{}}~ note: the member test69::C does not implement trait Streamable
    ~\lstbox{\textSFxi{}}~
    = when validating test69::bar -- (test_resources/type_union/test69.yr:(17,4))
    ~\lstbox{\textSFii{}}~~\lstbox{\textSFx{}}~~\lstbox{\textSFx{}}~~\lstbox{\textSFx{}}~~\lstbox{\textSFx{}}~~\lstbox{\textSFx{}}~~\lstbox{\textSFx{}}~~\lstbox{\textSFx{}}~
\end{lstlisting}

\subsubsection*{How to fix it}

Implement the trait in every member of the union -- or, for a union declaring a base, in the base itself.

%% Generated by tools/gen_error_codes.py from docs/errors/ of the compiler;
%% edit the compiler's pages or tools/error_themes.txt, then rerun it.

\clearpage
\section{Exceptions and scope guards}
\label{sec:codes:errors}

The book's chapter \textit{Error handling} specifies the rules behind these codes.

\begin{longtable}{@{}l p{0.78\linewidth} r@{}}
\toprule Code & Message & Page\\ \midrule \endhead
\hyperref[sec:codes:e4003]{E4003} & exception type \errhole{} is already caught by another catching pattern of type \errhole{} & \pageref{sec:codes:e4003}\\
\hyperref[sec:codes:e4012]{E4012} & catching pattern of type \errhole{} catches no exception & \pageref{sec:codes:e4012}\\
\hyperref[sec:codes:e4013]{E4013} & the pattern of a catch scope guard must start with a variable declaration or a field deconstructor & \pageref{sec:codes:e4013}\\
\hyperref[sec:codes:e4021]{E4021} & class type expected, not the class object instance type \errhole{} (i.e. the \& is useless) & \pageref{sec:codes:e4021}\\
\hyperref[sec:codes:e4054]{E4054} & scope guard does not catch exceptions of type \errhole{} & \pageref{sec:codes:e4054}\\
\hyperref[sec:codes:e4055]{E4055} & \errhole{} guard cannot be used when guarding a scope that cannot throw & \pageref{sec:codes:e4055}\\
\hyperref[sec:codes:e4056]{E4056} & \errhole{} scope guard value must be of type void, not \errhole{} & \pageref{sec:codes:e4056}\\
\hyperref[sec:codes:e4145]{E4145} & nothing to catch & \pageref{sec:codes:e4145}\\
\hyperref[sec:codes:e4149]{E4149} & class type \errhole{} does not inherit from exception type \errhole{} & \pageref{sec:codes:e4149}\\
\hyperref[sec:codes:e4206]{E4206} & the function \errhole{} might throw an exception of type \errhole{}, but that is not declared in its prototype & \pageref{sec:codes:e4206}\\
\hyperref[sec:codes:e4207]{E4207} & the method might throw an exception of type \errhole{}, but that is not covered by the method of the ancestor class & \pageref{sec:codes:e4207}\\
\hyperref[sec:codes:e4208]{E4208} & the prototype of the function \errhole{} declares a possible throw of an exception of type \errhole{}, but the function never throws it & \pageref{sec:codes:e4208}\\
\hyperref[sec:codes:e4209]{E4209} & throw statement is not within a function & \pageref{sec:codes:e4209}\\
\hyperref[sec:codes:e4210]{E4210} & cannot throw exception inside a scope guard & \pageref{sec:codes:e4210}\\
\hyperref[sec:codes:e4211]{E4211} & rethrowing an exception of type \errhole{} inside a scope guard is forbidden & \pageref{sec:codes:e4211}\\
\bottomrule
\end{longtable}

\errorcode{E4003}{exception type \errhole{} is already caught by another catching pattern of type \errhole{}}
\label{sec:codes:e4003}

Raised during \textbf{validation}, from \texttt{ValidateErrorMessage::\allowbreak{}ALREADY\_\allowbreak{}CAUGHT} (\textit{src/\allowbreak{}ymirc/\allowbreak{}semantic/\allowbreak{}validator/\allowbreak{}errors.yr}).

\subsubsection*{What it means}

Two catch patterns of the same \tokennolst{try} catch the same exception type, and the second can never run because the first already covers it.

\subsubsection*{How to fix it}

Remove the redundant pattern, or reorder them so the more specific type comes first: a pattern catching a base class shadows every pattern for a derived one written after it.

\errorcode{E4012}{catching pattern of type \errhole{} catches no exception}
\label{sec:codes:e4012}

Raised during \textbf{validation}, from \texttt{ValidateErrorMessage::\allowbreak{}CATCH\_\allowbreak{}NOTHING} (\textit{src/\allowbreak{}ymirc/\allowbreak{}semantic/\allowbreak{}validator/\allowbreak{}errors.yr}).

\subsubsection*{What it means}

A catch pattern names a type no exception thrown inside the guarded block can have, so the pattern can never run.

\subsubsection*{Example}

\textit{test\_\allowbreak{}resources/\allowbreak{}diagnostics/\allowbreak{}test130.yr}:

%% check: skip
\begin{lstlisting}[style=coloredverbatimError]
// a catch pattern for an exception the guarded block cannot throw
class First over Exception {
    pub self () {}
}

class Second over Exception {
    pub self () {}
}

class Other over Exception {
    pub self () {}
}

fn foo (i : i32)-> i32
    throws First, Second;

fn main () {
    let _ = {
        foo (0)
    } catch {
        Other () => { 1 }
        First () => { 2 }
        Second () => { 3 }
    };
}
\end{lstlisting}

\begin{lstlisting}[style=bashVerb, breaklines=true, escapechar=~]
Error[E4012] : catching pattern of type &(test130::Other) catches no exception
 --> test_resources/diagnostics/test130.yr:(21,15)
21  ~\lstbox{\textSFxi{}}~         Other () => { 1 }
    ~\lstbox{\textSFv{}}~               ^
    ~\lstbox{\textSFxi{}}~
    = when validating catching pattern with type &(core::exception::Exception) -- (test_resources/diagnostics/test130.yr:(20,7))
    = when validating test130::main -- (test_resources/diagnostics/test130.yr:(17,4))
    ~\lstbox{\textSFii{}}~~\lstbox{\textSFx{}}~~\lstbox{\textSFx{}}~~\lstbox{\textSFx{}}~~\lstbox{\textSFx{}}~~\lstbox{\textSFx{}}~~\lstbox{\textSFx{}}~~\lstbox{\textSFx{}}~
\end{lstlisting}

\subsubsection*{How to fix it}

Remove the pattern, or widen it to a type the block can actually throw. The types a block may throw are the ones declared by the \tokennolst{throws} of everything it calls.

\errorcode{E4013}{the pattern of a catch scope guard must start with a variable declaration or a field deconstructor}
\label{sec:codes:e4013}

Raised during \textbf{validation}, from \texttt{ValidateErrorMessage::\allowbreak{}CATCH\_\allowbreak{}PATTERN\_\allowbreak{}MULT\_\allowbreak{}OR\_\allowbreak{}VDECL} (\textit{src/\allowbreak{}ymirc/\allowbreak{}semantic/\allowbreak{}validator/\allowbreak{}errors.yr}).

\subsubsection*{What it means}

The pattern of a \tokennolst{catch} does not start with a variable declaration or a field deconstructor. A catch pattern has to bind the exception before it can test anything about it.

\subsubsection*{Example}

\textit{test\_\allowbreak{}resources/\allowbreak{}diagnostics/\allowbreak{}test22.yr}:

%% check: skip
\begin{lstlisting}[style=coloredverbatimError]
fn foo () throws AssertError;

fn main () {
    { foo (); } catch {
        1 => {}
    }
}
\end{lstlisting}

\begin{lstlisting}[style=bashVerb, breaklines=true, escapechar=~]
Error[E4013] : the pattern of a catch scope guard must start with a variable declaration or a field deconstructor
 --> test_resources/diagnostics/test22.yr:(5,9)
 5  ~\lstbox{\textSFxi{}}~         1 => {}
    ~\lstbox{\textSFv{}}~         ^
    ~\lstbox{\textSFxi{}}~
    = when validating catching pattern with type &(core::exception::assertion::AssertError) -- (test_resources/diagnostics/test22.yr:(4,17))
    = when validating test22::main -- (test_resources/diagnostics/test22.yr:(3,4))
    ~\lstbox{\textSFii{}}~~\lstbox{\textSFx{}}~~\lstbox{\textSFx{}}~~\lstbox{\textSFx{}}~~\lstbox{\textSFx{}}~~\lstbox{\textSFx{}}~~\lstbox{\textSFx{}}~~\lstbox{\textSFx{}}~
\end{lstlisting}

\subsubsection*{How to fix it}

Start the pattern with a binding, as in \tokennolst{err: \&SomeException =\textgreater{} \{ ... \}}.

\errorcode{E4021}{class type expected, not the class object instance type \errhole{} (i.e. the \& is useless)}
\label{sec:codes:e4021}

Raised during \textbf{validation}, from \texttt{ValidateErrorMessage::\allowbreak{}CLASS\_\allowbreak{}THROW\_\allowbreak{}PTR} (\textit{src/\allowbreak{}ymirc/\allowbreak{}semantic/\allowbreak{}validator/\allowbreak{}errors.yr}).

\subsubsection*{What it means}

A class \textit{type} was expected and the class \textit{instance} type was given -- the \tokennolst{\&} that turns a class into an instance reference is one level too many here.

\subsubsection*{Example}

\textit{test\_\allowbreak{}resources/\allowbreak{}aka/\allowbreak{}test15.yr}:

%% check: skip
\begin{lstlisting}[style=coloredverbatimError]
class A {
    pub self () {}
}

def B : &A;
\end{lstlisting}

\begin{lstlisting}[style=bashVerb, breaklines=true, escapechar=~]
Error[E4021] : class type expected, not the class object instance type &(test15::A) (i.e. the & is useless)
 --> test_resources/aka/test15.yr:(5,9)
 5  ~\lstbox{\textSFxi{}}~ def B : &A;
    ~\lstbox{\textSFv{}}~         ^
    ~\lstbox{\textSFii{}}~~\lstbox{\textSFx{}}~~\lstbox{\textSFx{}}~~\lstbox{\textSFx{}}~~\lstbox{\textSFx{}}~~\lstbox{\textSFx{}}~~\lstbox{\textSFx{}}~~\lstbox{\textSFx{}}~
\end{lstlisting}

\subsubsection*{How to fix it}

Remove the \tokennolst{\&}.

\errorcode{E4054}{scope guard does not catch exceptions of type \errhole{}}
\label{sec:codes:e4054}

Raised during \textbf{validation}, from \texttt{ValidateErrorMessage::\allowbreak{}EXCEPTION\_\allowbreak{}NOT\_\allowbreak{}CAUGHT} (\textit{src/\allowbreak{}ymirc/\allowbreak{}semantic/\allowbreak{}validator/\allowbreak{}errors.yr}).

\subsubsection*{What it means}

A scope guard is written on a block that can throw an exception the guard does not handle, so the exception would escape past the guard.

\subsubsection*{Example}

\textit{test\_\allowbreak{}resources/\allowbreak{}diagnostics/\allowbreak{}test31.yr}:

%% check: skip
\begin{lstlisting}[style=coloredverbatimError]
class A over Exception { pub self () {} }
class B over Exception { pub self () {} }

fn foo () throws A, B;

fn main () {
    { foo (); } catch {
        A () => {}
    }
}
\end{lstlisting}

\begin{lstlisting}[style=bashVerb, breaklines=true, escapechar=~]
Error[E4054] : scope guard does not catch exceptions of type test31::B
 --> test_resources/diagnostics/test31.yr:(7,11)
 7  ~\lstbox{\textSFxi{}}~     { foo (); } catch {
    ~\lstbox{\textSFv{}}~           ^     ^^^^^
    ~\lstbox{\textSFii{}}~~\lstbox{\textSFx{}}~~\lstbox{\textSFx{}}~~\lstbox{\textSFx{}}~~\lstbox{\textSFx{}}~~\lstbox{\textSFx{}}~~\lstbox{\textSFx{}}~~\lstbox{\textSFx{}}~
\end{lstlisting}

\subsubsection*{How to fix it}

Add a catch for the reported type, or declare the enclosing function as throwing it.

\errorcode{E4055}{\errhole{} guard cannot be used when guarding a scope that cannot throw}
\label{sec:codes:e4055}

Raised during \textbf{validation}, from \texttt{ValidateErrorMessage::\allowbreak{}EXIT\_\allowbreak{}SCOPE\_\allowbreak{}NO\_\allowbreak{}THROW} (\textit{src/\allowbreak{}ymirc/\allowbreak{}semantic/\allowbreak{}validator/\allowbreak{}errors.yr}).

\subsubsection*{What it means}

A \tokennolst{success} or \tokennolst{failure} guard was written on a block that cannot throw. Such a guard distinguishes the two ways of leaving the block, and a block that cannot throw only has one.

\subsubsection*{Example}

\textit{test\_\allowbreak{}resources/\allowbreak{}diagnostics/\allowbreak{}test32.yr}:

%% check: skip
\begin{lstlisting}[style=coloredverbatimError]
fn main () {
    {
        let _ = 1;
    } success {
    }
}
\end{lstlisting}

\begin{lstlisting}[style=bashVerb, breaklines=true, escapechar=~]
Error[E4055] : success guard cannot be used when guarding a scope that cannot throw
 --> test_resources/diagnostics/test32.yr:(4,7)
 4  ~\lstbox{\textSFxi{}}~     } success {
    ~\lstbox{\textSFv{}}~       ^^^^^^^
    ~\lstbox{\textSFii{}}~~\lstbox{\textSFx{}}~~\lstbox{\textSFx{}}~~\lstbox{\textSFx{}}~~\lstbox{\textSFx{}}~~\lstbox{\textSFx{}}~~\lstbox{\textSFx{}}~~\lstbox{\textSFx{}}~
\end{lstlisting}

\subsubsection*{How to fix it}

Use an \tokennolst{exit} guard, which runs however the block is left, or drop the guard. If the block was meant to throw, the call inside it is probably already caught.

\errorcode{E4056}{\errhole{} scope guard value must be of type void, not \errhole{}}
\label{sec:codes:e4056}

Raised during \textbf{validation}, from \texttt{ValidateErrorMessage::\allowbreak{}EXIT\_\allowbreak{}SCOPE\_\allowbreak{}VALUE\_\allowbreak{}TYPE} (\textit{src/\allowbreak{}ymirc/\allowbreak{}semantic/\allowbreak{}validator/\allowbreak{}errors.yr}).

\subsubsection*{What it means}

The body of a scope guard produces a value. A guard runs while a scope is being left, and there is nothing to receive what it would produce.

\subsubsection*{Example}

\textit{test\_\allowbreak{}resources/\allowbreak{}diagnostics/\allowbreak{}test33.yr}:

%% check: skip
\begin{lstlisting}[style=coloredverbatimError]
fn foo () throws AssertError;

fn main () {
    {
        foo ();
    } exit {
        1
    }
}
\end{lstlisting}

\begin{lstlisting}[style=bashVerb, breaklines=true, escapechar=~]
Error[E4056] : exit scope guard value must be of type void, not i32
 --> test_resources/diagnostics/test33.yr:(6,12)
 6  ~\lstbox{\textSFxi{}}~     } exit {
    ~\lstbox{\textSFv{}}~            ^
    ~\lstbox{\textSFii{}}~~\lstbox{\textSFx{}}~~\lstbox{\textSFx{}}~~\lstbox{\textSFx{}}~~\lstbox{\textSFx{}}~~\lstbox{\textSFx{}}~~\lstbox{\textSFx{}}~~\lstbox{\textSFx{}}~
\end{lstlisting}

\subsubsection*{How to fix it}

Make the guard body end in a statement rather than a value -- assign the value to something, or discard it.

\errorcode{E4145}{nothing to catch}
\label{sec:codes:e4145}

Raised during \textbf{validation}, from \texttt{ValidateErrorMessage::\allowbreak{}NOTHING\_\allowbreak{}TO\_\allowbreak{}CATCH} (\textit{src/\allowbreak{}ymirc/\allowbreak{}semantic/\allowbreak{}validator/\allowbreak{}errors.yr}).

\subsubsection*{What it means}

A \tokennolst{catch} was written on a block that cannot throw anything.

\subsubsection*{Example}

\textit{test\_\allowbreak{}resources/\allowbreak{}diagnostics/\allowbreak{}test57.yr}:

%% check: skip
\begin{lstlisting}[style=coloredverbatimError]
fn main () {
    {
        let _ = 1;
    } catch {
        _: &AssertError => {}
    }
}
\end{lstlisting}

\begin{lstlisting}[style=bashVerb, breaklines=true, escapechar=~]
Error[E4145] : nothing to catch
 --> test_resources/diagnostics/test57.yr:(4,7)
 4  ~\lstbox{\textSFxi{}}~     } catch {
    ~\lstbox{\textSFv{}}~       ^^^^^
    ~\lstbox{\textSFii{}}~~\lstbox{\textSFx{}}~~\lstbox{\textSFx{}}~~\lstbox{\textSFx{}}~~\lstbox{\textSFx{}}~~\lstbox{\textSFx{}}~~\lstbox{\textSFx{}}~~\lstbox{\textSFx{}}~
\end{lstlisting}

\subsubsection*{How to fix it}

Remove the catch. If a call inside was expected to throw, check its prototype: it may not declare the exception, or an inner catch may already handle it.

\errorcode{E4149}{class type \errhole{} does not inherit from exception type \errhole{}}
\label{sec:codes:e4149}

Raised during \textbf{validation}, from \texttt{ValidateErrorMessage::\allowbreak{}NOT\_\allowbreak{}AN\_\allowbreak{}EXCEPTION\_\allowbreak{}CLASS} (\textit{src/\allowbreak{}ymirc/\allowbreak{}semantic/\allowbreak{}validator/\allowbreak{}errors.yr}).

\subsubsection*{What it means}

A class is thrown, or caught, that does not derive from the exception base class.

\subsubsection*{Example}

\textit{test\_\allowbreak{}resources/\allowbreak{}diagnostics/\allowbreak{}test58.yr}:

%% check: skip
\begin{lstlisting}[style=coloredverbatimError]
class A { pub self () {} }

fn foo () throws A;

fn main () { foo (); }
\end{lstlisting}

\begin{lstlisting}[style=bashVerb, breaklines=true, escapechar=~]
Error[E4149] : class type test58::A does not inherit from exception type core::exception::Exception
 --> test_resources/diagnostics/test58.yr:(3,11)
 3  ~\lstbox{\textSFxi{}}~ fn foo () throws A;
    ~\lstbox{\textSFv{}}~           ^^^^^^ -
    ~\lstbox{\textSFxi{}}~
    = when validating test58::foo -- (test_resources/diagnostics/test58.yr:(3,4))
    ~\lstbox{\textSFii{}}~~\lstbox{\textSFx{}}~~\lstbox{\textSFx{}}~~\lstbox{\textSFx{}}~~\lstbox{\textSFx{}}~~\lstbox{\textSFx{}}~~\lstbox{\textSFx{}}~~\lstbox{\textSFx{}}~
Error[E4104] : symbol test58::foo has errors
 --> test_resources/diagnostics/test58.yr:(5,14)
 3  ~\lstbox{\textSFxi{}}~ fn foo () throws A;
    ~\lstbox{\textSFv{}}~           ------  error: class type test58::A does not inherit from exception type core::exception::Exception
 4  ~\lstbox{\textSFxi{}}~
 5  ~\lstbox{\textSFxi{}}~ fn main () { foo (); }
    ~\lstbox{\textSFv{}}~              ^^^
    ~\lstbox{\textSFxi{}}~
    = when validating test58::foo -- (test_resources/diagnostics/test58.yr:(3,4))
    = when validating test58::main -- (test_resources/diagnostics/test58.yr:(5,4))
    ~\lstbox{\textSFii{}}~~\lstbox{\textSFx{}}~~\lstbox{\textSFx{}}~~\lstbox{\textSFx{}}~~\lstbox{\textSFx{}}~~\lstbox{\textSFx{}}~~\lstbox{\textSFx{}}~~\lstbox{\textSFx{}}~
\end{lstlisting}

\subsubsection*{How to fix it}

Make the class inherit from the exception type reported in the message.

\errorcode{E4206}{the function \errhole{} might throw an exception of type \errhole{}, but that is not declared in its prototype}
\label{sec:codes:e4206}

Raised during \textbf{validation}, from \texttt{ValidateErrorMessage::\allowbreak{}THROWS\_\allowbreak{}NOT\_\allowbreak{}DECLARED} (\textit{src/\allowbreak{}ymirc/\allowbreak{}semantic/\allowbreak{}validator/\allowbreak{}errors.yr}).

\subsubsection*{What it means}

A function can throw an exception its prototype does not declare. Ymir checks that the throw list of a function covers everything its body can raise.

\subsubsection*{Example}

\textit{test\_\allowbreak{}resources/\allowbreak{}lit\_\allowbreak{}class/\allowbreak{}ctors/\allowbreak{}test7.yr}:

%% check: skip
\begin{lstlisting}[style=coloredverbatimError]
fn foo (t : bool = false)-> i32
    throws AssertError
{
    assert (t);
    12
}

record A {
    let a : i32 = 0;

    pub self ()
        with a = foo ()
    {}

}
\end{lstlisting}

\begin{lstlisting}[style=bashVerb, breaklines=true, escapechar=~]
Error[E4206] : the function test7::A::self might throw an exception of type core::exception::assertion::AssertError, but that is not declared in its prototype
 --> test_resources/lit_class/ctors/test7.yr:(11,9)
11  ~\lstbox{\textSFxi{}}~     pub self ()
    ~\lstbox{\textSFv{}}~         ^^^^
    ~\lstbox{\textSFxi{}}~ Note :
    ~\lstbox{\textSFxi{}}~  --> test_resources/lit_class/ctors/test7.yr:(12,22)
    ~\lstbox{\textSFxi{}}~ 12  ~\lstbox{\textSFxi{}}~         with a = foo ()
    ~\lstbox{\textSFxi{}}~     ~\lstbox{\textSFv{}}~                      ^
    ~\lstbox{\textSFxi{}}~     ~\lstbox{\textSFii{}}~~\lstbox{\textSFx{}}~~\lstbox{\textSFx{}}~~\lstbox{\textSFx{}}~~\lstbox{\textSFx{}}~~\lstbox{\textSFx{}}~~\lstbox{\textSFx{}}~~\lstbox{\textSFx{}}~
    ~\lstbox{\textSFii{}}~~\lstbox{\textSFx{}}~~\lstbox{\textSFx{}}~~\lstbox{\textSFx{}}~~\lstbox{\textSFx{}}~~\lstbox{\textSFx{}}~~\lstbox{\textSFvii{}}~~\lstbox{\textSFx{}}~
\end{lstlisting}

\subsubsection*{How to fix it}

Add the exception to the \tokennolst{throws} clause, or catch it inside the function.

\errorcode{E4207}{the method might throw an exception of type \errhole{}, but that is not covered by the method of the ancestor class}
\label{sec:codes:e4207}

Raised during \textbf{validation}, from \texttt{ValidateErrorMessage::\allowbreak{}THROWS\_\allowbreak{}NOT\_\allowbreak{}DECLARED\_\allowbreak{}OVER} (\textit{src/\allowbreak{}ymirc/\allowbreak{}semantic/\allowbreak{}validator/\allowbreak{}errors.yr}).

\subsubsection*{What it means}

An overriding method throws an exception the ancestor method does not declare, cf. \hyperref[sec:codes:e4186]{E4186}.

\subsubsection*{Example}

\textit{test\_\allowbreak{}resources/\allowbreak{}diagnostics/\allowbreak{}test67.yr}:

%% check: skip
\begin{lstlisting}[style=coloredverbatimError]
class A {
    pub self () {}
    pub fn foo (self) {}
}

class B over A {
    pub self () {}
    pub over foo (self) throws AssertError {
        throw copy AssertError ("x");
    }
}

fn main () { let _ = copy B (); }
\end{lstlisting}

\begin{lstlisting}[style=bashVerb, breaklines=true, escapechar=~]
Error[E4186] : the throwers of the overriding method test67::B::foo ()-> void are not compatible with the throwers of the ancestor method test67::A::foo ()-> void
 --> test_resources/diagnostics/test67.yr:(3,12)
 3  ~\lstbox{\textSFxi{}}~     pub fn foo (self) {}
    ~\lstbox{\textSFv{}}~            ^^^
 --> test_resources/diagnostics/test67.yr:(8,32)
 8  ~\lstbox{\textSFxi{}}~     pub over foo (self) throws AssertError {
    ~\lstbox{\textSFv{}}~                                -----------  error: the method might throw an exception of type core::exception::assertion::AssertError, but that is not covered by the method of the ancestor class
    ~\lstbox{\textSFxi{}}~
    = when validating -- (test_resources/diagnostics/test67.yr:(6,7))
    ~\lstbox{\textSFii{}}~~\lstbox{\textSFx{}}~~\lstbox{\textSFx{}}~~\lstbox{\textSFx{}}~~\lstbox{\textSFx{}}~~\lstbox{\textSFx{}}~~\lstbox{\textSFx{}}~~\lstbox{\textSFx{}}~
Error[E4104] : symbol test67::B has errors
 --> test_resources/diagnostics/test67.yr:(13,27)
13  ~\lstbox{\textSFxi{}}~ fn main () { let _ = copy B (); }
    ~\lstbox{\textSFv{}}~                           ^
 --> test_resources/diagnostics/test67.yr:(3,12)
 3  ~\lstbox{\textSFxi{}}~     pub fn foo (self) {}
    ~\lstbox{\textSFv{}}~            ---  error: the throwers of the overriding method test67::B::foo ()-> void are not compatible with the throwers of the ancestor method test67::A::foo ()-> void
   ...
 8  ~\lstbox{\textSFxi{}}~     pub over foo (self) throws AssertError {
    ~\lstbox{\textSFv{}}~                                -----------  error: the method might throw an exception of type core::exception::assertion::AssertError, but that is not covered by the method of the ancestor class
    ~\lstbox{\textSFxi{}}~
    = when validating -- (test_resources/diagnostics/test67.yr:(6,7))
    = when validating test67::main -- (test_resources/diagnostics/test67.yr:(13,4))
    ~\lstbox{\textSFii{}}~~\lstbox{\textSFx{}}~~\lstbox{\textSFx{}}~~\lstbox{\textSFx{}}~~\lstbox{\textSFx{}}~~\lstbox{\textSFx{}}~~\lstbox{\textSFx{}}~~\lstbox{\textSFx{}}~
\end{lstlisting}

\subsubsection*{How to fix it}

Catch it in the overriding method, or declare it on the ancestor.

\errorcode{E4208}{the prototype of the function \errhole{} declares a possible throw of an exception of type \errhole{}, but the function never throws it}
\label{sec:codes:e4208}

Raised during \textbf{validation}, from \texttt{ValidateErrorMessage::\allowbreak{}THROWS\_\allowbreak{}NOT\_\allowbreak{}USED} (\textit{src/\allowbreak{}ymirc/\allowbreak{}semantic/\allowbreak{}validator/\allowbreak{}errors.yr}).

\subsubsection*{What it means}

A function declares that it throws an exception its body can never raise. The declaration forces every caller to handle something that cannot happen.

\subsubsection*{Example}

\textit{test\_\allowbreak{}resources/\allowbreak{}lit\_\allowbreak{}class/\allowbreak{}ctors/\allowbreak{}test8.yr}:

%% check: skip
\begin{lstlisting}[style=coloredverbatimError]
fn foo (t : bool = false)-> i32 {
    t;
    12
}

record A {
    let a : i32 = 0;

    pub self ()
        with a = foo ()
        throws AssertError
    {}

}
\end{lstlisting}

\begin{lstlisting}[style=bashVerb, breaklines=true, escapechar=~]
Error[E4208] : the prototype of the function test8::A::self declares a possible throw of an exception of type core::exception::assertion::AssertError, but the function never throws it
 --> test_resources/lit_class/ctors/test8.yr:(9,9)
 9  ~\lstbox{\textSFxi{}}~     pub self ()
    ~\lstbox{\textSFv{}}~         ^^^^
    ~\lstbox{\textSFxi{}}~ Note :
    ~\lstbox{\textSFxi{}}~  --> test_resources/lit_class/ctors/test8.yr:(11,16)
    ~\lstbox{\textSFxi{}}~ 11  ~\lstbox{\textSFxi{}}~         throws AssertError
    ~\lstbox{\textSFxi{}}~     ~\lstbox{\textSFv{}}~                ^^^^^^^^^^^
    ~\lstbox{\textSFxi{}}~     ~\lstbox{\textSFii{}}~~\lstbox{\textSFx{}}~~\lstbox{\textSFx{}}~~\lstbox{\textSFx{}}~~\lstbox{\textSFx{}}~~\lstbox{\textSFx{}}~~\lstbox{\textSFx{}}~~\lstbox{\textSFx{}}~
    ~\lstbox{\textSFii{}}~~\lstbox{\textSFx{}}~~\lstbox{\textSFx{}}~~\lstbox{\textSFx{}}~~\lstbox{\textSFx{}}~~\lstbox{\textSFx{}}~~\lstbox{\textSFvii{}}~~\lstbox{\textSFx{}}~
\end{lstlisting}

\subsubsection*{How to fix it}

Remove the exception from the \tokennolst{throws} clause.

\errorcode{E4209}{throw statement is not within a function}
\label{sec:codes:e4209}

Raised during \textbf{validation}, from \texttt{ValidateErrorMessage::\allowbreak{}THROW\_\allowbreak{}NO\_\allowbreak{}FUNCTION} (\textit{src/\allowbreak{}ymirc/\allowbreak{}semantic/\allowbreak{}validator/\allowbreak{}errors.yr}).

\subsubsection*{What it means}

A \tokennolst{throw} was written outside any function.

\subsubsection*{Example}

\textit{test\_\allowbreak{}resources/\allowbreak{}diagnostics/\allowbreak{}test70.yr}:

%% check: skip
\begin{lstlisting}[style=coloredverbatimError]
lazy X : i32 = { throw copy AssertError ("x"); 1 };

fn main () { X; }
\end{lstlisting}

\begin{lstlisting}[style=bashVerb, breaklines=true, escapechar=~]
Error[E4209] : throw statement is not within a function
 --> test_resources/diagnostics/test70.yr:(1,18)
 1  ~\lstbox{\textSFxi{}}~ lazy X : i32 = { throw copy AssertError ("x"); 1 };
    ~\lstbox{\textSFv{}}~                  ^^^^^
    ~\lstbox{\textSFxi{}}~
    = when validating -- (test_resources/diagnostics/test70.yr:(1,6))
    ~\lstbox{\textSFii{}}~~\lstbox{\textSFx{}}~~\lstbox{\textSFx{}}~~\lstbox{\textSFx{}}~~\lstbox{\textSFx{}}~~\lstbox{\textSFx{}}~~\lstbox{\textSFx{}}~~\lstbox{\textSFx{}}~
Error[E4104] : symbol test70::X has errors
 --> test_resources/diagnostics/test70.yr:(3,14)
 1  ~\lstbox{\textSFxi{}}~ lazy X : i32 = { throw copy AssertError ("x"); 1 };
    ~\lstbox{\textSFv{}}~                  -----  error: throw statement is not within a function
 2  ~\lstbox{\textSFxi{}}~
 3  ~\lstbox{\textSFxi{}}~ fn main () { X; }
    ~\lstbox{\textSFv{}}~              ^
    ~\lstbox{\textSFxi{}}~
    = when validating -- (test_resources/diagnostics/test70.yr:(1,6))
    = when validating test70::main -- (test_resources/diagnostics/test70.yr:(3,4))
    ~\lstbox{\textSFii{}}~~\lstbox{\textSFx{}}~~\lstbox{\textSFx{}}~~\lstbox{\textSFx{}}~~\lstbox{\textSFx{}}~~\lstbox{\textSFx{}}~~\lstbox{\textSFx{}}~~\lstbox{\textSFx{}}~
\end{lstlisting}

\subsubsection*{How to fix it}

Move it into the function it belongs to.

\errorcode{E4210}{cannot throw exception inside a scope guard}
\label{sec:codes:e4210}

Raised during \textbf{validation}, from \texttt{ValidateErrorMessage::\allowbreak{}THROW\_\allowbreak{}SCOPE\_\allowbreak{}GUARD} (\textit{src/\allowbreak{}ymirc/\allowbreak{}semantic/\allowbreak{}validator/\allowbreak{}errors.yr}).

\subsubsection*{What it means}

A \tokennolst{throw} inside a scope guard would escape the guard uncaught. A guard runs while a scope is being left, so an exception leaving it would interrupt an unwinding already in progress.

\subsubsection*{Example}

\textit{test\_\allowbreak{}resources/\allowbreak{}scope\_\allowbreak{}guards/\allowbreak{}test10.yr}:

%% check: skip
\begin{lstlisting}[style=coloredverbatimError]
fn main () {
    {

    } exit {
        throw copy AssertError ("");
    }
}
\end{lstlisting}

\begin{lstlisting}[style=bashVerb, breaklines=true, escapechar=~]
Error[E4210] : cannot throw exception inside a scope guard
 --> test_resources/scope_guards/test10.yr:(4,7)
 4  ~\lstbox{\textSFxi{}}~     } exit {
    ~\lstbox{\textSFv{}}~       ^^^^
   ...
 --> test_resources/scope_guards/test10.yr:(5,9)
 5  ~\lstbox{\textSFxi{}}~         throw copy AssertError ("");
    ~\lstbox{\textSFv{}}~         ^^^^^
    ~\lstbox{\textSFii{}}~~\lstbox{\textSFx{}}~~\lstbox{\textSFx{}}~~\lstbox{\textSFx{}}~~\lstbox{\textSFx{}}~~\lstbox{\textSFx{}}~~\lstbox{\textSFx{}}~~\lstbox{\textSFx{}}~
\end{lstlisting}

\subsubsection*{How to fix it}

Catch the exception inside the guard. A \tokennolst{throw} in a nested block whose own \tokennolst{catch} handles every thrown type does not cross the guard and is accepted.

\errorcode{E4211}{rethrowing an exception of type \errhole{} inside a scope guard is forbidden}
\label{sec:codes:e4211}

Raised during \textbf{validation}, from \texttt{ValidateErrorMessage::\allowbreak{}THROW\_\allowbreak{}SCOPE\_\allowbreak{}GUARD\_\allowbreak{}RETHROW} (\textit{src/\allowbreak{}ymirc/\allowbreak{}semantic/\allowbreak{}validator/\allowbreak{}errors.yr}).

\subsubsection*{What it means}

A caught exception is rethrown from inside a scope guard, which would let it escape the guard, cf. \hyperref[sec:codes:e4210]{E4210}.

\subsubsection*{Example}

\textit{test\_\allowbreak{}resources/\allowbreak{}scope\_\allowbreak{}guards/\allowbreak{}test11.yr}:

%% check: skip
\begin{lstlisting}[style=coloredverbatimError]
fn foo ()
    throws AssertError;

fn main () {
    {

    } exit {
        foo ();
    }
}
\end{lstlisting}

\begin{lstlisting}[style=bashVerb, breaklines=true, escapechar=~]
Error[E4211] : rethrowing an exception of type core::exception::assertion::AssertError inside a scope guard is forbidden
 --> test_resources/scope_guards/test11.yr:(7,7)
 7  ~\lstbox{\textSFxi{}}~     } exit {
    ~\lstbox{\textSFv{}}~       ^^^^
   ...
 --> test_resources/scope_guards/test11.yr:(8,13)
 8  ~\lstbox{\textSFxi{}}~         foo ();
    ~\lstbox{\textSFv{}}~             ^
    ~\lstbox{\textSFii{}}~~\lstbox{\textSFx{}}~~\lstbox{\textSFx{}}~~\lstbox{\textSFx{}}~~\lstbox{\textSFx{}}~~\lstbox{\textSFx{}}~~\lstbox{\textSFx{}}~~\lstbox{\textSFx{}}~
\end{lstlisting}

\subsubsection*{How to fix it}

Handle the exception inside the guard rather than rethrowing it.

%% Generated by tools/gen_error_codes.py from docs/errors/ of the compiler;
%% edit the compiler's pages or tools/error_themes.txt, then rerun it.

\clearpage
\section{Templates, macros and compile-time evaluation}
\label{sec:codes:templates}

\begin{longtable}{@{}l p{0.78\linewidth} r@{}}
\toprule Code & Message & Page\\ \midrule \endhead
\hyperref[sec:codes:e3040]{E3040} & macro rule \errhole{} cannot be restricted to a list of values, only a built-in rule can be & \pageref{sec:codes:e3040}\\
\hyperref[sec:codes:e4030]{E4030} & cte assertion failed & \pageref{sec:codes:e4030}\\
\hyperref[sec:codes:e4031]{E4031} & cte assertion failed \errhole{} & \pageref{sec:codes:e4031}\\
\hyperref[sec:codes:e4110]{E4110} & validation of macro call failed & \pageref{sec:codes:e4110}\\
\hyperref[sec:codes:e4111]{E4111} & macro call with "\errhole{}" does not match the rule "\errhole{}" & \pageref{sec:codes:e4111}\\
\hyperref[sec:codes:e4112]{E4112} & macro validation is not complete, some tokens remain unassociated  \errhole{} \errhole{} \errhole{} & \pageref{sec:codes:e4112}\\
\hyperref[sec:codes:e4119]{E4119} & malformed \_\_pragma \errhole{} & \pageref{sec:codes:e4119}\\
\hyperref[sec:codes:e4124]{E4124} & reached the maximum number of cte iterations \errhole{} \textgreater{} \errhole{} & \pageref{sec:codes:e4124}\\
\hyperref[sec:codes:e4202]{E4202} & template specialization for \errhole{} fails with \errhole{} & \pageref{sec:codes:e4202}\\
\hyperref[sec:codes:e4203]{E4203} & template specialization fails with \errhole{} & \pageref{sec:codes:e4203}\\
\hyperref[sec:codes:e4234]{E4234} & macro symbol \errhole{} has no rule named \errhole{} & \pageref{sec:codes:e4234}\\
\hyperref[sec:codes:e4236]{E4236} & undefined template call for \errhole{} with \errhole{} & \pageref{sec:codes:e4236}\\
\hyperref[sec:codes:e4256]{E4256} & value of type \errhole{} is needed but unknown at compilation time & \pageref{sec:codes:e4256}\\
\hyperref[sec:codes:e4257]{E4257} & unknown length of expansion for type \errhole{} & \pageref{sec:codes:e4257}\\
\hyperref[sec:codes:e4258]{E4258} & unknown \_\_pragma expression \errhole{} & \pageref{sec:codes:e4258}\\
\hyperref[sec:codes:e4261]{E4261} & unresolved template parameter \errhole{} & \pageref{sec:codes:e4261}\\
\hyperref[sec:codes:e4271]{E4271} & template specialization expected a type not the value \errhole{} & \pageref{sec:codes:e4271}\\
\hyperref[sec:codes:e4305]{E4305} & \_\_pragma \errhole{} is not enclosed in a \errhole{} & \pageref{sec:codes:e4305}\\
\hyperref[sec:codes:e4316]{E4316} & \errhole{} is not an enum type & \pageref{sec:codes:e4316}\\
\bottomrule
\end{longtable}

\errorcode{E3040}{macro rule \errhole{} cannot be restricted to a list of values, only a built-in rule can be}
\label{sec:codes:e3040}

Raised during \textbf{declaration}, from \texttt{DeclareErrorMessage::\allowbreak{}MACRO\_\allowbreak{}NOT\_\allowbreak{}RESTRICTABLE} (\textit{src/\allowbreak{}ymirc/\allowbreak{}semantic/\allowbreak{}declarator/\allowbreak{}errors.yr}).

\subsubsection*{What it means}

A rule can be restricted to a list of values, written \tokennolst{\_\_key["null" | "in"]}, and the list is compared against the content the rule consumed. A rule the macro declares answers its rewritten body rather than that content, so a list cannot restrict it.

The rule head says so on its own, without a call, so the check runs when the macro is declared.

\subsubsection*{Example}

\textit{test\_\allowbreak{}resources/\allowbreak{}lit\_\allowbreak{}macro/\allowbreak{}test43.yr}:

%% check: skip
\begin{lstlisting}[style=coloredverbatimError]
/*
 * A list of values restricting a rule of the macro, which answers its
 * rewritten body and not the content it consumed
 */

macro named {
    self (v=entry["a"]) skips (" ") #{
        #(v)
    }

    fn entry (i=__ident) skips (" ") #{
        12
    }
}

fn foo ()-> i32 {
    named#{ a }
}
\end{lstlisting}

\begin{lstlisting}[style=bashVerb, breaklines=true, escapechar=~]
Error[E3040] : macro rule entry cannot be restricted to a list of values, only a built-in rule can be
 --> test_resources/lit_macro/test43.yr:(7,13)
 7  ~\lstbox{\textSFxi{}}~     self (v=entry["a"]) skips (" ") #{
    ~\lstbox{\textSFv{}}~             ^^^^^
 --> test_resources/lit_macro/test43.yr:(11,8)
11  ~\lstbox{\textSFxi{}}~     fn entry (i=__ident) skips (" ") #{
    ~\lstbox{\textSFv{}}~        -----  note: declared here
    ~\lstbox{\textSFii{}}~~\lstbox{\textSFx{}}~~\lstbox{\textSFx{}}~~\lstbox{\textSFx{}}~~\lstbox{\textSFx{}}~~\lstbox{\textSFx{}}~~\lstbox{\textSFx{}}~~\lstbox{\textSFx{}}~
\end{lstlisting}

\subsubsection*{How to fix it}

Restrict the built-in rule inside the rule that names it -- \tokennolst{fn entry (i=\_\_ident["a"])} -- or split the accepted values into one rule each and choose between them with \tokennolst{|}, which is ordered and selects the first that matches.

\errorcode{E4030}{cte assertion failed}
\label{sec:codes:e4030}

Raised during \textbf{validation}, from \texttt{ValidateErrorMessage::\allowbreak{}CTE\_\allowbreak{}ASSERT\_\allowbreak{}NO\_\allowbreak{}MSG} (\textit{src/\allowbreak{}ymirc/\allowbreak{}semantic/\allowbreak{}validator/\allowbreak{}errors.yr}).

\subsubsection*{What it means}

A \tokennolst{cte assert} evaluated to false during compile time execution, and no message was given.

\subsubsection*{Example}

\textit{test\_\allowbreak{}resources/\allowbreak{}diagnostics/\allowbreak{}test25.yr}:

%% check: skip
\begin{lstlisting}[style=coloredverbatimError]
fn main () {
    cte assert (1 == 2);
}
\end{lstlisting}

\begin{lstlisting}[style=bashVerb, breaklines=true, escapechar=~]
Error[E4030] : cte assertion failed
 --> test_resources/diagnostics/test25.yr:(2,9)
 2  ~\lstbox{\textSFxi{}}~     cte assert (1 == 2);
    ~\lstbox{\textSFv{}}~         ^^^^^^
    ~\lstbox{\textSFii{}}~~\lstbox{\textSFx{}}~~\lstbox{\textSFx{}}~~\lstbox{\textSFx{}}~~\lstbox{\textSFx{}}~~\lstbox{\textSFx{}}~~\lstbox{\textSFx{}}~~\lstbox{\textSFx{}}~
\end{lstlisting}

\subsubsection*{How to fix it}

Fix the condition, or the values it tests. Adding a message to the assert makes the next failure say what it was checking.

\errorcode{E4031}{cte assertion failed \errhole{}}
\label{sec:codes:e4031}

Raised during \textbf{validation}, from \texttt{ValidateErrorMessage::\allowbreak{}CTE\_\allowbreak{}ASSERT\_\allowbreak{}WITH\_\allowbreak{}MSG} (\textit{src/\allowbreak{}ymirc/\allowbreak{}semantic/\allowbreak{}validator/\allowbreak{}errors.yr}).

\subsubsection*{What it means}

A \tokennolst{cte assert} evaluated to false during compile time execution; its message is reported.

\subsubsection*{Example}

\textit{test\_\allowbreak{}resources/\allowbreak{}regressions/\allowbreak{}assert\_\allowbreak{}bool\_\allowbreak{}type/\allowbreak{}test3.yr}:

%% check: skip
\begin{lstlisting}[style=coloredverbatimError]
fn main () {
    cte assert (false, "Must fail");
}
\end{lstlisting}

\begin{lstlisting}[style=bashVerb, breaklines=true, escapechar=~]
Error[E4031] : cte assertion failed "Must fail"s8
 --> test_resources/regressions/assert_bool_type/test3.yr:(2,9)
 2  ~\lstbox{\textSFxi{}}~     cte assert (false, "Must fail");
    ~\lstbox{\textSFv{}}~         ^^^^^^
    ~\lstbox{\textSFii{}}~~\lstbox{\textSFx{}}~~\lstbox{\textSFx{}}~~\lstbox{\textSFx{}}~~\lstbox{\textSFx{}}~~\lstbox{\textSFx{}}~~\lstbox{\textSFx{}}~~\lstbox{\textSFx{}}~
\end{lstlisting}

\subsubsection*{How to fix it}

Fix the condition, or the values it tests.

\errorcode{E4110}{validation of macro call failed}
\label{sec:codes:e4110}

Raised during \textbf{validation}, from \texttt{ValidateErrorMessage::\allowbreak{}MACRO\_\allowbreak{}CALL\_\allowbreak{}VALIDATION\_\allowbreak{}FAILED} (\textit{src/\allowbreak{}ymirc/\allowbreak{}semantic/\allowbreak{}validator/\allowbreak{}errors.yr}).

\subsubsection*{What it means}

A macro call was expanded and validating the result failed.

\subsubsection*{Example}

\textit{test\_\allowbreak{}resources/\allowbreak{}lit\_\allowbreak{}macro/\allowbreak{}test25.yr}:

%% check: skip
\begin{lstlisting}[style=coloredverbatimError]
/*
 * A macro call made on a symbol that is not a macro
 */

fn notAMacro ()-> i32 {
    12
}

fn foo () {
    let x = notAMacro#{ 1 };
    x;
}
\end{lstlisting}

\begin{lstlisting}[style=bashVerb, breaklines=true, escapechar=~]
Error[E4110] : validation of macro call failed
 --> test_resources/lit_macro/test25.yr:(10,22)
10  ~\lstbox{\textSFxi{}}~     let x = notAMacro#{ 1 };
    ~\lstbox{\textSFv{}}~                      ^^  error: macro call is undefined for test25::notAMacro ()-> i32
    ~\lstbox{\textSFii{}}~~\lstbox{\textSFx{}}~~\lstbox{\textSFx{}}~~\lstbox{\textSFx{}}~~\lstbox{\textSFx{}}~~\lstbox{\textSFx{}}~~\lstbox{\textSFx{}}~~\lstbox{\textSFx{}}~
\end{lstlisting}

\subsubsection*{Notes}

A macro call that cannot be expanded reports why under this diagnostic, as a note carrying no code of its own:

\begin{itemize}
\item \texttt{macro \textless{}\textgreater{} has no accessible constructors in that context} (formerly E4228): the macro declares no constructor visible from the call.
\item \tokennolst{macro call is undefined for \textless{}\textgreater{}} (formerly E4232): the symbol called with \tokennolst{\#\{} is not a macro, as in the example above.
\item \tokennolst{undefined symbol \textless{}\textgreater{}} (formerly E4245): the rewrite of a macro rule names a variable its pattern does not capture.
\end{itemize}

\subsubsection*{How to fix it}

Read the errors reported under this one: they come from the expansion, and point either at the macro definition or at the arguments it was given.

\errorcode{E4111}{macro call with "\errhole{}" does not match the rule "\errhole{}"}
\label{sec:codes:e4111}

Raised during \textbf{validation}, from \texttt{ValidateErrorMessage::\allowbreak{}MACRO\_\allowbreak{}DOES\_\allowbreak{}NOT\_\allowbreak{}MATCH} (\textit{src/\allowbreak{}ymirc/\allowbreak{}semantic/\allowbreak{}validator/\allowbreak{}errors.yr}).

\subsubsection*{What it means}

A macro call does not match the rule it was tried against. Both the call and the rule are quoted in the message.

\subsubsection*{Example}

\textit{test\_\allowbreak{}resources/\allowbreak{}lit\_\allowbreak{}macro/\allowbreak{}test21.yr}:

%% check: skip
\begin{lstlisting}[style=coloredverbatimError]
/*
 * A macro call matching none of the branches of a '|' operator
 */

macro alt {
    self (v=(a=__literal | b=__ident) ";") skips (" ") #{
        #(v)
    }
}

fn foo ()-> i32 {
    alt#{ + ; }
}
\end{lstlisting}

\begin{lstlisting}[style=bashVerb, breaklines=true, escapechar=~]
Error[E4111] : macro call with "+ ; " does not match the rule "a=__literal | b=__ident"
 --> test_resources/lit_macro/test21.yr:(6,26)
 6  ~\lstbox{\textSFxi{}}~     self (v=(a=__literal | b=__ident) ";") skips (" ") #{
    ~\lstbox{\textSFv{}}~                          ^
   ...
 --> test_resources/lit_macro/test21.yr:(12,11)
12  ~\lstbox{\textSFxi{}}~     alt#{ + ; }
    ~\lstbox{\textSFv{}}~           ^^^^
    ~\lstbox{\textSFxi{}}~ Error[E2007] : unexpected +
    ~\lstbox{\textSFxi{}}~  --> test_resources/lit_macro/test21.yr:(1,1)
    ~\lstbox{\textSFxi{}}~  1  ~\lstbox{\textSFxi{}}~ + ;
    ~\lstbox{\textSFxi{}}~     ~\lstbox{\textSFv{}}~ ^
    ~\lstbox{\textSFxi{}}~     ~\lstbox{\textSFii{}}~~\lstbox{\textSFx{}}~~\lstbox{\textSFx{}}~~\lstbox{\textSFx{}}~~\lstbox{\textSFx{}}~~\lstbox{\textSFx{}}~~\lstbox{\textSFx{}}~~\lstbox{\textSFx{}}~
    ~\lstbox{\textSFxi{}}~ Error[E2007] : unexpected +
    ~\lstbox{\textSFxi{}}~  --> test_resources/lit_macro/test21.yr:(1,1)
    ~\lstbox{\textSFxi{}}~  1  ~\lstbox{\textSFxi{}}~ + ;
    ~\lstbox{\textSFxi{}}~     ~\lstbox{\textSFv{}}~ ^
    ~\lstbox{\textSFxi{}}~     ~\lstbox{\textSFii{}}~~\lstbox{\textSFx{}}~~\lstbox{\textSFx{}}~~\lstbox{\textSFx{}}~~\lstbox{\textSFx{}}~~\lstbox{\textSFx{}}~~\lstbox{\textSFx{}}~~\lstbox{\textSFx{}}~
    ~\lstbox{\textSFii{}}~~\lstbox{\textSFx{}}~~\lstbox{\textSFx{}}~~\lstbox{\textSFx{}}~~\lstbox{\textSFx{}}~~\lstbox{\textSFx{}}~~\lstbox{\textSFvii{}}~~\lstbox{\textSFx{}}~
\end{lstlisting}

\subsubsection*{How to fix it}

Compare the call with the rule. When a macro has several rules, each one that was tried is reported, so the diagnostic shows why none matched.

\errorcode{E4112}{macro validation is not complete, some tokens remain unassociated  \errhole{} \errhole{} \errhole{}}
\label{sec:codes:e4112}

Raised during \textbf{validation}, from \texttt{ValidateErrorMessage::\allowbreak{}MACRO\_\allowbreak{}REST} (\textit{src/\allowbreak{}ymirc/\allowbreak{}semantic/\allowbreak{}validator/\allowbreak{}errors.yr}).

\subsubsection*{What it means}

A macro rule matched a prefix of the call and left tokens over. A rule has to consume everything it is given.

\subsubsection*{Example}

\textit{test\_\allowbreak{}resources/\allowbreak{}lit\_\allowbreak{}macro/\allowbreak{}test7.yr}:

%% check: skip
\begin{lstlisting}[style=coloredverbatimError]
/*
 * A macro call matching a constructor, but whose content is not
 * entirely consumed by it
 */

macro one {
    self (a=__operand) skips (" ") #{
        #(a)
    }
}

fn foo () {
    let x = one#{ 1 2 };

    x;
}
\end{lstlisting}

\begin{lstlisting}[style=bashVerb, breaklines=true, escapechar=~]
Error[E4112] : macro validation is not complete, some tokens remain unassociated
==
2
==
 --> test_resources/lit_macro/test7.yr:(13,21)
13  ~\lstbox{\textSFxi{}}~     let x = one#{ 1 2 };
    ~\lstbox{\textSFv{}}~                     ^^
    ~\lstbox{\textSFii{}}~~\lstbox{\textSFx{}}~~\lstbox{\textSFx{}}~~\lstbox{\textSFx{}}~~\lstbox{\textSFx{}}~~\lstbox{\textSFx{}}~~\lstbox{\textSFx{}}~~\lstbox{\textSFx{}}~
\end{lstlisting}

\subsubsection*{How to fix it}

Extend the rule to cover the remaining tokens, or shorten the call.

\errorcode{E4119}{malformed \_\_pragma \errhole{}}
\label{sec:codes:e4119}

Raised during \textbf{validation}, from \texttt{ValidateErrorMessage::\allowbreak{}MALFORMED\_\allowbreak{}PRAGMA} (\textit{src/\allowbreak{}ymirc/\allowbreak{}semantic/\allowbreak{}validator/\allowbreak{}errors.yr}).

\subsubsection*{What it means}

A \tokennolst{\_\_pragma} expression is malformed -- the arguments it was given are not the ones the pragma takes.

\subsubsection*{Example}

\textit{test\_\allowbreak{}resources/\allowbreak{}pragma/\allowbreak{}test13.yr}:

%% check: skip
\begin{lstlisting}[style=coloredverbatimError]
/*
 * Malformed arities of the type producing pragmas.
 */

class A {
    pub self () {}
}

fn fieldTypeTooFew ()-> __pragma!field_type (&A) {
    12
}

fn unproxiedTooMany ()-> __pragma!unproxied (i32, i32) {
    12
}

fn superTooMany ()-> __pragma!super (&A, &A) {
    12
}

fn rootSuperTooMany ()-> __pragma!root_super (&A, &A) {
    12
}
\end{lstlisting}

\begin{lstlisting}[style=bashVerb, breaklines=true, escapechar=~]
Error[E4119] : malformed __pragma field_type
 --> test_resources/pragma/test13.yr:(9,25)
 9  ~\lstbox{\textSFxi{}}~ fn fieldTypeTooFew ()-> __pragma!field_type (&A) {
    ~\lstbox{\textSFv{}}~                         ^^^^^^^^
    ~\lstbox{\textSFii{}}~~\lstbox{\textSFx{}}~~\lstbox{\textSFx{}}~~\lstbox{\textSFx{}}~~\lstbox{\textSFx{}}~~\lstbox{\textSFx{}}~~\lstbox{\textSFx{}}~~\lstbox{\textSFx{}}~
\end{lstlisting}

\subsubsection*{How to fix it}

Check the pragma against its documented form; the name of the pragma is in the message.

\errorcode{E4124}{reached the maximum number of cte iterations \errhole{} \textgreater{} \errhole{}}
\label{sec:codes:e4124}

Raised during \textbf{validation}, from \texttt{ValidateErrorMessage::\allowbreak{}MAX\_\allowbreak{}LOOP\_\allowbreak{}ITERATIONS} (\textit{src/\allowbreak{}ymirc/\allowbreak{}semantic/\allowbreak{}validator/\allowbreak{}errors.yr}).

\subsubsection*{What it means}

Compile time execution ran past the iteration limit. Both the count reached and the limit are in the message. This is the guard against a compile time loop that never terminates.

\subsubsection*{Example}

\textit{test\_\allowbreak{}resources/\allowbreak{}diagnostics/\allowbreak{}test113.yr}:

%% check: skip
\begin{lstlisting}[style=coloredverbatimError]
// a compile time loop running more iterations than allowed
fn main () {
    cte for i in 0 .. 100000 {
        let _ = i;
    }
}
\end{lstlisting}

\begin{lstlisting}[style=bashVerb, breaklines=true, escapechar=~]
Error[E4124] : reached the maximum number of cte iterations 100000 > 1024
 --> test_resources/diagnostics/test113.yr:(3,9)
 3  ~\lstbox{\textSFxi{}}~     cte for i in 0 .. 100000 {
    ~\lstbox{\textSFv{}}~         ^^^
    ~\lstbox{\textSFxi{}}~
    = when validating test113::main -- (test_resources/diagnostics/test113.yr:(2,4))
    ~\lstbox{\textSFii{}}~~\lstbox{\textSFx{}}~~\lstbox{\textSFx{}}~~\lstbox{\textSFx{}}~~\lstbox{\textSFx{}}~~\lstbox{\textSFx{}}~~\lstbox{\textSFx{}}~~\lstbox{\textSFx{}}~
\end{lstlisting}

\subsubsection*{How to fix it}

Check the termination condition of the loop being evaluated at compile time. If the computation is genuinely that long, it belongs at run time.

\errorcode{E4202}{template specialization for \errhole{} fails with \errhole{}}
\label{sec:codes:e4202}

Raised during \textbf{validation}, from \texttt{ValidateErrorMessage::\allowbreak{}TEMPLATE\_\allowbreak{}SPECIALIZATION\_\allowbreak{}FAILS} (\textit{src/\allowbreak{}ymirc/\allowbreak{}semantic/\allowbreak{}validator/\allowbreak{}errors.yr}).

\subsubsection*{What it means}

A template could not be specialized for the given arguments; the reason is in the message.

\subsubsection*{Example}

\textit{test\_\allowbreak{}resources/\allowbreak{}templates/\allowbreak{}class/\allowbreak{}test4.yr}:

%% check: skip
\begin{lstlisting}[style=coloredverbatimError]
class A {

    pub self {T} (a : T)
        with self!{i32} (a)
    {
    }

}


fn main () {
    copy A (1);
}
\end{lstlisting}

\begin{lstlisting}[style=bashVerb, breaklines=true, escapechar=~]
Error[E4202] : template specialization for test4::A {T} fails with {T=> i32}
 --> test_resources/templates/class/test4.yr:(12,12)
12  ~\lstbox{\textSFxi{}}~     copy A (1);
    ~\lstbox{\textSFv{}}~            ^
    ~\lstbox{\textSFxi{}}~ Error : when validating test4::A::self!{i32}::self
    ~\lstbox{\textSFxi{}}~  --> test_resources/templates/class/test4.yr:(3,9)
    ~\lstbox{\textSFxi{}}~  3  ~\lstbox{\textSFxi{}}~     pub self {T} (a : T)
    ~\lstbox{\textSFxi{}}~     ~\lstbox{\textSFv{}}~         ^^^^
    ~\lstbox{\textSFxi{}}~     ~\lstbox{\textSFxi{}}~ Error[E4035] : infinite constructor redirection calls when calling test4::A::self!{i32}::self (a: i32)-> mut &(mut test4::A)
    ~\lstbox{\textSFxi{}}~     ~\lstbox{\textSFxi{}}~  --> test_resources/templates/class/test4.yr:(4,14)
    ~\lstbox{\textSFxi{}}~     ~\lstbox{\textSFxi{}}~  4  ~\lstbox{\textSFxi{}}~         with self!{i32} (a)
    ~\lstbox{\textSFxi{}}~     ~\lstbox{\textSFxi{}}~     ~\lstbox{\textSFv{}}~              ^^^^
    ~\lstbox{\textSFxi{}}~     ~\lstbox{\textSFxi{}}~     ~\lstbox{\textSFii{}}~~\lstbox{\textSFx{}}~~\lstbox{\textSFx{}}~~\lstbox{\textSFx{}}~~\lstbox{\textSFx{}}~~\lstbox{\textSFx{}}~~\lstbox{\textSFx{}}~~\lstbox{\textSFx{}}~
    ~\lstbox{\textSFxi{}}~     ~\lstbox{\textSFii{}}~~\lstbox{\textSFx{}}~~\lstbox{\textSFx{}}~~\lstbox{\textSFx{}}~~\lstbox{\textSFx{}}~~\lstbox{\textSFx{}}~~\lstbox{\textSFvii{}}~~\lstbox{\textSFx{}}~
    ~\lstbox{\textSFii{}}~~\lstbox{\textSFx{}}~~\lstbox{\textSFx{}}~~\lstbox{\textSFx{}}~~\lstbox{\textSFx{}}~~\lstbox{\textSFx{}}~~\lstbox{\textSFvii{}}~~\lstbox{\textSFx{}}~
\end{lstlisting}

\subsubsection*{How to fix it}

Read the reason: it is usually a constraint the arguments do not satisfy, or a parameter that could not be deduced.

\errorcode{E4203}{template specialization fails with \errhole{}}
\label{sec:codes:e4203}

Raised during \textbf{validation}, from \texttt{ValidateErrorMessage::\allowbreak{}TEMPLATE\_\allowbreak{}SPECIALIZATION\_\allowbreak{}FAILS\_\allowbreak{}SIMPLE} (\textit{src/\allowbreak{}ymirc/\allowbreak{}semantic/\allowbreak{}validator/\allowbreak{}errors.yr}).

\subsubsection*{What it means}

A template specialization failed, with the reason in the message and no specific symbol to attribute it to.

\subsubsection*{Example}

\textit{test\_\allowbreak{}resources/\allowbreak{}diagnostics/\allowbreak{}test79.yr}:

%% check: skip
\begin{lstlisting}[style=coloredverbatimError]
fn main () {
    let f = copy (|| => { 1 });
    f ();
}
\end{lstlisting}

\begin{lstlisting}[style=bashVerb, breaklines=true, escapechar=~]
Error[E4203] : template specialization fails with {}
 --> test_resources/diagnostics/test79.yr:(2,19)
 2  ~\lstbox{\textSFxi{}}~     let f = copy (|| => { 1 });
    ~\lstbox{\textSFv{}}~                   ^^
    ~\lstbox{\textSFxi{}}~ Error[E4250] : copy the lambda closure has no effect, no values are enclosed
    ~\lstbox{\textSFxi{}}~  --> test_resources/diagnostics/test79.yr:(2,13)
    ~\lstbox{\textSFxi{}}~  2  ~\lstbox{\textSFxi{}}~     let f = copy (|| => { 1 });
    ~\lstbox{\textSFxi{}}~     ~\lstbox{\textSFv{}}~             ^^^^
    ~\lstbox{\textSFxi{}}~     ~\lstbox{\textSFii{}}~~\lstbox{\textSFx{}}~~\lstbox{\textSFx{}}~~\lstbox{\textSFx{}}~~\lstbox{\textSFx{}}~~\lstbox{\textSFx{}}~~\lstbox{\textSFx{}}~~\lstbox{\textSFx{}}~
    ~\lstbox{\textSFii{}}~~\lstbox{\textSFx{}}~~\lstbox{\textSFx{}}~~\lstbox{\textSFx{}}~~\lstbox{\textSFx{}}~~\lstbox{\textSFx{}}~~\lstbox{\textSFvii{}}~~\lstbox{\textSFx{}}~
\end{lstlisting}

\subsubsection*{Notes}

A lambda or a closure whose body cannot be validated reports why under this diagnostic, as a note carrying no code of its own:

\begin{itemize}
\item \texttt{the type \textless{}\textgreater{} cannot be enclosed because it is not complete} (formerly E4052): a closure encloses a variable holding a lambda, whose type is only known once called.
\item \texttt{movable type \textless{}\textgreater{} cannot be enclosed within a closure function} (formerly E4053): a closure encloses an entity.
\item \texttt{\textless{}\textgreater{} is enclosed, and cannot be used as a lvalue} (formerly E4152): a closure assigns a variable it encloses by copy.
\item \texttt{a lambda function must be safe, but there are exceptions that are not caught} (formerly E4205): the body of a lambda can throw.
\end{itemize}

\subsubsection*{How to fix it}

Same as \hyperref[sec:codes:e4202]{E4202}.

\errorcode{E4234}{macro symbol \errhole{} has no rule named \errhole{}}
\label{sec:codes:e4234}

Raised during \textbf{validation}, from \texttt{ValidateErrorMessage::\allowbreak{}UNDEFINED\_\allowbreak{}RULE\_\allowbreak{}MACRO} (\textit{src/\allowbreak{}ymirc/\allowbreak{}semantic/\allowbreak{}validator/\allowbreak{}errors.yr}).

\subsubsection*{What it means}

A macro rule was named and the macro has no rule of that name.

\subsubsection*{Example}

\textit{test\_\allowbreak{}resources/\allowbreak{}lit\_\allowbreak{}macro/\allowbreak{}test26.yr}:

%% check: skip
\begin{lstlisting}[style=coloredverbatimError]
/*
 * A rule of a macro referring to a rule name that is neither a rule of
 * the macro, nor a built-in rule
 */

macro unknown {
    self (a=notARule) skips (" ") #{
        #(a)
    }
}

fn foo ()-> i32 {
    unknown#{ 1 }
}
\end{lstlisting}

\begin{lstlisting}[style=bashVerb, breaklines=true, escapechar=~]
Error[E4234] : macro symbol test26::unknown has no rule named notARule
 --> test_resources/lit_macro/test26.yr:(7,13)
 7  ~\lstbox{\textSFxi{}}~     self (a=notARule) skips (" ") #{
    ~\lstbox{\textSFv{}}~             ^^^^^^^^
   ...
 --> test_resources/lit_macro/test26.yr:(13,15)
13  ~\lstbox{\textSFxi{}}~     unknown#{ 1 }
    ~\lstbox{\textSFv{}}~               ^^
    ~\lstbox{\textSFii{}}~~\lstbox{\textSFx{}}~~\lstbox{\textSFx{}}~~\lstbox{\textSFx{}}~~\lstbox{\textSFx{}}~~\lstbox{\textSFx{}}~~\lstbox{\textSFx{}}~~\lstbox{\textSFx{}}~
\end{lstlisting}

\subsubsection*{How to fix it}

Check the spelling against the rules the macro declares.

\errorcode{E4236}{undefined template call for \errhole{} with \errhole{}}
\label{sec:codes:e4236}

Raised during \textbf{validation}, from \texttt{ValidateErrorMessage::\allowbreak{}UNDEFINED\_\allowbreak{}TEMPLATE\_\allowbreak{}CALL} (\textit{src/\allowbreak{}ymirc/\allowbreak{}semantic/\allowbreak{}validator/\allowbreak{}errors.yr}).

\subsubsection*{What it means}

A template was called and no specialization of it accepts the arguments. Both the template and the arguments are in the message, and the specializations tried are listed under it.

\subsubsection*{Example}

\textit{test\_\allowbreak{}resources/\allowbreak{}templates/\allowbreak{}explicit/\allowbreak{}test4.yr}:

%% check: skip
\begin{lstlisting}[style=coloredverbatimError]
fn foo {T : fn (i32)-> i32} (){}

fn main () {
    foo!{|x| => x + 1.0f} ();
    foo!{|x : f32| => x} ();
}
\end{lstlisting}

\begin{lstlisting}[style=bashVerb, breaklines=true, escapechar=~]
Error[E4236] : undefined template call for test4::foo {T: fn (_: i32)-> i32}() with {|x|-> _}
 --> test_resources/templates/explicit/test4.yr:(4,8)
 4  ~\lstbox{\textSFxi{}}~     foo!{|x| => x + 1.0f} ();
    ~\lstbox{\textSFv{}}~        ^
    ~\lstbox{\textSFxi{}}~ Error[E4203] : template specialization fails with {i32}
    ~\lstbox{\textSFxi{}}~  --> test_resources/templates/explicit/test4.yr:(4,10)
    ~\lstbox{\textSFxi{}}~  4  ~\lstbox{\textSFxi{}}~     foo!{|x| => x + 1.0f} ();
    ~\lstbox{\textSFxi{}}~     ~\lstbox{\textSFv{}}~          ^
    ~\lstbox{\textSFxi{}}~     ~\lstbox{\textSFxi{}}~ Error[E4224] : undefined operator + for types i32 and f32
    ~\lstbox{\textSFxi{}}~     ~\lstbox{\textSFxi{}}~  --> test_resources/templates/explicit/test4.yr:(4,19)
    ~\lstbox{\textSFxi{}}~     ~\lstbox{\textSFxi{}}~  4  ~\lstbox{\textSFxi{}}~     foo!{|x| => x + 1.0f} ();
    ~\lstbox{\textSFxi{}}~     ~\lstbox{\textSFxi{}}~     ~\lstbox{\textSFv{}}~                   ^
    ~\lstbox{\textSFxi{}}~     ~\lstbox{\textSFxi{}}~     ~\lstbox{\textSFii{}}~~\lstbox{\textSFx{}}~~\lstbox{\textSFx{}}~~\lstbox{\textSFx{}}~~\lstbox{\textSFx{}}~~\lstbox{\textSFx{}}~~\lstbox{\textSFx{}}~~\lstbox{\textSFx{}}~
    ~\lstbox{\textSFxi{}}~     ~\lstbox{\textSFii{}}~~\lstbox{\textSFx{}}~~\lstbox{\textSFx{}}~~\lstbox{\textSFx{}}~~\lstbox{\textSFx{}}~~\lstbox{\textSFx{}}~~\lstbox{\textSFvii{}}~~\lstbox{\textSFx{}}~
    ~\lstbox{\textSFii{}}~~\lstbox{\textSFx{}}~~\lstbox{\textSFx{}}~~\lstbox{\textSFx{}}~~\lstbox{\textSFx{}}~~\lstbox{\textSFx{}}~~\lstbox{\textSFvii{}}~~\lstbox{\textSFx{}}~
\end{lstlisting}

\subsubsection*{Notes}

The reason a specialization was refused is reported under this diagnostic, as a note carrying no code of its own:

\begin{itemize}
\item \texttt{closure functions cannot be used as template values} (formerly E4022): a value was given for a template parameter of delegate type.
\item \tokennolst{\textless{}\textgreater{} is not an array type} (formerly E4148): a \tokennolst{T of [U ; N]} parameter was given a type that is not a static array.
\item \tokennolst{\textless{}\textgreater{} is not a map type} (formerly E4158): a \tokennolst{T of [K =\textgreater{} V]} parameter was given a type that is not a map.
\item \tokennolst{\textless{}\textgreater{} is not a record type} (formerly E4159): a \tokennolst{record T} parameter was given a type that is not a record.
\item \tokennolst{\textless{}\textgreater{} is not a slice type} (formerly E4160): a \tokennolst{T of [U]} parameter was given a type that is not a slice.
\item \tokennolst{\textless{}\textgreater{} is not an entity type} (formerly E4161): an \tokennolst{entity T} parameter was given a type that is not an entity.
\item \texttt{template validation is not complete, unresolved: \{\textless{}\textgreater{}\}} (formerly E4201): more template arguments were given than the template has parameters.
\item \texttt{template specialization expected a value not the type \textless{}\textgreater{}} (formerly E4274): a type was given for a value parameter such as \tokennolst{N : i32}.
\item \texttt{the test of the template specialization failed} (formerly E4204): the \tokennolst{if} guard of the template evaluated to false.
\item \texttt{attribute \textless{}\textgreater{} is not usable in this context} (formerly E4221): the template carries an attribute its instantiation does not accept, such as \tokennolst{@abstract} on a template function.
\end{itemize}

When the template arguments are inferred from a call rather than written, the same notes are listed under the candidates of \hyperref[sec:codes:e4226]{E4226}.

\subsubsection*{How to fix it}

Read the listed failures: each says why one specialization did not apply.

\errorcode{E4256}{value of type \errhole{} is needed but unknown at compilation time}
\label{sec:codes:e4256}

Raised during \textbf{validation}, from \texttt{ValidateErrorMessage::\allowbreak{}UNKNOWN\_\allowbreak{}AT\_\allowbreak{}COMPILE\_\allowbreak{}TIME} (\textit{src/\allowbreak{}ymirc/\allowbreak{}semantic/\allowbreak{}validator/\allowbreak{}errors.yr}).

\subsubsection*{What it means}

A value is needed at compile time and is only known at run time.

\subsubsection*{Example}

\textit{test\_\allowbreak{}resources/\allowbreak{}regressions/\allowbreak{}cte\_\allowbreak{}if\_\allowbreak{}else\_\allowbreak{}if\_\allowbreak{}chain/\allowbreak{}test1.yr}:

%% check: skip
\begin{lstlisting}[style=coloredverbatimError]
in test1;

/**
 * `else if` continues the `cte if` chain, so its condition is evaluated at
 * compile time as well. A runtime comparison there is an error.
 */

fn check {R...} (x : [c8], lst : R)-> bool {
    cte if (R,)::arity > 1 {
        for z in lst if x == z {
            return true;
        }
    }

    else if lst == x {
        return true;
    }

    return false;
}

fn main () {
    let _ = check ("a", "b");
}
\end{lstlisting}

\begin{lstlisting}[style=bashVerb, breaklines=true, escapechar=~]
Error[E4256] : value of type bool is needed but unknown at compilation time
 --> test_resources/regressions/cte_if_else_if_chain/test1.yr:(15,17)
15  ~\lstbox{\textSFxi{}}~     else if lst == x {
    ~\lstbox{\textSFv{}}~                 ^^
    ~\lstbox{\textSFii{}}~~\lstbox{\textSFx{}}~~\lstbox{\textSFx{}}~~\lstbox{\textSFx{}}~~\lstbox{\textSFx{}}~~\lstbox{\textSFx{}}~~\lstbox{\textSFx{}}~~\lstbox{\textSFx{}}~
\end{lstlisting}

\subsubsection*{How to fix it}

Make the value a compile time constant, or move the construction needing it to run time.

\errorcode{E4257}{unknown length of expansion for type \errhole{}}
\label{sec:codes:e4257}

Raised during \textbf{validation}, from \texttt{ValidateErrorMessage::\allowbreak{}UNKNOWN\_\allowbreak{}LENGTH\_\allowbreak{}OF\_\allowbreak{}EXPANSION} (\textit{src/\allowbreak{}ymirc/\allowbreak{}semantic/\allowbreak{}validator/\allowbreak{}errors.yr}).

\subsubsection*{What it means}

An expansion was written over a type whose length the compiler cannot determine.

\subsubsection*{Example}

\textit{test\_\allowbreak{}resources/\allowbreak{}lit\_\allowbreak{}slices/\allowbreak{}test15.yr}:

%% check: skip
\begin{lstlisting}[style=coloredverbatimError]
in test15;

fn foo (_ : [c8], _ : i32) {}
fn foo (_ : i32) {}
fn foo (_ : [c8]) {}

fn foo {T, U...} (a : T, b : U) {
    foo ("[", a);
    cte if ((U,)::arity > 1us) {
        for i in b {
            foo (", ", i);
        }
    } else {
        foo (", ", b);
    }
    foo ("]");
}

fn main () {
    let mut c = 20;
    let a = copy [0 ; c];
    foo (expand a [0 .. c]);
}
\end{lstlisting}

\begin{lstlisting}[style=bashVerb, breaklines=true, escapechar=~]
Error[E4257] : unknown length of expansion for type mut [i32]
 --> test_resources/lit_slices/test15.yr:(22,10)
22  ~\lstbox{\textSFxi{}}~     foo (expand a [0 .. c]);
    ~\lstbox{\textSFv{}}~          ^^^^^^
    ~\lstbox{\textSFii{}}~~\lstbox{\textSFx{}}~~\lstbox{\textSFx{}}~~\lstbox{\textSFx{}}~~\lstbox{\textSFx{}}~~\lstbox{\textSFx{}}~~\lstbox{\textSFx{}}~~\lstbox{\textSFx{}}~
\end{lstlisting}

\subsubsection*{How to fix it}

Expand something whose arity is statically known -- a tuple, or a template variadic parameter.

\errorcode{E4258}{unknown \_\_pragma expression \errhole{}}
\label{sec:codes:e4258}

Raised during \textbf{validation}, from \texttt{ValidateErrorMessage::\allowbreak{}UNKNOWN\_\allowbreak{}PRAGMA} (\textit{src/\allowbreak{}ymirc/\allowbreak{}semantic/\allowbreak{}validator/\allowbreak{}errors.yr}).

\subsubsection*{What it means}

The name after \tokennolst{\_\_pragma} is not one the compiler knows.

\subsubsection*{Example}

\textit{test\_\allowbreak{}resources/\allowbreak{}pragma/\allowbreak{}test11.yr}:

%% check: skip
\begin{lstlisting}[style=coloredverbatimError]
/*
 * An unregistered pragma name is rejected, in value position as well as in type
 * position. Being a name the lexer knows(`pass_lazy` used to be listed in
 * PragmaKeys with no handler behind it) changes nothing : the registry is what
 * decides.
 */

fn valuePosition () {
    let _x_ = __pragma!does_not_exist (12);
}

fn typePosition ()-> __pragma!does_not_exist (i32) {
    12
}

fn passLazy () {
    let _x_ = __pragma!pass_lazy (12);
}
\end{lstlisting}

\begin{lstlisting}[style=bashVerb, breaklines=true, escapechar=~]
Error[E4258] : unknown __pragma expression does_not_exist
 --> test_resources/pragma/test11.yr:(9,15)
 9  ~\lstbox{\textSFxi{}}~     let _x_ = __pragma!does_not_exist (12);
    ~\lstbox{\textSFv{}}~               ^^^^^^^^
    ~\lstbox{\textSFii{}}~~\lstbox{\textSFx{}}~~\lstbox{\textSFx{}}~~\lstbox{\textSFx{}}~~\lstbox{\textSFx{}}~~\lstbox{\textSFx{}}~~\lstbox{\textSFx{}}~~\lstbox{\textSFx{}}~
\end{lstlisting}

\subsubsection*{How to fix it}

Check the spelling against the pragmas the compiler defines.

\errorcode{E4261}{unresolved template parameter \errhole{}}
\label{sec:codes:e4261}

where \tokennolst{\textless{}\textgreater{}} stands for the template parameter, as it is declared.

Raised during \textbf{validation}, from \texttt{ValidateErrorMessage::\allowbreak{}UNRESOLVED\_\allowbreak{}TEMPLATE} (\textit{src/\allowbreak{}ymirc/\allowbreak{}semantic/\allowbreak{}validator/\allowbreak{}errors.yr}).

\subsubsection*{What it means}

A template was used without being resolved to a concrete instantiation: nothing gave a value or a type to the named parameter, neither an explicit argument nor the arguments of the call. There is one error per parameter left unresolved.

\subsubsection*{Example}

\textit{test\_\allowbreak{}resources/\allowbreak{}templates/\allowbreak{}explicit/\allowbreak{}test6.yr}:

%% check: skip
\begin{lstlisting}[style=coloredverbatimError]
trait X {}
class A {
    pub self () {
        ;
    }
}

fn foo {Z impl X} (a : Z) {
    a;
}

fn foo {Z impl X} () {
    ;
}

fn main () {
    let a = copy A ();
    foo (a);
    foo!{&A} ();
}
\end{lstlisting}

\begin{lstlisting}[style=bashVerb, breaklines=true, escapechar=~]
Error[E4261] : unresolved template parameter Z impl X
 --> test_resources/templates/explicit/test6.yr:(12,9)
12  ~\lstbox{\textSFxi{}}~ fn foo {Z impl X} () {
    ~\lstbox{\textSFv{}}~         ^
    ~\lstbox{\textSFii{}}~~\lstbox{\textSFx{}}~~\lstbox{\textSFx{}}~~\lstbox{\textSFx{}}~~\lstbox{\textSFx{}}~~\lstbox{\textSFx{}}~~\lstbox{\textSFx{}}~~\lstbox{\textSFx{}}~
\end{lstlisting}

\subsubsection*{How to fix it}

Give the template its arguments, explicitly or in a form the compiler can deduce them from.

\errorcode{E4271}{template specialization expected a type not the value \errhole{}}
\label{sec:codes:e4271}

Raised during \textbf{validation}, from \texttt{ValidateErrorMessage::\allowbreak{}USE\_\allowbreak{}AS\_\allowbreak{}TYPE\_\allowbreak{}TEMPLATE} (\textit{src/\allowbreak{}ymirc/\allowbreak{}semantic/\allowbreak{}validator/\allowbreak{}errors.yr}).

\subsubsection*{What it means}

A template parameter declared as a type was given a value.

\subsubsection*{Example}

\textit{test\_\allowbreak{}resources/\allowbreak{}templates/\allowbreak{}implicit/\allowbreak{}test18.yr}:

%% check: skip
\begin{lstlisting}[style=coloredverbatimError]
record A {N : u32} {
    pub self () {}
}

fn foo {N} (a : A!{N}) {
    a;
}


fn main () {
    let a = A!{2u32} ();
    foo (a);
}
\end{lstlisting}

\begin{lstlisting}[style=bashVerb, breaklines=true, escapechar=~]
Error[E4271] : template specialization expected a type not the value 2u32
 --> test_resources/templates/implicit/test18.yr:(11,16)
11  ~\lstbox{\textSFxi{}}~     let a = A!{2u32} ();
    ~\lstbox{\textSFv{}}~                ^
    ~\lstbox{\textSFxi{}}~ Note :
    ~\lstbox{\textSFxi{}}~  --> test_resources/templates/implicit/test18.yr:(5,9)
    ~\lstbox{\textSFxi{}}~  5  ~\lstbox{\textSFxi{}}~ fn foo {N} (a : A!{N}) {
    ~\lstbox{\textSFxi{}}~     ~\lstbox{\textSFv{}}~         ^
    ~\lstbox{\textSFxi{}}~     ~\lstbox{\textSFii{}}~~\lstbox{\textSFx{}}~~\lstbox{\textSFx{}}~~\lstbox{\textSFx{}}~~\lstbox{\textSFx{}}~~\lstbox{\textSFx{}}~~\lstbox{\textSFx{}}~~\lstbox{\textSFx{}}~
    ~\lstbox{\textSFii{}}~~\lstbox{\textSFx{}}~~\lstbox{\textSFx{}}~~\lstbox{\textSFx{}}~~\lstbox{\textSFx{}}~~\lstbox{\textSFx{}}~~\lstbox{\textSFvii{}}~~\lstbox{\textSFx{}}~
\end{lstlisting}

\subsubsection*{How to fix it}

Pass a type, or declare the parameter as a value parameter.

\errorcode{E4305}{\_\_pragma \errhole{} is not enclosed in a \errhole{}}
\label{sec:codes:e4305}

Raised during \textbf{validation}, from \texttt{ValidateErrorMessage::\allowbreak{}NO\_\allowbreak{}ENCLOSING\_\allowbreak{}PRAGMA} (\textit{src/\allowbreak{}ymirc/\allowbreak{}semantic/\allowbreak{}validator/\allowbreak{}errors.yr}).

\subsubsection*{What it means}

\tokennolst{\_\_pragma!module}, \tokennolst{\_\_pragma!function} and \tokennolst{\_\_pragma!decl} name the symbol they are written inside of. Each walks the symbols enclosing the pragma outwards until it meets the kind it names -- a module, a function, or a class, record, entity, trait, enum or union -- and raises this error when it reaches the top without finding one.

\tokennolst{\_\_pragma!decl} is the one that usually raises it: a pragma written in a plain function of a module is enclosed in no declaration at all. \tokennolst{\_\_pragma!function} raises it wherever an expression is validated outside any frame, such as the value of an enumeration member.

\tokennolst{\_\_pragma!file} and \tokennolst{\_\_pragma!line} read the location of the pragma itself, and never raise it.

\subsubsection*{Example}

\textit{test\_\allowbreak{}resources/\allowbreak{}diagnostics/\allowbreak{}test149.yr}:

%% check: skip
\begin{lstlisting}[style=coloredverbatimError]
fn main () {
    let _x_ = __pragma!decl ();
}
\end{lstlisting}

\begin{lstlisting}[style=bashVerb, breaklines=true, escapechar=~]
Error[E4305] : __pragma decl is not enclosed in a declaration
 --> test_resources/diagnostics/test149.yr:(2,15)
 2  ~\lstbox{\textSFxi{}}~     let _x_ = __pragma!decl ();
    ~\lstbox{\textSFv{}}~               ^^^^^^^^
    ~\lstbox{\textSFxi{}}~
    = when validating test149::main -- (test_resources/diagnostics/test149.yr:(1,4))
    ~\lstbox{\textSFii{}}~~\lstbox{\textSFx{}}~~\lstbox{\textSFx{}}~~\lstbox{\textSFx{}}~~\lstbox{\textSFx{}}~~\lstbox{\textSFx{}}~~\lstbox{\textSFx{}}~~\lstbox{\textSFx{}}~
\end{lstlisting}

\subsubsection*{How to fix it}

Move the pragma inside the kind of symbol it names, or use one that is always available at that point:

%% check: skip
\begin{lstlisting}[style=coloredverbatim]
class A {
    pub self () {}

    pub fn foo (self) {
        let _x_ = __pragma!decl ();     // "main::A"
        let _y_ = __pragma!function (); // "main::A::foo"
        let _z_ = __pragma!module ();   // "main"
    }
}
\end{lstlisting}

\errorcode{E4316}{\errhole{} is not an enum type}
\label{sec:codes:e4316}

Raised during \textbf{validation}, from \texttt{ValidateErrorMessage::\allowbreak{}NOT\_\allowbreak{}AN\_\allowbreak{}ENUM} (\textit{src/\allowbreak{}ymirc/\allowbreak{}semantic/\allowbreak{}validator/\allowbreak{}errors.yr}).

\subsubsection*{What it means}

A template parameter written \tokennolst{enum T} constrains the kind of type that may be bound to it, the way \tokennolst{class T} or \tokennolst{tuple T} do: it accepts an enumeration and nothing else. The constraint is checked wherever the argument comes from -- written explicitly at the call site, or deduced from a parameter typed \tokennolst{T}.

\tokennolst{enum T} is the one kind that binds the wrapper rather than what it wraps: \tokennolst{T} is the enumeration, not the type its members are represented by, since the enumeration is what the constraint asked for.

What fails is the candidate, not the program: a template whose parameter cannot be bound is discarded, so this diagnostic is reported as the reason under the error naming the call, and the call fails only when no candidate is left.

\subsubsection*{Example}

\textit{test\_\allowbreak{}resources/\allowbreak{}type\_\allowbreak{}union/\allowbreak{}test64.yr}:

%% check: skip
\begin{lstlisting}[style=coloredverbatimError]
fn keep {enum T} (x : T)-> T {
    x
}

fn main () {
    keep (12);
}
\end{lstlisting}

\begin{lstlisting}[style=bashVerb, breaklines=true, escapechar=~]
Error[E4226] : the call operator is not defined for test64::keep {enum T}(x: T)-> T and {12: i32}
 --> test_resources/type_union/test64.yr:(8,10)
 8  ~\lstbox{\textSFxi{}}~     keep (12);
    ~\lstbox{\textSFv{}}~          ^
 --> test_resources/type_union/test64.yr:(3,4)
 3  ~\lstbox{\textSFxi{}}~ fn keep {enum T} (x : T)-> T {
    ~\lstbox{\textSFv{}}~    ----  note: candidate test64::keep {enum T}(x: T)-> T
    ~\lstbox{\textSFxi{}}~       ~\lstbox{\textSFii{}}~~\lstbox{\textSFx{}}~ error: i32 is not an enum type
    ~\lstbox{\textSFxi{}}~
    = when validating test64::main -- (test_resources/type_union/test64.yr:(7,4))
    ~\lstbox{\textSFii{}}~~\lstbox{\textSFx{}}~~\lstbox{\textSFx{}}~~\lstbox{\textSFx{}}~~\lstbox{\textSFx{}}~~\lstbox{\textSFx{}}~~\lstbox{\textSFx{}}~~\lstbox{\textSFx{}}~
\end{lstlisting}

\subsubsection*{How to fix it}

Pass an enumeration, or drop the \tokennolst{enum} constraint when the parameter is meant to accept any type.

%% Generated by tools/gen_error_codes.py from docs/errors/ of the compiler;
%% edit the compiler's pages or tools/error_themes.txt, then rerun it.

\clearpage
\section{Laziness, generators and comprehensions}
\label{sec:codes:lazy}

The book's chapter \textit{Laziness} specifies the rules behind these codes.

\begin{longtable}{@{}l p{0.78\linewidth} r@{}}
\toprule Code & Message & Page\\ \midrule \endhead
\hyperref[sec:codes:e4087]{E4087} & cannot construct a lazy value of type \errhole{} implicitly & \pageref{sec:codes:e4087}\\
\hyperref[sec:codes:e4108]{E4108} & validation of lazy closure function fails to generate value of type \errhole{} & \pageref{sec:codes:e4108}\\
\hyperref[sec:codes:e4109]{E4109} & the size of the list comprehension, constructed from type \errhole{}, must be determinable at compile time & \pageref{sec:codes:e4109}\\
\hyperref[sec:codes:e4249]{E4249} & the construction of a lazy value has no effect & \pageref{sec:codes:e4249}\\
\hyperref[sec:codes:e4291]{E4291} & yield statement is not within a generator & \pageref{sec:codes:e4291}\\
\hyperref[sec:codes:e4292]{E4292} & a yield statement must yield a value of type \errhole{} & \pageref{sec:codes:e4292}\\
\hyperref[sec:codes:e4293]{E4293} & cannot yield inside a try block, a catch or a scope guard & \pageref{sec:codes:e4293}\\
\hyperref[sec:codes:e4294]{E4294} & a generator cannot return a value, its values are yielded & \pageref{sec:codes:e4294}\\
\hyperref[sec:codes:e4295]{E4295} & a generator cannot throw exceptions & \pageref{sec:codes:e4295}\\
\hyperref[sec:codes:e4296]{E4296} & the parameter \errhole{} of a generator cannot be declared \errhole{} & \pageref{sec:codes:e4296}\\
\hyperref[sec:codes:e4297]{E4297} & a generator cannot hold a value of type \errhole{}, which has a destructor & \pageref{sec:codes:e4297}\\
\hyperref[sec:codes:e4299]{E4299} & yield in expects a generator, not a value of type \errhole{} & \pageref{sec:codes:e4299}\\
\hyperref[sec:codes:e4300]{E4300} & copying a generator comprehension has no effect & \pageref{sec:codes:e4300}\\
\hyperref[sec:codes:e4301]{E4301} & cannot \errhole{} out of the element of a generator comprehension & \pageref{sec:codes:e4301}\\
\hyperref[sec:codes:e4302]{E4302} & the size of a list comprehension filtered by if is not determinable at compile time & \pageref{sec:codes:e4302}\\
\hyperref[sec:codes:e4303]{E4303} & the filter of the list comprehension, constructed from type \errhole{}, must be determinable at compile time & \pageref{sec:codes:e4303}\\
\bottomrule
\end{longtable}

\errorcode{E4087}{cannot construct a lazy value of type \errhole{} implicitly}
\label{sec:codes:e4087}

Raised during \textbf{validation}, from \texttt{ValidateErrorMessage::\allowbreak{}IMPLICIT\_\allowbreak{}LAZY} (\textit{src/\allowbreak{}ymirc/\allowbreak{}semantic/\allowbreak{}validator/\allowbreak{}errors.yr}).

\subsubsection*{What it means}

A value was used where a lazy one is expected, without the \tokennolst{lazy} keyword. A lazy value defers its computation, which changes when the expression runs, so it is never built implicitly.

\subsubsection*{Example}

\textit{test\_\allowbreak{}resources/\allowbreak{}diagnostics/\allowbreak{}test41.yr}:

%% check: skip
\begin{lstlisting}[style=coloredverbatimError]
fn foo (lazy x: i32)-> i32 { x }

fn main () { foo (1); }
\end{lstlisting}

\begin{lstlisting}[style=bashVerb, breaklines=true, escapechar=~]
Error[E4087] : cannot construct a lazy value of type i32 implicitly
 --> test_resources/diagnostics/test41.yr:(3,19)
 3  ~\lstbox{\textSFxi{}}~ fn main () { foo (1); }
    ~\lstbox{\textSFv{}}~                   ^  note: for parameter lazy x: i32
    ~\lstbox{\textSFii{}}~~\lstbox{\textSFx{}}~~\lstbox{\textSFx{}}~~\lstbox{\textSFx{}}~~\lstbox{\textSFx{}}~~\lstbox{\textSFx{}}~~\lstbox{\textSFx{}}~~\lstbox{\textSFx{}}~
\end{lstlisting}

\subsubsection*{How to fix it}

Write \tokennolst{lazy} before the expression.

\errorcode{E4108}{validation of lazy closure function fails to generate value of type \errhole{}}
\label{sec:codes:e4108}

Raised during \textbf{validation}, from \texttt{ValidateErrorMessage::\allowbreak{}LAZY\_\allowbreak{}VALIDATION} (\textit{src/\allowbreak{}ymirc/\allowbreak{}semantic/\allowbreak{}validator/\allowbreak{}errors.yr}).

\subsubsection*{What it means}

A lazy value could not be built: validating the closure that defers the computation failed to produce a value of the expected type.

\subsubsection*{Example}

\textit{test\_\allowbreak{}resources/\allowbreak{}diagnostics/\allowbreak{}test109.yr}:

%% check: skip
\begin{lstlisting}[style=coloredverbatimError]
// a lazy value assigning a variable it encloses
fn main () {
    let mut x = 1;
    let y = lazy { x = 2; 1 };
    let _ = y;
    x = 3;
}
\end{lstlisting}

\begin{lstlisting}[style=bashVerb, breaklines=true, escapechar=~]
Error[E4108] : validation of lazy closure function fails to generate value of type mut i32
 --> test_resources/diagnostics/test109.yr:(4,13)
 4  ~\lstbox{\textSFxi{}}~     let y = lazy { x = 2; 1 };
    ~\lstbox{\textSFv{}}~             ^^^^   -  error: x is enclosed, and cannot be used as a lvalue
    ~\lstbox{\textSFxi{}}~
    = when validating test109::main -- (test_resources/diagnostics/test109.yr:(2,4))
    ~\lstbox{\textSFii{}}~~\lstbox{\textSFx{}}~~\lstbox{\textSFx{}}~~\lstbox{\textSFx{}}~~\lstbox{\textSFx{}}~~\lstbox{\textSFx{}}~~\lstbox{\textSFx{}}~~\lstbox{\textSFx{}}~
Error[E4095] : incompatible types error and error
 --> test_resources/diagnostics/test109.yr:(5,9)
 5  ~\lstbox{\textSFxi{}}~     let _ = y;
    ~\lstbox{\textSFv{}}~         ^   -
    ~\lstbox{\textSFxi{}}~
    = when validating test109::main -- (test_resources/diagnostics/test109.yr:(2,4))
    ~\lstbox{\textSFii{}}~~\lstbox{\textSFx{}}~~\lstbox{\textSFx{}}~~\lstbox{\textSFx{}}~~\lstbox{\textSFx{}}~~\lstbox{\textSFx{}}~~\lstbox{\textSFx{}}~~\lstbox{\textSFx{}}~
\end{lstlisting}

\subsubsection*{How to fix it}

Check the expression the \tokennolst{lazy} applies to; the errors under this diagnostic say what went wrong inside it.

\errorcode{E4109}{the size of the list comprehension, constructed from type \errhole{}, must be determinable at compile time}
\label{sec:codes:e4109}

Raised during \textbf{validation}, from \texttt{ValidateErrorMessage::\allowbreak{}LIST\_\allowbreak{}COMPR\_\allowbreak{}SIZE\_\allowbreak{}CTE} (\textit{src/\allowbreak{}ymirc/\allowbreak{}semantic/\allowbreak{}validator/\allowbreak{}errors.yr}).

\subsubsection*{What it means}

The length of a list comprehension is not computable at compile time, and the type it is being built into needs a statically known size.

\subsubsection*{Example}

\textit{test\_\allowbreak{}resources/\allowbreak{}for\_\allowbreak{}loops/\allowbreak{}lst\_\allowbreak{}compr/\allowbreak{}test12.yr}:

%% check: skip
\begin{lstlisting}[style=coloredverbatimError]
fn main () {
    let mut a = copy [1, 2, 3];
    let _ = [i for i in a];
}
\end{lstlisting}

\begin{lstlisting}[style=bashVerb, breaklines=true, escapechar=~]
Error[E4109] : the size of the list comprehension, constructed from type mut [i32], must be determinable at compile time
 --> test_resources/for_loops/lst_compr/test12.yr:(3,25)
 3  ~\lstbox{\textSFxi{}}~     let _ = [i for i in a];
    ~\lstbox{\textSFv{}}~                         ^  note: did you mean to enclose it within a copy?
    ~\lstbox{\textSFii{}}~~\lstbox{\textSFx{}}~~\lstbox{\textSFx{}}~~\lstbox{\textSFx{}}~~\lstbox{\textSFx{}}~~\lstbox{\textSFx{}}~~\lstbox{\textSFx{}}~~\lstbox{\textSFx{}}~
\end{lstlisting}

\subsubsection*{How to fix it}

Build a slice, whose length is a run time value, or make the bounds of the comprehension compile time constants.

\errorcode{E4249}{the construction of a lazy value has no effect}
\label{sec:codes:e4249}

Raised during \textbf{validation}, from \texttt{ValidateErrorMessage::\allowbreak{}UNECESSARY\_\allowbreak{}LAZY} (\textit{src/\allowbreak{}ymirc/\allowbreak{}semantic/\allowbreak{}validator/\allowbreak{}errors.yr}).

\subsubsection*{What it means}

\tokennolst{lazy} was applied where the value is used immediately, so nothing is deferred.

\subsubsection*{Example}

\textit{test\_\allowbreak{}resources/\allowbreak{}lazyness/\allowbreak{}test16.yr}:

%% check: skip
\begin{lstlisting}[style=coloredverbatimError]
fn main () {
    match (lazy 12) {
        x => { x; }
    }
}
\end{lstlisting}

\begin{lstlisting}[style=bashVerb, breaklines=true, escapechar=~]
Error[E4249] : the construction of a lazy value has no effect
 --> test_resources/lazyness/test16.yr:(2,12)
 2  ~\lstbox{\textSFxi{}}~     match (lazy 12) {
    ~\lstbox{\textSFv{}}~            ^^^^
    ~\lstbox{\textSFii{}}~~\lstbox{\textSFx{}}~~\lstbox{\textSFx{}}~~\lstbox{\textSFx{}}~~\lstbox{\textSFx{}}~~\lstbox{\textSFx{}}~~\lstbox{\textSFx{}}~~\lstbox{\textSFx{}}~
\end{lstlisting}

\subsubsection*{How to fix it}

Drop the \tokennolst{lazy}.

\errorcode{E4291}{yield statement is not within a generator}
\label{sec:codes:e4291}

Raised during \textbf{validation}, from \texttt{ValidateErrorMessage::\allowbreak{}YIELD\_\allowbreak{}NO\_\allowbreak{}GENERATOR} (\textit{src/\allowbreak{}ymirc/\allowbreak{}semantic/\allowbreak{}validator/\allowbreak{}errors.yr}).

\subsubsection*{What it means}

A \tokennolst{yield} appears in a function that is not a generator. Only a function introduced by \tokennolst{yield} can yield values; a lambda nested in a generator is a function of its own, and cannot yield on behalf of the generator either.

\subsubsection*{Example}

\textit{test\_\allowbreak{}resources/\allowbreak{}generators/\allowbreak{}test6.yr}:

%% check: skip
\begin{lstlisting}[style=coloredverbatimError]
// yield outside a generator
fn notGenerator ()-> i32 {
    yield 1;
    2
}
\end{lstlisting}

\begin{lstlisting}[style=bashVerb, breaklines=true, escapechar=~]
Error[E4291] : yield statement is not within a generator
 --> test_resources/generators/test6.yr:(3,5)
 3  ~\lstbox{\textSFxi{}}~     yield 1;
    ~\lstbox{\textSFv{}}~     ^^^^^
    ~\lstbox{\textSFxi{}}~
    = when validating test6::notGenerator -- (test_resources/generators/test6.yr:(2,4))
    ~\lstbox{\textSFii{}}~~\lstbox{\textSFx{}}~~\lstbox{\textSFx{}}~~\lstbox{\textSFx{}}~~\lstbox{\textSFx{}}~~\lstbox{\textSFx{}}~~\lstbox{\textSFx{}}~~\lstbox{\textSFx{}}~
\end{lstlisting}

\subsubsection*{How to fix it}

Introduce the function with \tokennolst{yield} instead of \tokennolst{fn}, its declared return type being the type of the yielded values.

\errorcode{E4292}{a yield statement must yield a value of type \errhole{}}
\label{sec:codes:e4292}

Raised during \textbf{validation}, from \texttt{ValidateErrorMessage::\allowbreak{}YIELD\_\allowbreak{}NO\_\allowbreak{}VALUE} (\textit{src/\allowbreak{}ymirc/\allowbreak{}semantic/\allowbreak{}validator/\allowbreak{}errors.yr}).

\subsubsection*{What it means}

A \tokennolst{yield} without a value hands nothing to the consumer of the generator. The type in the message is the one the generator declares.

\subsubsection*{Example}

\textit{test\_\allowbreak{}resources/\allowbreak{}generators/\allowbreak{}test7.yr}:

%% check: skip
\begin{lstlisting}[style=coloredverbatimError]
// a yield without a value
yield gen ()-> i32 {
    yield;
}
\end{lstlisting}

\begin{lstlisting}[style=bashVerb, breaklines=true, escapechar=~]
Error[E4292] : a yield statement must yield a value of type i32
 --> test_resources/generators/test7.yr:(3,5)
 3  ~\lstbox{\textSFxi{}}~     yield;
    ~\lstbox{\textSFv{}}~     ^^^^^
    ~\lstbox{\textSFxi{}}~
    = when validating test7::gen -- (test_resources/generators/test7.yr:(2,7))
    ~\lstbox{\textSFii{}}~~\lstbox{\textSFx{}}~~\lstbox{\textSFx{}}~~\lstbox{\textSFx{}}~~\lstbox{\textSFx{}}~~\lstbox{\textSFx{}}~~\lstbox{\textSFx{}}~~\lstbox{\textSFx{}}~
\end{lstlisting}

\subsubsection*{How to fix it}

Give the \tokennolst{yield} a value. To end the generator early, use \tokennolst{return;}.

\errorcode{E4293}{cannot yield inside a try block, a catch or a scope guard}
\label{sec:codes:e4293}

Raised during \textbf{validation}, from \texttt{ValidateErrorMessage::\allowbreak{}YIELD\_\allowbreak{}PROTECTED\_\allowbreak{}REGION} (\textit{src/\allowbreak{}ymirc/\allowbreak{}semantic/\allowbreak{}validator/\allowbreak{}errors.yr}).

\subsubsection*{What it means}

A generator is suspended at each \tokennolst{yield} and resumed by a jump back to it. A jump cannot re-enter a try block, a catch or a block guarded by a scope guard, as the handler of the region would not be installed again.

\subsubsection*{Example}

\textit{test\_\allowbreak{}resources/\allowbreak{}generators/\allowbreak{}test8.yr}:

%% check: skip
\begin{lstlisting}[style=coloredverbatimError]
// yield inside a try block
class Err over Exception {
    pub self () {}
}

fn mayThrow ()
    throws Err
{
    throw copy Err ();
}

yield gen ()-> i32 {
    {
        mayThrow ();
        yield 1;
    } catch {
        _: &Err => {}
    }
}
\end{lstlisting}

\begin{lstlisting}[style=bashVerb, breaklines=true, escapechar=~]
Error[E4293] : cannot yield inside a try block, a catch or a scope guard
 --> test_resources/generators/test8.yr:(13,5)
13  ~\lstbox{\textSFxi{}}~     {
    ~\lstbox{\textSFv{}}~     ^
14  ~\lstbox{\textSFxi{}}~         mayThrow ();
15  ~\lstbox{\textSFxi{}}~         yield 1;
    ~\lstbox{\textSFv{}}~         -----
    ~\lstbox{\textSFxi{}}~
    = when validating test8::gen -- (test_resources/generators/test8.yr:(12,7))
    ~\lstbox{\textSFii{}}~~\lstbox{\textSFx{}}~~\lstbox{\textSFx{}}~~\lstbox{\textSFx{}}~~\lstbox{\textSFx{}}~~\lstbox{\textSFx{}}~~\lstbox{\textSFx{}}~~\lstbox{\textSFx{}}~
\end{lstlisting}

\subsubsection*{How to fix it}

Move the \tokennolst{yield} out of the protected region, storing what it yields in a variable of the generator.

\errorcode{E4294}{a generator cannot return a value, its values are yielded}
\label{sec:codes:e4294}

Raised during \textbf{validation}, from \texttt{ValidateErrorMessage::\allowbreak{}RETURN\_\allowbreak{}VALUE\_\allowbreak{}GENERATOR} (\textit{src/\allowbreak{}ymirc/\allowbreak{}semantic/\allowbreak{}validator/\allowbreak{}errors.yr}).

\subsubsection*{What it means}

The declared return type of a generator is the type of the values it yields, and a call of a generator returns its handle. A \tokennolst{return} can only end the generator.

\subsubsection*{Example}

\textit{test\_\allowbreak{}resources/\allowbreak{}generators/\allowbreak{}test10.yr}:

%% check: skip
\begin{lstlisting}[style=coloredverbatimError]
// a generator returning a value
yield gen ()-> i32 {
    yield 1;
    return 2;
}
\end{lstlisting}

\begin{lstlisting}[style=bashVerb, breaklines=true, escapechar=~]
Error[E4294] : a generator cannot return a value, its values are yielded
 --> test_resources/generators/test10.yr:(4,5)
 4  ~\lstbox{\textSFxi{}}~     return 2;
    ~\lstbox{\textSFv{}}~     ^^^^^^
    ~\lstbox{\textSFxi{}}~
    = when validating test10::gen -- (test_resources/generators/test10.yr:(2,7))
    ~\lstbox{\textSFii{}}~~\lstbox{\textSFx{}}~~\lstbox{\textSFx{}}~~\lstbox{\textSFx{}}~~\lstbox{\textSFx{}}~~\lstbox{\textSFx{}}~~\lstbox{\textSFx{}}~~\lstbox{\textSFx{}}~
\end{lstlisting}

\subsubsection*{How to fix it}

Yield the value, then \tokennolst{return;} if the generator has to stop there.

\errorcode{E4295}{a generator cannot throw exceptions}
\label{sec:codes:e4295}

Raised during \textbf{validation}, from \texttt{ValidateErrorMessage::\allowbreak{}GENERATOR\_\allowbreak{}THROWS} (\textit{src/\allowbreak{}ymirc/\allowbreak{}semantic/\allowbreak{}validator/\allowbreak{}errors.yr}).

\subsubsection*{What it means}

The body of a generator runs while its consumer iterates over it, where the exceptions it declares could not be tracked. A generator method, declared by \tokennolst{yield} in a class, a record or a trait, is rejected the same way.

\subsubsection*{Example}

\textit{test\_\allowbreak{}resources/\allowbreak{}generators/\allowbreak{}test11.yr}:

%% check: skip
\begin{lstlisting}[style=coloredverbatimError]
// a generator declaring exceptions
class Err over Exception {
    pub self () {}
}

yield gen ()-> i32
    throws Err
{
    yield 1;
    throw copy Err ();
}
\end{lstlisting}

\begin{lstlisting}[style=bashVerb, breaklines=true, escapechar=~]
Error[E4295] : a generator cannot throw exceptions
 --> test_resources/generators/test11.yr:(7,5)
 7  ~\lstbox{\textSFxi{}}~     throws Err
    ~\lstbox{\textSFv{}}~     ^^^^^^
    ~\lstbox{\textSFxi{}}~
    = when validating test11::gen -- (test_resources/generators/test11.yr:(6,7))
    ~\lstbox{\textSFii{}}~~\lstbox{\textSFx{}}~~\lstbox{\textSFx{}}~~\lstbox{\textSFx{}}~~\lstbox{\textSFx{}}~~\lstbox{\textSFx{}}~~\lstbox{\textSFx{}}~~\lstbox{\textSFx{}}~
\end{lstlisting}

\subsubsection*{How to fix it}

Catch the exceptions inside the generator, outside of the region of any \tokennolst{yield}, cf. \hyperref[sec:codes:e4293]{E4293}.

\errorcode{E4296}{the parameter \errhole{} of a generator cannot be declared \errhole{}}
\label{sec:codes:e4296}

Raised during \textbf{validation}, from \texttt{ValidateErrorMessage::\allowbreak{}GENERATOR\_\allowbreak{}PARAM} (\textit{src/\allowbreak{}ymirc/\allowbreak{}semantic/\allowbreak{}validator/\allowbreak{}errors.yr}).

\subsubsection*{What it means}

The parameters of a generator are stored in its frame, which outlives the call creating it. A \tokennolst{ref} or \tokennolst{lazy} parameter refers to the frame of the caller, that may be gone when the generator is resumed. This applies to generator methods as well; the \tokennolst{self} of a record generator is copied into its frame, so a \tokennolst{mut self}, whose changes the caller would never see, is rejected.

\subsubsection*{Example}

\textit{test\_\allowbreak{}resources/\allowbreak{}generators/\allowbreak{}test12.yr}:

%% check: skip
\begin{lstlisting}[style=coloredverbatimError]
// a generator taking a reference parameter
yield gen (ref x: i32)-> i32 {
    yield x;
}
\end{lstlisting}

\begin{lstlisting}[style=bashVerb, breaklines=true, escapechar=~]
Error[E4296] : the parameter x of a generator cannot be declared ref
 --> test_resources/generators/test12.yr:(2,16)
 2  ~\lstbox{\textSFxi{}}~ yield gen (ref x: i32)-> i32 {
    ~\lstbox{\textSFv{}}~                ^
    ~\lstbox{\textSFxi{}}~
    = when validating test12::gen -- (test_resources/generators/test12.yr:(2,7))
    ~\lstbox{\textSFii{}}~~\lstbox{\textSFx{}}~~\lstbox{\textSFx{}}~~\lstbox{\textSFx{}}~~\lstbox{\textSFx{}}~~\lstbox{\textSFx{}}~~\lstbox{\textSFx{}}~~\lstbox{\textSFx{}}~
\end{lstlisting}

\subsubsection*{How to fix it}

Take the parameter by value.

\errorcode{E4297}{a generator cannot hold a value of type \errhole{}, which has a destructor}
\label{sec:codes:e4297}

Raised during \textbf{validation}, from \texttt{ValidateErrorMessage::\allowbreak{}GENERATOR\_\allowbreak{}MOVABLE} (\textit{src/\allowbreak{}ymirc/\allowbreak{}semantic/\allowbreak{}validator/\allowbreak{}errors.yr}).

\subsubsection*{What it means}

A generator can be abandoned while suspended, and its frame is then collected without its pending destructors ever running. A value whose type has a destructor cannot be held by a generator for that reason. This includes the \tokennolst{self} of a generator method, so an entity with a destructor cannot take the default generator method of a trait.

\subsubsection*{Example}

\textit{test\_\allowbreak{}resources/\allowbreak{}generators/\allowbreak{}test15.yr}:

%% check: skip
\begin{lstlisting}[style=coloredverbatimError]
// a generator holding a value with a destructor
entity Handle {
    pub self () {}

    __dtor (mut self) {}
}

yield gen ()-> i32 {
    let h = Handle ();
    yield 1;
}
\end{lstlisting}

\begin{lstlisting}[style=bashVerb, breaklines=true, escapechar=~]
Error[E4297] : a generator cannot hold a value of type mut test15::Handle, which has a destructor
 --> test_resources/generators/test15.yr:(9,20)
 9  ~\lstbox{\textSFxi{}}~     let h = Handle ();
    ~\lstbox{\textSFv{}}~                    ^
    ~\lstbox{\textSFxi{}}~
    = when validating test15::gen -- (test_resources/generators/test15.yr:(8,7))
    ~\lstbox{\textSFii{}}~~\lstbox{\textSFx{}}~~\lstbox{\textSFx{}}~~\lstbox{\textSFx{}}~~\lstbox{\textSFx{}}~~\lstbox{\textSFx{}}~~\lstbox{\textSFx{}}~~\lstbox{\textSFx{}}~
\end{lstlisting}

\subsubsection*{How to fix it}

Hold the value in the consumer of the generator, and pass the generator what it needs from it.

\errorcode{E4299}{yield in expects a generator, not a value of type \errhole{}}
\label{sec:codes:e4299}

Raised during \textbf{validation}, from \texttt{ValidateErrorMessage::\allowbreak{}YIELD\_\allowbreak{}IN\_\allowbreak{}NO\_\allowbreak{}GENERATOR} (\textit{src/\allowbreak{}ymirc/\allowbreak{}semantic/\allowbreak{}validator/\allowbreak{}errors.yr}).

\subsubsection*{What it means}

\tokennolst{yield in g} yields every value of the generator \tokennolst{g}, until it is over. The value after \tokennolst{in} must therefore be a generator handle, of type \tokennolst{yield-\textgreater{} T}; any other value, even an iterable one, has no values to hand over this way.

\subsubsection*{Example}

\textit{test\_\allowbreak{}resources/\allowbreak{}generators/\allowbreak{}test43.yr}:

%% check: skip
\begin{lstlisting}[style=coloredverbatimError]
// yield in only yields the values of a generator
yield all (a: [i32])-> i32 {
    yield in a;
}

fn main () {
    for x in all (copy [1, 2]) {
        x;
    }
}
\end{lstlisting}

\begin{lstlisting}[style=bashVerb, breaklines=true, escapechar=~]
Error[E4299] : yield in expects a generator, not a value of type [i32]
 --> test_resources/generators/test43.yr:(3,14)
 3  ~\lstbox{\textSFxi{}}~     yield in a;
    ~\lstbox{\textSFv{}}~              ^
    ~\lstbox{\textSFxi{}}~
    = when validating test43::all -- (test_resources/generators/test43.yr:(2,7))
    ~\lstbox{\textSFii{}}~~\lstbox{\textSFx{}}~~\lstbox{\textSFx{}}~~\lstbox{\textSFx{}}~~\lstbox{\textSFx{}}~~\lstbox{\textSFx{}}~~\lstbox{\textSFx{}}~~\lstbox{\textSFx{}}~
\end{lstlisting}

\subsubsection*{How to fix it}

Yield the values one at a time with a loop, \tokennolst{for x in a \{ yield x; \}}, or pass a generator producing them to \tokennolst{yield in}.

\errorcode{E4300}{copying a generator comprehension has no effect}
\label{sec:codes:e4300}

Raised during \textbf{validation}, from \texttt{ValidateErrorMessage::\allowbreak{}UNECESSARY\_\allowbreak{}GENERATOR\_\allowbreak{}COPY} (\textit{src/\allowbreak{}ymirc/\allowbreak{}semantic/\allowbreak{}validator/\allowbreak{}errors.yr}).

\subsubsection*{What it means}

A generator comprehension, \tokennolst{(yield x for x in it)}, builds no list: it creates a generator, which computes each value when it is iterated. There is nothing for \tokennolst{copy} to duplicate, unlike a list comprehension, where \tokennolst{copy} builds a slice instead of an array.

Only a \tokennolst{copy} applied to the comprehension itself is reported. A \tokennolst{dcopy} of a value merely holding one, \texttt{dcopy [(yield x for x in it) for \_ in 0 .. 2]}, does copy the slice, and simply leaves its generators alone.

\subsubsection*{Example}

\textit{test\_\allowbreak{}resources/\allowbreak{}generators/\allowbreak{}test57.yr}:

%% check: skip
\begin{lstlisting}[style=coloredverbatimError]
// a generator comprehension allocates its own frame, there is nothing to copy
fn main () {
    let _ = copy (yield x for x in 0 .. 3);
}
\end{lstlisting}

\begin{lstlisting}[style=bashVerb, breaklines=true, escapechar=~]
Error[E4300] : copying a generator comprehension has no effect
 --> test_resources/generators/test57.yr:(3,18)
 3  ~\lstbox{\textSFxi{}}~     let _ = copy (yield x for x in 0 .. 3);
    ~\lstbox{\textSFv{}}~                  ^
    ~\lstbox{\textSFxi{}}~
    = when validating test57::main -- (test_resources/generators/test57.yr:(2,4))
    ~\lstbox{\textSFii{}}~~\lstbox{\textSFx{}}~~\lstbox{\textSFx{}}~~\lstbox{\textSFx{}}~~\lstbox{\textSFx{}}~~\lstbox{\textSFx{}}~~\lstbox{\textSFx{}}~~\lstbox{\textSFx{}}~
\end{lstlisting}

\subsubsection*{How to fix it}

Drop the \tokennolst{copy}. To collect the values into a slice, use a list comprehension over the generator, \tokennolst{copy [x for x in (yield x for x in it)]}.

\errorcode{E4301}{cannot \errhole{} out of the element of a generator comprehension}
\label{sec:codes:e4301}

Raised during \textbf{validation}, from \texttt{ValidateErrorMessage::\allowbreak{}GENERATOR\_\allowbreak{}COMPR\_\allowbreak{}EXIT} (\textit{src/\allowbreak{}ymirc/\allowbreak{}semantic/\allowbreak{}validator/\allowbreak{}errors.yr}).

\subsubsection*{What it means}

The element of a generator comprehension, \tokennolst{(yield e for x in it)}, is not computed where the comprehension is written, but each time the generator is iterated, possibly long after the enclosing loop or function is over. A \tokennolst{break}, a \tokennolst{continue} or a \tokennolst{return} in it has nothing to leave: it cannot exit the loop or the function around the comprehension, and it cannot skip or end the values of the generator either.

A \tokennolst{yield}, or a \tokennolst{yield in}, is rejected for the same reason read the other way round: the comprehension yields the element and nothing else, so a \tokennolst{yield} written in it would add values to a generator that is not the element's to produce.

A \tokennolst{break} or a \tokennolst{continue} of a loop written inside the element itself is fine.

\subsubsection*{Example, of a \tokennolst{break}}

\textit{test\_\allowbreak{}resources/\allowbreak{}generators/\allowbreak{}test64.yr}:

%% check: skip
\begin{lstlisting}[style=coloredverbatimError]
// the element of a generator comprehension cannot break the loop enclosing the comprehension
fn main (args: [[c8]]) {
    loop {
        let _ = (yield { if (args.len > 2us) break; x } for x in 0 .. 12);
    }
}
\end{lstlisting}

\begin{lstlisting}[style=bashVerb, breaklines=true, escapechar=~]
Error[E4301] : cannot break out of the element of a generator comprehension
 --> test_resources/generators/test64.yr:(4,46)
 4  ~\lstbox{\textSFxi{}}~         let _ = (yield { if (args.len > 2us) break; x } for x in 0 .. 12);
    ~\lstbox{\textSFv{}}~                                              ^^^^^
    ~\lstbox{\textSFxi{}}~
    = when validating test64::main -- (test_resources/generators/test64.yr:(2,4))
    ~\lstbox{\textSFii{}}~~\lstbox{\textSFx{}}~~\lstbox{\textSFx{}}~~\lstbox{\textSFx{}}~~\lstbox{\textSFx{}}~~\lstbox{\textSFx{}}~~\lstbox{\textSFx{}}~~\lstbox{\textSFx{}}~
\end{lstlisting}

\subsubsection*{Example, of a \tokennolst{yield}}

\textit{test\_\allowbreak{}resources/\allowbreak{}generators/\allowbreak{}test68.yr}:

%% check: skip
\begin{lstlisting}[style=coloredverbatimError]
// the element of a generator comprehension cannot yield, it is the value the generator yields
fn main () {
    let _ = (yield { yield x; x * 2 } for x in 0 .. 3);
}
\end{lstlisting}

\begin{lstlisting}[style=bashVerb, breaklines=true, escapechar=~]
Error[E4301] : cannot yield out of the element of a generator comprehension
 --> test_resources/generators/test68.yr:(3,22)
 3  ~\lstbox{\textSFxi{}}~     let _ = (yield { yield x; x * 2 } for x in 0 .. 3);
    ~\lstbox{\textSFv{}}~                      ^^^^^
    ~\lstbox{\textSFxi{}}~
    = when validating test68::main -- (test_resources/generators/test68.yr:(2,4))
    ~\lstbox{\textSFii{}}~~\lstbox{\textSFx{}}~~\lstbox{\textSFx{}}~~\lstbox{\textSFx{}}~~\lstbox{\textSFx{}}~~\lstbox{\textSFx{}}~~\lstbox{\textSFx{}}~~\lstbox{\textSFx{}}~
\end{lstlisting}

\subsubsection*{How to fix it}

Test the condition outside the comprehension, or write the generator as a function, \texttt{yield f (it: …)-\textgreater{} T \{ for x in it \{ … yield e; \} \}}, where a loop, a \tokennolst{continue}, a \tokennolst{return} and a \tokennolst{yield} have their usual meaning. That is also the way to yield several values per element.

\errorcode{E4302}{the size of a list comprehension filtered by if is not determinable at compile time}
\label{sec:codes:e4302}

Raised during \textbf{validation}, from \texttt{ValidateErrorMessage::\allowbreak{}LIST\_\allowbreak{}COMPR\_\allowbreak{}FILTER\_\allowbreak{}SIZE} (\textit{src/\allowbreak{}ymirc/\allowbreak{}semantic/\allowbreak{}validator/\allowbreak{}errors.yr}).

\subsubsection*{What it means}

A comprehension filtered by \tokennolst{if} collects the values the filter keeps, and which ones those are is only known once the iteration ran. Its result is therefore a slice, appended to as the iteration goes, and a slice is allocated on the heap: the comprehension has to be enclosed in a \tokennolst{copy}.

Without the filter the same comprehension has the length of the iteration, so it can be an array on the stack, which is why dropping the \tokennolst{copy} is accepted there and rejected here.

A comprehension over a tuple is the exception: its iteration is unrolled during validation, so the filter is computed there and the length of the result is known. It stays an array on the stack, and needs no \tokennolst{copy} -- at the price of a filter that has to be a compile time value, or \hyperref[sec:codes:e4303]{E4303} is raised.

This is about the \textit{number} of values, not about their type: a filter whose test is not a boolean raises \hyperref[sec:codes:e4095]{E4095} instead, and a comprehension whose iteration itself has no compile time length raises \hyperref[sec:codes:e4109]{E4109}.

\subsubsection*{Example}

\textit{test\_\allowbreak{}resources/\allowbreak{}diagnostics/\allowbreak{}test146.yr}:

%% check: skip
\begin{lstlisting}[style=coloredverbatimError]
fn main () {
    let _ = [i for i in 0 .. 10 if i != 6];
}
\end{lstlisting}

\begin{lstlisting}[style=bashVerb, breaklines=true, escapechar=~]
Error[E4302] : the size of a list comprehension filtered by if is not determinable at compile time
 --> test_resources/diagnostics/test146.yr:(2,38)
 2  ~\lstbox{\textSFxi{}}~     let _ = [i for i in 0 .. 10 if i != 6];
    ~\lstbox{\textSFv{}}~                                      ^^
    ~\lstbox{\textSFxi{}}~
    = help: did you mean to enclose it within a copy?
    ~\lstbox{\textSFxi{}}~
 2  -     let _ = [i for i in 0 .. 10 if i != 6];
 2  +     let _ = copy [i for i in 0 .. 10 if i != 6];
    ~\lstbox{\textSFxi{}}~
    = when validating test146::main -- (test_resources/diagnostics/test146.yr:(1,4))
    ~\lstbox{\textSFii{}}~~\lstbox{\textSFx{}}~~\lstbox{\textSFx{}}~~\lstbox{\textSFx{}}~~\lstbox{\textSFx{}}~~\lstbox{\textSFx{}}~~\lstbox{\textSFx{}}~~\lstbox{\textSFx{}}~
\end{lstlisting}

\subsubsection*{How to fix it}

Enclose the comprehension in a \tokennolst{copy}, as the fix of the diagnostic suggests:

%% check: skip
\begin{lstlisting}[style=coloredverbatim]
let _ = copy [i for i in 0 .. 10 if i != 6];
\end{lstlisting}

\errorcode{E4303}{the filter of the list comprehension, constructed from type \errhole{}, must be determinable at compile time}
\label{sec:codes:e4303}

Raised during \textbf{validation}, from \texttt{ValidateErrorMessage::\allowbreak{}LIST\_\allowbreak{}COMPR\_\allowbreak{}FILTER\_\allowbreak{}CTE} (\textit{src/\allowbreak{}ymirc/\allowbreak{}semantic/\allowbreak{}validator/\allowbreak{}errors.yr}).

\subsubsection*{What it means}

A comprehension written between parentheses constructs a tuple, and a tuple has an arity fixed at compile time. Its iteration is unrolled during validation, so the filter choosing which elements the tuple holds is computed there too, and a test that is not a compile time value leaves the arity unknown.

What a filter can test therefore depends on what the iteration makes available at compile time:

\begin{itemize}
\item over a range, the value of the iteration is a literal, so it can be tested,
\item over a tuple, both the index and the element are, so both can be tested,
\item over a slice or an array, only the index is: the elements live in memory the comprehension does not read at compile time.
\end{itemize}

A comprehension written between brackets escapes the constraint only when its iteration is not unrolled: it then builds a slice at run time, where the filter is an ordinary test -- see \hyperref[sec:codes:e4302]{E4302}. Over a tuple it builds an array whose length is fixed at compile time, so the filter is computed at compile time there as well, and this error is raised on that form too.

\subsubsection*{Example}

\textit{test\_\allowbreak{}resources/\allowbreak{}diagnostics/\allowbreak{}test147.yr}:

%% check: skip
\begin{lstlisting}[style=coloredverbatimError]
fn main (args: [[c8]]) {
    let n = args.len;
    let _ = (i for i in 0 .. 5 if i != n,);
}
\end{lstlisting}

\begin{lstlisting}[style=bashVerb, breaklines=true, escapechar=~]
Error[E4303] : the filter of the list comprehension, constructed from type (..i32), must be determinable at compile time
 --> test_resources/diagnostics/test147.yr:(3,37)
 3  ~\lstbox{\textSFxi{}}~     let _ = (i for i in 0 .. 5 if i != n,);
    ~\lstbox{\textSFv{}}~                                     ^^  error: value of type bool is needed but unknown at compilation time
    ~\lstbox{\textSFxi{}}~
    = when validating test147::main -- (test_resources/diagnostics/test147.yr:(1,4))
    ~\lstbox{\textSFii{}}~~\lstbox{\textSFx{}}~~\lstbox{\textSFx{}}~~\lstbox{\textSFx{}}~~\lstbox{\textSFx{}}~~\lstbox{\textSFx{}}~~\lstbox{\textSFx{}}~~\lstbox{\textSFx{}}~
\end{lstlisting}

\subsubsection*{How to fix it}

Test something the comprehension knows at compile time, or iterate something whose length is not fixed at compile time, so that the result is a slice built at run time:

%% check: skip
\begin{lstlisting}[style=coloredverbatim]
let _ = copy [i for i in 0 .. 5 if i != n];
\end{lstlisting}

%% Generated by tools/gen_error_codes.py from docs/errors/ of the compiler;
%% edit the compiler's pages or tools/error_themes.txt, then rerun it.

\clearpage
\section{Unit tests}
\label{sec:codes:tests}

\begin{longtable}{@{}l p{0.78\linewidth} r@{}}
\toprule Code & Message & Page\\ \midrule \endhead
\hyperref[sec:codes:e3037]{E3037} & custom attribute \errhole{} is not valid on unit tests & \pageref{sec:codes:e3037}\\
\hyperref[sec:codes:e4286]{E4286} & a unit test taking parameters must produce them with the @parameters attribute & \pageref{sec:codes:e4286}\\
\hyperref[sec:codes:e4287]{E4287} & the @parameters attribute produces values for a unit test that takes no parameter & \pageref{sec:codes:e4287}\\
\hyperref[sec:codes:e4288]{E4288} & the parameters of a unit test are produced as \errhole{}, but a generator was expected & \pageref{sec:codes:e4288}\\
\hyperref[sec:codes:e4289]{E4289} & the parameters of a unit test are produced as \errhole{}, but the test takes \errhole{} parameters & \pageref{sec:codes:e4289}\\
\bottomrule
\end{longtable}

\errorcode{E3037}{custom attribute \errhole{} is not valid on unit tests}
\label{sec:codes:e3037}

Raised during \textbf{declaration}, from \texttt{DeclareErrorMessage::\allowbreak{}UNDEFINED\_\allowbreak{}ATTRIBUTE\_\allowbreak{}FOR\_\allowbreak{}UNITTEST} (\textit{src/\allowbreak{}ymirc/\allowbreak{}semantic/\allowbreak{}declarator/\allowbreak{}errors.yr}).

\subsubsection*{What it means}

A unit test is not a symbol other code can refer to, so the attributes shaping how a function is emitted or overridden mean nothing on it. \tokennolst{@parameters} is the only attribute it accepts.

\subsubsection*{Example}

\textit{test\_\allowbreak{}resources/\allowbreak{}unittest/\allowbreak{}test9.yr}:

%% check: skip
\begin{lstlisting}[style=coloredverbatimError]
@final
__test {
    assert (true);
}
\end{lstlisting}

\begin{lstlisting}[style=bashVerb, breaklines=true, escapechar=~]
Error[E3037] : custom attribute final is not valid on unit tests
 --> test_resources/unittest/test9.yr:(2,2)
 2  ~\lstbox{\textSFxi{}}~ @final
    ~\lstbox{\textSFv{}}~  ^^^^^
    ~\lstbox{\textSFii{}}~~\lstbox{\textSFx{}}~~\lstbox{\textSFx{}}~~\lstbox{\textSFx{}}~~\lstbox{\textSFx{}}~~\lstbox{\textSFx{}}~~\lstbox{\textSFx{}}~~\lstbox{\textSFx{}}~
\end{lstlisting}

\subsubsection*{How to fix it}

Drop the attribute.

\errorcode{E4286}{a unit test taking parameters must produce them with the @parameters attribute}
\label{sec:codes:e4286}

Raised during \textbf{validation}, from \texttt{ValidateErrorMessage::\allowbreak{}UNITTEST\_\allowbreak{}PARAMS\_\allowbreak{}NO\_\allowbreak{}PROVIDER} (\textit{src/\allowbreak{}ymirc/\allowbreak{}semantic/\allowbreak{}validator/\allowbreak{}errors.yr}).

\subsubsection*{What it means}

Nothing calls a unit test but the test runner, so a test taking parameters has no caller to fill them. The values it is run with come from the \tokennolst{@parameters} attribute, and a test declaring parameters without one could never run.

\subsubsection*{Example}

\textit{test\_\allowbreak{}resources/\allowbreak{}unittest/\allowbreak{}test5.yr}:

%% check: skip
\begin{lstlisting}[style=coloredverbatimError]
__test (value: i32) {
    assert (value > 0);
}
\end{lstlisting}

\begin{lstlisting}[style=bashVerb, breaklines=true, escapechar=~]
Error[E4286] : a unit test taking parameters must produce them with the parameters attribute
 --> test_resources/unittest/test5.yr:(2,9)
 2  ~\lstbox{\textSFxi{}}~ __test (value: i32) {
    ~\lstbox{\textSFv{}}~         ^^^^^
    ~\lstbox{\textSFii{}}~~\lstbox{\textSFx{}}~~\lstbox{\textSFx{}}~~\lstbox{\textSFx{}}~~\lstbox{\textSFx{}}~~\lstbox{\textSFx{}}~~\lstbox{\textSFx{}}~~\lstbox{\textSFx{}}~
\end{lstlisting}

\subsubsection*{How to fix it}

Add \tokennolst{@parameters (...)} above the test, giving it the generator of parameter sets -- calling the \tokennolst{yield} function producing them if that is where they come from -- or drop the parameters.

\errorcode{E4287}{the @parameters attribute produces values for a unit test that takes no parameter}
\label{sec:codes:e4287}

Raised during \textbf{validation}, from \texttt{ValidateErrorMessage::\allowbreak{}UNITTEST\_\allowbreak{}PROVIDER\_\allowbreak{}NO\_\allowbreak{}PARAMS} (\textit{src/\allowbreak{}ymirc/\allowbreak{}semantic/\allowbreak{}validator/\allowbreak{}errors.yr}).

\subsubsection*{What it means}

\tokennolst{@parameters} says what a test is run with, one run per parameter set its generator yields. A test taking no parameter has nowhere to receive those values, and would be run repeatedly with the same behavior.

\subsubsection*{Example}

\textit{test\_\allowbreak{}resources/\allowbreak{}unittest/\allowbreak{}test6.yr}:

%% check: skip
\begin{lstlisting}[style=coloredverbatimError]
@parameters ((yield i for i in 1 ... 3))
__test {
    assert (true);
}
\end{lstlisting}

\begin{lstlisting}[style=bashVerb, breaklines=true, escapechar=~]
Error[E4287] : the parameters attribute produces values for a unit test that takes no parameter
 --> test_resources/unittest/test6.yr:(2,2)
 2  ~\lstbox{\textSFxi{}}~ @parameters ((yield i for i in 1 ... 3))
    ~\lstbox{\textSFv{}}~  ^^^^^^^^^^
    ~\lstbox{\textSFii{}}~~\lstbox{\textSFx{}}~~\lstbox{\textSFx{}}~~\lstbox{\textSFx{}}~~\lstbox{\textSFx{}}~~\lstbox{\textSFx{}}~~\lstbox{\textSFx{}}~~\lstbox{\textSFx{}}~
\end{lstlisting}

\subsubsection*{How to fix it}

Declare the parameters the test is run with, or drop the attribute.

\errorcode{E4288}{the parameters of a unit test are produced as \errhole{}, but a generator was expected}
\label{sec:codes:e4288}

Raised during \textbf{validation}, from \texttt{ValidateErrorMessage::\allowbreak{}UNITTEST\_\allowbreak{}PROVIDER\_\allowbreak{}NOT\_\allowbreak{}A\_\allowbreak{}GENERATOR} (\textit{src/\allowbreak{}ymirc/\allowbreak{}semantic/\allowbreak{}validator/\allowbreak{}errors.yr}).

\subsubsection*{What it means}

\tokennolst{@parameters} must produce a generator, of type \tokennolst{yield-\textgreater{} T}: the runner drains it once, before any test runs, and registers one case per yielded parameter set, \tokennolst{\textless{}test\textgreater{}[k]} for the \tokennolst{k}-th one. A slice is not accepted, whatever its elements.

The runner selects a case by its index alone, so the generator must yield the same sets in the same order on every run for \tokennolst{-f "…[k]"} and \tokennolst{--resume} to rerun the same case.

The attribute takes the parameter sets themselves, not a way of obtaining them, so a generator function is called on the spot -- \tokennolst{@parameters (params ())}. Naming it without calling it raises \hyperref[sec:codes:e4084]{E4084} instead.

\subsubsection*{Example}

\textit{test\_\allowbreak{}resources/\allowbreak{}unittest/\allowbreak{}test7.yr}:

%% check: skip
\begin{lstlisting}[style=coloredverbatimError]
@parameters (copy [1, 2, 3])
__test (value: i32) {
    assert (value > 0);
}
\end{lstlisting}

\begin{lstlisting}[style=bashVerb, breaklines=true, escapechar=~]
Error[E4288] : the parameters of a unit test are produced as mut [mut i32], but a generator was expected
 --> test_resources/unittest/test7.yr:(2,14)
 2  ~\lstbox{\textSFxi{}}~ @parameters (copy [1, 2, 3])
    ~\lstbox{\textSFv{}}~              ^^^^
    ~\lstbox{\textSFii{}}~~\lstbox{\textSFx{}}~~\lstbox{\textSFx{}}~~\lstbox{\textSFx{}}~~\lstbox{\textSFx{}}~~\lstbox{\textSFx{}}~~\lstbox{\textSFx{}}~~\lstbox{\textSFx{}}~
\end{lstlisting}

\subsubsection*{How to fix it}

Yield the parameter sets, from a generator comprehension -- \texttt{@parameters ((yield i for i in [1, 2, 3]))}, the comprehension keeping its own parentheses -- or from a call to a \tokennolst{yield} function:

%% check: skip
\begin{lstlisting}[style=coloredverbatim]
yield params ()-> i32 {
    for i in 1 ... 3 {
        yield i;
    }
}

@parameters (params ())
__test (value: i32) {
    assert (value > 0);
}
\end{lstlisting}

\errorcode{E4289}{the parameters of a unit test are produced as \errhole{}, but the test takes \errhole{} parameters}
\label{sec:codes:e4289}

Raised during \textbf{validation}, from \texttt{ValidateErrorMessage::\allowbreak{}UNITTEST\_\allowbreak{}PARAMS\_\allowbreak{}ARITY} (\textit{src/\allowbreak{}ymirc/\allowbreak{}semantic/\allowbreak{}validator/\allowbreak{}errors.yr}).

\subsubsection*{What it means}

Each parameter set the \tokennolst{@parameters} generator yields fills the parameters of one run of the test. A test taking one parameter is fed the element itself; a test taking several is fed a tuple of that arity, unwrapped field by field.

\subsubsection*{Example}

\textit{test\_\allowbreak{}resources/\allowbreak{}unittest/\allowbreak{}test8.yr}:

%% check: skip
\begin{lstlisting}[style=coloredverbatimError]
yield params ()-> (i32, [c8]) {
    yield (3, "foo");
}

@parameters (params ())
__test (len: i32, text: [c8], ratio: f64) {
    assert (len > 0 && text.len > 0us && ratio > 0.0);
}
\end{lstlisting}

\begin{lstlisting}[style=bashVerb, breaklines=true, escapechar=~]
Error[E4289] : the parameters of a unit test are produced as (i32, [c8]), but the test takes 3 parameters
 --> test_resources/unittest/test8.yr:(6,2)
 6  ~\lstbox{\textSFxi{}}~ @parameters (params ())
    ~\lstbox{\textSFv{}}~  ^^^^^^^^^^
 7  ~\lstbox{\textSFxi{}}~ __test (len: i32, text: [c8], ratio: f64) {
    ~\lstbox{\textSFv{}}~         ---
    ~\lstbox{\textSFii{}}~~\lstbox{\textSFx{}}~~\lstbox{\textSFx{}}~~\lstbox{\textSFx{}}~~\lstbox{\textSFx{}}~~\lstbox{\textSFx{}}~~\lstbox{\textSFx{}}~~\lstbox{\textSFx{}}~
\end{lstlisting}

\subsubsection*{How to fix it}

Give the tuple a field per parameter of the test, or take as many parameters as the tuple has fields.

%% Generated by tools/gen_error_codes.py from docs/errors/ of the compiler;
%% edit the compiler's pages or tools/error_themes.txt, then rerun it.

\clearpage
\section{Intermediate language}
\label{sec:codes:yil}

\begin{longtable}{@{}l p{0.78\linewidth} r@{}}
\toprule Code & Message & Page\\ \midrule \endhead
\hyperref[sec:codes:e5001]{E5001} & failed to create YIL output directory \errhole{}, permission denied & \pageref{sec:codes:e5001}\\
\hyperref[sec:codes:e5002]{E5002} & failed to read YIL byte file \errhole{}, file not found or permission denied & \pageref{sec:codes:e5002}\\
\hyperref[sec:codes:e5003]{E5003} & failed to read YIL byte file \errhole{}, file format error & \pageref{sec:codes:e5003}\\
\hyperref[sec:codes:e5004]{E5004} & failed to write YIL byte file \errhole{}, permission denied & \pageref{sec:codes:e5004}\\
\hyperref[sec:codes:e5005]{E5005} & malformed bytecode & \pageref{sec:codes:e5005}\\
\hyperref[sec:codes:e5006]{E5006} & malformed bytecode, type table invalid & \pageref{sec:codes:e5006}\\
\hyperref[sec:codes:e5007]{E5007} & malformed bytecode, string table invalid & \pageref{sec:codes:e5007}\\
\hyperref[sec:codes:e5008]{E5008} & malformed bytecode, symbol table invalid & \pageref{sec:codes:e5008}\\
\hyperref[sec:codes:e5009]{E5009} & malformed bytecode, location table invalid & \pageref{sec:codes:e5009}\\
\hyperref[sec:codes:e5010]{E5010} & YIL byte file was created for a \errhole{} bits target, mismatch current target arch \errhole{} bits & \pageref{sec:codes:e5010}\\
\bottomrule
\end{longtable}

\errorcode{E5001}{failed to create YIL output directory \errhole{}, permission denied}
\label{sec:codes:e5001}

Raised during \textbf{lowering (YIL)}, from \texttt{SerializeYILErrorMessage::\allowbreak{}FAILED\_\allowbreak{}TO\_\allowbreak{}CREATE\_\allowbreak{}OUT\_\allowbreak{}DIR} (\textit{src/\allowbreak{}ymirc/\allowbreak{}lint/\allowbreak{}serialize/\allowbreak{}errors.yr}).

\subsubsection*{What it means}

The compiler could not create the directory it was told to write the YIL bytecode of the compiled package into. The path is reported in the message.

\subsubsection*{How to fix it}

Check that the parent directory exists and is writable, and that the path given to the YIL output option is the one intended.

\errorcode{E5002}{failed to read YIL byte file \errhole{}, file not found or permission denied}
\label{sec:codes:e5002}

Raised during \textbf{lowering (YIL)}, from \texttt{SerializeYILErrorMessage::\allowbreak{}FAILED\_\allowbreak{}TO\_\allowbreak{}READ\_\allowbreak{}BYTE\_\allowbreak{}FILE} (\textit{src/\allowbreak{}ymirc/\allowbreak{}lint/\allowbreak{}serialize/\allowbreak{}errors.yr}).

\subsubsection*{What it means}

The compiler could not read a \tokennolst{.yil} file it needs. The file is part of an already compiled package the current one imports; it is either absent from the include path or not readable.

\subsubsection*{How to fix it}

Check that the package providing the module has been compiled, and that its output directory is on the include path. A permission error on an existing file points at the file mode rather than at the path.

\errorcode{E5003}{failed to read YIL byte file \errhole{}, file format error}
\label{sec:codes:e5003}

Raised during \textbf{lowering (YIL)}, from \texttt{SerializeYILErrorMessage::\allowbreak{}FAILED\_\allowbreak{}TO\_\allowbreak{}PARSE\_\allowbreak{}BYTE\_\allowbreak{}FILE} (\textit{src/\allowbreak{}ymirc/\allowbreak{}lint/\allowbreak{}serialize/\allowbreak{}errors.yr}).

\subsubsection*{What it means}

A \tokennolst{.yil} file was read but its content is not valid YIL bytecode. Most often the file was produced by a different version of the compiler, or it was truncated by an interrupted build.

\subsubsection*{How to fix it}

Recompile the package that produced the file with the current compiler. A stale build directory left over from an earlier version is the usual cause.

\errorcode{E5004}{failed to write YIL byte file \errhole{}, permission denied}
\label{sec:codes:e5004}

Raised during \textbf{lowering (YIL)}, from \texttt{SerializeYILErrorMessage::\allowbreak{}FAILED\_\allowbreak{}TO\_\allowbreak{}WRITE\_\allowbreak{}BYTE\_\allowbreak{}FILE} (\textit{src/\allowbreak{}ymirc/\allowbreak{}lint/\allowbreak{}serialize/\allowbreak{}errors.yr}).

\subsubsection*{What it means}

The compiler could not write the \tokennolst{.yil} file of a compiled module. The path is reported in the message.

\subsubsection*{How to fix it}

Check the permissions of the output directory, and that the disk is not full.

\errorcode{E5005}{malformed bytecode}
\label{sec:codes:e5005}

Raised during \textbf{lowering (YIL)}, from \texttt{SerializeYILErrorMessage::\allowbreak{}MALFORMED\_\allowbreak{}BYTECODE} (\textit{src/\allowbreak{}ymirc/\allowbreak{}lint/\allowbreak{}serialize/\allowbreak{}errors.yr}).

\subsubsection*{What it means}

A \tokennolst{.yil} file is structurally invalid: the reader found a byte sequence that cannot be decoded at the position it had reached. As with \hyperref[sec:codes:e5003]{E5003}, this points at a file produced by another compiler version or at a truncated one.

\subsubsection*{How to fix it}

Recompile the package that produced the file. If it was produced by the current compiler and the error persists, the file is a compiler bug worth reporting with the code and the file.

\errorcode{E5006}{malformed bytecode, type table invalid}
\label{sec:codes:e5006}

Raised during \textbf{lowering (YIL)}, from \texttt{SerializeYILErrorMessage::\allowbreak{}MALFORMED\_\allowbreak{}BYTECODE\_\allowbreak{}TYPE\_\allowbreak{}TABLE} (\textit{src/\allowbreak{}ymirc/\allowbreak{}lint/\allowbreak{}serialize/\allowbreak{}errors.yr}).

\subsubsection*{What it means}

The type table of a \tokennolst{.yil} file could not be decoded. The bytecode carries the types of a module in a table the rest of the file refers to by index, and that table is invalid.

\subsubsection*{How to fix it}

Recompile the package that produced the file, cf. \hyperref[sec:codes:e5005]{E5005}.

\errorcode{E5007}{malformed bytecode, string table invalid}
\label{sec:codes:e5007}

Raised during \textbf{lowering (YIL)}, from \texttt{SerializeYILErrorMessage::\allowbreak{}MALFORMED\_\allowbreak{}BYTECODE\_\allowbreak{}STRING\_\allowbreak{}TABLE} (\textit{src/\allowbreak{}ymirc/\allowbreak{}lint/\allowbreak{}serialize/\allowbreak{}errors.yr}).

\subsubsection*{What it means}

The string table of a \tokennolst{.yil} file could not be decoded. Symbol names and literals are stored in a table the rest of the file refers to by index, and that table is invalid.

\subsubsection*{How to fix it}

Recompile the package that produced the file, cf. \hyperref[sec:codes:e5005]{E5005}.

\errorcode{E5008}{malformed bytecode, symbol table invalid}
\label{sec:codes:e5008}

Raised during \textbf{lowering (YIL)}, from \texttt{SerializeYILErrorMessage::\allowbreak{}MALFORMED\_\allowbreak{}BYTECODE\_\allowbreak{}SYMBOL\_\allowbreak{}TABLE} (\textit{src/\allowbreak{}ymirc/\allowbreak{}lint/\allowbreak{}serialize/\allowbreak{}errors.yr}).

\subsubsection*{What it means}

The symbol table of a \tokennolst{.yil} file could not be decoded: the frames, globals and constants the module exports are listed in a table that is invalid.

\subsubsection*{How to fix it}

Recompile the package that produced the file, cf. \hyperref[sec:codes:e5005]{E5005}.

\errorcode{E5009}{malformed bytecode, location table invalid}
\label{sec:codes:e5009}

Raised during \textbf{lowering (YIL)}, from \texttt{SerializeYILErrorMessage::\allowbreak{}MALFORMED\_\allowbreak{}BYTECODE\_\allowbreak{}LOCATION\_\allowbreak{}TABLE} (\textit{src/\allowbreak{}ymirc/\allowbreak{}lint/\allowbreak{}serialize/\allowbreak{}errors.yr}).

\subsubsection*{What it means}

The location table of a \tokennolst{.yil} file could not be decoded. Source locations are stored in a table so a diagnostic raised on an imported symbol can still point at its source, and that table is invalid.

\subsubsection*{How to fix it}

Recompile the package that produced the file, cf. \hyperref[sec:codes:e5005]{E5005}.

\errorcode{E5010}{YIL byte file was created for a \errhole{} bits target, mismatch current target arch \errhole{} bits}
\label{sec:codes:e5010}

Raised during \textbf{lowering (YIL)}, from \texttt{SerializeYILErrorMessage::\allowbreak{}MISMATCH\_\allowbreak{}ARCH\_\allowbreak{}POINTER\_\allowbreak{}SIZE} (\textit{src/\allowbreak{}ymirc/\allowbreak{}lint/\allowbreak{}serialize/\allowbreak{}errors.yr}).

\subsubsection*{What it means}

A \tokennolst{.yil} file was produced for a target whose pointer width is not the one being compiled for -- a 32 bit bytecode read by a 64 bit compilation, or the reverse. Pointer width changes the layout of every type in the file, so the bytecode cannot be reused across targets.

\subsubsection*{How to fix it}

Recompile the imported package for the target currently being built. A build directory shared between two targets has to be separated, one output directory per target.

%% Generated by tools/gen_error_codes.py from docs/errors/ of the compiler;
%% edit the compiler's pages or tools/error_themes.txt, then rerun it.

\section{Retired codes}
\label{sec:codes:retired}

A retired code named a message that was deleted from the compiler. It is
never reused, so an older compiler printing it still means the message
below.

\begin{longtable}{@{}l p{0.4\linewidth} p{0.45\linewidth}@{}}
\toprule Code & Catalogue entry & Message\\ \midrule \endhead
E1003 & \texttt{LexingErrorMessage::\allowbreak{}RETURN\_\allowbreak{}IN\_\allowbreak{}SINGLE\_\allowbreak{}LINE\_\allowbreak{}COMMENT} & line break encountered inside a single-line comment\\
E2003 & \texttt{SyntaxErrorMessage::\allowbreak{}MULTIPLE\_\allowbreak{}AUX\_\allowbreak{}CSTRS} & multiple constructor redirects are not permitted\\
E2005 & \texttt{SyntaxErrorMessage::\allowbreak{}UNDEFINED\_\allowbreak{}ATTRIBUTE} & undefined attribute \errhole{}\\
E2010 & \texttt{SyntaxErrorMessage::\allowbreak{}UNTERMINATED\_\allowbreak{}BLOCK} & unterminated block\\
E2011 & \texttt{SyntaxErrorMessage::\allowbreak{}USED\_\allowbreak{}AS\_\allowbreak{}IDENTIFIER} & \errhole{} cannot be used as an identifier\\
E3005 & \texttt{DeclareErrorMessage::\allowbreak{}DECL\_\allowbreak{}VARIADIC\_\allowbreak{}FUNC} & only external C functions may be variadic\\
E3008 & \texttt{DeclareErrorMessage::\allowbreak{}DTOR\_\allowbreak{}IN\_\allowbreak{}RECORD} & record types cannot define destructors\\
E3011 & \texttt{DeclareErrorMessage::\allowbreak{}NOT\_\allowbreak{}OVERRIDE} & method \errhole{} is marked override but does not override any base method\\
E3028 & \texttt{DeclareErrorMessage::\allowbreak{}UNEXPECTED\_\allowbreak{}IN\_\allowbreak{}CLASS} & unexpected declaration inside a class\\
E3031 & \texttt{DeclareErrorMessage::\allowbreak{}UNEXPECTED\_\allowbreak{}IN\_\allowbreak{}MACRO} & unexpected declaration inside a macro\\
E3032 & \texttt{DeclareErrorMessage::\allowbreak{}UNEXPECTED\_\allowbreak{}IN\_\allowbreak{}MODULE} & unexpected declaration\\
E4004 & \texttt{ValidateErrorMessage::\allowbreak{}ALREADY\_\allowbreak{}INIT\_\allowbreak{}BY\_\allowbreak{}CTOR\_\allowbreak{}REDIRECT} & field \errhole{} was already initialized by ctor redirection\\
E4022 & \texttt{ValidateErrorMessage::\allowbreak{}CLOSURE\_\allowbreak{}TEMPLATE\_\allowbreak{}LAMBDA} & closure functions cannot be used as template values\\
E4025 & \texttt{ValidateErrorMessage::\allowbreak{}CONFLIT\_\allowbreak{}DECORATORS} & conflicting decorators \errhole{} and \errhole{}\\
E4032 & \texttt{ValidateErrorMessage::\allowbreak{}CTE\_\allowbreak{}IGNORED} & \errhole{} values contained in a block scope are ignored during compile time execution\\
E4033 & \texttt{ValidateErrorMessage::\allowbreak{}CTOR\_\allowbreak{}CLASS\_\allowbreak{}DCOPY} & using a deep copy instead of a one level copy is prohibited when constructing class \errhole{}\\
E4034 & \texttt{ValidateErrorMessage::\allowbreak{}CTOR\_\allowbreak{}CLASS\_\allowbreak{}STACK} & creating an instance of class \errhole{} on the stack is prohibited\\
E4048 & \texttt{ValidateErrorMessage::\allowbreak{}DTOR\_\allowbreak{}SELF\_\allowbreak{}NOT\_\allowbreak{}MUT} & the self parameter must be mutable in the destructor\\
E4052 & \texttt{ValidateErrorMessage::\allowbreak{}ENCLOSE\_\allowbreak{}INCOMPLETE\_\allowbreak{}TYPE} & the type \errhole{} cannot be enclosed because it is not complete\\
E4053 & \texttt{ValidateErrorMessage::\allowbreak{}ENCLOSE\_\allowbreak{}MOVABLE\_\allowbreak{}TYPE} & movable type \errhole{} cannot be enclosed within a closure function\\
E4057 & \texttt{ValidateErrorMessage::\allowbreak{}EXTERN\_\allowbreak{}GLOBAL\_\allowbreak{}NO\_\allowbreak{}TYPE} & external global variable must be declared with a type\\
E4061 & \texttt{ValidateErrorMessage::\allowbreak{}FIELD\_\allowbreak{}METHOD\_\allowbreak{}ALIAS} & method \errhole{} is a field method\\
E4064 & \texttt{ValidateErrorMessage::\allowbreak{}FIELD\_\allowbreak{}METHOD\_\allowbreak{}COLLISION} & field method \errhole{} cannot be overlapped by a method with the same name \errhole{}\\
E4070 & \texttt{ValidateErrorMessage::\allowbreak{}FORWARD\_\allowbreak{}REFERENCE\_\allowbreak{}VALUE} & the value cannot be computed, as it depends on a forward reference\\
E4071 & \texttt{ValidateErrorMessage::\allowbreak{}FORWARD\_\allowbreak{}REFERENCE\_\allowbreak{}VAR} & the type cannot be inferred, as it depends on a forward reference\\
E4086 & \texttt{ValidateErrorMessage::\allowbreak{}IMPLICIT\_\allowbreak{}CONST\_\allowbreak{}REFERENCE} & ref of type \errhole{} must be const, but is mutable\\
E4098 & \texttt{ValidateErrorMessage::\allowbreak{}INCOMPLETE\_\allowbreak{}TYPE\_\allowbreak{}CLASS} & the type \errhole{} is not complete due to previous errors\\
E4120 & \texttt{ValidateErrorMessage::\allowbreak{}MAP\_\allowbreak{}KEY\_\allowbreak{}MUTABLE} & the type \errhole{} that is used as a map key must always be immutable\\
E4132 & \texttt{ValidateErrorMessage::\allowbreak{}MULTIPLE\_\allowbreak{}DECORATORS} & decorator \errhole{} is specified multiple times\\
E4139 & \texttt{ValidateErrorMessage::\allowbreak{}MUTABLE\_\allowbreak{}LAMBDA\_\allowbreak{}VAR} & a lambda variable cannot be mutable\\
E4148 & \texttt{ValidateErrorMessage::\allowbreak{}NOT\_\allowbreak{}AN\_\allowbreak{}ARRAY} & \errhole{} is not an array type\\
E4152 & \texttt{ValidateErrorMessage::\allowbreak{}NOT\_\allowbreak{}A\_\allowbreak{}LVALUE\_\allowbreak{}CLOSURE} & \errhole{} is enclosed, and cannot be used as a lvalue\\
E4158 & \texttt{ValidateErrorMessage::\allowbreak{}NOT\_\allowbreak{}A\_\allowbreak{}MAP} & \errhole{} is not a map type\\
E4159 & \texttt{ValidateErrorMessage::\allowbreak{}NOT\_\allowbreak{}A\_\allowbreak{}RECORD} & \errhole{} is not a record type\\
E4160 & \texttt{ValidateErrorMessage::\allowbreak{}NOT\_\allowbreak{}A\_\allowbreak{}SLICE} & \errhole{} is not a slice type\\
E4161 & \texttt{ValidateErrorMessage::\allowbreak{}NOT\_\allowbreak{}AN\_\allowbreak{}ENTITY} & \errhole{} is not an entity type\\
E4164 & \texttt{ValidateErrorMessage::\allowbreak{}NOT\_\allowbreak{}AN\_\allowbreak{}UNION} & \errhole{} is not a union type\\
E4170 & \texttt{ValidateErrorMessage::\allowbreak{}NO\_\allowbreak{}DEFAULT\_\allowbreak{}CTOR\_\allowbreak{}MOVE\_\allowbreak{}STRUCT} & type \errhole{} is an entity, but has no default ctor self()\\
E4174 & \texttt{ValidateErrorMessage::\allowbreak{}ONE\_\allowbreak{}ITER\_\allowbreak{}LOOP} & do while loop test is always false, the loop is always entered exactly once, so the branching construct is useless\\
E4191 & \texttt{ValidateErrorMessage::\allowbreak{}OVERRIDE\_\allowbreak{}PRIVATE} & cannot override private method \errhole{}\\
E4194 & \texttt{ValidateErrorMessage::\allowbreak{}PRAGMA\_\allowbreak{}FIELD\_\allowbreak{}NO\_\allowbreak{}DEFAULT} & field \errhole{} from type \errhole{} has no default value\\
E4200 & \texttt{ValidateErrorMessage::\allowbreak{}SUPER\_\allowbreak{}NO\_\allowbreak{}SELF\_\allowbreak{}CLASS} & the superclass proxy of \errhole{} is only accessible through self\\
E4201 & \texttt{ValidateErrorMessage::\allowbreak{}TEMPLATE\_\allowbreak{}REST} & template validation is not complete, unresolved: \{\errhole{}\}\\
E4204 & \texttt{ValidateErrorMessage::\allowbreak{}TEMPLATE\_\allowbreak{}TEST\_\allowbreak{}FAILED} & the test of the template specialization failed\\
E4205 & \texttt{ValidateErrorMessage::\allowbreak{}THROWS\_\allowbreak{}IN\_\allowbreak{}LAMBDA} & a lambda function must be safe, but there are exceptions that are not caught\\
E4215 & \texttt{ValidateErrorMessage::\allowbreak{}TUPLE\_\allowbreak{}MATCHER} & tuple matcher only applies to tuple values not \errhole{}\\
E4219 & \texttt{ValidateErrorMessage::\allowbreak{}TYPE\_\allowbreak{}HAS\_\allowbreak{}NO\_\allowbreak{}TYPEINFO} & temporary type \errhole{} has no type information\\
E4221 & \texttt{ValidateErrorMessage::\allowbreak{}UNDEFINED\_\allowbreak{}ATTRIBUTE} & attribute \errhole{} is not usable in this context\\
E4225 & \texttt{ValidateErrorMessage::\allowbreak{}UNDEFINED\_\allowbreak{}BIN\_\allowbreak{}OP\_\allowbreak{}TOK} & undefined binary operator \errhole{}\\
E4228 & \texttt{ValidateErrorMessage::\allowbreak{}UNDEFINED\_\allowbreak{}CTOR\_\allowbreak{}MACRO} & macro \errhole{} has no accessible constructors in that context\\
E4232 & \texttt{ValidateErrorMessage::\allowbreak{}UNDEFINED\_\allowbreak{}MACRO\_\allowbreak{}CALL} & macro call is undefined for \errhole{}\\
E4235 & \texttt{ValidateErrorMessage::\allowbreak{}UNDEFINED\_\allowbreak{}SUPER\_\allowbreak{}CALL\_\allowbreak{}CTOR} & no constructor of the super class is callable with the parameters \errhole{}\\
E4238 & \texttt{ValidateErrorMessage::\allowbreak{}UNDEFINED\_\allowbreak{}UN\_\allowbreak{}OP\_\allowbreak{}TOK} & undefined unary operator \errhole{}\\
E4241 & \texttt{ValidateErrorMessage::\allowbreak{}UNDEF\_\allowbreak{}DECORATOR\_\allowbreak{}TEMPLATE} & decorator \errhole{} is not applicable in template specialization\\
E4242 & \texttt{ValidateErrorMessage::\allowbreak{}UNDEF\_\allowbreak{}DECORATOR\_\allowbreak{}TYPE} & decorator \errhole{} is not applicable for types\\
E4243 & \texttt{ValidateErrorMessage::\allowbreak{}UNDEF\_\allowbreak{}DECORATOR\_\allowbreak{}VALUE} & decorator \errhole{} is not applicable for values\\
E4245 & \texttt{ValidateErrorMessage::\allowbreak{}UNDEF\_\allowbreak{}VAR} & undefined symbol \errhole{}\\
E4254 & \texttt{ValidateErrorMessage::\allowbreak{}UNION\_\allowbreak{}CTOR\_\allowbreak{}MULTIPLE\_\allowbreak{}FIELDS} & constructor of the union type \errhole{} initializes multiple fields \errhole{}\\
E4255 & \texttt{ValidateErrorMessage::\allowbreak{}UNION\_\allowbreak{}CTOR\_\allowbreak{}NO\_\allowbreak{}FIELD} & constructor of the union type \errhole{} initializes no field\\
E4274 & \texttt{ValidateErrorMessage::\allowbreak{}USE\_\allowbreak{}AS\_\allowbreak{}VALUE\_\allowbreak{}TEMPLATE} & template specialization expected a value not the type \errhole{}\\
E4276 & \texttt{ValidateErrorMessage::\allowbreak{}VAR\_\allowbreak{}DECL\_\allowbreak{}WITHOUT\_\allowbreak{}VALUE} & var declared without value, when necessary\\
E4277 & \texttt{ValidateErrorMessage::\allowbreak{}VAR\_\allowbreak{}DECL\_\allowbreak{}WITH\_\allowbreak{}NOTHING} & var declaration must at least have a type or a value\\
E4279 & \texttt{ValidateErrorMessage::\allowbreak{}VOID\_\allowbreak{}VAR} & cannot create a variable of type void\\
\bottomrule
\end{longtable}


